% ============================================================
%  R. Nielsen, Religionsphilosophie (Kjøbenhavn: Gyldendal, 1869)
%  TRANSLATION (English) --- Hans Halvorson
%  Translated from ./transcription.tex (Danish; complete, collated),
%  printed pp. 1--537. Method: ../../../TRANSLATION-PLAYBOOK.md.
%
%  Mirrors transcription.tex 1:1: every "% ---- printed p.N" comment and
%  every \opage{N} margin marker is carried at the corresponding English
%  words; all structure macros, footnotes and \emph (letterspacing) match
%  the Danish in number and order.
%
%  ---- FIXED TERMINOLOGY (hold constant) ----
%    Tro = faith; Viden = knowledge; Videnskab = science (Wissenschaft)
%    Erkjendelse = cognition; Troeserkjendelse = the cognition of faith
%    Troesprincipet = the principle of faith
%    ueensartet = heterogeneous (THE thesis word); eensartet = homogeneous
%    omvendt Videnskab = inverted science; Mysteriet = the Mystery
%    Aabenbaring = revelation; Tilegnelse = appropriation
%    Forestilling = representation; Begreb = concept; Idee = Idea
%    Bestemmelse = determination; Væsen = essence; Selvet = the self/Self
%    Personlighed = personality; Aand = spirit / the Spirit
%    Skabelse / Opholdelse / Styrelse = creation / preservation / governance
%    Forsyn = providence; Forsoning = atonement (reconciliation in Hegelian
%    self-reconciliation); Forløsning = redemption; Salighed = blessedness
%    Under / Mirakel = miracle; Kritik = critique; Selvmagt = self-power
%    Virkelighed = actuality; Forargelse = offense; Inderlighed = inwardness
%    Troesbekjendelsen = the Confession of Faith; Led = article
%
%  - Quotes „...“ -> ``...''. Latin/German kept in \textit as printed.
%    Greek copied verbatim (needs textalpha).
%  - Bible references anglicized in form (Joh. 3, 16 -> John 3:16); all
%    other bibliographic references kept exactly as Nielsen gives them.
%  - Printer's defects: the Danish transcribes them as printed; the English
%    renders the evident sense and logs each with a "% print:" comment at
%    the site. Apparent slips of the AUTHOR (e.g. wrong verse numbers) are
%    kept as printed, with a "% print: ... probably meant" comment.
% ============================================================
\documentclass[12pt]{book}
\usepackage[utf8]{inputenc}
\usepackage[T1]{fontenc}
\usepackage[english]{babel}
\usepackage[protrusion=true,expansion=true]{microtype}
\usepackage[margin=1.4in]{geometry}
\usepackage{setspace}
\usepackage{amsmath}
\usepackage{libertinus}
\usepackage{libertinust1math}
\usepackage{textalpha}   % polytonic Greek typed directly
\usepackage{enumitem}
\usepackage{fancyhdr}
\usepackage[colorlinks=true,linkcolor=black,urlcolor=black,
  pdftitle={Philosophy of Religion},
  pdfauthor={Rasmus Nielsen}]{hyperref}
\setlength{\headheight}{15pt}
\pagestyle{fancy}
\fancyhf{}
\fancyhead[LE,RO]{\thepage}
\fancyhead[RE]{\textit{Philosophy of Religion}}
\fancyhead[LO]{\textit{\leftmark}}
\setlength{\parindent}{1.5em}
\setlength{\parskip}{0pt}
\setstretch{1.3}

\newcommand{\parthead}[1]{%
  \clearpage
  \begin{center}{\LARGE\bfseries #1}\end{center}\vspace{2pt}%
  \phantomsection\addcontentsline{toc}{section}{#1}\markboth{#1}{#1}}
\newcommand{\lettersub}[2]{%
  \begin{center}{\large\bfseries #1}\par\vspace{3pt}{\large\bfseries #2}\end{center}%
  \phantomsection\addcontentsline{toc}{subsection}{#1\quad #2}\markboth{#2}{#2}}
\newcommand{\subhead}[1]{%
  \begin{center}{\large\bfseries #1}\end{center}%
  \phantomsection\addcontentsline{toc}{subsection}{#1}\markboth{#1}{#1}}
\newcommand{\parmark}[1]{\begin{center}{\bfseries \S\,#1.}\end{center}}
\newcommand{\runhead}[1]{%
  \medskip\begin{center}{\bfseries #1}\end{center}\medskip}
\newcommand{\greekrun}[2]{%
  \medskip\begin{center}{\itshape #1)}\ {\bfseries #2}\end{center}\medskip}

\usepackage{marginnote}
\newcommand{\opage}[1]{\marginnote{\footnotesize\textit{[#1]}}}
\newcommand{\apage}[1]{\marginnote{\footnotesize\textit{[#1]}}}
\begin{document}

% ============================================================
% TITLE PAGE — PDF 8
% ============================================================
\begin{titlepage}
\centering
\vspace*{3.0cm}
{\Huge\bfseries Philosophy of Religion\par}
\vspace{1.4cm}
{\large by\par}
\vspace{1.0cm}
{\LARGE\bfseries R.\ Nielsen.\par}
\vspace{3.0cm}
\rule{4cm}{0.4pt}\par
\vspace{2.0cm}
{\large\bfseries Copenhagen.\par}
\vspace{0.4cm}
{\normalsize Published by the Gyldendal Bookshop (F.\ Hegel).\par}
\vspace{0.3cm}
{\small\emph{I.\ Cohen's Printing House.}\par}
\vspace{0.6cm}
{\large 1869.\par}
\vspace{1.5cm}
{\small Translated from the Danish by Hans Halvorson\par}
\end{titlepage}

\frontmatter

% ============================================================
% PREFACE THESES — PDF 10 (unnumbered leaf)
% ============================================================
\thispagestyle{empty}
\vspace*{0.10\textheight}
\begin{quote}
\itshape\noindent
Faith and knowledge are as principles absolutely heterogeneous.

\medskip

The heterogeneity of the principles makes no breach in the universally
valid laws of thought, but shows the impossibility of any scientific
transition whatever, either from knowledge to faith or from faith to
knowledge.

\medskip

The long conflict between religion and science has its original ground
in an obscure conflation of the heterogeneous principles, and must
therefore cease when the conflation ceases.

\medskip

The proposition that truth in religion is not truth in science is no
longer a paradox, once the absolute boundary of knowledge coincides with
the mystery from which the cognition of faith proceeds and to which it
returns.

\medskip

It is these principal theses that the Philosophy of Religion in its
present form will develop and substantiate. Whether the development and
the substantiation have any validity, examination must decide.
\end{quote}

\vspace{1.2cm}
\begin{flushright}
Copenhagen, 28 December 1868.\par
\vspace{0.4cm}
R.\ Nielsen.
\end{flushright}

\tableofcontents
\mainmatter

% ============================================================
% BODY BEGINS — printed p.1 (PDF 14)
% ============================================================

% ---- printed p.1 (PDF 14) ----
\opage{1} \begin{center}
  {\Large\bfseries INTRODUCTION.}\par\vspace{4pt}
  \rule{2.5cm}{0.4pt}
\end{center}
\phantomsection
\addcontentsline{toc}{section}{Introduction}
\markboth{Introduction}{Introduction}

The relation of the philosophy of religion to religion cannot be
stated satisfactorily in general conceptual determinations; its
peculiar nature demands a peculiar treatment. What is said
religiously: Show me your faith by your works, must be said
philosophically: Show me your principle by its execution.

\subhead{Religion and Philosophy.}
\parmark{1}

If philosophy as a whole is a science, then the philosophy of
religion, insofar as it is to deserve the name of philosophy, must
naturally also be a science; but precisely the distinctive feature,
that it is a science of religion, gives it a character diverging
from the other sciences and stamps it with a mark of its own. What
this means will best be understood by a glance at the principal
forms the philosophy of religion has assumed in recent times.

The scientific form must be determined by the principle. The
fundamental question is: do religion and philosophy have the same
principle, or are the principles different? If the principle of
religion is different from that of philosophy, which principle is
science then to follow, that of philosophy or that of religion? Only
by an answer to
% ---- printed p.2 (PDF 15) ----
\opage{2} these questions can the concept and method of the philosophy of
religion be determined.

If one proceeds from the presupposition that religion and philosophy
have the same principle, the philosophy of religion becomes a
\emph{speculative science}; if one assumes that the principles are
indeed different, but that truth must nonetheless be judged by the
principle of knowledge, the philosophy of religion becomes a
\emph{critical science}. If, finally, one takes the principles to be
absolutely heterogeneous and lays the religious principle at the
foundation of the scientific elucidation, the philosophy of religion
becomes an \emph{inverted science}. The first of these points of
view is represented by Hegel; the second chiefly by L.\ Feuerbach;
the third is the one brought to application in the present
exposition.

\runhead{a) The Philosophy of Religion as Speculative Science.}

The object of religion as well as of philosophy is, according to
Hegel, the eternal truth in its objectivity, God and nothing but God
and the explication of God. Philosophy only develops itself in
developing religion, and in developing religion it develops itself.
Thus religion and philosophy then coincide into one. Philosophy is
in actuality worship, is religion, for it is the renunciation of
subjective fancies and opinions in the occupation with God.
Nonetheless the occupation with God in philosophy is a different
one, takes place in a different way and manner than in religion. In
religion God is the immediate object of consciousness; in philosophy
the consciousness of God is God's own consciousness of himself: the
difference lies in the form. Religion as religion is popular, is for
all, and for that very reason is referred to the figurative form of
representation; but such a form is inadequate to the content. The
form is finite, the content infinite: this is a disproportion. Can
this disproportion be removed; can the same content, without
% ---- printed p.3 (PDF 16) ----
\opage{3} changing, be presented in two entirely different forms? Here is the
difficulty, and on this difficulty Hegel founders. For where it is a
matter of the truth as truth, the Idea is so absorbed into the form
that the form cannot possibly be changed without the content being
changed. In the course of its movement of thought the speculative
philosophy of religion is imperceptibly led to a different result
from the one it originally aimed at. The movement is this: it begins
with science being to acknowledge and ground the truth of religion;
the truth of religion is determined by what is true in the content,
i.e.\ by the greater or lesser quantum of vital ideality that a given
religion, relative to the stage of formal development at which it
stands, is able to express; since now the form of religion is not
scientific, but the scientific form is the only perfect one, it is
brought to consciousness that even the perfect religion has an
imperfect form, and on account of the imperfection of the form can
express what is true only in an imperfect manner. That by which the
perfect religion is distinguished from philosophy is therefore a
form that must be sublated; but in the sublation of the form, what
is peculiar to religion is sublated, and therewith religion itself.
In losing its form religion at the same time loses its content, for
the truth of the content rests on the form. So it ends with this:
that philosophy, in order to be able to comprehend what is true in
religion, must comprehend religion itself as an untruth. The
negative relation between philosophy and religion, repellent from
the side of science, has come to expression with vigorous
consistency in Neo-Hegelianism.

\runhead{b) The Philosophy of Religion as Critical Science.}

While the speculatively positive tendency holds principally to the
content common to religion and philosophy and strives to develop
this true content of truth
% ---- printed p.4 (PDF 17) ----
\opage{4} into the true form, the negatively critical tendency, in
cognition of the difference in principle between religion and
philosophy, proceeds principally from the form peculiar to religion,
demonstrates its manifest conflict with reason, and so arrives,
through a demonstration of what is untrue in the form, at what is
untrue in the content. Immediately it looks as though religion had
to do with nothing but actual, if supernatural, facts; but reason
developed into knowledge comprehends that the supernatural facts are
not actual, and the actual ones not supernatural. It is not spirit
as spirit, but the ideal natural forces of the life of the soul,
that originally stir in the originally religious consciousness.
Religious representations are indeed not products of conscious
arbitrariness; on the contrary, they are owed to an unconscious
poetizing, a need of souls for self-objectification, a psychological
necessity by which imagination, overwhelmed by the Idea, has placed
the divine in man before man by placing it outside man. Religion is
human; it is the first elementary breakthrough of human reason; it
has its origin, not in God's love, but in man's passions. The dreams
of the heart, the stirrings of the mind, the longings of the soul
contain the basic elements out of which the Idea-inspired
imagination has formed all the gods of faith, the one Almighty as
well. This is confirmed by insight into the essence of monotheism as
well as of polytheism. The mighty, wrathful, vengeance-snorting
Jehovah of the Old Testament is, scientifically regarded, an
idealized reflection of the abstract, mighty spirit of the Jewish
people, of Israel's strong, law-bound and yet self-willed ego.
Christ in the New Testament is not an actual God-man, for an actual
God-man is a contradiction, but is nonetheless essentially a
man-god, the reconciling God of the wounded heart, of the anguished
soul, a God for the suffering. For imagination and feeling the ideal
presents itself immediately;
% ---- printed p.5 (PDF 18) ----
\opage{5} on this rests not only the wish itself but also its fulfillment.
Man, with his confining natural limitation, wishes for power; if he
had power as he has will, he would feel blessed and free; but alas,
unfortunately he is powerless; so in his powerlessness he feels the
need of divine assistance. The divine assistance is immediately
determined thus: that a God who has power over nature arbitrarily
intervenes in the course of things and helps by a miracle. In the
miracle is intuited the unlimited will, freed from the compulsion of
nature, a will that is able to annihilate all resistance, satisfy
all wishes and hear all prayers. The miracle is the ground of faith,
the way out for hope, the consolation of imagination, the holy fairy
tale of the pious disposition; when a God has grown so old and weak
that he can no longer work miracles, he does best to resign, to
cease to be God, and likewise a religion that cannot continually
show miracles must cease to be religion. Here is the turning point
where the conflict between religion and philosophy culminates.
Religion holds to the miracle, philosophy to reason. The miracle is
in contradiction with reason; when reason awakens to the cognition
of the laws of actuality, the miracle disappears, and with the
miracle religion as well. Thus philosophy, as negative critique, has
then, by its own assurance, accomplished the dissolution of
religion.

By the negative path, however, a positive result has been reached.
The result is that religion, rejected and scorned by ``science'', is
referred to its own principle and thereby called upon to rely on its
own strength and assert its independence over against science. The
decisive break between religion and science can then, on these
terms, benefit religion just as well as science. The absolute
heterogeneity of religion with all objective knowledge and science
has been demonstrated with pene-
% ---- printed p.6 (PDF 19) ----
\opage{6} trating sharpness and victorious superiority by S.\ Kierkegaard.

The fundamental question from which S.\ Kierkegaard proceeds is
this: \emph{To what extent can the truth be learned?} Here one then
meets what Socrates in the Meno calls a contentious proposition: that
a man cannot possibly seek what he knows, and just as impossibly can
seek what he does not know; for what he knows he cannot seek, since
he knows it, and what he does not know he cannot seek, for he does
not even know what it is he is to seek. Socrates removes the
difficulty by holding that all learning and seeking is only a
recollecting, so that the ignorant one merely needs to be reminded
in order, by himself, to call to mind what he knows. The truth is
thus not brought into him, but was in him. This is what we may call
the human explanation. From this point of view ``the moment'', the
entry of the divine into time, the miracle, receives no significance.
``For the Socratic view every man is himself the center, and the
whole world centers only in him, because his self-knowledge is a
knowledge of God.'' He who learns the truth learns it by himself;
the teacher is only the occasion; the moment at which the truth
dawns is, in relation to the truth that is, and that inheres in the
essence of man, a vanishing and to that extent indifferent point in
time. ``The point of departure is a nothing; for in the same moment
that I discover that I have known the truth from eternity without
knowing it, in that same instant that moment is hidden in the
eternal, taken up into it in such a way that I, so to speak, cannot
even find it, even if I sought it, because here there is no here and
there, but only an \textit{ubique et nusquam}.'' It is otherwise
from the exclusively religious standpoint. Here man is presupposed
to be, not in the truth, but in untruth. If now he who through the
guilt of the race and through his own guilt has lost the truth is to
come to the cognition of truth, namely that cognition of truth
% ---- printed p.7 (PDF 20) ----
\opage{7} which leads to blessedness, then it becomes necessary that God not
only enlighten the one in darkness about the untruth, but also, by
himself entering into time and existing as a man, give the man who
is in untruth the condition for attaining the truth. Through God's
self-communication in time the moment, the definite point in time,
and therewith the miracle, receives absolutely decisive significance.
The miracle is the paradox of the understanding. By holding fast to
the miracle and setting up the paradox, S.\ Kierkegaard has secured
religion against the encroachments of science and raised a bulwark
against the incursions of negative critique; but then he has also,
with decisive aim at his ``introductory problem, not to Christianity,
but to becoming a Christian'', come to a halt at the decision that
Christianity is not a doctrine but an existence-communication, and
an existence-communication expressed in the problem: ``that the
individual's eternal blessedness is decided in time through the
relation to something historical, which is furthermore historical in
such a way that there is taken up into its composition that which by
its essence cannot become historical, and thus must become so by
virtue of the absurd''.

According to the negative critique it is science that declares
religion an absurdity and thereby means to dissolve it; according to
the Kierkegaardian settlement it is religion that, in spite of
objective knowledge, declares itself an absurdity and thereby once
and for all excludes science. From this it is seen that religion and
philosophy are not merely heterogeneous, but absolutely
heterogeneous. With this presupposition, then, is a philosophy of
religion possible?

\runhead{c) The Philosophy of Religion as Inverted Science.}

The exclusive relation to God, determined by the miracle, which is
religion's own inescapable presupposition, science cannot
appropriate. This presupposition is
% ---- printed p.8 (PDF 21) ----
\opage{8} the incomprehensible, the unfathomable, and can be acknowledged on
one condition only, namely that human knowledge acknowledge a
barrier, an absolute boundary, at which it must stop. What in other
regions of cognition shows itself contradictory and unreasonable is
always something lower, above which science can freely raise itself;
it can discuss this lower, dissolve it, and by its dissolution
continually advance further. But religion, what is contradictory for
knowledge, is not something lower but something higher, for religion
is the highest ideal power of life; where this power, by virtue of
its original essence, comes forward with ``an absurdity'', there the
way is objectively impassable, there reason finds the passage
barred, there science must turn aside.

But what is absurd for scientific thinking, ``the paradox'',
nevertheless falls within the thinking consciousness, as surely as
it falls within religion; from which it follows in turn that what is
unthinkable for science can be thinkable in religion. There is an
essential difference between overstepping the boundary of knowledge
and overstepping the boundary of thought, between falling outside
knowledge and falling outside thought. The difference between the
thinkable and the unthinkable is here to be determined by the
principle from which thinking proceeds. That the paradox falls
outside knowledge is certain; but does it therefore fall outside
thought? In order that thought may strike upon the paradox, thinking
must of course itself bring the paradox forth; the paradox is
unthinkable only by being thought; its dialectical knot is an
intertwining of absolutely heterogeneous threads of thought.
Existentially the knot is insoluble; but it by no means follows that
it must also be so intellectually. The intellectual solution is
conditioned on the heterogeneous threads being unraveled and
separated from one another. Science would not be able to apprehend
and determine its principle of knowledge if it
% ---- printed p.9 (PDF 22) ----
\opage{9} were not able, by the opposition between what can be known and what
cannot be known, to characterize, if only negatively, the principle
of faith that is impenetrable to knowledge. But if it is possible
for science to separate the absolutely heterogeneous principles,
then it is of course also possible for it to separate the thoughts
proceeding from these principles, possible therefore to distinguish
between thoughts of knowledge and thoughts of faith, between
concepts of knowledge and concepts of faith. With the peculiar task
of clarifying the relation of refraction between the heterogeneous
thoughts, the philosophy of religion is most nearly to be determined
as a \emph{boundary science}.

The philosophy of religion does not, however, remain standing at
general determinations of the boundary between faith and knowledge;
it cannot possibly let religion develop its content without letting
a whole system of concepts of faith come to development. But the
religious content of faith is, according to the principle, a
supra-scientific content; here then the contradiction enters, that a
supra-scientific content is taken up into science. If now the form
of science were actually confused with the true original form of the
content and introduced and maintained in place of it, the
contradiction would certainly be insoluble. But this is by no means
the case. Scientific reflection has precisely the pliancy that it
can bring the form of knowledge to reflect another and opposite
form. For just as clear day can bring the eye to forget the gleam of
the light and look at the objects, so reflection too can bring
thought, in knowledge, to look at faith and forget knowledge. The
philosophy of religion is so far from forcing a foreign form upon
the religious content that its development of form consists, on the
contrary, in freeing it from the admixtures of immediately
individual representations, images, feelings, utterances of will,
etc., with which in the life of faith itself it is always grown
together. Insofar as reflection, as regards form as well as content,
sublates knowledge through knowledge,
% ---- printed p.10 (PDF 23) ----
\opage{10} in aiming to clarify not a scientific but a supra-scientific truth,
the relation between the truth and science is essentially inverted,
and the philosophy of religion is in this respect to be
characterized as an \emph{inverted science}.

What is decisive for every science is the grounding. Now, to be
sure, any thorough development of a content whatever can be regarded
as a peculiar development of the matter, thus as a development in
which the matter unfolds its moments and grounds itself. But in all
direct sciences the principle of the grounding of the matter is at
bottom one with the universal principle of knowledge, so that the
grounding of the matter coincides point for point with the
self-grounding of the science in question. In the inverted science,
on the other hand, the grounding is inverted. The fundamental
principle here is so far from being absorbed into the principle of
knowledge that the two absolutely heterogeneous principles, that of
the miracle and that of reason, everywhere repel each other and in
their negative relation determine each other negatively. Of this
negative relation, with the reciprocal determinations corresponding
to it, a clear knowledge can now certainly be given; but the
universal knowledge, hovering above the opposed principles, is at
the present stage no longer able to take the part of its own
principle; the inverted science has precisely the task of putting
faith into force and bringing knowledge itself to bow before it. As
it goes with the grounding principle, so it must go with the
grounding. The inverted science is neither to conceal nor to deny
the universally valid grounds of knowledge, but it is to bring to
clear insight how the self-confirmation of religion takes place
precisely through the grounds of knowledge consistently dissolving
themselves and giving place to the absolutely self-valid grounds of
faith. By grounding the conversion and dissolution of the grounds of
knowledge into grounds of faith, the philosophy of religion grounds
its own scientific dissolution and, on these terms, acquires
% ---- printed p.11 (PDF 24) ----
\opage{11} significance as a science that offers itself up to religion,
\emph{a science that sublates itself in doctrine of faith}.

\subhead{Religion and Mythology.}
\parmark{2}

On the relation in which philosophy places itself to religion
depends the scientific conception of the opposition between true and
false religion, between religion and mythology. Speculative reason
makes mythology into religion; the critical understanding makes
religion into mythology; only in inverted knowledge can the
separation in principle between religion and mythology be carried
through. We shall consider these points of view more closely.

\runhead{a) Mythology Is Essentially Religion: The Speculative Conception.}

With the presupposition that the truth in religion and the truth in
science are not two mutually heterogeneous truths but one and the
same absolute truth, the speculative philosophy of religion must
consistently acknowledge religion in every mythology, and regard the
mythical clothing as a husk that can easily be removed. Religion is
spirit, and for spirit as spirit the finite representations are
indifferent. To fathom the essence of spirit is to fathom the
essence of religion; in order to fathom the essence of religion,
philosophy develops the concept of religion. The concept of religion
is in its generality a genus-concept, but according to its
particular modes of appearance a species-concept. The Idea of
religion has divided its content and laid out its moments in a
plurality of more or less one-sided religions, in order through
these to gather itself out of
% print: "Adpredelsen" (for "Adspredelsen")
dispersion and come
% ---- printed p.12 (PDF 25) ----
\opage{12} to completion in the one, true, perfect religion. ``It has
been'', says Hegel, ``the work of spirit through millennia to carry
out the concept of religion and make it the object of
consciousness.'' The task is to show how the human spirit,
beginning in complete naturalness, has been able step by step to
raise itself above the nature-bound condition to ever greater
freedom and at last to reach the turning point at which spirit
comprehended itself as spirit and expressed the concept in the
exemplary existence that is to be appropriated and actualized by all
in the form of community. Hegel's systematics in the philosophy of
religion is characteristic in this respect. From \emph{nature
religion}, the stage at which consciousness, in a fermenting unity
of nature and spirit, intuits the divine, through the
\emph{religions} of \emph{spiritual individuality or free
subjectivity} --- those of majesty, of beauty, and of purposiveness
--- the development advances to the apprehension and determination
of the religion that is clear in itself, revealed, corresponding to
the concept of revelation, \emph{the absolute religion},
Christianity.

According to this view all religions are related, all sprung from
one root, all branches on one trunk. The old distinction between the
nature religions and the revealed religion is overlooked by Hegel,
even to such a degree that the ``religion of majesty'' in Judaism is
--- for the sake of the concept --- placed lower than the ``religion
of beauty'' among the Greeks and the ``religion of purposiveness''
among the Romans. What now follows from this? That the Jewish
religion is mythology just as well as the Greek and Roman! But if
Judaism is mythology, then Christianity is also mythology: the
absolute religion is of the same origin as the relative ones. With
the absolute religion, then, the mythical has not disappeared; on
the contrary, it has become transparent to spirit as something
belonging to the peculiar form of religion.
% ---- printed p.13 (PDF 26) ----
\opage{13} Only when the concept represented in religion is comprehended in
knowledge, that is, when religion is transformed into philosophy,
does the mythical disappear.


``The absolute religion,'' says Hegel, ``is the revealed. Religion
is the revealed only when the concept of religion is itself \emph{for
itself}, that is, when its concept has become objective to itself,
not in a restricted, finite objectivity, but such that it is,
according to its concept, objective to itself. We thus have two:
consciousness and the object; but in the religion filled with itself,
the revealed religion, which has grasped itself, the content is
itself the object, and this object, \emph{the self-knowing essence},
is \emph{spirit}. Here for the first time is spirit as such the
object, the content of religion, and spirit is only for spirit. In
being content, object, it gives itself the other side of subjective
consciousness, and thereby a finite appearance. It is religion that
is filled with itself. This is the abstract determination of this
Idea, or religion is in actuality Idea. For Idea in the
philosophical sense is the concept that has itself for its object,
i.e.\ has existence, reality, objectivity, is no longer the inner or
subjective, but objectifies itself, only that its objectivity is at
the same time its return into itself, or, insofar as we call the
concept purpose, the fulfilled, accomplished purpose, which is
consequently objective.''\footnote{Religionsphilosophie II.
Berlin 1840. S.\ 192--93.}

But when the content of religion in this manner coincides with the
Idea of philosophy, when religious cognition is made one with the
absolute knowledge of speculation, who does not see, even at first
glance, the consequence? Religion demands faith in its supernatural
facts, holds fast to existential immediacy, % print: "Umiddelharhed" (h for b)
and will not let itself be transformed into philosophy; philosophy,
which has transformed the sacred
% ---- printed p.14 (PDF 27) ----
\opage{14} facts into mythical representations, is sure of knowledge,
declares all positive religion to be mythology, and gives up faith.

\runhead{b) Religion Is Mythology: The Critical Conception.}

The conflict between religion and philosophy began among the Greeks
very early, and as for the Romans, two augurs in Cicero's time could,
as is well known, not look at one another without laughing. To be
sure, religion furnished the original foundation for the life of
spirit among the Greeks as well as among the Romans; but in both it
also becomes evident that the form of religious representations
satisfies spirit only at the first, childlike stage of a people's
life: the law of development is that consciousness, with increasing
reflection and advancing culture, outgrows religion. When Socrates
had discovered the dialectic of subjectivity, set self-knowledge as
the goal of philosophy, and to that extent introduced ``new gods,''
popular religion could no longer withstand the attacks of the
philosophers. Even if the concept of myth itself was not developed
in antiquity with the clarity and depth that has been achieved in
our own time, it nonetheless soon became sufficiently evident, as
regards the lives of the gods and the deeds of the heroes, that
``the divine,'' to use an expression of Strauss, ``could not have
happened in such a way that what thus happened could not be
divine.'' This dissolution, which explains the myths but annihilates
the gods, is not a spiritual regression but, in the name of thought,
of reason, of self-consciousness, an essential advance. The
mythological is the unconsciously illusory; but this plainly
illusory element religion cannot put off without laying bare its
innate weakness and exposing itself to ridicule. The supernatural
phenomena in the New Testament are just as contrary to reason as the
events in the Greek war of the gods. Religion and mythology are
synonymous concepts, and the difference between lower and
% ---- printed p.15 (PDF 28) ----
\opage{15} higher, more imperfect and more perfect religions is only a
difference between lower and higher, more imperfect and more perfect
mythologies.

The higher a mythology stands, the more it is filled with essence and
spirit, the longer it will naturally hold out alongside the
consciousness of culture; but precisely the circumstance that the
most spirit-filled of all myths, the myth of the Risen One, has been
able to carry culture through eighteen centuries and to hold out
alongside the human spirit, until consciousness at last got the
better of all obscure presuppositions and came to itself, in that it
``as the free universal that is for itself, became transparent to
itself,'' precisely this shows what it really is that in the
``revealed'' religion becomes the revealed. It is not a fundamental
difference between myth and revelation that comes to light; no, what
now comes to light is something entirely different, it is the great
objective truth that the human spirit, not the spirit of the
individual human being, but the human spirit in its principle, is
its own divinity, is in the religious self-duplication both subject
and object, as finite spirit subject, as spirit in inner infinity
object. This is no blasphemous self-deification; it is, after all,
the discovery of the Absolute, of the Idea as Idea, of the
divine-human unity of the subjective and the objective.

By carrying through the principle of knowledge, philosophical
critique has consistently dissolved religion as religion---for to
want to preserve the religious when the principle of religion has
been dissolved is of course nonsense---and it has thereby at the same
time reached the boundary line at which the powers of life react:
the religious asserts its right, faith steps forward against
knowledge, and the heterogeneous principles measure themselves
against one another.
% ---- printed p.16 (PDF 29) ----
\opage{16} \runhead{c) Religion Is Absolutely Heterogeneous with Mythology: The Antirational Conception.}

There is, in the stance of the critical as well as of the
speculative philosophy toward religion, an ambiguity which must from
first to last dazzle the observer and confuse his conception; this
ambiguity lies in the use of the word truth. That the eternal,
absolute truth is only one, and that the one divine truth cannot
have fallen into discord with itself or have divided itself into two
mutually conflicting truths, is certainly a presupposition as
unambiguous as it is incontestable. But when one then without further
ado overlooks the difference between a divine and a human knowledge
of the truth, overlooks that what is absolute in God's knowledge must
be relative in man's, overlooks the difference in principle between
religious and scientific cognition, the ambiguity arises. It looks
as though philosophy actually entered into religion's own
presuppositions, let these presuppositions develop their religious
consequences, and thereby arrived at conflicting, mutually canceling
results. Indeed, if it actually went that way, then it could truly
be said that philosophical critique let religion itself carry out
its own dissolution. But that is not how it goes. Instead of
beginning with the religious principle, science on the contrary
begins with the principle of knowledge. By being applied to
religion, the principle of knowledge consistently becomes a
religion-denying principle. When it then appears as though critique
proves the self-dissolution of religion, this is only a semblance.
The truth is that critique gets out what it has itself put in; it
begins by putting in the denial of religion and ends by getting the
denial of religion out: the dissolving critique moves in a circle.
For what is it, then, that science in principle posits as the
mythical? It is precisely the relation to God
% ---- printed p.17 (PDF 30) ----
\opage{17} peculiar to religion. To clarify and determine the relation to
God is to clarify and determine the fundamental conditions decisive
for the life of spirit. Now science ought surely to know that these
fundamental conditions present themselves differently in the domain
of faith than in the domain of knowledge. If it is truly the earnest
intent of science to get the life of spirit itself clarified and
understood from the ground up, then it must, at least for the time
being, be able to show the self-denial of disregarding its own
principle and letting religion have its say.

That the principle of religion is absolutely heterogeneous with the
principle of science can be seen from the grounding. All grounds of
knowledge are in the deepest sense grounds of necessity; the grounds
of religion, on the other hand, are grounds of will. God's will is
the Absolute of religion. If it were, as the thoughtless assert, a
contradiction that science should come to know of a principle
absolutely heterogeneous with knowledge, then it would have to be
precisely the negative critique that has entangled itself in the
contradiction. What does this critique want? To dissolve positive
religion by virtue of knowledge. Why does it want to dissolve
positive religion? Because its original presuppositions are
absolutely incompatible with all science. Now if science cannot come
to know of the presuppositions incompatible with science, then
critique is, after all, dealing with something it does not know
what it is, and its whole undertaking is meaningless; if it asserts
that it knows fully what it is that it dissolves, then it has thereby
conceded that science can come to know of what is intolerable in
knowledge. But that the principle of religion is absolutely
heterogeneous with the principle of knowledge means, after all,
nothing other than that this principle, peculiar in itself, must,
when it is to be conceived and applied as a principle of knowledge,
become intolerable to science. And why? Because the principle of
religion is God's unconditioned will.

To know God's will, to give oneself up to God's will,
% ---- printed p.18 (PDF 31) ----
\opage{18} and to do God's will, is the one and all of religion. But an
unconditioned will no science can assimilate. To be sure, will is
not for itself without essence and Idea; but by the content of the
Idea being posited in the form of will, and determined by will, the
thoughts of essence become thoughts of will, and the truth becomes
personal. In that the relation to God is posited as a relation of
will, there is expressly posited a fundamental difference, and with
it a boundary line, between the religious and the scientific, the
personal and the impersonal truth. The impersonal truth is cognized
by necessity of the Idea; it is inherent in the essence of man and
cannot be lost without man losing his reason and ceasing to be man.
The personal truth, the truth of will, is on the other hand identical
with freedom and can be lost when freedom is lost, i.e.\ when the
relation to God is disturbed. If the philosophy of religion is truly
to begin with the principle of religion, then it must begin, not
with a knowledge \emph{of}, but with a knowledge \emph{about} two
absolutely heterogeneous principles.

Originally---so runs the word intolerable to science---man was placed
in the incipient relation to God in such a way that he had therein
all the presuppositions and conditions for the cognition of truth
and of freedom. Through the misuse of will, and thus through the
misuse of freedom, he has lost freedom, has sunk down into unfreedom
and with it into personal untruth. In untruth he has lost the
cognition of God's will, of God as God, and thereby at the same time
the condition for cognizing himself according to his eternal
destiny. As fallen man's cognition of God has become illusory, so
too has his self-cognition become illusory. The capacity for general
cognition of essence is preserved, but the cognition of will is
lost: reason is preserved, but spirit is darkened. Now if the
relation to God rested on man's being able, of himself, by his own
spiritual impulse, to place himself in relation to God, then it
would also rest with man himself whether, after
% ---- printed p.19 (PDF 32) ----
\opage{19} having stepped out of the relation, he would step in again and
restore what he had forfeited. But that is impossible. It is not man
who began by placing himself in relation to God, but God who from
the beginning placed himself in relation to man; if the broken
relation is to be restored, it must come about by God, through an
express act of will, once again taking man to himself, thus by his
\emph{revealing himself} to man. This is---we repeat it---the
fundamental thought of religion. What the knower, from the objective
standpoint of a contemplation hovering above, will judge of this
thought, whether he will find it reasonable or unreasonable,
acceptable or unacceptable, attractive or repellent, is inessential;
the essential thing is that science come to know of it, take note
of it, and see how it develops in religion.

By the fundamental thought stated here, the difference between
revelation and myth, between religion and mythology, is irrefutably
established. This difference is determined not by outward but by
inward marks, not by historical reports but by the testimony of the
Spirit. But how, then, is the testimony of the Spirit determined? If
the testimony of the Spirit were nothing other than the religious
individuals' immediate utterance of their own inner, individual
experiences, then the Spirit would be without principle, and
religion not to be distinguished from mythology. For what a
multiplicity of theophanies, of proclamations of will, counsels,
guidances, commands, threats and promises from gods to men is not
contained in the ancient legends of the gods! This mixture of all
manner of imaginings, this motley swarm of more and less ideal
representations of God, cannot dazzle the mind that holds itself
convinced that religion is the life of spirit, and that life has its
principle. The relation to God which in one and all satisfies the
demands of spirit, of eternal life, which fixes the goal,
% ---- printed p.20 (PDF 33) ----
\opage{20} points out the way and offers the means, is and must be the true
one; every other that deviates from it, the untrue. Here, then, a
critique is required that surely distinguishes the true relation to
God from the untrue.

As the principle is, so must the critique be. Now since the
principle, as principle of will, the absolute principle of freedom,
belongs not to science but to religion, it follows from this that
the critique which is to be able to discern religion from mythology
in spirit and in truth must be, not scientific, but religious.

What, then, has science to do with this critique? It has to procure
the formal conditions; it has to bring out all essential
determinations, remove all irrelevant considerations, ward off all
confusing mixtures of thought, secure the consistent progress of the
movement, and come to know of the peculiar self-development of the
religious principle.

The relation to God is a relation of spirit to spirit. On this,
science can in general be at one with religion; but when the
relation is to be determined more precisely, the separation sets in.
For the relation has two main sides, the side of the Idea, which is
grounded in the necessity of reason, and the exclusive determination
of will, which is grounded in freedom: the former is the abstract,
for the side of the Idea can be brought out without will being
cognized; the latter, on the other hand, is the concrete, for will,
whether human or divine, cannot be cognized in truth without being
cognized in the Idea. As little as spirit is an Idea universal in
itself without will, just as little is it an empty, arbitrary will
without Idea: the will of spirit is will of the Idea, spirit itself
the absolute unity of Idea and will.

In the mythical relation to God, freedom is paralyzed, the element
of will corrupted. The corruption shows itself, from the divine as
well as from the human side, in egoism having stepped into the place
of the will of spirit;
% ---- printed p.21 (PDF 34) ----
\opage{21} the mythical gods are egoists just as much as men are. The
difference, however, is this, that only the human egoism is real,
the divine invented and fantastic. Then reason awakens and discovers
the fantastic; the will of the gods vanishes with the gods
themselves: only the necessity of the Idea remains. From the
standpoint of revealed religion the nature religions cannot possibly
be taken to belong to the kingdom of truth; they are on the contrary
testimonies to a relation to God sunk into untruth, a preceding
apostasy from the true God. The mythical gods are, in the language
of religion, idols, false gods; the unity of spirit and nature with
which mythology begins is a false unity. Reason can shed light on
the false gods, discover the laws of actuality and drive out
superstition; but the ambiguity of the Idea, the false unity of
nature and spirit, reason left to itself cannot possibly clear up,
and that for the good natural reason that the necessity of reason
and the necessity of the Idea essentially coincide. Reason can
cognize a world order in the connection of things, can in the state,
in art, in science unfold and clarify the system of Ideas; but
spirit as spirit, the eternally free spirit working as absolute will,
it cannot explain. Spirit as spirit has only one explanation, and
that is its own, unconditionally free self-revelation. Reason is so
far from being able to sublate the mythical that, by misjudging the
will of God and letting freedom in an ambiguous manner be born of
necessity, it is rather itself mythical in essence and ground. This
is the fundamental thought of religion; whether religion will be
able to develop it and defend it against scientific attacks, we
shall see; but this much is certain: it is the fundamental thought
of religion, and science must come to know of it.

In that the principle of will is posited as spirit's own divine
principle, the conception of what is peculiar to religion
% ---- printed p.22 (PDF 35) ----
\opage{22} must become, not rational, but antirational. The antirational is
by no means that which is in itself contrary to reason, but that
which stands against reason with so peculiar a determinateness that
the reasonable man, on closer reflection, must see the
unreasonableness of wanting to comprehend it. The opposition of
religion and reason is by no means determined by religion's supposedly
either breaking with the general laws of thought or denying reason
the right to test the objective truths grounded in the necessity of
the Idea; no, the opposition is determined by the principle peculiar
to religion, the principle of will. A scientific grounding of the
principle of will is impossible; but for the same reason a
scientific refutation of this principle is then also impossible. How
antirational the principle is becomes especially evident when we
consider that its essence is the miracle itself, that the absolute
principle of will is precisely the principle of the miracle.

When the talk is of the supernatural and ``unnatural'' in religion,
which offends reason, attention usually fastens on the miracles in
the outward, the external signs and works, as though it were they
alone that gave offense. In this one overlooks the inward
supernatural principle. And when the talk then is of this principle,
one again overlooks only too easily that the supernatural principle
is the principle of will itself. So-called sound reason finds it
objectionable to deny God a will, but it does not find it
objectionable to deny religion its miracle. If---so one
reasons---man, who is only a finite rational being, has a free will,
however imperfect, then the Godhead, the infinite rational being,
must necessarily have a perfectly free will. What is offensive is
thus placed, not in the Almighty's having a rational will, but in
an infinitely rational will's being supposedly able to resolve upon
something so unreasonable as to perform a miracle. But the
% ---- printed p.23 (PDF 36) ----
\opage{23} categorical expression of the unconditioned freedom of the
almighty will is the miracle. That God as God can will, can resolve,
can in a single determinate work utter a single determinate thought
of will, is precisely the miracle of miracles, the fundamental
miracle: religion, the relation to God itself, is the fundamental
miracle.

The difference between religion and mythology is, then, in brief
this: religion expresses the true, almighty principle of will,
mythology the false, impotent one; religion is known by the true
miracle, mythology by the false, merely imagined miracle.

\subhead{Revelation and Faith.}
\parmark{3}

If the God of spirit, that is, of will, were not revealed to himself,
the will of spirit certainly could not become revealed to man
either; but what is decisive for the concept of revelation is
nevertheless not, as Hegel supposes, that spirit as spirit
\emph{becomes} revealed, but that spirit, self-determining,
\emph{makes} its will revealed: revelation, as an act of freedom, is
expressly an act of will. On this rests, from God's side, the
indissoluble unity of revelation with the miracle, from man's side,
its indissoluble unity with faith. Just as faith becomes divinely
determined by revelation, so revelation in turn becomes humanly
determined by faith: revelation is a concept of faith, faith a
concept of revelation. The same Spirit that works revelation works
faith as well; and just as revelation becomes spiritless when its
objective elements are separated from faith, so faith too becomes
spiritless when the believing subjectivity tears itself loose from
revelation. Since the unity of the two, however, is not a merely
formal but a real unity, and thus a unity of opposition, it will be
necessary, for the
% ---- printed p.24 (PDF 37) ----
\opage{24} clarification of the matter, to illuminate the relation from its
different main sides: from the historical, through the opposition
between revelation and sacred history; from the metaphysical,
through the opposition between revelation and reason; from the
psychological, through the appropriation of revelation in faith.

\runhead{a) Revelation and Sacred History.}

As an act of will of God himself, revelation is not the
``manifestation'' of a speculative Idea, but divine action, a work
of the Spirit, a sacred fact. Revelation is, according to its
principle, historical; only it matters that the historical not be
brought under a false light. By applying the same measure to sacred
as to profane history, critique easily manages not only to shake
the credibility of the biblical narratives, but also to set up a
divorce between revelation and faith that must annihilate the
significance of both. A consideration corresponding to this is the
following: historically there appears a difference between
immediate and mediate communication of what is revealed. The
immediate communication is made not to all, but only to some. These
few, the elect, who so to speak receive revelation at first hand,
themselves experience the wondrous. But of what a man thus
experiences he has, after all, a direct, immediate certainty; and to
direct, immediate certainty there corresponds knowledge, not faith.
I do not believe that I exist, for I know it; I do not believe that
I actually see and hear, and that there are visible and audible
things, for I know it. So with the elect: they do not \emph{believe}
that a revelation has taken place, for they \emph{know} it. The
multitude, on the other hand, the non-elect, that is, those who do
not themselves experience the wondrous, but, by relying on reports
of what others have experienced, receive revelation at
second hand,
% ---- printed p.25 (PDF 38) ----
\opage{25} and thus properly historically, become the believers in the
proper sense. For to believe in a revelation one has expressly
received is, after all, no great matter, but to believe in a
revelation one has not received is a strenuous task. In this way
revelation and faith are made to fall outside each other:
\emph{blessed are they that have not seen, and yet believe}. By
falling outside faith, revelation becomes outwardly objective, and
thus spiritless: by falling outside revelation, faith, with an empty
inward whose content must be sought in something outward, becomes
formally subjective, and thus spiritless. This two-sided
spiritlessness, however, concerns not religion, but critique alone.

In order to satisfy the demand of spirit, sacred history takes up and
unites absolutely heterogeneous elements. God proclaims his will to
the human being; the proclamation takes place immediately, in a
finitely determinate manner. By the immediate determinateness the
supernatural wants to be known; but the determinate immediacy, in its
mode of expression, clings at the same time to the natural, to what is
sensuous in the phenomenon; content and form, the inner and the outer,
stand in decisive opposition: the content is supernatural, the form
natural; the inner is supersensuous, the outer sensuous. This is a
contradiction. Now if the content were commensurable with the form,
and the relation between the inner and the outer an immanent relation
of necessity, the contradiction would naturally have to be resolved in
knowledge. But the content is precisely incommensurable with the form;
the relation between the inner and the outer is by no means grounded
in necessity, but on the contrary posited with divine freedom. The
contradiction is that the supersensuous is to be sensed, what is in
itself invisible seen, the inaudible heard; such a contradiction must
be resolved in faith. Only faith in the religious sense possesses the
happy unity of infinite reflection and childlike immediacy, the
wonderful
% ---- printed p.26 (PDF 39) ----
\opage{26} gift of knowing the supersensuous by the sensuous, and of joining the
natural together with the supernatural. Faith binds together extremes
of the holy and the sinful, of divine omnipotence and human impotence:
by faith, and only by faith, is the human being set in a relation to
God that is at once finite and yet infinite, as spirit to spirit.

We learn it from sacred history. By immediate revelation Abraham
receives the command to sacrifice his son Isaac. Now if such a command
could actually be apprehended with the same immediate, sensuous
certainty with which one can apprehend a human command, the whole
narrative of Abraham's obedience would be without religious
significance. What Abraham does, anyone else in his place could and
would have done. For who is it, after all, that commands? Jehovah
himself, the almighty Creator! Here, then, on the presupposition of
sensuous certainty, there can be as little thought of a possibility of
evading the command as of any danger in fulfilling it. What indeed did
the man venture who knew within himself that the God who spoke was his
friend---when, mark well, he knew it with sensuous, undeceivable
certainty? The Almighty might, after all, merely wish to let Isaac
% print: "døer" for "døe".
die in order to raise him up again, or, if it did not please him to
raise the dead one, give his chosen one a thousand sons in compensation
for the one. No, the religious significance lies precisely in this,
that Abraham, by the nature of the case, could not possibly have any
sensuously decisive certainty, and for just that reason had to have
his assurance in faith. He believed he heard Jehovah's voice; he
believed that the command, although, incredibly enough, it called for
the sacrifice of Isaac, proceeded from God, and not, as the natural
heart must after all find far more credible, from dark fancies or from
an evil spirit. --- If we turn to the Gospel, the contemporaries needed
of course only to use their natural senses in order to see and hear
% ---- printed p.27 (PDF 40) ----
\opage{27} the man Jesus; but to discover in the Son of Man \emph{Christ, the Son
of the living God}, natural seeing and hearing were not sufficient:
\emph{blessed art thou, Simon, son of Jonas, for flesh and blood hath
not revealed it unto thee, but my Father which is in heaven}
(Matt.\ 16).

Just as little as an outer revelation can acquire religious
significance without faith, just so little can faith acquire religious
significance without a corresponding inner revelation. If the inner
revelation is lacking, accounts of the supernatural will either be
rejected with unbelief or be accepted in superstitious fashion
according to a merely historical faith. For all straightforwardly
historical faith in sacred history is superstition. Sacred history
differs from other history, namely, in this, that it wants to be
believed not with, but against, all probability. On account of the
improbability, what is historical in revelation is, as regards its
acceptance, not straightforwardly supporting but on the contrary
repelling. The highest that can be attained by historical testimony is
only the probability that the witnesses themselves believed in the
improbable. The improbable, which is otherwise a mark of the
unhistorical, is in sacred history precisely the sign of the
religious, insofar as it is the sign of the miracle.

The Gospel accounts of Christ's resurrection by no means make the
resurrection itself probable, any more than they give any perfect
assurance of its actuality. This is no objection against the
trustworthiness of the accounts; on the contrary, they could not be
more trustworthy than they actually are, but the trustworthiness
dwells in the repellent. Without the historical testimonies religious
faith would lack its supernatural presuppositions, determined by
determinate facts, and with them its proper foundation; but it is
nonetheless true, for all that, that it is not the outward events,
improbable in themselves,
% ---- printed p.28 (PDF 41) ----
\opage{28} but what is inward in the events, the revealed, which, revealing itself
anew to the gaze of faith, sublates pastness and makes the divine
present. The outer revelation is, on condition of faith, expressly set
down in sacred history, in order that the historical may in turn, on
condition of faith, be converted into inner revelation.

\runhead{b) Revelation and Reason.}

The opposition between revelation and reason is designated, briefly
and sharply, as an opposition between unconditioned freedom and
absolute necessity, between truth in faith and truth in knowledge. To
want to veil what is conflicting and contentious in this opposition is
in vain. People have appealed to the fact that truth must be one and
undivided, in all respects in agreement with itself; but as soon as
one sets about stating where the highest criterion of truth lies, in
revelation or in reason, the conflict breaks out.

When the question of the truth of revelation is made dependent on the
judgment of reason, science comes forward with objections, each more
loud-voiced than the last. Revelation in the religious sense is an
unthinkable concept, equally self-contradictory whether we look to the
essence of the human being, or of the world, or of God.

``If we have regard,'' says a profound thinker, ``merely to human
nature in itself, we convince ourselves that revelation
presupposes a receptivity in the human spirit which entirely conflicts
with the essence of spirit. For since the essence of spirit consists
in activity, it can be in suffering only insofar as it is at the same
time in acting. But through a revelation, an immediate action of the
highest being upon the human spirit, nothing is left over for the
latter but an absolute suffering; for the highest being is absolutely
active,
% ---- printed p.29 (PDF 42) ----
\opage{29} but the necessary correlate of absolute activity is absolute
passivity.'' (Schelling). Regarded from this side, then, a revelation
is impossible. Revelation shows itself just as impossible when we
consider the world, the natural order of things. Revelation is a
miracle, but the concept of miracle sublates the concept of nature.
Here no petty criticism is needed to find out the difference between
reasonable and unreasonable, necessary and superfluous miracles; the
point is that a single miracle annihilates the whole reality of nature.
For what is nature? A revelation of the Idea necessary to reason; if
not, then a sensuous, fantastic shadow-being. If things have a natural
ground and a necessary connection, then the miracle is impossible; if
the miracle is possible, then the ground of nature is hollow, the
connection loose, and universal reason a misty fancy. A miracle such
as, for example, that a staff is changed into a serpent and the serpent
back into a staff is enough to show how a nature-less, unconditioned
will can play with all conditions, treat finite causes as meaningless
occasional causes, abolish the law, and rob nature of its reality.

The same absurdity shows itself with regard to God's own essence.
Religion teaches, to be sure, that God is unchangeable; but if we
accept its stories of revelation, then God must in himself be very
changeable. A comparison between present and past shows clearly
enough that in the course of the ages a remarkable change must have
taken place in his relation to the world and to human beings. The same
God who formerly spoke, and spoke almost incessantly, and worked so
many miracles that the Old and New Testaments can hardly contain them,
behold, he has in later times become entirely silent and works no
miracles at all. From the standpoint of revelation this is
inexplicable. If it is not the purpose of the miracles that they
should strengthen faith and combat unbelief,
% ---- printed p.30 (PDF 43) ----
\opage{30} what then is the purpose? But if a strengthening in faith is the
purpose: who then will explain to us why the miracles fail to appear?
If there was ever a time when human beings might truly need to hear an
actual word of God, an almighty word from God's own mouth, and to see
an actual work of God, it would surely have to be this present time;
for just as the soil after long drought and heat needs refreshing and
fructifying rain, so must pious hearts, in the scorching sun of doubt
and of culture, at last need refreshing and fructifying miracles. If
faith in a supernatural revelation is to be kept alive, the
unchangeable God must from time to time change himself to a remarkable
degree.

But is a revelation, then, even necessary? Religion wants to ground
the necessity on the universal corruption of the human race through
the Fall; but here, if anywhere, there appears a striking disproportion
between means and end. The corruption, it is said, is universal, and
the need for the pretended help consequently universal as well; and
yet, although all have become through the fall equally guilty, equally
helpless, equally needy, the help has nonetheless by no means extended
itself to all. Originally it was only the Jews whom God came to help by
special revelations; the heathen, the great majority, had to help
themselves as best they could by the use of reason. But to help some
and pass over others, what is that other than divine partiality,
groundless arbitrariness?

That this is on the whole ``a reasonable line of thought'' no one will
surely deny: reason, as reason, when it is to make its own principle
of cognition valid, cannot judge otherwise. But when reason cannot
judge otherwise, what use is it then to want to conceal that revelation
is and must be an offense to reason.
% ---- printed p.31 (PDF 44) ----
\opage{31} If, conversely, the matter is made dependent on the verdict of
revelation, then the judgment will go straightforwardly against reason,
and it may easily seem as though reason's right were being violated.
Reason supposes that it can explain, comprehend, and see through
spirit; it is this supposition that religion rejects as a fancy. For
revelation alone explains spirit and opens its depths: revelation is
the sun of spirit. Just as the sensuous eye is dazzled when it stares
into the sun, so reason is dazzled when it would stare at revelation.
Reason does not know what spirit is, and cannot know it; for it does
not know what a holy, almighty will is, and cannot know it. That
reason can take up into its system of universally valid concepts the
concept of will too, is not to be denied; by deifying the universal
will bound to the necessity of the Idea it can even suppose that in the
concept it actually acknowledges God's will. But the acknowledgment
that reason grants to the divine will by robbing it of the
unconditionedness of self-determination and raising it to the rank of
``actual concepts'' has no more significance than the acknowledgment a
heathen emperor once showed Christianity by taking Christ up among his
household gods. It is of no use to want to satisfy reason, which
inquires according to grounds of necessity, with assurances that God's
will is one with God's essence, and God's essence and will one with
the eternal truth: every time the will would express itself in
\emph{this} or \emph{that} commandment, in \emph{this} or \emph{that}
work, and thus as will, the grounds of necessity vanish; the
conditions disappear in the unconditioned; the sting of
unconditionedness wounds reason, and wounded reason complains of
arbitrariness. But when reason cannot comprehend the unconditioned
freedom of the almighty will, and thus cannot comprehend the highest
presupposition of revelation: how then should it
% ---- printed p.32 (PDF 45) ----
\opage{32} be in a position to master religion, penetrate into the relation to
God, and comprehend faith?

Only on one condition can God's unconditioned will be cognized; but
this condition, although at bottom one, is two-sided: from the divine
side revelation, from the human side faith. What is first and decisive
in revelation is not what God wills of the human being in this or that
matter, but that God expressly wills something of the human being, has
a will toward the human being, and thus, unconditionally: that God has
will. Revelation is the divine act of freedom, the ideal action in
which spirit originally proclaims itself as spirit. Faith, like
revelation, is not an involuntary stirring or an unclear feeling, but
an inner action, a first act of freedom, by which the will of spirit
comes into being in the human being, and by which the human being
essentially becomes spirit. What is first and decisive in faith is not
this: that the human being, concerning this or that, immediately
assures himself of what God wills, but this: that the human being is
to will God's will, and thus unconditionally: that the human being is
to \emph{will}, is to cognize that the relation to God is a relation of
will, the relation in which alone there is spirit and freedom and
eternal life.

Into this relation reason cannot possibly penetrate; it looks as
though it had cut itself off from all grounds of essence: just because
it takes up and encloses all objective grounds in subjectivity itself,
in the ground of freedom, it looks in the eyes of reason as though it
were groundless. Or on what, indeed, should faith be grounded? On
revelation, and thus on a miracle. And on what, then, should the truth
of revelation be grounded? On faith, and thus on a miracle that is
itself grounded on a miracle. As an unconditioned relation of will,
the relation to God, when grasped in its innermost peculiarity, is an
anti-rational relation; by wanting to rationalize what is in itself
anti-rational, reason only betrays its own blindness. In the
% ---- printed p.33 (PDF 46) ----
\opage{33} highest concerns of spirit, for all absolute decisions of the will,
reason is blind, blind to the Fall and blind to the restoration; blind
to the Fall, for the secret of the Fall lies in the will; blind to the
restoration, for the secret of the restoration lies in the will.
Reason would fathom wisdom, and yet with all its wisdom cannot cognize
God in his wisdom; for all the secrets of wisdom hide themselves in the
will; the Idea of wisdom, the Idea of God's will, is revealed only to
the will. So religion judges ``blind reason.'' From this standpoint it
can be understood that it is not reason, but its overreaching, that is
an offense to the faith of revelation.

The many lengthy and seemingly endless disputes between revelation and
reason can be brought to a final decision only when they are led back
to the first point of departure, namely the opposition between the
principles of freedom and of necessity, between grounds of will and
grounds of essence. Every decisive ground of will always draws the
universal into singularity, sublates necessity by free choice, and
makes the ground of essence insufficient; in the sufficient ground of
essence, by contrast, the single is dissolved in the universal,
necessity is posited with the given conditions, and the ground of will
is unrecognizable. Both of these reason can see; but for that very
reason it can also see that the ground sublated in the will, which
gives the action as action the stamp of arbitrariness, cannot possibly
be a ground of reason. In this connection it can at the same time be
seen how idle and empty is the thinking that fancies it can remove the
separation between the heterogeneous principles by the assurance that
freedom and necessity must absolutely be one in the absolute essence.
For as soon as it is asked what it actually is that makes the absolute
essence absolute,
% ---- printed p.34 (PDF 47) ----
\opage{34} then, as already shown, revelation again posits freedom, reason
necessity, and the conflict begins all over again.

Here there is no other way out than to let the two absolutely
heterogeneous principles furnish the centers of two absolutely
heterogeneous groundings, two absolutely heterogeneous spheres of
cognition. Now if it is just as impossible to derive necessity from
freedom as freedom from necessity, and at the same time impossible to
deny the original fundamental unity of the opposed principles: then the
fundamental unity must necessarily be conceived as a Mystery
impenetrable to human knowledge.

This Mystery is no paradox. For the paradox appears only when the
knowing consciousness one-sidedly takes the side of reason and wants,
by straining the understanding to the utmost, to comprehend revelation
from the principle of reason. By acknowledging the Mystery we
acknowledge the inner infinity of both spheres of cognition. All
cognitions developed from the principle of necessity have full
validity in their own sphere, and reason to that extent an
incontestable right to investigate everything, test everything, and
develop the content of its knowledge infinitely. Likewise all
cognitions developed from the principle of freedom have full validity
in their sphere, and the faith of revelation for its part an
incontestable right to let its thinking test everything and judge of
everything; in this respect reason cannot set any limit whatever to
the consistent development of the content of faith.

But, it is asked, when revelation and reason now come into conflict
over the same, the selfsame actuality, the same, the selfsame facts:
what then? Why, then it follows from this that the same actuality must
show itself differently, be judged differently, and be assigned a
different significance, according as it is seen either in the light of
revelation or in the light of reason. That revelation and reason
contradict each other is then so far from being a proof
% ---- printed p.35 (PDF 48) ----
\opage{35} either against the religious cognition of truth of revelation or
against the scientific cognition of truth of reason, that it is much
rather a necessary consequence of the absolute heterogeneity of the
principles.

% print: "Aabenbaringen Tilegnelse" for "Aabenbaringens Tilegnelse".
\runhead{c) The Appropriation of Revelation in Faith.}

How entirely exactly revelation and faith correspond to each other is
seen expressly from the fact that the concepts themselves, in their
reciprocity, are defined by each other. What, indeed, is the revealed?
Just that which alone can be apprehended by faith, cognized, preserved,
and explained in faith; that whose truth alone can correspond to the
truth of faith. And faith? Why, just as sight is determined
specifically by the visible, hearing by the audible, thought by the
thinkable, so faith must be determined by what the natural sense
regards as the incredible, although it is nonetheless, by its nature,
precisely what is in truth credible, the only thing that in the most
proper and strictest sense of the word can be believed, namely what God
himself has revealed. By this fundamental determination every attempt
to involve reason in the concerns of religion is cut off once and for
all. But by this there is at the same time laid upon the philosophy of
religion the demand that it show how the relation between revelation
and faith can be developed from the relation to God itself, and
thereby open for us an insight into the psychology of faith.

If we consider the two main sides of the relation to God, the divine
and the human, each by itself, then the first points specifically to
the outward, the latter to the inward, in the miracle of revelation;
the opposition between the outward and the inward is here so prominent
that it grounds a distinction between the manifest and the hidden
miracle.

The relation to God is posited by a divine action; the eternal takes
its beginning in time, the supra-
% ---- printed p.36 (PDF 49) ----
\opage{36} historical becomes historical. To reason this is an offensive
contradiction, and the contradiction is sharpened to the utmost when
the gaze fixes on a fact such as this, that God at a given time has
revealed himself in a single human being, a human being in flesh and
blood; this human being, who was born, has suffered and has died, is
actually to be God! It looks, after all, as though contradiction itself
had become incarnate. If the eternal has in this way become historical,
then, says the understanding, incited by reason, it has become
something that expressly conflicts with its nature, and must therefore
have become so ``by virtue of the absurd.'' --- But what is absurd in
itself is \textit{eo ipso} what is impossible in itself, and what is
impossible in itself cannot be believed. Here, then, a difference must
show itself between what is impossible in the thoughts of reason and
what is impossible in itself. If the self-revelation of the eternal in
time is actually to belong to what is impossible in itself, then such
religious expressions as divine work, divine action, divine will must
necessarily be reckoned among what is impossible in itself. One of
two: either God is will, and then he can also become incarnate; or God
cannot become incarnate, and then he is not will either. The miracle as
phenomenon is something outward, but the cause of the phenomenon, the
will itself, is something inward. Offense at the outward, at the
manifest miracle, is thus an offense at the inward, the hidden miracle,
an offense at the absolutely religious presupposition that the almighty
will is God, that God's will is free, that God can act.

In the relation to God determined by revelation and faith, the free,
almighty action by which God comes to meet the human being is,
unconditionally, the first thing. In the vision, in sign and speech,
the communication takes place immediately. This vision, this sign,
this word is addressed not to human beings in general but to this
determinate human being, concerns not human affairs in general but
this
% ---- printed p.37 (PDF 50) ----
\opage{37} or these determinate affairs: revelation comes from infinity like a
lightning bolt that strikes down into the finite. The vision, the sign,
the word, this supernatural thing, takes place by an immediate
breaking-through of the natural; the miracle is manifest, and the
emphasis falls on the outward. The faith that immediately appropriates
such an outward revelation must, according to rational psychology,
consistently be conceived as superstition. Were the possibility of an
outward revelation to be counted as belonging to the reasonable order
of things, then, as was stated above, such revelations would surely
have to be able to take place at one time just as well as at another,
and fall to the lot of one human being just as well as another; but
anyone who in the actual concerns of life appeals to visions and
revelations is regarded, if not outright as a deliberate deceiver, then
at least as a superstitious enthusiast.


If we now turn to the next point, namely the human side of the
relation to God, the matter takes on an opposite appearance. The
external vanishes with its extraordinary influences; the divine will
transforms itself into a universal will, determinable by the Idea of
the good and yet supernatural; God withdraws; everything takes place
within the believer: the manifest miracle gives way to the hidden one.
From the standpoint of rational psychology the verdict on such a state
of soul is characteristically pronounced in Strauss's critique of
Schleiermacher's faith in the historical Christ. ``As a Christian,''
says Schleiermacher, ``I feel the union of my lower and my higher
self-consciousness facilitated, or my piety furthered. This furthering
cannot proceed from myself, since, on the contrary, nothing proceeds
from me but a stoppage of the religious life.'' Strauss asks: ``How do
you know that? You, that is, the modern-Christian I; for the
old-fashioned believing I did indeed know it from the same source from
which it knew that of itself it knew nothing true, that of itself it
could do
% ---- printed p.38 (PDF 51) ----
\opage{38} nothing good: it knew both from its fundamental presupposition, the
foreignness of the Spirit to it and its going out from it. But whence
does the modern I, which has come to consciousness of itself and has
become master of itself, get the means to make such a separation?
Concerning something external and accidental, such as a piece of news,
a historical notice, I can indeed know that I do not have it from
myself but have learned it from others; but here it is a matter of two
opposite movements in the innermost depths of the human heart, of
which only the one is supposed to have its origin there, while the
other has been led into it from without.'' The religious disposition
thus becomes recognizable by ``orthodoxy's childish sum, by which man
first \textit{per subtractionem} rids himself of all good, in order
then \textit{per additionem} to let it flow back to him as a gift of
grace.'' The hidden miracle is the unclear consciousness's own hidden
contradiction: such a reflection must consistently dissolve faith and
end in unbelief.

Faith is determined by the extremes which, in consequence of the nature
of the relation to God, must form between the external and the
internal. The external, the divine influence in its immediate
particularity, is a necessary condition for awakening the internal;
but as external it is only a condition, only a means; the internal and
its appropriation is the end. The vision is an external, but the
essence of the Almighty an internal; the word that is heard is an
external, but the thought in the word, the divine thought, an
internal; the command an external, but the will that commands an
internal. Now faith differs from superstition in that through the
external it grasps and appropriates the internal; and it differs from
knowledge in that the internal which it grasps is supernatural, hence
an internal that is not determined by its external directly or by
virtue of rationally necessary reflection: were faith to engage in
objective reflections on essence, it would have to end in unbelief.

% ---- printed p.39 (PDF 52) ----
\opage{39} Since superstition and unbelief both fall outside faith, they call
each other forth reciprocally. Superstition wants sensuous certainty
about the supernatural, a sign for every case, an omen for every
event, a command of God for every action; it clings to the external.
Understanding and reason then readily agree that such representations
of the supernatural are unnatural and dishonor the divine. So signs
and omens are rejected, and moral actions are brought under moral
laws. With the moral laws follows the necessity of the Idea; but the
necessity of the Idea is like a cloud that hides the face of God:
``God has no will, there is no God''---that is the declaration of
unbelief, and this declaration is empty and comfortless. The empty,
God-forsaken mind then reaches once more for supernatural
presuppositions, wants sensuous certainty once more: unbelief in turn
gives birth to superstition, and so on to infinity. From faith as
faith, on the other hand, superstition can never be born; for faith
converts the outward into its corresponding inward and lives by the
Spirit; nor can unbelief ever gain power over faith; for faith has
inwardness, and inwardness rests upon supernatural facts, upon free,
divine presuppositions that withstand necessity.

If, now, the internal is posited as the originally first, yet in such
a way that it is at the same time conditioned by a sublated external,
then surely a criterion of truth, an infallible mark of true and
false, must be demonstrable in such an internal; this criterion must
be determined by the relation between internal and external that
corresponds to the principle of revelation. Where, then, is the
criterion of truth to be sought? Not in feeling; for when feeling
conflicts with feeling, according as moods change, a criterion for
feeling is precisely what is needed here%
\footnote{That the essence of religion cannot be placed exclusively in
feeling, even if we take feeling in its substantial universality as
the feeling of absolute dependence, Martensen has clearly shown.
``When Schleiermacher designated the religious feeling as a feeling
of absolute dependence, he thereby, like the mystics, designated piety
as a \emph{God-suffering} state, a state in which man in the innermost
depths of his heart feels himself touched by the power in which we
live, and move, and have our being, a holy pathos in which man feels
himself to be the vessel and dwelling of the Godhead. But as this
description is reminiscent of mysticism, so it is itself mystical;
for it remains in indeterminate twilight \emph{which} this absolute
power is on which I feel myself dependent, whether it is the
impersonal Absolute, whether it is fate, or whether it is the ethical,
the holy, the good power. Only in the last case can the God-suffering
state, can the feeling of absolute dependence, also be a liberating
and uplifting feeling.'' Den christelige Dogmatik. Indledning
S.\ 10.\par
How difficult Schleiermacher finds it to make the distinction between
the essence of God and of the world properly recognizable is evident
also from what he says (``Der christliche Glaube''. Berlin 1835.
1.\ Bd.\ S.\ 169) about ``the divided and the undivided unity''; for
the unity of the world is indeed divided, but, seen from the
standpoint of the immanent Absolute, it is nevertheless also
undivided, and the impossibility of a feeling of absolute dependence
in relation to the world is thus by no means proved.}.
Still less in the imagination, since the imagination
% ---- printed p.40 (PDF 53) ----
\opage{40} can be moved just as easily by false as by true promptings. Nor does
it help to seek the criterion in universal thinking, since the
movements of thought can be just as diverse as the principles by which
it is moved. The criterion must, according to the nature of the
relation, be sought originally in the will.

As a general act of will, faith can be determined from the ground up
only by the will itself, namely by the fundamental will. In the
fundamental will the innermost in man, the human self, is immediately
touched, inspired, and moved by the innermost in God, by the divine
Self, by the will of God.

But when the separation of true and false is made dependent upon the
will, must it not then become arbitrary? By no means! The will of
which we speak here is not the egoistic formal will choosing among
finite, instinctive impulses,
% ---- printed p.41 (PDF 54) ----
\opage{41} but the will of the spirit itself, the will of freedom that
penetrates all the powers of the mind, all feelings, representations,
and thoughts, and gathers up the whole life of the soul, the will by
virtue of which the inner man can and must will nothing but his own
salvation. The awakening of this will of the spirit is conditioned
exclusively by the influence of the divine will; it therefore rests
upon a hidden miracle; hence may be seen its relation to historical
facts of manifest miracles; it is an internal that, in the sense of
faith, reflects its external. The unity of the will of the spirit
% print: "Aandvilliens" (no linking s), against "Aandsvillien" elsewhere in the paragraph.
with feeling, imagination, and thought gives, in the very deepest
sense of the word, self-understanding, a self-understanding that rests
solely upon an inward understanding with God. For the scientific
testing of true and false, the philosophy of reason sets up \emph{the
principle of conceptual contradiction}; for the personal testing
of true and false---for self-examination before God---the philosophy
of religion sets up \emph{the principle of
self-contradiction}\footnote{By conceptual contradictions we are
thinking here not merely of contradictions between purely objective
concepts of knowledge, such as a round square, a body falling outside
space, etc., but also, and chiefly, of such contradictions as arise
from a confusion of religion and science, when (cf.\ the lectures: Om
Hindringer og Betingelser for det aandelige Liv i Nutiden) one draws
scientific consequences from religious presuppositions, and religious
consequences from scientific presuppositions. Instead of genuine
concepts of knowledge and pure concepts of faith one then obtains a
series of bastard concepts, all of which are naturally full of
contradictions. That a scientific contradiction, moreover, taken
formally, can also be conceived as a self-contradiction, and
conversely every self-contradiction as a conceptual contradiction,
goes without saying, and is stated expressly only to spare
``verbal critique'' the trouble of making formal discoveries where
discoveries are superfluous.}.

In self-understanding before God is salvation, in self-contradiction
before God is perdition; these are the outermost poles of the life of
the spirit; about these two poles, in the inwardness of faith, all
feelings,
% ---- printed p.42 (PDF 55) ----
\opage{42} all fear, all hope, all the thoughts of the heart must move. The
peculiarity of the relation to God now shows itself in this, that
everything is at once so immediate and yet infinitely reflected, so
simple and yet, once reflection properly begins, infinitely intricate.
Faith is of a childlike nature: what is hidden from the wise and
prudent is revealed to babes; but faith is nevertheless first and last
of a spiritual nature; its essence is spirit, and ``the Spirit
searcheth all things, yea, the deep things of God.''

\subhead{The Word and the Spirit.}
\parmark{4}

But ``the deep things of God'' are not searched in the Spirit on the
same terms as the depths of the Idea in reason; for the Idea of reason
is sunk deep in necessity, but in God inwardness is free, infinitely
deepened in self-determining. To an unfree internal corresponds an
unfree external, and the utterance takes place with necessity; thus
essence must necessarily have its appearance, the real ground its
existence, the soul its body, etc. The free internal, by contrast, can
give itself a true expression only in a free external, and the
utterance takes place with freedom. The unfree internal can, under
given conditions, be forced forth to revelation in its external;
therefore nature can, if we may so put it, have no secrets from a
reason that in its researches penetrates ever deeper and always goes
further. But the free internal cannot be searched out in this way.
That a human being under certain circumstances lets himself be taken
by surprise and comes, against his will, to betray his inner self
shows only that his freedom is imperfect; for God's inner being,
freedom itself in its inner infinity, no reason, however profound,
shall bring to betray itself. What is to be known of God must be
revealed by God, revealed
% ---- printed p.43 (PDF 56) ----
\opage{43} in the Word; for the Word is the external of the spirit, the free
utterance of the free internal.

\runhead{a) The Originality of the Word.}

The significance of the word for the spirit, the free internal, is
seen psychologically in the fact that language is at once both
condition and consequence of man's self-consciousness. The liberation
of thought into conscious thought, the liberation of the will into
determining will, is possible only through language. The liberation is
twofold: partly a theoretical one, through the formation of objective
concepts and their development into science; partly a practical one,
through the inner clarification of everything that belongs in a
special sense to the life of the will. Through this twofold liberation
the word acquires a twofold significance: a theoretical one, namely as
word of knowledge, which speaks of impersonal truth, and a practical
one as word of will, to express the truth personally. That the word of
will presupposes another and freer opposition between the internal and
the external than the word of knowledge is plainly enough evident from
facts such as these: he who scientifically establishes a proposition
will not demand that anyone should simply believe him on his word, but
he demands that every thinking person who is able to follow the proof
shall see the thought become manifest in the word, and the internal
shine through the external; he, on the other hand, who speaks of his
disposition, his free personal inner being, can lay his hand on his
heart and expressly demand to be believed on his word, since the
internal of which he speaks always has something hidden that cannot
possibly be resolved directly into its external. What holds humanly of
the spirit of man holds divinely of the Spirit of God. ``What man
knoweth the things of a man, save the spirit of man which is in him?
Even so the things of God knoweth no man, but the Spirit of God.''

In the doctrine of the original unity of the Word with God, revelation
and reason---one would think---ought
% ---- printed p.44 (PDF 57) ----
\opage{44} easily to be brought into agreement. Speculative knowledge cognizes
the necessity of an inner self-objectivation in the Absolute as
spirit; such a self-objectivation is possible only by virtue of the
Word from eternity, the original Word, the Logos. The Logos is the
unity of essence and its revelation in divine thinking, the true
identity of reason and the Word. Through the Word reason is God; in
reason itself the Word is divine. ``In the beginning was the Word, and
the Word was with God, and the Word was God''---this evangelical
prologue is then so far from arousing offense among philosophical
thinkers or containing anything objectionable to philosophy that it
rather looks like a religious philosopheme, a borrowing from
speculative philosophy. And yet this is a great misunderstanding.
Reason and revelation are, as far as the principle is concerned, just
as much at odds in the doctrine of the Logos as on any other point
whatever. The Logos of which philosophy speaks is, if we disregard
the figurative and the mystical-symbolic, essentially that which runs
through all true science; the Logos doctrine, objectively understood,
is the apotheosis of speculation; in the speculative Logos doctrine
science glorifies its own principle of reason immanent in the Idea.
There is indeed a liberation of the internal here, insofar as the
Idea, the internal, is thought through the word and comprehended in
the word; but the liberation is only formal. One begins with essence
and ends with essence; the self-objectivation of the Idea into spirit
cannot rise above the necessity of thought (\textit{libera
necessitas}): ``man's knowledge of God is God's knowledge of
himself''; this is the unveiled secret of the word of knowledge, the
true meaning of the objectively scientific Logos doctrine.

From the standpoint of revelation the movement of thought is exactly
the reverse. Religion too, indeed, presupposes that the result and the
presupposition of self-objectivation, God's eternal freedom in the
Spirit through the Word, is grounded in God's
% ---- printed p.45 (PDF 58) ----
\opage{45} essence and insofar falls under the immanent necessity of divine
knowledge itself. But essence, ground, necessity, all categories
indispensable for an objective grounding in knowledge, are taken up
by the thought of freedom and enter into the will, without its being
possible to state the intermediate determinations necessary for the
concept of knowledge. Just as the philosophy of reason with its
rational Logos doctrine is not able to get from the necessity of the
Idea to freedom, so neither is the philosophy of religion with its
antirational Logos doctrine able to get from freedom to the necessity
of the Idea immanent in God's essence and exhausting all objective
grounds. Nor is that its task; for the task is to bring out, hold
fast, and develop the principle of revelation, the divine principle of
will: here there is no question of the Word as word of knowledge, but
as word of will, the unconditioned word of freedom, the word of faith.

For the word of knowledge freedom is a mystery; for the word of will
necessity is. The fundamental unity of freedom with necessity, the
conceptual unity of the word of knowledge and the word of will, is the
Word itself; the Word that was in the beginning with God and
presupposes revelation. To revelation corresponds mystery; both
categories belong essentially to the determination of the free
internal. The God of revelation is the God of inwardness, and
inwardness is infinitely deep.

In inwardness (\textit{ad intra}) God is revealed to himself through
the Word. For reflection the presupposition presents itself thus: the
revealed Word is, in reflected immediacy, different from the God who
reveals himself in the Word; God as subject has in the Word his inner
object; self-deepening is self-objectivation: \emph{the Word is with
God}. But the Word could not possibly reveal God's essence to God
himself if in its \emph{difference} from God it were not originally
one with God; the inner object is essentially itself subject;
self-objectivation is self-deepening: \emph{the Word
% ---- printed p.46 (PDF 59) ----
\opage{46} is God}. This sounds speculative, so speculative that it almost looks
as though the philosophy of religion were now in real earnest about to
penetrate into the ground of inwardness and search the deep things of
God. But here is the limit. The features of the Logos doctrine
adduced serve precisely to show, not the possibility, but the
impossibility of developing any concept of knowledge whatever from the
hidden essence of the divine self-objectivation. Questions about the
metaphysical Logos-substance, about God's second I in the Logos, etc.,
betray abundantly the vanity of wanting to make manifest what has not
been revealed to men.

The limit between what is revealed only to God in the Word and what
God has revealed to men through the Word is a limit of essence, which
categorically marks out the Mystery and sublates mysticism. The
proposition: the Spirit searcheth all things, yea, the deep things of
God, is the irrefutable fundamental principle of free, hidden
inwardness, not a proposition of knowledge, but unconditionally a
proposition of faith; for ``the things of God knoweth no man, but the
Spirit of God.'' The proposition: ``God hath revealed them unto us by
his Spirit,'' holds exclusively for believers, and is therefore
likewise not a proposition of knowledge, but a proposition of faith.
The proposition of faith posits the communication of the Spirit as an
act of will, an unconditioned act of freedom, which cannot be checked
by any science whatever.

As further proof that what is revealed through the Word is a
communication of free inwardness, this too serves: that in the
language of religion there is talk of God's counsels, God's
foreordained counsel, predestination, etc.; for such categories as
counsel, resolve, and predestination are characteristic, not of
necessity, but of freedom, not of the thoughts of knowledge, but of
the thoughts of the will. The word of revelation is a word of will and
as word of will the object of faith.
% ---- printed p.47 (PDF 60) ----
\opage{47} \runhead{b) The Trustworthiness of the Word.}

That the word of faith is, by its principle, heterogeneous with the
word of knowledge can be seen, when the question concerns the truth
communicated through the word, in assurance. From the comparison
between religious and scientific assurance we must of course at once
exclude that domain of knowledge where an objective grounding leads
directly to incontrovertible certainty, and confine ourselves to a
consideration of that \textit{cognoscendi genus} which expressly
requires a special regard both to the free inner being of the one who
communicates and to that of the one who receives, namely the
historical. In the assessment of historical reports the opposition
between religion and science shows itself chiefly in this, that
science, as indicated above, makes trustworthiness depend upon
probability, whereas trustworthiness in the religious sense is
independent of probability.

``The experience,'' says Hume, ``that a testimony usually agrees with
the fact, and the experience that this latter occurs only very rarely,
lead to a contest between two arguments, of which the stronger keeps
the field, yet in such a way that it loses as much of its strength as
the strength of the opposing argument amounts to. If, then, a Cato
told me that he had seen a Fabius act rashly, I should have to
calculate whether there were more cases in which, the conditions being
given, I had seen Fabius do something that approached imprudence, or
Cato say something that bordered---not exactly on being a lie, but,
what makes it in high degree easier for an erroneous report to be
possible, on being an incorrect conception and assessment of what was
given: and if one now assumes that the grounds of probability
\emph{for} the correctness of the statement are $= 1$, the grounds
\emph{against} $= \frac{9}{10}$, then the statement would indeed keep
the right, but would have only the slender surplus of a
one-tenth preponderance in probability.''
% ---- printed p.48 (PDF 61) ----
\opage{48} But how unreasonable it is to want to transfer such a calculation of
probability to the trustworthiness of revealed facts will strike the
eye when we consider that the fundamental presupposition to which all
such facts must be traced back, namely the presupposition that God
has word and will, that God has spoken to men, that God therefore can
actually speak, belongs to the most unreasonable, the most improbable
thing that any knower can imagine. To believe the improbable is
nevertheless by no means the same as to believe the impossible. The
improbable must be believed; the impossible cannot be believed! The
possibility of faith is conditioned precisely by a difference in
principle between the improbable and the impossible.%
\footnote{Even mathematics makes this difference recognizable. Since
probability is stated by a quotient of favorable and possible cases,
it makes an essential difference whether the quotient is
$\frac{o}{a} = 0$ or $\frac{a}{\infty} = 0$; the first denotes
impossibility, the latter only an improbability bordering on
impossibility.}

The trustworthiness of the Word has been disputed from the standpoint
of reason especially with regard to those acts of revelation in which
God directly addresses speech to man in such a way that the divine
word becomes one with the human. God's speaking to the patriarchs, the
Mosaic legislation, the word to the prophets must necessarily appear
to reason, striving after concept and insight, as something
unreasonable, improbable, indeed impossible. It is not primarily what
is said, but that it is said and the manner in which it is said, that
arouses offense. As explicable as it is that a human being can speak
with a human being, and the thinking human being speak with himself in
silence, just so inexplicable is it that a human being should be able
either to speak to God or to be spoken to by God. Conversation between
God and man, a language through whose words God himself audibly and
clearly communicates to man his
% ---- printed p.49 (PDF 62) ----
\opage{49} thoughts and counsel, reason reflecting by the understanding must
count among what is inconceivable in itself, the absurd. Were an
explanation to be possible, analogies would have to be demonstrable
in the universally human life of consciousness, in the sphere of inner
experiences; but whatever analogies, whatever inner proclamations of
the divine in man, whatever aesthetic, ethical, and intellectual
inspirations of the Idea one may appeal to, the likeness dissolves
into complete unlikeness. The only analogue that can be applied here
is the religious consciousness's own mythical self-duplication,
whereby its speaking with itself is transformed into a speaking with
God; but this is no analogue, it is the thing itself: the divine
communication is mythical, the Word is myth.


These objections of reason, however, are so far from weakening the
religious credibility of the revealed word that the believer, who
knows the principle from which they proceed, can find them very
reasonable and yet altogether meaningless. What according to its
nature must be determined by freedom cannot possibly be derived from
necessity. The question of the credibility of the word coincides with
the question of the reality of the relation to God: it is the
relation to God in its exclusive peculiarity on which the credibility
rests.

In his relation to God man has --- by God's word and will --- become a
self in the spiritual and eternal sense; his self-understanding can
therefore be grounded only in his understanding of God. But man's
understanding of God is possible only in that God himself makes
himself understandable to man; and God can make himself
understandable to man only through the word. How the word in outward
immediacy could sound forth to the elect, of this no explanation is
given, nor, by the nature of the case, can any be given. The one
thing on which everything turns is the transformation by which the
outward, the revealed miracle can become the means and condition for
the awakening of
% ---- printed p.50 (PDF 63) ----
\opage{50} the spirit, for the liberation of the inward: the explanation lies
in the appropriation. Directly, i.e.\ in sensuous immediacy, no word,
whether in outward or inward revelation, can be held fast by itself as
a word of God; such an immediacy conflicts, as already said, with the
essence of spirit. The understanding must be reflected in
self-understanding, and the reflection is steered by the will. What
matters is to will, as the word says, to understand God in the
spiritual and eternal sense, to will to become revealed to oneself
down to one's innermost being by becoming revealed to God.
Expressions such as ``God has said'', ``God says'', must in the
appropriation be transformed into: ``I believe that God has said'',
``I believe that God says''. Only by this transforming can man, in
God's word and will, obtain unconditioned, decisive presuppositions
for his human word, his human will: the believer must presuppose
God's word by himself positing it in faith; the believing
presupposing is in this sense a self-positing. Yet neither can this
self-positing be held fast immediately. For what is humanly posited
by the human word is indeed posited as that which has already
beforehand been posited by God in God's word. It is by the positing
and transposing of the divine into the human and conversely ---
immediate in feeling, finite in representation, but infinite in
reflection --- that man becomes able to bind his word to God's word,
to find his will moved by God's will, and to cognize that secret of
the spirit, that the God of the word can speak credible words, that
in the inwardness of faith there is such a thing as speaking with
God.

In what sense the word of revelation can be said to concern man and
to make man's understanding of God dependent on his original
self-understanding is then seen expressly from the fact that the word
with which revelation is
% print: "uldendes" (initial f dropped) for "fuldendes".
completed is not merely a word through a man to men, but a word
which, in its fullness of Godhead, its unity with God, is so human
that it is one with
% ---- printed p.51 (PDF 64) ----
\opage{51} man, that man is the word itself, and the word man. Here again is a
turning point where faith and knowledge, philosophy of religion and
rational speculation seem to flow together.

That the doctrine of the Logos belongs to religion just as much as to
philosophy has already been shown. In the eternal Logos no finite
separation of essence and will, of nature and spirit, can possibly be
thought; the Logos, the Word from the beginning, is, as the word of
essence and of will, one and the same word. Now if revelation
expressly teaches that the Word became flesh, and that the Word which
became flesh and dwelt among us is the Word itself, the Word which in
the beginning was with God and was God: does it not then follow that
the perfect revelation must be assumed precisely to hold fast the
unity of essence and will, necessity and freedom, life and light in
its divine Word? ``In it is life, and the life is the light of man''!
But when the word of life is one with the word of light, and the word
of faith consequently one with the word of knowledge, is it not then
a self-contradiction to wish to separate faith and knowledge as
heterogeneous, indeed as absolutely heterogeneous principles?

What is dazzling in this consideration is easily removed. For the
principle of knowledge, which keeps to essence, ground and necessity,
does not thereby mean to deny freedom; still less does the principle
of faith, which keeps to freedom, mean to exclude necessity: faith and
knowledge can in general easily agree in admitting an original
unity-in-difference of necessity and freedom. But when a development
is to take place, when it is a matter of method, when the question of
the starting points decisive for cognition presents itself, the
principles become heterogeneous. From necessity as starting point
there is no way to a complete grounding of freedom, and from freedom
as starting point there is no way to a thorough insight into
necessity: both ways lose themselves in the Mystery. The divine word,
whose validity is
% ---- printed p.52 (PDF 65) ----
\opage{52} grounded in knowledge, is heterogeneous with the divine-human Word of
which revelation speaks. The revealed Word, the Word which became
flesh and dwelt among us, is, according to the grounds of knowledge,
not credible at all. How indeed could it be credible, when the
``assumption of flesh'' is in its phenomenon so improbable, in its
concept so contradictory! That the word of knowledge must deny the
word of revelation is self-evident, since by virtue of the grounds of
necessity it must deny the miracle. For that the Word became flesh is
a miracle, not one among many miracles, but the miracle itself: the
revealed Word is the revealed miracle. According to the grounds of
faith, on the other hand, the Word is credible as the miracle; for the
grounds of faith are grounds of freedom, and the revealed Word is in
the very deepest sense revealed freedom. Contradiction of freedom is
self-contradiction in the spirit; doubt of freedom is despair in the
spirit. Faith in the revealed Word thus coincides with faith in
revealed freedom: only the Spirit of freedom can bear witness to the
credibility of the Word.

\runhead{c) The Witness of the Spirit.}

If the principal proposition of the philosophy of religion, that the
principles of faith and of knowledge are absolutely heterogeneous,
% print: "til at literærhistorisk Beviis" ("at" for "et").
really needs a literary-historical proof, then a more striking one
can hardly be given than that which emerges from the dispute for and
against the credibility of the word of the Bible. Precisely those
narratives on whose credibility the religious consciousness must lay
the very greatest weight, because they present phenomena of
revelation that are of the most essential significance for the
understanding of the relation to God --- the narratives, for example, of
Paradise and the Fall in the Old Testament, of Christ's birth and
resurrection in the New Testament --- precisely these narratives
science denies all credibility, and does so, be it noted, not out of
contrariness or ill humor, but by virtue of
% ---- printed p.53 (PDF 66) ----
\opage{53} objective grounds of proof, in all soberness and after the most
careful examination. Now in the assessment of historical facts a
stronger opposition between two principles can hardly be conceived
than when the same narrative which according to the one is declared
absolutely credible must according to the other be regarded as
altogether incredible.

For a true understanding of the historical foundation of revelation,
however, it is not enough that faith free itself from objective
science; what is required here is a deeper-going investigation of the
relation between the outward and the inward, between the historical
and the religious, the vanishing and the abiding, thus an
investigation of the relation between the word and the Spirit.

To build one's faith on outward historical testimonies is spiritless;
it is, in superstition, to cling to the outward and forget the inward.
But to separate one's subjective faith from its objective
presuppositions is also spiritless; for if disregard for the
historical is to mean disregard for everything outward, then the
hidden miracle is only an imagining when the revealed one is denied;
and if conceit goes so far that the individual with his inward
experiences regards himself as an isolated organ of revelation, then
we have fanaticism.

In the Old and the New Testament there occur phenomena of revelation
whose peculiarity corresponds so entirely to the notions and ways of
thinking of particular ages that the forms of the divine's appearing
must seem to be directly adapted to the childlike imagination and
power of comprehension of the men of antiquity. Revelation in Holy
Scripture is not confined to a proclamation of God's will through the
word; it is pictures of the world in great and wondrous strokes that
are unrolled before the sight. The events take place in two kingdoms,
a natural and a supernatural. The supernatural kingdom is populated
with angels,
% ---- printed p.54 (PDF 67) ----
\opage{54} good and evil; here, in opposition to the kingdom of light, is a
kingdom of darkness, in which the devil with his demons plays the
leading role, etc.: does this whole apparatus now belong to the
essence of revelation, is all this taken up into what abides? Here
critique is necessary!

But as the principle is, so must the critique be. The critique
required here is not a critique of knowledge but a critique of
faith: the difference is characteristic. In the first place, the
critique of knowledge will lay weight on the outwardly objective in
its separation from that part of the content which is assumed to
derive from the narrator's subjective conception and presentation; in
consequence of this there can then be a question only of the actuality
of the natural, for all traits of the supernatural must be set apart
by themselves and regarded as unactual. When a narrative from sacred
history has thus been dissolved into its component parts, there can
naturally be no talk here of any revealing word: the word of God,
which was to inspire the whole, sinks down to a plain narrative, an
old legend of greater or lesser ideal value. The critique of
knowledge is negative; it denies the word of revelation because it
denies its spirit. In contrast to this, the critique of faith, which
presupposes that the presentation of a content of the spirit must
always and essentially be a work of the spirit, will grasp the
narrative according to its inward totality; it will see the inward
reflected in the outward, the supernatural in the natural, precisely
in such a way that the presentation and what is presented fuse
together into one single ideal and yet real actuality, actuality of
spirit, actuality for the spirit. Thus the word about Paradise is by
no means the same as the narrative about Paradise; nor is the word,
the word revealed in the Spirit, about Christ's resurrection to be
confused with the collective result of all the Gospel narratives. By
the word as God's word, emphasis is laid not merely on what
% ---- printed p.55 (PDF 68) ----
\opage{55} is presented, but also on its being presented precisely thus in its
fundamental character. That a narrative is to be subjected to a
critique of faith will then mean: it is to be grasped in its
entirety as the expression of a definite phenomenon of faith that
reaches deep into the life of the spirit; and the question now is
whether what according to the narrative is referred to an outward
revelation can find confirmation in the inward one. The legend of
Paradise, for example, is not to be believed as an outward history
because it stands in the Bible; but the legend that stands in the
Bible is to be believed because all the essential traits in the
apparently outward history gather into a total expression of the word
of the Spirit, into an inward history to whose validity the Spirit
bears witness. Christ's resurrection is not to be believed because
the Gospel narratives present it as an actual fact; but the word about
the historical, actual and yet supernatural fact is to be believed as
a revealed word because in the believer's inward being there is an
actual and yet supernatural fact through which the Spirit bears
witness to the word.

By putting side by side narratives of unlike kinds one will easily
convince oneself of the necessity of critique. For no great
acuteness is required to see that the word about Christ's
resurrection stresses actuality in a different way from either the
word about Paradise or the word about the devil, and yet all of this
has its actuality. Now if the critique of faith, with the witness of
the Spirit, is not to become arbitrary, it must necessarily have its
corrective. This corrective cannot immediately be placed in the
believer's subjective consciousness; for when consciousness is held
fast in its immediate state, then the one believes one thing, the
other another; the one, for example, believes in the devil and appeals
to the witness of the Spirit; the other, on the contrary, does not
believe in the devil and likewise appeals to the witness of the
Spirit; the one believes in the possibility
% ---- printed p.56 (PDF 69) ----
\opage{56} of an eternal damnation according to the witness of the Spirit; the
other asserts its impossibility, naturally also according to the
witness of the Spirit: where now is the corrective? In the principle
of faith! For the principle of faith is one with the principle of the
relation to God, is the Spirit's own principle, and only by virtue of its
own principle can the Spirit bear witness.

\subhead{The Principle of Faith: The Method of the Philosophy of Religion.}
\parmark{5}

Religion is the work of the spirit. But, it is said, science too is
surely a work of the spirit: can the spirit then be at strife with
itself? We answer: the spirit as the Spirit of truth is not at strife
with itself; the strife presupposes that besides the Spirit of truth
there must also be a spirit of lies. This spirit of lies is just as
powerless in science as in religion; its power is conditioned solely
by the fact that the human spirit, in religion as in science,
misunderstands and misuses itself. To assure oneself that science and
religion belong to opposite spheres of spirit, one need only reflect
on the difference in principle that obtains between the
\emph{knowing} and the \emph{witnessing} spirit. The more the
principle of knowledge prevails in all its rigor, the more unfitting
it is to speak of a special witness of the Spirit. The Spirit does not
bear witness with our spirit either that the square on the hypotenuse
is equal to the sum of the squares on the legs, or that all bodies are
heavy, or that Quirinius was not governor in Syria under the Emperor
Augustus when the taxing took place. Impersonal truth transforms
testimony into a ground of knowledge and is absorbed in the grounding.
But with personal truth the case is the reverse; the relation to the
Spirit is here not a relation of necessity but unconditionally a
relation of
% ---- printed p.57 (PDF 70) ----
\opage{57} freedom: in the relation of freedom the Spirit is witnessing; its
witness is in the hidden miracle.

It is easier to misunderstand the witnessing spirit than the knowing
one. By grasping religious truth in a spiritless way, the believing
consciousness comes to hover between two extremes, according as it
places the personal decision chiefly in God's essence or in man's. If
the decision is spiritlessly shifted over to the divine side, then the
% print: "blindthen" as one word (for "blindt hen").
word of revelation is blindly a word of authority: God himself or his
representatives must then expressly dictate to those under tutelage
what and how much they are to believe; hereby faith sinks down to a
passive, immediate acceptance of the articles of faith prescribed by a
divine word of power: subjectivity is hollow, the believing soul in
itself empty, a \textit{tabula rasa}. If the decision is spiritlessly
shifted over to the human side, then faith --- whether its inwardness
otherwise comes to rest in a mystical twilight of gentle feelings or
is goaded in passion into an incessant running of one's head against
``the paradox'' --- becomes a purely subjective matter; it is then
more or less indifferent \emph{what} is believed, the whole emphasis
is laid on \emph{how} it is believed: the finite, spiritless
separation of \emph{What} and \emph{How} then leads to the well-known
sophism of subjectivity, that everyone is saved by his faith. In order
to draw the extremes mentioned under the principle of faith, and
thereby under the witness of the Spirit, we consider separately how
the content of faith proceeds subjectively from faith itself, how it
is determined objectively by the Confession of Faith, and what then
follows from this for the method of the philosophy of religion.

\runhead{a) Faith and Its Content.}

Faith is in the strictest sense determined by the relation to God. In the
relation to God man's freedom is one
% ---- printed p.58 (PDF 71) ----
\opage{58} with his absolute dependence, and faith, the most decisive act of
freedom, is in consequence the essential expression of dependence. In
faith nothing can be accepted and nothing taken up except what the
believing consciousness itself posits and presupposes: to believe in
God is to posit and presuppose God; this is the inwardness of faith,
its subjectivity and freedom, that the believer can believe only what
faith posits. But when freedom is thus held fast in inward infinity,
it is at the same time seen that faith itself and that in which it
believes do not let themselves be separated: as faith in God is, so
too is the God in whom one believes; the God of faith is, subjectively
understood, a production of faith. It is to this side of the matter
that negative critique exclusively keeps; but in doing so it
overlooks that the believer's infinite freedom is precisely determined
by his absolute dependence. Were the God in whom one believes nothing
other than a product of believing subjectivity, then it would quite
rightly be the highest of freedom that both God and faith should
vanish and man end by relating himself absolutely, through himself, to
the Absolute. But this is the secret of faith, that with regard both
to the hidden and to the revealed miracle it produces God only in that
it is itself produced by God: that faith produces God denotes the
infinite freedom, that it is produced by God the absolute dependence.
So inwardly pervasive is the dependence that he whose faith is
assailed by doubt must ascribe the doubt to himself and pray God for
faith. The question of negative critique: whence do you know, you
modern I, that all the good, all holy powers, all purity in the soul,
all the love of the spirit and all its hope come from God, and only
the evil, the weakness, the impurity, the selfishness from yourself?
% print: "omkrives" for "omskrives".
can then be rewritten thus: whence do you know, you modern I, even
yet what faith is and what it is to believe, since the critique that
dissolves faith does not know it?

% ---- printed p.59 (PDF 72) ----
\opage{59} What holds in the relation to God of the divine Self holds also of the
human self. The self which through faith is set in relation to God is
not what religion calls ``the sensuous man''
(ἄνθρωπος ψυχικος), for ``the sensuous man grasps not the things that
belong to the Spirit of God'': the sensuous man does not know what
faith is and cannot believe. The self which in faith produces God is
the spiritual self, the self which precisely through its relation to
God is produced by God, and therefore infinitely needs God. The self
that functions in knowledge has the certainty of knowledge and
possesses itself knowingly; but the self that functions in faith has
only the certainty of faith and possesses itself only believingly. By
functioning in two absolutely heterogeneous systems of cognition, the
one and the same human self becomes a double self, a knowing and a
believing one. But the doubling is no self-division; it is a movement
about two different centers, that of the necessity of the Idea and
that of freedom; it would be self-division only if the content of
truth that is to be sought in the one of the two spheres were also to
be sought in the other.

Now in order to see more closely how the content of faith develops in
and with faith itself, we fix our attention on the changes the
relation to God undergoes as revelation step by step assumes its three
fundamental forms. It is self-evident that the development takes
place by virtue of the principle of faith, and that the intellectual
grounding must be turned against reason, thus inverted.

At the first stage of revelation the relation to God is to be
characterized by the fact that the opposition between God and man is
stressed abstractly as an opposition between Creator and creature.
The Creator has to command, the creature to obey; that is the
condition of the spirit. With the relation to God begins man's
self-contradiction. He cognizes that God is holy and his will
righteous; he cognizes himself
% ---- printed p.60 (PDF 73) ----
\opage{60} bound to unconditional obedience. But the spiritual man, the
believing self that grants that God is right, is essentially different
from the natural man, whose selfish will, bound by drives, does not
grant that God is right; here are two men in one man, or rather: the
one and the same man has from the ground up become divided against
himself: this is the self-contradiction, and the self-contradiction is
the negative mark of the relation to God. The immediate cognition of the
self-contradiction, which is involuntarily felt as a contradiction
against God, a violation of his holiness, is the cognition of sin.
Faith as faith does not abolish the self-contradiction. On the
contrary! The self-contradiction, and with it the cognition of sin,
increases the more living and inward faith is. In the determination of
man's essence one sees rightly how different the explanation of faith
is from that of reason.

The man divided against himself reason must expressly refer to the
universal human essence that holds in individuals; the contradiction
between the spiritual and the sensuous man is a contradiction of
concepts, which is to be resolved by bringing the necessary
fundamental unity of human nature to clear consciousness. By virtue
% print: "I Kraft of" ("of" for "af").
of the fundamental unity, the lower elements in man's nature must then
be regarded as the spiritual life's own necessary presuppositions:
drives, inclinations and passions are natural elements of reason;
they are rational in their disposition, inasmuch as they are destined
to be mastered by the rational will that cognizes itself as the
universal and the universal as itself; the contradiction then comes
about in that the sensuously individual, the lower, opposes itself to
the higher, the universal of cognition, which alone ought to determine
the rational will. The consciousness of the contradiction that
obtains --- a consciousness of the strife between the actual man and
his ideal --- can now indeed in a certain sense be called a
consciousness of sin, insofar as it is a consciousness that
% ---- printed p.61 (PDF 74) ----
\opage{61} the ideal, the holy, has been violated: conscience, the involuntary
rising up of the ideal nature against the encroachment of the selfish,
shows this clearly enough. But this self-contradiction lying in the
consciousness of sin dissolves and vanishes in proportion as the
contradiction between the lower and the higher dissolves into an ever
clearer, ever purer and stronger self-assurance of the inner harmony
and fundamental unity of human nature. Without sin man, immediately
regarded, cannot be, for sin and strife indeed arise precisely from
natural immediacy, the particular selfishness which again and again
would draw the individual away from what is in itself rational, from
what is universally valid and necessary in self-consciousness. But
once the essence of sin has in this way been seen through, the
self-conscious man has taken sin upon himself, has essentially and
from the ground up himself sublated sin, in that he has resolved to
strive and suffer for the right of the ideal, has found and posited
the ideal good in his own self, has thereby come to himself, is in
truth reconciled with himself, and can by this reconciliation
accomplish his own self-liberation, his true redemption.


But what is the consequence of this rational explanation? That
revelation vanishes, that the divine will vanishes, that faith
ceases to be a principle!

When the divine will vanishes, the principle of will ceases to hold
unconditionally, and freedom perishes in a necessity developed into
its concept. Instead of letting the self-contradiction fade away in
an objective, scientific resolution of objective contradictions,
faith lays redoubled weight on sin as sin and drives the
self-contradiction to the extreme. The whole emphasis is laid on the
divine will; only God is free in an eternal, unconditioned, blessed
sense: only the eternally free God can and will set man free ---
yet, be it noted, on the condition, the inviolable stipulation of the
Spirit, that man both can and will, in faith in
% ---- printed p.62 (PDF 75) ----
\opage{62} God, do God's will. This man can now do in faith, and yet cannot;
this man now wills in faith, and yet does not will it; this
self-contradiction man understands in faith, and yet does not
understand it: faith as faith testifies as loudly to the latter as to
the former. If, then, the self-contradiction, so understood, is
posited by a revelation, it follows that it can also be resolved
only by a revelation.
% print: "Sevmodsigelsen" (dropped l) for "Selvmodsigelsen".

What is characteristic of the second stage of revelation is that a
new man enters into the relation to God; this, then, is the
inviolable fundamental condition under which alone the
self-contradiction can be resolved. In the eyes of reason the
condition looks like an enormous contradiction. The new man is not an
abstract ideal but a man in flesh and blood, an actual man. And yet
this man is to differ from all other actual men in this, that he does
God's will in each and every thing, that he is from birth pure and
without sin; but an actual man without sin, what is that but an empty
notion? According to the rational explanation, sin belongs so
necessarily to the nature of man that only a fancy inflamed by ideal
longings can form for itself an exemplar in whom there is no sin, and
take this abstraction of the imagination for a sinless, an actual
man. It is evident here how little reason is able to grasp the
mystery of sin. If sin necessarily belongs to actual human nature,
then by that very fact it must cease to be sin; it can then at most
be only an imperfection. The imperfection necessary by nature can,
under given conditions, become the occasion of transgression, guilt,
crime; but sin in the proper sense it cannot be. For what is sin in
the proper sense, if it is not the consciousness of the innermost
contradiction of the self, the incomprehensible self-contradiction
that man both can and yet cannot do God's will?
% ---- printed p.63 (PDF 76) ----
\opage{63} That sin must be a sickness in man's nature is recognizable from
its being a sickness in the will.

By virtue of the principle of will, faith now reasons thus: if the
sick human nature can exist, then the healthy must be able to all the
more; if the old man is actual, the new must be so still more.

The new man, who in all things does God's will, is the object of
faith, the exemplar of believers, but not himself a believer: this
too gives offense. That the healthy one, it is said, is an exemplar
for the sick is explicable, inasmuch as everyone who feels himself
sick must reasonably wish to become like the healthy one; but that
the exemplar for believers is not himself a believer, how can that
be? We see it in the following way. There is in faith an
approximation to absolute personal certainty, but the approximation
is infinite; the limit, perfect assurance, faith can never reach, for
when the limit is reached, faith is abolished. If, then, it lies in
the nature of faith to strive toward sublating itself, it will be
understandable from this that the exemplar for believers is not a
believer but one who knows --- not one who knows scientifically, but
one who knows personally. In order to be able to do God's will
unconditionally, the exemplar must know God's will, must in every
moment of life and on every occasion have an absolutely infallible
certainty of what it is God wills; but such an infallible certainty
points to a peculiar relation to God.

In the believer's relation to God, sin is an outstanding account
between them; the nearer a man comes to God, the more immediately he
feels sin. The consciousness of sin, which is a sign of the drawing
near, is then at the same time an expression of the distance; the
believer who becomes conscious of God's holy presence becomes by that
very fact conscious of his distance, i.e.\ of his infinite difference
from God in will and essence: it is sin that sets discord between
% ---- printed p.64 (PDF 77) ----
\opage{64} man's self-consciousness and his God-consciousness. From the new
man's relation to God sin is excluded; herein lies that the
exemplar's human will must essentially and fundamentally be one with
God's will. The exemplar's consciousness of himself is identical with
his God-consciousness in the sense that he knows personally his
essence to be the revelation of God's essence, his will the
revelation of God's will: sin is taken away, and the relation to God
posited in absolutely original purity. God is now revealed to men not
merely through a man but as a man, as this man; not merely through
the Word but as the Word itself, the Word that was in the beginning.

The perfect man in divine self-certainty is absolutely different from
all other men. Just as through sin a yawning abyss is set between God
and men in general, so now, after the difference between the sick and
the healthy one, between the sinful and the sinless one, has become
manifest, a yawning abyss is in particular set between believers and
the exemplar. The separation is here all the harsher, as it is now
placed within human nature itself. That the exemplar is actual is
seen from this, that he actually sacrifices himself, suffers, is
crucified and dies; that the actual exemplar, by doing God's will,
reveals God himself is seen from this, that, having suffered for the
truth and given his life for righteousness, he rises from the dead
and enters into glory. The believer cannot without self-contradiction
give up such an exemplar, but neither can he without
self-contradiction think of imitating such an exemplar: how is this
contradiction to be resolved? The healthy one must heal the sick, the
sinless one must come to the help of the sinful: revelation must, if
the principle of will is not to be given up and the whole end in
meaninglessness, assume a new fundamental form.

% ---- printed p.65 (PDF 78) ----
\opage{65} Only at its third stage is revelation completed; its distinguishing
mark is completion itself: in completion it shows what it originally
is, what it was from the beginning and will continue to be, a
proclamation of will in the Word, a communication of freedom from
spirit to spirit. How inwardly the Word of faith, the Spirit of faith,
and faith itself are united in the principle may be briefly
elucidated.

At the first stage of revelation the Word is still one-sided: it is a
word from God to men about divine and human things; but the Word from
God is not yet itself God, the Word to men not yet itself man. In a
corresponding way the Spirit is one-sided. It is the Spirit from God
that awakens the spirit in man; but the Spirit from God, ``the power
of God'', is not itself God; the Spirit over man works in man, but
does not yet allow itself to be grasped as man's proper self. Moved
by the Spirit, man believes in God's Word, cognizes his sin, and
acknowledges that God is in the right. Forsaken by the Spirit, man is
seized by doubt and, in despair at not being able to come to terms
with the God of the Spirit, rebels against God. Then the Spirit
returns, chastises man and brings him once more to cognize his sin.
The movement from sin through sin to sin is essentially the circular
movement in which the life of faith repentantly maintains itself;
this circular movement takes up breakthroughs of hopeful bright
moments, blessed understanding, tried faith and fulfilled promise as
intermediate states, through which the educating Spirit of wrath and
mercy leads its believers to an ever deeper cognition of sin. At the
second stage the Word is present as God and actual as man. In this
actual divine-human presence the Spirit is personally one with the
Word. But the Word has become flesh; the inward of revelation
coincides immediately, for the believing view, with its outward. He
in whom ``all the fullness of the Godhead dwells bodily'' stands, for
% ---- printed p.66 (PDF 79) ----
\opage{66} those who behold him and listen to his words, as the one who alone
has the Spirit abiding in himself. The Spirit is in believers as if
it were outside them; if it were not in them, they could not believe,
and if it did not outwardly come to meet them, they could not make
the exemplar the object of faith. But the Spirit as spirit cannot
immediately be an object: the outward must be transfigured into its
inward; therefore the exemplar, the Word revealed in the flesh, goes
away, that the transfiguration may come. In the outward revelation
the Spirit was personally one with the Word; in the inward, therefore,
the Spirit comes personally to believers by means of the Word. The
coming of the Spirit is not an immanent metaphysical alteration of
consciousness, but a divine act of will, a holy fact, whereby the
outward miracle converts itself into the inward: the Spirit comes as
``Holy Spirit'', as God himself, to stir and work in believing souls.
In the new man the Spirit was personal; the new man went away in the
outward, but in order to come again in the inward. Now the resolution
of the self-contradiction can be accomplished. The self-contradiction
is that the believing man is a double man, a sensuous one who cannot
believe and a spiritual one who lives in faith. So long as the self
is, so to speak, ambiguously divided between the two, so that it is
continually drawn to both sides, man in his hovering unity of
consciousness is a ``wretched man''. Now, on the other hand, after
the Spirit has come with the new man, and the unity of consciousness
has been expressly posited in the spiritual self, the self that,
though human, yet becomes and is transformed into unity with the
divine, in the measure that the believer himself becomes and is
transformed into likeness with the exemplar, the self-contradiction
will vanish.

From these features it can provisionally be seen that the content of
faith develops in and with faith.
% ---- printed p.67 (PDF 80) ----
\opage{67} \runhead{b) The Confession of Faith.}

The spirit in science, the knowing spirit, demands cognition but not
confession; the witnessing Spirit of revelation, on the other hand,
requires both cognition and confession. In laying emphasis on
testimony and confession we lay precisely the emphasis of freedom on
the personal. Confession too shows how categorically the outward
corresponds to the inward. When the exemplar, or, to put it now
plainly historically, when Christ had gone away and the Spirit had
been poured out upon the apostles, they began to speak to the people
about what had happened. Their speech was at once a testimony and a
confession; a testimony inasmuch as they personally --- even at the
risk of their lives --- explained and confirmed what they had heard
and seen; a confession inasmuch as the confirmation was not placed in
objectively universally valid proofs, but in a voluntary utterance,
prepared for scorn and opposition, of what stirred within them. The
apostles did not testify of themselves but of Christ; the community
whose founding they prepared was not to believe in them but, through
their word, in Christ, was with them to confess Christ. But to
confess Christ is impossible without knowing something about him and
acknowledging him as something: for cognition in faith a confession of
faith is expressly required.

The word in the Confession of Faith is not merely a word saying that a
revelation has taken place, but the word of revelation itself, the
word on which the life of the congregation is built, the word in which
Christ still continues to reveal himself personally to each
individual in particular who will enter the congregation.

The Confession of Faith presents what has been revealed in pure
fundamental features; every article, every word has decisive
determinacy and acts with divine authority; the
% ---- printed p.68 (PDF 81) ----
\opage{68} relation between Christ and the believer which is thereby founded is,
in an eternal sense, a relation of selfhood, a relation of self to
self. The question: whence is the Confession of Faith, and who
prescribed it? is a question of life for the congregation. The
question is historical; but there is a difference between the history
of revelation and all other history: historical-critical inquiries
into the origin of the Confession of Faith acquire religious
significance only insofar as they are guided by the principle of
faith. In this connection too it holds that historical actuality, seen
in the light of revelation, is actuality of spirit, actuality for the
spirit.

If the matter is taken outwardly, it must seem a spiritless quarrel to
want to wrangle over the historical origin of the Apostolic
Confession of Faith: whether it was composed by the Apostolic Fathers
or by the apostles or expressly prescribed by Christ himself; and
equally spiritless is it to idolize an old formula to such a degree
that, with a care as scrupulous as though it were a matter of
salvation, one searches and searches in order to find out by written
testimonies whether in all its several articles and parts, down to jot
and letter, in all congregations and all translations through all
centuries, it has remained unaltered. But why, indeed, is the quarrel
spiritless, unless because one overlooks the Spirit, and overlooks
the Spirit because one overlooks the principle.

In consequence of the principle the question must run: what does he
believe and confess who believes and confesses that it is the
teachers of the Church --- the Apostolic Fathers or the oldest
Fathers of the Church --- who with the authority of the Spirit
composed the original articles of faith and prescribed them to the
congregation for inviolable acceptance? And the answer runs: he
believes and confesses that the priests of the Church are the
spiritual lords of the congregation, and that Roman Catholicism is the
only consistent ecclesiastical system. Further: what does he believe
and confess who believes and confesses that the
% ---- printed p.69 (PDF 82) ----
\opage{69} apostles, with the infallibility which inspiration exclusively
assured to them, gave the Confession of Faith that bears their name its
literal expression and thereby founded the first congregations? The
answer runs: he believes and confesses that the Spirit itself has
raised a dividing wall, an insurmountable barrier between Christ and
his congregation; for he believes and confesses that the apostles, by
their exclusive inspiration, stand as mediators and middlemen between
the congregation and Christ, so that the way to Christ goes, not
through the Word of the present Christ himself, but, on account of his
temporary absence, solely through the apostolic word. Finally: what
does he believe and confess who believes and confesses that it is
Christ himself who, as well through the apostles as through the
spokesmen of the congregation in all subsequent times, utters the
words of the Confession of Faith, so that in invisible yet living
presence he reveals himself, not mediately and impersonally through
the words of others, but immediately and personally through his own
Word, the Word with which, in his own divine-human personality, he
speaks to each one in particular who will hear him, invites each one
who will come to him, accept his terms and make a covenant with him?
The answer runs: he believes and confesses what, by virtue of the
eternal principle of faith and confession, alone can be believed and
confessed.

With the different points of view regarding the originality of the
Confession of Faith there naturally follows a correspondingly
different conception of the form and content of the confession. The
\emph{Catholic} view, which puts the priestly word on a level with the
apostolic word, can --- even if, following the legend, one has each of
the twelve apostles insert his article and so produce the confession
pieced together from twelve articles --- regard the apostolic articles
only as the first dogmatic draft, to which the whole series of dogmatic
determinations and symbols owed to the infallibility of the councils
and the popes attaches itself with equally indisputable
% ---- printed p.70 (PDF 83) ----
\opage{70} authority. The \emph{Protestant} view, which, as regards
infallibility, either breaks off with the apostolic age or, for the
sake of catholicity, concedes to the oldest ecumenical councils, as
also to the age of the Reformation, a symbol-forming authority, but
otherwise makes Holy Scripture the source, norm and touchstone
(\textit{lapis lydius}), must on the whole regard all confessions of
faith and symbols as compendia of the doctrine of Scripture, so that
only the word of the Bible and its right interpretation properly
furnish the infallible, sufficiently comprehensive foundation on which
the congregation can build its faith. Common to pope-believing
Catholicism and Bible-believing Protestantism is the fundamental view
that Christianity is a factual, juridical, objective doctrine, and
faith the subjective acceptance of this doctrine.

If the Confession of Faith is considered from the point of view of an
objectively juridical contract of doctrine, then there are certainly
essential objections to be made against both content and form. The
Apostolic Confession of Faith is a far too small dogmatics. Why keep
exclusively to just these expressions, these formulas, when it is
after all well known that many chief doctrines are not even touched
upon in this extract? ``As the Apostolic Symbol, historically
considered, is a post-apostolic product, so too, by its whole inner
character, it is insufficient to be the highest \emph{critical} norm in
the Church. Every one of its words would be unintelligible if we did
not have a richer source to resort to in order to explain them. Hence
we find, too, that the Fathers of the first three centuries never
detached tradition from Scripture, and Irenaeus, so often appealed to
for the rule of faith, himself calls Scripture \textit{columna et
fundamentum ecclesiæ}. It is clear, too, that without Scripture we
would be but poorly supported by the Apostolic Symbol. Although it is
the baptismal symbol, it gives us not the slightest information about
the sacramental significance of baptism; and with the complete
% ---- printed p.71 (PDF 84) ----
\opage{71} confession of the Apostolic Symbol a conception of baptism can very
well be combined which sees in this sacrament a merely symbolic
ceremony. Just as little does it give any information about the
Lord's Supper. The same holds of the important doctrine of
justification by faith, a doctrine whose fundamental significance
surely no one among us has the courage to dispute. Even the doctrine
of the person of Christ is set forth so indeterminately that both
Arians and Socinians have been able to profess the words of the
Symbol, as indeed the Socinians continually appealed to their
agreement with the Apostolic Symbol in order thereby to legitimate
themselves as good Christians''.\footnote{Den christelige Dogmatik af
Dr.\ H.\ Martensen. S.\ 43--44.}

But the way of viewing the matter changes when the principle of faith
is held fast and carried through. As regards, first, the assertion
that the Apostolic Symbol, historically considered, is supposed to be
a post-apostolic product, it will, by the nature of the case, be just
as impossible to produce a positive historical proof of this as it
will be impossible, from written testimonies, to produce a
historically positive proof that the Symbol is word for word from
Christ. The word of the Confession of Faith is what Grundtvig calls
``a living word'', the word of life ``from the Lord's own mouth''. Oral
tradition ``is the way of the Word, a living stream from Christ
through the apostles to the congregation''. To demand written
testimonies for what is to be sought only in oral tradition is a
contradiction; and to bring proof against an oral tradition by
adducing the lack of written documents is likewise a contradiction.
Both contradictions coincide in the one: \emph{to want to build one's
religious faith on objective historical knowledge}%
\footnote{The impossibility of ``building one's blessedness'' on
objective certainty about anything whatever objectively historical
S.\ Kierkegaard has strikingly shown in the first part of
``Afsluttende uvidenskabelig Efterskrift'', where he treats the
objective problem of the truth of Christianity. ``Those who build on
Scripture have in this respect no greater security than those who
build on the Confession of Faith''.
% print: no opening quotation mark before "Naar Skriften betragtes" (closing mark after "subjectivt" is printed).
``When Scripture is regarded as the secure stronghold that decides
what is Christian and what is not, then the task is to secure
Scripture historically-critically. One then treats here of the several
writings' belonging to the canon, their authenticity, integrity, the
authors' axiopistia, and one places a dogmatic guarantee in
inspiration. When one sees the English at work on the Tunnel: the
enormous exertion of force, and how then a small accident can for a
long time disturb the whole --- then one gets a fitting notion of this
whole critical undertaking. What time, what diligence, what glorious
talents, what distinguished learning have here from generation to
generation been requisitioned for this marvel. And yet here a little
dialectical doubt that touches the presuppositions can suddenly, for a
long time, disturb the whole, disturb the underground way to
Christianity that one has wanted to construct objectively and
scientifically, instead of letting the problem come to be as it is:
subjectively''. \dots\ ``This'', he concedes, ``Grundtvig had rightly
seen, that the Bible could not possibly hold out against the
intruding doubt''; but then he does Grundtvig an injustice when,
misunderstanding the Grundtvigian talk of the historical, he supposes
that Grundtvig too wants to have objective security beforehand, before
he can believe subjectively. The historical of which Grundtvig speaks
is not an outward objective to which a subjective can then without
further ado be added, but the living unity of subjective and
objective, thus, in the language of reflection, precisely the unity
of principle. Characteristic in this respect, too, is that while
S.\ Kierkegaard reproaches Grundtvig for taking the historical as
objective, his other opponents reproach him, on the contrary, for
taking it as subjective. The fact is that S.\ Kierkegaard, as
Martensen has quite rightly felt but from his standpoint could not
counteract, has, for all his dialectic, run aground in subjectivism:
he has, as no one else, brought it to evidence that subjectivity in
the believer is the truth, but, with his gaze incessantly fixed on the
believer, has overlooked that faith has its intellectual principle,
that this principle is one with the principle of revelation itself,
and that the principle of revelation is just as originally objective
as subjective.}. If we seek the vital nerve of the Grundtvigian
% ---- printed p.72 (PDF 85) ----
\opage{72} view, it must unconditionally be placed in this, that faith in the
Confession of Faith coincides in principle with faith in its origin
from the object of faith, thus with faith in the Christ personally
present in the words of the confession.

% ---- printed p.73 (PDF 86) ----
\opage{73} That the Apostolic Symbol, regarded as a sum of objective doctrinal
propositions, cannot be ``critical norm in the Church'' must be
conceded; but that Holy Scripture, regarded as a document of
revelation in the historically objective sense, cannot be critical
norm either must surely also be conceded, in view of the phenomena of
spirit that the controversies of centuries have brought to light. Or
have Arians and Socinians perhaps not been able, for their one-sided
doctrine of the person of Christ, to profess the words of Scripture
just as easily as those of the Symbol? The indispensability of Holy
Scripture for the development of doctrine surely no one will call
into doubt; but the point is that the objectivity of knowledge can be
reached just as little by means of Scripture as by means of the
Symbol. In the feeling that what is objective for faith is
essentially different from what is objective in knowledge, one indeed
also demands expressly that the word of Scripture be taken up with a
believing mind, that in the interpretation of Scripture a Christian
consciousness must necessarily be presupposed. But the believing,
Christian mind, the Christian consciousness, cannot be critical norm
either; for, not to speak of each individual among believers having
his own mind, every community of faith has its Christian
consciousness, and this consciousness is in turn different at
% print: "Bevidtshed" (transposed letters) for "Bevidsthed".
different times. It is of no help to refer the one who seeks a
critical norm from the Symbol to Scripture, from Scripture to
consciousness, from consciousness, by way of tradition and the
impression of typically given differences in existing church
communities, back again to Scripture: for where is the center, the
point of departure, the point of conclusion, where is the corrective?
It is, one says, the total impression, which must be
% ---- printed p.74 (PDF 87) ----
\opage{74} made clear through determinations of the particular, while the
particular is in turn determined by the total: in totality the
subjective coincides with the objective, in the total impression we
have the testimony of the Spirit. But if the whole is now not to go
round in circles before one's eyes, and the total impression become
just as indeterminable as the Spirit, and the Spirit in turn just as
indeterminable as the total impression: then there must be in all
this a principle: the principle is Christ himself, the Word from the
beginning.


From the principle everything must be derived, by it everything must
be determined, to it everything must be led back. The principle is not
a confession of faith, but it bears the Confession of Faith and has in
it its short, clear, pithy expression. The principle is not a holy
Scripture, but it runs through the holy Scripture and furnishes the
critical norm according to which Scripture must be interpreted. The
principle is not a Christian consciousness, neither the individual's
nor the community's; but it moves and determines the Christian
consciousness, and consciousness must bow beneath it and let itself be
corrected by it. Only gradually, as the various communities of faith
learn to direct their faith by the principle of faith, and not the
principle of faith by the diversity of the confessions, will religious
life win self-conscious unity, form and shape.

\runhead{c) The Method of the Philosophy of Religion.}

That the content of the philosophy of religion must essentially be the
content of faith is self-evident; but the content of faith is
originally a content of life, and its forms are peculiar forms of
life; whereas the philosophy of religion is able to develop the
content only through intellectual forms: the question then is how the
intellectual forms relate to the forms of life. This question, which
has already been touched on, if from another side, at the beginning of
the Introduction, is now at its close to be
% ---- printed p.75 (PDF 88) ----
\opage{75} answered by a statement of the method of the philosophy of religion.

The philosophy of religion cannot build immediately on any confession
of faith; for the form of the Confession of Faith is a form of life,
intended for the living congregation; but it can and ought to develop
the concept of faith and thereby state intellectually its categorical
significance for the life of the congregation. In the intellectual
form the religious concepts are, one and all, concepts of confession,
that is to say: concepts whose validity is personally decided by
confession.

The three principal concepts of revelation, \emph{the Father, the Son
and the Spirit}, are not concepts of knowledge but concepts of faith,
immediate concepts of confession. The philosophy of religion, which
finds such concepts of confession present in the congregation, then
has the task of translating them into the intellectual form and, by
virtue of the principle, of developing them to clarity and full
transparency. The clarity to which the religious concepts can be
developed depends methodically on the certainty with which the
separation of the two spheres of cognition, that of faith and that of
knowledge, is maintained, and on the deduction and grounding which is
determined exclusively by the relation to God. God's essence in itself the philosophy of religion cannot clarify, for ``God
dwelleth in the light which no man can approach unto''; but God's
essence reflected in and through the relation to God (the relation to
men) it ought to strive to clarify, for in the revealed relation to
God no ambiguity can be thought.

The relation to God is, as we have already seen, a relation of will;
the philosophy of religion, as philosophy of will, is absolutely
heterogeneous with the philosophy of reason, the philosophy of essence
necessitated by the Idea. But cognition of will does not exclude
cognition of essence. Now since will always presupposes essence, and
essence in the highest sense is will, it is seen from this
that the opposition of the principles must rest exclusively on the
method. Just as the philosophy of essence everywhere
% ---- printed p.76 (PDF 89) ----
\opage{76} proceeds from essence and through everything returns to it, so the
philosophy of will, in all problems, proceeds from will and returns to
will. The same holds of the opposition between conceptual
contradiction and self-contradiction; this opposition likewise gains
its significance solely through the method. Seen \textit{in
abstracto}, any conceptual contradiction can easily be presented as a
self-contradiction, and conversely; but according to the methodical
separation between the two regions, the principle of
self-contradiction becomes decisive for faith, that of conceptual
contradiction for knowledge. Where, then, there is a choice between
one of two things, either to place the corrective in the conceptual
contradiction without heeding the self-contradiction, or conversely,
there the necessity of a methodical orientation shows itself. If a
problem is solved in such a way that the conceptual contradiction,
contradiction in the objective sense, is normative, while the
self-contradiction recedes and loses itself in the Mystery, then the
solution belongs to science; if, conversely, a problem is solved in
such a way that what is normative lies in the self-contradiction,
contradiction in the subjective sense, while the conceptual
contradiction is resolved only by a resolution into the Mystery, then
the solution belongs to religion.

Only in the method can the opposition between rational and
antirational, between direct and inverted grounding, find its
application and its proper significance. We illustrate this by some
examples. There are found in ``the Christian Dogmatics'' numerous
statements of which we may say that in and for themselves they could
be both true and untrue, so that their truth or untruth depends
solely and entirely on the grounding.

It is said of creation: ``The world has a double beginning, a
cosmogonic and a creative, a natural and a supernatural. --- The world
must therefore in every moment of its existence be regarded as
\textit{natura} or as organic self-development, and as
\textit{creatura} or as pro-
% ---- printed p.77 (PDF 90) ----
\opage{77} gressive revelation of the divine \emph{will}, and it is the one only
because it is the other.'' This statement is ambiguous insofar as
there is presupposed in human knowledge a point of unity in which the
two dissimilar concepts, creation and nature, find their grounding
equally. Such a point of unity is not given. Between creation and
nature lies a dark, unfathomable deep; here lies the Mystery. The
concept of creation is originally a concept of faith, an unconditioned
concept of will, and as a concept of will it cannot possibly furnish
the basis for a rational development of the concept of nature. The
concept of nature is originally a concept of knowledge, a concept of
necessity resting in the Idea as Idea, and as a concept of necessity
it cannot possibly furnish the basis for the antirational development
that belongs to the concept of creation.

It is said of the Son: ``As the Mediator between the Father and the
world it belongs to the Son not only to live his life in the Father,
but also in the world. As `the heart of God the Father' he is at the
same time the eternal `heart of the world,' through which the divine
life streams into the creation. As the Logos of the Father he is at
the same time the eternal world-Logos, through which the divine light
streams into the creation. He is the ground and source of all reason
in the creation, be it in man or angel, in Greek or Jew.'' This
statement is ambiguous insofar as it is presupposed that from a
cognition of faith of the Logos, supernatural in God and as God, there
could be given a direct transition, necessary to reason, to a
cognition of knowledge of the natural world-reason immanent in things
themselves. Such a transition is not given. Between the Logos in
revelation and reason in the world lies the Mystery, the great riddle
of the world, in which the actual reciprocal determinations of freedom
and necessity are hidden from the eye of man.

It is said of the Spirit: ``The Holy Spirit is distinguished from all
world-spirits not only by the fact that all world-
% ---- printed p.78 (PDF 91) ----
\opage{78} spirits are shut up under sin, but also by the fact that the
world-spirits, according to their original determination, pursue
wholly other ends. The cosmic spirits, be they the spirits of peoples
or the spirits of the state, be they the spirits of art or the spirits
of science, have only the end of founding and spreading their
\emph{kingdom}, while human \emph{individuals} count with them chiefly
as mere means to this end. The Holy Spirit, on the other hand, will
create a kingdom of sanctified, of blessed \emph{individuals}; its end
is the salvation of the soul.'' This statement certainly gives the
impression of being both clear and true, so long as we consider it in
general without asking about the grounding. But when the question of
the grounding is methodically set in motion: how does the statement
look then? Is it really true --- and it is indeed really true --- that
``the world-spirits,'' to which ``the spirits of science'' also belong,
according to their original, mark well: their original determination,
pursue wholly other ends than ``the Holy Spirit,'' ends so opposed to
those of the Holy Spirit as are the making of human individuals into
means for an objective, impersonal kingdom, and the creating of a
kingdom of sanctified, of blessed individuals, hence a kingdom
personal in the eternal sense: then it is indeed
% print: the quotation opened at „den Hellig-Aand is never closed; the single “ closes only the inner quotation.
hereby stated that ``the Holy Spirit is not merely something different
from, but heterogeneous --- absolutely heterogeneous with the
world-spirits, heterogeneous --- absolutely heterogeneous with `the
spirits of science.'\,'' But when a content of the Spirit is to be
methodically developed, then the proposition holds irrefutably: as the
spirit is, so must the principle of development be; and as the
principle of development is, so must the grounding be.

If it is held fast unconditionally that the relation to God is a
relation of will, then a methodical separation of the philosophy of
reason and the philosophy of religion can never be avoided. We see
this from the doctrine of the Trinity. The development of this
doctrine in the philosophy of religion
% ---- printed p.79 (PDF 92) ----
\opage{79} must correspond exactly to a development of the relation to God as
determined by revelation; by departing from this principle, by wishing
to regard God as God in objective relations, partly inwardly in the
eternal deep, partly outwardly in interaction with nature and the life
of the world, one comes at once outside the boundary of faith, is
dazzled by science and goes astray in speculation. Within the
categorically determined boundary of faith, on the other hand, it will
become apparent that God's relation to man also sheds a light on his
relation to himself, that the relation of will gives a mirroring of
the relation of essence, and that it is precisely a methodical
misreading of the reflection of essence in the revelation of will that
has occasioned so many broodings and confused so many heads.

In the relation to God we have three terms: God in his relation to
man, man in his relation to God, and the world, the actual middle term
between God and man. Now insofar as each of the terms of the relation
can be illuminated by the whole relation, the manifoldly nuanced
relation of will in which God, through the world, relates himself to
man may well be able to cast an illumination back upon God's inner
essence. But the main thing is that the principle of will is and
remains the decisive one.

In the doctrine of the Father, God is conceived as the author of all
existence, as the being on whose will all things are absolutely
dependent. It is the opposition between God in himself and all else
that is not God on which the emphasis falls. Now if this opposition is
considered objectively in general, conceptual propositions arise of
which it is not easy to say whether they have their origin from
religion or from speculation. That God in himself cannot with full
freedom have the world for his opposite unless he thereby at the same
time has himself for his opposite must be granted. The divine
self-consciousness
% ---- printed p.80 (PDF 93) ----
\opage{80} requires not only its own I but also its own Thou, and can be thought
to objectify its world only in such a way that therein it at the same
time objectifies itself. Now if God's objectification of the world has
his own self-objectification as its presupposition, then it follows
quite correctly from this that God's free, conscious unity with
himself must be an inner, infinite unity of opposition, a unity
developed in spirit as spirit. The proposition: \emph{God is spirit},
can to that extent, in a formal respect, very well be developed into
the proposition: \emph{the divine essence is trinitarian}.

But is this, then, the religious Trinity of the Confession of Faith?
To make each of the three moments in the concept of God's
self-objectification into a person by itself, the divine I into the
Father, the divine Thou into the Son, and the divine unity of both
into the Holy Spirit --- what is that but an unclear confusion of
subjective representations with objective concepts, a personification
of universal forms of essence, of which one may with reason say that
the images belong to faith, but the thoughts are borrowed from
knowledge: the voice is Jacob's voice, but the hands are the hands of
Esau.

If the unclarity is to be removed, it must be removed methodically.
Faith in the triune God is of a wholly different origin and has a
wholly different character from the objectively metaphysical doctrine
of the Trinity. The religious concept of the Trinity is a concept of
faith, to be determined expressly by the moments contained in the
revealed relation to God: the principle of faith is itself a principle
of the Trinity. Herein lies this: that no moment in this concept can
gain significance for doctrine if it does not prove to have
significance for life. Could anyone show that a faith in the only true
God, with all the demands, presuppositions and consequences contained
in the essence of faith, were possible without faith in the triune
God, then the doctrine of the Trinity would once and for all have to
be regarded as
% ---- printed p.81 (PDF 94) ----
\opage{81} abolished. But if now the converse holds, that faith in the true God,
when the relation to God is developed on all sides to full clarity, is
and must be one with faith in the triune God: then the doctrine of the
Trinity is typical of the method of the philosophy of religion. To
forestall a confusion of scientific and religious objectivity, we
therefore take the \emph{doctrine} of the Trinity expressly in the
sense of \emph{faith} in the Trinity, and state the division of the
philosophy of religion as: \emph{Faith in the Father, Faith in the
Son} and \emph{Faith in the Spirit}.

\begin{center}\rule{5cm}{0.4pt}\end{center}

% ============================================================
% ---- printed p.82 (PDF 95) ----
\opage{82} \parthead{Faith in the Father.}
\lettersub{A.}{The Essence of the Father.}
\parmark{6}

When the question is of the author of all existence, science and
religion at once turn in different directions: science looks down
into necessity, downward into the ground, and demands of reason that
it comprehend the Absolute; but religion looks up toward freedom and
enjoins faith in the heavenly Father. In order that true understanding
between faith and knowledge may be attained, the conflicting views
must first of all be brought to so complete an expression that the
heterogeneity of the principles can thereby come to light, and the
opposite results show themselves to proceed from opposite
presuppositions. The conflict will concern the interior of the
divine, its immanent moments, and the peculiar conceptual form of the
self-subsistence developed through the moments. Against the impersonal
universal of knowledge, faith then sets the fatherly Self; against
knowledge's dissolution of all determinations of inwardness in the
undifferentiated Absolute, faith sets a living sum of fatherly
attributes; against the negative form of the Absolute, faith sets
positively the personality of the Father. Since, however, the
conflict is waged, not between a knower who denies faith and a
believer who rejects knowledge, but
% ---- printed p.83 (PDF 96) ----
\opage{83} essentially between faith and knowledge, in other words: not between
individuals but between principles, then --- as each of the principles
develops and moves freely in its own sphere, and the boundary lying in
the Mystery is acknowledged from both sides --- the conflict, carried
through, will be able to lead to a reconciliation.

\runhead{a) The Self in the Father.}
\parmark{7}

In practical as well as in theoretical respects a great many
deliberations have been undertaken for and against the possibility of
proving God's existence. He who is to prove that God exists must
surely know beforehand what God is; for to wish to prove of something,
one knows not what it is, that it exists, is thoughtless. Now since
God is not recognizable by anything other than himself, and cannot be
known by himself without being cognized as existing, the cognition of
God is surely immediate, and every attempt to prove his existence must
then seem to be as superfluous as it is meaningless. And yet the
certainty of God's existence is by no means immediate; it is, on the
contrary, infinitely reflected. This will become evident when we
consider that a man can cognize God only in proportion as he cognizes
himself: as the man is, so is his God. The proposition that every
judgment about the divine Self, however objectively it may sound, is
essentially reflected in a judgment about the human self, is a main
proposition which sets a separation between faith and knowledge. We
hold fast to this main proposition in that, as regards the reality of
the representation of God, with reference to the so-called proofs, we
conceive \emph{the divine Self} first from the standpoint of
knowledge, then from that of faith, and determine the category of the
Mystery as the boundary of faith and knowledge.
% ---- printed p.84 (PDF 97) ----
\opage{84} \greekrun{α}{The Divine Self Conceived from the Standpoint of Knowledge.}

If it stands firm that the judgment about the divine Self depends on
the judgment about the human self, then the question of how the
knowing subjectivity conceives and must conceive itself in science is
surely one of the first and most important questions. That in all
objective knowledge the essential thing is to cognize the universal,
the necessary, what is independent of all subjective individuality,
has often been emphasized. But the self that in the form of knowledge
is to cognize the universal must through abstraction, so to speak, die
away from its individuality and make itself one with the universal;
in order to grasp the matter in its purity, it must transform itself
into a mirror for the matter. The contemplating self is then absorbed
into universal thinking; the contemplating, thinking, judging self
becomes, in its contemplating, thinking and judging, lost to itself to
such a degree that it determines nothing, can determine nothing and
dares determine nothing except that to which it is necessarily
determined by the nature of the matter. Lost in the matter, absorbed
in the inquiry, the investigator may even be seized by so complete a
self-forgetfulness that he has at last forgotten both the time and the
place and contingent actuality, just as Archimedes, absorbed in his
problems, had indeed forgotten both the time and the place and
contingent actuality when he said to the Roman soldier: \textit{noli
turbare circulos meos}.

It is the first inescapable demand of knowledge that the knowing self
must transform itself into abstract essentiality, into intellectual
selfhood; selfhood is the reflected unity of the individual and
particular with the universal; only in the universal form of
intellectual selfhood can the self identify itself with the objective
and think objectively. He whose principle of life is to be absolutely
objective, then, knows nothing
% ---- printed p.85 (PDF 98) ----
\opage{85} else than that the self itself is the objective, ``the universal,''
``the free universal being for itself,'' and whatever else the phrases
are called with which empty verbal philosophy can dazzle an abstracted,
abstracting subjectivity.


As it goes with the determination of the knowing self, so also with
the determination of the content of knowledge. Here the question is
not primarily what, in a finite sense, can become the object of
scientific investigation, the many outward or inward facts; here the
question is of the absolute content of knowledge, i.e.\ of the Idea
as Idea. That the Idea, in all its objectivity, is also selfhood and
has in itself the moment of subjectivity, is expressed in the concept
by its being said to be subject-object, the absolute unity of the
subjective and the objective. In immediate objective form the Idea is
nature, in subjective form it is spirit: the absolute identity of the
subjective Idea with the objective is God. The proofs of God's
existence now have this significance: to set the Idea in motion and
to show through what transitions and ``mediations'' the thought of
God gives itself reality.

These proofs, whose number modern philosophy in particular has
reduced to three: \emph{the cosmological}, which infers from ``the
contingent being, which cannot bear itself'' to ``the true being,
necessary in and for itself'' --- the proof \textit{ex contingentia
mundi}; \emph{the teleological}, which proceeds from the
purposiveness of the world and infers its wise author; and finally
\emph{the ontological}, which proceeds from the concept of God itself
and infers its reality --- these proofs, as they are found treated
and developed especially in Hegel,\footnote{Vorlesungen über
die Beweise vom Daseyn Gottes.} serve, each by itself as well as in
systematic connection, to confirm what was said above, that knowledge
determined by the Idea is
% ---- printed p.86 (PDF 99) ----
\opage{86} as little able to raise itself from the world to God as from
necessity to freedom.

The thought in the cosmological proof is this: not because the world
is must we think of God as being, but because no true being belongs
to the world must we seek the true being above the world, in God. To
think finite being is to free it from the form of singularity and
contingency and to comprehend it as a universal being, necessary in
and for itself, which is distinct from that first: as God. That this
necessary being, however, can be separated from the contingent only
by abstraction, negative critique has quite rightly seen. ``What
does this mean other than'' --- observes Strauss\footnote{Troeslære
I. S.\ 334.} --- ``that the world shows itself to ordinary
representation as an aggregate of finite things, contingent to one
another, and then as similar laws; comprehending cognition negates
these things as singularities subsisting for themselves, and rises to
the universal unity which just as much posits them out of itself as
it takes them back again into itself, i.e.\ which relates to them as
substance to accidents''? Thus: God is the substance of the world,
its own indwelling necessity.

With the teleological proof, according to Hegel, substance was to
develop itself into the free concept, and the light of freedom to
rise over necessity; but the freedom that is won through the immanent
dialectic of means and ends is only necessity explained by its own
concept. It is organic nature, the organizing principle, the end of
ends, the life in all that lives, that is conceived as God; but such
a God is indeed precisely one with the world, is the life of the
world, the indwelling necessity of living nature.

As for the ontological proof, finally,
% ---- printed p.87 (PDF 100) ----
\opage{87} its result is plainly enough stated. ``Those who are convinced'',
says Hegel, ``that the absolute essence in thought is not thought
itself, speak constantly of God as a beyond for consciousness, and
of thinking him as having a concept of him whose existence or
actuality is something distinct from that concept; just as when we
think or represent to ourselves an animal or a stone, our concept of
it or our representation is something quite other than the animal
itself''. The result is that God thinks himself only as he is thought
by man, that God is the spirit in all spirits, the concept in all who
comprehend, the thinking in all who think --- not the willing Self,
but only a universal selfhood, the logical necessity of the thinking
being.

The worldview that corresponds to the concept of God developed by
these proofs is, as is well known, pantheism: in pantheism God,
without a self, is only universal selfhood. However much Hegel labors
to get beyond Spinozism and pantheism, in this sense he nevertheless
does not get from the spot. True enough, his concept of God is far
more richly developed than Spinoza's; but the principle is the same;
the development is formally logical; the whole logical development
of the Idea takes place within the bounds of pantheism. As man's self
is, so is his God; the knowing self dissolves itself into the
selfhood of the Idea, for the God of knowledge is the selfhood of the
Idea. But this intellectual dissolution of the self has yet another
consequence. Once the divine Self has been dissolved into the Idea of
knowledge, the Idea immediately coincides with natural actuality and
is lost, in this its immediacy, as Idea. ``Only the actual, the
sensuous, is; the Idea as Idea is not; universal selfhood is the
empty, selfless universal'': this is the watchword of atheism. For
atheism the selfless object, the atom in material immediacy, is the
eternal first, the unoriginated.
% print: the 1869 print has a comma after "det Uoprundne" before the capital "Atheismen".
% ---- printed p.88 (PDF 101) ----
\opage{88} Atheism is a kind of naturally necessary consequence of pantheism, a
dispiriting, but scientific consequence.

``It is'', says Martensen, ``the sorry fame of Young Germany to have
brought the negating consequences of pantheism into a system, for
which they have sought to gain entry among the multitude. Instead of
the intellectual-poetic, logical-mystical worldview of Schelling and
Hegel, we have now of late obtained a quite ordinary \textit{système
de la nature}. Under the hands of this generation the
mystical-poetic dust has gone from the wings of the pantheistic
butterfly; the once so glittering creature now shows itself only in
a prosaic nakedness, or at least only with a death's head on its
wings. For it is now proclaimed without circumlocution: there is no
God; the name `God' is a tedious name by which nothing clear can be
thought.''

\greekrun{β}{The Divine Self Conceived from the Standpoint of Faith.}

There is much talk, far and wide, of the necessity of sublating the
dualism of the soul and securing the unity of human self-consciousness
from the ground up: this talk is empty so long as it is not made
evident whence the dualism springs, how essential it is, and on what
conditions it can be sublated. Can the self that denies God, and the
self that puts the Idea in God's place, and the self that believes in
God, really be said to be three modifications of one and the same
human self, or have we here three kinds of men, each with his own
self, or can the three kinds of self perhaps exist in one and the
same man? The question may perhaps seem strange; but its answer is
necessary to the elucidation of the matter. Man's judgment about God
is essentially a judgment about his own self; only in himself, within
himself and from himself can he know God. The sensuous self denies
God; the
% ---- printed p.89 (PDF 102) ----
\opage{89} rational self dissolves itself into selfhood and transforms God
into Idea: only the spiritual self can believe in God. When three
different kinds of self are spoken of here, it will readily be
understood that the talk here is not of three immediately existing
things, but of three possible points of concentration for
self-consciousness, three centers decisive for the life of the will.

The unity of the human essence is originally a unity of self. All
typical, organic formative drive in all living organisms points to a
forming and working that is subject to the law of selfhood, the
fundamental law of self-development. But self-development is
impossible without a self-determining, however wholly unconscious.
Every self-determining indicates a willing, but such a willing is not
yet will. Only in the ensouled organisms, the animals, does the
fundamental willing indicated in the formative drive begin to take on
the stamp of instinctive will. The higher concentration is so far
from being a consequence of animal sense-consciousness that the
latter, in its peculiarity, is on the contrary a consequence of the
former. By nature, i.e.\ in the first immediacy of his existence, man
certainly has instinctive will in common with the animals, just as he
indeed also has senses, the rudiments of understanding, etc.\ in
common with the animals; but since the principle of soul in man
expressly distinguishes itself from the animal one in that it is,
from the very first beginning, a principle of self-consciousness, so
too the instinctive will in man is from the beginning a will of
self-consciousness, a fundamental will that produces itself by
producing its consciousness. Self-consciousness filled with sensuous
representations has in the instinctive will its selfish center, its
first governing point of unity: to this corresponds the sensuous
self.

By acknowledging the universal and being moved by the Idea, human
life undergoes an essential change, which consists chiefly in this,
that the governing point of unity, so
% ---- printed p.90 (PDF 103) ----
\opage{90} to speak, moves and places itself within the Idea: the instinctive
will does not exactly disappear, but it is taken up into the will of
the Idea, and man appears as a rational self. But the will of the
Idea is ambiguous; the universal fills the self and consumes the
self, dissolves it into selfhood. Only when such a
self-concentration takes place that the governing center is placed
in God, so that man's will becomes the organ of God's will, is man
transformed: he has become a new self, a spiritual self.

Now, that man's consciousness cannot at one and the same time revolve
about three equally fixed centers is understandable; but that one and
the same man can find himself in transitional states in which the
life of the soul divides itself so that now the instinctive will, now
the will of reason, now the will of spirit has the upper hand either
wholly or in part, is a fact. We see it in ``God-consciousness''.

Everywhere, in all lands where revealed religion has exercised any
social influence, there will be a great number of people who fancy
that they really believe in God's existence, and yet this faith of
theirs is not much more than a pious fancy. The open assertion of
atheism: there is no God, they reject with horror and aversion; it
seems to them as though such an assertion offended the dignity of man
and made existence meaningless. Whence does this come? It is the
half-rational, half-spiritual self that feels offended. But do they
not, then, really believe in God's existence? By no means! It is only
so long as the representation of existence can coincide with a hazy
representation of absence that the existence of God hovers before
them. ``God exists, but he is in heaven''; that is to say: God is
existent, but infinitely absent. If, on the other hand, the
representation of the existent God is to be put into force, if his
existence is to mean his inner immediate presence,
% ---- printed p.91 (PDF 104) ----
\opage{91} how many will not then fall away? That the divine Self, the willing
God, is actually immediately near to every single man, that a man in
the midst of the thousand concerns of existence shall not be able to
stir, not have a single feeling, not a thought so hidden, not a
desire so secret, but that the existent and precisely therefore
infinitely present God has at every moment his eye fixed on this
single man, and in every single concern asks in silence whether he
will do his will: such a representation is not to be endured! The
natural self reacts from the natural ground of the mind, and the
natural self has the support of ``sound reason''. Man's half-sensuous
and half-rational self here fuse together into a doubting self, and
this doubting self, otherwise at odds with itself about life's
highest concerns, is now at one with itself on this, that the divine
Self is far too importunate.

What then does religion want? Does it want to root out human nature,
condemn reason, tear up the sensuous and rational self by the roots,
in order to plant a nature-less, unnatural and reason-less self in
its place? Far from it! Such an interpretation is no explanation; it
is a slander. What religion demands is faith, faith in God, and first
of all faith in his existence. The sensuous self cannot believe; and
therefore religion demands that the sensuous shall, by strengthened
will, be mastered by the spiritual. The universally rational self
cannot believe; and therefore religion demands that man shall, in the
unconditionally decisive concentration of the will, set his
spiritual, eternal self higher than the universal selfhood of the
Ideas of reason: faith in God's existence is exclusively conditioned
by the concentration of consciousness in the eternal self.

So then --- it is said --- from the standpoint of faith no proofs at
all are needed for God's existence: God has revealed himself, has
shown himself to man as he
% ---- printed p.92 (PDF 105) ----
\opage{92} is; with this an end must once and for all have been made to the
otherwise endless proving! --- Yet it depends on how one understands
the matter. For although revelation has its immediate just as well as
faith does, this immediate is in its ground infinitely reflected. We
see this from the essence of the Word and of the Spirit.

\emph{Proof from the Word.} When, in the scientific sense, there is
talk of offering proof of something, be it a fact, a doctrine or any
universal truth whatever, the conduct of the proof falls essentially
within the subjectivity of the one who proves. The matter to be
proved is objective, rests in objective being, behaves passively and
is in so far indifferent to its proof: this holds even of the
so-called self-movement, the dialectical movement of the matter
itself, which is especially, and with reason, commended as an
advantage of Hegel's immanent
% print: the 1869 print has "immancnte".
method. It is otherwise when the talk is of God and God's existence.
Here it does not hold that what is to be proved is, in objective
being, indifferent to its proof.
% print: the 1869 print has a comma for the full stop after "sit Beviis".
The thinker in the philosophy of religion cannot treat the divine
Self as one treats an immanent Idea in science; he cannot, by any
dialectic whatever, conjure this Self forth from its eternal
concealment, charm it into an irresistible movement of thought and
place it under the necessity that it must absolutely show itself
openly. It is not man who proves to himself and others that God must
necessarily exist; no, it is God himself, the free God, who proves
his existence to man. The proof is that God speaks to man: the proof
is God's Word. The Self and the Word correspond to each other, reflect
and explain each other: take the Self away, and the Word will vanish;
take the Word away, and the Self will vanish. In his Word, and only
in his Word, has God expressly, unambiguously proved his own
existence. To believe in God's existence is to believe in God's Word;
to doubt God's Word is to doubt God's
% ---- printed p.93 (PDF 106) ----
\opage{93} existence.
% print: the 1869 print has a comma for the full stop after "Tilværelse".
What God ``spoke unto the fathers by the prophets'' is not, in this
connection, the decisive thing; the decisive thing is that God has
spoken.

\emph{Proof from the Spirit.} From the reflected unity of the Word and
the Self it follows that the Word from God can as little be held fast
in objectively existing immediacy as God himself. The human word,
once it has been uttered, has an objective existence independent of
the speaker, an existence which, even if the speaker afterwards
regretted his word, cannot be recalled. But it is not so with the
divine Word; it does not tolerate outward facticity: it dies in
immediacy; its objective existence in Scripture and tradition is
without significance when it is separated from the speaking God. Once
a man has clearly uttered what dwells in him, and in what is thus
uttered has given his self a spiritual presence, that is enough: then
let him keep silent, let the word go its way and do its work. With
the Word of God it is another matter: the divine Word cannot work,
cannot be read, cannot be heard, cannot be grasped or understood
except in and through and with the divine speaking.

As regards the manner of revelation, a distinction must certainly be
made between the outward and the inward, between the audible and the
inaudible Word; but it is a distinction between the form of
immediacy and that of reflection, which in no way entitles us to
separate the Word from the speaking. In audible speaking the Spirit
is Word; in silent speaking the Word is Spirit. If the same Word
that comes to consciousness from without through Scripture and
tradition does not, in the moment of apprehension, come forth in
consciousness from within, so that from both sides it is cognized as
coming from God, it cannot possibly be apprehended as God's Word: the
emphasis falls not on the Word but on the Self that speaks the Word.

% ---- printed p.94 (PDF 107) ----
\opage{94} From this it is seen that the proof of God's existence must be
absolutely heterogeneous with any scientific proof whatever of any
scientific truth whatever. Scientific truth is understood from the
word, and it is in so far indifferent who speaks the word. Once a
doctrine has been proved once and for all, then it is proved, and
the matter is settled: the proposition stands fast, and one can, when
occasion arises, point to the proof. But the proof of God's existence
cannot be concluded in this way; it is a proof of the Spirit; it
lives by the Spirit in the Spirit, and inspiration must, like
respiration, be ceaselessly continued. If we now call the proof of
the Spirit, in its inner infinite continuation, a proof of reflection,
then what is peculiar in the reflection is recognizable in the
infinite reciprocal determination of the outward and the inward, the
audible and the inaudible Word. Without the outward, audible Word,
the Word God spoke of old ``unto the fathers by the prophets'', the
existent God would transform himself into Idea, and the divine Self
be dissolved into metaphysical universal selfhood; without the
inward word of thought, inaudible in itself, the outward Word would be
obscure and dead, God's existence immediate absence, and assurance
lacking. The assurance that there is in truth an outward Word from
God is, in every moment of reflection, inseparable from the assurance
that there is in truth an inward Word from God, and conversely; the
assurance of God's existence is conditioned by self-assurance, an
inner assurance of one's own existence, and conversely.

\emph{The Proof of Faith.} ``I know that God exists, as surely as I
know that I myself exist''; these words can be uttered with so
immediate a strength of conviction that they make an impression as
though the speaker really had a perfect assurance and were, as
regards God's existence, entirely and wholly clear with himself.
Nonetheless the utterance is ambiguous, and that both
% ---- printed p.95 (PDF 108) ----
\opage{95} when the question is of the kind of conviction and when it is of its
degree. F.\ H.\ Jacobi's assurance: ``The watchword of myself and my
reason is not I, but more than I, better than I, an entirely Other,
God --- I am not and will not be, if he is not'', furnishes a telling
example. For reason is not, as Jacobi supposes, a peculiar immediate
sense for the personal God. In order to make this view prevail, he
had to set a divorce even between reason and science and show the
impossibility of bringing the understanding to reason. By this
nothing is accomplished. The immediate certainty of God's existence
is not a certainty of knowledge but of faith; nor is the assurance,
in its immediacy, as decisive as the lyrical-pathetic expression
leads one to suppose. About sensuous objects the one who senses can
say: I am as sure that these things exist as that I myself exist; if
one cannot believe one's own eyes, what then can one believe? The
assurance is immediate.

About universal truths of reason the thinker can say: I am as
assured of the truth of this proposition as of the truth of my own
thinking; he who doubts thought must doubt thinkingly, and precisely
thereby acknowledge thought: the assurance is reflected. Should
anyone be able to endure doubting the reality of sensuous things in
general, he would have to contradict ceaselessly the testimony of the
senses; but so monstrous a contradiction in the sensuous self is
unnatural and intolerable: positive immediacy will soon gain power
over negative reflection, so that the doubt must stop and the senses
hold good with assertoric certainty. In expressions such as: to
believe one's senses, to believe one's own eyes, etc., the word
`believe' is taken in a meaning which is as different from the religious one as
a sensuous thing is different from God. As regards the truths of
reason, a reasoning understanding will easily be able to raise doubt
% ---- printed p.96 (PDF 109) ----
\opage{96} about the validity of thought; a man of sound sense can then,
without refuting the skeptic's objections, well say: I believe, i.e.\
I rely immediately on the reality of thought. But faith in thought is
only a provisional faith, a faith that thought itself will sublate;
for scientific thought is not content with faith, it demands
knowledge. Faith in science is as different from faith in religion
as the fulfilling of a law of nature is different from a miracle.

If, then, it is to hold in principle that a man in faith always has
the same certainty of the divine Self as he has of his own, it must
above all not be forgotten that the human self spoken of here is not
the sensuous, nor the rational, but the spiritual. The spiritual self
has an existential immediacy, in which it resembles the sensuous, but
a reflected essentiality, in which it resembles the rational. Now if
a man could be as immediately sure of his spiritual as of his
sensuous self, he would also have as immediate an assurance of God's
existence as of the world's: the cosmological and the teleological
proof could, on this presupposition, obtain full validity for the
believer too. But that is impossible. For whereas the concentration
into the sensuous self takes place as of itself, the concentration
into the spiritual demands, on the contrary, the very utmost
exertion. Could reflecting subjectivity concentrate itself in the
reflection of faith with the same incontrovertible assurance of its
spiritual self as it can, in the reflection of knowledge, expand
itself in an assurance of selfhood saturated with the content of the
Idea: then the proof of faith, which proceeds from the self-revealing
God, would stand on an equal footing with the ontological proof,
which proceeds from the universal concept of God. But neither can
this parallelism be carried through, unless it is conceived
antithetically;
% ---- printed p.97 (PDF 110) ----
\opage{97} for whereas scientific consciousness
% print: the 1869 print has "Bevidthed".
surrenders itself in self-expansion in order to be filled,
determined and moved by the Idea in its objective universality, the
religious consciousness must on the contrary summon all its strength
of will and collect itself, so that in unconditional concentration
into a willing self it may be touched and inspired by God's will.
The rational self can certainly fix itself in the service of the Idea
so as to live for the good; but that this self-fixing, corresponding
to rational ethics, is heterogeneous with the religious is seen
expressly from this, that the relation to the Idea in which it places
man to himself is so one-sidedly absolute that at its maximum it
excludes God.


The spiritual self, the self that alone can, will and must believe in
God's existence, is the self that cognizes that it owes its existence
to God, the self that understands in faith that, with its hope of life
and of blessedness, it would sink together and vanish like a dream if
God were not. The assurance of God's existence and the assurance of the
absolute validity of the self are in principle one and the same
assurance.

\greekrun{γ}{The Fatherly Self: The Mystery.}

That science begins by proving God's existence and ends by proving that
God's existence does not admit of proof is, of course, quite in order.
But when, instead of seeking the reason in the natural limitation of
human knowledge, one flatly declares the first and most essential
principal proposition of religion to be untrue and misuses science to
prove God's non-existence, the presumption is conspicuous. That God
exists expresses briefly and definitely that the self is absolute, or
that the Absolute is a self. If we gather what has been said above
about selfhood and knowledge into one survey, it becomes evident that
science is as little able to prove the impossibility as the necessity
of an existing God:
% ---- printed p.98 (PDF 111) ----
\opage{98} by proving the necessity, it dissolves the self into selfhood,
transforms God into Idea, and thus proves the necessity of a God who is
not God; by proving the impossibility, it entangles itself in the
contradiction that God, whose existence, as surely as he exists, must
necessarily fall outside knowledge, cannot exist because his existence
does not fall within knowledge. How far the philosophy of reason is
from being able in truth to disprove God's existence is seen from the
ambiguous position of pantheism within science itself. What holds of
rational ethics holds of pantheism: it is a spiritual half-measure.
Pantheism makes the Idea into God; that is both too much and too
little. Too much! For the pure Idea cannot be God. The Idea has in
nature no self-reflection; it cannot think its own thoughts, cannot
comprehend its own working. Its being and working in actual nature is
the being and working of the forces of nature: nature is actual, the
Idea is only a metaphysical abstraction. Materialism is wrong when it
attacks religion, for religion lies outside science; but it is right
when it attacks formal idealism, and with it pantheism, with sharp
weapons, for idealism and pantheism fall within science. Too little!
For if the Idea is to be able to work as Idea in its own form and to
master the material, its divine selfhood must be concentrated into
ontological subjectivity, and then we have a logical theism.

That the Idea, when it is developed on all sides, takes on a
logical-theistic form,\footnote{``The Logic of the Fundamental Ideas''
is laid out on that basis.} is a proof that the conflict between
philosophy and religion is not a conflict between the laws of thought
but between the principles. The divine Self in the ontological
subjectivity of the theism of knowledge is one with, and yet absolutely
different from, the
% ---- printed p.99 (PDF 112) ----
\opage{99} fatherly Self of religious theism that is confessed by the
believers. Where does this leave us? That the divine Self both admits
of acknowledgment and yet does not admit of acknowledgment in science;
that the Self that is comprehended in knowledge and the Self that is
confessed in faith are both the same and yet not the same; that the
same divine Self is thus both homogeneous and yet heterogeneous
--- absolutely heterogeneous! --- with itself: herein there is
certainly contradiction, indeed contradiction upon contradiction, when
we detach the word from the matter. By clinging to bare words the
emissaries of verbal critique could here, no doubt, look forward to a
rich harvest. But the point is that the word cannot be separated from
the matter. From the standpoint of the development of the matter the
contradictions at once take on another aspect. The ontological Self is
the limit of knowledge, the fatherly Self the limit of faith; the fixed
point from which faith goes out and the highest point to which
knowledge can reach here coincide in one point, and the contradictions
signify that this point, which in separating unites two heterogeneous
spheres of cognition, is a limit-point. That the contradictions both
proceed from the limit and dissolve in the limit signifies that the
limit is absolute, that here there is a Mystery.

And what, then, is the Mystery? If we could content ourselves with
abstract determinations of the abstract, general explanations of the
general in all generality, then it would be an easy matter to get the
Mystery explained. We would then need only to turn to Hegel: Hegel,
after all, has already explained the Mystery and dissolved it in the
objective concept. God in the concept is the speculative Idea. ``This
speculative Idea is opposed to \emph{the sensuous}, and then also to
\emph{the understanding}: it is therefore a \emph{secret} for the
sensuous mode of consideration and also for the understanding. For
both it is a \textit{μυστὴριον}, namely with respect to what is the
rational in it. The nature of God is
% ---- printed p.100 (PDF 113) ----
\opage{100} not a secret in the ordinary sense, least of all in the
Christian religion; there God has made himself known, shown what he is,
there he is revealed; but it is a secret for sensuous conceiving and
representing, for the sensuous mode of consideration and for the
understanding.''\footnote{Religionsphil. 2. Thl. S.\ 253.}

If we could content ourselves with a mystical notion of the
inconceivable, the unnameable, a \textit{ὑπερουσιον} that withdraws
itself from every possible determination of thought, then we would soon
be finished too. But with such a thing, indeterminable in itself,
revelation, and with it the relation to God, is dissolved. The
relation to the unnameable would then itself become unnameable; the
Mystery is faith itself and faith the Mystery: this is meaningless.

But if the philosophy of religion is now set the task of determining
what the Mystery is and what is in the Mystery, the contradiction
seems to be insoluble. To determine what the Mystery is is, after all,
precisely to abolish the Mystery. ``Insofar as the Mystery is, it
cannot be determined; insofar as it is determined, it cannot be'': such
a contradiction may well amuse a beginner in dialectic; for it takes
little to amuse children. But now to the matter!

\emph{The Mystery is, in the fatherly Self, the original, unoriginated
unity of essence and will, of necessity and freedom.}

The unity itself in abstract generality can be seen and comprehended
from the standpoint of knowledge as well as from that of faith; in the
general unity there is nothing mysterious. But when the unity is to
have a content --- and it is the development of the content on which
everything depends --- then knowledge begins with essence and faith
with will, knowledge with necessity and faith with freedom; from the
two different points of departure the two methods move in two
converging lines. That the converging lines must
% ---- printed p.101 (PDF 114) ----
\opage{101} have a point of contact is incontestable; but that the point
of contact falls on the other side of the limit of knowledge as well as
of the cognition of faith can be proved: it falls in the Mystery.

Between essence and will, more precisely, between necessity and
freedom, there lies a limit which only the self, the fatherly Self,
alone can posit and remove. Were the limit to be determined from
essence, by developing the grounds of necessity, then the grounding
would have to end by making the determination of will a necessary
determination of essence; the will would in this way become only a
semblance of essence, freedom a reflection of necessity. To ground
freedom and will would then be the same as to explain away both will
and freedom. Were the limit, conversely, to be determined from the
will, by developing the grounds of freedom, then the grounding would
have to end by making the determinations of essence themselves into
determinations of will; but what then becomes of necessity? It
vanishes in freedom, is consumed by freedom: the ground of essence is
hollowed out, and arbitrariness becomes absolute. What good does it do,
then, to appeal to the trivial commonplace that in the divine will
there is no difference between essence and will, in the divine freedom
no difference between necessity and freedom? The philosophy of reason,
which in all its groundings must on principle use grounds of
necessity, oversteps its limit and makes itself ridiculous by
inserting immediate notions of free actions and adding assurances that
many of the divine manifestations of essence are at the same time
manifestations of will, that the will as unconditioned
self-determination necessarily belongs to the higher necessity, and so
forth. % print: comma for full stop after "o. s. v"
In a corresponding way the philosophy of revelation, which in all its
groundings must on principle use grounds of freedom, would overstep
its limit if, instead of contenting itself with the general
% ---- printed p.102 (PDF 115) ----
\opage{102} presupposition that the moments of necessity belong
essentially to freedom, it began to go into particulars and to make
distinctions between what God does out of pure freedom, what he does
out of pure necessity, and what he must do half out of freedom, half
out of necessity. To add grounds of necessity to grounds of freedom is
equally unreasonable in religion and in science; to identify grounds of
necessity with grounds of freedom is equally ruinous for religion and
for science. But when the distinguishing, the inner articulation, of
the grounds of necessity and of freedom is determined solely by the
divine principle, the fatherly Self, and is and must therefore be
equally inaccessible to science and to religion, then there is given
herewith both what the Mystery is and what is in the Mystery.

\emph{The Mystery is, in the fatherly Self, the original, unoriginated
unity of the Idea in the form of necessity and the Idea in the form of
freedom, the original, unoriginated unity of nature and spirit.}

That nature and spirit must presuppose each other ad infinitum is easy
to comprehend as long as there is talk only of the general; it follows
directly from what has been said about the relation of freedom to
necessity. The ground of the all-unity of God and the universe, of
spirituality and corporeality, which speculation has now so often
ploughed through lengthwise and crosswise, was already unfolded in a
naive and vivid manner by Jacob Böhme. ``When man sees the ocean of
air above the earth, he sees only stars and clouds of water, and then
he thinks that there must be another place where the Godhead shows
itself with its heavenly and angelic regiment. He will absolutely
separate the ocean of air with its regiment from the Godhead; for in it
he sees nothing but stars, and the regiment between them is fire, air
and water. So he thinks: this God has created thus out of nothing
according to his purpose; how could God
% ---- printed p.103 (PDF 116) ----
\opage{103} be in that being, or how could it be God himself? He
constantly imagines that it is only a house in which God rules and
dwells with his spirit. There might well be many who would say: what
sort of a God would that be whose body, being and power consisted in
fire, air, water and earth? Behold, you incomprehending man, I will
show you the true ground of the Godhead. If all this being is not God,
then you are not God's image; if there is a foreign God, then you have
no part in him. For you are created from this God and live in him, and
he constantly gives you of himself power, blessing, food and drink;
all your knowledge, too, rests in this God, and when you die, you are
buried in this God.''\footnote{Aurora c.\ 22, 35, 23, 1 fl.\
Jvnfr.\ Strauss. Den christel.\ Troesl.\ I. S.\ 481--82.} It is this
theme that more recent philosophy has worked through and richly varied
in the doctrine of the Absolute; but speculation has not been able to
bring it to an actual unveiling of the ``\textit{Mysterium magnum}.''
The task is to deduce: either spirit from nature or nature from spirit;
but whichever of the two ways one chooses, it ends --- in the Mystery.

The philosophy of reason strikes out on the first way; it wants in
actuality to derive spirit from nature. The All, the Absolute, the Idea
indicates the point of departure. First the whole totality of
determinations of essence is developed to its concept in the logically
absolute Idea. The logical Idea is the unity of essence of nature and
spirit. The universal, which includes within itself all differences,
separates itself from itself, releases its elements of difference and
converts itself into the particular: herewith natural existence is
preformed. The particular, however, is no parceling out; the
difference is posited in order to be sublated again; it is posited as
a difference that is yet no difference, it is posited as sublated.
% ---- printed p.104 (PDF 117) ----
\opage{104} The totality in the particular and the totality in the
universal are the same totality. The universal, which primordially
divides itself into the particular, then sublates the primordial
division again and closes together with itself in the individual:
herein is seen a preformation of spirit. The next thing is that the
unity of nature and spirit preformed in the logical Idea is developed
into actual opposition. This happens through particularity breaking
through, through the separating-out becoming earnest. The kingdom of
nature is the kingdom of otherness, and just for that reason of
externality and particularity. The Idea begins in immediate existence
with scattered elements, moments cast outward, and only by passing
through the various intermediate stages of mechanical, physical and
organic formations can it gradually gather itself out of dispersion
and concentrate itself into individuality in man. In human
self-consciousness the light breaks through; in and with man the Idea
is posited as spirit; yet still only in finite opposition to nature,
thus still only as finite spirit. Nature and finite spirit stand
opposed to each other as particulars; here is a primordial division
which the universal must again sublate. But the universal that
sublates the primordial division has now itself been developed and
articulated through the primordial division. The concrete universal,
which has originally presupposed the opposition of particulars through
which it closes together with itself into the individual, is spirit as
absolute spirit: God. The development of the world, ``regarded
\textit{sub specie æternitatis},'' is thus God's self-development; for
the development of the world is nothing other than the Absolute, the
spirit, which ``eternally beholds itself in the mirror of all the
finite spirits that arise and pass away again.''

If the philosophy of religion chooses the other way, begins with the
priority of spirit and seeks to derive nature from spirit, then other
difficulties tower up. God as spirit is God in the eternal, pure
% ---- printed p.105 (PDF 118) ----
\opage{105} light of his self-knowledge and will. As the spirit
absolutely free in itself, God is \textit{actus purus}, a thinking
that thinks itself, a willing that works itself, a comprehending that
eternally comprehends itself. But this thinking of thinking, working of
working and comprehending of comprehending is altogether empty if in
God's essence there is nothing determinable either to think or to
comprehend. One then supposes in God's being a ground, so that his
thinking and working can get a content; herewith a being is introduced
over which God does not dispose. Just as God has not been able with
free will to make himself God, but must from eternity have found
himself as he is and taken himself as being, so God cannot himself
dispose over the promptings that come from the ground of his essence,
the visions, images and representations that involuntarily rise up
and present themselves before his gaze. God is, humanly speaking,
infinitely gifted; but he finds the gifts --- wonderful gifts! --- in
his interior, and has not arbitrarily endowed himself; in other words:
the divine spirit, from which all nature was to be derived, is itself
dependent on a divine nature.

If one now tries to explain the absolute dependence of the divine
nature on spirit by letting the immediate determinations of nature,
the inner material, if we may so say, of immediacy, be for God as real
possibilities, a stuff of ideality which the will can arbitrarily
shape, from which it can at pleasure choose what it will, repress what
it will, organize and develop what it will: then the last state becomes
worse than the first. For if weight is seriously to be laid on the
reality of the possibilities, then no possibility in nature can be real
without having its immediate being in a totality of previously
existing elements. Thus in outer, actual nature water has its
possibility in oxygen and hydrogen, the plant its determinate
possibility, its germ, in the seed, and so forth. The possibilities in
God's spirit would thus have to
% ---- printed p.106 (PDF 119) ----
\opage{106} presuppose in this way a nature already existing behind
spirit; and yet this nature, with its preformations, was to be derived
from spirit? If one evades the difficulty by urging that the
possibilities are only divine representations, ideal productions of
spirit itself, then these productions are at once reduced to being
represented representations, empty, imaginary imaginations,
natureless possibilities; but from natureless possibilities no nature
can be derived.

Thus the matter stands. From the opposition between actual nature and
finite spirit the philosophy of reason wants to derive the actuality
of absolute spirit; but that is impossible. Absolute spirit is
essence and remains essence; an abstract essence, for all its
concretions belong to finitude; an unfree, metaphysical essence, for
it lacks the eternal Self, in which alone it can centralize itself into
will and freedom! that is, into actual spirit. That negative
philosophy knows, with a certain cleverness, how to veil the
contradiction by pregnant phrases is certain; but that it is not able
to derive spirit from nature is also certain.%
\footnote{The following statement of Strauss (Den christl.\ Troesl.\ I.
S.\ 452 flg.) is characteristic in this respect. ``One will not,'' he
says, ``be able to do other than praise the account which the younger
Fichte sketches of the Hegelian doctrine of God, when he says: Hegel
does not teach that God is mere substance, nor that he is the
colorless indifference of the subjective and the objective \dots\ the
absolute fluctuation of the opposition eternally posited and at once
sublated and reconciled again. He is never exhausted in one of these
self-formations, but has infinite superiority over each of them, which
are thus posited in him as something ideal. God is thus here the
eternal intuition of himself in the other, the infinite creation as
infinite subject-objectivity. As I said, this account must be called
correct, and Hegel himself expressed himself in almost the same way
when he wished to present his doctrine of God in popular form. But all
the more must one wonder when Fichte continues: just because the
Hegelian Absolute is the actuality and unrest of the absolute process,
it still lacks the kernel, the resting center, the inner unmoved
clarity of identity, the eternal bond of the self-grasping unity ---
the calmly resting eye must first be implanted in the progressing
process. Has not Fichte himself said immediately before, of the
Hegelian Absolute, that just in its infinite becoming-other it is one
and self, subjectivity; that in the eternal positing of itself the
opposition was just as eternally sublated, or, in Hegel's words, that
its movement in itself is just as much absolute rest? How then should
the Hegelian Absolute first be able to receive unity, rest, selfhood
from Fichte the Younger, when by his own admission it is already in
possession of all this? To what end the artifice of wanting, although
it has natural eyes, to set some more into it, which, once they are
set in, can after all only be glass eyes?'' --- In order that spirit
should not need glass eyes, ``God's personality must be thought not as
individual personality but as all-personality; instead of ourselves
personifying the Absolute, we must learn to comprehend it as that which
personifies itself in the infinite.''}
% ---- printed p.107 (PDF 120) ----
\opage{107} In order to assert the independence of finite spirit, to
raise it above perishability and assure it of an eternal life,
religion makes the absolutely free spirit, exalted above the kingdom of
nature, into the eternal prius. But if the philosophy of religion,
on this account, were to aim at deducing nature from spirit, it would
have to end by transforming nature into a semblance and making spirit
natureless. The principial unity of nature and spirit is a divine
unity of self. That the unity is, is not the Mystery; but the unity
itself in its peculiarity, the unity of determinations of nature that
can only be in God's nature, as surely as the Absolute itself is God,
and of determinations of spirit that can only be in God's spirit, as
surely as God is the absolute Self, the unity of origin, the
primordially fatherly unity of nature and spirit, is the Mystery.
% ---- printed p.108 (PDF 121) ----
\opage{108} \runhead{b) Fundamental Determinations in the Self of the Father.}
\parmark{8}

That the Mystery lies in the divine Self seems at first glance bound to
make all true cognition of God impossible. How, one asks, should it be
possible to develop any concept of God when that which is to constitute
the pervasive unity and center of the determinations of the concept is
an unknown X? How, indeed, should a concept of the self be separable
from the self's own concept? The self, after all, does not sit hidden
behind its marks of essence, as that Greek painter sat listening behind
his pictures, waiting for the judgments of others; nor is the self a
subject on which one who reflects can hang predicates, as one can hang
titles on a man or clothes on a pole. What the self is in itself, that
it is in its self-manifestations, its essential functions: the thinking
self in its thinking as thinking, the willing self in its willing as
willing, and so forth. The self has only predicates that it gives
itself, and in which it gives itself: if the self is mysterious, the
same holds of the predicates, and conversely.

There is in this consideration a mixture of the true and the false,
which we shall throw light on in what follows. In the first place, we
must then remark that the unity of the self and its functions that is
appealed to is ambiguous. Certainly the self as thinking is one with
its thoughts, as willing one with its acts of will; but the self
concentrated in itself is nonetheless a being different from its
thoughts, manifestations of will, and so forth; if the thinker were
nothing other than a thinking, the one who wills nothing other than
willing, then the thoughts would have to be able to think themselves,
and the manifestations of will to will themselves, as when in objective
% ---- printed p.109 (PDF 122) ----
\opage{109} logic the concept is said to comprehend itself. The next
thing to be emphasized is this: that the unity-in-opposition of the self
in itself and its forms of appearance grounds a distinction between the
inner unity of the predicates and their differences corresponding to
the forms of appearance. In order to clear up the opposition between
the philosophy of reason and the philosophy of revelation on this
point, we first direct attention to the reduction of the divine
attributes which ``modern science'' has carried through with a
peculiar consistency, and then pass on to develop the fundamental
determinations in the divine Self that follow directly from the
principle of faith.


\greekrun{α}{The Attributes of the Absolute: Dissolution of the Divine Properties.}

The scholastic treatment of the doctrine of faith is distinguished by
an acute use of formal logic, insofar as the principle of
contradiction can subsist with a finite combining of representations
and concepts and a corresponding distinguishing between pure identity
and abstract difference. But precisely the one-sidedness of the
understanding with which the principle of contradiction is held fast,
where what is at stake is the development of an absolutely divine
content, brings it about that the principle comes into conflict with
itself on essential points. We see this in the treatment of ``the
properties of God.'' The finite is both to hold its ground against,
and yet to dissolve itself in, the infinite; the properties are to be
both mutually different and yet without difference one in the divine
One; in their mutual difference they are both to hold (realism) and
not to hold (nominalism) for God himself: the contradiction is
glaring.

Already the category of \emph{property}, applied to God, makes the
treatment precarious; for property corresponds to thing, just as thing
to property. The thing that has several
% ---- printed p.110 (PDF 123) ----
\opage{110} properties is itself the existing unity of the properties, and the
properties are the inner constitution of the thing, laid out in
existing difference. If the thing is extended, red, hard, square, etc.,
then extension, color, hardness, figure appear, despite their union,
as different properties; the thing is one, the properties another: the
immediate plurality is resolved into the equally immediate unity. This
immediacy is nonetheless reflected; as soon as the reflection is
noticed, the contradiction sets in. It is the same thing, namely the
thing, that is extended, red, hard, etc.; thus: this extended thing is
red, is the red; this red is extended, is the extended; the red is ---
in the thing --- the same as the extended, and yet --- in the thing ---
different from the extended; the same is and yet is not the same: the
contradiction is in full swing. And yet the natural understanding does
not notice the contradiction; this comes from the fact that the
immediate glance at the thing so easily discovers: in the difference
the unity, in the unity the difference.

But when the talk is of God and the divine properties, there is no
supporting immediacy; identity and difference go to infinity, and the
finite understanding is confounded by the contradiction. We shall now
see how the negative critique, the Straussian in particular, has made
use of the contradictions in the concept of God and dissolved the
properties.

\emph{Attributes of the divine being: omnipresence and eternity.} As
the one unrestricted by space, God is omnipresent. That God's being
everywhere must not be conceived as bodily extension, so that half of
God was in half of the world and the whole God in the whole, is
self-evident. Concerning the manner in which the divine essence is in
all things, it is taught: ``not as a body, nor yet as a spirit, but in
a divine, wholly
% ---- printed p.111 (PDF 124) ----
\opage{111} incomprehensible manner''; to which Strauss remarks: ``The alleged
incomprehensibility is as a rule the ecclesiastical embellishment of a
contradiction. Either --- it might be objected --- the world and every
part of it is outside God, and he is then limited; or, if in its
individual parts it is not outside God, then it is mingled with him
and is itself God, or at any rate God's body.'' If, in order to avoid
so pantheistic a conclusion, one would, with the Socinians, separate
God's \textit{præsentia operativa} from his \textit{præsentia
essentialis}, restrict the latter to heaven and conceive the former as
\textit{actio in distans}, then a new contradiction arises: instead of
an actual, God now receives only a possible omnipresence. The
possibility is that God can act immediately where and when he will;
but the possibility itself is impossible to comprehend. In order to be
in all places, God must be in space, for all places are in space; in
order to be in all places in space, the divine being must be exalted
above space, independent of space, nowhere in space: the contradiction
is that omnipresence becomes a nowhere-presence, and conversely.

The resolution of the contradiction is sought in a causality that
conditions both space itself and everything in space and is absolutely
spaceless. This causality is, according to the speculative view, not
God but the Idea. The Idea ``makes space and everything in space its
presupposition in order to issue from it as spirit; the Idea is
omnipresent in all spaces and in everything that fills space, by
virtue of its immanent dialectic as life and law.''

Exalted above the bounds of time, God is eternal, his being therefore
without beginning and end. To bring out the difference between eternity
and endless time, one lets succession vanish for God, so that he is
thought to unite past and future in an eternal present. Eternity is
wholly present at once, but
% ---- printed p.112 (PDF 125) ----
\opage{112} time is not. But now the contradiction breaks through. For if God's
eternity ``has no parts,'' nothing preceding and nothing following,
then God must of course constantly do what he has once done, and must
already have done what he will do in the future. If he has once
created the world out of nothing, then he will go on continuing to
create the already created world out of nothing. Indeed, he must even
create it and destroy it at once. For if he will one day in time
destroy the world, then he has of course already destroyed it; he
destroyed it when he created it.

The contradiction is resolved when the attribute of eternity, instead
of being referred to a personal God, is referred to substance.
``Eternity in itself, or the eternity of substance, is nothing
unthinkable, but on the contrary a necessary concept: it is the unity
in the difference of the moments of time; when time in thought is
deprived of beginning and end as its represented points of support, it
collapses of itself into the indifference of eternity, as the ground
from which it just as constantly issues without end; eternity and time
are related as substance and its accidents.''

\emph{Attributes of the divine understanding: omniscience and wisdom.}
God's properties are conceived in likeness to the properties of the
human soul, and yet the unlikeness is greater than the likeness. At the
very question of the divine understanding we come up against the
dilemma: ``Are consciousness and will to be different in God as in
human beings, or are they to be the same? It does not satisfy to say,
with the popular dogmaticians, that understanding and will in God are
not different powers but one and the same infinite power, which
expresses itself partly in cognizing, partly in acting according to
that cognizing: this is still entirely the relation between human
faculties of the soul, which likewise
% ---- printed p.113 (PDF 126) ----
\opage{113} issue from one fundamental power into different activities. But
such a difference between thinking and willing and acting is a
relation of finitude, which cannot be assumed to take place in God.''
By extensive acquaintance with the given, the human spirit is
\emph{knowing}; by purposive use of the given, it is \emph{wise}.
Were it possible to distinguish in God between omniscience and wisdom,
then it would also have to be possible to distinguish between
something that God found given and something he has himself produced;
but such a separation does not hold. Were God to be said to be
knowing, then within God's knowledge a difference would have to be
made between the discursive and the intuitive, indeed between the
faculty of cognition and the act of cognition, between an intuiting of
the present, a remembering of the past, and a foreseeing of the
future, between a knowledge of the actual and of the possible; but the
more one strives to make these differences clear to oneself, the
clearer it becomes that they vanish in the undifferentiated One, the
eternal Now, in which the Absolute is. How little wisdom befits God
Spinoza has already shown when he says: ``If God acts for the sake of
an end, then he desires something that he does not have. If God has
once created the world, be it for whatever end it may, then he has
beforehand lacked that which he intended to attain by the creation of
the world.'' Wisdom is moreover halved in that it must leave it to
another property, namely love, to determine the final end, so that it
can reserve to itself alone the determining of the means. ``But the
means is something that is not willed for its own sake but merely for
the sake of the end, something that would not be willed if another
thing, which is willed, could have been attained without it; since,
then, to attain the end by a means is the same as to attain it by
something in which it is not attained, or by not attaining it, it
follows that having to use
% ---- printed p.114 (PDF 127) ----
\opage{114} means for the attainment of ends is a manifest limitation and
finitude.''

\emph{Attributes of the divine will: omnipotence, holiness,
righteousness.} An absolute will is posited in God over against an
absolute understanding; but two such absolute concepts must, whether
they are otherwise represented as falling outside each other or as
coming together into unity, absolutely annihilate each other. What is
will? A striving after something one does not have, thus a
determination of finitude which presupposes a limit, a negation. Were
God now to have a will, it would have to be thought of as ``a craving
that is eternally stilled, a seeking that has always already found, a
desire that is inseparably one with the attainment,'' i.e., as a will
that is no will. The contradiction leaps to the eye when one reflects
on the difference between power and will. God is almighty and can do
everything he wills, but cannot possibly will everything he can do,
e.g., not what is bad. Can God, then, choose between different
possibilities? God could, it is said, have created infinitely more than
he has actually done; for all possibilities cannot possibly become
actual. ``But if it is to be incompatible with God's omnipotence
actually to will and to create everything he can, then he not only
does not will everything he can, but he cannot will everything he can
will, a contradiction to which Spinoza has already drawn attention.''
God is \emph{holy}, that is, ``he wills what he ought, and ought what he
wills.'' The relation between the will and the law is essentially to be
expressed by an ought. When one speaks of God's unchangeable resolve to
will to act in accordance with his perfections, the possibility is
manifestly implied that God could also will and act otherwise.
``God himself could not be called good if evil were not for him too a
metaphysical possibility.'' % print: "methaphysisk"
But a possibility of evil
% ---- printed p.115 (PDF 128) ----
\opage{115} is incompatible with the concept of the Absolute. ``For,'' says
Strauss, ``if evil is nothing other than a being's reflection into
itself and not at the same time into the Absolute, then the
possibility of evil is consequently present only in a being that is
not the Absolute; for in the latter, reflection into itself is of
course at the same time reflection into the Absolute. The case is
similar with that ought; for it enters only for a being that merely
has the Idea of the good as an other without being it itself. Thus the
concepts `the Absolute' and `the Holy' will not let themselves be
united; but he who holds fast to the Absolute dissolves holiness,
which is something only in a being placed in relations; and he who
takes holiness seriously encroaches upon the Idea of absoluteness,
since it is obscured by the faintest shadow of a possibility of being
otherwise than it is.''

God is \emph{righteous}, that is, the divine essence is a
contradiction. Owing to the transitive side of the concept of
holiness, the concept of God's holiness is easily confused with the
concept of God's righteousness, which in turn is usually divided into
legislative (\textit{justitia legislatoria}) and judicial righteousness
(\textit{justitia judicialis}). But the divine holiness, which demands
the good of human beings, and the divine righteousness, which
prescribes to them the laws for what is good, cannot with reason be
kept apart from each other. We thus retain only the judging or
retributive righteousness, and the question becomes: partly, how the
divine essence is determined by this property; partly, how that is
determined which God ordains on the ground of it. As regards the first,
both punitive righteousness in itself and in particular the
manner in which the Old Testament speaks of God's wrath, indignation,
and vengeance
% ---- printed p.116 (PDF 129) ----
\opage{116} have aroused such offense that many Gnostics have even regarded the
righteous God of the Old Testament (\textit{δημιουργὸς δίκαιος}) and
the good God of the New Testament (\textit{θεὸς ἀγαθὸς}) as two
different gods. As regards the second, namely the rewards and
punishments that God ordains according to his righteousness, the
conception of them proves to stand in remarkable agreement with the
naturally advancing course of development of the human spirit. It
begins with a childlike belief in finite retribution, reward and
punishment in the external, and it ends with the rational insight that
God himself cannot be righteous, but that the divine righteousness must
be conceived as ``the causality according to which there is posited in
the intelligent individual a connection between its acting and its
feeling, which neither the opposition of another's acting nor a
contradiction between the outward condition and the mode of action is
able to dissolve. That every individual's outward condition, too, comes
to correspond to his moral acting is a task whose solution humanity
always accomplishes only approximately; but that too is only the
accidental, and the remainder of inequality that always remains cannot
be separated from the concept of the finite.'' (Strauss).

\greekrun{β}{The Significance of the Dissolution for Science.}

If the critique's dissolution of the divine properties had been aimed
exclusively at showing the insufficiency of the formal application of
the principle of contradiction in a comprehending of the Absolute, and
thus at establishing the invalidity of the objective, scholastic
theories, no objection could with reason be made against it. But the
critique is not held to this character: it is not the form of the
development but the reality of what is developed, not the theoretical
concept of God but God himself, whose dissolution is here expressly
aimed at.
% ---- printed p.117 (PDF 130) ----
\opage{117} The critique is evidently a tendentious critique, and the tendency has
a perceptible influence on the method. Far from consistently setting
speculative thinking against unspeculative, the critique avails itself
of two kinds of thinking, a speculative and a ratiocinative, two kinds
of logic, according as it can best use the one or the other.
Everywhere that the Absolute's own metaphysical self-opposition is in
question, it is excellently at home with the unity of identity and
difference and lays the greatest weight on the dialectic of the
determinations of difference; but in the analysis of God's properties
it introduces another standard, makes a formal use of the principle of
contradiction, denies the determinations of difference validity, and
reasons scholastically. The divine understanding and the divine will
are either identical or different; if they are identical, they are not
different; if they are different, they are not identical: this is so
unspeculative a piece of reasoning that it can be applied with just as
good effect to a dissolution of the concept of the Absolute as to a
dissolution of the concept of God. In one place we hear Strauss assure
the younger Fichte that the Hegelian Absolute does not first need to
acquire unity, rest, selfhood, since it is already in possession of
all this; it is merely a matter of ``comprehending rest in motion, the
abiding in change,'' in other words: ``thinking speculatively.'' In
another place, however, where the talk is of God's eternity, the same
author assures us, on the basis of an entirely different logic, that
``the unity of self-consciousness, conditioned by the alternation of
outer objects and inner states, is necessarily itself drawn into
time,'' that ``a self-consciousness that always remains equal to
itself'' can as little be ``an actual self-consciousness as a single
unchanged tone could be heard.'' What now becomes of the rest, the
rest that was to be comprehended in unity with motion? What becomes
of the
% ---- printed p.118 (PDF 131) ----
\opage{118} abiding, selfhood itself, which was to keep itself unchanged in all
changes?

By the arbitrary, uncritical manner in which the critique of
dissolution at its pleasure identifies changes in the self with the
self's own change, it undermines the logical foundation of selfhood,
abandons ideality, and violates the Ideas, first the metaphysical and
then the ethical. Divine knowledge falls away; the Absolute may not, as
an absolute, have any knowledge of itself, since knowledge
categorically falls under relations and belongs to finitude; but how,
we ask, can there then be a true knowledge of the Absolute? One of two
things: either the knowledge that is to have the Absolute as its
content must itself be absolute, and the Absolute consequently be able
to have a knowledge of itself; or, if the form of knowledge is
relative but the content of knowledge absolute, then there is of
course a conflict in knowledge between form and content. If the
eternal essence has no knowledge of itself because there can be no
eternally knowing selfhood, how then does it stand with the eternal
thinking, the eternal understanding, and the eternally working
concepts? A thinking that precedes all thinking, an understanding that
has not understood itself and that no one has yet understood, a system
of world-concepts, active, operative, realized in actual things,
concepts that cannot have their origin in an infinite comprehending,
since the Absolute neither knows of itself nor understands nor
comprehends itself, and still less in any finite comprehending, since
they work independently of and prior to all comprehending individuals
--- such conceptual fantasies from the twilight of speculation endure
no critique; they vanish before science like creatures of the
underworld before the light of day.

If every knowing subjectivity is temporal, and consequently all
knowledge temporal, what then becomes of eternal truth? As little as
truth in itself is a knowledge
% ---- printed p.119 (PDF 132) ----
\opage{119} without being, just as little is it a being without knowledge. To make
being eternal but knowledge temporal, being absolute but knowledge
relative, is to violate and deny the Idea of truth. The critique's
Absolute does not have truth in itself, is not itself truth; for it
does not have knowledge in itself, is not itself knowing. Once the Idea
of truth has been denied from the ground up in metaphysics, it is easy
to foresee what fate awaits such religious-ethical Ideas as goodness,
holiness, and righteousness. Since the selfhood in the Absolute, the
selfhood that was to be comprehended as rest in motion, as the abiding
in change, has become so watered down in abstraction and thinned out
in speculation that it cannot even bear a knowledge, how then should
it be able to bear a willing and concentrate itself in eternal will?
But when all will is relative, and the Idea of the good is an Idea of
will, then it of course follows incontestably that the Idea of the
good is itself without eternal validity. The critique's Absolute is
therefore not the good; it does not have the good in itself, for it
has no will in itself; for the same reason, then, it does not have
evil in itself either: it is as exalted above good as above evil, an
indifferent abstractum that does neither evil nor good.

Under these presuppositions it can be understood that an abrupt
transition can be made from pantheism to atheism, from objective
idealism to materialism. In the ethical sense the pantheist and the
atheist stand at roughly the same level; the latter can be just as
great an ethicist as the former. The pantheist becomes virtuously
indignant when one reminds him that human virtue is too weak if it is
not strengthened by the consciousness of God's holiness, that the
wavering righteousness of human beings cannot do without the point of
support that lies in the thought of God's eternal, retributive
righteousness; the Idea-intoxicated pantheist will hear nothing of a
need for divine
% ---- printed p.120 (PDF 133) ----
\opage{120} assistance. The good is to be done solely for the good's sake ---
strange, since the good has ceased to be absolute. ``Happiness is
identical with virtue and only the other side of it; to be unhappy,
i.e., to feel oneself unhappy, is the true immorality.'' But the
atheist can of course speak just as forcefully, just as urgently, of
the worth of virtue. The atheist too does the good solely for the
good's sake; he too recognizes that happiness is identical with virtue,
etc. The difference between pantheism and atheism is \textit{in
abstracto} merely a logical difference of modification, \textit{in
concreto} a difference of mood dependent on individual dispositions.
The pantheist prefers to conceive existence from a metaphysical, the
atheist from a physical standpoint; ethics acquires a finer,
religious-mystical tinge when it is presented in the light of
pantheism, a bolder and more immediately practical cast when it is
set forth atheistically; that is the whole difference.

Insofar as pantheism and atheism appear with scientific energy and are
phenomena within science, they have scientific justification and can,
by their vigorous one-sidedness, contribute mightily to furthering the
development of the Ideas. For the Ideas will not be cowed; they react
with heightened strength when abstract thinking would rob them of the
self-validity which they owe solely to the eternal, knowing and
willing Godhead: pantheism has its scientific counterweight in logical
theism. But the correctives in science are and must be scientific:
logical theism has no immediate point of support in religious theism.

\greekrun{γ}{Religious Theism: Fundamental Determinations in the Fatherly Self.}

If a smooth, continuous transition from logical to religious theism
were possible, then we would indeed have to regard science and religion
as different in kind, since
% ---- printed p.121 (PDF 134) ----
\opage{121} in religion a weight is laid on the self that does not acquire
significance in science; but there could be no talk of absolutely
heterogeneous principles. But the matter stands otherwise. Let one lay
as strong an emphasis as one likes on the selfhood concentrated in
ontological subjectivity, the ontological concentration will never be
able to give us the unconditionally free self, centered in absolute
will, that is the God of religion and the object of faith. When the
self is posited with divine freedom, everything is changed. Between
essence and will, necessity and freedom, a chasm opens that cuts off
the way from knowledge to faith and from faith to knowledge: the
union is in the Mystery, in the unfathomable depth: every mother's
soul that leaps into the depth perishes. ``And yet God has revealed
himself!'' Yes, by positing the relation to God in faith. When the
relation to God is given, then a possibility conditioned by it, of the
cognition of God, of the development of the concept of God as a
concept of faith, is given at the same time. It is these conceptual
determinations that we here conceive as fundamental determinations in
the fatherly Self.


The objections of critique against the divine attributes, that they
presuppose at one and the same time both a likeness and an unlikeness
between God and man, are meaningless from the standpoint of faith. Nor
can it be said with good reason that the Mystery brings either
obscurity or uncertainty into the cognition of God; for the boundary
between what can and shall be cognized and what neither shall nor can
be cognized is determined categorically by the principle itself. As
regards the mutual connection of the fundamental determinations, it
must be expressly emphasized that they cannot be derived from one
another in the same way as immanent determinations in a concept of
knowledge; they must, each of them, be deduced from the relation to
God; the deduction is conditioned by self-cognition, and the grounding
conversely.

The fundamental question of faith, what we must in the proper sense
% ---- printed p.122 (PDF 135) ----
\opage{122} call its vital question, is the question of \emph{eternal life}. It
is not the concept of eternity in objective metaphysical essentiality,
not a logical insight into the universal relation between selfhood and
eternity, that matters here: here eternity is thought of only insofar
as what is expressly thought of here is the eternal in the self, the
eternity of the self. In the light of this eternity the reciprocal
determination of the divine and the human is seen in inner infinity.
The cognition of faith is not grounded by man first assuring himself
in one objective way or another of God's eternity, the eternity in the
divine Self, and inferring from there to a disposition for eternity in
the human self; for this is the condition of faith, that all cognition
of God is conditioned by self-cognition. Nor is the cognition of faith
grounded by man assuring himself in an exclusively subjective way, in
himself and by himself, of his disposition for eternity, and inferring
from there to God's eternity; for this is the condition of faith, that
true self-cognition rests originally and essentially on the cognition
of God. Not in God's essence by itself, not in man's essence by
itself, but in the relation between the essential will of God and of
man, in the relation to God, must the thought of eternity be grounded.
\emph{God is eternal.} But the fundamental determination of eternity
in the divine Self cannot present itself to faith as a resting
determination of essence; it contains a demand: that what is eternal
in the absolute sense is the Self, that the Self is the eternal, this
\emph{shall} be believed, for it \emph{can}not be known.

Over against eternity in the self there presents itself to objective
contemplation another eternity, that of the world, the selfless
eternity of temporal nature, an eternity for knowledge. The strife
between the perishable and the imperishable which thereby arises in
man's consciousness is well suited to show what it means that the true
thoughts of eternity are essentially thoughts of will. If the thought
% ---- printed p.123 (PDF 136) ----
\opage{123} becomes objective and the will weak, then universal reflection goes
into league with sensuous immediacy, then the improbable becomes the
measure of the impossible, then the eternal is not God, God not
eternal, then the selfless eternity triumphs. In the strife of the
self with the selfless, the believer comes to reflect deeply on his
own impotence, but precisely thereby also comes to believe in God's
omnipotence. \emph{God is omnipotent.} Herewith the actuality of the
world and of nature is neither denied nor reduced to semblance. On the
contrary! Increased weight has, if possible, even been laid on actual
nature, since its existence is, so to speak, sanctioned by omnipotence
itself. Independence is ascribed to the world; but the independence is
derived: the actual world subsists before God by God's power; posited
by this power it is something, but over against this power it is
essentially nothing. Omnipotence is self-power, the Self is
omnipotent: with this presupposition life in the temporal can be lived
for the eternal. Since omnipotence is posited as divine self-power,
hence as ideal, spiritual power, the law of selfhood, the law of
self-development, is in the deepest sense cognized as the fundamental
law of existence: herewith eternity has triumphed. The believer must
indeed know God's omnipotence by his own impotence, but in such a way
that this impotence, which in the inwardness of faith nevertheless has
omnipotence itself as its principle, is transformed into a becoming,
ever-growing self-power.

As a life of the spirit which wrestles with actuality and fights its
way forward through temporality, eternal life is a life in becoming
consciousness. Restriction of consciousness is restriction of spirit;
development of consciousness is expansion and elevation of spirit:
understanding of the conditions of life presupposes self-understanding
and conversely; cognition of the conditions given for development is
conditioned by self-cognition, and conversely. The more living,
inward, energetic self-development is, the more painfully
% ---- printed p.124 (PDF 137) ----
\opage{124} will the darkening of consciousness be felt. The limits of nature show
themselves here from a new side: the dark, unconscious powers stand
threateningly against the becoming, easily darkened
self-consciousness. This contradiction is resolved when the self, with
its finite knowledge, presupposes in faith the eternally Knowing One
as its principle: \emph{God is omniscient.} That God's omniscience
must also for faith be a knowledge of the essence of things, hence an
objective, all-embracing knowledge, is self-evident; but in this
generality the thought is unfruitful. That God knows all and is
acquainted with all things, the past and the future as well as the
present, is in the philosophy of religion an assertoric proposition.
Such a proposition can first set thoughts of doubt in motion when it
is carried over into direct science, where it gives occasion for
critical and dialectical discussions. What is decisive for religion is
that faith in God's omniscience, determined by the relation to God,
expressly aims at the believing self. In states where the individual
finds himself forgotten by others, overlooked, misjudged by others,
the consciousness of God's omniscience comes into force: ``no one is
forgotten, for the Omniscient forgets no one; no one is overlooked and
no one misjudged, for the Omniscient overlooks no one, misjudges no
one''; but chiefly, however, in such crises where man strives with
himself in order to come to self-understanding: to examine oneself by
oneself, to judge oneself and become completely manifest to oneself
is in the finite sense impossible; but to have the Omniscient as
one's co-knower, to examine oneself before God, to judge oneself as
one who is judged by God and to become completely manifest to God, is
in the spiritual, infinite sense possible, for God is omniscient.

The life of faith, with its immediate presuppositions, is infinitely
reflected; the reflected unity of reflection and immediacy is seen in
the relation to God. In the relation to God, God is present. The
presence is immediate and real,
% ---- printed p.125 (PDF 138) ----
\opage{125} for it is God himself, not the concept of God, to which the religious
person stands in relation; but the presence is nonetheless reflected
and ideal, insofar as the divine Self is not seen as a finite object
but is cognized as the principle of the self, the principle of the
believing self. In the relation to God the world is the great,
far-reaching matter outstanding between God and man. Since man's
relation to the world must everywhere be determined by his relation to
God, God's relation to man must itself contain determinations of God's
relation to the world. Immediately it looks as if only the world were
the eternal; but faith sublates the immediacy and sets God, the
Eternal One, himself above the world. Immediately it looks as if
nature were the omnipotent; but faith sublates the immediacy and sets
God, the Omnipotent One himself, above nature. To finite reflection it
appears as though all actual cognizing and knowing must have finite
conditions and belong to finitude, but faith transforms finite
reflection into infinite, and sets God, the Omniscient One himself,
exalted above all finite knowledge. Immediately and to finite
reflection it looks as if the world, and only the world, in all things
in all places were the only thing present, but faith, which sublates
this finite both in immediacy and in reflection, discovers the God
present in the world; for \emph{God is omnipresent.} God's
omnipresence expresses the infinite reciprocal determination of ideal
and real presence which gives the relation to God its ideal reality.
The ideal reality consists in this, that the divine Self is the
principle, and the principle in turn the divine Self. That ontological
subjectivity is everywhere effectively present in its other can, from
the standpoint of knowledge, be made evident in logical theism; but
this metaphysical omnipresence the philosophy of religion is not to
develop. The presence as principle with which we are dealing here is
% ---- printed p.126 (PDF 139) ----
\opage{126} the presence of God as the principle at work in all believing souls:
it is the principle of God itself that is the principle of faith, and
as principle has ideal presence in the human self and, for the self,
in the world. Different from and yet identical with the ideal presence
is the real presence, which shows itself in this, that the principle
which awakens consciousness of God does not awaken a thought of God in
general, but makes the living God himself present in thought. ``Whither
shall I flee from thy presence?'' As little as a man can flee away
from himself, just as little can the believing self either flee away
from, or by any finite change of place find itself removed from, the
Omnipresent One. If the ideal presence is one-sidedly held fast by
itself, then the relation to God loses its energy: faith loses its
immediate presupposition, and the relation to God is transformed by
reflection into a relation to the thought of God. If the presence is
held fast in immediate, unreflected reality, then the Spirit is
offended against. Faith is transformed into superstition;
superstition makes an appointment with God or, what is the same
thing, imagines that God makes an appointment with man at a
particular place, is to be found and worshipped only at some external
holy place or other, be it on Gerizim or in Jerusalem.

% print: "guddommellge" (l for i)
That the divine Self is principle means, in other words, that it is
causality, infinitely free causality, absolute will. The
pantheistic-atheistic critique, which denies the Godhead all immanent
determinations of relation and places the relative outside the
absolute, treats the absolute and causality in a very ambiguous
manner. This becomes especially noticeable in the treatment of the
attributes that correspond to understanding and will, wisdom and
goodness, holiness and righteousness. Although the absolute itself is
too absolute to have either understanding or will, either wisdom or
goodness, it is nonetheless to be the principle of understanding and
will in man,
% ---- printed p.127 (PDF 140) ----
\opage{127} the cause of all the wisdom and goodness, the holiness and
righteousness, that is in the self-consciousness of finite beings. But
when the absolute is thus held fast as a causality to which the
categories that determine its effects cannot be applied, then surely
the effects become so heterogeneous with the cause itself that this
cause, as a Kantian ``Ding an sich'', must fall outside all
comprehension. The understanding is finite, but the principle of the
understanding is infinite; the principle of the understanding cannot
itself be understanding; wisdom, goodness, holiness belong only to
relativity and finitude, but the absolute causality whose effects they
are belongs to infinity: the cause of wisdom is not itself wise, the
cause of goodness not itself good, the cause of holiness not itself
holy, etc.\ And what then is the reality that is thereby violated? It
is the self, the human as well as the divine! Such Ideas as wisdom,
goodness, holiness and righteousness are, in the very strictest sense
of the word, Ideas of the self, Ideas of the will, every determination
in them originally a self-determination, every thought a thought of
the self, a thought of the will.

Nothing more thoughtless can be conceived than a half-ethicizing
metaphysics that would give the categories of the will a religious
tinge by letting the finite self assimilate the absolute and yet have
its existential root in the relative; one might just as well try to
let the square assimilate the circle and yet remain a square. How
should the finite self, bound to sheer relativities, be able to
appropriate the absoluteness of essence! If the assimilation is to
consist in the finite self more and more putting off finitude, being
sublimated into selfhood and being absorbed into the infinity of pure
thought, then surely it ends in a metaphysical suicide. If the
assimilation is to consist in the relative subject, with ever-growing
energy of self-determination, drawing the absolute
% ---- printed p.128 (PDF 141) ----
\opage{128} into itself, so that this at last pervades all feelings, thoughts and
actions: what then has happened? If the assimilation has actually been
carried through, then a metaphysical miracle has happened: then the
absolute has become a self, and the self absolute. But the absolute
must, according to the presupposition, remain heterogeneous with the
relative, and the relative with the absolute: the assimilation cannot
be carried through. Without being able to be carried through, the
assimilation is thus to be striven for; it is a striving toward an
unattainable end: the assimilation itself is this striving. This
striving then perhaps has its bright moments, when the self is near,
quite near to being the absolute, and the absolute near, quite near to
being the religiously ethicizing self, and then again its equally
black moments, when the one striving perceives that the eternal and
the temporal must part, because they can never be united. What
contradictions! The deeper one penetrates into the fundamental
relation between the absolute and the self thus determined, the more
clearly the irreconcilable discord shows itself.

From the standpoint of religion the discord is removed. The divine
causality, the principle of all understanding and will, all wisdom and
goodness, all that is holy and righteous in the finite, is in inner
infinity itself understanding and will, is wisdom and goodness,
holiness and righteousness themselves. Thus determined, the principle
is one with the Self, the fatherly Self, from which all selfhood and
self-development has its origin. This fatherly Self is not an
externally represented subject that keeps itself behind finite
predicates; its predicates are determinations of totality, concepts
rich in content, fundamental forms filled with the Idea. The
totalized unity of the fundamental determinations is not the
difference-less and yet difference-pregnant, grayly ambiguous
absolute, which is capable of nothing and yet is to be capable of
everything; understands nothing and yet is to give understanding of
% ---- printed p.129 (PDF 142) ----
\opage{129} everything, does nothing and yet metaphysically gets the blame for
everything. The identity of the predicates, the unity that comprehends
itself in all concepts, wills itself in all acts of will, is at one
and the same time different from and identical with its own
fundamental determinations: the fatherly Self is personal.

\runhead{c) The Personality of the Father.}
\parmark{9}

The concept of personality is in the absolute sense a concept of
faith, not a concept of knowledge. To judge by everyday experience it
might well seem as if the cognition of personality concerned faith and
knowledge in equal degree. Human individuality, which is to make its
self count \textit{in concreto} as personality, must surely above all
have confidence in itself, strength of will and courage for life; but
self-confidence, courage, resoluteness are categories that recall
faith, immediate and practical assurance. Further, the individuality
that relies on itself must know itself and the conditions of its life.
Self-cognition presupposes and demands cognition of the actual
surrounding world, of the essence of things, the course of events,
etc.\ The completed reflection of self-cognition and cognition of the
world perfects itself in theory and becomes a cognition of knowledge.
Personality knows with itself what it itself is; it will thus not
only be able to comprehend what personality is, but even to cognize
the concept of personality as the highest and richest concept, uniting
all concepts in itself. --- This obscurity ceases when the question of
the unity of personality with the absolute, of the absolute
personality, is to be answered.

It lies in the nature of the case that science, when it is to develop
the concept of personality, must seek its elements in the individuality
of man and lay out the development
% ---- printed p.130 (PDF 143) ----
\opage{130} psychologically; but from this it follows in turn that a concept of
personality which is originally laid out on psychologically restricted
presuppositions cannot possibly be driven beyond the limitation lying
in its outline without being dissolved. The contradiction hidden in
this the elder Fichte has completely exposed. ``You ascribe to God'',
he says,\footnote{Ueber den Grund unseres Glaubens an eine göttliche
Weltregierung. (Fichtes og Niethammers phil.\ Journal VIII, I.
S.\ 16).} ``personality and consciousness. What then do you understand
by personality and consciousness? Surely that which you have found in
yourselves, learned to know in yourselves, and designated by this
name? But that you in no way think, or can think, this to yourselves
without limitation and finitude, the slightest attention to your
construction of this concept can teach you. By ascribing that
predicate to this being you therefore make it into a finite one, into
a being that is like you, and you have not, as you wished, thought
God, but only multiplied yourselves in thought.'' --- Hegel, it is
true, has sought to develop the concept of personality into a concrete
concept of infinity; but the development is ambiguous: according to
its universal essentiality personality is infinite and absolute, but
in actuality finite and relative. ``A new epoch'' --- so it runs ---
``has begun in the world; it now seems that the world-spirit has
succeeded in freeing itself from all alien objective essence and in at
last comprehending itself as absolute spirit. The struggle of finite
self-consciousness with absolute spirit, which appeared to it to stand
outside it, is nearing its end. Finite self-consciousness has ceased
to be finite, and thereby, on the other side, absolute
self-consciousness has attained the actuality which it did not have
before. The whole previous history of the world in general, and the
history of philosophy in particular, present merely this struggle, and
seem to have come
% ---- printed p.131 (PDF 144) ----
\opage{131} to the goal where the absolute self-consciousness, of which it is a
presentation, ceases to be an alien one, where therefore spirit is
actual as spirit.'' From this concept of spirit, which is now for the
future to count as the absolute concept of personality, it shines
forth clearly enough how matters stand with the personal in the
concept of personality.

The unreasonableness of wanting to develop the religious concept of
God's personality metaphysically strikes the eye when one reflects on
such immediate categories of selfhood as ``God's heart'', ``God's
wrath'', ``God's mercy'', etc.\ Such pathological, anthropopathic
fundamental categories belong originally to religion, but they cannot
be developed and grounded in science. The absolute of knowledge,
which is to be the principle of all that lives, is not itself living,
does not have life in itself; it is to express itself in pious
feelings, stir in all hearts, but itself it is without feeling,
without heart, absolutely heartless. The ground that knowledge offers
for the concept of personality is, when the anthropological vanishes,
in itself impersonal; from this fundamental contradiction descend all
the particular contradictions that are adduced as proofs against the
personal God. The fault lies not with science, but with the knower who
demands of science what is against its essence.

As a concept of faith the concept of God's personality must
necessarily presuppose the Mystery. The God of religion is the God of
life, a living God in free inwardness. The personal in God is the
Self; as a living, feeling, seeing, hearing, thinking, willing Self,
this personal element is, in the totality of immediately reflected
moments of selfhood, infinitely concrete. It looks human and is yet
absolutely superhuman; here is a contradiction for knowledge, a
conceptual contradiction. What does religion now do with this
contradiction? It does not defy the contradiction;
% ---- printed p.132 (PDF 145) ----
\opage{132} it resolves it, but on the terms of faith. In the relation to God the
divine must be known by the human: the concepts of revelation are
divine-human. In the divine-human concepts the unity of divine and
human is a unity of opposites; thus, for example, the concept of sin
is determined by the concept of holiness: \emph{sin} is in man, not in
God, and yet the concept must be determined divinely;
\emph{holiness} is God's, not man's, and yet holiness, insofar as it
is to be determined in relation to sin, must be determined
divine-humanly. From this standpoint the anthropopathic categories
that hold for God's personality must be conceived. That the living God
personally has feeling for man, religion expresses by saying that God
has a heart; this is the most inward and richest expression of
fatherliness. Immediately man knows no other heart with personal
feelings than his own: here, then, the beginning is made entirely
humanly. Reflection, however, soon discovers that God cannot possibly
have a heart like a man; the expression God's heart then becomes a
figurative expression. God really has no heart, reflection then says;
but the strife of immediacy with reflection is sublated in faith; for
faith, with all its immediacy, is infinitely reflected. That the
expression is figurative must be granted: God's fatherly heart is an
image of reality, drawn upon the background of the Mystery. The organ
of divine feeling is, like the unity of nature and spirit,
uncognizable, and veils itself from man's gaze in an impenetrable
darkness. Thus: how God can personally feel, can be moved to
compassion and kindled to wrath, of this revelation gives no
enlightenment whatsoever; all questions belonging here are dismissed
in religion as irrelevant to faith. But that God can feel, can be
moved, can personally be stirred for man's sake, that is the
fundamental thought of revelation, the thought of life that breathes
in religion.
% ---- printed p.133 (PDF 146) ----
\opage{133} The assurance is in faith. The more living faith is, the greater is
the warmth of inwardness with which the relation to God is cognized.
The feelings that are awakened in the believing self are in the
personal sense divine-human: to the strong feeling of guilt and
transgression corresponds an equally strong feeling of God's punishing
righteousness and wrath. The strong feelings express the living
personal interaction between the human and the divine Self in
reflected-immediate inwardness. Here the question is not of the
difference between the figurative and the non-figurative; here the
question is of the personal reality of the feelings, of the truth and
strength of inwardness. That the images under which the now stern and
wrathful, now mild and gracious Judge can show himself personally to
the penitent are manifoldly colored and changed in the moved mind is a
matter of course.


If we disregard the individual and reflect on the relation to God in
its generality, the concept of the all-fatherly personality will be
seen to develop along with the great phases of revelation; with
regard to this development we may conceive the ideal of personality
as the ideal of majesty, of love, of blessedness.

\greekrun{α}{The Ideal of Majesty.}

It cannot be sufficiently repeated, stressed, and impressed upon the
mind that the doctrine of religion is from first to last, in the sense
of the will, a doctrine of selfhood, a doctrine of the self; the
relation to God is a relation of selfhood, a relation between the
absolute and the relative, the original and the derived self; its
grounds are grounds of selfhood, its assurance self-assurance. Only
under this presupposition can there be talk of God's personality. The
kernel of personality is the self. If the self is finite, relative,
and perishable, then so is the personality. The finite personality
may well have, in the finite world, a relatively higher worth
% ---- printed p.134 (PDF 147) ----
\opage{134} than other finite things; but unconditioned self-worth it cannot
possibly have, inasmuch as it has no unconditioned self. The self that
in the relation to God is absolutely unconditioned is God himself in
eternal might and majesty; the world in all its greatness is only a
revelation of the sole might of the Eternal, the infinite self-might
of the fatherly Author. As the one who holds the whole world in his
power, the divine Self is infinitely opposed to the existing world,
and in this opposition is raised above the world: God's personality,
in the relation to God the ideal of all personality, is seen in the
light of this its revelation as the ideal of majesty.

As the ideal of majesty, God with his personal sole might does not
stand in the first instance in relation to the individual human being
or to the individual things in the world; for the individual things
and the individual human being vanish before his gaze as a nothing.
``What is man, that thou art mindful of him?'' It is the universe,
the immeasurable whole, that presents itself to God. And yet this is
still only the one side of majesty. If God, as he surveys the world,
were so lost in the contemplation of the whole that for the whole he
forgot the individual, then he would after all be, as it were,
absorbed and impressed by outward greatness, and would not be the
absolutely Exalted One. From the other side, then, majesty shows
itself in this, that the great infinite whole is itself a nothing to
the Almighty. But when the greatest is a nothing just as much as the
least, then God's eye can rest on the least as well; when the total is
something vanishing just as much as the particular and individual,
then the Eternal in might and majesty can be mindful of the particular
just as well as of the universal, of the individual just as well as of
the total. Thus the divine Self fixes its gaze on the human; thus God's
personality becomes the principle of man's personality; thus God is
``mindful of man, makes him a little lower than the angels, crowns
him with glory and honor.''

% ---- printed p.135 (PDF 148) ----
\opage{135} It is not thoughts of knowledge but thoughts of faith, personal
thoughts of impression, that the ideal of majesty will set in motion.
The human consciousness that gets no further than bowing before the
powers of nature is overwhelmed by outward greatness, loses the
impression of itself, intoxicates itself with fantasies about the
infinite, and dreams of liberation in the great nothing. This is the
Indian pantheism of nature, a standpoint where individuality is a
barrier and the self a torment; the thought of personality has not
yet come into being. Different from the oriental pantheism of nature
is the pantheism of spirit in the Occident, a fundamental philosophical
view that is just as pervasive in the present as it was in antiquity,
and can develop just as well out of the dissolution of monotheism as
out of that of polytheism. At this stage the self, by virtue of the
absolute, will be sufficient unto itself; the thought of independence
has been awakened; man cognizes himself as a personality coming into
being, an actual one. But this personality is kernel-less, hollow, and
false; for the self is the kernel of personality, and the pantheistic
self, which can be blown out like a candle, is vain. Only impressions
of the ideal of majesty can in faith awaken the human self and give
its personality uprightness. It is a peculiar stamp that the ideal of
majesty impresses on the relation to God. The believing self has,
superficially seen, a deceptive likeness to the egoistic self; they
resemble each other in all immediacy in this, that both love life and
fear death. But the superficial likeness turns into decisive unlikeness
when we consider that the life the believer loves is not the temporal
but the eternal life, and the death he fears not the natural but the
spiritual death. This must be held fast if we are to understand the
personal thoughts of impression.

The pantheistic self immerses itself in the universal; it develops
personality in the temporal on temporal terms; through this
self-development it wants to be filled with the
% ---- printed p.136 (PDF 149) ----
\opage{136} absolute; with a transitory eternity in the Idea it wants to satisfy
itself and prepare its self-dissolution. In conflict with the great
life of the world, i.e.\ with actual events and laws, it practices
resignation, and can practice resignation on a great scale.
Resignation is this, that the individual will makes itself into the
universal will, so that it essentially wills only what is universally
necessary, universally rational. Since now the objectively universal,
according to the fundamental view, is the selfless impersonal, the
pantheistic self is from the ground up prepared for its dissolution
and reconciled with objective necessity.

It is otherwise with the believing self. It is indeed prepared for
death, natural death, but it is so far from being prepared for an
essential and eternal self-dissolution that the thought of such a
dissolution must appear to it the deepest self-contradiction. In
order to cancel the self-contradiction it then raises its gaze toward
the ideal of personality, and receives an impression that there is in
the self a power that overcomes death. Thus the religious person
develops a disposition to eternal personality. Instead of proceeding
to self-dissolution he proceeds to self-affirmation; but it is a
self-affirmation that must stand its test.

For the sake of personal self-affirmation the believer will see the
finger of God in all dispensations, and in all dangers trust in God's
aid. But by virtue of the ideal of majesty something else befalls him.
God's thoughts are far above man's thoughts, his will infinitely
exalted above man's will. Far from finding either the way as he had
imagined it, or help such as he had pictured it in human fashion, the
believer finds barriers everywhere and the way out blocked. Necessity
reigns; no regard is taken for the freedom of the individual; help
from God fails to come; the one who struggles in
% ---- printed p.137 (PDF 150) ----
\opage{137} powerlessness must help himself as best he can: here all goes dark
before the eye, here the believer must learn to resign himself.

Infinite resignation is at this stage doubly difficult to practice.
The presupposition is that nothing, not even the most insignificant,
happens without God's will, and yet the effects of this will are not
immediately recognizable; the presupposition is that the Almighty can
help in all dangers, and yet even in the extremity of mortal danger
not a trace of personal help is to be discovered. It is as if God in
all his majesty found this present actuality too mean, far too
insignificant to concern himself with. Here resignation is what is
called for. If the believer cannot carry resignation through, then he
is lost, then doubt triumphs; faith turns over into unbelief, and it
ends in despair. But if he can carry resignation through, then he
triumphs doubly; for then he understands that the deepest
self-affirmation can demand that the finite personality succumb
sensuously, that a pressure in the temporal conditions an impression
of eternity from the Eternal. Resignation is then not a compromise
with necessity but a consecration to freedom. The expression of the
fact that the consecration takes place, that the self rises up under
the heavy pressure, is a distinguishing mark by which religious
resignation differs essentially from merely rational resignation.
Rational resignation ends in silence, in mute submission; but
religious resignation, which in its relation to God has the impression
of personal majesty, ends in adoration with thanks and praise.

As the ideal of personal majesty, the Almighty is the object of fear
and reverence. The self that fears the impersonal powers of the world
is cowed and unfree; it gives up its personality by bowing before the
impersonal: the greater the fear, the greater the subjugation. But the
believing self that fears God gains in fear an impression of the ideal
of personality and feels itself set free,
% ---- printed p.138 (PDF 151) ----
\opage{138} liberated from selfish, finite, temporal fear: the more inward the
fear, the greater the liberation.

One human being may well look up with admiration to another, and
recognize that the impressions owed to a noble, great, superior
personality are beneficial. But if the person who, for all his
superiority, is still finite becomes the object of unconditioned
reverence, deification, and adoration, then the adorers have in the
personal idol an essential barrier to their own personal self; the
blinder the deification, the more fatal the barrier: it is slavish and
degrading to place the ideal of personality in something external. It
is otherwise when the human personality bows with holy reverence
before God. More unconditionally than an oriental slave can cast
himself in the dust before his mighty, silent despot, must the
believing self bow before the high ``All-disposing One, Creator of
heaven and earth''; but religious reverence is, in its inner
boundlessness, infinitely different from slavish reverence. In slavish
reverence the unfree self succumbs to an alien, arbitrary power, and
thereby proves that it has its ruling principle outside itself, that it
has lost its personality; but the believing self that in adoration
bows before the Father in heaven gains itself by giving itself up; for
it bows humbly and yet personally free before the Author of its own
personality.

\greekrun{β}{The Ideal of Love.}

Judging by the ideal of majesty, however, the opposition between God
and man is so great that the distance becomes infinite. The unity of
essence that expresses itself in the relation to God, the relation of
selfhood determining from both sides, is so abstract that a personal
understanding must seem unattainable. Man indeed has self-being in
likeness to God; but the self in almightiness and the self in
powerlessness do not agree. The personal self must have a
% ---- printed p.139 (PDF 152) ----
\opage{139} natural foundation; this holds just as much for the divine as for the
human personality. But that God's nature is other than man's shines
forth clearly from the fact that in actual life there is ceaseless
strife between God's will and man's. The relation is always strained:
God and man stand, so to speak, on strained terms with each other. God
is wroth with man, and man trembles before the wrathful God; it is a
painful relation. So in the Old Testament: Jehovah's countenance is
terrible, whoever sees this countenance must die; Jehovah's judgments
are terrible, he visits the iniquities of the fathers upon the
children unto the third and fourth generation, his righteousness is
as a flaming sword, his wrath a consuming fire. The true personal
relation is a relation of love; but between sinful man and the holy
God a relation of love is impossible.

According to the ideal of majesty, considered by itself, man relates
to God as creature to Creator; at this stage the Mystery of love
cannot be revealed: the relation of nature, the holy fundamental
relation, is hidden. God cannot love man when he cannot recognize
himself in man, his own true image, and man cannot love God when he
cannot cognize in God his personal nature, the exemplar of his
essence. With the actual, true man comes understanding. The true man
has human nature, is of human descent, ``born of a woman''; this self
stands, according to its origin, to that extent on an equal level
with the human self in general. But the true man is, in an exclusive
sense, also of divine origin, has as God's Son divine nature: in this
self God truly sees himself, the brightness of his glory, the express
image of his essence (\textit{Alter Ego}): the true man is personally
one with the true God.
% ---- printed p.140 (PDF 153) ----
\opage{140} In the true man, then, God can love human beings, as in his own Son
he originally loves himself; and human beings in turn can learn to
love God, inasmuch as in the Son of Man, who is eternally God's Son,
they can recognize and love the true divine-human Self.

This is the Mystery of all-love as of all-personality: that God the
almighty has as his opposite not merely the world, the selfless other,
and within it man, the powerless self, but has, in eternal
self-opposition, himself as his opposite. If this self-opposition
were only abstract and formal, then God's self-love would be only a
divine self-seeking love, and self-seeking love would thereby be set
up as the supreme principle of personality. But the self-contradiction
in this is as glaring as the blasphemy. For the self-seeking one, who
makes not only things but every self outside him into a means, into an
impersonal other, has precisely in the impersonal other a barrier to
his personality. The Mystery is that what from eternity constitutes
the otherness in the divine self-opposition is not a merely outward
world, for that is too little; not a finitely becoming humanity, for
that is too little. Even an inward world, a world of Ideas in which
the Idea of humanity forms the center, is too little, far too little:
\emph{the self-opposition requires a self.} A divine one! Yes, but not
tautologically, for then the opposition becomes formal and love
egoism. A human one! Yes, but not a finite, powerless one; for in such
a one God cannot cognize himself as God, cannot love himself. This is
the Mystery, that only the divine-human Self, God's other Self, can
express the unity of love in the one personal God: God is love.

To see what lies in this, we need only reflect on the
self-contradiction in which the essence of love is dissolved when
God's personality is denied, and
% ---- printed p.141 (PDF 154) ----
\opage{141} the physiologically limited individuality of human nature is to
constitute the sole foundation of personality. This is not now to be
understood as though the personal were to coincide immediately with
the merely natural. On the contrary! The self that is disposed for free
self-development is spiritual, its nature not merely sensuous but also
spiritual: it has within it a lower and a higher nature, and personal
development is determined precisely by the higher nature coming to
master the lower. The self with its natural spiritual individuality
stands originally in individual relations to other individualities
likewise naturally spiritual; it has within it a double natural
impulse: to hold on to itself and to give itself up; both are united
in love. Thus the love between man and woman is natural and yet
spiritual. It is the self's deepest need to give itself up in another
self, unequal and yet of like kind; for through this self-surrender it
gains itself in the full sense: it gains personality. Those united in
love cognize that their individual union is borne by the personally
universal, that this universal is something divine; they cognize that
Eros is a deity, i.e.\ they cognize the Idea. From the spiritual
natural love in the erotic, personal self-development advances further
to free, conscious love of the Idea: personality is filled and moved
by the concretely universal Ideas of art and science, of the state
and the fatherland. The more vigorously individuality lives for the
universal, the greater is the self-affirmation of its inner individual
independence. But the demands of personality are first satisfied in
the proper sense only when the individual learns to cognize the Idea
of the good as the universally valid Idea of its essence, of its own
better nature, and so learns to love the good. In love, the true,
active love of the good, personality and the personal life then find
their highest consummation. Not good actions singly,
% ---- printed p.142 (PDF 155) ----
\opage{142} but the love of the good expressed in good actions, love of what is in
truth the absolute, consummates personality. Just as man can love the
good inwardly, in spirit and in truth, because the good is one with
his ideal nature, so he can also love righteousness inwardly, in spirit
and in truth, and through his love of the absolute in righteousness
become righteous; can love holiness, and through his love of the
absolute in holiness be sanctified. Only by letting go of God and
grasping the absolute can the human personality be perfected in love.

The hollowness of this whole way of looking at things catches the eye
once the gaze has fixed itself on the essence of love, the
naturally determined spiritual relation of self to self. The goal of
love is the most inward, deepest self-satisfaction. The more inward
the need and the stronger the desire, the richer, warmer, more blessed
the satisfaction. What, then, does the self need? Another self, equal
yet unequal, with a nature of like kind, yet of unlike kind! The more
essential the likeness within the pervading unlikeness, the more
opposed the two natures are within the fundamental unity, the richer
the satisfaction. The temporal self cannot love cold eternity; the
relative subject cannot love the subjectless universal; unless it were
to be a thinned-out love such as that with which the pure thinker
loves the absolute in pure thinking with pure being. Personality can,
in the absolute sense, love only the personal. The universal love of
the Idea is ambiguous. It indeed stands higher than the immediate,
instinct-bound love of an individual, insofar as it is spiritual and
free, reflected in knowledge and will; but it stands lower insofar as
the Idea, this universal selfhood without a self, cannot return love.
Satisfaction in love is conditioned by love returned: only in the inner
in\-%
% ---- printed p.143 (PDF 156) ----
\opage{143}finitude of love returned, of mutual love, is love perfected. From
this can be seen the difference between loving God and loving an Idea.
To love the Idea as Idea is to love the impersonal, whose secret is
that love, without love returned, ends in the dissolution of the self;
for the Idea as Idea and the self as self are, in their opposition, not
of the same nature. But to love God as God is to love the eternally
rich personal Self, the fatherly Author of all personality; it is to
love, in love returned, him who first loved us; it is, in spirit and
in truth, to love oneself.

\greekrun{γ}{The Ideal of Blessedness.}

It will always appear to the consciousness that reflects on the
universal as though it were at bottom more rational to think of God
according to the ideal of majesty than according to that of love. As
love itself, the all-good God is surely the God of the heart, and as
the God of the heart the heavenly Father. But when the heavenly Father
is to be conceived as Lord of all the world, the almighty Creator of
heaven and earth, and thus in truth as the God of the spirit, then
the spirit is cold, and it seems as though the representations of the
God of the spirit and the God of the heart would not let themselves be
united in one concept. The reason for this is, after what we have
elucidated concerning the Mystery, explicable enough. As the unity of
will of might and knowledge, God according to the ideal of majesty is
the Eternally High, the Father of lights, ``with whom is no
variableness, neither shadow of turning.'' In the God of the spirit
not even the remotest hint of natural oppositions is to be discovered.
As the unity of love of Father and Son, on the other hand, the divine
essence is an essence doubled in its holy ground of nature; from this
doubling spring countless oppositions: it is as though the
naturally necessary oppositions overwhelmed the free unity. If it were
now possible to begin with the spirit itself and develop nature from
it, or conversely to be\-%
% ---- printed p.144 (PDF 157) ----
\opage{144}gin with nature in God and from it develop the spirit, then it would
also be possible to develop the ideals of majesty and of love from one
and the same concept; but both are, according to the Mystery, equally
impossible. For the same reason a rich, intuitive development of the
ideal of blessedness is then also impossible. To say in few words what
blessedness is, or rather what is to be understood by the word
blessedness: that by it is to be understood the life perfect in
itself, the life of all-love in its inner, holy light, its eternal
peace and rest, is an easy matter; but what is then gained by this?
In the essence of blessedness the deepest oppositions, the strongest
tensions, are at once both posited and yet equalized. What, humanly
speaking, is a blessed moment? To grasp the moment itself we must
consider what goes before it. It is a need so strong, a thirst so
burning, a longing so inward, a desire so irresistible, that if
satisfaction failed to come the torments of hell would be unbearable:
that the need is sated, the want stilled, and the delight of
refreshment drunk in deep draughts, that is precisely the
blessedness of the moment. This holds not only of the sensuous-psychic
but also of the spiritual. There is a sensuous hunger, but there is
also a spiritual hunger. If there were no spiritual hunger, there
would be no spiritual satiety either, and if spiritual hunger could
not be still more tormenting than sensuous hunger, then satiety would
be dull and flat, and there would be no blessedness in the spirit.

But to make the ideal of blessedness intuitive on such a foundation is
equally impossible in faith and in knowledge. If all tensions are
equalized, all oppositions leveled, then eternal peace comes with
eternal stillness; but this stillness is without movement, in the
infinite peace there is no life, in the boundless uniformity no warmth,
no riches of inwardness. If, conversely, emphasis is laid
% ---- printed p.145 (PDF 158) ----
\opage{145} on the immediate elements of difference in natural oppositions, then
there comes not only life and stir but also, what Jacob Böhme saw in
powerful intuition, ferment and strife and torment and wrath and
anguish and tension into the eternal essence of all-love; but where
then is blessedness?


The ideal of blessedness is categorically an ideal of revelation or, what
expresses precisely the same thing, an ideal of faith. Of what in the human
sense we must call a blessed moment, everyone knows, after all, that it can
be cognized only either by being immediately experienced and thus enjoyed
really, or, if without actually being experienced it is to be enjoyed
ideally in a vivid presentation, that it is then not science but poetry
that is able to do it justice. On the other hand, it now seems to have been
entirely forgotten that the concept of blessedness in the religious sense
cannot be a concept of knowledge, but belongs to faith alone. And yet it is
surely explicable enough that the ideal of blessedness, the ideal of perfect
personal life, is not capable of making an immediate impression on the
merely knowing consciousness --- of metaphysicians who become blessed in the
Absolute we will not speak --- but for the life of faith in struggle and
suffering the ideal has the most decisive significance.

Faith is, as now sounds almost like a figure of speech, a matter of
blessedness. That blessedness is the goal can then also be expressed by
saying that life in God is the goal. The ideal of blessedness thus
coincides with the ideal of God, i.e., with the perfect ideal of
personality. It is again the fundamental relation between nature and
spirit on which everything depends. Personality can as little be perfected
in natureless spirituality as in immediate naturalness; all essential
misjudgments of the ideal of blessedness accordingly also show themselves as
practical deviations from the ideal of personality. We see it both in
moralism and in mysticism. ``Two
% ---- printed p.146 (PDF 159) ----
\opage{146} things,'' says Kant, ``fill me with the highest admiration: the starry
heavens above my head and the moral law in my breast.'' In that the moral
law is referred to the lawgiver, the ideal of majesty is made to count, but
in an entirely one-sided way. How one-sided the conception is can be seen
from the fact that the moral law, the law of freedom, of personality, is
set in so original a conflict with the promptings of nature, the higher as
well as the lower, that the universal must in principle consume the
individual; the ``happiness'' which is to be attained in a life to come,
when God's justice at last evens out the disproportion between worth and
fate, is owed to an arrangement which concerns personality itself only
insofar as, according to the principle, it must necessarily end by
annihilating both duty and virtue. To an opposite extreme the mystics have
gone, especially those who, in an overflowing feeling of the unity of
essence of God and man, have quite forgotten the ideal of majesty over that
of love; in the mystical relation of love man has in this way become as
indispensable to God as God to man; here the well-known verse fits:
``Ich weiss, dass ohne mich Gott nicht ein Nu kann leben; Werd ich % print: "ein Nun"
zu nicht, er muss von Noth den Geist aufgeben.'' In this twilight of
inwardness, nature and spirit have then so completely dissolved each other
that the light of freedom, and with it of personality, has gone out both
for God and for man.

In order to grasp the reciprocal relation between love and blessedness, we
must lay weight on freedom. The goal of personality is freedom; but man's
freedom demands self-development, and self-development demands struggle.
While reason and humanity will always secure for themselves the highest
thing of freedom by setting the Idea above the struggling self, religion
reverses the relation and sets the struggling self above the Idea. By
setting the self above the Idea, religion by no means wishes to flatter
human self-love; on the contrary! it wishes precisely to
% ---- printed p.147 (PDF 160) ----
\opage{147} make an end of self-love; it wishes to purge egoism out from the root.
In that it lays decisive weight on the eternity of the self and sets up the
blessed God as exemplar, it can rightly demand a thoroughgoing sacrifice of
the temporal. Self-sacrifice through renunciation, pain and death is then
an actual testing of the freedom of the self, a confirmation of its
personal, indelible independence, a preparation for blessedness. ``Be ye
perfect, even as your Father which is in heaven is perfect''; the perfect
life of personality is blessedness: be ye blessed!

But the struggle of temporality stands sharply against the peace of
eternity, natural life against spiritual. If this opposition is now
conceived in such a way that the natural is without further ado identified
with the sinful, while natureless spirituality is made one with the holy,
then we have asceticism. The ascetic, in blind zeal for a natureless
spirituality, will break the staff over nature as nature. The rich
immediacy of existence, what delights the eye and refreshes the mind and
gladdens the heart, yes, the heart itself, if it were possible, he would
gladly see uprooted, extirpated, annihilated, for the sake of blessedness;
but what sort of blessedness is that? The natureless spirit never becomes
blessed. And wherein lies the sinful? It is in the strictest sense not
nature that has made the will sinful; it is conversely the will that has
made nature sinful. Spirit as spirit is not the enemy of nature; were
spirit the enemy of nature, then there would be no hope of blessedness, % print: "Saligbedshaab" (b for h)
then the Father in heaven would himself not be blessed.

Just as blessedness presupposes love, so it also presupposes personal
independence: love and independence are united in the Spirit. This again
throws light on the relation of love to freedom. If love is considered
one-sidedly by itself, then the opposed personalities stand, as it were,
on the point of dissolving and melting together in immediate devotion; if
free\-%
% ---- printed p.148 (PDF 161) ----
\opage{148}dom is considered one-sidedly by itself, then independence becomes so
hard that the independent ones, like monads, close themselves off from one
another. The intensified opposition is a hint that in freedom itself there
must be a need for a strong uniting middle. If each of the independent ones
is free and strong by personal will, then the higher power that is to unite
them must also be free and strong by personal will: this power is the
Spirit. As merely universal, even if rich in Ideas, selfhood, the Spirit
cannot possibly unite self with self when each self, in its wealth of
inwardness, is unconditionally independent. If the particular wills have a
stronger self than the will of the Spirit, then the opposition will also
be stronger than the unity: each individual will then retain a surplus of
isolating self-will, a doubtful remnant of Idea-egoism, which brings
freedom into conflict with love and disturbs blessedness. Weakness in the
Spirit is inseparably one with a weakness in freedom, an enfeeblement in
personality. The power of the Spirit is in the will; its blessedness is to
found a kingdom of personally opposed wills upon the difference of freedom
in the unity of love. For the will of the Father and of the Son are in the
Spirit one and the same will, and precisely thereby two wills opposed in
love itself: they could not be one true free will if they were not two,
and they could not be two if they were not one in the Spirit's own
absolute self-will.

That God, according to the ideal of majesty, is raised above all
oppositions means in the language of religion that he masters all
oppositions; but that he masters all oppositions means in turn that with
freedom in love he enters into all the oppositions, those of temporality
and sinfulness as well; that he, in other words, ``lives a double life, a
life in himself in undimmed peace and self-sufficiency, and a life in and
with his creature, in which he submits himself to the conditions of
finitude, yes, lets his power be re\-%
% ---- printed p.149 (PDF 162) ----
\opage{149}stricted by man's sinful will. In this life of God together with his
creature, the biblical concepts of a divine sorrow, wrath and other
determinations which manifestly place a restriction on the divine
blessedness find their application. But this restriction is sublated in
his inner life of perfection, wholly independent of the creature, and in
the victory-assured foresight of the fulfillment of his counsels. In the
outer chambers --- we therefore say with the old theosophers --- there is
sorrow, but in the inner chambers there is nothing but
joy.''\footnote{H.\ Martensen. Den christelige Dogmatik. S.\ 105.}

\lettersub{B.}{The Works of the Father.}
\parmark{10}

To speak directly scientifically of the works of the Father is in itself
just as contradictory and meaningless as to speak directly scientifically
either of the Trinity or of blessedness or of the personality of the
Father. But just as rational science must be said to be within its good
right when it declines to treat the concepts of faith, so the philosophy of
religion, the inverted, antirational science, must be said to be within its
good right when, in consequence of the principle of faith, it expressly
aims at putting the categories of revelation into force. The philosophy of
religion can to that extent at its ease, without any particular exertion,
continue the fight against a marrowless, immature metaphysics which has
not yet learned to heed the limit of knowledge and acknowledge the Mystery;
but in its endeavor to keep the principles of religion and of science
apart, the inverted science must take good care that it
% ---- printed p.150 (PDF 163) ----
\opage{150} does not itself, in the struggle against the encroachment of scientific
thinking, come into conflict with the laws of thought.

There is a law of thought that forbids placing determinations of finitude
in the Absolute. From this it has then been derived that, since God is the
Absolute and determinations of finitude cannot be transferred to God, God
cannot be said to act or to do definite works either; particular actions,
definite works are determinations of finitude: from this standpoint, then,
it would be a pure self-contradiction to speak of God's works.

What is logically groundless in this objection is easy to make clear.
Without determinations of finitude there could be no determinations of
infinity at all: every determination of the infinite is a reflex of a
corresponding determination of finitude sublated in reflection. That God
stands in a relation of infinity to the world thus by no means excludes
the moments of finitude corresponding to this relation, but on the
contrary takes them up into itself. So much as regards the abstractly
logical. If we next lay stress on the fact that the God of religion is not
a selfless absolutum, but the absolute Self in the unity of will of
knowledge and power, of power and wisdom, then the determinations of
finitude contained in the relation of infinity will expressly correspond to
such categories of personality as action and work. The personality that
acts works, after all, out of its inner totality and concentrates its
essence in the single, definite act of will. God's works reflect at one and
the same time the inner, infinite totality in God himself and the outer
infinite, and precisely thereby finite, totality in the world. Such works
of totality are: \emph{creation, preservation and governance}%
\footnote{The formal reciprocal determination of finite and infinite can
best of all be made clear by a philosophical elucidation of what in
mathematics is called the analysis of the infinite. (Cf.\ ``Mathematik og
Dialektik.'') But once insight into this reciprocal determination has been
won, it can in turn contribute from many sides to clearing up the relation
between necessity and freedom, nature and spirit, the world and God. The
interconnection of the world (the expression of necessity) is infinite;
the working of the divine will (the expression of freedom) is infinite. But
if one then wishes to conclude from this that necessity and freedom are,
without any difference, one in the Absolute, one concludes in philosophy
just as falsely as if in mathematics one assumed without further ado
$\frac{\infty}{\infty} = 1$ and $\infty - \infty = 0$. He who, from the
circumstance that we are not able to demonstrate the reciprocal action
between natural causes and God's will in the particular case, would
conclude that there is no reciprocal action, that the divine causality of
will must consequently drop out, concludes in philosophy just as rashly as
he who, from the indeterminacy of the forms $\frac{0}{0}$, $\infty \cdot
0$, $0^{0}$, $\infty - \infty$, would conclude that these indeterminate
forms cannot be transitional forms of determinate functions and as such
have determinate values. In all natural effects there are differential
effects of divine will; but it is one thing to acknowledge the cooperation
of the preserving and governing will in general, another to be able to
determine the functions.\par
The scholastic perplexities with \textit{concursus}, \textit{gubernatio}
and \textit{conservatio} spring from two mistakes: 1) unknown functions are
treated as known, and 2) it is forgotten that the continuous transition
from finite natural effects to finitely determined effects of freedom takes
place through differential effects.}.
% ---- printed p.151 (PDF 164) ----
\opage{151} \runhead{a) Creation.}
\parmark{11}

The worldview according to which alone there can be talk of creation is
the one taken up into the relation to God itself, the religious one.
Herein lies the fact that there can and must be another view of the world,
falling outside the relation to God, namely the empirical-rational one,
according to which there is not in the proper sense talk of creation, but
only of nature. Nature is a concept of knowledge, creation a concept of
faith: between nature and creation lies the Mystery. An
% ---- printed p.152 (PDF 165) ----
\opage{152} objection such as this --- that in religion there is also talk of nature,
both of the nature of the world and of God, and in science of creation, in
geology, e.g., of periods of creation, and that the concept of nature must
consequently just as well be capable of being called a concept of faith as
the concept of creation a concept of knowledge --- is without significance.
What matters here is not words and designations, but the methodical
development and grounding of the content of the real concepts. In what now
follows we shall treat first the scientifically naturalistic view, then
the immediately religious view, and finally investigate how matters stand
with the question of principle on both sides.

\greekrun{α}{The Scientifically Naturalistic View: Cosmogony.}

From nothing comes nothing: this universally logical proposition receives
in the doctrine of nature a confirmation of its own through the
rational-empirical consideration of the coming to be and passing away of
things. Anaxagoras already pointed out to the Hellenes that they wrongly
assumed a becoming. Things form themselves by combinations of eternal
fundamental parts, and they are dissolved again when the fundamental parts
are separated. The old doctrine of eternal fundamental parts, atoms, has
found recognition and further development in modern science: the doctrine
of the eternity of atoms, i.e., of matter, must now presumably be regarded
almost as an axiom in the natural sciences. With this the doctrine of
natural forces and causes stands in the closest connection. The atom acts;
matter has force; force and matter cannot be separated: the material
forces have their possibility in the properties of matter. From this it
follows that everything produced in nature must have natural causes; the
realm of nature is an infinite system of presuppositions, conditions and
consequences; possibilities, dispositions, real grounds are, in their
manifestations and modes of appearance, subject to unchangeable laws:
everything that
% ---- printed p.153 (PDF 166) ----
\opage{153} happens happens solely through the fulfillment of the laws; this holds
just as inflexibly of the formations in organic as in inorganic nature.

In organic nature the formations do indeed assume a peculiar stamp of
self-formation, and philosophy has to that extent reason to characterize
the cause of life as an Idea-cause, a self-cause working in matter. But
while it is still unclear to the physicist how something that in itself
is as immaterial as an Idea-cause must be can act upon matter and subject
to itself the elementary material causes, it is on the other hand both
clear and certain to him that the living organism rests on a peculiar
cooperation of a system of physical forces. That plan and type show
themselves in organic dispositions and formations must be granted; the
living organism presents a totality of means and ends. But the
teleological points at no point beyond necessity; here in this ingenious
order there is no thought of consciousness, in the many works and
mysterious workshops no trace of any free master-workman; no spirit and
will, no fatherly Self does any work here: motherly nature does everything
from blind and yet never erring drive, with wonderfully sure instinct.
Only when every notion of the cooperation of a free, creating will has
been removed will one understand that the world cannot have any finite
beginning in time, that its natural self-formation can, by filling out an
eternal time, unite time with eternity. Natural development proceeds on
the whole slowly: no intermediate link may here be skipped, no secondary
causes, however insignificant they may seem, be lacking with their
conditions. The developmental history of the terrestrial globe must be
reckoned in billions, and the formation of our planetary system perhaps in
billions times billions of years. An almighty Creator could well have
finished all this in less than six days; but a miraculously accelerated
production
% ---- printed p.154 (PDF 167) ----
\opage{154} would then also have reduced the whole natural development to
semblance. How should natural seed be able to come forth from created,
i.e., not natural, plants, or natural offspring from created, i.e., not
natural, animals? The natural must have a natural origin: to wish to
derive the natural world from a non-natural cause is a contradiction.

We have in these few strokes indicated the character of the naturalistic
extreme, and shall now briefly give a characterization of the
supranaturalistic or biblically religious one.

\greekrun{β}{The Biblically Religious View: Creation in Six Days.}

``In the beginning God created the heaven and the earth''; herewith it is
declared: partly that the production of the world is an absolute act of
freedom, an actual work of God; partly that the coming into being of the
finitely determined world is to be referred to a definite point of
beginning lying in time. It is the finite determinateness of the
representation, the immediately intuitable, on which weight is laid here.
``And the earth was without form, and void; and darkness was upon the face
of the deep. And the Spirit of God moved upon the face of the waters.'' In
an entirely childlike, naive manner the waste and empty earth is here set,
as a foundation already produced by the creation and yet natural,
elemental, externally over against the Creator, while the externality in
this opposition is again sublated by the fact that the Spirit, the
substantial unity of the almighty thoughts and forces, moves upon the
waters. ``And God said, Let there be light: and there was light. And God
saw the light, that it was good.'' This presentation is characteristic in
its particulars. The almighty, fatherly Self reveals itself in this work
acting with free inwardness; it appears human and is yet in essence
divine. That the creating will must work in unity with thought and the
word is, as surely as the act of creation is to be a
% ---- printed p.155 (PDF 168) ----
\opage{155} true act of freedom, understandable. First, then, light exists in
God's inner being, according to its concept in thought through the word;
next the word passes from the inner as a word of almighty power over into
outer actuality: the actual light, the realized concept, is also a word of
God, a word of power in the external. Finally --- so finitely are the
moments distinguished --- God looks and sees that what has been produced
actually and completely corresponds to what was thought; herewith the
beginning has taken place: the morning of creation has dawned.

In agreement with the act of creation here described all the following
ones take place. Step by step the work advances, and it is completed with
the creation of man. Since man, created in God's image, is the crown of
creation and destined to rule over the earth, it is a matter of course
that the earth is regarded as the center of the universe, and that sun,
moon and stars are created to shine over the earth. This conception of the
structure of the world and its purpose, certainly very restricted compared
with the widened horizon of the present, is from the standpoint of
religion nonetheless perfectly correct insofar as it is essentially
determined by the relation to God: the cognition of the world and the
cognition of God intervene in each other in an infinite manner. But no
further than that.

From a Bible-believing standpoint much has been made of the fact that the
order in which the work of creation advances agrees on the whole with the
natural sequence of stages for the lower and higher, elemental and organic
productions, and it has thus been sought to vindicate historical truth for
the narrative. But already Celsus found it unreasonable that the Mosaic
history of creation spoke of days before the sun was created, so that the
interpreters had to help themselves as best they could, either by assuming
an alternating expansion and contraction of the light, which before its
concentration into firmly bounded luminous bodies was still an ether, or
by taking refuge in allegories. --- The description of Adam's creation,
that he is formed of clay (red earth), % print: "(rød Jord,)" (comma inside the parenthesis)
% ---- printed p.156 (PDF 169) ----
\opage{156} that Jehovah breathes the breath of life into his nostrils, etc., cannot
give any picture belonging to natural history, any more than the later
narrative of the creation of woman from the man's rib. It must indeed be
granted that the doctrine of nature is not itself able to set up any
probable hypothesis about the appearance of man on earth; but it may for
all that be fully entitled to reject a mode of presentation that is
incompatible with its scientific character. In order to defend themselves
against the objections of critique on the ground that God, according to
Gen.\ 1:27, creates man and woman at the same time, but according to
2:7, 21 ff.\ first the man alone and later the woman, the apologists had
either to assume that the first account was merely a summary statement of
what is depicted more circumstantially in what follows, or, with
Augustine, to distinguish between the production of man according to
possibility and his actual appearance. This shows sufficiently that the
literal interpretation cannot be carried through.

That the historical here does not belong to the essence of faith is seen
from the fact that the particulars, represented one by one, cannot
possibly acquire any influence on determining the relation to God more
closely. The Creator does not become more almighty because he uses only
six days to create the world in than if he had used six thousand years;
the decisive thing is that he creates. The Creator does not become more
free, not more loving and good, because he creates the man first and the
woman afterwards, than if he had created them both at once; the decisive
thing is that God created man in his image, that he ``created them male
and female.''

Yet one misgiving still remains. ``Whence do we know that God has created
the world, except from God's word? But God's word about creation is,
after all, precisely inseparable from the history of creation. If we
cannot believe the history as actual, then neither can we believe in God's
word;
% ---- printed p.157 (PDF 170) ----
\opage{157} if we can believe in God's word, then we can also believe in the
historical actuality of the account.'' To meet misgivings of this kind we
need only recall the distinction developed in the Introduction between
profane and sacred history. The history of creation is surely not to be
conceived as a piece of supernatural natural history: for if it really
were to be so, then we should ask in vain for the testimony of the Spirit.
All natural history is profane history, just as much as general world
history; concerning the visible things of natural history the Holy Spirit
gives no testimony. But the history of creation is history of revelation;
the actuality of the history of revelation is actuality of the Spirit,
actuality for the Spirit.


\greekrun{γ}{Creation and Cosmogony: The Mystery.}

That the existing world is a work of God's unconditionally free
omnipotence is the fundamental thought of revelation, the kernel of the
whole doctrine of creation, the one thing that can be believed and must
unconditionally be believed. The assurance that God has actually created
this world does not depend on the accounts in the first book of Moses;
it stands and falls with the relation to God, it lies in the essence of
faith. Here again it holds: the improbable must be believed, but the
impossible cannot be believed. Could science in truth prove the
impossibility that the world actual in nature can have been created by
God, then faith and knowledge could not be united in one consciousness;
then it would be impossible for the one who knows to believe. But
knowledge with its cosmogonic concept of nature is just as little able
to disprove creation as faith with its religious concept of creation is
able to infringe upon the concept of nature. From this standpoint we
continue the negotiations between religion and science.

As regards, in the first place, science, it is not physics --- physics
is within its good right --- but metaphysics, the one-sided and false
metaphysics, that
% ---- printed p.158 (PDF 171) ----
\opage{158} has done its part to undermine faith in creation. It is
chiefly three categories, \emph{power}, \emph{matter} and
\emph{disposition}, about which the investigation turns.

\emph{Power.} In his presentation of ``the religion of majesty''
Hegel\footnote{Philos.\ d.\ Relig.\ II Theil. S.\ 55--56.}
has apparently asserted and acknowledged the creative power in its
unity with goodness. ``God can in the true sense be creator only as
infinite subjectivity; thus he is free \dots\ only what is free can
have its determinations over against itself as something free. \dots\
This going-forth from one another, whose totality is the world, this
being is goodness. But the being of the world is only the being of
\emph{power}, or the positive actuality and independence of the world
is not its \emph{own} independence, but \emph{the independence of
power}.'' --- One would now suppose that the independence of power must
be sought in the divine Self exalted above the world, that is, in the
free, willing God himself; but no! ``The world must \dots\
% print: two-point ellipsis ". ." in the 1869 print.
with regard to power be represented as something \emph{broken} in
itself: the one side is the manifoldness of difference, the infinite
richness of existence; the other side is then the substantiality of
the world; but this does not belong to the world itself, it is
\emph{the identity of essence} with itself.'' This essence, whose
identity with itself is nevertheless not supposed to be the world
itself, is, when we look more closely, none the less the world's own
essence, the form of its own being existing for itself. The pantheistic
mark is that the Self, the creating Self, ``the infinite subjectivity,''
is posited only in order to be extinguished again, comes only in an
intimation only to vanish again into the universally necessary
``identity of essence with itself.''

The ambiguity is removed in logical theism\footnote{Cf.\
Grundideernes Logik I og II.}. Here the ontological subjectivity, the
substantial
% ---- printed p.159 (PDF 172) ----
\opage{159} selfhood, is reflected into itself in such a way that the
infinite power can be seen to be centralized into self-power. Just as
otherness stands against selfhood, so power-of-another stands against
self-power. There is in the essence of power a doubling which must be
rightly impressed upon us. Actual power must have an opposite,
% print: "Modsætnig" (for Modsætning).
something that it can master: the weaker that is which is mastered,
the weaker is the power; the stronger that which is mastered, the
greater is the power. Already from this one sees how matters stand
with God's omnipotence and the actuality of the world. It is not the
case, as a metaphysics enervated by abstractions supposes, that the
divine self-power would make the independence of the world into
semblance; for if God's power had only a world of semblance to master,
it would sink down to a mere power of semblance. But the fact is that
the relative independence of the world is posited by the absolute
self-power as the latter's own immanent opposite. Instead of standing
as an external resistance, a limit to omnipotence, the actual
independence of the world over against God expresses, on the contrary,
its absolute dependence on God. The identity of this dependence and
this independence gives us the concept of God's infinite
over-power: God's infinite over-power is one with his unconditioned
self-power.

\emph{Matter.} The most essential hindrance to assuming a creation
would presumably lie in the doctrine of the acknowledged eternity of
matter. Could science actually prove that the atoms with their present
properties are in themselves absolutely independent, then faith in
creation would undeniably have to be given up. But what science would
conduct such a proof? Natural science, which is now properly the one
that holds to the atoms and makes the atomic theory fruitful, cannot
conduct the proof. What does the empirical doctrine know about the
atoms? It knows that they are the ultimate limits which scientific
analyses have so far reached; it recognizes that these limits are
movable,
% ---- printed p.160 (PDF 173) ----
\opage{160} that substances formerly regarded as indivisible have since
proved to be divisible; it shows the atoms, not in hard, absolute
closedness, but in relations that amply demonstrate their mutual
infinite relativity. It is not the material point in its material
immediacy that is the unchangeable; its unchangeableness lies in its
substantial character as a point of function; through their
interaction the discrete points of function become moments of bodily
continua.\footnote{Cf.\ Forelæsninger over ``Philosophisk
% print: "Propædentik" (for Propædeutik).
Propædeutik'' fra Universitetsaaret 1860--61. S.\ 143--177.} The
rapidly advancing science of chemistry is precisely in the course of
leading the expert, through an atomic theory distinctly developed in
all directions, back to the old assumption of the infinite formability
of matter, of its being according to possibility; but, to be sure, in
such a way that this formability conditioned by discreteness, this
law-determined being according to possibility, is essentially different
from the Aristotelian. No, it is not physics but an extremely crude,
superficial metaphysics that postulates the eternity of the atoms.

To get the atoms reduced to moments of discreteness is, however, not
enough: we must further consider the relation of the atoms, and thus of
matter, to power. The old dispute for and against a creation out of
nothing betrays in equal degree an unclarity in the disputants about the
essence of matter as well as of power. God can just as little be thought
to have created the world out of a pure nothing as out of a matter found
ready to hand. Instead of nothing and matter we posit the reality of
omnipotence. The nothing of which there is talk here is matter,
conceived as the immediate other of power, as the first universal and
yet finite opposite within power itself: that matter is nothing means
that it has not a glimmer of independent being from itself, but that
all its being is from God's power; that nothing is as a stuff, a
% ---- printed p.161 (PDF 174) ----
\opage{161} matter, means that the opposite which self-power has posited
as its other has immediate existence precisely through power and from
power. Immediate existence in the form of otherness is the issue of
power in an infinite swarm of discrete points, and thus we have the
atoms.

\emph{Disposition.} The mythical cosmogonies of antiquity are in
general so grown together with corresponding theogonies that the latter
even play the chief role; these speculative sketches of fantasy explain
to us neither nature nor creation. It is otherwise with the cosmogonic
hypotheses of modern science; they are grounded on observations, they
are strictly naturalistic. When Kepler had discovered the great laws of
the universe and Newton had developed ``the mathematical principles of
natural philosophy,'' the foundation was given. It is characteristic,
however, that Newton for his own person, without a trace of skepticism,
united the physical and the religious, the naturalistic and the
supranaturalistic, conception with one another in an entirely immediate,
almost mechanical way, without wavering in faith in the omnipotence of
the Creator. There is in this personal phenomenon, no doubt, something
more than the merely personal. If we consider the system of the world,
most immediately our solar system, in Newton's theory, then it is indeed
so that all motions can be derived from the two fundamental
determinations, the inertia of matter and the law of attraction; but in
order to explain the system in its given actuality, a plurality of
particular presuppositions and conditions was nevertheless necessary.
These particular conditions --- the initial velocity, the intensity and
direction of the impulse, etc.\ --- were apparently independent of one
another; the wonderful concurrence of the conditions, although from the
standpoint of mechanics it looked like pure chance, gave a reflection of
plan and purpose and eternal wisdom and still left a place for the
religious conception. This place, however, became ever narrower as the
surmise of a fundamental presupposition enclosing all particular
presup-
% ---- printed p.162 (PDF 175) ----
\opage{162}positions grew and at last found an expression in the
Kant-Laplace hypothesis. According to this hypothesis, the primal nebula,
contracting and rotating as it contracted, contained all the
presuppositions requisite for the natural self-formation of the solar
system. How satisfactory the hypothesis has been in a cosmogonic respect
for the author of
\textit{Exposition du systême du monde} is expressed tellingly enough in
the answer he is said to have given Napoleon when the latter asked why
he had not taken creation into account: \textit{Sire, je n'avais
pas besoin de cette hypothèse.}

Can the Laplacean hypothesis now demonstrate the impossibility of a
free creation? Far from it! There are ambiguities in the primal nebula
that easily escape the physicist's attention. The hypothesis presupposes
that the formation of the world-system in question has a definite
beginning in time. What, then, precedes the formation beginning in
time? The matter distributed in immense space! But if this matter is
to have existed from eternity, then it becomes inexplicable how, after
having been in an eternal rest, it has at last come into motion.
Further, the rotating primal vapor must not only have contained the mass
out of which, besides the sun as central body, planets and satellites
could be formed so that the purely mechanical laws of motion were
fulfilled; it must also, as a glance at the history of the earth
teaches, have contained in proportionate quantities all the diverse
substances on whose physical, chemical, physico-chemical and organic
interactions all subsequent formations, organic as well as inorganic,
depend, and from which they can be derived with necessity: the primal
nebula must, in other words, have contained a closed system of
world-possibilities, an inner manifoldness of dispositions toward
dispositions toward a progressive total development. But this looks
all too nebulous. To explain such a disposition there is required
% ---- printed p.163 (PDF 176) ----
\opage{163} in truth more than nebula: a rich sum of real principles and
Ideas must be hidden in this nebula. But Ideas and nebulae are
incompatible unless they are united in the ideal-real power of the
Godhead, in the ontological subjectivity, the self-power rich in Ideas,
which according to the logic of theism alone can ground a disposition
of the world.

Through the categories power, matter and disposition we are thus led
to one and the same point of selfhood; but here too is the outermost
limit of knowledge: the articulating determinations of selfhood lose
themselves in the Mystery.

If we next proceed from the presupposition of religion and begin with
the creating will, then there is in the concept of creation, conceived
as a concept of faith, certainly nothing that stands in conflict with
the knowledge that cognizes its limit. The free beginning, the free
disposition of the world corresponding to the Idea of wisdom, goodness
and love, and therewith God's free power over this present world, faith
can unconditionally assert without in any way encroaching upon science.
But from the ideal presuppositions of freedom it is impossible to
develop and ground the corresponding formations of nature, impossible to
make any rational transition from creation to cosmogony: all the
intermediate determinations necessary for articulation lose themselves
in the Mystery.

``Could God have refrained from creating the world?''\footnote{Dogmatics
has done ill to take up such problems of knowledge.} ``Can the whole
universe be thought to have a beginning in time?'' ``Is one to think,
in the creation of the heavens and the earth, of the whole universe or
only of a part of it?'' Such questions betray themselves plainly enough
as boundary questions for science just as much as for religion.

When it is asserted in the name of religion that God could have
refrained from creating the world, then no
% ---- printed p.164 (PDF 177) ----
\opage{164} other thought can be connected with this than that God has
brought forth the world with freedom (\textit{solo liberrimæ voluntatis
suæ imperio}, Quenstedt); but the grounds of freedom, as we have seen,
do not exclude the necessity of the Idea, only that the necessity has
been taken up into the self-determination of the will. When science
emphasizes that the concept of God becomes empty if the content of the
world, the inner together with the outer, is thought away, that God's
omnipotence must necessarily presuppose an actual world, since the real
revelation of power must necessarily posit and presuppose the other in
which power can be revealed: then this necessity cannot possibly exclude
freedom. Necessity states only the universal, that a world must be
presupposed; but \emph{which} world? That depends unconditionally on
freedom.

The question of the beginning of the world, the universal, first,
very first beginning, is in itself a nebulous question, just as
nebulous as the notion of the boundless universe. The universe is not a
finitely determinate world, but a system of worlds, or rather a system
of world-systems. If we now say with Augustine: \textit{Deus fecit
omnia simul}, then this timeless \textit{fecit} cancels every finitely
determinate act of creation and shows us the relation between God and
the universe as a universal fundamental relation. Against such a
fundamental relation nothing can be objected from the standpoint of
knowledge. Now inasmuch as science lays weight on the inner necessary
connection of the world-systems, so that each system, and each world,
comes into existence in its time and in its place precisely in
accordance with the conditions lying in the infinite whole, it must
expressly lay weight on the fact that every particular world has its
particular beginning, a beginning with a particular disposition. But,
religion adds, a particular disposition is a disposition of the Idea
determined with divine freedom; thus nature, the nature of this world,
with its ends, conditions
% ---- printed p.165 (PDF 178) ----
\opage{165} and means, is originally subordinated to the spirit, and the
blind selfless powers are dependent on the eternally seeing self-power.

And where, then, is the limit of the world of which religion speaks?
In the outward there is no definite limit; the limit is in the inward.
The difference between the religious and the scientific picture of the
world is, in brief, this: that the one who knows sees nature everywhere
in the creation, while the believer sees the creation everywhere in
nature. The religious gaze would go through the heavens, yes, through
``the heaven of heavens,'' were it possible, just as much as the
scientific; but while the knowing contemplation goes from the natural
to the natural, from the necessary to the necessary, out into the
boundless, the believing view goes through the natural to the
supernatural, through the kingdom of necessity to the kingdom of
freedom, through the creation to the Creator. In the relation to God
the outward greatness of the world's power is reduced to a sensuous
measure for the inward greatness of self-power, and the heaven of nature
to a symbol for that of the spirit. Therefore the language of religion
speaks of God in the heavens, of the Almighty, Creator of heaven and
earth, of angels and powers that surround the Creator's throne.

Here, then, there is no conflict between faith and knowledge, but there
is indeed a conflict between believers and knowers, when both, each from
their own side, disoriented as much in the fundamental requirements of
faith as in those of knowledge, confuse what is absolutely heterogeneous
and by this confusion bring about a confusion of thought most ruinous
and deplorable in theoretical as well as in practical respect.

\runhead{b) Preservation.}
\parmark{12}

The one-eyed consciousness which, if it gets sight of knowledge, will
explain everything, faith too, scientifically,
% ---- printed p.166 (PDF 179) ----
\opage{166} or, if it has turned toward faith, will derive all truths,
the scientific ones too, from the presuppositions of faith, can never
come to true insight into the relation between nature and creation.
What can be seen in this relation must be seen in preservation. The
religious view is that the Creator preserves his creation; the
scientific, that nature preserves itself. Religiously understood,
preservation is to level out the chasm between nature and creation; but
regarded in the light of science, preservation in this sense is just as
incomprehensible as creation. So-called sound human understanding is
inclined to separate creation from preservation in such a way that
creation is a miracle, preservation on the other hand a natural matter.
The first origin of the world, of all things, no one can explain; here
it is permissible to assume a miracle relegated to primeval time; but
once the world has come into being and been wisely arranged as a great,
well-ordered whole, things act by their own powers; everything then
goes its natural course, so that no continued divine cooperation can be
thought of, still less an immediate and arbitrary intervention on God's
part. That such a separation between a miraculous beginning and a
natural continuation of the life of the world is thoughtless must be
admitted. As the beginning is, so must the continuation be; as the
origin is, so must the development be. We convince ourselves of this by
considering \emph{creation} and \emph{preservation, preservation} and
\emph{natural process}, and \emph{the Mystery of preservation.}

\greekrun{α}{Creation and Preservation.}

The unity in principle of creation and preservation has been developed
from the standpoint of knowledge, in a manner as forceful as it is
one-sided, by Malebranche. The fundamental thought of faith, that God is
the free, almighty cause of the world, Malebranche transformed into a
thought of knowledge, with the consequence that in the world of things
and causes God is the only efficient
% ---- printed p.167 (PDF 180) ----
\opage{167} cause \textit{(causa efficiens)}\footnote{Similar views occur
already in the Middle Ages, e.g.\ among the Arabic philosophers
\textit{Meddaberim}.}. That finite things are taken for efficient causes
happens through a deception of the senses, or rather through a deception
of the understanding brought about by sensuous will. The one thing cannot
of itself effect changes in the other: the ball that strikes another, so
that motion arises through the impact, is not in truth the efficient
cause of the motion; the will that seems to move hand and foot is not in
truth the efficient cause of the movements. And why not? Because an
efficient cause, according to Malebranche, can be thought only as a
creating cause. A creating cause is a god. Were created things to be
efficient causes, they would have to act as created gods; but created
gods are impossible: God cannot create efficient causes, for God cannot
create gods. When it nonetheless appears to the natural understanding as
if the world of things were a world of naturally efficient causes, this
phenomenon undeniably requires an explanation. The explanation is that
natural causes are only apparent causes, finite transitional links for
the true divine cause, and thus, for God, sheer occasional causes. When
the one ball strikes the other, God takes occasion from the impact to
call forth a motion; when the will wants to move hand or foot, God takes
occasion from this to effect the movement. That the one efficient,
almighty cause is perfectly independent of all occasioning causes is
expressed thus: God can, when he will, refrain from using the occasion,
so that the natural cause calls forth no effect at all, or even use it
so arbitrarily that an effect arises which seems to stand in open
conflict with the nature of the cause.

% ---- printed p.168 (PDF 181) ----
\opage{168} That a world of such occasional causes is not able to subsist
by itself or preserve itself is comprehensible. Created out of nothing,
the occasional cause is in itself a nothing; it is not enough that it
cannot act without God, it cannot even exist without God: in the moment
God withdrew his preserving power from it, it would vanish into nothing.
What holds of the individual holds of the whole; created out of nothing,
the whole world, if the creating omnipotence withdrew, would instantly
collapse into nothing; creation and annihilation \textit{(creatio
rerum} and \textit{annihilatio rerum)} correspond exactly to one
another; preservation is in this sense just as miraculous as creation:
preservation is a continued creation.

By transforming concepts of faith into concepts of knowledge,
Malebranche has, to be sure, with a certain consistency asserted the
metaphysical reality of creation and preservation; but that this reality,
so asserted, far from satisfying science, undermines the foundation of
science is self-evident. In preserving the world, God undeniably observes
the laws his wisdom has found serviceable; the preservation of the order
once introduced depends on the whole on the observance of law. But the
observance of the laws depends solely on God's will. Although God only
seldom departs from what has once been established, since his thoughts
are unchangeable and his wisdom perfect, yet the Almighty can, at
whatever moment he will, abolish any law whatsoever when higher
considerations require it; in other words, the world order that is
maintained is maintained solely and alone by God's arbitrary wisdom.
Such a world may with a certain right be called actual, since it is
posited and preserved by divine power; but its actuality is groundless;
it has no ground in itself, since it has no efficient causes in itself:
groundless actuality is
% ---- printed p.169 (PDF 182) ----
\opage{169} natureless. But a natureless actuality, an actual world
without nature, is to science an intolerable contradiction.


But does the concept of preservation, developed in this way, then
satisfy faith? If that were the case, we should here find ourselves
in a very grave conflict between faith and knowledge. The reckoning
would then be this: it is here shown that a clear, unambiguous and
consistent development of the concept of preservation stands in open
contradiction with the concept of nature. Faith in preservation, as
irreconcilable with the acknowledgment of an actual nature, is thus
contrary to nature. But faith in preservation is inseparable from
faith in creation, and the concept of creation is consequently also
contrary to nature. What is contrary to nature is identical with the
impossible; to believe in creation and preservation is thus the same
as to believe in what is in itself impossible, the absurd.

The whole difficulty vanishes, however, when we bethink ourselves
that a concept of faith developed in the form of knowledge can never
satisfy faith. Faith does not engage with the consequence that a
creation out of nothing should make nature into a semblance; the
religious consciousness cannot, without self-contradiction, own a
metaphysics which, on the occasion of the murderer's will, makes God
the actual murderer, and in the sinner's carnal lust regards the Holy
One himself as the immediately operative cause. The religious
consciousness acknowledges the reality of nature and distinguishes
between a corruptible and an incorruptible, a sinful and a holy
nature; but between nature and omnipotence lies the Mystery.

In his zeal to get something scientific out of faith in God's
omnipotence, the pious Malebranche neglected to develop the concept of
omnipotence in logical theism. The self-doubling that lies in power
entails, as we have seen, a discharge into immediate discreteness,
from which
% ---- printed p.170 (PDF 183) ----
\opage{170} the existence of the atoms can be explained. Although the atoms are
only punctual negations in the infinite sole power, and thus a
% print: "den uendelig Enemagt" (for "uendelige")
nothing over against the self-power, they are nevertheless, as the
over-power's own elements of opposition, actually something, so that
the derived reality which must be ascribed to each individual atom is
an immediate, material reality qualified by the self-power. But that
every atom has its qualified reality means, in other words, that it is
acted upon by other atoms and in turn acts upon others. To ascribe
punctual actuality to the atoms without acknowledging their active
relations and mutual interaction is meaningless. With the interaction
of the atoms the reality of nature is founded. Every acting atom is a
thing-cause: the infinite self-cause acts point by point through
nothing but conditions and thing-causes. From this are explained
conditioned necessity, the validity of the laws of nature, and the
established order of the world.

\greekrun{β}{Preservation and Natural Process.}

For the physical view, the atoms are the last or, if one will, the
first causes of all natural things. From the properties, forces and
interactions of the atoms the formations of nature can be
sufficiently interpreted, all composition and separation, all
coming-to-be and passing-away exactly explained. In what does the
nature of a thing consist, if not in the conditioned necessity with
which it expresses itself through its properties, maintains itself in
interaction with other things, is dissolved into its elements and
again proceeds from them? What is the nature of granite? Its
consisting in a mechanical composition of quartz, feldspar and mica.
Granite cannot come forth immediately at omnipotence's bidding,
because the mediatedness of its coming-forth belongs precisely to its
nature; nor can it be preserved immediately by divine omnipotence,
because it lies in its nature
% ---- printed p.171 (PDF 184) ----
\opage{171} to endure as long as its constituents endure in their mechanical
combination, but not a moment longer. What is the nature of water? To
come forth through the chemical combination of oxygen and hydrogen.
Water arises naturally, maintains itself naturally, and passes away,
that is, is dissolved, naturally. If the chemist has oxygen and
hydrogen with the requisite conditions, then no creative omnipotence
is needed to produce water; if the chemist has the conditions for
separating the water, then no preserving omnipotence can save it from
destruction. What holds of chemical formations holds also of organic
ones: oxygen, hydrogen, carbon and nitrogen act just as surely in
accordance with their nature when a formation of cells takes place as
when a formation of water, of carbonic acid, of ammonia, etc.\ takes
place. It is true that the chemist does not command the conditions for
a formation of cells as he does those for a formation of water, and
that the living organism cannot be produced from elementary
constituents as one can produce a salt from an acid and a base; but
it stands no less firm for that, that the organic processes are
natural processes just as much as the chemical ones. Formation,
maintenance and dissolution, or coming-forth, preservation and
destruction, are only different sides of one and the same process.

It is clear from the above indications that there is no scientific
means of bringing the physicist beyond the position he has chosen.
Objections such as these: that physics cannot explain the origin of
the atoms, and thus not the unchangeableness with which they are
maintained either; that from the physical interaction of the atoms
one cannot explain the formation of cells, and thus cannot explain
either the origin of life or the division of living organisms into
species and genera, and so cannot explain either the emergence of the
species from the genus or of the genera from still more general
forms, and consequently cannot
% ---- printed p.172 (PDF 185) ----
\opage{172} explain the ground and the limit of the preservation of species and
genera, and so forth --- such objections do not touch physics and make
no impression whatever on the physicist. The standpoint of physics is
essentially empirical; the peculiar phenomena, the existing
formations, are not to be constructed \textit{a priori;} they are to
be taken as they occur in nature. It is enough for the physicist that
he finds nature before him; he does not deduce it: it is because it
is, it is of necessity; he reduces bodily formations to atoms, but the
atoms themselves he does not deduce: they are because they are; they
are in nature. The same holds of living organisms. The physicist does
not explain either why or whence life is; he finds life before him,
he investigates the living, searches out the conditions and laws of
the phenomena of life, and assures himself that life is natural, is in
nature, from nature. --- Beyond the limits of the inquiries determined
by the empirical standpoint physics does not go and will not be able
to go, even if it continues them for millions of years.

The certainty and definiteness with which physics everywhere brings
forward physical causes easily gives it the appearance that this
science stood in contradiction with religion. In principle this is a
misunderstanding. Physics says that all effects of nature pass through
natural causes; it keeps to the material side of causality, but it
does not say that what is material in the causes is the only thing
operative; it says that all vital functions depend on material causes,
but it does not say that the cause of life itself is a material cause.
Although physics itself has to do only with material causes, their
interactions and productions, it naturally cannot have anything
against metaphysics introducing other, more spiritual, immaterial
causes (principles, Idea-causes); only that this take place in such a
way that
% ---- printed p.173 (PDF 186) ----
\opage{173} none of the physical presuppositions is violated. If someone, for
example, were to assert that the principle of life, as an Idea-cause
independent in itself, was able to dispense with the material causes,
or to take the materials by force and, in defiance of chemical
affinities, use them in the service of the organism: then physics
would have to protest; but that the principle of life can, as an
Idea-cause, so work through the material causes that it develops the
possibilities lying in the affinities of the materials into
organically living interactions --- against this it can object
nothing. When the spiritualist assures us that matter is only a
semblance and atomic theory a fiction, the physicist steps forward as
a naturalist; but if it can be proved in logical theism that the
acting atoms in all their actuality are originally posited by
omnipotence, and must, in order actually to function as actual atoms,
be borne and preserved by omnipotence: then physics cannot have the
least objection to this.

The proposition of natural science, that the world preserves itself,
thus does not, when the spheres of cognition are duly kept apart,
stand in contradiction with the religious proposition that the world
is preserved by God: the natural process, conceived as a concept of
knowledge, does not exclude faith in preservation, any more than the
concept of faith, preservation, excludes scientific insight into the
natural process. What is it that physics wants to investigate? All
natural causes, all thing-causes and conditions actual and possible in
nature! But that nature is a system of causes and conditions running
on to infinity is surely presupposed by religion too, as surely as it
presupposes nature at all. And what is it, then, that religion wants
to hold fast? It is surely the truth that the Creator's all-preserving
power acts in, by and through all possible natural causes, all
material causes with conditions, intermediate causes and intermediate
causes of intermediate causes to infinity. Faith in the preservation
of the world by God can thus no more restrict the scientific
% ---- printed p.174 (PDF 187) ----
\opage{174} cognition of natural causes or set these any finite limit, than an
investigation of natural causes continued to infinity can restrict the
assurance of the preserving self-power or set any limit to the
inwardness of faith in this sense.

\greekrun{γ}{The Mystery of Preservation.}

That we speak here of a Mystery of preservation must not be
understood as though this were something particular by itself in
opposition to the Mystery in the fatherly Self, the Mystery of the
Trinity, the Mystery of creation, and so forth; the Mystery of
revelation is in itself only one: the Mystery is the absolute limit of
knowledge and faith; the Mystery is the hidden unity of the absolutely
heterogeneous principles; the Mystery is the dissolution of dualism.
As the content of faith develops from the principle of revelation, it
meets with a corresponding content developed from the principle of
reason; it looks as if religion and science once more had a common
subject, a common task, so that the contradiction and the strife
threaten each time to begin all over again. The science of reason has
its judgment about God and God's works just as religion does; so long
as the heterogeneity of the principles has not been made clear and
the circles of cognition kept apart, the judgments must conflict and
the contradictions be glaring: reconciliation is brought about only by
the demonstration of the determinations of limit, so that the
contradictions dissolve in the Mystery.

What holds of God's personality, his freedom reflected into itself,
holds naturally also of his freedom as the causality of the world,
i.e.\ of creation, and of the same freedom as the bearer of the world,
i.e.\ of preservation; but since the movements of thought on both sides
receive a relatively new point of departure, and the untying of the
knot a new tinge, the Mystery is reached by relatively different
paths, and for that reason it also appears each time in a
corresponding light of its own.

% ---- printed p.175 (PDF 188) ----
\opage{175} To the scientific view, the world presents itself, as regards both
its origin and its maintenance, in quite another light than to the
religious. Science knows no free creation, but only a maintenance
necessary by nature. Great and mighty facts, impressions of
experience, laws, causes and intermediate causes to infinity unite to
confirm the fundamental scientific view and to make everything that
religion teaches about creation and preservation as improbable as
possible. In order to disarm the misunderstanding that improbability
should be the same as impossibility, so that faith in creation and
preservation came to be understood as faith in what is in itself
impossible and absurd, logical theism has stepped in between, and by
transforming the conceptual contradictions into determinations of
limit has dissolved the impossibility into the Mystery. To believe in
creation and preservation is thus not to believe in impossibilities
and absurdities, but in the manifestations of the self-power, such
works of God as have for their grounds grounds of freedom, grounds of
will from the unfathomable depth of the Godhead.

Without regard to the possible objections of science, and long before
these came into being, revelation has spoken its word about creation,
and the religious consciousness, in the inwardness of faith, has found
inspired words for the praise of the all-preserving power, which can
shake the mountains in their foundations, can roll up the heavens like
a garment, can draw back its breath so that everything living dies.
What is peculiar to the religious view of the world must in this
connection be especially emphasized. It surely looks almost as though
the religious view of nature were just as sensuously objective as the
observation corresponding to science; it looks as though the believer
saw God himself in outward nature, in the mighty sea, where he lays
the beams of his chambers in the waters, saw the flash of his eye in
the lightning, heard his voice in the thunder, his rolling chariot in
the storm-
% ---- printed p.176 (PDF 189) ----
\opage{176} clouds, when he ``rides upon the wings of the storm''. This,
however, is not so; the objectively natural is everywhere only an
occasioning cause, only a means and point of transit for what is
subjective in the spirit. That God does not leave himself without
witness, in that he ``gives us rain from heaven and fruitful
seasons'', is a proposition of faith which does not rest either on a
meteorological consideration of the rain or on an agronomic
calculation of the crop of the field: the assurance is unconditionally
a self-assurance in faith. That the principle of faith, conceived as a
principle of cognition, is heterogeneous with the principle of
knowledge must strike the eye when we consider how impossible it is to
make any transition from the representation that it is God who gives
rain to a scientific explanation of the formation of rain, or to make
the cognition that God ``causes the grass to grow for the cattle'' the
basis of a further-reaching botanical interpretation of the growth of
plants. The religious view, which everywhere in nature discovers ``the
Father's works'', is absolutely heterogeneous with the scientific,
which discovers ``the unchangeable laws''.

That it is, however, not knowledge alone but faith also that finds its
absolute limit in the Mystery of preservation is self-evident, and
can, for further confirmation, be seen from a host of questions which
must be left unanswered from the side of the philosophy of religion as
well. How the creative and the preserving power are properly related
to each other in God himself cannot be fathomed. Doubtless it is at
bottom one and the same power; but this one and the same power
nevertheless acts in two entirely different ways: in creation, to
produce the world; in preservation, to let what has been produced
develop and produce itself. Where, now, does the creative activity end
and where does the preserving activity begin? The creative and the
preserving will is quite certainly at bottom one and the same will;
but it acts actively in creation, pas-
% ---- printed p.177 (PDF 190) ----
\opage{177} sively in preservation. That the active has a moment of the passive
in it, and conversely, there is no doubt; but where is the limit
between the active and the passive, between the producing and the
bearing energy of will? It is said expressly in the history of
creation that God rested on the seventh day. That the finite
specification of time cannot be derived from the principle of faith
has already been made clear; but the difference between creation and
preservation as two divine works, indicated by the rest, has validity
for faith. That there is a limit, and thus a difference of essence,
between creation and preservation must be presupposed; but where is
the limit? The doctrine of the relation of creation to preservation
has been connected, and with good reason, with the doctrine of the
relation of individuals to the race, a relation which is of
pervasive significance in the determination of original sin. ``It
is''
--- so it runs\footnote{H.\ Martensen. Den christelige Dogmatik.
S.\ 146--47.} --- ``the truth of traducianism that every human
individual is an offspring of the natural activity of the race, as
this is conditioned by the peculiarity of peoples, families and
parents. But it is the truth of creationism that the general natural
activity by which the race propagates itself and new souls are
formed, that this mysterious natural activity is the organ and means
of God's individualizing creative activity, that every human
individual being is thus a new revelation of the divine will, which
here prepares for itself a peculiar form for its image \dots\ Even if
it remains a mystery how, in the secret workshop of human formation,
nature and creation pass over into each other, how the creative
activity and the natural conditions mutually limit each other \dots\
yet every individual must be regarded at once from the point of view
of traducianism and of creationism,
% ---- printed p.178 (PDF 191) ----
\opage{178} or as continuation and link in the series and as a new, original
point of beginning''. If this is to become science, such a conception
of the Mystery is altogether too vague.

If it actually stood firm that the difference between being
``continuation and link in the series'' and being ``a new and original
point of beginning'' corresponded exactly to the difference between
the preserving and the creative activity, then a determination of
concept fruitful for objective knowledge would have been discovered in
the Mystery: we should then be able to develop new concepts of
knowledge from a concept of faith; but that is impossible. Of animal
individuals too it holds, surely, that the living individual is both
formed from a new point of beginning and at the same time as
continuation and link in the series. If the animal individual can come
forth without the individualizing principle of life entering into
activity, then --- the matter regarded in the light of knowledge --- a
new animal species must also be able to come forth without a special
principle of species (Darwin's hypothesis). If it stood firm that the
relation between natural and creative activity could be determined by
the relation between less and more gifted individuals, then the
Mystery of creation would be reflected in so many rays of light from
anthropology, history and psychology that the mysterious would sink
down to a vanishing obscurity, a pure relativity. But the individual
richly gifted by nature is not in a deeper sense a creature of God
than the less gifted one. If the genius of a general, combined with
political genius, is to be a fruit of God's special creative activity,
then the same must hold of mathematical genius, poetic genius, the
genius of an actor, etc.\ When the great gifts of nature are without
further ado regarded as great gifts of creation, the small gifts of
nature must of course be regarded as small gifts of creation; but to
make all gifts of nature into gifts of creation is precisely to efface
the limit between natural and creative activity. To regard the
dif-
% ---- printed p.179 (PDF 192) ----
\opage{179} ference between worldly and religious natures as the distinguishing
mark is just as unreliable; for if an inborn difference between
religious and non-religious natures could actually be conceived, then
the Creator would surely exclude from religion those who did not have
an inborn religious nature. The marks of difference that faith sets in
this respect are subjective, not objective. The Mystery effaces all
marks of knowledge, all objective limits, and precisely thereby
designates the absolute limit between faith and knowledge.

\runhead{c) Governance.}
\parmark{13}

Preservation is the divine function of power which furnishes the
middle between natural and creative activity. Between the three
different functions the difference is essential, and yet every finite
limit has vanished. ``God creates all things'', ``Nature produces all
things'', ``God preserves all things'': in this threefold generality
the functions melt together. On what the fusion rests, and how the
determinations of difference develop out of the mysterious fundamental
unity, is now to be more closely investigated. Attention fixes first
on the thing-causes, then on the Ideas, but yet first and last on the
almighty will. Of the thing-causes we saw in the foregoing that they
are originally moments for omnipotence, realities within its own
self-doubling, that they are, in other words, all created. Further,
we saw that the atoms, these original thing-causes, are on the one
hand points of transit for omnipotence and to that extent occasioning
causes, on the other hand existing realities, and thus operative
causes. If now the thing-causes are posited in the moment of points
of
% ---- printed p.180 (PDF 193) ----
\opage{180} transit, the emphasis falls on creation; if they are posited in the
moment of operative causes, the emphasis falls on natural activity; if,
finally, they are posited in the moment of created independence
resting on omnipotence, the emphasis falls on preservation. In
preservation the divine power is posited as the medium passive with
respect to the natural causes, a medium which, as regards essence, can
in a certain sense be said to be what space is with respect to the
phenomenal. Just as space does not itself act to change the distances
or dimensions of bodies, and yet is absolutely indispensable in the
determination of dimensions and distances, so the preserving,
all-bearing power, the medium of the real concepts and of the forces,
likewise does not change the reality and mode of action of the natural
causes, but is nevertheless, as the all-inspiriting, all-pervading
medium of the concepts and forces, absolutely indispensable to the
activity of the natural causes.

If, now, the preserving activity, which keeps to the universal, cannot
decide anything regarding the individual and the particular, then the
determinations of particularity must, so to speak, be divided between
the thing-causes and the creative will. ``If something new is to be
produced, the creative will must set itself in activity; if special
formations are to take place, special causes and conditions must be
presupposed''. But a finite division in this sense is only a
parceling-out and piecing-together of what is in itself infinite, and
thus a contradiction. The contradiction is resolved by the
intervention of the Ideas. The creative will produces everything
through Ideas; the finite natural causes point back to the
generative, forming ground of nature, the \textit{natura naturans}
impregnated with Ideas: the concrete unity-in-difference, rich in
content in itself, of creative will and conditioning natural causes
must, as regards the great whole, be sought in the Idea of the world.
The Idea is in principle a teleological uni-
% ---- printed p.181 (PDF 194) ----
\opage{181} ty of means and ends. The Idea of the world unfolds its content in a
system of Ideas, in the plan of the world; the creative will, acting
through the Ideas, and thus in the form of wisdom, aiming at the goal,
then reveals itself as the governing will.


It is, however, so far from being the case that faith could give any
objective knowledge of the divine governance that the opposition of the
principles shows itself the moment the development begins. It is a fact
that the governance of the world proceeds through a system of Ideas; in
the system of Ideas there is unity of necessity and freedom; but when it
is expressly asked in which form the unity is to be held fast, in that
of necessity or that of freedom, then faith and knowledge stand sharply
against each other; for knowledge holds to the identity of necessity
and freedom in necessity, but faith in freedom. The consequences of this
appear when we consider, from the standpoint of science,
\emph{governance and the course of the world}, and, from the standpoint
of religion, \emph{the providence revealed in governance} and \emph{the
personal faith in the governing providence.}

\greekrun{α}{Governance and the Course of the World.}

If the actual whole, comprising nature and history, which we in a
special sense call the world, means nothing other than the
self-unfolding of the Absolute through finite oppositions, then the
course of the world is objectively comprehensible and governance is easy
to explain. Here there can be said to be in the great process of
development both an end and yet not an end. The Absolute constantly
goes out from itself and, through the intermediate stages of the
development, constantly returns to itself. If emphasis is here laid on
the difference between the Absolute's beginning and resulting
development, then the Absolute in the result of its development is end,
the intermediate links are means, and the conception is teleological.
If, on the other hand, emphasis is laid on the Absolute's absolute
identity
% ---- printed p.182 (PDF 195) ----
\opage{182} with itself under all forms of development, at all stages of
transition, in all final results, then the development is an infinite
process, then there is here, in relation to the Absolute, no progress,
no purpose and no final end, but the Absolute first and the Absolute
last in the eternal circuit of the movements.

As it goes with the comprehensive whole, so also with the particular
parts. If we consider the formations of nature and bring out the
opposition between organic forms at lower and higher stages, in
particular the opposition between animal and man, then the ascending
movement takes on a teleological appearance. By manifoldly winding
detours, through the most diverse preparations and approaches, the Idea
has known how to work its way up from obscurity to clarity, from
dispersions in the elemental through totalizations in the organic,
through rising concentrations, to that summit of individualization where
the universal knows itself as the subject and the Absolute becomes
conscious of itself. What wisdom, what plans and purposes: say now that
there is not a governance here! And yet it is a small matter to put all
this teleology, these means and aims, aside again as soon as, by their
decisive claims of selfhood and the omnipotence of will reflected in
them, they begin to be a nuisance. What, after all, is animal
individuality? A reciprocal determination, thoroughgoing at its stage,
of the individual with the particular in the universal. The individual
and the particular must necessarily be as it is, because the universal
expressed in it necessarily is as it is; there is in the aim of these
organic formations no aim, in these apparent means no means: the finite
difference of means and aims dissolves in the totalized unity of the
Absolute into an indifference that obliterates all aims.

This ambiguity plays a great role where the talk is of the aim of human
life, of the purpose for the
% ---- printed p.183 (PDF 196) ----
\opage{183} individual, for society and for states, of the epochs and periods of
world history. One extols the right of reason and the might of reason
and the cunning of reason; the individual is an end in himself, the
state is an end in itself; reason is in league with freedom, and freedom
is the goal. From this teleological point of view one is then also
easily reconciled with unfreedom, injustice, evil. For evil, after all,
has its ground only in the finitude of the world, and finitude is
inseparable from actuality. A world without limits and negativity would,
as both philosophers and philosophical Church teachers have recognized,
be the emptiest, most monotonous and most imperfect world one could
think of, ``a painting without shadows, a piece of music without deeper
tones among the higher, a speech without pauses, a verse without short
syllables among the long, a poem without antitheses, a drama without a
villain, a sheer one and the same without multiplicity''.

But since no self is absolute under this refraction of oppositions,
inasmuch as the Absolute is selfless, and the self in each individual
human being is gnawed by the worm of corruptibility, this whole drama
of freedom and self-development is an empty play: the serious matter is
that each individual being is a point of passage for the universal, and
the significance of the ephemeral individualities is necessarily
determined by their position and place in the whole. In the Absolute's
``divine thinking and willing each individual is posited with precisely
only that measure of perfection which it actually has. When we see a
blind man, we compare him with others who see, or with his own former
condition when he could still see, --- likewise when we find someone
carried away by sensuous desire: and therefore we judge that both are
contrary to his nature, the former as physical, the latter as moral
evil. But on the contrary, once the divine will or the concatenation
of causes has brought about
% ---- printed p.184 (PDF 197) ----
\opage{184} blindness for one individual and awakened sensuous desire in another,
then in the first case sight, in the second, at this moment, the
rational constitution belongs as little to the essence of such an
individual as the former to that of a stone, the latter to that of an
animal: the supposed privation is thus only negation, which cannot be
subject to any blame''\footnote{Spinoza. Cf.\ Strauss. Den
christl.\ Troesl.\ II. S.\ 307--8.}.

The philosophy of the Absolute has, as regards the development of form,
undeniably worked its way far beyond the Spinozistic abstractions and
knows, in a quite different and penetrating manner, how to bring out the
particular, acknowledge the conflicts of freedom, and assert the
significance of national individualities, of ages, of heroic passions
and of great characters in world history; but as regards the point of
selfhood, the religious-ethical cardinal point, pantheistic speculation
has not come a hair's breadth further than Spinoza. ``The government of
the world itself'' --- says Strauss --- ``is not to be regarded as a
determination of the course of the world by an understanding standing
outside the world, but as the reason immanent in the cosmic forces and
their relations themselves: and the fact that the destiny of mankind, of
the nations and of individuals stands under the guidance of providence
can only have the meaning that, on account of the universal dominion of
spirit over nature, the development of the human race on the whole
proceeds according to its concept, and that the contingency of the
individual's action and of the natural occurrence constantly evens
itself out again into universal necessity''.

Still the old theme, the same empty selfless essence! The play has been
rewritten, the exposition is richer, the catastrophes more exciting, the
costumes more brilliant and the scenery more splendid; but the plot is
the same, the untying of the knot the same: the
% ---- printed p.185 (PDF 198) ----
\opage{185} inadequacy of the individual wills for the universal of spirit, the
swallowing up of individuals and races by necessity, the corruptibility
of the self and its half-hour eternity in the Absolute --- that is the
result of the struggle, the outcome of the course of the world, the
share of governance in the whole.

\greekrun{β}{The Revealed Word Concerning the Governance of Providence.}

With creation and preservation the foundation is given for the
governance of the world. Since the creating and preserving will bears
and permeates the activity of nature in an infinite manner, the ground
of necessity and the natural independence of the world can subsist
together with the hidden miracle of preservation, and no proof drawn
from science can establish the impossibility of the governing
cooperation of an omnipotent, omnipresent will. The final separation
between what happens by natural causes and what is immediately effected
by God has the inner infinite unity-in-difference as its presupposition
and is determined solely by revelation and faith. Through everything in
the world, the least together with the greatest, God can place himself
in a personal relation of will to man, and man to God. It is not the
event according to its immediate objectivity, not the outwardly actual
fact, that determines the relation; it is conversely the relation that
determines the conception of the fact and gives it its significance:
not what happens, but how what happens is interpreted and understood,
is the main thing. From this it follows that the revealed word
concerning governance must, from an objective point of view, take on an
arbitrary appearance, and the representation be colored by the
imagination. On account of the absolute heterogeneity of the principles
this cannot be otherwise.

A sensuous natural world, an actual order of things, is given; this
foundation is included in the immediate presupposition of all religious
judgments. What is
% ---- printed p.186 (PDF 199) ----
\opage{186} peculiar to faith is that it refers the natural facts now to the
creating, now to the preserving and, on the whole, to the governing
will. When, then, it is expressly said in the revealed word that God
determines each individual human being's birth and span of life, that
he draws man forth from the mother's womb and shelters him from
childhood on (Ps.\ 22:10 ff.), that man's days are numbered and written
in his book (Ps.\ 139:16), etc., then there is hereby not given the very
slightest contribution to a scientific determination of the relation
between the physiological working of nature and the scholastic fixing
of God's cooperation (\textit{concursus}), but a correct statement of
what religion and faith are.

If natural science could actually prove that man's natural coming into
being was conditioned by such forces and causes as cannot possibly be
subject to the omnipotent will, then this one thing would be enough to
make faith impossible. By an intermediary independent in nature, such a
dividing wall would be raised between the divine and the human self that
an essential relation to God could not be conceived. By such an
intermediary God would have to cease to be omniscient, omnipotent,
i.e.\ cease to be God; nor could man then be set free by the unfree
Deity, but would have to put up, from the beginning of his life onward,
with being bound to obscure selfless powers of nature.

To presuppose God's free omnipotence --- and that, after all, is what
the believer does --- to be in infinite need of God's liberating
self-power --- and that, after all, is what the believer does --- and then
nevertheless to assume that the human self, according to its natural
origin, is separated from God, sprung from a root torn loose from God's
knowledge and will, is an intolerable self-contradiction. It is this
self-contradiction that faith sublates, and that with a power of
inwardness all the greater, the purer and stronger it is. It is with
this
% ---- printed p.187 (PDF 200) ----
\opage{187} strength of inwardness that the revealed word speaks of the human
self's original springing forth from the divine; it is to the dependence
of the formation of man on the creating will that the Psalmist bears
witness when he says: ``I will give thanks to thee, for I have come into
being in a wonderful manner. My bones were not hidden from thee when I
was formed in secret, when I was curiously wrought; thine eyes saw me
when I was yet unformed'' (Ps.\ 139:33).
% print: "(Ps. 139, 33)" -- as printed (clear in KB and Bodleian; possibly authorial); the psalm has 24 verses; probably meant Ps. 139:14-16.

But although the word of revelation thus refers the disposition of each
human individuality to the creating self-power that governs all things,
it is by no means overlooked that the disposition is an actual natural
disposition, and as a natural disposition must presuppose natural
connection and be subject to natural conditions. The reality of the
natural conditions is, as regards this point, not merely included in
the immediate presupposition, but comes, through the cognition of sin,
to play a leading role in particular.

It has already been shown that the cognition of sin rises with the
cognition of God, so that the intensity of the one of these cognitions
can be regarded as the gauge of the strength of the other; it has
further been pointed out that the cognition of sin lays bare the
self-contradiction that man both can and yet cannot take the whole guilt
upon himself. Now, as the cognition of sin steps in separatingly between
the holy and the sinful self, how then is the penitent before God to
explain to himself what we may here call the inner infinity of sin, of
the sinful essence? To make the creating will the author of sin is
blasphemy; to excuse oneself with the imperfections of one's contingent
individuality is frivolous and irreligious; to take the whole guilt upon
oneself in the sense that the arbitrarily free will, which commits sin
because it wills to, can also make itself free of sin whatever moment it
wills, is brazen superficiality: what, then, remains? That the penitent
must take the whole guilt upon himself, place the cause of sin in the
will,
% ---- printed p.188 (PDF 201) ----
\opage{188} but in such a way that what is sinful in the will is referred to its
natural ground. ``I was born in iniquity, and my mother conceived me in
sin'' (Ps.\ 41)
% print: "(Ps. 41)" -- as printed (clear in KB and Bodleian; possibly authorial); probably meant Ps. 51:5 (51:7 in Danish numbering).
--- in these words the thought of original sin is expressed. Original
sin is not willed by God; but it falls under the divine governance.

That the concepts of faith are in principle absolutely heterogeneous
with the concepts of knowledge shines forth from the propositions set
up here: man's individual disposition is determined by creative
activity, and man's natural disposition is sinful. Should these two
propositions be conceived scientifically, they must yield a conceptual
contradiction that cannot be resolved. The insolubility of the
contradiction does not lie in the fact that a self-cause works together
with finite natural causes; for then the contradiction can be resolved
by reducing the self-cause to an immanent Idea-cause: the solution must
then be sought in the identity of the material organism and the
organizing principle. Nor can the insolubility be placed in the fact
that the will of God stands outside or over against the finite causes
and, by holding these in absolute passivity, transforms them into
non-causes; for the whole difficulty with an absolute \textit{agens}
standing outside is an easily seen-through chicanery of the negative
critique, which, by the deterring notion of ``a God standing outside,''
exempts itself from any closer treatment of any theistic problem
whatsoever that goes deeper. No, the insolubility lies in the fact that
the creating will, which cooperates in the formation of each individual
human being's natural disposition, is a holy will that wants sin away,
and yet expressly --- according to the thought of original sin ---
refrains from cleansing the natural disposition from the ground up: this
contradiction cannot be resolved in objective knowledge. If one says:
the divine will cannot cleanse the natural disposition, then the natural
disposition forms a limit, and we fall back into the contradiction
touched on above. If one says:
% ---- printed p.189 (PDF 202) ----
\opage{189} the connection of nature does not tolerate a thoroughgoing cleansing of
each individual disposition by itself, then one places the reality of
nature in sinfulness, and we then have the contradiction in another
form. If one says: God will not cleanse the natural disposition, then
one appeals to God's arbitrariness; but grounds of arbitrariness are not
scientific. From the standpoint of faith the contradiction is resolved,
namely in the form of self-contradiction: the solution is in the
expression of the cognition of sin.

With this, faith in the governing will has been characterized; the
wisdom that governs the world cannot be justified by objective reasons;
that is against its essence and nature: it is personal, absolutely
personal; the last reasons, the deepest motives, are hidden in the
fatherly Self --- in the Mystery.

All reasons of governance are, on God's side, reasons of will, and thus,
as regards understanding, on man's side reasons of faith. To dismiss the
scientific debates concerning the relation between God's foreknowledge
and men's free will, between the divine plan of the world and the evil
in the world, between God's unchangeableness and the effects of men's
prayers, etc., will in the present context be superfluous: according to
the standpoint chosen they fall away of themselves. On the other hand,
there is another side of the doctrine of providence which the philosophy
of religion must illuminate in particular; it presents itself where the
talk is of definite dispensations, particular aims and corresponding
means. The statements belonging here are of very different kinds, have
on the whole a different stamp in the Old than in the New Testament,
and must, as regards their universal religious validity, be judged
according to the principle of faith.

In the election of the people of Israel, its wonderful deliverance from
Egyptian bondage, its legislation, its wandering through the desert, its
conquest of the
% ---- printed p.190 (PDF 203) ----
\opage{190} promised land, features of a governing providence are seen so great,
so clear, so significant, that the people of Israel has rightly been
called ``the typical people.'' But if we go through the believing
statements about this in detail, it is unmistakable that the strong
images, the colorful descriptions, the interpretations offered of
natural phenomena and historical events bear witness to the changing
moods, the religious conceptions and subjective judgments that are
peculiar to that people with its prophets and priests. To conceive of
the plagues of the land --- pestilence, locusts, drought, failure of
crops, the calamities of war, slavery and exile --- as chastisements of
the punishing providence, and of deliverance, victory and prosperity as
signs of the grace of providence, is at once so popular and so
religious that we cannot on the whole conceive of faith in providence
otherwise. That elements of a lower kind mingle in the religious
conception, eudaemonistic expectations, egoistic secondary notions and
the like, is explicable from the sensuousness of the people and the
one-sidedness of the preparatory standpoint.

In order not to misinterpret the details we must constantly remember
that the immediate judgments of faith are judgments of mood, a
circumstance that can very well be united with what is fundamentally
religious in their content. Thus, e.g., when Jehovah says in the prophet
Amos (3:6): ``Shall evil befall a city, and the Lord hath not done
it?''\ then the thought is by no means that Jehovah's will is the cause
of evil as evil; but the thought of governance is that fortune and
misfortune, evil and good, all dispensations are so wholly subject to
God's will and might that every intermediate cause, every finite author,
vanishes as a nothing when the talk is of standing in the way of
Jehovah's thoughts and plans.

The doctrine of providence is in its concept just as anti-rational as
the doctrine of creation; which is why it is also characteristic that
the doctrine of providence in the New Testament can be so
% ---- printed p.191 (PDF 204) ----
\opage{191} spiritually pure without being rational. That God clothes the lilies
of the field and feeds the birds of the air is, from a rational point of
view, only a lyrical-poetic expression; but that from the standpoint of
the Gospel it has another meaning is seen from the application: ``If God
so clothe the grass of the field, which today is, and tomorrow is cast
into the fire, how much more shall he not clothe you, O ye of little
faith!''

\greekrun{γ}{Faith in the Governing Providence.}

The opposition characterized here between the scientific and the
religious conception of the governance of the world confirms anew the
fundamental proposition set up concerning faith and knowledge as
absolutely heterogeneous principles; we now come to an answer to the
following two main questions: 1) Can the one who knows believe in a
governing providence? and 2) what significance does faith in providence
have both for the one who knows and for the one who does not know? The
answer to both questions is already contained in the presupposition of
principle that the relation to God is an eternal relation of selfhood,
and the goal personal liberation, personal self-development.

\emph{1) The one who knows can, without coming into contradiction with
science, believe in the governing providence.} By identifying the one
who knows with science, it is an easy matter to tighten the opposition
between faith and knowledge into caricature. ``Science cognizes no
personal God, no creation, no preservation by a free omnipotent will,
and hence no free omnipotent governance of the world either, and yet the
one who knows is to believe in all this: he is thus to know that a
personal God is unthinkable, and yet believe in God's personality; to
know that the notions of a creating, preserving, governing will lead
into insoluble conceptual contradictions, and yet, in spite of these
actual contradictions, in spite of the open
% ---- printed p.192 (PDF 205) ----
\opage{192} breach of the principle of contradiction, do violence to his whole
spiritual being, persuade himself and others that he can actually
believe in what he himself has cognized to be manifest contradictions!''
From what we have said before it can now be seen that these distortions
rest on sheer surreptitious assumptions and false presuppositions. It
is, as we have seen, a subtle surreptition that science should actually
be able to prove either the impossibility of God's personality or the
impossibility of any work of God whatsoever corresponding to it. One
confuses two entirely heterogeneous results of knowledge with each
other: being able to comprehend the possibility, and being able to prove
the impossibility --- a crude, unscientific confusion! It is a false
presupposition, a subtle surreptition, that the one who knows should
have no need of cognitions of any other kind than such as can be
grounded in science. Science is not a personal being, but the one who
knows, however much he may, in objective self-forgetfulness, forget the
ideal demands of personal life, is and remains a personal being. If the
one who knows awakens to personal self-cognition, then the very first
thing he will see is this, that science, according to its principle,
cannot possibly give him decisive personal answers to decisive personal
questions. It is, when we consult science, neither comprehensible nor
probable that the Absolute itself should be a personal being, but it is
possible --- yes, it is possible! It is, from the point of view of
knowledge, neither comprehensible nor probable that there should be an
Omniscient One who could see everything with full clarity, an
omnipotent Self who was able to master, determine and guide everything
down to the least, the very least details, so that even all the hairs
of the head are numbered, so that not even a sparrow falls to the ground
without the will of the Father who is in heaven; it is incomprehensible,
it is improbable, but --- it is possible. By settling personality's
account with
% ---- printed p.193 (PDF 206) ----
\opage{193} science and presupposing the principle of faith, the one who knows gets
the possibility out on all sides. More than the possibility he does not
get; more than the possibility he does not need.


\emph{2) Faith in providence has essentially the same significance for the one who knows as for the one who does not know.} In the course of personal life the present always draws its most essential significance from what is to come; the actuality on which everything here chiefly depends is a future actuality; but a future actuality is always a possibility. The acting individuality lives, in actuality, on nothing but possibilities. The way in which he conceives the conditioned possibilities of actuality depends on how decisively he acts, and on how he understands himself in the moment decisive for the action, whether his self-understanding falls outside or inside the relation to God.

If the self-understanding falls outside the relation to God, God as God being denied, then the highest concentration the subject can attain is a self-concentration in the Absolute. Were it possible to drive the concentration to the point where the self became absolute, then the corresponding action in time would not merely hold for eternity but would even be an act of eternity, whereby the one acting had made himself God. But to make oneself God, or, what is the same thing, to make one's own self absolute, is impossible. The Absolute does not itself act; the acting self is not absolute: as soon as the action is accentuated, the Absolute and the self part from each other. Where, then, is governance here, what is it that governs? To govern is, acting with an eye to the goal, to cross through conditions and possibilities: the Absolute cannot act, and so it cannot govern either. The Absolute has possibilities enough---indeed it has all possibilities---it releases them into the
% ---- printed p.194 (PDF 207) ----
\opage{194} manifoldness of contingencies, it gathers them under the unity of necessity, but that does not help the one who acts. Under the pressure of necessity, surrounded by accidents and contingencies, the acting subject must on its own account take on the possibilities, must itself set its goal and determine its way: the acting subject must govern itself. And by what shall it steer? To steer by the Absolute instead of by a relative goal is like a skipper wanting to steer by the whole of space instead of by a star. To say that the Absolute is in the one acting, is in the possibilities, in the purposes, is meaningless; for the secret is that the Absolute is so immanent in all relations of finitude, in all acting subjects and actions, that precisely because of this pure immanence it is outside the whole of it. So the self governs itself finitely and goes under in infinity.

It is otherwise when the self-understanding falls inside the relation to God. The religious individuality has its decisive concentration in faith itself: faith is the action of actions, the inward, first, absolute action. The believing self does not through the act of faith itself become absolute, but from the absolute Self it receives the principle of action and of governance as its first mover. That the absolute Self can be a principle of action lies in its being able, in the most original sense of the word, itself to act; for the same reason it can be a principle of governance; as principle in the relative it is what it is in itself as absolute. The absolute Self can act. For what is it to act? It is to make the will immediately the causality of actual changes in an actual world: the absolute Self is one with the absolute will. The absolute Self can govern. For what is it to govern? It is to make the will \emph{mediately} the causality in a system of conditions and means corresponding to a certain end. The objection that an acting and governing God is a
% ---- printed p.195 (PDF 208) ----
\opage{195} contradiction, since God, by acting and governing, would have to be drawn down into finitude and thereby himself be made finite, is superficial. The acting and governing God does indeed himself enter into finitude, enters, so to speak, into the terms of finitude; but since this happens through action and governance, it happens in such a way that finitude is at one and the same time both upheld in its reality and permeated by infinity. The unity of the self with itself through all actual oppositions is unity of will,% print: comma missing after „Villieseenhed“
{} the actual unity of will is unity of action. Through the unity of action, then, the divine Self can unite the relative with the absolute, can work in all that is finite and yet, as the one infinite Self, be exalted above the finite.

The immediate expression of the divine action, God's work, is the miracle, the manifest miracle. The principle of the miracle is, as already shown, the unconditioned principle of will. That the will as will can become the causality of actual changes in nature, this and this alone is what is miraculous in all miracles. Through the miracle, according to what has been made clear in the foregoing, nature is by no means either sublated or botched; for the causality of will does not annihilate the causes of things, it only transforms them: the miracle, seen from the side of nature, is an effect of the interaction of transformed natural causes. Conceived as divine action, the miracle expresses that the natural, what in its elementary naturalness is external and finite, is permeated and transfigured by the spirit in an inward, infinite way. The miracle makes it manifest that God has a will.

The essential expression of governance is the will's use of intermediate causes in their immediate finitude and relativity. In the miracle the intermediate causes seem quite to vanish before the immediately intervening will; in governance,% print: stray semicolon „i Styrelsen; synes“
{} conversely, a mediately intervening
% ---- printed p.196 (PDF 209) ----
\opage{196} will seems quite to vanish before the intermediate causes. Both forms of the will's activity must arouse offense in the unbeliever. That the miracle arouses offense as contrary to reason---the contrariety to reason lies, namely, in the fancy that the intermediate causes have been annihilated---is understandable; but governance, which with its intermediate causes looks so reasonable, how can governance arouse offense? For the opposite reason, namely that the governing will disappears altogether for the observer fixed on the intermediate causes, so that it becomes a will that is no will. It is infinity from the other side that here tips the balance. When a finite power of will governs, it is always noticeable in one way or another how it intervenes here and there in the course of events. The finite will, however great the shrewdness and wisdom of the one governing may be, can never so omnipresently permeate means and intermediate causes that it would not at some point need to intervene externally; but this the divine will can do: it governs in an inward, infinite way, and therefore it is as though it did not exist at all, as though there were no governance.

The miracle and governance contrast with each other and yet hold together in infinite interaction: the miracle is governing, governance is miraculous; the miracle is the manifest, governance the hidden miracle. The expression of the unity of miracle and governance is given in the revealed word: ``\emph{With God all things are possible.}'' With this word the believer appropriates the principle of the governance of all things as the principle of personal self-development. That all things are possible with God means that the ideal tasks of freedom can, with the aid of governance, be solved in infinitely many ways within the existing order of the world. Necessity is the firm foundation of freedom, and is to that extent necessary for freedom; but in actual existence necessity is
% ---- printed p.197 (PDF 210) ----
\opage{197} also a limit to freedom. The more relative the demand of freedom is, the more natural is the self's agreement with the finite limits and conditions: a reasonable man expects only the probable, not the improbable, sets himself to attain the reasonable, not the extraordinary. The higher the flight, on the other hand, that personal development takes, the more ideal the demand of freedom becomes, the sharper becomes the conflict between the probable and the improbable; the fulfillment of the ideal demand will in actuality always prove improbable: then the task is to will the possible with enthusiasm, however improbable this possible may be. The inward, unconditioned liberation of the self---unconditioned insofar as the personality is to subject all finite conditions to itself---must therefore drive the one acting forward in a direction that has the current of probabilities against it: the one who acts thus wills the possible in spite of the improbable; it almost looks as though he willed the impossible, for in faith he wills unconditionally his freedom and blessedness. So placed, the individuality disposed toward freedom, be he otherwise one who knows or one who does not know, will surely learn in self-understanding to understand the significance of the words: with God all things are possible.

It must be emphatically impressed that, as regards faith in the governing providence, the talk is not of the actual but of the possible; thus not of what is finite in actuality but of the infinite, for what is infinite in actuality is the possible. With God all things are possible: such a faith in the governance of the world seems at first glance to verge on the fabulous; but the limits of actuality, the entanglements of life, the pressure of necessity and the impersonal powers are included in the presupposition, and the presupposition therefore contains guarantees that faith does not become fantastic. Here there is no question of what governance \emph{will}
% ---- printed p.198 (PDF 211) ----
\opage{198} do in this or that case; here the talk is only of what it \emph{can} do. This distinction between what governance can do and what governance will do is of infinite significance for the life of faith.

The believing self that ventures out on the sea of improbabilities must, seeing itself surrounded by dangers, be just as prepared to perish temporally as the unbelieving self; it must, as has already been said, practice infinite resignation. But the difference between the resignation determined by the Absolute and that which corresponds to faith in governance shows itself in this, that the former is only a reconciliation with necessity, the latter by contrast a personal deepening in freedom. Let them both, then, the one who knows, living on the Absolute, and the believer, who has placed himself under the fatherly providence, let them both suffer the same thing outwardly, both go under outwardly in the same way: here there is nonetheless an infinite difference inwardly. The one who knows, who goes under in the Absolute, goes under because the selfless dissolves the self, because personality is finite and only the power of necessity infinite, because the individual, a merely temporary appearance of the universal, must vanish in death. The believer, on the other hand, who lives under providence, goes under in no danger, however threatening and apparently unavoidable it may be, because it follows unavoidably from necessity, but because governance wills it. And why does governance will it? Because it is thus the best, is what serves this personality best. If it is not personally the best, then the believer is fully assured that he will not perish temporally either. Therefore the believing resignation is so inseparably bound up with hope. Hope holds unconditionally to the eternal, but also takes up the temporal conditionally. As hope is, so is fear: ``Fear not them which can only kill the body, but cannot kill the soul''; ``yea, even all the hairs of your head are
% ---- printed p.199 (PDF 212) ----
\opage{199} numbered'': these words were given by the revealed providence as provision for the way to those who were going out precisely to meet such as both could and would kill the body.

The reciprocal determination of fear and hope receives, in the inwardness of faith, its deepest explanation in prayer. ``With God all things are possible'': this is the fundamental thought of prayer when it is seen in connection with governance. That man, under all finite dispensations, in fortune and misfortune, in sorrow and joy, has essentially to do not with the external dispensations but with the all-governing power, God's eternal, unchangeable will; that the infinite task of freedom, of personality, is: under all possible dispensations to understand the governing will and in this believing understanding to understand oneself---that is the truth, the personal truth, at whose attainment prayer aims. The one who prays presupposes that the governing will, for which all things are possible, both can and will help, not \textit{in abstracto} but \textit{in concreto}, not merely in general with the eternal, but for the sake of the eternal also in the temporal, the finite, the present actual. The one who prays does not want, in a selfish sense, to try to bring about a change in God's will---if that is so, he violates the principle of faith and is not a believer; if that is so, he violates the revealed word, for he then posits some finite help or other as the only possible one, although the word expressly says that with God all things are possible---the one who prays does not in his hoping forget infinite resignation: he prays to the Unchangeable; but the Unchangeable is not the Absolute of necessity. The Unchangeable is the heavenly Father---``the ideal of majesty!''; yes, but also the ideal of love and personal freedom, in love and freedom the Unchangeable.
% ---- printed p.200 (PDF 213) ----
\opage{200} \lettersub{C.}{The Kingdom of the Father.}
\parmark{14}

At the expression ``the Kingdom of the Father'' thought is involuntarily led to a comprehensive worldview, a coherent cognition of the fundamental relations of existence. It must above all not be forgotten that the worldview the philosophy of religion here means to develop is not the scientific but the religious one, and that the development is aimed especially at setting the opposition between the cognition that falls inside and the cognition that falls outside the relation to God in the clearest light.

It is the human world, with the great turning points of the life of the spirit, that steps into the foreground. The object of consideration is on the whole common to religion and science. The state of things on earth when man appeared, the beginning of human life, the transition from the natural to the cultural state, the separation of the peoples, the difference between the people of the myths (of the Ideas) and the people of revelation (of the spirit), the education of the people of the spirit until the absolute turning point, ``the fullness of time,'' is reached: these facts, which speak for the governance of the world, may well be said to lie as close to the science of reason as to the philosophy of faith; but precisely because the object, although outwardly taken it is the same, nevertheless shows itself so very different in the light falling on it from the opposed standpoints, precisely because the world of spiritual events becomes quite another when it is viewed in the light of science than when it is presented in the light of religion, we must lay weight on the alternation of principles and the heterogeneity of standpoints.
% ---- printed p.201 (PDF 214) ----
\opage{201} \runhead{a) Ideal and Actuality.}
\parmark{15}

In the actual relation to God, the one corresponding to man's given condition, the life of the spirit moves around the two cardinal points, sin and redemption. Now since redemption can have no significance except in relation to sin, the question of the origin and essence of sin is one of the very first and deepest fundamental questions of religious consciousness. From what has already been made clear about the relation between the cognition of God and the cognition of sin, it follows that the answer to the question must be sought in revelation. To clarify what sin is, how sin came into the world, and its consequences, sacred history begins with a presentation of what we may call the ideal natural side of the relation to God, passes from there to a presentation of man's breach with the ideal in the relation, and then gives the hard actuality, contrary to the ideal and just for that reason so oppressive and burdensome, an interpretation that corresponds to the original, undimmed consciousness of sin. It is thus the \emph{ideal of innocence}, the \emph{temptation} and the \emph{Fall}, with the ensuing \emph{curse}, that we are to consider.

\greekrun{α}{The Ideal of Innocence.}

And God said: ``Let us make man (Adam) in our image, after our likeness \dots\ So God created man in his own image, in the image of God created he him.'' Seen under his ideal determination, man is thus an image of God. Wherein the image of God properly consists is not said outright, but is nevertheless not difficult to find out when we note Jehovah's words to man. Man is to rule over the earth: ``replenish the earth, and subdue it''; and man is to obey God's commandment: ``of
% ---- printed p.202 (PDF 215) ----
\opage{202} every tree of the garden thou mayest freely eat: but of the tree of the knowledge of good and evil, thou shalt not eat of it''; by the first of these determinations man is conceived as a rational being, by the latter as a religious being. It thus belongs to man to develop not merely in his relations to the world but originally according to his relation to God, for---``in the day that thou eatest thereof thou shalt surely die''---the relation to God is what is decisive. Created in the image of God, after the ideal of love in the image of the eternal Son, man is ``the free, personal unity of spirit and nature, a \emph{spiritual soul}, which is not, like the natural souls, captive in corporeality, but is destined freely to reveal itself through its body as a temple of the spirit''\footnote{Martensen. ``Den christelige Dogmatik.'' S.\ 140.}.

To the ideal description of man in his originality now corresponds the whole surrounding, the Garden of Eden, the animals, which Jehovah brings to man to see what he will call them, etc.\ On children and childlike souls such a description must involuntarily make the impression of something actual, something that once existed but has now, alas, vanished. And this is indeed actually so, if we take the ideal actuality for what it in truth is, actuality of the spirit, actuality for the spirit.

But the naive fusion of the ideal with what is actual in the external sense has produced misunderstandings and confusions which, if they were to persist, must end by undermining religion. On the presupposition that the facts of sacred history must absolutely have the same external actuality as those of profane history, the biblical presentation of Paradise has been taken, in a spiritless sense, for an externally actual history. In this way faith and knowledge have been
% ---- printed p.203 (PDF 216) ----
\opage{203} fused together, and out of the fusion their mutually irreconcilable strife has arisen.

In the name of faith the explanation runs thus: If the story of Paradise is not actual, then there has never been any actual Paradise; if there has never been any actual Paradise, then no actual Fall has ever taken place either; if no actual Fall has taken place, then the whole of revelation is false and faith a fancy. Now we know that revelation cannot possibly be false; therefore we know that an actual Fall must have taken place; therefore we know that at the beginning of history there must have been an actual Paradise in an actual place on earth.

In the name of knowledge the explanation runs thus: There has never been any actual Paradise; there has never existed in nature either a tree of life or a tree of the knowledge of good and evil; no such Fall as the one the Bible speaks of has ever happened; revelation is therefore false, and faith a fancy.

Both parties are agreed in identifying the principles of faith and of knowledge; both are agreed that what is true in religion must also be true in science, and conversely; both are therefore agreed in drawing scientific consequences from religious presuppositions and religious consequences from scientific ones. But what, then, is the result of all this agreement? Disagreement! It stands fast in science that the biblical Paradise has never been actual, so fast that natural science would instantly have to undermine its own foundation if it admitted the possibility of the opposite; and it stands fast in religion that the ideal state of innocence has reality, so fast that faith cannot possibly give up this reality without giving up itself. To settle the strife there is only one way out: to cognize the absolute heterogeneity of the principles. The actuality at which revelation aims in
% ---- printed p.204 (PDF 217) ----
\opage{204} the description of its Paradise is not the actuality that natural science sets out to investigate: the two kinds of actuality, the ideal and the material, belong to two quite different regions.

Already from an aesthetic standpoint it must be recognized that it is not nature in its outward material immediacy that appeals to man, but that everything depends on the mood it awakens, and so on man's inner being. Not what is actually seen and heard and felt, but the eye with which one sees, the ear with which one hears, the heart with which one feels, is what lends nature its spiritual luster. We find it in order that the poet, in order to open to us an insight into the essence of nature, so to speak creates nature anew, so that we may rightly see how it is created; for we understand that he produces illusion only because the illusion is even more truthful than actuality. Is Paradise, then, an illusion? By no means! As high as religion stands above poetry, so high stands the religious explanation of nature above the merely aesthetic one.

The believing mind resists with reluctance conceiving the narrative of Paradise as a philosophical myth, a poetic clothing of a profound thought; and from the standpoint of revelation it cannot be conceived so either: the content of the narrative has taken on the stamp of actuality and means to count as full actuality. But then it also becomes a matter of becoming conscious of what kind of actuality it is, so that one does not confuse the higher with the lower.

The lower actuality is sensuous; it stands and falls with the testimony of the senses. He who says: I cannot believe the story of Paradise unless I may assume that the trees in that garden had the same external actuality as the trees in any other garden, builds his faith on sensuous ground; his faith is superstition.
% ---- printed p.205 (PDF 218) ----
\opage{205} Yet here is an objection of a more serious kind. ``If the serpent that seduced Eve was not externally actual,'' one might say, ``then Eve, who was seduced, was not externally actual either, and so neither was Adam, nor Cain and Abel: our `first parents,' with the descendants whose names are read on the oldest genealogy of the human race, are thus sublimated into ideal beings of fantasy without contact with the so-called lower actuality; but that such a conception conflicts both with the nature of the case and with the express testimony of Scripture is incontestable.'' This objection, however striking it must seem at first glance, rests on a misapprehension, as thoroughly ruinous as it is pervasive, of the original position in which, according to the principle, the historical and the religious stand to each other in the revealed word.


Religion is of the people; the history of revelation goes from the beginning
together with the history of the people, namely as this history was conceived
by the people chosen for revelation. But it is so far from being the case that
historical faith grounds religious faith that the opposite holds: it is
precisely religious faith that governs, explains, and grounds historical faith.
Originally the two heterogeneous elements are fused together; but after the
conflict between faith and knowledge has been driven to its extreme, the moment
has arrived when a critical separation must be carried through: the critique of
knowledge has already done its part; it remains for the critique of faith also
to do its part.

The ideal of innocence, whose actuality is made visible in Paradise, is
precisely the ideal of actuality that forms the bright background of the
cognition of sin, the presupposition without which sin cannot be actual. The
difference between the aesthetic and the religious ideal is here conspicuous.
The aesthetic ideal belongs to the imagination alone; it satisfies
contemplation, but makes
% ---- printed p.206 (PDF 219) ----
\opage{206} no demand upon the will; it requires no realization; in illusion it has its own
unactual actuality, its completion and conclusion. The religious ideal, on the
contrary, is in the very strictest sense of the word an ideal of actuality; for
it is first and last an ideal of the will, and the ideal of the will demands its
realization. Sin is guilt; the sinful will is guilty. But the will can have
become guilty only through an actual break with its ideal, i.e.\ through an
actual sublation of its own innocence. Original innocence is therefore an
actual, inescapable presupposition of actual guilt, and for that very reason
contains the actual, inescapable demand that sin must be atoned for, that guilt
must one day be blotted out.

The ideal of innocence thus stands at one and the same time both before and
behind actual human life: before, as the final goal toward which the guilty one
yearns; behind, as the point of departure without which his yearning and
anxiety would be in vain. Inwardness in the cognition of actual guilt and
inwardness in the cognition of original innocence are in principle the same
inwardness, the same actuality: it is actuality of spirit, actuality for the
spirit.

Insofar as actual human life is marked by guilt, the sinful condition comes
immediately to present the given as actuality, namely the untrue actuality,
while the ideal of innocence, the actuality true in itself, is pushed back, has
assumed the form of something unactual, and presents possibility. Man in his
state of innocence is man at the beginning of his life, conceived according to
his inner true essence, i.e.\ according to his ideal as actual man. Both
sides---that the beginning actuality is ideal, and that the ideal of the
beginning is actual---the revealed word emphasizes with equal force. Revelation
speaks a double language, but without ambiguity. It speaks of the actual
begin\-%
% ---- printed p.207 (PDF 220) ----
\opage{207}ning of human life, thus of actual human beings, of the first human beings,
``our first parents''; but it speaks also of the original, the ideal-actual in
this actual beginning; then it speaks in images, in symbols; then the veil of
corruptibility is drawn aside; then the holy vision is seen: nature transfigured
in the light of the Spirit, the springs of peace with the tree of life, Adam and
Eve in the garden of Eden.

\greekrun{β}{Temptation and Fall.}

``Now the serpent was more subtle than any beast of the field which the Lord
God had made, and it said to the woman: Is it so, that God hath said, Ye shall
not eat of every tree of the garden? And the woman said to the serpent: We may
eat of the fruit of the trees of the garden; but of the fruit of the tree which
is in the midst of the garden, God hath said, Ye shall not eat of it, neither
shall ye touch it, lest ye die. And the serpent said to the woman: Ye shall not
surely die; for God doth know that in the day ye eat thereof, then your eyes
shall be opened, and ye shall be as God, knowing good and evil. And the woman
saw that the tree was good for food, and that it was pleasant to the eyes, and
a tree to be desired to make one wise, and she took of its fruit, and she did
eat, and she gave also to her husband with her, and he did eat.'' (Gen.\ 3:1--6).
This is in brief the ``supra-historical'' history of the Fall, a history that
yields no gain for knowledge, but all the richer a gain for faith.

It is the woman, according to the ideal the most innocent of all women on
earth, who is here appointed to make manifest how sin came, and still continues
to come, into the world. Innocence is a two-sided concept; one can say of
paradisal innocence both that it is and that it is not destined to disappear.
Insofar as innocence designates the first child\-%
% ---- printed p.208 (PDF 221) ----
\opage{208}like immediacy of the life of spirit, and this is indeed precisely the case
with Eve's innocence, it is certainly destined to be sublated. Before the
temptation sets in she is without any guilt, and therefore innocent. Innocence
shows itself in this, that the personal is still a possibility; the duality in
the natural-spiritual self has not been awakened to conscious separation: the
religious and the selfish elements still slumber in tranquil unity. She is
innocent; her soul is pure, a mirror for the dawning of the Spirit in the
incipient relation to God; but in her innocence she is nonetheless a natural
self, and in her natural self she has a hidden egoism. Thus the child, in all
its innocence, is a little egoist. This hidden, still unconscious egoism is not
in itself sinful; it is only the natural disposition of the natural
will, and the natural disposition of the will is necessary as the foundation for
the will of the spirit. The ideal determination of the natural will is to be
awakened by the Spirit to self-denying obedience. If the will obeys, then the
ambiguity of innocence is sublated in the sense that guilt is excluded and the
tested faith, with self-conscious freedom, has been set in its place. If, on the
contrary, the will refuses the Spirit obedience, then egoism comes forth from
its hiding place and sublates the ambiguity by making the innocent one guilty.

For a closer determination of the relation between ideal and actuality we must
begin with the relation to God. The relation to God corresponding to the
principle of revelation is actual; it is indeed a relation between the working
God and an actual human soul; Eve is to that extent to be regarded as an
actually existing individuality. But this actual individuality is more than a
single individual; the ancestral mother is, essentially seen, the individuality
in all individuals yet to come, individuality according to its first ideal
fundamental determination; she is \textit{instar omnium}. All attempts to
explain psychologically the spiritual condition of ``the first human beings,''
i.e.\ of the first human individuals, must
% ---- printed p.209 (PDF 222) ----
\opage{209} fail even for the most acute, for they rest on a misunderstanding. The ideal
presentation that is to reveal to us what from the very beginning dwells
essentially and originally in the soul must go far beyond the immediate
condition of empirically limited individuals. Here one cannot ask how Eve in
her first naturalness could have understood Jehovah's word and grasped the
Spirit's command so clearly and purely; such considerations of finitude stem
from an objectively reasoning understanding that wants to know something. Here
there is nothing for knowledge, but everything for faith. Eve's soul is the
soul in its first ideal relation to God; but in its first ideal relation to God
the soul can understand the Spirit, understand the thought of the holy will
deeply and purely and clearly; for without the possibility of such an inner
understanding faith would lack its principle, and the relation to God
% print: Gudsforhøldet (wrong sort ø for o)
would remain in eternal ambiguity.

For a closer determination of the relation between ideal and actuality we must
also take note of the holy symbolism: it is the tree of knowledge and the
serpent that capture our attention. That the tree of knowledge is as little an
external, tangible actuality as an empty allegory is easy to see. The slumbering
life of the spirit is to be awakened to self-activity; it is awakened by the
Spirit through the word, it is awakened through the law: the commandment of the
law is a prohibition. But a prohibition cannot, after all, be given unless there
must in actuality be something that is prohibited. In the reciprocal
determination of the prohibition and the prohibited, the relation between the
actual and the ideal shows itself here. The prohibition against eating the
fruit of a definite, named tree is so unambiguous, so immediately
comprehensible, that not even the most innocent can mistake it. The actual
definiteness, then, is stressed as strongly as possible. And yet the ideal
significance goes far beyond any limited particular whatever. For what
% ---- printed p.210 (PDF 223) ----
\opage{210} kind of fruit is it, then, whose enjoyment is here forbidden? The forbidden
fruit! The object of the prohibition coincides to such a degree with the
prohibition itself that the universal is: Thou shalt not do what is forbidden!
The tree of the knowledge of good and evil is the object of the trial; but since
the knowledge of good and evil depends precisely on the outcome of the trial,
the object of the trial is essentially a moment in the trial itself: in such
visible forms does the symbolism of revelation express the first proclamations
of the Spirit.

Between the tree of knowledge and the woman the serpent suddenly appears. It is
not enough here to ask: what does the serpent \emph{mean}; we must expressly ask:
what \emph{is} the serpent? If it holds of the tree of knowledge that it
\emph{is} what it means, then the same must surely hold of the serpent. It lies
in wait by the tree and---it lies in wait for the woman. At once, at the first
glance, the serpent shows itself to be a contradiction; and indeed it is a
contradiction: it is the author of contradictions, contradiction itself from
the beginning. In the first place it is a contradiction that the serpent is an
animal and yet a spirit. That it is an animal is plainly said; that it is a
spirit we hear from its speech. What, now, does this animal spirit, this
spiritual animal, mean; whence does it come, what is it? We answer: it comes
from the hidden depths of natural life; it rises up from the ground of the
soul; it is the egoism hidden in paradisal innocence, which in this symbol
projects its shadow-outline. This is seen from the manner in which it comes. It
comes, as it seems, of itself: no one has sent for it; and yet it does not come
of itself, for the spirit has conjured it forth, the spirit of contradiction
that the prohibition awakens. That is why it is so characteristic that, as if
to make sure, it begins by asking whether it is really so that God has forbidden
man to eat of all the trees in the garden.
% ---- printed p.211 (PDF 224) ----
\opage{211} Here, then, we have the prohibition reflected in egoism. Through the reflection
that rebels against the prohibition, egoism becomes spirit. In elemental
naturalness egoism is animal; with animal unruliness it works in the strong,
raw, savage natural man. Were it to be presented in this character, it would
surely have to be symbolized by the wolf, the tiger, the panther rather than by
the serpent. But in the form of rawness it could not possibly make its
appearance in Paradise. It comes forth, then, under the symbol of cunning,
falsehood, and craftiness, and takes the innocent one by surprise. Just as the
word from God is clear to Eve, so now is the corresponding word too, an echo of
self-will from the selfish interior, a clear and definite contradiction of the
word of God. It happens before Eve's eyes; it takes place in the soul, and the
soul knows of it, and yet does not know of it. The ambiguity is at its height;
that means that it must disappear: so, then, the temptation begins.

``Let no man say when he is tempted, I am tempted of God; for God cannot be
tempted with evil, neither tempteth he any man. But every man is tempted, when
he is drawn away of his own lust, and enticed.'' (Jas.\ 1:13--14). These last
words some have wished to use to explain the serpent away, as if the serpent
meant that the temptation came from
without\footnote{S.\ Kierkegaard. Begrebet Angest. S.\ 43.}, but without any
ground. With the serpent the temptation by no means comes from without, but on
the contrary from within: precisely for that reason the external symbolized in
the serpent is necessary. All temptation presupposes reflection; but the
reflection in the first temptation has this peculiarity, that it is
involuntary, that it takes place as if of itself. Certainly Eve is tempted by
her own lust; but she does not yet know the subjective essence of lust: if she
knew it, she would not be innocent. Innocence shows itself
% ---- printed p.212 (PDF 225) ----
\opage{212} precisely in this, that the hidden egoism, the natural selfishness of human
nature, emerges as a being different from herself; the egoism is in herself,
and yet is not her personal self, for the personal self is in the will that is
now to concentrate itself in the choice, whether she chooses obedience or
disobedience. It is the immediacy in the tempting reflection that gives the
temptation so objective, and precisely thereby so innocent, a tinge. What is
tempting in the phenomenon forces itself upon Eve as a fact; she sees that the
tree is good for food, and pleasant to the eyes, and then---yes, then it is a
moment of perdition, a single now, and then she becomes guilty.

But how, now, does it really stand with this guilt, this very first sin? That
Eve has fallen is certain; but in what does her guilt consist? She is, after
all, taken by surprise, beguiled; before she has had time to collect herself and
offer resistance, she has already fallen. Her fall is like that of one who
suddenly stands at an abyss and then grows dizzy and plunges down. And yet she
has become guilty, as essentially guilty before God as any soul in the world can
become. The moment in which she falls is precisely the moment in which she could
have been saved, namely the moment when faith was to come into being for her. In
the untested immediacy of innocence faith is still only, so to speak, conceived,
but not born; it can be born only in temptation, for temptation awakens doubt,
and faith is the inner act that strikes down doubt and holds fast to God. This
act is the fundamental act of the will of the spirit: it is a moment that
encloses a whole eternity. If the moment passes without the act of faith taking
place, then the selfish has conquered; the temptation works like a sorcery; the
eye darkens, and even if the one grown dizzy recovers her sight in the next
instant, it is another sight, the sight of desire upon the false ideal: the self
is transformed. What transformation takes place in Eve in the moment of
temptation
% ---- printed p.213 (PDF 226) ----
\opage{213} is seen from this, that without hesitation she believes that it is God who is
the liar, and the serpent who speaks the truth.

In excuse of Eve people have pleaded her inexperience, her lack of firmness, her
guilelessness, thus her innocence. But however much such grounds of finitude may
signify where the talk is of the psychological weaknesses of a limited
individual, they signify nothing here, where the talk is of the ideal collision
of the will of the spirit and the fundamental egoism in the depths of
individuality. Instead of diminishing the guilt, the preceding innocence will on
the contrary heighten it. How, in the spiritual sense, should the inexperienced
gain experience except by setting faith against doubt and resisting temptation?
By succumbing to a first preliminary temptation and becoming provisionally
guilty, she could surely not possibly gather strength to preserve her innocence
in the decisive trial. The only thing that could still be said in excuse of the
fallen one would then have to be that the first temptation was too great, was
beyond her strength. But it is not so. Eve has in her innocence all the
presuppositions, all the help that a human being can have, for meeting the first
temptation; the temptation itself is of such a nature that if she were able to
resist any, it would have to be this one. She lacks nothing, indeed she has
everything in abundance: she may ``eat of all the trees of the garden'' except
the one. What tempts her is not something natural in itself, but, when we look
more closely, something so heaven-storming, so spiritually rebellious, that it
could, one would think, be tempting only to a Titan, not to an innocent woman.
For what is it that tempts her? A forbidden fruit, people say, an apple
perhaps, nothing else! Yes, it looks childish and insignificant enough; thought
can almost be revolted by the notion that such
% ---- printed p.214 (PDF 227) ----
\opage{214} a fuss is made over so slight a thing, that the whole human race should be
punished so severely for a piece of apple. That this is a spiritless
misunderstanding anyone can see. If one says: she is tempted by the forbidden;
the forbidden, whatever it may be, always tempts; that lies in human nature;
\textit{in vetitum ruimus, petimusque negata}---then one has reduced the highest
to a triviality. No, what tempts her is nothing less than the ideal-egoism of
wanting, by transgressing God's commandment, to become like God:
\textit{Eritis sicut Deus.} That this main point is not sufficiently heeded has
its ground chiefly, no doubt, in this: that with all manner of mystical
secondary notions about what the serpent is---perhaps the masked devil---one
lets the serpent keep the thought and Eve thoughtlessly reach for the fruit. But
such a conception conflicts not only with the spirit of the narrative, but also
with its letter. For the woman, who sees that it is ``a tree to be desired to
make one wise,'' enters precisely into the thought, the serpent's thought; she
reaches for the fruit expressly in order to gain understanding, the
understanding of good and evil that is to make her equal with God. To a higher
spiritual majesty Adam himself could not rise; he therefore follows Eve in
silent admiration; thinks as she does, and does as she does: thus sin came into
the world.

Is this, now, to be comprehended? In faith! yes---but not in knowledge.

From the standpoint of knowledge the task would be to explain the Fall by
grounding the act of will from which guilt proceeds; but that is impossible.
Whether one (Schelling, Steffens), in order to explain sinfulness, assumes an
``intelligible act'' taking place outside time, whereby the individual has
predetermined himself to evil; or one (Hegel) derives sinfulness from the
contradiction that obtains between man's concept and his immediate naturalness;
or one (Fichte) conceives evil
% ---- printed p.215 (PDF 228) ----
\opage{215} as the \textit{vis inertiæ}, the original sloth, according to which the ego
wants to remain in the given state of nature instead of grappling with its
limits and working its way forward to freedom: the knot is equally insoluble.
The question that remains unanswered in all these attempts is: what reality
does the will have? This question science cannot answer, since will must in the
last instance be explained by will, by a fact of self-determination which, in
spite of all reasons and counter-reasons, rests not on universal necessity of
thought, but on the willing self. Were scientific grounding to be possible, one
would have to have the strength of the will as a magnitude known by itself, the
resistance of nature likewise as a definite magnitude by itself, and then, by a
kind of psychological differential calculus, state the point of equilibrium and
find out the conditions under which the preponderance would have to fall either
on the side of the resistance of nature or on the side of the will. But the
strength of the will can in no case be stated by any finite known magnitude
whatever. For the utmost exertion of the will there is no finite, fixed limit.
Even if there were such a limit, sin and guilt still could not be explained from
it. For guilt cannot be derived from the resistance having been stronger than
the will; one would then have to subtract the quantum of inhibition that
indicated the excess; for the will does not become guilty because it has not
\emph{been able} to do the impossible, but because it has not \emph{willed} the
possible. The relation of the will to motives cannot be scientifically
determined. If one lets the motives determine the will, then the will's
self-determination ceases; it then becomes merely the tongue on the balance of
the motives; it evaporates into a semblance. If, conversely, one lets the will
determine itself without motives (indifferentism), or at any rate determine the
motives itself, then it must either determine itself blindly, and thus quite
involuntarily, or have in itself the motive according to which it decides for
or against
% ---- printed p.216 (PDF 229) ----
\opage{216} the other motives; but by having the decisive motive in itself, it becomes
incommensurable for scientific grounding. When predeterminism wants to explain
sin from an act of will preceding all sinfulness, it is right insofar as it
wants to derive sin from a sinful will instead of deriving the sinful will from
sin; but it falls into contradiction with itself when it comes to explaining how
it is thinkable that a will preceding sin and sinfulness should be able to be so
sinful as to determine itself to sin. Sin is a concept of faith and never
becomes a concept of knowledge: the origin of sin cannot be explained in
science.

From the standpoint of faith the task is soluble, namely in the sense that all
the moments corresponding to the principle can be cognized from the revealed
word. The ideal of innocence shows that the fundamental will, with its two-sided
disposition, has in the relation to God the possibility of subordinating to
itself the natural self-will and determining itself as will of the spirit, but
also the possibility of neglecting the decisive moment and letting self-will
rule. Outside the relation to God, sloth, \textit{vis inertiæ}, is no doubt the
only acceptable, though insufficient, ground of explanation when the question
is of the preponderance of the sensuous over the spiritual: the first sin is
then only an impersonal act, merely \textit{peccatum omissionis}. But within the
relation to God the choosing self is so strongly accentuated that
\textit{peccatum omissionis} is precisely to be regarded as the forming
personality's own fundamental act, its first \textit{peccatum commissionis}.
The innocent one is tempted so that personality may come into being: the tempted
self is \textit{eo ipso} choosing. In the relation to God corresponding to the
ideal of innocence the choosing self has not only the obligation but also the
condition to concentrate itself into will of the spirit and put faith into
force. When the self nonetheless does not do so, the reason is that in its
selfishness it will not do so. The choosing
% ---- printed p.217 (PDF 230) ----
\opage{217} self does not will sin immediately, for it wills only its self-will; and yet
this is a delusion, for to will its self-will is, according to the given
prohibition, expressly to will sin: the self sins, therefore, because it wills
to sin; its first act of personality is a sin; it does not sin from sinfulness,
but it produces sinfulness by its sin, and sin by its personal will.


\greekrun{γ}{The Curse.}

The first sin is a free act of consciousness, insofar as it is a personal act
of will; and yet through the act itself a transformation takes place in
consciousness. Before the act, sin is still only a possibility; after the act
it is an actuality; the act itself is the free translation from possibility to
actuality; by this translation sin acquires so altered an aspect that it has
not merely become something else in the deed, but has even, by this its
alteration, altered the whole actual surroundings. ``And the eyes of them both
were opened, and they knew that they were naked; and they sewed fig leaves
together, and made themselves girdles'' (Gen.\ 3:7). In these few words the
psychological effect of the Fall is compressed. In its first innocence the
spirit is dreaming; the dreamer is a child who does not see actuality as it
actually is: the dreamer must be awakened. But spirit is awakened by spirit, in
that it awakens itself. When, then, it awakens itself under the self-power of
self-will, it wakes in sin to consciousness of guilt; it wakes to see what the
child could not see --- its nakedness, and it is ashamed. It wakes to
self-reflection, to work; but this work is not the free deed of freedom, the
first bold step on the way of culture; it is on the contrary an involuntary
expression of the bad conscience, a wretched endeavor to hide the nakedness
that sin has laid bare. ``And they heard the
% ---- printed p.218 (PDF 231) ----
\opage{218} voice of Jehovah walking in the garden in the cool of the day, and Adam and
his wife hid themselves from the presence of God Jehovah amidst the garden.''

What, according to what we have said above, holds of the state of innocence and
the Fall, holds also of the punishment and the curse: the revealed word acquires
its meaning only when the content of the account is conceived as an object of
faith and judged from the standpoint of faith. Two diametrically opposed modes of
view, the historical-literal and the philosophical-mythical, have each in its
own way contributed in equal measure to confusing the conception.

How impossible it is to carry through the historical-literal interpretation may
well be assumed, at the general level of cultivation of the present day, to be
self-evident. We cite merely one example. ``In the day that thou eatest
thereof,'' says Jehovah (Gen.\ 2:17), ``thou shalt surely die.'' Taken
literally, Adam should accordingly have died on the same day he transgressed
God's command. Was this threat then fulfilled; did death then come on the same
day? --- ``Assuredly death came on the same day!'' --- remarks Irenaeus ---
``For the Fall took place, together with the creation, on a Friday, and Adam
died on this day 930 years later. Or, more exactly still, on the very same day,
namely by divine reckoning, according to which 1000 years are $=$ 1 day
(Ps.\ 90:4); for Adam did not, after all, live to be 1000 years old.''

That the mythical interpretation can come to terms with the words of the text in
a far easier and more natural manner is understandable; and this circumstance
has indeed in recent times done its part to efface the boundary of principle
between myth and revelation. ``In the divine pleasure garden, whose cultivation
and keeping had been entrusted to the first human beings, these had permission,
it is true, to eat of all the other trees; but one Jehovah had
% ---- printed p.219 (PDF 232) ----
\opage{219} reserved for himself and the other divine beings, namely the tree from whose
fruits the Elohim obtained their higher insight. \dots\ By threatening death as
an inevitable consequence, Jehovah had therefore tried to deter the humans from
partaking of this fruit; but the serpent, which on account of its cunning saw
through the real motive of this prohibition, revealed it to the woman, who now,
through the enticing hope of an insight equal to God's, and through the
beautiful appearance of the fruit, let herself be led into partaking of it, and
straightway thereafter also seduced her husband to do so. At once the food of
the gods had its effect in the humans; they awoke from the animal
unconsciousness of the difference between good and evil, the seemly and the
unseemly. \dots\ With displeasure Jehovah, who jealously watches over his
prerogatives, remarks that the humans have gained a knowledge like that of the
gods; he cannot let such an overstepping of the bounds set for them go
unpunished. \dots\ The most important thing is that man, so that he may not also
seize the Elohim's other prerogative, immortality through (a continued)
partaking of the fruits of the tree of life, is driven out of that place, the
garden of Eden.''\footnote{Strauss. Den christl.\ Troeslære II. S.\ 15--16.}

It lies in the nature of the case that the narrative itself is not in a position
to decide which of the two modes of explanation is the only right one. What
turns the scale here is the principle alone. What ``God Jehovah'' said:
``Behold, the man is become as one of us, to know good and evil'' (Gen.\ 3:22),
can certainly from a mythical standpoint be conceived as a divine envy
(\textit{θειον Φθονερον}), but it can and must none the less, from the standpoint
of revelation, be conceived as an expression of the righteous wrath of the
Spirit. How necessary it is that the conscience which wakes in sin, veiled and
stupefied by sin,
% ---- printed p.220 (PDF 233) ----
\opage{220} be awakened by a word that reveals the holy and the righteous, can be seen in
Adam and Eve. They have a bad conscience; they want to hide from the presence of
the Holy One. At the moment they committed the trespass they had quite forgotten
the Omniscient and took no notice at all of the Omnipresent; but hardly is the
deed done before there is a stirring in the conscience, and Jehovah is present:
they hear his voice ``walking in the garden in the cool of the day.'' And yet
selfishness is so strong, and the stupefaction so great, that when the voice
calls and demands an account, the answer is not given in repentance with
acknowledgment of sin; no, the answer is given with evasion and excuse: each of
the guilty excuses himself by shifting the guilt onto another, the man onto the
woman, the woman onto the serpent. This system of excuse is a soul-corrupting,
spirit-darkening system of stupefaction. With sin the system of stupefaction
came into the world; but for that very reason the stern word of holiness, the
punishing word of righteousness, has also been heard in the world. If there is
any point in sacred history where it becomes truly manifest that the judgment
of conscience must in principle be a judgment of God, it is surely here.

The serpentine lies in the stupefaction. ``And God Jehovah said unto the
serpent, Because thou hast done this, thou art cursed above all cattle, and
above every beast of the field; upon thy belly shalt thou go, and dust shalt
thou eat all the days of thy life.'' So the judgment begins, then, with a curse.
The lightning stroke of the curse, however, does not strike man immediately ---
had it struck immediately, man would have been crushed --- no, the lightning
strikes at man's foot; it strikes, but it does not kill: it strikes something
human, but in the symbol. It is natural egoism, the selfishly sly in man's
essence, the false principle of will, which has brought and still brings sin
into the world; it is the principle of selfhood, actual and yet, in spiritual
essentiality,
% ---- printed p.221 (PDF 234) ----
\opage{221} ideal, that is immediately struck by the curse. With vital vividness the
actual animal --- in natural-historical actuality certainly not accountable,
but through cunning and venom naturally alarming --- is stamped as the
cunningly stinging hereditary enemy of the race, as a symbol crawling in the
dust.

As regards the penal sentence upon man, it becomes evident to an unprejudiced
eye that the lot of the woman and of the man corresponds straightforwardly to
the burdens that actual existence, according to the factually given world
order, lays upon both. Either, so it seems, there can properly be no question
here of punishment, or, if the burdens imposed are really to count as
punishment, it must be assumed that another world order in actuality preceded
this present one. Now since the latter assumption is flatly rejected by science,
and that not on account of an unbelief which argues from the improbable, but on
account of a cognition of the world so penetrating that everyone who knows sees
the impossibility: the conflict between faith and knowledge is once more on the
point of breaking out. The curse on the serpent alone is enough to show where
the literal explanation leads. Not to speak of the animal's lack of
accountability, even if the devil himself had spoken through it, the supposition
that the serpent originally had feet leads to absurdity and justifies the
objection from natural science that the curse itself is incorrect and merely
rests on the Jewish prejudice that the serpent actually eats dust. If it were in
truth impossible to carry through a critical separation between a merely
external actuality and actuality for the spirit, then the principle of
revelation would really be dissolved; then faith would have to give up. But in
this connection too it is seen that the conflict is imagined, not actual.

% ---- printed p.222 (PDF 235) ----
\opage{222} It is not primarily an external change in nature but an internal change in man
that is in question. The transition from innocence to guilt coincides here with
the transition from ideal to actuality; as man is in earnest about to make the
acquaintance of stern actuality, he comes, according to the penal sentence, to
begin with what is still far sterner, the cognition of his trespass. The
cognition of the trespass and the cognition of the stern conditions of actuality
cannot, from the standpoint of revelation, stand as two cognitions indifferent
to each other. If that were so, the consciousness of sin would soon be stifled
under the pressure of actuality. Instead of being dissatisfied with himself, man
would readily excuse himself with the imperfection of existence and cast the
blame on actuality, on its burdens and sufferings, its pains and bitter toil.
So that this shall not happen, the punishing word sounds to man: this you have
deserved; such a hard and burdensome actuality you have brought upon yourself by
your guilt; you have deserved more than that, you have deserved death. The task
for faith, then, is not to speculate on how actual nature would have looked if
Adam had not sinned; rather, the task of faith is here, as everywhere, a task of
inwardness; the task is to cognize original sin and guilt so deeply that it can
in earnest be understood that the burdens of actuality are deserved.

But --- it is asked --- is not something explained away hereby that nevertheless
belongs essentially to faith, namely ``death and the corruptibility of the whole
creation''? We answer: No! here not one jot, not one letter is explained away of
that which in truth is, or ever has been, an object of faith. As regards, in the
first place, the corruptibility of nature, science will not deny that there is
corruptibility in nature. The question, then, is only in which light
corruptibility is to be
% ---- printed p.223 (PDF 236) ----
\opage{223} regarded, that of necessity or that of freedom: here, then, we again have the
absolutely heterogeneous principles.

Under the point of view of faith, and so under that of freedom, the Pauline view
(Rom.\ 8:19--23) is unconditionally correct, but precisely for that reason not
scientific. When the ideal of freedom furnishes the presupposition for the
judgment on the corruptibility of nature, then the consideration takes its
starting point in the believing consciousness; then corruptibility comes to mean
a condition in which the one consecrated to freedom, awaiting ``the adoption, the
redemption of our body,'' is bound; then nature echoes the sigh that presses
forth from the breast of man; then we cognize ``that the whole creation groaneth
and travaileth in pain together until now.''

As regards, next, the relation between sin and death, it ought not to be passed
over in silence what science has established by striking, incontrovertible
facts: that in the physical sense there was pain and death on earth long before
the appearance of man. But it must be emphasized on the side of faith that it is
not death as a physical fact but the consciousness of physical death on which
it essentially depends. For the animal death does not exist; for when the animal
is, death is not, and when death is, the animal is not: only for a mortal,
self-conscious being does death exist as death, or, only in the consciousness of
death is death. It can, then, be understood in the language of science too that
only with man, the self-conscious being that cognizes death as death, does death
come into the world; on the other hand it cannot be said on any scientific
ground that death is a punishment for sin, that ``death is the wages of sin.''
This, by contrast, can be said with full right in religion. The decisive thing
is that man is guilty, that his original consciousness of spirit has become a
consciousness of guilt,
% ---- printed p.224 (PDF 237) ----
\opage{224} so that he who does not know within himself what guilt is does not know, in
the religious sense either, what spirit is.

Just as man is not merely by nature, and so far with necessity, disposed to work
and bear the burdens of actuality, but in consequence of sin is expressly
sentenced to work in the sweat of his face and to regard the thorns and thistles
that the earth bears as a punishment --- indeed as a gracious punishment, in
that the curse, instead of annihilating him in his inner being, has fallen upon
the outward: so man is not merely born naturally with the germ of death in his
body, but in consequence of the transgression is spiritually sentenced to dread
death as the dark enemy, obliged to bear as a punishment for sin the
overwhelming impression of a sensuous nature that confirms Jehovah's word:
``dust thou art, and unto dust shalt thou return.''

This is the essential beginning of human life: the ideal vanishes in sin;
actuality begins with guilt. The lost ideal is to be reached again; that is a
promise. But the way goes forward, not back: the cherub with the flaming sword
guards the entrance to the tree of life.

\runhead{b) World-Ideals.}
\parmark{16}

The religious worldview is not, and will not be, what in the language of
cultivation is called many-sidedly objective. It is not the inner reciprocal
struggles of the working forces, not the refraction of individual wills against
the universal of reason, not the gradual, conditioned shaping of ideas in
national peculiarities, states and civilized societies, about which the revealed
word will instruct man: every enlightenment tending in that direction it leaves
to reason. It is a direct consequence of the absolute heterogeneity of the
principles that the cognition of the world
% ---- printed p.225 (PDF 238) ----
\opage{225} which reason can develop and ground is not imparted by revelation, while
conversely the view of the world that reason essentially lacks is opened up in
religion. A glance at sacred history, as it begins in the Old and ends in the
New Testament, readily shows us what this means. In the Old Testament, where the
history of revelation constantly goes together with the history of the people,
the outlook is no doubt on the whole turned outward. The conception of the
affairs of the Jewish people combines naturally with a corresponding conception
of the foreign nations, which from so many sides intervened in this people's
history. In this respect there is no lack of descriptions either in Moses or in
the prophets. But of a calm, comprehensive survey of world-historical relations,
of a penetrating rational understanding of the position and significance of the
individual kingdoms for one another and for the development of the human spirit
as a whole, there is here no thought.

We must, as regards the narrowness of the horizon, make a critical separation
between the limitation that lay in the passionate particularism of the Jewish
people and its lack of humane cultivation, and the limitation that lies in the
religious principle itself. It is a task for the critique of knowledge to find
out what is historically factual in the Jewish legends and descriptions; but it
is the task of the critique of faith, by consistently holding fast the
religious principle, to bring to light what is characteristic of revelation in
the Old Testament worldview. As regards the New Testament, the particularistic
worldliness that obscures the view of what is foreign does, to be sure, fall
away, as the eternal light of the Spirit rises over the world; but then at the
same time the gaze becomes so turned inward, the principle of revelation so
overreaching, and the view of the world as a whole so otherworldly, that the
opposition between the rational and the antirational appears in full.

% ---- printed p.226 (PDF 239) ----
\opage{226} The religious worldview is and will be antirational. This is seen from the
principle, the principle of freedom. The humane view of the world, too, posits
freedom as the goal of development. But when the opposite of freedom, unfreedom,
is brought forward, the character shows itself. For the humane consciousness
unfreedom is only a naturally necessary boundness, a provisional limitation of
reason, and liberation a constantly advancing civilization. For the religious
consciousness unfreedom is sin; ``sin is the ruin of the peoples,'' it is
bondage itself and the cause of bondage: liberation, then, is not a
self-emancipation of the human spirit but a work of God, a redemption. Under this
point of view, decisive for faith, the actual world order must be regarded. What
religion teaches us about the world-ideals we cognize from what it teaches us
about: \emph{the fallen race}, about \emph{the gods of the world}, and \emph{the
prince of this world.}

\greekrun{α}{The Fallen Race.}

That the history of the human race must have begun with violent outbreaks of the
selfish forces of individuals can, in good agreement with what still goes on in
the midst of the world of culture, be understood without revelation. Heathen
popular legends, too, tell of bloody strife, of violence and murder from the very
first beginning. What is peculiar to the Hebrews' cycle of legends, on the other
hand, is this: that the revealed word does not stop at the natural side of the
crimes, but lays bare sin as sin and shows how the original relation to God has
been disturbed. Seen from the side of the ideal, death comes into the world as a
punishment for sin; seen in actuality, death comes into the world as an effect
of sin: the first death is a murder, the first murder a fratricide.

It begins with a religious act; Cain brings Jehovah his offering, just as Abel
does. But how selfish
% ---- printed p.227 (PDF 240) ----
\opage{227} Cain is, to what degree the serpent principle has poisoned his mind, comes to
light when the act of offering is at an end (Gen.\ 4:4). In envy of Abel and
resentment against Jehovah the smoldering egoism works its way forth in Cain on
two sides: the dark look, the downcast countenance betray the glow that burns
within him. First the relation to God reacts through conscience: ``Why art thou
wroth, and why is thy countenance fallen?'' Beware, beware! that is the
admonition, and in the admonition sin is brought to consciousness as sin. But so
strong is the selfish essence in self-will that the resistance which the
admonishing word has set against the sinful thoughts in order to drive them back
serves precisely to heighten the tension with which they suddenly break forth,
and the crime is done. Then comes the after-effect within; God's voice is heard:
``What hast thou done? The voice of thy brother's blood crieth unto me from the
ground.''

It is, in the character of actuality, a powerful manifestation of the Fall. The
curse is repeated, but in such a way that it now falls, not on the symbolized
principle, but on the head of the guilty one; and besides, at the same time, in
such a way that in actuality it strikes him in the inner, not in the outer, hits
him in his relation to God, not in his relation to the world. The cognition of
the trespass brings Cain --- not to repentance, but to despair: ``My iniquity is
greater than I can bear.'' And yet, seen outwardly, it is as though the despair
could pass and the deed be pushed back into oblivion. Cain must flee, that is
unavoidable; but in fleeing from home he flees after all in a way from the crime
as well. So it passes over. The thundercloud of judgment drifts past, the storm
of the spirit subsides; Jehovah sets a mark upon Cain.

If we regard Cain from the point of view of world-consciousness, he is more than
a fugitive criminal, a common murderer. There is in this figure something
mighty, for
% ---- printed p.228 (PDF 241) ----
\opage{228} he is marked by the spirit; something that arouses a shudder, for he is
branded. But the very mark he bears on his forehead is like a coat of arms with
which he makes his way. Everyone who meets him at once gets out of the way: so
he presses through and breaks himself a path. If we take the childlike account
immediately according to the letter, one could almost be tempted to think that
Jehovah had come to terms with Cain, had settled with him on easy conditions, so
that Cain should be free of further punishment if only he ``went out from the
presence of Jehovah'' and from his land. But if we attend to the thought of
revelation, another light falls on the whole.

In the land ``on the east of Eden,'' where Cain settles, everything goes as he
wishes. The family grows; Cain builds a city and becomes the father of a numerous
race, as vigorous and active as it is clever and inventive. Tillage, fixed
property, crafts and arts proceed from this race; the life of culture begins:
there is work in the workshops, there is ``hammering in brass and iron''; one
hears ``harps and flutes.'' Amid this busy life, rich in exertions and
enjoyments, the thought of Jehovah is wholly pushed back; so far is the
obscuring of the relation to God from being a hindrance to the development of
the forces and the successful progress of the works that precisely the opposite
is the case. The hidden despair, the secret gnawing of conscience, contains
precisely a spur to outward-going activity; for outward-going, grandiose activity
is a powerful means of stupefaction, a means that the ever more reflective
selfishness is even able to employ until the self is completely secularized and
its original relation to God extinguished.

In this respect Lamech is a truly remarkable figure beside Cain. For Cain's
worldliness despair furnishes a hidden foundation; for that the mark of Cain
should let itself be effaced by work is not to be thought of.
% ---- printed p.229 (PDF 242) ----
\opage{229} With Lamech it is another matter. He does not seem to carry any dark secret.
His essence is worldliness, but a worldliness in which, without any thought of a
relation to God, he feels quite at home. Outward power and personal repute are
his gods. He
% print: "valgt vig Kain" ("vig" for "sig")
has chosen Cain for himself as his exemplar, but enjoys at the same time the
satisfaction of far surpassing the exemplar. That Cain is a mighty man, a man
who can defend himself and take retaliation if anyone dares to injure him,
Lamech concedes; but Lamech --- with such sons and such wives: ``Adah and
Zillah, hear my voice!'' --- Lamech himself is after all quite a different man.
If Lamech murders man or boy, he can defend it; but let anyone come near him
himself and try to insult him, then he shall learn what vengeance is: ``For
sevenfold shall Cain be avenged, but Lamech seven times seventy times.'' In
measuring himself against Cain, Lamech has naturally forgotten the religious
measure; it does not occur to him at all that anything can be said against his
view of life: that a man avenges himself when he has the power to do so is,
after all, quite in order.


The fallen race is a race of Titans. When the consciousness of God is
darkened while great divine powers nonetheless continue to stir in
individuals, the mythical representation of a natural mingling of the
divine and the human comes to the fore. What this means is clear from
the heathen, and especially the Greek, myths of the heroes. Such myths
naturally find no place in sacred history; and yet there is here a
trait of the mythical-heroic which must be brought out in particular,
to shed light on the religious worldview. ``And it came to pass, when
men began to multiply on the face of the earth, and daughters were
born unto them; and the sons of God saw the daughters of men, that
they were fair, and
% ---- printed p.230 (PDF 243) ----
\opage{230} they took them wives of all which they chose \dots\ There were giants
in the earth in those days, and also after that, when the sons of God
came in unto the daughters of men, and they bare children to them; the
same became mighty men which were of old, men of renown'' (Gen.\ 6:1--4).

In reading of this race of giants we receive an impression somewhat
like the one we have when we hear Nestor in Homer begin his speech to
Achilles and Agamemnon by reminding them that he had lived with
stronger men, with a Caeneus, Dryas, Pirithous, Theseus, men ``with
whom none of the mortals now living would venture to enter into
contest.'' That the sons of Elohim in the ancient Hebrew legend have,
however we may interpret the narrative, a mythical tinge, cannot be
denied. But this hint of the mythical, far from obscuring the
principle, serves precisely to set the opposition between myth and
revelation in the clearest light.

With the Greek, whose consciousness is essentially mythical, the
legends of the heroes give the imagination rich material for detailed
elaborations and ingenious poetic inventions. In the Hebrew legend
there is no trace of anything of the kind; what is said here is said
briefly; for what is really to be emphasized is that ``the wickedness
of man was great in the earth,'' and that it was this wickedness that
aroused the wrath of Jehovah: it is the offense against the Holy One
himself on which the stress falls. The Greek legends, too, tell of the
wrath of Zeus and of the flood in which he let the human race perish,
save for a single pair. But the difference between the wrath of
Jehovah and that of Zeus is like the difference between Jehovah and
Zeus. The ground of Zeus's wrath is not the violation of a holy will,
but an equivocal resentment, because the Titan Prometheus, and not he
himself, had kindled in man a spark from heaven. Geologically
considered, the Noachian flood has, to be sure, no greater extent than
the flood of Deucalion; but its significance is other from the
religious standpoint than
% ---- printed p.231 (PDF 244) ----
\opage{231} from the mythical. Mythically, the catastrophe ends with a new race of
men springing up, according to the oracle's answer to Deucalion and
Pyrrha, from the hard stones of the earth, a hard, unbending, defiant
race, \textit{et genus durum experiensque laborum;} whereas according
to the revealed word it ends with a renewal of the relation to God,
with a holy promise: the rainbow is seen in the cloud, and Jehovah
makes his covenant with Noah.

The pictures of life that the religious history of the people unfolds
here are to be regarded not merely as counter-images but also as
after-images of the archetype corresponding to the ideal. From this
standpoint too the thought of faith shows itself to be a thought of
life running through the whole. When man, by crimes that cry to
heaven, breaks with the Holy, turns away from God and gives himself
up to the movement of the world, the religious principle reacts; the
Spirit is angered, and the wrath is revealed in the punishment. But
what is punitive in the punishment is not felt, because of the
stupefaction; selfishness, the egoism inspired by the principle of the
serpent, renews its outbreaks with growing force over an ever wider
compass. At last the Spirit withdraws; it is as if God at length had
to give up, as if the governing hand were too weak. Already as a
prelude to the Flood, Jehovah says: ``My Spirit shall not rule among
men for ever'' (v.\ 3), and (v.\ 6) it is said expressly: ``It
repented Jehovah that he had made man on the earth, and it grieved him
at his heart.'' However figurative this language may be, it is
nevertheless not mythical; what is stressed here is not a divine
rashness but the holy wrath of the Spirit; and the wrath of the Spirit
is, in the sense of revelation, sufficiently grounded.

When actual dispensations, however hard, are no longer conceived
religiously, and so also cannot be understood as judgments of
punishment upon sinners, then
% ---- printed p.232 (PDF 245) ----
\opage{232} --- this is the condition of freedom --- the self-willed race must be
left to the consequences of the presuppositions that it itself posits
in its spiritual darkness. That is the one side of the matter; the
other is that there are always found in the fallen race individuals
who are receptive to the workings of the Spirit, and in whom the
relation to God can be renewed. What happens with Noah and his family
at the Flood is a foreshadowing of what will later happen with
Abraham.

The race of Noah too, which grows and multiplies after the Flood, is
subject to the workings of the principle of the serpent; it too tries
its hand at storming heaven; it too arouses the wrath of the Spirit.
But as the Mosaic legend leaves the general primeval history of the
human race and is about to pass over to telling of the patriarchs, an
event is presented which is characteristic where revelation is
concerned.

``Is the whole human race descended from one pair, or does each of the
different races have its own origin?'' The treatment of this question
again recalls the difference between faith and knowledge. What
importance the question has in science may be seen from the dispute
between the polygenists, who derive the origin of the race from
several different pairs, and the monogenists, who explain the
differences found among the human races partly by ``heredity'' and
partly, and chiefly, by the ``medium,'' counting under the medium not
only the physical conditions --- climate, humidity, heat, means of
nourishment, etc.\ --- but also the intellectual and moral
surroundings. This matter will, now that science has once emancipated
itself from the authority of the Bible, be pursued and decided without
the intervention of religion. But just as science, by virtue of its
empirical-rational character, is entitled to carry through its
investigations without letting anything be prescribed to it by faith,
so the philosophy of
% ---- printed p.233 (PDF 246) ----
\opage{233} religion, by virtue of the antirational nature of faith, is entitled to
conceive the matter from its own standpoint without consulting the
science of reason. ``And the whole earth was of one tongue and of one
speech''; it is the building of the tower of Babel that is alluded to
here.

One of the most magnificent phenomena of worldliness is the
cooperation of peoples in multitudes; to bring forth by common labor
an astounding great work, a uniting center, a monument that could
defy the vicissitudes of the earth and raise the names of its master
builders to heaven --- that was the ideal the peoples in the vale of
Shinar set themselves to realize; with this lofty goal in view they
began to build the tower of Babel. Jehovah's displeasure at this is
depicted in a way that undeniably makes it easy for negative critique
to explain the whole as a myth.

For one sees plainly, critique will say, how Jehovah, impressed by the
high tower, actually begins to grow afraid of the growing might and
power of men. ``Behold, the people is one, and they have all one
tongue; and this they begin to do: and now'' --- if it goes on like
this --- ``nothing will be restrained from them, which they have
imagined to do'' (Gen.\ 11:6). To set a dam against the swelling
storm-flood, the embittered Jehovah, jealous of his sole rule, knows
no better counsel than to ``go down and confound their language, that
they may not understand one another's speech.'' This measure,
moreover, is quite profoundly devised. The ``Almighty,'' anxious for
his over-power, could easily, to be sure, have destroyed both the
tower and the city by lightning and earthquake; but what good would
that do? So long as the human race held together, and the
high-striving multitude formed a closed mass, they could easily begin
upon another, still more magnificent and dangerous undertaking; by
confounding the languages,
% ---- printed p.234 (PDF 247) ----
\opage{234} bringing misunderstanding into the crowd and scattering its forces, the
governing wisdom succeeds, on the other hand, once and for all in
setting a limit to the presumptions, and moreover in doing it so
cleverly that it looks for all the world as though the tower and the
plan collapsed of themselves: \textit{divide et impera}, that is the
policy of governance, a golden rule for all who govern. What is
negatively edifying gains still more weight when, on the
presupposition that we have here an actual event --- which is
scarcely the case --- one goes more closely into the connection of the
matter and sees how the supernatural governance turns the natural
causes upside down. Seen from the standpoint of actuality, every human
undertaking, be it ever so magnificent, has its essential limitation
in itself. By overstepping the limit, ambition only betrays its
impotence and sets a limit to itself: this holds of all high-flying
building plans, the Babylonian one included. The top of the tower was
to ``reach unto heaven''; that was an exaggeration, and the
exaggeration fell before the limit. The whole ambitious race imagined
itself to be agreed on one single giant work of honor and renown; that
was an illusion. The higher ambition that unites men for great
exploits is inseparable from the lower, which smolders in individuals
and calls forth discord and confusion. The actual causal connection,
grounded in reason, is therefore this: in its youth the human race
had a swelling energy for deeds, united with a gigantic,
heaven-storming imagination that drove it on into the unlimited. But
exaggerations and presumptions have served precisely to make the
indwelling limit of reason manifest. The point is not that the
multitude of people on the earth would have been agreed among
themselves to build the infinite tower if the divine will had not
caused the disagreement by a miracle; the point is rather that the
germ of disagreement
% ---- printed p.235 (PDF 248) ----
\opage{235} lay hidden beneath the apparent agreement, and that its breaking
through satisfied the divine reason. The miracle is not the cause of
the confusion of tongues, any more than the many different tongues are
the cause of the dispersion of the peoples; but the reverse: the
dispersion of the peoples, the original and natural diversity of the
tribes, is the cause of the diversity of tongues, and the diversity of
tongues is the most essential cause of that separation and dividedness
which, by preventing the masses from running together into one single
great uniform crowd, rationally conditions an all-round development of
powers and dispositions.

The rational interpretation of negative critique is, from the
standpoint of knowledge, irrevocable. Were the religious conception to
go in a rational direction, then in working against critique it would
only be working against itself. But the difficulty vanishes when the
word of revelation is allowed to be what it in truth is: antirational.
The antirational is recognized at once by the inverted causality.
Suppose, to go to the extreme, that the polygenists are right, that
natural actuality is therefore as opposed as possible to the word of
revelation: what follows from this? That religion is nonetheless
proved right in the spiritual sense! We already have the suggestion of
it in the humane view. For let the marks of race be ever so
prominent, and the roots of the great families of languages ever so
dissimilar, one thing still stands firm: that man in all races is
man, a rational being disposed for self-conscious development.
Rationally seen, the human race is but one.

This essential unity, however, receives a quite different and higher
significance through the religious presupposition that man, as man,
is created in the image of God; before this original stamp of unity
of the spirit, every natural difference must be regarded as vanishing.
While
% ---- printed p.236 (PDF 249) ----
\opage{236} Judaism otherwise, by its sharp opposition to heathenism, expresses
itself particularistically, as regards the original equality of men
before God it not only lets the whole race descend from one pair of
parents --- for an analogue to this is found also in several
mythologies, for example in the Norse --- but also holds to its
presupposition so steadfastly in thought that the most essential
consequences of all are later drawn from it.

Now when religion sees in the human race one single great family from
the beginning, and yet wants to explain the fact that this family has
in the course of time become estranged from itself, has divided into
tribes and nations that do not understand one another's tongues, then
the explanation must be conditioned by the inversion of causality.
The preceding dispersion is then regarded not as the cause of the
confusion of tongues; rather the confusion of tongues, a consequence
of discord in mind and thought, is regarded as the cause of the
dispersion of the peoples, and so as a punishment. Inner discord,
confusion and division belong precisely to burdensome actuality, and
every burden of actuality must be borne as a punishment for sin.

The particular guilt corresponding to the particular punishment is
made visible in the presumptuous attempt at building the tower. On
this point too causality is inverted. Rationally, the cause of the
enterprise's coming to a halt lies in its own unreasonableness, and
the divine displeasure is to be regarded as an effect of this
unreasonableness; antirationally, on the contrary, the cause of the
halt lies in the divine disapproval, and the effect of this is the
unreasonableness, the infinite confusion, in which the enterprise
dissolves. From this inversion of causality it is seen at the same
time that the antirational does not stand in logical conflict with the
rational, but only signifies the absolute heterogeneity of the
religious principle with the scientific.
% ---- printed p.237 (PDF 250) ----
\opage{237} The natural intermediate causes are by no means displaced by the
supernatural cause of will; but the latter appears as immediately
effective in that the former are, at the decisive turning point,
reduced to vanishing moments. ``The tongues are confounded, and thereby
God's will is done; for the fallen race in its presumption wills
something sinful, and the confusion of tongues is a punishment for
sin.'' In this way the antirational can indeed, in all ambiguity, be
combined with the rational. But revelation has no liking for the
ambiguous; in its language it means expressly to remove the ambiguity;
that is why, passing over the intermediate causes, not rhetorically but
categorically, it emphasizes exclusively the divine cause of will
wherever the thought of governance demands it.

From this it can be understood that the effect which in one connection
is referred to a natural intermediate cause may, in another connection,
from another standpoint, have to be referred to God's will. The
difference between God's effecting will (\textit{voluntas efficiens})
and his permitting will (\textit{voluntas permittens}) is, in the form
of knowledge, meaningless, or rather: contradictory; but for faith it
is essential and significant: the difference is antirational, not
rational. We see it from the confusion of tongues.

In what sense the confusion of tongues or the multiplicity of
languages, regarded from its shadow side --- the essential significance
of the abiding
differences is not under discussion here --- is a punishment for sin
becomes evident when we look to the final end. For the final end is
that division, enmity and confusion shall vanish; the final end is
that the God of Israel shall be the God of all peoples; the final end
is that the middle wall of partition shall be broken down, and that
men, despite their many dissimilar thoughts and purposes, shall
understand one another in the understanding of the one purpose that is
exalted in the eternal sense. Just as the race is one by its fall, so
shall it
% ---- printed p.238 (PDF 251) ----
\opage{238} also be one by its restoration; it is the governing causality that
scatters, and it is the governing causality that gathers; but in both
respects the causality is antirational: it wills not only the reunion
on account of the dispersion, it wills also the dispersion with a view
to the reunion; it wills not only the restoration for the sake of the
fall, it wills also that the fall become manifest for the sake of the
restoration. ``For God hath concluded them all in unbelief, that he
might have mercy upon all'' (Rom.\ 11:32): the fallen, falling race
falls not outside but inside the Kingdom of the Father.

\greekrun{β}{The Gods of the World.}

It lies in the nature of the world to be ambiguous: it is dependent,
for it is created by God, and yet independent, for it has its ground
in itself; it is dependent, for it is preserved by God, and yet
independent, for it is borne and moved by its own powers; it is
dependent, for the aim and final end of its movements are determined by
the governing will, and yet independent, for aims and means develop out
of its original disposition. The ambiguity doubles and grows when the
gaze fixes on the essential center, the human self, which with its
spiritual tasks is placed precisely midway between God and the world.
The ambiguity of the world is now not merely outside man but in man,
for the human self begins in natural ambiguity. If the relation to God
could simply be strengthened by the weakening of the consciousness of
the world, then the choice between God and the world might well become
recognizable enough; but the difficulty is that the relation to God
becomes empty when the content of the world falls away; the difficulty
is that the self which is to renounce the world must precisely take a
vigorous hand in the world, work in the world against the world, and
so know the world both in itself and outside itself. The difficulty
increases when we consider that
% ---- printed p.239 (PDF 252) ----
\opage{239} the self, developing under great impressions of the world, must take
the first step on the way of the spirit with the acknowledgment of
guilt, and the next with decisive self-denial. To acknowledge one's
guilt when one has already, beforehand, in the finite sense, with full
freedom done something and made oneself guilty of something, is
difficult enough already because of the stupefaction; but to have to
begin with the acknowledgment of guilt --- not in this or that, but of
essential guilt --- in order to be able to begin with God, is yet far
more difficult. To deny oneself a good in the world, be it this or
that allurement, is already hard enough; to force oneself to infinite
resignation and to acknowledge all relative goods as indifferent
(\textit{Adiaphora}) is harder still; but to resign everything in the
world in such a way that infinite resignation means a new beginning
with the world, an unconditional readiness to use the world and do
God's will in the world --- this task is, as it seems, so superhuman
that reason's eyes swim when it rightly considers it. Religious
self-denial is, spiritually understood, certainly the very deepest
self-affirmation; but the living cognition of this can in actuality
only arrive when the self-denial has taken place. That the cognition
which should come first, in order to bring about the self-denying
act, comes only afterwards, is an expression of the originality of the
relation to God and of the inner infinity of faith; but a direct
transition from the relation to the world to the relation to God
there cannot be.

From what has been set out here it can be understood that there must
be a conflict between reason and religion about the essence and
significance of the world, a conflict that can be settled only in
consequence of the absolute heterogeneity of the principles.

For a rational consideration, the actual world is rational. The
rational is seen in the laws of things, i.e.\ in the unity of the
individual with the universal, for the things represent the
individual, the laws the universal. The rational
% ---- printed p.240 (PDF 253) ----
\opage{240} is that the universal is higher than the individual. Now if the world as
a whole is rational, then man, the human self, ought surely to be
rational too; but it becomes rational only by cognizing itself, as an
individual being, to be the universal, and by making itself an organ
of the universal. This instinct of reason --- to let the universal
raise itself step by step and master the individual --- underlies the
formation of great national communities, the formation of states,
history and civilization. Science has its strength precisely in
investigating, developing and presenting these many formations in
their characteristic diversity and mutual connection, constantly, in
the midst of particulars, with its gaze fixed on the universal. The
rational self, which is to work in all directions for the universal
and therein to cognize and recognize its own essence, finds in manifold
suffering and acting rich occasion both for self-denial and for
self-affirmation; but the question is how far it can go, how high it
can reach. It is the gods of the world that must lead the way in this
connection. The fundamental significance of these gods changes
remarkably according as they are regarded in the light of revelation
or in the light of reason. From the standpoint of revelation they are
all without distinction idols, a Zeus and Athena just as much as a
Moloch, Baal and Astarte, visible signs of outrageous crime. The crime
is that the false gods are set in the place of the true God. Conceived
from the standpoint of humanity, on the other hand, the mythical
deities are bearers of great world-ideas. The figure of the god itself
has the higher determination of being the protecting symbol while the
Idea works its way through to actual existence and deposits a form for
the universal.

The highest universal, in the form of spiritual actuality, is the
state. The state as state is the supreme power, the ruling deity on
earth; the state is the great projected human essence; it is the
% ---- printed p.241 (PDF 254) ----
\opage{241} reflected splendor of the glory of reason, the express image of its
essence. The highest position in the highest universal of actuality is
the position of the ruler in the state; only here does the rational
self find its proper determination, its true sphere of activity: a
magnificent prospect, a wide horizon! The universal, however, cannot
coincide so immediately with the individual that a single one should
have all the power and all the others be slaves. This form of state is
only a first crude sketch. The task is that power be distributed, and
society organized in such a way that all, each with his limitation and
according to his titles, can be given a share in power and so far also
become holders of power. If only this task were solved, the state
would be the temple of the deity, or rather the kingdom of the deity,
and the deity present in its kingdom would then, as the universal and
by virtue of the universal, permeate both arts and sciences and the
customs of the people and the activity of the citizens and the ethical
consciousness of the individual: this is an ideal of the world whose
realization must be striven for through millennia, and which can yet be
reached only by approximation.


With the world-ideal set out here, religion cannot possibly be in
agreement with reason. It does not lie in the principle of faith that
religion should either deny the state or in any respect fail to
recognize its actual significance; but the deification of the state and
of power that springs from the principle of reason is what repels the
consciousness of God and makes it comprehensible that sacred history,
far from engaging in any pragmatic presentation of the history of
states, rather holds fast to its antirational standpoint, protests
against the gods of the world, and holds judgment over those in power.
In order to consider what significance this has from the standpoint of
faith, we choose a couple of examples from the prophets.

Isaiah sees in a vision the fall of Babylon (ch.\ 14) and describes the
pride of the fallen Chaldean king. ``And thou
% ---- printed p.242 (PDF 255) ----
\opage{242} saidst in thine heart: I will ascend into heaven, I will exalt my
throne above the stars of God, I will sit upon the mount of the
congregation, in the sides of the north.'' This description of a
tyrant ``who smote the peoples in wrath with a continual stroke, who
ruled the nations in anger'' is religiously full of pathos, aglow with
holy passion, but in its fundamental character it is not in the least
exaggerated. In the Babylonian ruler self-deification goes so
essentially together with the deification of power that the egoism of
the ideal shines through. It is not the conceit, not the boastful words
of a high-flying imagination, on which the emphasis falls; such
outbursts of an overflowing Oriental imagination, which point rather to
a state of spiritual drunkenness than to a methodical lust for rule and
deliberate wickedness, cannot heighten the imputability; no, what is
outrageous is the egoism of the ideal itself. This outrageous thing the
critique of reason cannot strike, since it is itself stricken by the
egoism that idealizes itself through the universal. The reflecting
world-consciousness is, to be sure, not selfishly naive like an Eastern
despot, but it nevertheless bows deeply before outward greatness and
practices the deification of power with every Nebuchadnezzar.

Religion does not begrudge the ruler his power, but it punishes the
encroachments of ambition and judges the deification of power. Religion
is the only authority that can do so, since it is the only one that
sees through the self, the only one that has a measure of height for
power. If the holy will is thought away, then the egoism of the ideal is
entirely justified; if the almighty Self goes out, then the selfish
deification of power is in order. What reason cannot do, the revealed
word can: step forth against presumptuous greatness and make history a
mirror of the impotence of the mighty. It is not selfish hatred, it is
the authority of the Spirit, it is Jehovah's judgment, that breathes in
the mighty words with which the prophet describes
% ---- printed p.243 (PDF 256) ----
\opage{243} the annihilation of the fallen prince. He descends into the realm of
shades, the pale heroes rise, the dead kings stand up from their
thrones, and trembling they speak: ``Art thou also become weak as we?
art thou become like unto us? Thy majesty is brought down to the
underworld, the noise of thy harps. Now shall the worms spread under
thee, and the maggots cover thee.'' --- True enough, a gifted poet
could give an equally shattering portrayal; but there is this great
difference between the poet and the prophet, that the former only
affects the imagination, while the latter, by virtue of the principle
of faith, asserts the claim of the Spirit and affects the will.

``World history is the world's judgment''; when these words are
understood scientifically, they are understood objectively: there are
in history crises and catastrophes of which we can say that each of
them has been like a day of judgment over a given age; but the judgment
is only relative and ambiguous, for the Judge of the world does not
appear. There is a picture in Daniel (ch.\ 5) that makes really vivid
what it means that a day of judgment draws near. It is not ambition
strong in deeds that swings the scourge over sighing nations; it is
another form of glory that here comes to revelation: luxury, spiritual
slackness, stupefying enjoyment and contempt for the holy have with
noisy festivity gathered a circle around the gods of the world. It is
King Belshazzar, who makes a feast for his lords, ``a thousand in
number,'' and he drinks wine before the thousand. --- ``And when he had
tasted the wine, he commanded to bring forth the vessels of gold and
silver which his father Nebuchadnezzar had taken out of the temple in
Jerusalem, that the king and his lords, his wives and his concubines,
might drink therein \dots\ And they drank wine, and praised the gods
of gold and of silver, of brass, of iron, of wood and of stone.''

% ---- printed p.244 (PDF 257) ----
\opage{244} The critique of reason will surely not, as regards the total
impression, find a single feature in this picture so rich in character
that goes against probability. On the contrary! Here there is, if not
exactly clear signs, then at least likelihood that a kingdom under such
governance, with such a governor, might be near its downfall. In
actuality it will depend on the circumstances. Should the mighty kingdom
suddenly be shaken and collapse, the whole upheaval would naturally be
comprehensible from causes, conditions and intermediate causes which
clearly show that the governing power is the world-reason itself. For
that a God who felt offended because the guests drank from the holy
vessels of his altar should intervene as an arbitrarily acting cause is
mythical. --- But what here, for the rational world-consciousness, is
mythical and in that respect inessential, the word of revelation
expressly posits as what is in the highest sense essential, decisive;
religion closes its eye to the intermediate causes in order to see the
holy cause all the more clearly: religion is antirational; that is seen
in the moment of tension.

``In the same hour came forth fingers of a man's hand, and wrote over
against the candlestick upon the plaster of the wall of the king's
palace: and the king saw the hand that wrote. Then the king's
countenance was changed, and his thoughts troubled him \dots'' And none
among the wise men of Babylon on whom the king called, neither the
enchanters nor the Chaldeans nor the soothsayers, could read him the
enigmatic writing. The enigmatic writing is a word of judgment which,
often repeated, still repeats itself each time that the deification of
power and of self is to come to an end. \textit{``Mene, Mene, Tekel,
Upharsin''}: so it is often written with fingers of an invisible hand;
but it is seldom that the writing is read aright; for only a prophet
can read it. --- ``Is this, then, actuality?'' Yes, actuality of the
spirit, actuality for the spirit.
% ---- printed p.245 (PDF 258) ----
\opage{245} \greekrun{γ}{The Prince of This World.}

If the world is an ambiguous being, it goes without saying that the
world-spirit must be an ambiguous spirit. Expressions such as
zeitgeist, the spirit of the time, party spirit, national spirit and
the like are reminder enough of the ambiguity. In the science of
reason, however, the category world-spirit is used as the designation
for the human spirit developed to full consciousness, a spirit that has
its purposes and its final end in itself, and it is presented as the
deity of the Absolute revealing itself in history. That the
world-spirit, so conceived, is more than a fantastic personification of
an essentiality indeterminable in itself, that it actually corresponds
to a determinate concept, can easily be shown. The Idea in the form of
consciousness is spirit. Now, society as society, the people as people,
humanity as humanity can by no means be said to have unity of
consciousness in itself, or in a literal sense to be conscious
of itself: the unity of consciousness is distributed among the
individuals singly. But as the individuals are awakened to the
cognition of the essence of society, they are at the same time
awakened to the cognition of the universal power that governs society.
Out of the many individual consciousnesses a common consciousness forms
itself, and this common consciousness, penetrating all individuals,
then transforms itself for reflection into a consciousness of the
universal, an essential universal consciousness with retroactive
objectivity. What content the universal consciousness is to receive
depends, moreover, on the circle of tasks, interests and Ideas which
the age sets in motion each time. The spirit impregnated with Ideas
acts on the individual wills as a universal will moving the whole; in
great turning points the universal spirit can even at times rise up
with so decisive a stamp of objective unity of will that it gives the
impression of a willing self, an invisible
% ---- printed p.246 (PDF 259) ----
\opage{246} ruler. And yet this spirit ruling in the times is not a willing
self; it is only a general selfhood, a selfless willing; this selfhood
and this willing are, with all their objectivity, only the mirroring
and reflection of what moves subjectively in the individuals. The
world-spirit in its concept is at once both objective and subjective,
but the unity of the objective and the subjective is an ambiguous
unity; the ambiguity is mythical: the world-spirit is a mythical being.

Now, as religion sets its holy principle of will sharply against the
principle of the world, it is a matter of course that the struggle must
become hard. A mutual hatred develops: love of the world is hatred of
God, and conversely; this conflict is so thoroughgoing precisely
because it goes through the innermost core of the self. The self moved
by the world-spirit must, in order to remain true to its purpose, be
willing to sacrifice itself for an Idea, and above all for the most
comprehensive subjective Idea of the egoism of the ideal, the Idea of
honor; for just as the Idea of power is the most objective, so the Idea
of honor is, among all the world's Ideas, the most subjective: if the
self cannot attain power, then it must at least preserve its honor. Now
if the goal of honor were unconditionally inward, the claim of honor
would be so inward, so absolutely decisive a claim of selfhood that
the self would therein comprehend its own eternity and posit itself and
its own glorification as the final end of existence. That such a
self-deification must be regarded by religion as rebellion against God
is a matter of course. But that the world-spirit, too, does not
tolerate a deification of this nature is then also evident. For it is
not the individual as individual, but the universal revealing itself in
the totality, on whose development and progress everything depends. If
the self that works for the Idea is to come to a true understanding
with the world-spirit, then it must unconditionally live for the
universal, comprehend the universal, not in metaphysical abstraction,
% ---- printed p.247 (PDF 260) ----
\opage{247} but in actual existence as power; it must deify power. Over against
actual power the self as self is a nothing. Even if at a given time it
comes to power and, through the union of power with honor, can unite
the deification of power with self-deification, the fundamental
contradiction is nevertheless only veiled, but not resolved: in the
moment of crisis the contradiction breaks forth, and the self sinks
together into nothing. When religion now comes forward against this
with its urgent admonition: ``Love not the world, neither the things
that are in the world! If any man love the world, the love of the
Father is not in him'' (1 John 2:15), then we can understand that it is
not a dark ascetic failure to recognize the validity of the world's
Ideas and the humane tasks of actual life, but care for souls, concern
for the self, that breathes in this admonition.

So long, then, as we remain standing at the world-spirit in its mythical
ambiguity, it will not be difficult to follow the religious train of
thought. But now revelation goes a step further. In order to unveil
the symbolized serpent-principle and illuminate the darkness in which
evil keeps its origin hidden, it leads the egoism of the will back to
the evil will, to the egoist himself in supernatural magnitude.

The evil will in its highest unity of consciousness is the evil self:
the evil will is spirit, the mightily ruling spirit of the world-spirits
--- ``the prince of this world''! Here the antirational is once more
conspicuous; and here again is a summons to examine the relation between
what can be ascribed to the immediacy of representation and what
essentially follows from the principle.

Religion often speaks in images of the universal powers; the psalmist
sees in the storm winds the angels of Jehovah (Ps.\ 104:4), and the
water in the pool of Bethesda receives, according to the evangelical
explanation (John 5:4), its healing power from an angel's going down at
certain seasons into the pool. If it
% ---- printed p.248 (PDF 261) ----
\opage{248} is now unmistakable that personifications of both good and evil
powers occur in Holy Scripture, what certainty have we then that
everything said about the spirit of darkness is not likewise to be
understood figuratively, so that expressions such as ``Satan,'' ``the
devil,'' ``the god of this world'' and the like are to be taken as
personifications of the essence of wickedness, of evil in the world,
but by no means to be taken literally of a personal evil spirit? What is
said in the New Testament about the devil, that he falls from heaven
like lightning and yet desires to sift the disciples as wheat, that he
is already bound and yet goes about as a roaring lion, and so on, is
actually so vague and figurative that Schleiermacher must seem to be
right when he\footnote{``Der christliche Glaube''. I. Band.
Berlin 1835. S.\ 209--22.} rejects the doctrine of the devil as not
belonging to the Christian doctrine of faith. And yet even the most
acute explaining-away of the devil's personality will not satisfy the
Christian consciousness. ``It is clear,'' says Martensen, ``that the
statements of Scripture about the devil, about the \emph{enemy}, also
speak of \emph{more} than a principle, of more than a universal evil
\emph{willing}, namely of an actual personal \emph{will}, although
indeed now the one, now the other side, now the generally spiritual,
now the personal, may be the one that comes forward more strongly.
Christ's struggle with the devil is indeed a struggle against a
universal principle, but it first receives its full significance when
it is at the same time a personal \emph{will} that he repels and
overcomes.''\footnote{Den christelige Dogmatik.
S.\ 202--5.} So it is. Scientific critique does not tolerate the
devil; if he flees from knowledge, he is arrested by faith; if he makes
a sally from faith, he is attacked again by knowledge: so the devil
himself must keep guard on the border between faith and knowledge.

% ---- printed p.249 (PDF 262) ----
\opage{249} The spirit of darkness is essentially a tempting power; temptation is
the fundamental function on which its superiority rests; it attacks no
one except through temptation, it conquers no one except by temptation:
the devil is the tempter. But what is it to be tempted by the devil?
Every man is tempted ``when he is drawn away of his own lust, and
enticed''; that is the general psychological explanation of temptation;
it fits every human being as an individual in the race, every human
being burdened with sin and guilt; but from participation in sin and
guilt there are two ideal exceptions: pure innocence, which is not yet
spirit, and the sinless will, which is personally spirit. To neither of
these can the general psychology of temptation be applied. The innocent
one is indeed tempted by its desire, but not in the form of desire,
that is to say: not with consciousness that what tempts is precisely
the self's own impure desire; for in the same instant as consciousness
of this enters, the self, as we have seen, is impure and innocence is
lost. --- The sinless will cannot possibly be tempted by its own desire,
for it is tempted with full consciousness as spirit; were it now
conscious under temptation that it was being tempted by its own desire,
then it would indeed be conscious of its sinfulness and could not be
sinless. And yet the sinless will can and must be tempted, as surely as
it is subject to human conditions. Sinlessness is determined in
actuality by the sinless one's freeing himself from sin by excluding
sin, keeping sin away: sin will thus press in from without, i.e.\ from
another will, and must be kept out by an express act of will. Eve is
tempted by the serpent, but Christ is tempted by the devil; on these
two extremes attention must fix itself if we are to understand the word
of revelation about ``the tempter.''

In order to be tempted by spirit, man must himself be spirit; only in
the same measure as the individuality that
% ---- printed p.250 (PDF 263) ----
\opage{250} is tempted is conscious of its will of the spirit will it be able to
refer the tempting promptings of lusts, images, representations and
thoughts to the evil will; the less firmly the self is established in
the will of the spirit, the more it is either affected by the
immediately sensuous or, in the form of general selfhood, streamed
through by universal Ideas, the less will it be able to apprehend the
tempting principle as will. As regards sensuous allurements, the
individual must with a little self-reflection easily notice that what
tempts lies in desire; as regards beguiling promptings of Ideas and
false ideals, the gifted man with his strongly working reflection will
discover that what tempts really comes from dark possibilities, from an
obscure egoism of Ideas; he will cognize that he is tempted by spirit;
but it is still only the world-spirit in all its veiledness: the evil
will, the tempting self, does not appear. It is otherwise with Christ's
temptation; in this connection we shall reflect only on the tempter.

It is, from the standpoint of faith, highly characteristic that
Christ's temptation in the wilderness does not precede but falls at a
point of time that follows immediately upon his baptism in the river
Jordan. His messianic coming forward is thus a fully matured resolve;
he has the consecration; he is in the Spirit, absolutely concentrated
in will of the spirit. What in the world should now be able to shake
this holy, pure, resolved will? Nothing! What illusion, formed of
billowing air, of the juggling vapors of the wilderness, should be able
to confuse the thoughts of this resolved one and becloud this gaze that
sees through every deception? Nothing! And yet in the wilderness, out of
this nothing, in this waste, out of this emptiness, a tempter comes
forth with a temptation. It must then be a temptation of a peculiar
kind, a temptation to which rational psychology knows nothing
corresponding. To be tempted is, at this summit, not to be ensnared,
beguiled
% ---- printed p.251 (PDF 264) ----
\opage{251} and made to waver; but it is: in redoubled selfhood, as spirit, to be
attacked by spirit, so that thought stands against thought, will
against will, self against self. This redoubling in selfhood is the
deepest crisis of freedom. The tempting self is not Christ's self; far
from it! it is as opposed to the Holy One as possible. And yet Christ
takes up both the tempter and the temptation into his consciousness;
for if he did not do so, he could not be tempted. The tempter can only
``come to him'' because he, in his infinite superiority, lets him come;
and he lets him come for the sake of freedom, of revelation and of
salvation: only upon the tempter's fall can the Messiah found his
kingdom.

That Christ lets himself be tempted means that he compels the tempter
to reveal himself: it is the decisive defeat of the evil will that its
self-contradiction comes to light, that it becomes completely revealed.
From the standpoint of knowledge not only the tempter but also the
tempting word are to be brought under the category: myth; but in
consequence of the principle of faith the story of the temptation is
precisely antimythical. The antimythical is known by the clear
opposition of the tempter to the tempted, of self to self; through this
opposition the revealing word receives its significance in the course
of the temptation. The tempter makes his attack in such a way that each
utterance strikes precisely at the contradiction, the conceptual
contradiction which, from the viewpoint of a negative critique, lies in
this: to be God-man. ``If thou be the Son of God, command that these
stones be made bread.'' To hunger as a man and yet be as God, who by a
word can turn stones into bread, what a mythical contradiction! But
just as the contradiction is to have its tempting effect and bring the
tempted one into self-contradiction, the antimythical is heard: ``Man
shall not live by bread alone, but by every word that proceedeth out of
the mouth of God.'' But a contradiction it is nonetheless, a mythical
self-contradiction, to know oneself as the one who could command
legions of angels, and notwithstanding
% ---- printed p.252 (PDF 265) ----
\opage{252} to know oneself as the one who, over against the world's
misrecognition, contempt and mockery, is to be the powerless one who
cannot make himself a phenomenon, because he cannot for his own
glorification avail himself of supernatural means. ``If thou be the
Son of God, cast thyself down'' --- from the pinnacle of the temple ---
``for it is written, He shall give his angels charge concerning thee,
and they shall bear thee up in their hands.'' \dots\ If thou art really
the one whom the Father will glorify, then let the glorification begin
at once, that the world may see it and marvel. --- Against a mythical
glorification that bursts all the laws of actuality the Messiah now
sets the antimythical resignation: thou shalt not selfishly challenge
the supernatural power, ``thou shalt not tempt the Lord thy God.''

But it is nonetheless a contradiction, an intolerable
self-contradiction, that he should be the servant of all who knows
within himself that he is equipped to be the lord of all and to lay the
whole world under his foot. He who is as God in himself must surely
also have the right to rule as a God and to deify himself. The
self-contradiction should thus, if this is the tempting thought, find
its resolution in immediate self-deification. This turn of the matter
seems, however, at first glance to conflict with the text. According to
the story of the temptation (Matt.\ 4:8, cf.\ Luke 4:6)
% print: "Luk, 4, 6" -- comma for the abbreviation point
it is not, after all, immediate self-deification to which the tempter
urges the Messiah; what he literally urges him to is, on the contrary,
an infinite and, as it seems, utterly outrageous self-abasement. ``All
these things'' --- all the kingdoms of the world and their glory ---
``will I give thee, if thou wilt fall down and worship me.'' A more
foolish demand, then, in the opinion of critique, the devil could not
make; \textit{ergo} --- this is the scientific verdict of critique ---
not only the clothing of the temptation, but the temptation itself too,
is a myth.

The absolute heterogeneity of the principles, of faith and knowledge,
% ---- printed p.253 (PDF 266) ----
\opage{253} must surely here, if anywhere, be conspicuous. The story of the
temptation, which in the eyes of knowledge has driven the mythical to
its uttermost point, is, according to the principle of faith, precisely
to be capable of being said to have cast light into the innermost
recesses of selfhood, into those corners of the spirit from which the
mythical proceeds. It is the relation between the personal will and the
false, tempting world-ideal that is to be cleared up once and for all.
The temptation is at its height when the opposition of the tempting
ideal and the will is at its height. The temptation ends when all
circumlocutions fall away. Had Christ had the tempting principle within
himself, had he, like another Hercules, stood at the crossroads with
ambiguous impulses within him, then the power tempting through the
world-ideal could have wrapped itself in a mythical nimbus, and the
mythical would have been recognizable by the tempter's making
circumlocutions. The mythical temptation could then not possibly begin
with the tempter---however many thousand times he might have imagined
that he really had what he, plainly understood, does not have, all the
kingdoms of the world and their glories to give away---openly and
without ceremony demanding that the Messiah, ``the Son of God,'' should
fall down and worship him. No, he would have had to begin the other way
round; he would rather have had to begin with the promise that if the
Son of God, the Messiah, would only understand his task practically,
then not only all the kingdoms of the world, but he himself too, the
prince of this world, would fall down and worship the almighty Messiah.
Could the tempted one let himself be beguiled by this, then the tempter
would have conquered; then it would follow of itself that it would not
be the prince of this world who fell down and worshiped the Messiah,
but conversely the Messiah who would have to fall down and worship the
prince of the world, victorious through the temptation.


Before the pure, self-transparent will such evasions are impossible.
What the tempter has to say, he must say
% ---- printed p.254 (PDF 267) ----
\opage{254} openly and plainly, for the ambiguous, the veiled, is instantly
seen through. Precisely the circumstance that every rejoinder on the
tempter's side is so plain and unreserved that guile itself is forced
into candor brings it about that each single attempt, especially the
last, comes to look so presumptuous and yet so poor, so simple and
almost insane. But precisely thereby the secret in all temptation is
laid bare: temptation comes in the last instance from the evil self;
its innermost meaning is essentially meaningless; the decisive victory
is conditioned by its obscure, self-contradictory essence being
completely seen through; the decision is that the pure will, the true
self, with resolute energy banishes the tempter: ``Get thee hence,
Satan! for it is written: thou shalt worship the Lord thy God, and him
only shalt thou serve''.

Once the mythical veil which the ambiguous world-spirit spreads over
the tempting principle, and thereby over all temptations, has been
torn apart, it will also become apparent that the selfish being which
tempted Eve was a veiled appearance of the tempting power, the selfish
self that wished to try its strength with Christ. The tempting spirit
can now be seen in reflected light as ``the old serpent'', as ``the
liar'', ``the murderer from the beginning''.

Immediately and directly, no sinful human being on earth has ever
stood \textit{vis à vis} the tempter himself, the prince of this
world. To wish to know oneself immediately and personally tempted by
the tempter himself is a blasphemous presumption, the presumption of
wishing to set oneself equal with Christ and on one's own account to
pass through the deepest crisis of freedom; for only in the deepest
crisis of freedom can the evil self stand objectively clear before the
holy Self, the evil will be objectively determined against the good.
It is not the case that the power of the devil, together with the
devil himself, is already sufficiently revealed
% ---- printed p.255 (PDF 268) ----
\opage{255} before Christ comes, so that Christ has only the task of overcoming
the well-known devil; no, the truth is that Christ alone makes the
essence of the devil revealed and, precisely by making it revealed,
breaks the devil's power. According to the principle of revelation,
faith in the devil can acquire essential significance only through
faith in Christ; all propositions about the devil, in the New as in
the Old Testament, are propositions of faith, which look pitiful when
they are transformed into propositions of knowledge and treated
cosmologically.

And how, then, do matters stand with faith in the devil and his
tempting power? ``Every man is tempted, when he is drawn away of his
own lust, and enticed.'' This is temptation as a psychological fact.
But this fact, about which reason and revelation are agreed, would
hide its ground in the night of doubts, the eternal darkness of
ambiguity, if it did not receive light from the revealed word. The
more relative and veiled the tempting thing is, the more bewitching is
the temptation; the bewitchment vanishes only when the temptation,
briefly and definitely, without any evasions, can in faith be traced
back to the tempting principle, the selfish self, the evil will. Just
as faith in Christ is unconditionally positive, so faith in the devil
is unconditionally negative, a faith that aims exclusively at
renunciation. ``Dost thou renounce the devil and all his works and all
his ways''? That is the question; this pithy, short, authoritative
sacramental word is without evasion, without ambiguity.

The ideals of the world, like the Ideas of the world, are in and for
themselves neither good nor evil, in the personal sense neither true
nor false; they become so only through the relation in which they are
set to the holy will. The kingdom of the world acquires its ideal
significance only by furnishing, with the conditions of actuality, a
natural and yet spiritual foundation for the Kingdom of God.
% ---- printed p.256 (PDF 269) ----
\opage{256} \runhead{c) The Ideal of the Kingdom of God.}
\parmark{17}

How the original relation to God was obscured and cancelled has
already been shown; we shall now see how it is renewed. The fallen
race is to be raised up again; so a Kingdom of God is founded on
divine and yet human terms. It is a covenant which God establishes
with man, a relation of Self to self, of will to will, from which the
fundamental character of revelation shines forth. The Kingdom of God
thus founded is actual and yet supra-actual, present and yet
constantly to come. The ideal enclosed in actuality rises forth step
by step and demands its actualization; but at each following step the
present becomes ever more unsatisfying, and the strife between ideal
and actuality more keenly felt. The thought of life is great and
blessed, but it is followed by a curse and satisfies with blessing
only insofar as it nourishes itself on a promise. The inversion, that
progress is regress and regress in turn progress, gives the education
of the chosen people so peculiar a stamp that the whole development of
spirit noticeably distinguishes itself from rational development of
culture. We survey the movement from the three points of view:
\emph{election, legislation} and \emph{promise.}

\greekrun{α}{Election: The Chosen People.}

Religion is of antirational origin; the thought of governance cannot
be grounded in knowledge: the grounding is inverted, as is seen from
the election of Israel. In order to answer the question why God chose
the Jewish people as his own people above all others, appeal has been
made to the piety of its patriarchs and to its natural peculiarity,
its individuality stamped to a wondrous degree. In order
% ---- printed p.257 (PDF 270) ----
\opage{257} to bring out still more the peculiar significance of natural
determination, some have even wished to demonstrate a division in the
human race from the beginning, according to which the pious natures,
``the children of God'', were noticeably distinguished from the
ungodly, as Abel from Cain. This view is without objective validity.
There is indeed a difference between pious and impious individuals;
but the difference is so individual that it cannot ground any
separation of peoples. It is said that Seth was set in Abel's place
(Gen.\ 4:25); one would therefore expect to see the stock divide into
the good race of Seth and the evil race of Cain; but against this
stands, in the genealogies, the fact that the opposition is erased
again. According to ch.\ 5 all Adam's descendants without distinction
are Sethites, unless, that is, one will suppose that Chanoch (Enoch),
who ``walked with God'', was a son of Cain, and Noah, who ``found grace
in the eyes of Jehovah'', a son of Lamech. Indeed, Lamech --- if
ch.\ 5:25 and ch.\ 4:18, through a mingling of the genealogies, refer
to the same Lamech --- Lamech himself, with all his worldliness, seems
at last to have felt a need for divine comfort, a sign that his
arrogance was more youthfully naive than reflected in defiance.
Through Noah, then, Abraham himself is of Cain's line: the Israelites
are Cainites. Judging by such facts, if they really are facts,
``religious natures'' must therefore be found in roughly equal measure
in all peoples.

But even if it is so, as we willingly concede, that the Jewish
national character is in a peculiar sense disposed toward religious
development, this fact can in any case not be something accidental for
governance. Whence, if we hold fast to the standpoint of religion, do
the peoples have their various natural dispositions, if not from the
Creator? Now if the reason why the Jewish people is created with a
peculiar physiognomy is to be sought in its being created to be God's
people, and the reason why it becomes God's people through election is
to be sought in its being created with the dispositions needed for
this, then the grounding goes in a circle; this circular grounding
leads from knowledge to faith. We hear (Gen.\ 12:1) Jehovah say to
Abraham: ``Get thee out of thy country, and from thy kindred, and from
thy father's house, unto the land
% ---- printed p.258 (PDF 271) ----
\opage{258} that I will shew thee''; but nothing precedes that could inform us
what it is that makes Abraham, precisely, worthy of such a calling
rather than Nahor and Haran. We may ascribe to Abraham as many hidden
advantages as we please, but nothing stands about them in the text.
Abraham's election must then be ascribed just as unconditionally to
God's free choice as it must be ascribed to God's free choice that
Jacob was later preferred to Esau (Rom.\ 9:11--13). With election as
the expression of God's unconditioned will, the character of what
follows is already indicated.

Here, at the founding of the Kingdom of God, there is from the very
beginning the sharpest opposition of an almighty will and natural
conditions, and yet both factors in turn interlock in so infinite a
manner that the unity loses itself in the Mystery. Just as election
itself has taken place in faith and on the terms of faith, so its
historical actuality, too, can confirm itself only for a believing
consciousness. The history of the Kingdom of God is in a double sense
a history of faith, namely a history about believers and a history for
believers. The series of natural intermediate causes is broken through
by the miracle, and while probability at every moment runs against the
ideal, the improbable happens, and faith triumphs.

Jehovah's promise to Abraham runs: ``I will make of thee a great
nation, \dots\ and I will bless them that bless thee, and curse him
that curseth thee: and in thee shall all families of the earth be
blessed''. Thus expectation is at its highest; but time passes, the
fulfillment keeps them waiting, and comes only when it can happen only
against all probability. The son on whom the promise rests, Abraham's
only one, on whose life it depends whether the patriarch shall really
become a great nation, or only leave behind him illegitimate children
and so after all die a childless old man --- him Jehovah demands as a
sacrifice;
% ---- printed p.259 (PDF 272) ----
\opage{259} it is as though everything were to go backward again. Here there
must be faith with hope against hope, in possibility against
probability. And faith triumphs; the family % print: "Sandsynlighed, Og" (comma for full stop)
grows into a great nation, but --- is this the chosen people, the
people over whom the blessing rests? Why, they are slaves, nothing but
slaves! Far from the land of the fathers, the land Jehovah himself gave
Abraham, Isaac and Jacob for inheritance and possession, the people of
the promises toil in slavery under a foreign yoke: the apparent
progress is a sorrowful regress.

Here, humanly seen, there is no prospect of salvation; but precisely
this signifies that salvation is near, that it comes from the God of
the fathers. From the burning bush in the desert of Midian Jehovah
speaks to Moses, and Moses stands forth before Pharaoh, and Pharaoh is
terrified by sign upon sign, and the deliverance is wondrous. ``In the
third month, when the children of Israel were gone forth out of the
land of Egypt, the same day came they into the wilderness of Sinai
\dots\ And Moses went up unto God, and Jehovah called unto him out of
the mountain, saying: Thus shalt thou say to the house of Jacob, and
tell the children of Israel: Ye have seen what I did unto the
Egyptians, and how I bare you on eagles' wings, and brought you unto
myself. Now therefore, if ye will obey my voice indeed, and keep my
covenant, then ye shall be a peculiar treasure unto me above all
people: for all the earth is mine. And ye shall be unto me a kingdom
of priests, and an holy nation'' (Exod.\ 19:1--6). So the covenant is
concluded, so the promise stands firm after all; so the slavery was
only for the sake of deliverance, so the sorrowful regress was after
all essentially progress.

But what progress? The trial is long and the way narrow: there is
hunger in the desert, complaints and cries of woe are heard; in
bitterness and defiance the people rebel against Moses and against
Jehovah; indeed, so harsh is the opposition of ideal and actuality
that the earth must open and swallow the rebels, and a whole
generation die away, and its bones
% ---- printed p.260 (PDF 275) ----
\opage{260} be scattered in the desert, without any eye having yet seen the
promised land.

In the history of peoples, contrasts of fortune and adversity,
exaltation and abasement, fluctuations of progress and regress will
everywhere be demonstrable, while the movements on the whole
nevertheless lead toward the goal, the goal that is set for the
particular people and designates its calling. In this respect it
depends not merely on what a people actually experiences or has
experienced in outward things; it depends also, and chiefly, on how it
has experienced these outward things and understood what it
experienced. Only through understanding is what is outwardly actual
transformed into inward actuality, into actuality of spirit, actuality
for the spirit. In this respect a people's legendary history is not
merely just as essential as, but even more essential than, everything
the critique of knowledge extracts as bare facts; for in legendary
history the outward has become inward, and the inward in turn
outward. Legendary history is the history of the spirit of the
people, a history that wishes not merely to be remembered but to be
essentially recollected. But from this it follows that the history of
the spirit in every peculiar people must in principle correspond to
the spirit of the people; and from this it follows in turn that the
history of the chosen people, according to its whole actual character,
in the outer as in the inner, must be just as different from any
heathen national history whatever as the word of revelation in its
principle is different from myth.

\greekrun{β}{Legislation: The Theocracy.}

The Mosaic legislation is the spiritual, universal foundation of the
Kingdom of God. What is peculiar to this legislation is that it is
immediately and without evasion traced back to God himself; already
from this it is seen that it is intended for believers, not for
knowers. The Kingdom of God is not an
% ---- printed p.261 (PDF 276) ----
\opage{261} ordinary kingdom of Ideas but in the strictest sense a kingdom of
will: the almighty Self, the holy will, will be absolute sole ruler
in this kingdom. Not only the words of the Law but also the tables of
the Law are originally from Jehovah himself; ``the tables were the
work of God, and the writing was the writing of God, graven upon the
tables'' (Exod.\ 32:16). The opposition between Judaism and
heathendom is thereby designated as sharply as possible. In spiritual
heathendom it is the Ideas, the intermediate forms rich in content and
dazzling by their divine content, that place themselves between God
and man and obscure the consciousness of God; in Judaism these
intermediate forms are pushed back, even to the degree that it looks
as though they had quite vanished. This removal, not to say exclusion,
of the Ideas, which is so characteristic of revelation, is a condition
for the divine Self being able, in reflected immediacy, to penetrate
and move the human, a condition for God being able inwardly to breathe
upon man, to speak to man, and in the Spirit through the word to make
himself understood by man.

But precisely through this peculiar self-communication of the divine
will in the commandments of the Law a unity-in-difference of the divine
and the human is posited which impresses the imagination and repels
the understanding.

What, for the immediately believing view, could be more exalted or
more awe-inspiring than a beginning such as this: ``And God spake all
these words, saying: I am Jehovah thy God, which have brought thee out
of the land of Egypt, out of the house of bondage. Thou shalt have no
other gods before my face. Thou shalt not make unto thee any graven
image, or any likeness of any thing that is in heaven above, or that
is in the earth beneath, or that is in the water under the earth. Thou
shalt not bow down thyself to them, nor serve them: for I, Jehovah thy
God, am a jealous God,
% ---- printed p.262 (PDF 277) ----
\opage{262} visiting the iniquity of the fathers upon the children unto the
third and fourth generation; and shewing mercy unto thousands of
generations of them that love me and keep my commandments'' (Exod.\
20:1--6). Thus the holy Self has a voice; where this voice resounds in
consciences, it awakens earnest thoughts; where this voice of the
spirit is heard, it shakes the ground of the mind, so that the gods of
the world, the new with the old, must collapse.

But the same thing which in the divine speaking conditions the
exalted brings with it at the same time such a finitization,
parceling out and multiplication of commandments and prescriptions
that reflective understanding, which is not attentive to the principle
itself, must take offense. That prescriptions for the arrangement of
the tabernacle, insofar as it is a matter quite generally of the
fundamental thought, are traced back to Jehovah we can indeed
understand from the standpoint of religion; but the more the
prescriptions go into particulars and are yet immediately put forward
as Jehovah's commandments, the more pettily exact everything is
prescribed, the more the human shows through, the more we miss the
intermediate determinations of the Idea, and the more offensive the
whole becomes. Or what man of reflection can read the following
command (Exod.\ 25:1 ff.): ``And Jehovah spake unto Moses, saying:
Speak unto the children of Israel, that they bring me an offering: of
every man that giveth it willingly with his heart ye shall take my
offering. And this is the offering which ye shall take of them: gold,
and silver, and brass, and blue, and purple, and scarlet, and fine
linen, and goats' hair, and rams' skins dyed red, and sealskins, and
acacia wood, oil for the light, spices for anointing oil'' etc.\ ---
who, we ask, can read such things without being involuntarily seized
by a secret dread that negative critique may after all be right, that
the spirit of revelation is spiritless, and the whole holy work of
legislation a work of religious politics?

% ---- printed p.263 (PDF 278) ----
\opage{263} Immediate historical faith in sacred history is superstition; this
holds just as well of the facts through which the highest and deepest
shines forth as of the prescriptions that concern the least, the most
petty of all. If faith in the revealed word is to be put into force,
then superstition must go; if superstition is to go, then the
principle must be looked to. That religion is antirational does not
mean that religion is spiritless; but spiritless it becomes, more
spiritless than the most spiritless of all mythologies, when there is
immediate, literal faith in a God who in an outward, finite manner
dictates outward, finite prescriptions. It is this dictating that
forms the hard point, the knot, in the Old Testament revelation. But
the knot is untied when we reflect on the principled relation between
the outward and the inward, between the immediacy of religious
representation and the thought of faith. For immediate beholding the
matter stands thus: that Jehovah outwardly and sensuously shows
himself to Moses and speaks to Moses as one man can speak to another.
If this way of seeing is to be defended, then it must be defended by
defying all thinking; if faith in this sense is to bid defiance to
thought, then the defense is easy to conduct: ``Religion is
thoughtless; he who believes dare not think''. It is a poor pretext
that there is some thinking here and there and back and forth about
religion, when the foundation itself is thoughtless and
thoughtlessness is fixed fast in the principle.

To begin with thought is to begin with the spirit, that is: not to
begin with abstract reflections that despise the outward, but with the
thought of revelation itself, the thought of will, the holy
fundamental thought; it is: through the outward to seek the inward and
in the light of the inward to clear up the outward. The almighty Self
is not for Moses something outward, but originally and essentially
something inward; Jehovah is invisible and becomes visible only to
% ---- printed p.264 (PDF 279) ----
\opage{264} the prophet's vision: Jehovah is the God of the spirit, Jehovah is
spirit. But spirit is will, and the all-determining will reveals
itself in thought and speech. The inner speaking cannot be explained by
rational psychology; for here it holds that religion is antirational.
But that the inner speaking in the prophet's spirit is determined by
an infinite conversion of divine thoughts into human ones, of divine
words into human language, can be cognized in faith; for this ideal
doubling expresses precisely the essence of faith. Since, now, the
holy thoughts are thoughts of will, and the revealed word is
proclaimed in commandments and prescriptions, it follows from the
principle of faith itself that Jehovah's prophet cannot engage in any
finite distinguishing between Jehovah's thoughts and his own thoughts,
between Jehovah's commandments and the commandments devised by
himself; the moment he permitted himself such a distinction, the
Spirit would leave him: he could not possibly, in the prophetic sense,
come forward as Jehovah's servant. Mightily moved by the Spirit, the
legislating prophet must proclaim every commandment in Jehovah's name,
with authority from Jehovah.

When laws and institutions are immediately determined by God's will,
the state must be conceived as God's state, and God himself be
believed to be the people's king. The Idea of the state then does not
come to any particular development; the civil laws become religious,
and the religious civil. With this fusion of the holy and the secular,
the ideal and the actual, which is the distinguishing mark of
theocracy, it accords that the prescriptions for cultus and customs
are put forward so definitely and so specially that they can be
immediately complied with. In that the divine will gives itself
expression in the many finite commandments and prohibitions, it
acquires a concrete content and enters into many-sided reciprocal
action with the human will. In order to be able to obey God's will,
the Israelite must know what it is God wills in each given case; this
he comes to know in the Law. What significance the individual
prescription
% ---- printed p.265 (PDF 280) ----
\opage{265} has in itself, whether it is necessary or unnecessary,
rational or irrational, significant or insignificant --- about that the
believing Israelite dare not reason: the decisive thing is that the
prescription is Jehovah's commandment, that it is found in his Law.
That there must now also be a special order which watches over the
promulgation and observance of the Law, that the Levitical priesthood,
in other words, corresponds to the theocracy, is understood of itself.


But although the Law consists in a sum of finite precepts, it is nonetheless
in its essence spirit, and must be understood spiritually. This determines the
position of the prophets in relation to the priests. While the priests have to
preserve what is already given in revelation, uphold what has been handed down,
and see to it that the Law is fulfilled according to the letter, the prophets
are organs of the continuing revelation; they therefore look not primarily at
the outward but at the inward, and take the Law not according to the letter but
according to the spirit.

The life of the spirit now has a double resistance to overcome: it meets
resistance in the lower, sensuous nature, and it meets resistance in the egoism
of the understanding. In the first respect the temptation to idolatry, in the
latter respect righteousness by works, is what predominates. The sterner the
spirit is, the more it provokes to recalcitrance: the colder the high Invisible
One is, the warmer are a Moloch and the other foreign gods. The continuing
strife between the people and Jehovah is a natural consequence of the dualism in
man's essence, and an eloquent testimony that it is not the people that has
chosen Jehovah for its God, but Jehovah that has chosen the people for his
people. In righteousness by works we have a typical expression of what spiritless
spirit means. It is a demand of the spirit, the first, elementary condition of
spirit, that the law of the commandments shall be observed; only in the
observance of the commandments does reverence have its actual expression. But the
commandments nonetheless concern primarily what is outward;
% ---- printed p.266 (PDF 281) ----
\opage{266} the outward fulfillment of outward precepts is in practice possible
without essential change in what is inward. The State of God is then conceived as
a civil state, God's laws as merely civil laws, the Ruler, Jehovah, as a civil
king. This is the inherent one-sidedness of theocracy. That through the outward
the inward is aimed at, that conditional obedience in what is outward is only a
sign of the inward, unconditional obedience---the obedience that breaks egoism
and bends man's will in each and every thing to God's will---this is overlooked.
Now if it were merely crudeness and thoughtlessness that brought the people to
forget the inward over the outward, then spiritlessness would be immediate, and
could essentially be removed by instruction and cultivation. But with a
revelation as its presupposition, spiritlessness cannot be immediate; it is the
serpent-principle at work, that egoism places the fear of God in the outward in
order to keep the inward for itself. As the fear of God becomes egoistic, egoism
becomes God-fearing; from this fusion there issues not only hypocrisy, but here
is bred a fanaticism, a zeal glowing in selfish self-ignition for the
commandments of the Law, for holiness, righteousness and blamelessness in what is
outward. It is spiritlessness; but this spiritlessness is not naive: it is a
crime against the spirit. Against the egoism reflected in a crime against the
spirit, an egoism of godliness which can well unite even bloody misdeeds with its
holiness of works, the merely instructive, quietly admonishing word is too weak;
then the prophetic voice is raised. ``To what purpose is the multitude of your
sacrifices unto me?'' saith Jehovah. ``I am full of the burnt offerings of rams
and of the fat of fed beasts, and in the blood of bullocks and of lambs and of he
goats I delight not \dots\ Your new moons and your appointed feasts my soul
hateth; they are a burden unto me, I am weary to bear them. When ye spread forth
your hands, I will hide mine eyes from you.
% ---- printed p.267 (PDF 282) ----
\opage{267} Yea, when ye make many prayers, I will not hear you; your hands are
full of blood'' (Isa.\ 1:11--15).

The prophetic judgments are penal judgments. The chastening of the people, its
spiritual education, is conditioned precisely by the religious reciprocity
between these judgments and their execution. If we disregard the judgments, then
the humiliations and chastisements which, after the division of David's kingdom,
befall first Israel and then Judah would not be qualitatively different from
such misfortunes and miseries as also befall heathen peoples. If, conversely, we
abstract from the actual chastisements and conceive the penal judgments merely as
pious threats against the ungodly, outbursts of mood from a subjective holy
wrath, then the prophets sink down into ordinary preachers of repentance, without
higher divine authority; the word of revelation dissolves into a poetic imagery
of religious enthusiasm. Only when the execution of the penal judgments is
cognized in the corresponding chastisements, so that these receive their
spiritual significance precisely by being in truth the execution, can the
punishing word be cognized as a prophetic word. Thus it was cognized in faith by
the people of Israel. By humbling itself under the mighty hand of Jehovah, the
believing Israel felt itself not only bowed down under the fatherly rebuke, but
also raised up by the promise attached to its conversion.

\greekrun{γ}{The Promise: The Messianic Ideal.}

On the basis of opposition, the supramundane joins itself in theocracy so
immediately and directly with the mundane that the strife between ideal and
actuality must strike the eye plainly. For the keeping of the Law the very
richest blessings are promised, but against transgression of the Law the most
terrible curses are directed as well. ``If ye walk in my statutes, and keep my
commandments,
% ---- printed p.268 (PDF 283) ----
\opage{268} and do them,'' saith Jehovah (Lev.\ 26:3 ff.), ``then I will give you
rain in due season, and the land shall yield her increase \dots\ and ye shall
eat your bread to the full, and dwell in your land safely \dots\ And ye shall
chase your enemies, and they shall fall before you by the sword \dots\ And five
of you shall chase an hundred, and an hundred of you shall put ten thousand to
flight, and your enemies shall fall before you by the sword'' \dots\ So worldly
is the in itself supramundane Kingdom of God that the rewards held out by Jehovah
must appeal in the very highest degree to the egoism of the people. Equally
sensuously deterrent, again, are the threats added. ``But if ye will not
hearken unto me, and will not do all these commandments; and if ye shall despise
my statutes, and if your soul abhor my judgments \dots\ and break my covenant, I
also will do this unto you: I will visit you with terror, with consumption and
with burning fever, that shall consume the eyes and cause sorrow of heart, and ye
shall sow your seed in vain, and your enemies shall eat it'' \dots; in this
spirit curse follows upon curse, each more terrible than the last. What holds of
the blessing holds of the curse: it appeals to what is selfish in man. Herein
there is certainly a contradiction, indeed a fundamental contradiction; but the
contradiction does not strike revelation, it strikes only human nature, whose
egoism becomes manifest precisely through the relation to the Holy.

The Kingdom of God will and shall be a kingdom of actuality; so corresponding
conditions are set: the blessings aim immediately at what is actual, just as the
threats do. The actual, in sensuous presence, is something earthly and tangible.
It is from the earthly-tangible that one must begin, when one is to begin with a
crude, recalcitrant, but vigorous people and yet begin in the spirit. The spirit
cannot from the beginning make comprehensible to the
% ---- printed p.269 (PDF 284) ----
\opage{269} sensuous man what in the eternal sense must be understood by
actuality, actual blessing and actual curse; so it begins with the temporal, for
the temporal too is actual. This double-sided essence of actuality provokes and
deceives egoism, and sets the powers of the spirit in motion. ``If ye walk in my
statutes, and keep my commandments''---that is the ideal demand, which what is
egoistic in man involuntarily resists; yet what follows? ``Then I will give
you,'' etc.; here in the added promises there is something for egoism, and yet
the promises cannot possibly be understood egoistically without being
misunderstood. When, then, reflection is set in motion, and man asks himself:
does it then always go well with the pious? How does Jehovah fulfill his
promises?---then conscience is awakened and raises the religious question: how
are Jehovah's commandments fulfilled? ``There is none that doeth good, no, not
one.'' It is the problem of retribution that presses forward and, in its
continued, anguished working-through, brings it about that the cognition of sin
grows deeper and deeper, and that the actuality of reward and punishment
acquires for believers an ever more spiritual significance.

How the spiritual, ideal conception of the Kingdom of God works its way forth
out of the first actual presuppositions is clearly seen from the transformation
that the fundamental theocratic thought underwent in the course of time, a
transformation which, far from being a product of human arbitrariness, was in
the strictest sense of the word a work of the Spirit.

Originally Jehovah is the King of the people, and the community, with its civil
institutions, is immediately Jehovah's State. But as regards the organization
and governance of the State, the chasm between the human and the superhuman, the
earthly and the heavenly, was bound gradually to show itself so great and deep
that the need for an
% ---- printed p.270 (PDF 285) ----
\opage{270} actual mediating link could be cognized. The development of this
need, called forth by the strife of the ideal with actuality, went hand in hand
with the development of the fundamental religious consciousness and
corresponded to the prophetic view. So the people chose for itself an earthly
king; but the earthly king was anointed by the prophet with Jehovah's spirit,
and now comes forth in his ideal light as the actual mediating link between the
Lord in heaven and his kingdom on earth, as Jehovah's son, the Anointed, the
Messiah.

On the king, the head of the people, the promise now rests; the greatness, might
and glory of the kingdom are conditioned by the king's corresponding to his
ideal. But just here the strife between ideal and actuality is made visible. Only
one among the kings was so anointed with spirit that all the signs of the
promise could be recognized in him: that was David. For the glory of Solomon, a
worldly reflection of David's spiritual greatness, was a brilliant phenomenon
that came and vanished.

The Son of David who was to establish the Messianic kingdom, at once holy and
yet worldly, had then also to be the son of the Lord of Hosts, a mighty
war-prince, to whom Jehovah would give the heathen for an inheritance, that he
might ``break them with a rod of iron and dash them in pieces like a potter's
vessel'' (Ps.\ 2), and yet a Prince of Peace, who ``hath delight in the fear of
Jehovah,'' who ``judgeth the poor with righteousness'' and ``reproveth with
equity for the afflicted of the land'' (Isa.\ 11:4 ff.). As the heavenly is in
this way drawn down into the earthly, this earthly is by the same token raised
continually more and more up toward the heavenly; actuality thus gradually takes
on an ever stronger stamp of the supra-actual, of symbolized actuality of the
spirit. When the Messiah ``smiteth the earth with the rod of his mouth, and with
the breath of his lips slayeth the wicked,'' then ``the wolf shall dwell with the
lamb, and the leopard shall lie down with the kid,'' and the child
% ---- printed p.271 (PDF 286) ----
\opage{271} ``shall play at the lizard's den.''---``They shall not do evil nor
work ungodliness in my holy mountain, for the earth is full of the knowledge of
Jehovah, as the waters cover the sea.''

Essential progress is conditioned by regress. The more things went backward with
actuality, the outwardly actual kingdom, the more they went forward with the
ideal conception of the Messiah and his kingdom. What is not to be seen on earth,
the prophet then sees in the clouds of heaven. ``I saw in the night visions, and,
behold, with the clouds of heaven came one like a Son of man, and he came to the
Ancient of days, and was brought near before him. And there was given him
dominion, and glory, and a kingdom, that all peoples, nations, and languages
should worship him. His dominion is an everlasting dominion, which shall not
pass away; his kingdom shall not be destroyed'' (Dan.\ 7:13--14). The Messiah is
thus the Son of David and yet the Son of God, an actual man and yet a
supernatural being. By the Messianic ideal a view of the union of the divine and
the human nature is prepared. Jehovah's countenance is softened, for the Father
loves his Son. But it is only a preparation. Ideal and actuality, the heavenly
and the earthly, do not approach one another directly: the extremes are strained
to the utmost. Only when the representation of a Son of God who is also to be a
Son of man has taken up so many ideal presuppositions that the seal of the
mysteries can be broken does the turning point arrive, the turning point of
eternity, which is the fullness of time.

\begin{center}\rule{5cm}{0.4pt}\end{center}

% ---- printed p.272 (PDF 287) ----
\opage{272} \parthead{Faith in the Son.}
\lettersub{A.}{The Essence of the Son.}
\parmark{18}

That Jesus of Nazareth, the son of Mary, is the Savior of the world and the Son
of God means, in other words, that Jesus of Nazareth is the sign of
contradiction. ``This child,'' says Simeon of the Child in the Temple (Luke
2:34), ``is set for the fall and rising again of many, and for a sign which shall
be spoken against.'' The sign of contradiction has, as shall be treated more
closely later, the peculiarity that each must understand it in proportion as he
understands himself; it is the problem of self-understanding that this sign would
awaken. If Jesus of Nazareth is the Savior of the world, then his figure must be
seen in faith, transfigured by the Spirit; if, on the other hand, the figure of
Christ is regarded from the standpoint of knowledge, then the negative critique
has already sufficiently shown what historically comes out of the evangelists'
accounts. There can be no greater proof of the absolute heterogeneity of the
principles than the change that takes place in the Savior's image according as
the light falls in through the temple door either of science or of religion. The
improbable must be believed, but the impossible cannot be believed. It ought
therefore first to be investigated how far the science of reason is able to
prove the impossibility
% ---- printed p.273 (PDF 288) ----
\opage{273} that Jesus of Nazareth, this individual man, can have been conscious
of himself as the Savior of the world, and how the divine-human
self-consciousness, if its possibility is maintained, is to be conceived in
faith; we bring this investigation under the consideration of: \emph{The Self in
the Son}. It ought next to be investigated how the ideal marks of essence by
which alone the Savior of the world can be recognized, conceived as marks of
faith, withdraw themselves from the dissolving critique of ``the modern
consciousness''; we treat of this under: \emph{Fundamental Determinations in the
Son's Self}. It ought finally to be investigated how far and in what way it is
possible to cognize the relation between God and Christ, between the One who
reveals himself and the Revealer, as an original relation of essence, an eternal
fundamental relation between the Father and the Son; we gather together what
belongs here under the point of view of: \emph{The Personality of the Son}.

The problem of personality is at once the fundamental problem of humanity and of
religiousness; it must then at length become evident where this problem is to
find its final, proper solution, in knowledge or in faith, in reason or in
revelation.

\runhead{a) The Self in the Son.}
\parmark{19}

What lies originally in the concept of a Savior of the world, that he must be
``both true man and true God, and both in union,'' the old teachers of the
Church, Athanasius for example, have already brought out with clear
consciousness and full emphasis. ``For unless his nature had been of like kind
with ours, he could not have redeemed it; without true Godhead he could not have
made us divine; and without a living bond between both
% ---- printed p.274 (PDF 273) ----
\opage{274} sides he could not have established any connection between God and
men.'' Here the following three questions now press themselves upon us:
α) in what sense has the evangelical Christ been conscious of himself as a human
self? β) in what sense has he been conscious of himself as a divine Self? and
γ) within what limitation can the unity of the divine and human self be
cognized as divine-human unity of self?

\greekrun{α}{The Human Self.}

That Jesus Christ, who in the Gospels by preference calls himself ``the Son of
Man,'' was an actual man, is something that scarcely anyone in the present day
calls into doubt. For that Christ should have been an Aeon emanated from the
heavenly Pleroma, who had not become man in flesh and blood but had only taken
on an apparent body, with which he had also been crucified without suffering---such
gnostic fairy tales no longer find credence. Clement of Alexandria does indeed
find it a ridiculous opinion ``that the body of Christ, which was sustained by
the holy power of the Logos, should have needed food for its nourishment, since
on the contrary it was only in order to forestall docetic opinions about his
person that Jesus had taken food, although he himself was equally insensible to
impressions of pain and pleasure, toward which one is passive''; but not to
speak of the fact that such a forestalling of docetic opinions is
self-contradictory, since it confirms precisely what it was meant to forestall,
this whole way of looking at things is equally at odds both with the
evangelical presentation and with the concept of a Savior of the world. If now
the emphasis is to fall on being an actual man---on being born, growing, being
affected by sensuously pleasant and unpleasant impressions, etc.---then weight
must first and foremost be laid on answering the question how
% ---- printed p.275 (PDF 274) ----
\opage{275} the development of the human individual properly stands in relation
to the human self.

The human self begins in unconsciousness. That an actual child should be able to
be born into the world with fully developed self-consciousness is a
contradiction from which faith recoils just as much as knowledge. The Gospel of
the Infancy of Jesus, in which it is related that Jesus spoke in the cradle and
said to Mary his mother, ``I am Jesus, the Son of God,'' is certainly
apocryphal, not only by the judgment of knowledge but also by that of faith.
With all the supernatural signs with which the Christ-child is found surrounded,
the canonical Gospel narratives nonetheless hold fast to the fundamental thought
that the child in actuality began as a child, that it grew, in soul as well as
in body, a quiet growth---that, in other words, an actual natural
self-development took place. If, therefore, in the present connection there
were talk only of the actual man, then, as regards Christ's human self, there
could not possibly be any dispute. But here, surely, the talk is not of the
merely actual man, but expressly of the \emph{true} man. The \emph{actual} man
is to be the \emph{true} man, and this \emph{true} man an \emph{actual} man:
here we have the strife again between faith and knowledge.

In his ``Presentation of the Christian Doctrine of Faith''\footnote{II Bd.\
S.\ 191--194.} Strauss, from the standpoint of knowledge, with clear survey and
penetrating acuteness, has defended his Christology against the attempts to
uphold the ecclesiastical representations. We set down this defense, for the
sake of what follows, as a rounded whole. ``It was surely not at all
unthinkable---it was said---that the Creator had originally given human nature
such a disposition that the universal human determination, i.e.\ the
religious-moral one, could at a particular point be attained in
% ---- printed p.276 (PDF 289) ----
\opage{276} perfection. As such a mere thinkability says nothing, there is also
nothing to say against it; but the conclusion which is now drawn, namely that if
one denies that thinkability one must assert that it belongs to man's essence to
be imperfect, and that human nature would be sublated by religious
perfection---this conclusion need only be better expressed: namely thus, that it
is essential for the individual man, or for the phenomenon of man \textit{an
sich}, to be imperfect, in order to say exactly the same as we assert. But the
manifestation of a man perfect in religious respect, it is now said, was
necessary if the religious life of the community was to reach its completion.
The whole religious process of development of mankind strives toward an
all-commanding founder of the true universal Church, whose religious life cannot
be further heightened and made complete. One easily sees that what was first to
be proved is here presupposed. That all other forms of religion are destined to
pass over into the Christian Church can, so long as this transition has not
actually taken place, be inferred only from the dignity of the founder of
Christianity; but one cannot conversely infer the latter from something that as
yet stands only as an empty pretension. But there dwells in man, it is further
remarked, a deep longing for the completion of life: now if he is unable to
actualize the Idea of mankind in himself, it nonetheless affords him essential
satisfaction when he finds it actualized outside himself, in another individual
of his race---or imagines it actualized; and we do not see how the argument,
lame by nature, could defend itself against receiving this twist. Yet one comes
with still another objection, and in it the polemic of the Schleiermacherian
tendency against the speculative Christology meets with that of the Hegelian
school: if there has
% ---- printed p.277 (PDF 290) ----
\opage{277} never actually existed a man perfect in religious respect, then the
concept of mankind never becomes perfectly actual; so long as the Idea has not
% print: "blevetfuld stændig" (misplaced word space), rendered as "fully realized".
been fully realized in one individual, it is a mere ought and an empty ideal.
At the root of this lies the most unclear confusion of two quite different
representations of the actualization of a generic concept. By it one can
understand either that the genus and each individual, or that the genus and the
totality of individuals, exhaust and cover one another. The former one does not
find in mankind (as indeed not in any genus at all, since this concept contains
the opposite): very well, then one must content oneself with the latter. That one
has no wish to do, and so one will take a middle way: in the one individual,
Christ, the genus is to have actualized itself in the first manner, but in the
rest of mankind not. The human race thus becomes impotent as regards the
overwhelming majority of its individuals, and in addition there comes a rent in
the circuit of its actualization such as one finds in no other genus: namely the
difference between one individual in whom the genus satisfies itself, and such
in whom it does not find its satisfaction. If these latter are not to appear
rejected and cast off, then the excellence of that one individual must in one
way or another redound to their benefit: which is not only the doctrine of the
Church but also Schleiermacher's, since according to him God sees mankind only
in Christ. But when God thus, in order to find the Idea of mankind perfectly
actualized in mankind, must as it were reckon together Christ's merit and the
defects of the rest of men, then after all it becomes the second manner in which
a generic concept can be actualized, namely such that it becomes completely
actual only in the totality of its individuals, which mutually supplement one
another; only in the thought-
% ---- printed p.278 (PDF 291) ----
\opage{278} less combination with the first kind of actualization, which is to
take place in one individual, while the other takes place through all the rest
in their relation to that one. This combination is contrary to reason for the
reason that the genus, if it could actualize itself completely in one
individual, would not torment itself with dispersing itself into a multiplicity
of such individuals and into their development in time; but it would exist only
in that individual and as that individual, as genus-individual; just as God, if
he could immediately be an individual personality, would have been able to spare
the whole enormous apparatus for the production of personalities (the creation
of the world and world history).''


Though from a scientific standpoint there may be a good deal to object to in
the \textit{nervus probandi} of the arguments with which Strauss here combats
his opponents, it is nevertheless evident that on the ground of knowledge there
is no foothold for one who would forcefully repel the well-aimed attack of the
negative critique. Consider the intelligent resoluteness with which the
spokesman of ``modern science'' concedes ``that it is essential to the
individual human being, or to the phenomenon of man \textit{an sich}, to be
imperfect''; note the objective unconcern with which ``the argument lame by
nature'' that is drawn from ``man's deep longing for the completion of life''
is pushed aside: and one will understand that religion must absolutely be given
up when its antirational principle is to be rationalized.

Strauss appeals to the fact that the race and the single individual do not
``cover each other'', since it lies necessarily in the concept of the race that
it can be exhausted only in the totality of its individuals. From this point of
view the human race as race is identified without further ado with any animal
genus whatever, and even so the consideration is lame. If the human race as
race is to be placed in
% ---- printed p.279 (PDF 292) ----
\opage{279} line with the animal genera, then it comes to stand not on an equal footing
with them, but lower than them all. For of genera and species in the animal
kingdom it really holds that the genus is covered by the species, and the
species, each in turn, by the corresponding specimens. It is the normal health
of the specimens, and by no means their number, on which this covering
essentially depends. A single specimen of a sound and tolerably faultless horse
is without further ado congruent with the species, whereas a plurality of
faulty specimens, of which one, e.g., has spavin, another glanders, a third the
staggers, and so on, by no means contribute by their number either to
exhausting or to covering the species. In the animal world there is no
essential discord between species and specimen, and therefore there is no
incommensurability either between the genus and its individuals: the genus is
exhausted in the totality of the individuals precisely in the corresponding way
as, through the species, it is exhausted in each single specimen. The
one-sidedness that naturally clings to the specimens in their singularity is,
when we look at the matter in principle, not at all remedied by multiplication.
If the one-sidedness is placed in the difference of sex, two specimens would be
enough to represent the whole species. This comes simply from the fact that in
the animal individual there is no conflict between ideal and actuality.

But conflict between ideal and actuality is precisely characteristic of every
single human being, and just thereby, in the highest sense, characteristic of
the relation between the human race and its individuals. That no human being is
perfect, that according to the natural order of things there is not a single
true and perfectly sound human being on earth --- on this reason and religion
will easily agree; but when it is to be determined more precisely wherein the
imperfection consists and whence the untruth stems, they at once disagree. How
deep the disagreement goes is seen already from this, that while religion holds
to the actuality of
% ---- printed p.280 (PDF 293) ----
\opage{280} \emph{the one true man}, reason holds to its impossibility. The whole dispute
turns essentially on the self. Religion, which asserts the self-validity of
every single human being by unconditionally asserting the validity of the self,
places the imperfection in the self and the untruth of the self in sin.
Conceived as sin, the untruth is not merely a general contradiction in the
essence of the self, but a contradiction of the will, in every human individual
the innermost, deepest self-contradiction. Now if it lies in the nature of
every contradiction that it must demand a resolution, then this must hold in
the very strictest sense of self-contradiction. That in the race One has
appeared who has presented himself as the true man, as the one who is himself
from the ground up free of self-contradiction, and thereby able to loose the
individuals from the chains of contradiction, and able, on the terms of faith,
to save everyone from the gnawing untruth and to make it possible for him to
\emph{become} a true man: this presupposition is so far from containing any
contradiction that everyone who contradicts it precisely makes himself doubly
guilty of self-contradiction. If the individuals and the race are to ``cover''
each other, this can happen only through the sinless man.

It is, however, highly interesting from the standpoint of knowledge to observe
the cunning with which the negative critique maneuvers when it is a matter of
getting every testimony to Christ's sinlessness explained away. ``Christ's
sinlessness, and that not merely as the realized possibility of not sinning
(\textit{potuit non peccare}), but as the impossibility of being able to sin
(\textit{non potuit peccare}), had,'' remarks Strauss, ``in the old
ecclesiastical system its firm dogmatic ground in this, that in Christ the
human nature, which was apt for sin, had no personality of its own at all, but
was only taken up into the personality of God's holy Son. For Schleiermacher
this grounding falls away together with the ecclesiastical Christology.
% ---- printed p.281 (PDF 294) ----
\opage{281} Just as little will he rely on the exegetical proof, and in this he does
right. For what probative force would it have, if Paul, of whom it is not known
that he knew Jesus personally, after his conversion ascribed sinlessness to him
in ever so definite terms? But instead, he moreover says only in a rhetorical
antithesis (2 Cor.\ 5:21) that God has made him who knew no sin to be sin, or a
sin offering, for us, that we might become righteous in him; wherein, as a
definite thought, there lies only this, that Jesus' violent death may not be
regarded as punishment for a transgression of his own. No more lies in
1 Pet.\ 1:19, where he is called a lamb pure and without spot, and
Heb.\ 7:26, where he is called a holy, innocent high priest; only that here the
notions of Levitical purity have also mixed themselves in \dots\ Yet even if
these statements had a less typical and dogmatic form, even if they had the
most definite historical form, and were moreover taken from genuine writings of
Jesus' immediate disciples, they still could not, even if they would, say more
than that well-known testimony of Xenophon concerning Socrates
(\textit{Οὐδεις πωποτε Σωκρατους οὐδεν ἀσεβες οὐδε ἀνοσιον
οὐτε πραττοντος εἰδεν οὐτε λεγοντος ἠκουσεν}), from which, however, no one has
yet wished to conclude that that heathen was sinless.
% print: "ποποτε" for πωποτε.
For as sinful men Jesus' disciples could never be certain whether something in
him which, precisely on account of their sinfulness, did not appear to them as
sin, was not nevertheless sin in actuality. Only the sinless one can give a
valid testimony to sinlessness, and consequently Jesus, since he is supposed to
have been the only one of this kind among men, could alone know that he was so.
This consciousness he is now also supposed to have expressed in the well-known
question (John 8:46); but even leaving aside that
\textit{ἁμαρτια} % print: "άμαρτια" (acute for the rough breathing).
here, according to the context, must rather mean error, and that, even if it
were said in the sense of sin, it could only have been said of such sins as
could be observed, Jesus'
% ---- printed p.282 (PDF 295) ----
\opage{282} own assertion cannot prove his sinlessness, since we should indeed already
have to be certain of this beforehand in order to be sure that by that
assertion he did not deceive himself or
others''\footnote{Strauss. Anf.\ Skr.\ II. S.\ 155--56.}.
% print: no closing quotation mark for the Strauss quotation opened on p.280.

Critique has here brought out a decisive point, a truth that cannot be
impressed enough, namely: that of Christ's sinlessness there is no objective
knowledge. Sinlessness is a concept of faith. The objection that nothing is
known about that which, by the nature of the case, cannot be an object of
knowledge thus falls away. The sinless one can know with himself that he is
without sin; but his own testimony to this must be received in faith: faith is
in this connection not a surrogate for a lacking knowledge; the certainty of
faith is, in relation to its object, self-certainty, and self-certainty is more
infallible than any straightforward historical fact whatever. That Jesus Christ
knew himself to be the true man, free of all sin, is not to be proved from
single passages of Scripture; it is to be seen and cognized from his whole
appearance, from the authority with which he speaks, from the decisive weight
he lays upon himself and his relation to the Father, from the infinite
superiority with which he, ``meek and lowly in heart'', invites all ``that
labor and are heavy laden'' to come to him, that souls may find rest. In
truth, if he who calls himself the Savior of the world did not know with
himself that no one either in heaven or on earth could convict him of any sin,
then he has been caught in a terrible self-deception. For the nearer a man
comes to the ideal without realizing it, the more clearly he sees the distance;
this holds of every ideal, but especially of the holy one.
% print: "Idèal" (grave accent) against "det Ideale" just before.

% ---- printed p.283 (PDF 296) ----
\opage{283} It is a religious-psychological fact that the consciousness of sin grows with
the consciousness of God, until sin is removed; therefore the great religious
characters, not only the Reformers but also the apostles, not only a Luther but
also a Paul, have so vivid a cognition of sin that people at a lower stage of
consciousness cannot rightly understand them. This cognition of sin is not
among the things about which truly religious characters are glad to keep
silent, over which they prefer to spread a veil; on the contrary! it belongs
precisely to what they speak of with emphasis, so that they speak loudly
without becoming loud-voiced: for as repulsive as it is to hear a preacher of
repentance, exalted by holy imaginings, speak conceitedly of himself as the
greatest sinner and, as further proof, shamelessly peddle his secret frailties
and private wretchednesses in order to titillate the pious curiosity of his
hearers, just as essential and categorically necessary is it that the cognition
of sin sound along in religious discourse, and that so penetratingly that
everyone must feel that this speaker speaks from within.

Now when this stands firm, what then are we to judge of him who calls himself
the Savior of the world? One of two things: either this man is in truth without
any sin, or else he is the greatest sinner! That this is said without any
exaggeration, everyone will see. Even if Christ had in actuality been the most
religious of all religious men, if he had nevertheless not been sinless, then
his word and work are a monstrous crime of spirit against God and men; for he
has indeed spoken only of the sin of others, but not of his own. Even if it
could be proved exegetically, which it cannot, that none of the utterances
attributed to him speak of sinlessness in the strictest sense of the word, that
is still not enough. Had he really not been sinless, we should have to expect
% ---- printed p.284 (PDF 297) ----
\opage{284} that in his whole appearance, in all his discourses, on every occasion, there
would be found a confession that he himself was originally a sinner, precisely
the one who in his human self had the deepest cognition of sin, since he had
the truest cognition of God and by it came nearest to the ideal. Now when this
confession is nowhere to be found, when neither followers nor enemies have been
able to discover the least hint of such a self-confession, but on the contrary
have heard from his own mouth, both to edification and to offense, that he was
the truth itself, the King of truth, the way and the life and the truth: then
he has, with or without consciousness, concealed a truth and violated the
truth. Here there can be no talk of innocent self-deception. If Christ on this
one point deceived himself and others, then he is not a King of truth but of
lies, the incarnate serpent-principle, the greatest deceiver who has lived on
earth. Though knowledge might not find such a deception improbable, faith
nevertheless finds it impossible.

When Christ's sinlessness is referred to its possibility \textit{(potuit non
peccare)}, it is chiefly the human side of freedom that is brought out. Herein
lies an invitation to compare Christ with Adam, the original human essence in
the individuals of the race. Of Adam too it holds that, by virtue of the ideal
of innocence, he could have avoided sinning \textit{(potuit non peccare)}.
Every attempt to ground the necessity of Adam's sin, of the continually
repeated fall in the individuals destined for freedom, is, according to what we
have seen, only an attempt to explain sin away. Since sinlessness in the human
self begins as possibility, temptation must enter in order that the possibility
may differentiate itself and pass over into actuality. Adam was tempted and
fell because he \emph{willed} to fall. He willed to fall because he did not
originally will to resist the temptation; for
% ---- printed p.285 (PDF 298) ----
\opage{285} he could indeed have resisted it: \textit{potuit non peccare}. That the old
Adam now afterwards explains it otherwise, namely that he could not resist the
temptation, that he really would not have fallen if he could have resisted, is
an equivocal evasion, an unclear mixture of excuse and self-accusation, which is
invalid, and invalid above all for the reason that it comes afterwards. Christ
was tempted and stood firm in the temptation; he conquered because he willed to
conquer: he did not sin because he did not will to: to that extent the same
conditions hold for Christ as for Adam.

If, however, Christ's sinlessness were merely to express the realization of
that whose possibility human nature, regarded from the standpoint of the ideal,
bears within itself, then Christ, the true man, could well be said to be the
moral exemplar of the race, the awakener of the powers of will slumbering in
the individuals, but by no means the Savior of the world. For Christ to be the
world's Savior, the Redeemer of the race, his sinlessness must have another
ground than Adam's innocence: freedom in him must essentially be conceived from
the divine side \textit{(non potuit peccare)}. The qualitative difference
between Adam's and Christ's state of innocence, between Adam's and Christ's
temptation, in a word: between Adam's and Christ's self, consists originally in
this, that whereas the germ of selfhood in Adam is naturally selfish, the germ
of selfhood in the Christ-child is, on the contrary, supernatural, a
Logos-germ, which through the Spirit has entered into interaction with the
natural matter and has organized the body. How a principle of life can in
general form and fashion a bodily organism shall not be investigated here; here
it shall only be pointed out, on behalf of faith, that the organizing activity
of the Logos-principle is, from the standpoint of natural science, just as
comprehensible and just as incomprehensible as the activity of any other
principle of life in nature. That
% ---- printed p.286 (PDF 299) ----
\opage{286} the Logos-principle works as an organizing principle of life means that the
Self in the Son is from the germ onward not merely natural but also holy, not
merely human but also divine.

\greekrun{β}{The Divine Self.}

The essential difference between Christ, the true man, and every other human
being is decisively marked by this, that the true man has himself presented
himself as a true God, as God's own Son, the Son who has come down from heaven,
as divine Self. This difference of essence shines forth, among other things,
from the fact that he remembers a pre-existence: ``before Abraham was, I am''
(John 8:58). Such a recollection of pre-existence must, as far as human
individuals are concerned, necessarily either lose itself in abstract
speculations or dissolve into mythical fantasies. With mythical fantasies
nothing is accomplished. What is appealing in Plato's beautiful philosophical
myth of the pre-existence of souls lies by no means in the depiction of the
souls' existence with the gods preceding birth, but in the ingenious, conscious
and yet naive visualization of the a priori in every psychical-rational self.
How great a blunder is committed when someone seeks to explain Christ's
pre-existence to himself by holding to the representation and taking the
Platonic myth as basis, we learn from Origen. In Origen's doctrine of
pre-existence attention is particularly fixed on the self-being of the
pre-existent souls in objective immediate singularity. The Logos, in contrast
to God himself, is held fast in so self-enclosed an individual existence that
he is even represented as a god by himself, a subordinate god
(\textit{δευτερος θεος}). For in Origen the Son is ``not the highest God, but
less than the Father, the second God, and stands in a certain way midway
between the nature of the Original and of the Non-Original. He does not
cognize the Father so
% ---- printed p.287 (PDF 300) ----
\opage{287} perfectly as the latter cognizes himself; the world is not created by him, but
through him by a higher principle than he, by the Father, and just as little
may one pray to the Son, who indeed himself prays, but through him to the
Father''. --- ``Origen first,'' says Strauss, ``was led by the necessity, partly
of finding a subject for Christ's bodily and mental sufferings, in general for
the finite form of his inner life as well, outside the unchangeable and
unbounded Logos, partly of finding a connecting middle term between this and
the material body, to distinguish definitely a human soul in Christ from his
higher nature; a distinction which, in connection with his theory of
pre-existence, also served him to answer the question how and why just this
soul had come into such a union with the divine Logos. From the beginning,
i.e.\ already in pre-existence, it was supposed to have clung to the Logos in a
quite different way (from free impulse, as goes without saying in Origen's
system), and through this devotion it was supposed to have been absorbed into
him ever more and more, and to have become one spirit with him --- to be sure
only in the way that, according to 1 Cor.\ 6:17, everyone who cleaves to the
Lord is one spirit with him; and thus Christ and the Logos were determined as
two persons who are one only as far as the moral is
concerned''\footnote{Saml.\ I Bd.\ S.\ 380 flgd.\ med II Bd.\ S.\ 84--85.}.

It is the way of objective knowledge in the doctrine of faith that we have here
trodden with Origen. That this way is impassable when a solution of the problem
of pre-existence is the goal, surely no one will deny; the result to which
investigations in this direction, despite all orthodox speculations and
explanations, must consistently lead is by no means even approximately an
understanding of the relation between God the Logos and the man Jesus, but on
the contrary so
% ---- printed p.288 (PDF 301) ----
\opage{288} bottomless a misunderstanding that the critique of knowledge need only follow
the debates of the history of dogma in order step by step to point out the
contradictions and triumph over the dissolution. To shelter oneself against the
flood tide of the consequences behind a dike of Scripture passages is
impossible. For Scripture passages there is always an interpretation, and as
the principle of cognition is, so will the interpretation be. Those who would
rightly impress upon themselves that the so-called Scripture principle is,
strictly speaking, no principle at all, need in this connection only note the
exegetical tactics with which, e.g., the Socinians went to work in order to
explain away Christ's pre-existence, a tactic which Strauss has known how to
appraise excellently; one need only cast a glance at the sample he gives of an
interpretation which Faustus Socinus has given of John 8:58; 17:5 and of the
Johannine prologue\footnote{``When Jesus in John 8:58 says to the Jews:
\textit{ἀμην, ἀμην λεγω ὑμιν, πριν Αβρααμ γενεσθαι ἐγω εἰμι}, then, first,
according to F.\ Socinus, \textit{πριν Ἀβρααμ γενεσθαι} can just as well mean:
before Abraham came into being, as: before he came into being;
% print: "γενεθαι" for γενεσθαι, in both places.
% as printed: both clauses read "blev til"; Strauss's contrast is presumably "was" / "shall be".
secondly, these words must not be connected with \textit{ἐγω εἰμι}, but with
\textit{λεγω ὑμιν}; and thirdly, \textit{ἐγω εἰμι} must not be understood of
mere existence, but of the dignity of the Christ: Verily, before Abraham comes
into being, I say unto you that I am the Messiah. But Abraham had indeed not
merely come into being but also passed away many hundreds of years before; how,
then, can it be set down as something future that he is to come into being? The
man with the name Abraham, answers Socinus, had to be sure come into being long
ago, but that had come into being, or come to pass, for the sake of which God
had given him this name, namely that he should become the father of many
nations (Gen.\ 17:5),
% as printed: no negation ("men det var blevet til"); the sense requires "had not yet come into being".
i.e.\ according to the Apostle's interpretation (Rom.\ 4:11) \textit{πατηρ
παντων των πιστευοντων δι᾽ ἀκροβυστιας}, since the gospel had not yet been
% print: "δί" (acute over the iota) for δι᾽.
offered to these, but was first to be so on account of the unreceptiveness of
the Jews. To these, therefore, Jesus here gives the earnest admonition that
they should acknowledge him as the Christ before the Kingdom of God is taken
from them and given to the Gentiles''.

``At John 17:5 F.\ Socinus struck out on another path, which Sam.\ Crell and,
later, the rationalists also used for the passage just mentioned. How can one
be so blind, he asks, as not to see that Christ here means to say nothing other
than that it had already been resolved with God before the creation of the
world to bestow on him in due time the glory for which he now prays''?

``Yet a tightly tied knot of exegetical difficulties awaited the sword of
Socinian hermeneutics in the Johannine prologue. We shall see whether it has
managed to cut through it, or whether it has not rather got hacked to pieces in
it. --- V.\ 1. \emph{In the beginning}. This must be understood according to
the context. Just as Moses (Gen.\ 1:1) meant by this expression only the
beginning of the things whose origin he there intends to speak of, and those
were the things of the visible, not those of the invisible world, so John too
can here mean only the beginning of that of which he intends to treat, namely
the gospel. Not before anything at all was created, but before the new world of
Christianity took its beginning, \emph{was} already --- \emph{the Word}. Just
as the Baptist is metonymically called a voice, so Christ is called the Word as
the proclaimer of the divine word or the gospel, and it is said of him that in
the beginning, i.e.\ at the time when the Baptist called the people of Israel
to repentance, he was already present and chosen by God to proclaim his will.
But the Word --- \emph{was} in the beginning \emph{with} God; i.e.\ before
Jesus became known to the Jews through the Baptist, his higher determination
was known to God alone; but this determination was of such a kind that it could
with justice be said, in the wider sense established above: \emph{the Word was
God}. --- V.\ 3. \emph{All things were made by it}, etc.: all the spiritual
and divine effects which Christianity produced in the world have their ground
in Jesus as the proclaimer of the divine word. --- V.\ 4. \emph{In it was
life}, % print: "V 4" (no stop).
inasmuch as everyone who follows Jesus will attain eternal blessedness and be
raised by him to immortality; he \emph{was the light of men}, in that by his
teaching and his example he gave them light. --- V.\ 10. \emph{It was in the
world, and the world was made by it}: i.e.\ Christ appeared among men,
% print: "blaudt" (turned n) for "blandt".
who through his activity were transformed in a moral respect. --- V.\ 14.
\emph{And the Word became flesh} --- on the contrary, Christ, the proclaimer of
the divine word, \emph{was} already flesh or man at the time spoken of, and so
therefore must \textit{ἐγενετο} be translated, as it also can be with justice;
and it has been added so that one should not take the being from which such
great effects proceeded for a superhuman one; but precisely this word, which was
meant to forestall the dreams of a higher nature in Christ, the devil caused men
to misunderstand''. Anfr.\ Skr.\ I Bd.\ S.\ 410--413.}.
% ---- printed p.289 (PDF 302) ----
\opage{289}


Only when the recollection of pre-existence is regarded from the standpoint
of the self's relation to the truth will its
% ---- printed p.290 (PDF 303) ----
\opage{290} significance show itself; but then the difference between a knowing
and a believing conception will also be conspicuous. For the knower it holds
that the self and the truth correspond to each other, not \textit{in
concreto} but \textit{in abstracto}, not in actuality but in essentiality,
in all generality. Through the truth the knowing self, purified of accidental
individuality, goes back into pure selfhood and thereby into pure eternity.
Recollection of pre-existence is in this sense a metaphysical dissolution of
all existence: the knowing selfhood does not exist, it \emph{is}; it
recollects its eternal being by viewing itself \textit{sub specie
æternitatis}. The Socratic thesis that all knowing is a recollecting is, says
S.\ Kierkegaard, ``an intimation of the beginning of speculation; but for
that very reason Socrates did not pursue it either; essentially it became
Platonic \dots\ For Socrates, therefore, the thesis that all knowing is a
recollecting signifies \dots\ two things: 1) that the knower is essentially
\textit{integer}, and that for him there is no other difficulty with regard
to the cognition of the eternal truth than this, that he exists, which
difficulty is so essential and decisive for him that it signifies that
existing, inward deepening in and through existing, is the truth; 2) that
existence in temporality has no decisive significance, because there is
constantly the possibility of taking oneself back into eternity by
recollection, even though this possibility is constantly cancelled by the
fact that inward deepening in existing
% ---- printed p.291 (PDF 304) ----
\opage{291} fills up time. --- Thus subjectivity, inwardness, is the truth; is
there now a more inward expression for it? Yes, if the saying: subjectivity,
inwardness, is the truth, begins thus: subjectivity is untruth. \dots\
Speculation also says that subjectivity is untruth, but says it in exactly
the opposite direction, namely toward the claim that objectivity is the
truth. Speculation determines subjectivity negatively toward objectivity.
% print: „Subjectiviten“ for Subjectiviteten
The other determination, on the contrary, obstructs itself in wanting to
begin, which is precisely what makes inwardness far more inward. Socratically,
subjectivity is untruth if it will not grasp that subjectivity is the truth
but, e.g., wants to be objective. Here, on the other hand, subjectivity, in
wanting to begin to become the truth by becoming subjective, is in the
difficulty that it is untruth''\footnote{Afsluttende uvidenskabelig
Efterskrift. S.\ 151--153.}.

The relation between the self and the truth is quite different in rational
ethics from what it is in revealed religion. The demand of existence is
indeed as decisive for ethical subjectivity as for religious; but insofar as
the universal-ethical truth is rational, the self in its striving constantly
proceeds from the presupposition that it is originally in the truth, and
that its self-development is therefore a development of truth, whereas the
religious self in its relation to God must cognize that it is essentially in
untruth, insofar as it is in sin. The rational-ethical subjectivity, whose
troubles and exertions for the truth's sake are owed only to the temporal,
not to the eternal, can therefore, if it will, withdraw from temporality,
rest in the thought of eternity, and find relief by identifying itself,
through a metaphysics of the recollection of pre-existence, with ``the
Absolute''. For the religious self such relief is impossible. Self-cognition
is in
% ---- printed p.292 (PDF 305) ----
\opage{292} faith inseparable from the cognition of sin; the cognition of sin,
the self-contradiction concentrated in an inescapable demand of existence,
signifies that the strife between the self and the truth, with all the
conflicts belonging to it, concerns not merely the temporal but, in the very
strictest sense, the eternal; for the eternal threatens. For the believing
consciousness the eternity lying behind is closed, and every recollection of
pre-existence unconditionally cut off. But Christ, whose consciousness is the
light of revelation, Christ could not possibly be the true, the sinless Self,
had he not been able to recollect his pre-existence.

``I am the truth''; ``before Abraham was, I am'': these two sentences
correspond so completely to each other that if the first fits Christ, then
the second does too. It is not the universal Idea of truth, but the holy,
personal truth, the truth that is eternally one with the true Self, that is
spoken of here. In that Christ personally knows himself to be the
truth revealed as man, he knows himself \textit{eo ipso} to be God revealed
to men, the divine Self, the Son of God from eternity. That in his discourses
he nevertheless presents himself as the one who speaks the truth, bears
witness to the truth, and so again makes a distinction between himself and
the truth, is comprehensible when we think of the existential opposition
between phenomenon and essence, between the self as existing individual and
the inner infinity indwelling the self, which contains all ``the fullness of
the Godhead''. Moved by the Spirit, a prophet too can speak the word of truth
and bear witness to the truth; but as existing self the prophet is still only
a human organ for the divine truth and therefore cannot, even with the
highest inspiration, make himself one with the truth; if the inspired one for
a single moment confused the human with the divine and
% ---- printed p.293 (PDF 306) ----
\opage{293} gave himself out as the truth, he would be condemned. Therefore none of the
prophets in the Old Testament has anywhere hinted at any recollection of
pre-existence or dared to say: ``before Abraham was, I am''. A prophet too
can indeed speak with divine authority; but when he uses his highest
authority, he speaks not in his own name but in the name of Jehovah. But
Christ, the divine Self, speaks with the highest authority in his own name,
speaks from God in speaking from himself: ``Verily, verily I say unto you''.

The most telling expression of Christ's acting as divine Self we have in his
word to the paralytic: ``Thy sins be forgiven thee''. The scribes, who in him
saw and heard only the man, said within themselves: ``This man blasphemes
God'' (Matt.\ 9:2--3). This was, when we enter into their way of thinking, at
bottom quite correct. That one man cannot in his own name pronounce the
forgiveness of sins to another, insofar as both are sinners, goes without
saying; if it happened, it would seem as parodic as if two prisoners on the
way to the examining magistrate were to take turns comforting each other with
assurances of acquittal. Even if the one knew himself innocent of the crime
charged against him, he could still not possibly acquit the other. Just as
little could Christ, even if he knew himself sinless, as man impart
forgiveness of sins to any other man. To present oneself as the one who in
his own name has authority and power to forgive sins is to present oneself as
the divine Self, as the one who is before Abraham was, i.e.\ as the eternal
Judge.

The critique of knowledge has made a great to-do about the great discrepancy
that is supposed to exist between the Synoptics and the author of the fourth
Gospel in the presentation of Christ's person. Now, that not only in words
and
% ---- printed p.294 (PDF 307) ----
\opage{294} mode of expression, but also in what we may call the fundamental
representations and characteristic categories of the philosophy of religion,
a difference shows itself, must be granted: John has a deeper eye for the
Logos, for the Son's eternity, for his going forth from the Father, than the
Synoptics. But this discrepancy does not touch the principle of faith%
\footnote{``The first three Gospels conceive Christ's divine glory
essentially under the prophetic standpoint of eschatology, or, if we may so
express ourselves, under the standpoint of his post-existence. \dots\ But in
this conception the thought is latently present which John clearly
expresses: that he who is the \emph{Last}, who in his coming is to be over
all in creation, over all in heaven and on earth, must also be the `First',
must be before all creatures; that he to whom we must in this sense ascribe
post-existence must also have pre-existence.''
% print: the outer quotation is never closed (no mark after „Præexistens.“)
H.\ Martensen. ``Den christelige Dogmatik''. S.\ 247--48.}. The Johannine
Gospel marks no transition whatever from the immediately religious
representation to speculation. The speculative element that people have
wanted to foist on John is a semblance: the Johannine principle is from first
to last a principle of will; the relation between the Father and the Son is
unconditionally a relation of will, and Christ's recollection of
pre-existence as antirational as possible.

When Christ has healed the sick man at the pool of Bethesda (John 5:1 ff.),
and the Jews, who take offense at this, ``sought the more to kill him,
because he not only had broken the sabbath, but also called God his own
Father, making himself equal with God'';
% print: the quotation is not closed after „Gud lig;“
there follows a declaration in which the Savior expressly presents himself as
the one sent by the Father, as the Judge of the world; for ``the Father
judgeth no man, but hath given all judgment unto the Son: that all men should
honour the Son, even as they honour the Father''. The unity of essence of the
Father and the Son of which the discourse treats is, as anyone can see, by no
means rationally grounded in the metaphysical
% ---- printed p.295 (PDF 308) ----
\opage{295} necessity of the divine self-objectification, but has its
antirational ground in the religious-ethical unity of freedom, in the unity
of self, the original unity of will; for, testifies the holy voice whose
testimony is true, ``I seek not mine own will, but the will of the Father
which hath sent me''.

\greekrun{γ}{The Divine-Human Self: The Mystery.}

That unity of self rests essentially on a self-doubling conditioned by an
other is a principle whose validity modern philosophy has made so evident
that it can no longer be called into doubt. But just as the unity of self,
when it is to be more closely determined, is of different kinds, so too the
self-doubling is of different kinds. There is, to keep to self-consciousness,
an essential difference between the merely formal and the real doubling. The
formal one is easy to see through. To become conscious of oneself is to
posit the self at once both as subject and as object; the doctrine of formal
self-objectification already belongs among the worn-out topics and will not
be further developed in this context. It is otherwise with the existential
reality of self-doubling. On this important point great obscurity still
prevails in philosophy, and for the simple reason that the metaphysics of
abstraction, which one-sidedly looks at the universal, all too easily makes
short work of turning subjectivity into a theory and reducing the fundamental
act of self-objectification to a general act of thought. If this
one-sidedness is already harmful to the philosophy of reason, which thereby
imperceptibly sinks down into an empty philosophy of words, it is still more
harmful, indeed utterly ruinous, to the philosophy of religion.

In the philosophy of religion the unity of self must unconditionally be
conceived as personal unity of will, and the corresponding doubling posited
not as a mere act of thought but originally as an act of will. How matters
stand in this respect with the existential doubling of the self,
% ---- printed p.296 (PDF 309) ----
\opage{296} and how the doubling depends on the difference in the self's
centralization, has, as far as the consciousness of sin is concerned, already
been touched on in the Introduction. Such concepts as repentance and regret
would indeed lack all deeper reality if the pain of repentance were not
grounded precisely in the fact that the same existing self cognizes itself as
a double self; for the self that has posited the transgression and the self
that would blot it out are the same and yet not the same. Only a superficial
regret can demand a changed action without demanding a changed self.

The deeper the self, the richer the inwardness, the more peculiar the
phenomena of self-doubling that may arise. The mystic's self, e.g., may be
now as in heaven, now as in hell; may, without his therefore in a finite
sense going out of himself and madly imagining that he has several selves,
live through wondrous crises, in that the ideality of the self, through
successive and diverse concentrations, makes a series of transformations of
consciousness possible. Where rich spiritual gifts and religious enthusiasm
are united, a sober observer will always find it extremely difficult to get
to the bottom of another's inner life. We can imagine a man, e.g.\ a poet, so
richly gifted that his productivity is, literally speaking, a riddle even to
himself; we can imagine a spark of holy fire falling into such a soul and
kindling it into flame; but to measure the depths of subjectivity, to unveil
the secret of productivity and discover the hidden wellspring of enthusiasm,
this we cannot do. We are taken by surprise, we are almost gripped by fear of
madness, when a hitherto unknown inwardness wells forth, incalculably
swelling, when a religious enthusiasm that knows no bound suddenly breaks out
of its concealment and would give itself vent in transfiguration. The more
peculiar the relation to God is, the more it withdraws from every control of
a scientific kind. Truth and self-deception,
% ---- printed p.297 (PDF 310) ----
\opage{297} wisdom and derangement can enter into the most diverse
combinations in a religiously moved interior and deposit the most wondrous
phenomena: the more gifted the individuality is, the harder it is to see
through its ground of life and bring the development of its character under
universally valid determinations. The one moved in the Spirit can be measured
only by his own measure, and the stronger, deeper and more original the
movement, the more exceptional the standard. Even when self-reflection, with
a psychologically observing, scouting gaze, has been driven to an extreme,
and a religious-ethical formation of character carried through with strict
vigilance and impressive consistency, the inwardness in the relation to God
cannot possibly be expressed in direct communication. ``I would'' --- says one
wondrously gifted, just where he is to speak of his relation to God, of what
has taught him to marvel at God, his love, and at what a man's impotence is
able to do with his assistance --- ``I would, more resolved than that king who
cried: my kingdom for a horse, and blessedly resolved as he was not, give
everything, my life included, to find what it is more blessed for thought to
find than for the lover to find the beloved: the expression'' --- the
expression of thanksgiving --- ``so as then to die with this expression on my
lips. And lo, it offers itself: the thoughts as enchanting as those fruits in
the garden of the fairy tale, so rich, warm, inward, the expressions so
soothing to the need of thought in me, so cooling over the hot longing --- it
seems to me that if I had a winged pen, yes, if I had ten, I would still not
be able to keep pace quickly enough with the riches that offer themselves.
But as soon as I take my pen in hand, I cannot for the moment stir it, as one
speaks of not being able to stir a foot''\footnote{S.\ Kierkegaard.
Synspunktet for min Forfattervirksomhed en ligefrem Meddelelse. Rapport til
Historien. Kbhvn., 1859. S.\ 51.}.

% ---- printed p.298 (PDF 311) ----
\opage{298} So long, however, as the relation to God keeps to a derived
stage, so that the religious individuality essentially has its life in an
appropriation of the revealed word, without itself daring to present itself
as a special, independent organ of revelation, analogies from the
universally human will still be able to furnish orienting points of support,
and a rational psychology, although the hidden miracle of faith will always
withdraw from scientific inquiry, will be able to make negative
contributions in determining such limits as may not be overstepped, if
actual phenomena of faith are to be distinguished at all from phenomena of
eccentric fanaticisms and uncanny aberrations of spirit. If, on the other
hand, an exclusive, independent organ of revelation appears, e.g.\ a prophet,
then the mystery in such a self-consciousness will correspond categorically
to the revelation. But where mystery and revelation correspond categorically
to each other, so that the mystery halts knowledge while what is revealed
enriches faith, there every standard taken from humanity and rational
psychology is inapplicable. Of penetrating cognition of a self-consciousness
of a different kind it holds first and last: only like is understood by like.
If already a prophet's religious self-consciousness, by virtue of what has
been entrusted to him, is qualitatively different from every other believing
consciousness, then Christ's self-consciousness must surely be in an
absolute sense different even from the prophet's. The Son, the Only-begotten,
knows himself --- realiter --- to be God, and knows himself at the same time
--- realiter --- to be man: this self-doubling is in its divine-human
reality an impenetrable Mystery.

Nevertheless the spokesmen of faith and of unbelief alike, each from his own
side, have with great importunity attempted to work their way through the
Mystery. The confusion of thought has moreover been increased by the fact
that, instead of holding the task fast in the spiritual form of
self-consciousness, hence of self-certainty and ideality, people have had it
% ---- printed p.299 (PDF 312) ----
\opage{299} transformed into a piece of philosophy of nature, a
physiological-hyperphysiological task, by introducing the doctrine of
Christ's two natures.

The starting point of the christological conceptual aberrations must
originally be sought in the speculations on the Logos. However different the
representations were from the beginning, there was nevertheless general
agreement in regarding the Logos as an independent being. The representation
of a distinct Logos-self stood firm as an indisputable presupposition when
the doctrine of the Logos as the uncreated Son, begotten of the Father from
eternity by necessity, was declared orthodox. So the task was to explain to
oneself how the eternal, almighty, omniscient and omnipresent Logos-self
could at a given point in time take on human nature and let itself be born
as a man. This task is insoluble. Here, for the objective view, a double
Logos-self takes shape, one in God and another in man; the former is from
eternity, the latter has a beginning in time; the former is being, the
latter becoming; the former is almighty, the latter suffering; the former
remains in heaven while the latter walks on earth: instead of one Son of God,
speculation, however reluctantly, gets two Sons of God,
% print: „Sønnner“ for Sønner
one who has not become man and one who has become man. If we next consider
the Logos-self that has become man, the difficulty repeats itself with
respect to the self of the man. If God the Logos has actually himself become
man, a true, complete man, then in the God-man there must surely be either a
double self, a Logos-self, for this man is the Son of God, and a human self,
for this Son of God is a true man; or else the Logos-self must exclude the
human self. But whether one now assumes that in the God-man there is a double
self, or that the Logos-self has displaced the human self and put itself in
its place, the result is equally absurd.
% ---- printed p.300 (PDF 313) ----

\opage{300} When the limits of the cognition of faith are overstepped and the
inseparable unity of the two separated natures is to be exactly fixed, the
critique of knowledge at once comes forward with objections such as these.
``If the Son of God, an almighty being, the ruler of the world, eternal,
equal to himself in actuality, was to be united with a finite man who had
come into being, then only two ways were open: either the God had to
contract himself into human finitude and submit to human becoming, so that
he had actually slumbered unconscious in the mother's womb along with the
man, after birth developed successively, and for his lifetime renounced
participation in the government of the world; but thereby his divinity was
sublated, and the Church, which needed it, consequently could not go this
way. So there was then the other way: while the man developed successively
from embryo to adult, the God remained equal to himself in his absolute
actuality; while the man slept unconscious in the ship, the God was
omniscient; while the man hung on the cross and lay in the grave, the God
governed all things. This is the general basis of the ecclesiastical
representation; but thereby, to be sure, the unity of Christ's person would
be completely torn apart. Therefore his human nature was also made to share
in the prerogatives of the divine; thereby the still more glaring
contradiction was posited, that one was now to think of this one human
nature on the one hand as raised above successive becoming, as omniscient and
almighty, and yet on the other hand as successively becoming and limited in
power and knowledge. This latter it never actually was either, but
constantly only the former, though in a hidden way, said the one party: ---
and thus it dissolved the whole human life of Christ into a deceptive
jugglery. It did possess those prerogatives of the divine nature, said the
other party, but as a rule abstained from using them: --- so Christ was
after all sometimes, as
% ---- printed p.301 (PDF 314) ----
\opage{301} man, with the human structure of the brain etc., omniscient and
almighty; indeed, his human nature could just as well already in the
mother's womb, consequently as an embryo, have known everything, governed the
world --- it was only the fault of his free will when he did not do
it.''\footnote{Strauss Troesl.\ II.\ Bd.\ S.\ 117}


Once the doctrine of faith has crossed over onto the preserves of objective
knowledge, it is too late to take refuge in the incomprehensible and appeal
to the Mystery when one is struck by the arrows of critique; for the boundary
has been overstepped. The Mystery is the boundary, the categorically decisive
boundary for the thinking of faith as well as for that of knowledge. The
boundary is determined by the concepts' own consistent self-dissolution. That
the unity of the self in Christ with the eternal Logos-self is a unity of
principle, faith must hold fast; but how Christ's consciousness was able to
convert this unity of principle into a concrete personal unity of self is and
remains an impenetrable Mystery. What help is it, then, to take refuge in
phrases such as these: ``Here there are not two Sons of God, but one Son of
God; with Christ no other Son is added to the Trinity; the whole movement takes
place within the circle of the Trinity,'' and so on, when it is nevertheless
impossible in the objective sense to think, or even to form a representation
of, how the Logos-self in God and the Logos-self in man could, despite the
separation, be united within the circle of the Trinity? Or is there really any
thinker who will seriously undertake to make clear to us a doubling such as
this, that one and the same self can at one and the same time be both
omniscient and yet not omniscient, both almighty and yet suffering? If there
cannot be two Sons of God here, then still less can the one Son of God be
thought to have divided himself into two half Sons. It is surely a monstrous
notion that the one half Son, the one half Self of the Son, ``the pure
Logos of the Godhead,'' should
% ---- printed p.302 (PDF 315) ----
\opage{302} pervade ``the kingdom of nature in all-filling presence,'' while
the other half Son, the other half Self of the Son, could only at a particular
time be present in a particular place, be it in Nazareth or in Jerusalem. What
help is it to want to correct the old dogmaticians, who, in order to maintain
the unity of self, represented to themselves that the divine properties were
communicated by \textit{communicatio idiomatum} --- to the human nature in all
infinity, when there is not the slightest scientific prospect of getting a
hair's breadth further with ``the rightly understood concept of the divine
self-limitation'' (\textit{κενωσις})? Just one example: Christ foretells the
end of the world and adds (Mark 13:32): ``But of that day and hour knoweth no
man, neither the angels which are in heaven, neither the Son, but the Father
only,'' and Quenstedt explains the passage: ``I, Christ, as man, now in the
state of humiliation make no use of the divine omniscience, which I
nevertheless possess, so that I now in actuality do not know the time of the
last day.'' How will one correct Quenstedt in an objective scientific doctrine
of faith? If one says that Christ in his state of humiliation actually did
not know the day and % print: "Dagen eg Timen" ("eg" for "og").
the hour, while ``the pure Logos of the Godhead'' in heaven knew both day and
hour, then we surely get a double Logos-self, ``a Christ with two heads,'' the
one with limited, the other with unlimited knowledge.

Now if it were possible to cut away the spongy excrescences of speculation
with a single stroke, or at least to weed out the weeds of the science of
faith as one weeds thistles with a sharp iron, then it would undeniably be the
most practical method to declare ``the God-man'' ``a paradox'' and the paradox
``an absurdum.'' That is the Kierkegaardian way out, which, as regards the
inward deepening of the problems of existence, certainly has its right. But
the thought of faith, as thought, has its right too. The
% ---- printed p.303 (PDF 316) ----
\opage{303} religious conception of the God-man must, as surely as it is to be
at one with an essential self-understanding, necessarily presuppose a
believing conception of the divine-human Self; but a believing conception is
conditioned by the thought of faith; now if the thought of faith is broken off
just as abruptly as the thought of knowledge, then one surely pulls up the
wheat together with the tares, and things look doubtful for the weeding.

S.\ Kierkegaard reduces the religious conception of the God-man to the
conception of a divine-human incognito. The God-man is God, but chooses to
become the individual man. ``This is the deepest incognito or the most
impenetrable unrecognizability that is possible; for the contradiction between
being God and being an individual man is the greatest possible, the infinitely
qualitative one. But it is his will, his free resolve, and therefore an
almightily maintained incognito. Indeed, by letting himself be born, he has in
a certain sense bound himself once and for all; his unrecognizability is so
almightily maintained that he himself is in a way in the power of his
incognito, wherein lies the literal \emph{actuality} of his purely human
suffering, that this is not a mere appearance, but in a certain sense is the
over-power of the assumed unrecognizability over himself.''\footnote{Indøvelse
i Christendom af Anticlimacus. 2det Oplag. 1855. S.\ 141.}

Whoever rightly reflects on what is said here will see that there is a
difference between a mystery and an absurdum. The Mystery dissolves the
contradiction by dissolving the contradictory concepts. In an absurdum, on the
other hand, we have a hard knot of fixed, undissolved and yet incompatible
conceptual determinations, of eternal contradictions, indissoluble in
themselves. ``That an Almighty One'' ``lets himself be born,'' ``that an
Almighty One voluntarily posits an incognito so almightily that the Almighty
himself is in the power of his incognito''---such screaming contradictions
belong not to the Mystery but to the pure
% ---- printed p.304 (PDF 317) ----
\opage{304} absurdum. The reserving expressions ``in a way,'' ``in a certain
sense,'' give no understanding in the context cited. The author must surely be
right when he adds: ``It is a strange kind of dialectic, that he (the God-man)
almightily --- binds himself; and he does it so almightily that he actually
feels himself bound, suffers under the consequence of his having lovingly and
freely resolved to become an individual man.'' Indeed, it is a strange kind of
dialectic that the divine-human Self, the Self in which all cognition of God,
all true self-cognition is to be grounded, is not merely a sign that is spoken
against by others, but essentially a sign that contradicts itself.

The divine-human Self of the Son is no paradox; but it is a Mystery.

\runhead{b) Fundamental Determinations in the Self of the Son.}
\parmark{20}

\emph{Jesus Christ, God's only-begotten Son}, is not an object of knowledge
but the object of faith. As the object is, so must the objectivation be; just
as the objects of knowledge presuppose an objectivation in knowledge, so the
objects of faith must presuppose an objectivation in faith. Now as surely as
the principles, that of faith and that of knowledge, are absolutely
heterogeneous, so must the objectivation of faith also be absolutely
heterogeneous with that of knowledge. % print: "Objectiverering" (doubled -er-) for "Objectivering".
The absolutely heterogeneous lies, as we have seen, not immediately in the
difference of the object, but essentially in the transformation which the
differently constituted objectivation brings about. Now all objectivation in
faith has this peculiarity about it, that self-understanding and understanding
of the object condition each other in inner infinity. The sinful one must know
himself by the sinless One, and conversely; the one darkened in temporality
must know himself by the One transfigured in the light of eternity, and
conversely;
% ---- printed p.305 (PDF 318) ----
\opage{305} this is objectivation in faith. The like and the unlike, the common
and the different, can only through such an objectivation come forth and
obtain their full significance. This, then, is the likeness, the common
element: that the Christ-self, the holy Logos-self, which ``sinks itself into
the womb of humanity as possibility, as holy seed'' and lets itself be born of
a woman, begins from the germ of its becoming its development in time in a
manner corresponding to every natural human self; while the unlikeness and
difference is that the selfish self, lost in finitude, forgets eternity,
whereas the holy Self, striving in finitude, unceasingly remembers its
eternity and makes manifest the eternal life of the self. This is the
likeness, the common element: that the Christ-self, the holy Logos-self, does
not have its life immediately of itself, but of God in its relation to God, in
a manner corresponding to every believing human self; but the unlikeness and
difference shows itself in this, that the true original relation to God, the
Son's life of the Father in the Father, is his eternal life in himself, while
the believer can have a derived life in himself only by appropriating the life
in the Son and through the Son of the Father. This is the likeness, the common
element: that the man Christ, the object of faith, is as an existing man with
human limitation a single self-luminous point in the outer infinite universe,
in a manner corresponding to every man developed to consciousness; but the
unlikeness, the thoroughgoing difference, is that Christ, the original light
of truth, of life, is the sun of spirits for time and eternity, while the light
of eternity that is kindled by faith in man's soul is a derived light, a ray
from the sun of spirits.

\greekrun{α}{The Self of the Son Is Logos.}

Originally the gaze of faith is to be fixed on Jesus of Nazareth, appearing as
an actual man and yet as Messiah.
% ---- printed p.306 (PDF 319) ----
\opage{306} By hearing his words and seeing his works the disciples at once
received a double impression. With John it begins as if the understanding
came of itself; it begins with ideal expectations: ``Hereafter ye shall see
heaven open, and the angels of God ascending and descending upon the Son of
man'' (John 1:52). But that the contrast between the outward lowliness and the
inward glory worked strongly on the first disciples, and that it fermented and
labored in them before they could come to a clear and secure understanding,
of this the Gospels contain many and eloquent testimonies. It came to a
breakthrough in Peter when he (Matt.\ 16:16), to the Master's question, gave
the celebrated answer: ``Thou art the Christ, the Son of the living God''; but
how unsteadily the flaring light burned, and how easily the understanding
could go backward during the walk with Christ, Peter himself shows by his
example to the very last. Only after the Risen One had entered into his glory,
and the Spirit had been poured out upon his witnesses, did the image of the
Transfigured One stand living before them with all its divine-human features.
If we then gather the utterances of the Evangelists and Apostles about the
divinity of Christ into one total impression, they all agree in acknowledging
the Son of man as the Son of God, as the Firstborn of all creation, as Logos,
as the brightness of God's glory and the express image of his essence
(Heb.\ 1:3); this is the first fundamental determination in the Self of the
Son: \emph{the Son is Logos}.

If we now take note of the way by which the fundamental determination has been
reached, we learn what it means that the object of faith will be objectified
in faith. For the starting point of cognition is not set in God's eternal
essence, in the Trinity \textit{ad intra}; but the starting point is in God
revealed as man, in the living personal impression of the person of Christ.
One does not begin with God's immanent self-objectivation, the great problem
of necessity
% ---- printed p.307 (PDF 320) ----
\opage{307} in speculation; but one begins with the existing Self, with the
pure, the holy will, which alone is able to awaken the life of truth, of
freedom and of love in believing souls: one begins not rationally, but
antirationally.

Therefore in the apostolic objectivation a series of determinations is also
left out which scientific reflection constantly seeks but never finds, and
these are the metaphysical intermediate determinations, precisely those
determinations of essence which were to raise faith up to knowledge and give
the doctrine of faith the stamp of objective science. It is said briefly and
simply that the Logos was in the beginning with God and was God; but of the
original relation of the Logos-self to the God-self there is not even the
faintest hint. It is said briefly and simply that the Logos became flesh, but
about the relation between the self in the eternal Logos and the self in the
Logos become man not a word is said; and why not? Because the metaphysical
problem of the Incarnation did not exist at all for the Evangelist. The more
personally the matter is taken, the plainer it becomes that the metaphysics
of the problem falls outside faith. Just imagine how it would look if one of
the disciples, for example Peter, had turned metaphysical along the way and
had put the following question to Christ: ``Master, tell me honestly, how can
you, who are now the Son of God on earth, bodily present among us, at the same
time be the Son of God in heaven, the Almighty, Omnipresent One, who together
with the Father governs all the world---are you able to do this
\textit{κατα κρυφιν} or perhaps \textit{κατα κενωσιν}?''---and one will get
an impression of blasphemy, an impression that Christ would sternly have
rebuffed the questioning Peter, as on a given occasion he rebuffed the
admonishing Peter: ``Get thee behind me, Satan! thou art an offense unto me.''
Should a clever man of metaphysics find the rebuff natural on the ground that
``Christ, who was not a meta%
% ---- printed p.308 (PDF 321) ----
\opage{308}physician and consequently not yet fully transparent to himself in
the universal, naturally had to avail himself of his authority and snap at the
disciple because he was unable to answer his scientific question'': then an
enlightenment of knowledge of that nature will certainly not darken the
believer.

Scholastic orthodoxy, however, is not content with a religious development of
the content of religion; it demands metaphysically exact intermediate
determinations, it objectifies in knowledge, it determines the
indeterminable, it violates the Mystery.

How characteristic in this respect is the controversy between Arius and
Athanasius! They both proceed from the presuppositions of faith, they are both
agreed that the content of faith must allow itself to be thought, and so they
both also come to commit the same blunder, of confusing thoughts of faith with
thoughts of knowledge: we hear it at once in the fundamental question. Is the
Son \emph{created} or \emph{begotten} of the Father, is the Logos from God
with freedom or with necessity? To answer this question exactly, one would of
necessity have to have the two concepts, creation and begetting, perfectly
clear, and also be able to make the concept of God so objectively transparent
that the difference between a \emph{created} and a \emph{begotten} God became
entirely evident. But how impossible it is, from the standpoint of faith, to
comprehend God's inner, eternal self-revelation is noticeable already at the
beginning of the controversy. When Arius asserts that the Son is created, but
with the advantage over the other creatures that these all, the Holy Spirit
not excepted, are in turn created through the Son, it is evident that he does
not comprehend what he himself is saying. There are certain notions, in his
opinion offensive, of an emanation of the Logos from God, a division of God, a
fusion of the essence of the Father and of the Son, and the like, that he
wants to ward off; but that he has not made the concept of creation clear to
himself, and still less seen
% ---- printed p.309 (PDF 322) ----
\opage{309} through what lies in such concepts as a \emph{created God}, a
\emph{created Creator}, is unmistakable. But the same must be said of
Athanasius. Instead of seeing that the blunder of Arius consists precisely in
his demanding metaphysical intermediate determinations where they cannot be
given, and therefore arbitrarily setting up propositions that confuse the
religious representations, Athanasius himself treads the same path as his
opponent, only that he wants his own metaphysics, not his opponent's,
sanctioned. That Athanasius comprehends what he himself sets up as little as
Arius does can easily be proved. According to Athanasius the Son is to have
been begotten of the Father from eternity; but not to mention that the
category of begetting, as the Arians indeed amply emphasized, points
especially to a sensuous-bodily natural production, a division of the
Father's essence, the concept of a begetting from eternity is surely heavily
burdened with contradiction. To be begotten from eternity was accordingly,
for the Arians too, the same as not to be begotten. ``If the Son is not to be
begotten, since he already was, which is a contradiction, then he cannot
possibly have existed before his begetting.'' The eternal begetting is a
beginning without beginning, a production of something already produced, a
coming-to-be of that which does not come to be, but is. So Athanasius and his
friends had at last to admit that the begetting of the Son took place in an
incomprehensible manner (\textit{ἀῤῥητως και ἀνεκδιηγητως}), that in other
words the concept of begetting, applied to God, is an incomprehensible
concept%
\footnote{That some of the Church Fathers at least had an inkling that the
defenders of orthodoxy were forced by the inventions of the heretics onto an
unnatural path can be seen, among other things, from the following utterance.
\textit{Sufficiebat quidem credentibus Dei sermo} \dots\
\textit{Sed compellimur hæreticorum et blasphemantium vitiis illicita
agere, ardua scandere, ineffabilia eloqui, inconcessa præsumere.}
\textit{Hilarius de trinitate. II, 1.} Jvfr.\ Gieseler I, 2, S.\ 43.
--- So far is it from being possible to say that the \textit{Symbolum
Nicænum} clarified and strengthened the life of the congregation, that the
Greek Church, as Grundtvig has with full right emphasized, by making this
Symbol its foundation and, in accordance with its Alexandrian principle of
Scripture, exchanging faith in the living word of Christ for faith in
dogmatic metaphysics, brought about precisely the spiritual self-enfeeblement
under which it suffers to this day.}.
% ---- printed p.310 (PDF 323) ----
\opage{310}

So then: first one violates the Mystery by forcing metaphysical intermediate
determinations into the Confession of Faith; when it then turns out that these
intermediate determinations end in sheer contradictions, one appeals to the
Mystery --- the Mystery one has oneself violated.

\greekrun{β}{In the Self of the Son Is Life.}

``The Father hath given to the Son to have life in himself''; herein lies that
Christ gives the world a new life, in that he introduces a new Self into the
world. From the standpoint of ecclesiastical metaphysics this Self has
admittedly been mistreated to such a degree that negative critique has done
religion an essential service by making the mistreatment manifest. According
to the scholastic-objective doctrine of the Church, no personality of its own
is to be ascribed to Christ's human nature; but --- this is a Straussian
objection --- Jesus' humanity would then surely be mutilated, and he could not
have had any human spirit at all if, on account of the union with the Logos,
no personality could develop from his human nature. In that case no human
soul could develop either, nor any human body. ``For since it is evidently an
effect of the human spirit enclosed as a germ in the human embryo that it
develops into an ensouled body of definite shape and constitution, an
organism of quite another kind would have had to form for the Son of God when
he
% ---- printed p.311 (PDF 324) ----
\opage{311} as the germ of personality joined himself to a human embryo (just
as the gall wasp, when with its sting it bores into the branch of a rosebush,
produces no rose, but the moss-covered concretion that the common people in
Germany call Schafapfel)\footnote{Strauss. Troesl.\ II.\ S.\ 94.}.''

These and similar objections lose their significance when, in accordance with
the principle of faith, we place the qualitative difference between the
Christ-self and every other human self in sinlessness alone. The sinless Self
is, as Logos-germ, of divine nature, holy and pure as God's own Self. That the
interaction between the supernatural Logos-principle and the bodily matter in
a human embryo is an impenetrable Mystery is granted; but then we ask: what
does science have with which it can prove, not the improbability, for that is
understood, but the impossibility of such an interaction? Facts, as in the
fine example of the gall wasp and the rosebush, it surely does not have at
all. ``But concepts!'' Yes, empty formulas for the benefit of verbal critique,
meaningless contradictions such as: that a supernatural principle is not a
natural one, that Christ is not a man like other men, and so on, conceptual
contradictions which all creep forward under the mystification that we know
what we nevertheless do not know, have knowledge of what we nevertheless do
not have knowledge of.

As the true man and Son of God, Christ lives in himself only by living in God,
and lives in God in living in himself: he does not have a religion, for he is
religion. In religion itself the relation to God is the pure, ideal relation
of will; religion itself is life, eternal life. That man in his relation to
God is absolutely dependent on God, the believer can learn from Christ. For
although the Son testifies of himself that he and the Father are one
(John 10:30), yet he also testifies most
% ---- printed p.312 (PDF 325) ----
\opage{312} emphatically to his dependence, his absolute dependence on the
Father. He has nothing of himself, cannot speak of himself and can do nothing
at all of himself (John 5:19); but all that he is, he is of the Father, and
all that he speaks, and that he does, and that he is able to do, all of it is
given him by the Father. But in this absolute dependence he has \textit{eo
ipso} his absolute freedom, his inner infinite independence. His life of the
Father is his own life, his original living in himself. ``For as the Father
hath life in himself, so hath he given to the Son to have life in himself, and
hath given him authority to execute judgment also, because he is the Son of
man'' (John 5:26--27). The relation to God is, to speak in the language of the
schools, an infinitely rich relation of immanence; but while we lay weight,
decisive weight, on the immanence, we must not forget that this holy relation
of immanence is a relation of freedom, a relation of will in principle,
therefore --- antirational!

The divine-human life, the eternal life of religion, is life in the Logos; to
this Logos-life, doubled in itself, there now corresponds the doubled Self,
the Self in God and the Self in man, but in such a way that the Self in God,
by the living contrast to man as man, is one with the Father revealing himself
in the Son, while the Self in man, by the living contrast to God as God, is
the true man, the Son of God, who says of himself: ``He that hath seen me hath
seen the Father'' (John 14:9; 12:45). Objectified in faith, the life of Christ
is thus an original double life; to this double life there corresponds, as we
shall now see, the categorical doubling of the Logos-Self.

In order to do the Father's will on earth, the Son must have knowledge of the
Father's will, which is in heaven. Whence, then, does he as man have this
knowledge? He has it from himself; for the human Self in him, the
% ---- printed p.313 (PDF 326) ----
\opage{313} Logos-principle developed in time, is the revelation of the
Logos-Self in God, the eternal Self in which the Father has all knowledge and
will. And yet the Son does not have his knowledge immediately or once and for
all of himself; for in this sense he has nothing, nothing at all, of himself:
he has it only in that it is given him and in every moment is being given him;
that is to say: he has it in his relation to God: he has it in religion, and he
is religion itself. Both moments in the fundamental determination, both that
he has his knowledge of himself and that he has it from the Father, are
accordingly also presented in the evangelical objectivation in such a way that
the union and the separation come to view equally. He is the knower of hearts,
who sees through both Pharisees and Sadducees and ``sees their thoughts''; but
that he is so because it is given him is recognizable from this, that
although as man he is not omniscient, he nevertheless knows, when it is
needed for the completion of the revelation, what he could not possibly know
on human terms (John 11:4), while such knowledge in turn, where it is not
needed (Mark 13:52), is also not given him. % print: "Mr. 13, 52" (Mark 13 has 37 verses); probably meant Mark 13:32, quoted on p.302.
To draw a scientific inference from his human to his divine knowledge, or
conversely, is impossible, since his divine-human knowledge moves in and with
the revelation itself as a living personal knowledge.


In order to do the Father's will on earth, the Son must have a power like
the Father's in heaven. The revelation of the will in the sign of power is
the miracle: the Son could not reveal the Father's will in deed if he
could not perform the miracle. Whence, then, has he the power? Not from
himself immediately, as an actual human being; for an actual human being
cannot, in consequence of the limitation of human nature, be omnipotent. He
has it from the Father, through the eternal Logos; he has it in his
relation to God, inasmuch as it is given to him (John 11:41). And yet he
has
% ---- printed p.314 (PDF 327) ----
\opage{314} it in himself as his own power, a Logos-power appropriated in the
human self, where it is needed: he helps because he wills to help, he
heals because he wills to heal. From the momentary revelations of power,
therefore, no scientific conclusion can be drawn to the denial of his
human suffering, any more than from his human suffering any scientific
conclusion can be drawn to the denial of his divine power. He who cannot
understand in faith that the Son's power is not an omnipotence inherent in
his human nature, but a living power of will, a power which he neither can
nor will have except as the Father gives it to him, a power of religion and
of prayer, a power in the relation to God, must go to Gethsemane and hear
the one who prays: ``Father, if it be possible''!

Just as little as the miracles that Christ performs can shed any light on
his actual human nature, just so little can the miracles that happen to
him himself give any light that falls outside faith. He walks upon the
sea, he is transfigured on the mountain; but he who would conclude from
the walking upon the sea that his human body was, by its nature, raised
above the laws of gravity, would go astray just as surely as he who would
infer from his transfiguration on the mountain the constitution of the
body with which he was crucified. The momentary transformations,
foreshadowings of the last great, concluding transformation in which he
enters into glory, are as incomprehensible for faith as for knowledge: the
concepts belonging here dissolve into the Mystery, the Mystery in which
the unity of nature and creation is hidden.

Instead of holding to the evangelical exemplar and conceiving the relation
between the Christ-Self and the Logos-Self indwelling in the Father as the
revelation of the ideal relation to God, the living religion, ``the life,
the eternal, which was with the Father and was revealed to us,'' orthodoxy
early struck into another, an objectively scientific
% ---- printed p.315 (PDF 328) ----
\opage{315} track, and, without heeding the categorical boundary-line of the
Mystery, has set out by exact determinations of character to make plain
the union of Christ's two natures as a union of two existing substances.
What disputes and confusions, what a swarm of contradictions, what
monstrosities of the understanding this doctrine of union has brought into
the world in heterodox pains with orthodox birth-pangs, is well known. The
orthodox theory of \textit{unitio personalis} and \textit{status unionis},
with the consequences flowing from it, cannot be improved by paraphrases or
mediating alterations; by its sharply stamped form, as the ripe fruit of
scholasticism, it has precisely fallen due to critical dissolution.

Concerning Jesus' words on the cross: ``My God, my God, why hast thou
forsaken me''? Strauss remarks that these words befit only a human being
who is not personally united with God, so that he whose personality is
constituted precisely by this union could not speak thus until his
personality had become another; and the critique follows: ``Now in the
doctrine of the Church the person of Christ is formed by the Son of God,
while the human nature in him, as impersonal, is only taken up into that
divine person. That outcry, then, could come only from the Son of God or
the Logos, since there is no person in Christ besides him, and he must
have called upon the Father, by whom he had therefore felt himself
forsaken. But how will one be able to represent this to oneself without
tearing the Trinity asunder and abolishing the Son's divinity? Therefore it
is not the Son of God or Christ's divine nature that is supposed to have
felt itself forsaken by the Father, but Christ's human nature by the
divine, and at the same time (since \textit{derelictio} --- as
\textit{opus ad extra} --- is supposed to be common to the three Persons)
by the whole Trinity. But the human nature, which was in itself
impersonal, could not have any self-feeling of its own either, and
consequently could not cry out
% ---- printed p.316 (PDF 329) ----
\opage{316} those words either; rather the Son of God in Christ must have done
so, and he must therefore, in the name of his human nature, have
complained to himself of what he himself had done, or, since his human
nature was not a subject distinct from him, have forsaken himself; and so
he would more properly have cried out: why have I, rather than: why hast
thou forsaken me''?\footnote{Troesl.\ II Bd.\ S.\ 119--20.}

Such positive and negative unnaturalnesses fall away of themselves when
the relation between the Father and the Son, with the corresponding
doubling of life in the Logos-Self, is conceived as a fundamental
religious relation; we can then, from the standpoint of religion --- not
from that of science and speculation --- conceive ``the divine-human
states'' as exemplary designations for the alternation of states which is
lived through by believers in outward as well as inward experiences. ``For
not only in the outward states of the saints, but also in their inward
states and moods of soul, there must be an alternation of
\textit{κρύψις} and \textit{Φανέρωσις}, of obscurity and transfiguration,
emptiness and fullness. \dots\ As surely as Christ had to be tempted in
all things as we are, yet without sin, so surely must there also be in his
inward life this alternation of states of obscurity and transfiguration,
of emptiness and fullness. Although he was never alone, but the Father was
always with him, although the fullness of the Godhead dwelt bodily in him,
he nevertheless knew not only such states in which the fullness broke
through the earthly limitation and the Only-begotten rested blessedly at
the Father's bosom, but also such states in which the fullness of the
Godhead had withdrawn into the hidden depths of his being, was only in the
\emph{ground} of his existence, while in his actuality of soul he had to
fight the fight of patience''\footnote{H.\ Martensen. ``Den christl.\
Dogm''. S.\ 298--99.}.
% ---- printed p.317 (PDF 330) ----

\opage{317} \greekrun{γ}{The Self of the Son Is the Light of the World.}

``As long as I am in the world, I am the light of the world'' (John 9:5).
Christ is the light of the world because he is the life of the world: as
the life is, so also is the light. Now as surely as the Christ-life, the
holy Logos-life, is a life in God, the eternal life of religion in the
ideal relation to God, so surely is the light that goes forth from Christ
to illuminate the world also the holy light of religion, which must not be
confused with that of science. Where, however, a great and comprehensive
religious contemplation of the world is concerned, the confusion lies near
at hand. At first glance it looks as though the evangelists and apostles,
in their various utterances about the Son as Logos, as the Firstborn of
all creation, as the Mediator between the Father and the world, had really
given us the foundation for an objective scientific cosmology resting on
the doctrine of Christ. What can be adduced from this point of view has
already been set forth in ``the Christian Dogmatics.'' The Son of God, it
says, as the Logos of the Father, who is at the same time the eternal
world-Logos, has in his pre-existence already prepared for himself the
conditions for his revelation of love, his becoming man in the fullness of
time, the revelation in which he first completely actualizes his concept
as the \emph{mediating} God. His revelation, which has not merely
historical, not merely religious and ethical, but metaphysical
significance, cannot be determined by sin alone, for sin has not entered
according to a metaphysical necessity; it is cosmologically determined: he
could become the Savior only because it is his eternal concept to be the
Mediator \textit{(Etsiamsi homo non peccasset, Deus tamen incarnatus
esset)}. Only as the world-redeeming and world-completing Mediator can the
Son in truth be cognized as the self-revelation of the divine Logos. For
the world-completing
% ---- printed p.318 (PDF 331) ----
\opage{318} principle cannot be different from the world-creating one, through
which all things have come to be. Thus, then, Christ, the Logos become
man, is not merely the center of the human world but the center of all
creation. For him who stands as personal eternity in the midst of the
times, indeed in the midst of all creation, the sensuous difference
between past and future has only vanishing significance; for all the ages
of the world, all aeons, move about him as about their all-determining
center.

The profound utterances in which the author of the Christian
Dogmatics\footnote{Cf.\ ``Læren om Sønnen''. S.\ 243--44; 265--68;
273--74.} has presented the Only-begotten as the life of the world and the
light of the world we shall willingly grant our recognition, now that the
dark sides of the matter have already been sufficiently brought out in
another connection\footnote{Den gode Villie som Magt i Videnskaben.
S.\ 76--110.}; but we ask: is it not after all necessary to secure the
holy truths by securing their grounding? And, when grounding is in
question, is it not then necessary to reflect on the difference between
rational and antirational grounding, on the difference between the light
of science and that of religion, on the difference between objectifying in
knowledge and objectifying in faith? To a separation of this nature the
word of Scripture cannot directly guide us. The biblical worldview
corresponds, also with regard to outward actuality, immediately to the
believing conception of Christ as the center of all creation; nor, as far
as the light is concerned, can there be any question here of a divorce
between religion and science, for to the apostolic consciousness science
did not exist at all. But in the present time, when science from so many
sides and in so many ways makes it felt
% ---- printed p.319 (PDF 332) ----
\opage{319} that it exists; in the present time, when actuality in outward nature
has changed because the consciousness of it has changed, and the whole
universe has received another shape; in this present time it is
unavoidably necessary to distinguish between world and world, between life
and life, between light and light.

In order to see the necessity of a thoroughgoing separation between
religion and science, we need not make science partial to unbelief; we may
well let it proceed from the presupposition that everything revelation
teaches about ``the mediating God'' as ``the center of all creation'' has
its validity, its full truth, in the universe of the spirit and hence, to
that extent, also in the universe of actuality: and yet, when the question
here is of light and illumination, of ground and grounding, science has
its misgivings and will and must make its objections. Christ says: ``All
power is given unto me in heaven and on earth.'' Now when these words are
explained to mean that ``all heavenly and earthly powers, all forces in
nature and history, find in him their liberating center and are
subservient to the kingdom whose head he is,'' then the question is: how
is this to be understood scientifically? If Jesus of Nazareth, the son of
Mary, is the only Christ in the whole universe, then the human race on
this earth must presumably also be the universe's only human race. Or, if
there are on other globes also races that stand in need of redemption,
then there are indeed several possibilities. Either the redemption of all
races on all globes must be included in the redemption which has taken
place once for all on earth; or Jesus of Nazareth must wander from globe
to globe and on each globe let himself be born anew of that globe's Mary,
become a child again and grow up in the form of a servant and begin the
work of redemption over again; or else each globe must have its own
Christ, so that there will be as many Christs as there are fallen races.
But if there are several Christs, indeed perhaps
% ---- printed p.320 (PDF 333) ----
\opage{320} an incalculable number, how then can we understand that our Christ
must be precisely the highest among them all, the only true
over-Christ, ``the center of all creation, the second Adam, in whom the
heavenly and the earthly, the invisible and the visible, the forces of all
creation, angels, dominions and powers, find their all-comprehending
conclusion''? If we say: there can be only one Christ in the universe,
since in all rational races there is essentially only one race, in all
falls essentially one sin, in the work of redemption valid for the whole
world essentially one redeeming activity, then we hold exclusively to the
principle, and can to that extent appeal with reason to the world-creating
and world-redeeming Logos-principle. But the principle is one thing, the
self another.

Between the Logos-Self in God and the Logos-self in the true, sinless
human being the difference must surely become infinitely great when we
take it cosmologically; for an orthodox Protestant who would say that it
is Jesus of Nazareth who created the world would surely express himself
just as unscientifically as a simple Catholic who would hold that it is
the Virgin Mary who created the world.

Once one has placed oneself at the standpoint of science and drawn the
outlines of a christological cosmology, one ought not to shrink back from
difficulties lying near at hand or excuse oneself from further explanation
of obscure points. If Christ is the light of the world, and the light
scientific, then the religious cosmology must be capable of being placed in
objectively clear illumination. That this has not happened, and by the
conditions of cognition cannot happen, is an eloquent proof that the
religious cosmology neither is nor will be scientific.

From the standpoint of faith the objections of science fall away as of
themselves. According to the principle of faith,
% ---- printed p.321 (PDF 334) ----
\opage{321} no objective proof is to be conducted that Christ, by virtue of the
concept ``the mediating God,'' must absolutely be the only possible Savior
of the world and in consequence of this the true Savior for us; were such
a proof to be conducted, then the grounding would certainly have to be
rational, and lack of rational grounding would mean lack of light. But the
grounding is inverted. Here one begins with the inner self-assurance, the
assurance that Jesus of Nazareth, the true, sinless human being, is the
only true Savior for us human beings, that, as Paul expresses it, ``there
is no salvation in any other,'' and then one sees from this that he,
objectified in faith, must also be the Savior true in himself, the Son of
God, the Only-begotten of the Father, whose redeeming, saving power is
cognized not only on earth but also in heaven.

\runhead{c) The Personality of the Son.}
\parmark{21}

The corresponding fundamental features are already indicated in the
exposition of the Father's personality. But just as the Self in the Father
first receives its full illumination through its relation to the Self in
the Son, so too the Son's human existence, with the revelation of
personality corresponding to it, will first be able to furnish the real,
life-filled foundation for the development of the concept of personality
which, because it corresponds to the Son only by corresponding to the
Father, and conversely, expresses the personal unity-in-opposition of the
Father and the Son. We bring this development of the concept under the
three points of view: \emph{the metaphysics of personality, the revelation
of personality, and the eternal life of personality}.
% ---- printed p.322 (PDF 335) ----

\opage{322} \greekrun{α}{The Metaphysics of Personality.}

The improbable is to be believed; but the impossible cannot be believed:
to secure ourselves against the absurdity of impossibilities we take our
point of departure in knowledge. ``Two,'' says Hegel, ``cannot be One,''
and in this he is certainly right. His explanation is, in brief, this. The
category \emph{One}, the abstract One, is, as logic shows, a bad category.
Now as far as personality is concerned, every person is indeed something
rigid, brittle, independent, being-for-itself; but it nevertheless lies in
the character of the person, still more in that of the subject, to sublate
its isolation, its \emph{separateness}. Ethical life, love, is: to give up
one's particularity, one's particular personality, and to extend it to
universality; so that the identity of the one with its other shows itself.
In acting rightly toward another, I treat him as identical with me. In
friendship, in love, I give up my abstract personality and thereby win it
as concrete. The truth in personality is precisely this: to win it through
this sinking in, this immersion in the other. When personality is as
\emph{undissolved}, then one has \emph{evil}; for the personality that
will not merge in the divine Idea is evil. In the divine unity personality
is posited as dissolved; only in appearance is the negativity of
personality different from that by which it is sublated. The Trinity is
brought under the relation of \emph{Father, Son} and \emph{Spirit}; this
is a childlike relation, a childlike, natural form. The understanding has
no such category, no such relation, as could fittingly be compared with
this; it must then be brought to consciousness that it is only figurative:
the Spirit does not clearly enter into the relation. Love would have been
more fitting; but the true is the Spirit. The abstract God, the Father, is
the universal, the eternal, all-embracing, total particularity.
% ---- printed p.323 (PDF 336) ----
\opage{323} We are at the stage of the spirit; \emph{the universal here includes
everything in itself}; the other, the Son, is the infinite particularity,
the appearance (the phenomenon); the third, the Spirit, is individuality as
such; all three are the Spirit. In the third, we say, is the Spirit of
God, but this is also presupposing; the third is also the
first\footnote{Cf.\ Philosophie der Religion. Berlin 1840. S.\ 239--40.}.

That such a foundation is far too thin for a doctrinal edifice to be
erected upon it to the honor of the personal God of the Trinity, anyone
will easily see. The divine Self, the real kernel, is entirely dissolved,
and the whole reduced to a trilogical unity of a formal process in the
ideality of the spirit. The slipperiness of the Hegelian doctrine of the
Trinity Strauss has demonstrated with sharp acuity of understanding.
Concerning Hegel's rejection of the application of number: ``Three cannot
be One; of three as a number there is no question here at all,'' he
remarks: ``Against what, then, is one not to dare to bring number to bear?
against the Hegelian or against the ecclesiastical Trinity? To do so
against the former has not yet occurred to anyone, because moments which
dialectically pass over into one another are not stable enough to permit
any counting; but in the doctrine of the Church these moments, by Hegel's
own declaration, stand fast as subjects existing for themselves, and
consequently they not only can, but must, be counted. Against the doctrine
of the Church, then, what Hegel calls bringing in the categories of
finitude has its full right, as he himself has acknowledged in other
places (e.g.\ Phænomenol.\ S.\ 577 f.). Why, then, censure it here and
thereby prepare for the orthodox a delusive joy, for the rationalists an
undeserved vexation?''

The whole metaphysical-critical-dialectical dispute for and against God's
personality, whether personality is otherwise conceived trinitarianly or
antitrinitarianly, is a mock
% ---- printed p.324 (PDF 337) ----
\opage{324} fencing, for the good reason that the concept of personality in the
divine sense, i.e.\ applied to God, is not a concept of knowledge at all,
but a concept of faith. On the other hand --- and this always has its
significance when the question is of possibility --- a metaphysical
concept of subjectivity or selfhood can be developed which shows that
concrete selfhood essentially consists in the totalized syllogistic unity
of a threefold self\footnote{Cf.\ ``Grundid.\ Logik''. II.\
S.\ 344--84.}.

\greekrun{β}{The Revelation of Personality.}

There have, as is well known, been vehement disputes over whether the
Trinity was to be found in the Old Testament or not. An evident
misunderstanding on both sides, both on the part of those who deny it and
of those who assert it, lies at the root of these disputes. That Jehovah
has revealed himself as an almighty, omniscient, holy Self, as the unity
of will of knowledge and power, cannot be doubted from the standpoint of
faith. On the other hand it can in equal degree be disputed whether
Jehovah in the Old Testament revelation is in the proper sense of the word
a \emph{personal} God, and whether he is a \emph{triune} God. For in the
concept of God personality and trinity belong together so inwardly that
the one of these categories cannot possibly be developed without the
other.

It has already been shown in the doctrine of the Father's personality that
the fatherly Self originally presupposes a self-doubling. In formal
respect this self-doubling is so essential, belongs so ineluctably not
only to the concept but even to the bare notion of a self, be it otherwise
human or divine, that not a word can be said about the self, indeed that
the self itself cannot say a word, without the self-doubling being
noticeable. ``And God said: Let there be light!'' Here, in this
% ---- printed p.325 (PDF 338) ----
\opage{325} sentence, there is not merely a difference between the speaker and
the word, but there is at the same time a difference in the self, a
logical self-separation. In distinguishing himself from the word he
speaks, the speaker becomes conscious of his speaking. But to be conscious
of one's speaking is impossible without being conscious of oneself as
speaking. Whom, then, is it that the speaking God addresses when he says:
Let there be light! He speaks of the light, but he does not address the
light; he speaks in himself to himself; but to speak to oneself is, in
speaking, to double oneself, to be at once both subject and object for
oneself. ``And God said: Let us make man in our image, after our
likeness''; in this form of speech the self-separation is clearly
expressed. The orthodox Church teachers certainly missed the point of
view when, relying on the plural form, they thought to see in it a
revelation of the Trinity; it is of course not the plural form
\textit{(Pluralis majestatis)} but the self-separation in the speaker's
inner being on which the emphasis falls. Now since the self-separation
both proceeds from and returns to the unity of the self, and consequently
rests on the indissoluble unity of doubling, there is undeniably shown
herein a logical adumbration of the essence of the Trinity, and the
spokesmen of heterodoxy are wrong in wanting to make the mode of
expression an empty grammatical form, since it is incontestable that the
grammatical form has its proper ground in the logical. But that there can
be no question here of trinity in a concrete and personal sense is evident
to such a degree that the demonstrated logical adumbration of the unity of
doubling, when more closely considered, is not even sufficient to justify
the doctrine of Jehovah's personality.

Personality presupposes not merely a self, but also a totality
of original fundamental determinations, a richness of content, an articulated
inwardness in the self. To inwardness and to a self-absorption rich in content
belongs a system
% ---- printed p.326 (PDF 339) ----
\opage{326} of inward-going redoublings, a world of powers and thoughts, a
world of Ideas, whose reality corresponds exactly to the reality of the
self-unity's own fundamental redoubling; but of such an inwardness,
such a self-absorption rich in content and unfathomable, there is no trace
in the revelation of Jehovah. Inwardness, concrete inwardness rich in Ideas, is
on the whole not the strong side of Judaism. If Judaism is set against
paganism, then it has inwardness, elements and presuppositions of
true personality; but if it is set against Christianity, then it lacks
inwardness and true personality. This is easy to understand
when one considers that the stage of Judaism is the stage of the Spirit, but in the
form of the Law. Such stirrings of sin and repentance as give, e.g., a
Davidic self inwardness and a personal unity of redoubling can
throw light on the essence of Jehovah only insofar as he is the holy,
law-giving Self who is angry with the sinner, punishes the sinner, but
has mercy on the penitent; these stirrings, if we may so call them,
on Jehovah's side are not inward-going but outward-going,
since they concern not the Almighty himself but man: from this
it is seen that at the stage of the revelation of the Law, the Old Testament,
there can be no conditions for religious insight into the
``depths'' of the Godhead, or clear features of revelation \textit{ad intra}.
Jehovah is spirit, invisible, holy, almighty will; but the spirit, with all
its grace, mercy and loving-kindness, is cold, for here inwardness
is lacking; the high, almighty Self reveals its will in the Law,
but the will is hard and severe, it contracts into the self's
arbitrary commandments and keeps the grounds of personality hidden.

In confirmation of a revelation of the Trinity, and especially of an
unambiguous designation of the personality of the Son, appeal has furthermore
been made to theophanies in which Jehovah and the angel of Jehovah
(e.g.\ Exod.\ 13:21 and 14:19)
% ---- printed p.327 (PDF 340) ----
\opage{327} show themselves to be both the same and yet not the same: the angel is
then the Son, Jehovah the Father. The truth in this is that every
revelation of God presupposes an inner distinction between the one
revealing himself, God in essence, and the form of revelation, God in his
phenomenon; but how fleeting and formal the distinction is shines forth
precisely from the way in which the two designations are interchanged.

In the New Testament God has inwardness, and consequently also
personality; the Idea of this inwardness, this personality, is owed
neither to Plato nor to Philo, but solely to the Word revealed in Christ.
That the author of the fourth Gospel makes use of the
pregnant category Logos in order to express the fundamental relation between God
in himself and God in his revelation by no means entitles us to the
conclusion that John has taken up either Platonic or Philonic
subsidiary notions into his Gospel. On the contrary! it goes with the
category \textit{λογος} as with the categories \textit{ζωη} and
\textit{Φως}; they are chosen because they occur in the language; but their
meaning does not conform to an already given usage:
the Spirit of truth breathes a new life into them, so that they undergo a
rebirth.
Thus we come once again to the problem of the
double Logos-self, and thereby to the question of the Son's
personality.

To make the solution of the problem possible we must remind ourselves beforehand of
what all know and yet so easily forget: that the spirit gives life,
but the letter kills. It is the principle of revelation, not the
literal holding fast to the Evangelist's words, to which we have to
hold. Were we to hold literally, punctiliously and grammatically to
the evangelical text, we should have to deny the spirit and deify
\textbf{Absurdum}.
For what does the text say literally? It
says that \emph{the Word was with God}, and that \emph{the Word was God}; it
says, therefore, that God
% ---- printed p.328 (PDF 341) ----
\opage{328} in the beginning was with God: it says, therefore, when we put the
letter in the place of the spirit, that in the beginning there were two Gods. The
literal twoness comes clearly forward in the sentence: \emph{the Word
became flesh}; for it is not the God with whom the Word was, but the
Word which was with God, and so the second of the two Gods, that became
flesh. If finally we make precise what is immediately literal in the
sentence: the Word (the second God) became flesh, what then does it say? It
does not say that the Word, God, assumed the form of flesh --- such a
paraphrase departs far too much from the literal and spiritless ---
no, it says that the Word, which was God, became flesh, so that flesh
stepped into the place of God: here the literal goes together with the
absurd, and the absurd together with the blasphemous. For the sake of what
follows it is necessary to have brought this out sufficiently.

If we interpret the Evangelist's words in the spirit of the Gospel, then we certainly
need do no violence to the text in order to reach the fundamental thought: for
the fundamental thought is explained by revelation itself. In the beginning was
the Word, the Word of the Creator, the Word by which all things were made: in and
through this his creating Word God is eternally revealed to himself. The Creator
is God, God in himself; the creating Word is God, God in his inner
revelation: this is God with God, God's unity with himself in
original self-redoubling. The self-unity is personal, for the
opposition is spiritually real.

The spiritually real opposition in the Creator's Self is, however,
different, qualitatively different, from the externally real
opposition between \emph{Creator and creature}, between \emph{God} and
\emph{world}. Over against the outer opposition the inner one vanishes; it
withdraws from contemplation, since with all its reality it is spiritual.
Opposed to the manifoldly conditioned unity of connection of the actual world,
God in himself is infinitely concentrated
self-unity. At this stage of revelation it comes easily and
naturally to the religious
% ---- printed p.329 (PDF 342) ----
\opage{329} view to hold fast to faith in \emph{the only true God} as
Creator of the world; while all talk of God's trinitarian essence only
looks like offensive speculation.

But the revelation of personality has not yet been completely
set forth; the Evangelist adds: ``And the Word became flesh and dwelt
among us,'' and by this he points to the new revelation, namely that
by which the personal God, and so God's personality, first of all
becomes revealed to men. That which reveals is again the Word,
the Word which was in the beginning. The Word itself has not altered its
essence, but its revelation has become another. For here we must
distinguish between the revealing Word and the means by which it
becomes revealing. From the standpoint of creation it is the world that is
the means; from the standpoint of the incarnation it is the God-man. Religiously
understood it is then clear enough that the world, the natural, actual
world, stands in a quite different relation to the Word than the God-man does.
The world is, according to what we have learned from creation, a nothing for
the Word of omnipotence and yet an actual, existing something: this
unity-in-difference is impenetrable to thought, opaque to
view: the means is, in other words, externally finite,
inadequate. The God-man on the other hand, the true, sinless man, is in
his temporal difference from the Word essentially one with the eternal Word;
the unity-in-difference of the means of revelation and revelation itself is so
transparent, so intuitable, that the Word's means is the Word's form
(John 1:1--2). The revealing Word in the God-man, the
actual Logos-man, itself becomes revealed. In creative unity with
the means the Word is life; in the transparency of the means for the Word, the Word is
light. The Word therefore does not give, with the world, the great whole, as means,
a more perfect revelation of God, and in the God-man with his individual limitation
a less perfect one, but conversely: the revelation which the Word
began in crea\-%
% ---- printed p.330 (PDF 343) ----
\opage{330}tion it completes in the incarnation. If we now remember that the talk is of
the Word from the beginning, that it is the same Word which in
becoming man completes what it began in creation, and that
the God-man, seen in the light of revelation, is the Word, then we understand
also the Evangelist's testimony: ``He was in the world, and the world was
made by him, and the world knew him not. He came unto his own, and
his own received him not.''

In making Christ, in expressions full of life, one with the
Creator-Word from the beginning, the Evangelist evidently places the unity in that which
reveals, \emph{the Word}, not immediately in the existing self.
It is the confusion of the unity of revelation with immediate
self-unity that has occasioned the unhappy speculation that it
is the second person in the Godhead, despite all the protests of orthodoxy
orthodoxly understood, the second God, who, in accordance with express agreement,
consultation, with the first God, has descended from heaven, has
humbled himself and assumed the form of flesh. Against this
absurdum the religious view sets an express opposition
between \emph{God and Christ}. Personality in Christ has for
this view a quite different stamp from that of the Creator of the world, the
almighty, omnipresent God's personality. God's personality
rests in its infinite self-concentration on a universal basis, Christ's
personality on the other hand on an individual basis: the opposition
between personality in the universal form and personality in the
individual form is of decisive significance for the faith of revelation.

So then it becomes a task of the philosophy of religion to state
how the personal relation between God and Christ is to be thought
more precisely. It is first and last a relation of revelation: the way of
cognition in this determinacy goes not from God to Christ, but from
Christ to God. Only in faith in Christ's personality is it
% ---- printed p.331 (PDF 344) ----
\opage{331} possible to cognize God's personality. The revelation of the
personal opposition is grounded in the personal unity;
the unity of revelation is the Logos; the Logos-self therefore belongs at once
both to God's and to Christ's personal self-revelation. In the Logos
and by means of the Logos God is from eternity personally revealed to himself
as God; in the Logos and by means of the Logos Christ is personally
revealed to himself as God-man; in the Logos and by means of the Logos
the personal opposition between God and Christ is expressly
revealed as a personal unity; only by the Word become flesh
being set in personal unity with the world-creating Word has God in his Word
become reconciled with the world: ``God was in Christ, reconciling the world
unto himself'' (2 Cor.\ 5:19). The Logos-self in God and the Logos-self in
Christ is a self redoubled in itself:
% print: first "Logos selvet" as two words, second "Logosselvet" as one
in this redoubling
revelation is grounded; but the redoubling presupposes the
principial unity: on the personal self-unity of the redoubling
revelation is grounded. Thus the relation between God and
Christ enters, for the view, into a new stage and presents itself as
a personal relation between \emph{the Father and the Son}. By means of
the Logos-self in God, God is the Father of Jesus Christ; by means of the Logos-self
in Christ, Jesus Christ is the only-begotten Son of God.

\greekrun{γ}{The Eternal Life of Personality.}

``This is life eternal, that they might know thee, the only true God, and
him whom thou hast sent, Jesus Christ'' (John 17:3). For the sharp
separation of the eternal and the temporal, propositions are sometimes set up
such as these: the eternal \emph{is}, but does not \emph{become}; the temporal
\emph{becomes}, but \emph{is} not. But however much significance such a
separation may have for abstracted thinking, it is nevertheless easily seen
that it cannot hold for life; for in this way abstraction
takes life away from both the
% ---- printed p.332 (PDF 345) ----
\opage{332} eternal and the temporal. Life is just as little a being without becoming
as, conversely, a becoming without being. Life is, in its peculiar being-for-itself,
essentially the unity of being and becoming. If now being designates the eternal, and
becoming the temporal, then it can rightly be said that life as life
is the unity of the eternal and the temporal.

Even natural life has a stamp of eternity, insofar as it
has a stamp of selfhood. Its development is self-development, and in the
becoming of self-development it preserves, under all formations and
transformations, a typically peculiar being: the temporal is that it
ceaselessly changes; the eternal is that in the stream of changes
it preserves the fundamental stamp of its essence unchanged. The smallest blade of grass
at the foot of the mountain has, as a living organism, more eternity in it than
all ``the everlasting mountains'' together; but it also has more eternity in
it than all the streaming rivers together. That life is the unity of
being and becoming means that in all its immediacy it is
reflected, that its organic self-activity is an ideal-real
circular motion, a motion of form bent back into persistence. To this
corresponds the opposition between the inner and the outer. The plant is the outer,
the life in the plant the inner; however much the inner goes together with
its outer, the living opposition of inner and outer is nevertheless clearly
expressed in the fact that the outer, in nourishment, growth and
formation, is constantly determined by the inner. Life moves from itself
through its other back into itself, from its inner through its
outer back into its inner: it is an image of eternity. The
determining inner lies in the germ as a peculiar disposition; life
consists in the actual unfolding of the disposition, but the unfolding is
as a self-revelation: it is the essence of life that the eternal wills
to reveal itself in the temporal.

Now just as temporal life already has something eternal in it, so
eternal life too originally has something
% ---- printed p.333 (PDF 346) ----
\opage{333} temporal in it. It belongs to the life of the spirit as much as to the life of nature
to move through oppositions. The richer and deeper the life of the spirit
is, the freer is its inwardness; but the freer the inwardness is,
the greater the significance it is able to give to its corresponding
outer. If the free inner were indifferent to its outer, if
eternal life had nothing whatever to do with temporality, then the
relation between God and the world would make no sense. It is indeed
not God who is dependent on the world, but the world that is absolutely
dependent on God; nonetheless God, by creating the world, has
brought forth temporality, just as he also, by preserving the world,
bears temporality and, by governing the world, works in temporality. But
that the Eternal himself brings forth the temporal, bears temporality and
works in the temporal, is surely an eloquent testimony that
eternal life must have something to do with the temporal.

It is, says religion, not out of necessity but out of goodness and
love, and so with full freedom, that God has created the world. What
does this mean other than that freedom, goodness and love, which
belong to God's eternal life, entail precisely that eternal life wills to
move through the temporal? Now since God's freedom, goodness and
love are one with his personality, we can understand from this
that it must belong to God's personal life to impart
life to the world. It is a fact that temporal life arises in time on
natural conditions; but it is none the less
untrue that any life whatsoever could be derived from
any temporal composition whatsoever; the least germ of life is of
eternal origin: all life in the world is from God. Nor can any
temporal life be annihilated by anything temporal whatsoever; it can
be developed in temporality, on temporal terms, but --- it follows
incontrovertibly from the
% ---- printed p.334 (PDF 347) ----
\opage{334} law of selfhood\footnote{Cf.\ Forelæsninger over phil.\ Propædeut.\
1860--61.% print: "Proprædeut."
} --- when its self-development is ended, it returns, enriched
with the inner result of its development, to its source, to
eternity, to God, whence it came.

And yet there is a difference, a qualitative difference, between a temporal and an
eternal life, as surely as there is a difference, an infinite qualitative difference,
between God's life and the world's life, a qualitative difference between life in
God and life in the world. The eternity of life rests on its essential
selfhood. The lower selfhood stands, the more corruptible is
life; the higher selfhood stands, the more it approaches the
incorruptible. In personal life selfhood is a self conscious of itself,
determining itself: the more personal life is, the more clearly
eternity shines through.

To become conscious of one's personal life and to become conscious of one's eternal
life are at bottom one and the same. But to become in truth conscious of one's
personal life is possible only in faith; for it is
possible only by becoming conscious of the personality of the revealed God.
To know eternal life is to know the only true God;
but life can only be known by life; to know God's life is to live in God,
from God; but to live in God is to live from God's revelation.
Revelation is life; God himself lives only in his revelation. Therefore
it belongs not merely to a higher development of life, a clearer
conception of the eternity of life, but it is life itself, eternal life,
to know the revealed God, to know him personally in him whom he
sent, Jesus Christ.

So long, however, as it is still only a matter of personality, life and
revelation in general, the consideration can easily move
through crossings of thoughts of faith and of knowledge with so hovering a gait
that it almost looks as if the opposition of faith and knowledge
% ---- printed p.335 (PDF 348) ----
\opage{335} were here soared over. This holds, however, only so long as one speaks
generally about the general. As soon as the problem, the problem to
whose solution the introductory remarks point, is drawn forth and
held fast sharply, the crisis too sets in. For the problem is this:
\emph{how is the eternal life of personality in the Son's relation to
the Father essentially to be thought?} When this question is set in
motion, faith and knowledge arm themselves anew for battle. Even if in
knowledge one begins ever so piously with the presuppositions of faith, the
method corresponding to the grounding must nevertheless necessarily bring with it
consequences that deny faith. We see it already in Jacob Böhme.
The struggle of the powers of life in the divine essence, that
wondrous tension and stir of the elements of spirit through which the eternal
life of the all-personality must work its way up from darkness to light, so that ``God's hidden will,
which in itself is only one, may become revealed,'' was indeed precisely the object
of the profound theosophist's wildly fermenting speculations. From him
stems the Idea of a becoming God; only a becoming God is a living
God. Eternal being is in this life of God so absorbed into temporal
becoming that the eternal itself becomes principle and possibility, but never
full actuality: the becoming God is never finished. On the
Böhmean foundation Schelling worked further and attempted to
develop God's personal unity out of an originally presupposed twoness. His
words against Jacobi are well known: ``So long as no
actual twoness is cognized in God, and so long as no limiting, negating
power is opposed to the affirming one, so long will the denial of a personal
God be scientific honesty.'' But even from the dawning light
of Schelling's ``primal ground'' the Godhead's way to perfect
self-revelation is infinitely long: the becoming God is never
finished.
% ---- printed p.336 (PDF 349) ----

\opage{336} Now if God as God, because of the infinite becoming of his life, can never
be finished with the formation of the person of the Father, then the one who knows
in faith has still one way out left, namely: to give up the
fatherly personality, personality on the universal basis,
to let God as God sink down into the indeterminable unity of the substantially
absolute, in order then to hold fast to personality on the
individual basis, the true personality in the person of Christ. As
representative of this peculiar way of seeing we name
Schleiermacher. Schleiermacher's ``Son of God'' is a son without a father, a
fatherless son, so to speak, in the divine sense; in the human sense
he has on the other hand, as a kind of compensation for the want of the heavenly one,
received an earthly father. What is contradictory in Schleiermacher is that he
wishes, within science, to free religion from science.
Schleiermacher has very rightly seen that ``the ecclesiastical formulas concerning
the person of Christ\footnote{Der christl.\ Glaub.\ II.\ Bd.\ S.\ 49.}
might need a critical treatment.'' He acutely
draws attention to the impropriety of designating ``the Redeemer's
divine nature from eternity, before its union with the
human, in such a way that this union does not appear at all as a
moment co-constitutive of the person Jesus Christ, but much rather
merely as an act of this person himself.'' How little the
Schleiermacherian Christology, despite the critical fineness and
dialectical sharpness with which it is carried through, has however
been able to withstand ``the daily swelling stream and
storms of critique'' can be seen in Strauss; how little it, with its otherwise
unmistakable religious merits, satisfies the ecclesiastical
view can be seen from utterances such as this: ``Since Schleiermacher
proceeds from the opposition between sin and redemption, he conceives
the Redeemer as the second Adam, the perfected creation of human
% ---- printed p.337 (PDF 350) ----
\opage{337} nature, and cognizes his God-consciousness
as the expression of the actual presence of the Godhead in him, sees him
as exemplary for all times and generations, as the head of a new
communal life that continually receives of his fullness. But the defects of
this Christology may be referred to two main points, to
the divine side and to the worldly side of Christ's existence. As regards
the divine side, the concept of the Son's \emph{pre-existence} is lacking here,
and with it also the concept of a descending revelation of love of
the eternal Logos. It is only the impersonal Godhead that in Christ
attains personality. As regards the worldly side of Christ's
existence, this Christ has only religious and ethical, but not
cosmic significance; he stands indeed as the perfected creation of human nature,
but not as the conclusion of \emph{all creation},
and the biblical
passages which declare that all principalities and powers and forces are
subject to him must be content with a merely moral
interpretation.''\footnote{H.\ Martensen. Den christl.\ Dgmt.\
S.\ 277--78.}


This, in outline, is the result of the objectively scientific movement
of thought. We shall now try, from the standpoint of faith, how far it
is possible, by means of the inverted grounding of the philosophy of
religion, to approach the fundamental view of the Church.

It must then once more be insisted that science knows nothing and can
know nothing of God's personality, since this concept lacks every
scientific foundation. Since here, as regards the cognition of
personality, one begins in faith, one begins here unconditionally with
Christ as the object of faith; thus not immediately with the universal
personality, but with the personal individuality of the God-man. Were
the objective consideration of science really to be carried through
here, then, according to what we have shown in the
% ---- printed p.338 (PDF 351) ----
\opage{338} foregoing, the knower would have to think he discovered two Sons of
God, one in heaven and one on earth. So we should then have not three
but four Persons in the Godhead: 1) the Father, 2) the eternal Son, by
whom the world is created, 3) the Son of Man, born of the Virgin, and
4) the Holy Spirit: a confusing sight! Whence comes the confusion? From
an optical illusion! The knower of faith sees double, much as one
speaks of a drunken man seeing double. The knower of faith looks at one
and the same time into two absolutely heterogeneous mirrors, that of
reason and that of faith, that of metaphysics and that of revelation;
thereby the image of the one Son becomes two images: the image of the
eternal Son, the universal personality, appears in the mirror of
metaphysics, while the image of the Son of Man, the individual
personality in flesh and blood, is beheld in that of revelation. To
remove the optical illusion we must get rid of the metaphysical
mirror, or at least see to having it placed in such a way that no more
than one image can possibly appear.

The image that is seen in the mirror of revelation is Christ's
individuality with the pure, holy features of the God-man. How this
personal individuality, which came to be in time, is related to the
eternal universal personality, exalted above temporal becoming, is
then the question. Is the Son who appeared in time perhaps the eternal
Son's representative or deputy among men, \textit{vicarius filii
æterni}, much as the Pope is \textit{vicarius Christi} for the
Catholics? By no means! For an individual or a society of individuals
an individual can be a deputy; but the eternal Son is no more an
individual than the eternal Father: the eternal Son is one with the
Father in universal personality. But a universal personality has no
deputy and can have none. Instead of being conceived as deputy for the
Eternal, the Son of Man in time must be conceived as the Eternal's
% ---- printed p.339 (PDF 352) ----
\opage{339} Revealer. This gives an entirely different view. The deputy always
stands in an external, finite relation to the one or ones he
represents; but the Revealer is one in being and will, in life and
personality, with the one he reveals. When Christ remembers his
pre-existence, he knows himself to reveal the eternal Son; but in his
revealing of the Son he does not represent the Son, but is the Son, the
Son who reveals the Father.

The eternal Son, the Logos-Self in God, is God's invisible
Self-revealer \textit{ad intra}; but the Son of Man, the Logos-Self
become man in time, is God's visible Self-revealer \textit{ad extra}.
Whence, now, do we know the revelation \textit{ad intra}? Solely and
only from the revelation \textit{ad extra}! How do we know that the
revelation \textit{ad intra} is a revelation of personality? We do not
know it; but we cognize it in faith, thus solely by virtue of the
revelation \textit{ad extra}. Faith, with its antirational aim, sees
the inner revelation only through the outer, sees the eternal Logos
only through the Logos become man, the universal personality only
through the individual. Therefore the eye of faith cannot possibly get
two images, one of the eternal Son, another of the Son revealed in
time. Between God and the Word, between the One who is revealed and the
One who reveals, i.e.\ between the Father and the Son, there forms a
religious, personal opposition, an opposition of revelation for time
and eternity; but between the eternal Son, the Father's Revealer
\textit{ad intra}, and the Son of Man, the Revealer \textit{ad extra},
no religious, personal opposition can possibly form. The inward-going
and the outward-going revelation cannot be set up side by side in such
a way that an actual separation of two revelations appears here; for
the same reason, neither can any personal separation appear between
the Revealer \textit{ad extra} --- the Word that became flesh --- and
the Revealer \textit{ad intra} --- the Word from the beginning --- the
Word \textit{ad intra} and
% ---- printed p.340 (PDF 353) ----
\opage{340} the Word \textit{ad extra} are, in the sense of revelation, the same
Word: the two Revealers are in personal working one and the same.

But what, then, does the doubling signify? It signifies that the
individual personality, which passes through temporal development, has
its origin from the universal personality living in eternal being and
working. The self-doubling signifies that the eternal Revealer has in
the fullness of time doubled himself. The Son of Man comes in addition,
that is true; but he does not come in addition as a finite helper; he
grows forth from the Logos-seed, his life in temporality is a personal
continuation of the Logos-life from eternity; he does not go forth
from the Logos, for his personal Self is indeed precisely the Logos-Self
become man; but he goes forth from the Father. Then he completes his
life in temporality and returns to his source; but he does not return
to the Logos; the visible Son does not return to the invisible; in the
light of revelation they are, after all, both one essence, one life,
one will; but he returns to the Father, for it is the Father, the
Invisible, whom he has revealed to the world. --- ``I came forth
(John 16:28) from the Father and came into the world; again I leave the
world and go to the Father.''

\lettersub{B.}{The Works of the Son.}
\parmark{22}

That the Visible reveals the Invisible, that the Son reveals the
Father, is recognizable in the essence of the Son; but now the Son's
essence doubles itself in such a way that
% ---- printed p.341 (PDF 354) ----
\opage{341} the Visible too is in itself invisible, that the Son can reveal the
Father only by revealing himself in words and works. This doubling
follows directly from the double-sided nature of individual personality
itself. In personal life there is not merely an opposition of inner and
outer; but on account of free inwardness the separation is at the same
time so thoroughgoing that the outer, individuality in its phenomenal
immediacy, can just as well be said to conceal as to speak for and
explain the inner. The richer, clearer and freer the inwardness, the
more significant the opposition. In the God-man, whose human exterior
at one and the same time both conceals and reveals the divine
interior, the opposition is at its extreme: it is this opposition of
essence and phenomenon, of the glory of the Godhead and the form of a
servant, strained to the utmost and thereby straining, that
categorically makes the Son of Man a sign of contradiction.

The point of view for the sign of contradiction is, however, a double
one, since it is to correspond most immediately either to the life of
faith or to the thought of faith, either to religious existence or to
the conceptual development of the philosophy of religion: these two
points of view must not be confused.

From the side of existence and in ``the situation of contemporaneity''
S.\ Kierkegaard has, with edifying emphasis, set the sign of
contradiction in relation to ``the possibility of offense,'' and has
treated in particular: 1) ``the possibility of offense that does not
relate itself to Christ as Christ (the God-Man), but to him as a
simple, individual human being who collides with an established
order''; 2) ``the essential possibility of offense in the direction of
loftiness, that an individual human being speaks or acts as if he were
God, declares himself to be God, thus in the direction of the
determination God in the composite God-Man,'' and 3) ``the essential
possibility of offense in the direction of lowliness, that the one who
gives himself out to be God shows himself to be
% ---- printed p.342 (PDF 355) ----
\opage{342} the lowly, poor, suffering, and finally powerless human
being.''\footnote{Indøvelse i Christendom af Anti-Climacus.
Kbhvn.\ 1855. S.\ 93 flgd.} By thus making clear what is categorically
posited in the sign of contradiction, the author shows that a
scientific conception of what the God-man is, a direct historical
knowledge of Christ, is a contradiction, a sheer impossibility,
--- that, in other words, all immediate historical faith in the sacred
history is superstition --- and in this we must unconditionally grant
that he is right. ``The God-Man is the sign of contradiction, and why?
Because, Scripture answers, he was to reveal the thoughts of hearts.
Does, then, all the modern talk about the speculative unity of God and
man, all that which merely regards Christianity as a doctrine: does it
bear the remotest resemblance to the Christian? No \dots\ The God-Man
is an individual human being --- not a fantastic unity which has never
existed except \textit{sub specie æterni}; and he is anything but a
lecturer who teaches directly for rote-learners or dictates paragraphs
for stenographers; he does just the opposite, he reveals the thoughts
of hearts. \dots\ And only the sign of contradiction can do that: it
draws attention to itself, and then it presents a contradiction. There
is a something which makes it impossible not to look --- and lo, in
looking one sees as in a mirror, one comes to see oneself, or he who
is the sign of contradiction sees straight into one's heart as one
stares into the contradiction.''

But the sign of contradiction is not for that reason in itself a
contradiction; the improbable must be believed, but the impossible
cannot be believed: the philosophy of religion has the task of
developing the principle of faith and freeing the concepts from every
semblance of contradiction. Here, then, an express distinction must be
made between the contradictions that naturally accompany
% ---- printed p.343 (PDF 356) ----
\opage{343} the phenomenon, and the contradictions which, both from the side of
the dialectic of paradox and from that of the critique of knowledge,
are alleged to lie in the essence. Perhaps on this occasion it might
occur to verbal critique to come forward with an objection such as
this: ``It is granted that there are contradictions in the phenomenon;
but since the phenomenon is nothing other than the essence's own
appearing, the contradictions of the phenomenon are precisely
appearances of the contradictions hidden in the essence itself; the
essence is thus in itself just as contradictory as the phenomenon, and
the \emph{absurdum} which one means to avoid by shifting the
contradictions onto the phenomenon is by no means removed.'' To see how
matters stand with this sort of fine critique, we need only take note
of a concept such as, for example, misrecognition. Misrecognition is a
contradiction in the phenomenon; a contradiction in the phenomenon is
at the same time a contradiction in the essence. The contradictory is
an absurdum, something impossible in itself; \textit{ergo:}
misrecognition is impossible. --- There is a metaphysical fineness in
thinking which comes about through the thinker's avoiding touching the
problems, and it recalls that fineness of tone which was conjured up
by the virtuoso's bow not touching the strings at all.

The works of the Son, the divine-human acts in which Christ has
revealed his personal essence to the world, may, from the standpoint
of the philosophy of religion, be comprehended under the three
categories: \emph{the Fulfillment of the Law}, \emph{the Atonement}, and
\emph{the Redemption}.

\runhead{a) The Fulfillment of the Law.}
\parmark{23}

Christ's relation to the Law is double-sided. As the Son of Man, born
of a woman, he is born under the Law and ``made subject to the Law'';
as the Son of God he is
% ---- printed p.344 (PDF 357) ----
\opage{344} exalted above the Law, is Lord not merely of the Sabbath but of the
Law, himself Lawgiver. The double-sidedness shows itself also in this,
that he both confirms and abolishes, both fulfills and sublates the
Law. The one does not stand in conflict with the other: the
contradiction is only apparent. Christ is not under the Law in spite of
being above the Law, but he is above it just because he is under it;
he does not abolish the Law in spite of confirming it, but he
abolishes it by confirming it, he sublates it by fulfilling it.
% print: "jnst" (turned u) for "just"
To see more closely how matters stand in this regard, we then seek an
elucidation of \emph{the Spirit of the Law}, of \emph{the Law of the
Gospel}, of \emph{the End of the Law}.

\greekrun{α}{The Spirit of the Law.}

How Judaism in Christ's time had become an external, spiritless affair,
an established order which, following the direction of the Pharisees
and the scribes, held to a hair-splitting interpretation of the letter
of the Law, to traditions and arbitrary ordinances, and in this
spiritlessness deified itself, is well known. That Christ's appearance
is marked from the beginning by a break with this established order is
likewise well known, and has, in our time especially, been brought
home with convincing force and clarity by S.\ Kierkegaard. That it now
matters not only that there is a break with a spiritless established
order, and that the break is decisive, but that it also matters how
the break is made, in what spirit the break takes place, every
thoughtful person must grant. Among the works in which Christ has
revealed the divine-human essence of his personality belongs, then,
this one too: that he has taught us how, in spirit and in truth, one
is to break with an established order that defiantly claims to have
the laws on its side. ``Think not that I am come to dissolve the Law
or the Prophets: I am not come to dissolve, but to fulfill. For verily
I say unto you: till heaven and earth pass away,
% ---- printed p.345 (PDF 358) ----
\opage{345} not even the least letter or one tittle shall pass from the Law, till
all be fulfilled'' (Matt.\ 5:17--18). It is a quite peculiar
beginning. Christ knows within himself that what is at stake is a
struggle of life and death with an externality that has on its side
not merely actuality, with the venerability of what is inherited, but
also the semblance of upholding the religious-moral order instituted
by God himself, by holding to the forms and enforcing the laws; he
knows that this external, seemingly godly affair is an ungodly deceit,
that against this deceit, as ungodly as it is official, there must
take place a revolution that leaves not one stone upon another of the
old building: for that, precisely for that reason, he begins by
insisting on the validity of the Law. In breaking outwardly with the
established order, he breaks at the same time inwardly with
everything that in any way resembles willfulness, arbitrariness,
lawlessness; and he makes this break so abruptly, so thoroughly, that
no one who wishes to emancipate himself from the commandments of the
Law shall find support in him. ``Whosoever therefore shall break one
of these least commandments, and shall teach men so, he shall be
called the least in the kingdom of heaven: but whosoever shall do and
teach them, shall be called great in the kingdom of heaven'' (Matt.\
5:19). It is thus not in order to sublate, but in order to enforce the
Law, that he breaks with what the established order
understands by the Law.

But what, then, does Christ understand by the Law and the commandments
of the Law? We hear it from his words. It is precisely the remarkable
thing about the way in which he interprets the Law in the Sermon on
the Mount, that he places its essence neither in the disposition of
mind alone nor in the commandment alone, but in that union of
commandment and disposition which shows precisely that obedience must
always be the essential foundation of religious inwardness. ``Ye have
heard that it was said by them of old time: Thou shalt not kill, and
whosoever shall kill shall be in danger of the judgment'' (Matt.\
5:21). Had Christ now placed
% ---- printed p.346 (PDF 359) ----
\opage{346} the essence of the Law in the disposition alone, he would surely have
said: we do not need this commandment; for it is the essence of the
self, the innermost longing of the soul, to love righteousness and in
the spirit of righteousness to love one's neighbor as oneself. But
instead of abolishing the old commandment, he adds new commandments,
commandments so severe that not even a Draco himself, who wrote in
blood, could write them so severe: ``Whosoever shall say to his
brother, Raca, shall be in danger of the council: but whosoever shall
say, Thou fool, shall be in danger of hell fire'' (Matt.\ 5:22). To
take this in all indefiniteness as ideal demands is just as spiritless
as to cling pettily to the words Raca and fool. What is
characteristic is that the ideal demand attaches to definite
commandments, commandments which, in expressly aiming at the
disposition, attack the selfish point by point, keep account of the
individual feelings, thoughts and actions, and precisely thereby, as
true commandments of inwardness, bring about a thoroughgoing working
over of the actual interior. To hold to the disposition alone is, in
all indefiniteness, to hold to the universal, and precisely thereby to
grant to contingencies, to the contingent mood, the contingent notion,
the contingent word, an all too arbitrary latitude. But to lose oneself
in a swarm of commandments of inwardness, an ascetic anatomizing of the
inward particulars in the psychic fibers of every entanglement, would
lead not merely to spiritless mechanism but to endless self-torment,
to scrupulosity and despair. Therefore Christ too, through his
commandments of inwardness, lays emphasis first and last on the
disposition; on mercy: Blessed are the merciful; on peaceableness:
Blessed are the peacemakers; on purity of heart: Blessed are the pure
in heart. The commandments of inwardness, which serve to cleanse the
mind, are sunk into the ground of the mind in order to surface again
when self-examination begins. The Law is the universal essentiality of
inwardness, of the inwardness that is thoroughly purified through
religious obedience to the
% ---- printed p.347 (PDF 360) ----
\opage{347} ideal demands of the commandments; as a reflected, active unity of
commandment and essentiality that moves the will, the Law is spirit.

From the spirit of the Law the judgment upon the established order
must be derived. What matters then first and foremost is in what
character the established order is held; for there is a difference, a
qualitative difference, between a spiritual and a worldly, between a
religious and a civil established order. And as there is a difference
between established order and established order, so there is also a
difference between law and law. The political-civil established order
is not merely something external, but has by its nature also a
relatively original right to be something external, to maintain itself
as something external, and to defend itself against attack by external
means. The laws on which such an external established order grounds
its right are then not pure laws of inwardness either, but on the
contrary laws of fact, laws of compulsion, laws of right, not in the
proper sense laws of righteousness. With an established order of this
nature, a political-civil established order, Christ, properly
speaking, will have nothing to do; he will neither tear it down nor
build it up nor improve it; he is personally independent, absolutely
independent of it, of every finite regard, even the faintest
side-regard to it. As his independence of it is absolute, so too his
dependence on it, his submission and obedience, is in a certain way
unconditional, insofar as he does not set himself up in outward
resistance, offers no physical resistance at all, stands in league with
no outward-going selfish striving whatsoever. No rebel against the
civil established order can use his name; but let us then, for truth's
sake, not forget either that no defender of a political-civil
established order can use his name.

The established order against which Christ turns with his whole
personal authority and the strongest expressions of holy wrath is a
religious, ecclesiastical established order, which,
% ---- printed p.348 (PDF 361) ----
\opage{348} after the ideal demands of the divine Law have long been forgotten,
has transformed the law of righteousness into an external positive law
of right, the original, incorruptible right of religion, of the
disposition, of the inwardness of faith, into a juridical church law.
It is for the Law's sake, in the spirit of the Law, that he breaks with
this established order; it is in the name of righteousness that he
attacks the pharisaic right. ``For'' --- this For refers expressly to
the admonition not to ``break one of the least commandments of the
Law'' and ``teach men so'' --- ``I say unto you: except your
righteousness exceed that of the scribes and Pharisees, ye shall in no
case enter into the kingdom of heaven,'' (Matt.\ 5:20). Were positive
church law to be made to prevail, then Christ would simply be in the
wrong, and the Pharisees, with all juridical formalities, would easily
get him condemned as a transgressor of the law; but the fact is that
Christ does not know and will not know any positive juridical church
law whatsoever, because he does not recognize and will not recognize
any form whatsoever of externality as entitled, by laws of compulsion, to
set a limit to the law of inwardness. Does Christ, then, condemn
positive juridical right? By no means! Take away the religious motives
and put civil ones in their place: and surely no objection will be
derivable from Christ's gospel against a civil ordering of external
ecclesiastical affairs. Christ does not break the external bonds
without giving full compensation by strengthening the internal ones. To
loosen all external bonds without the least guarantee that the internal
ones become the firmer --- whether this happens under one pretext or
another, whether it happens by open rebellion or by crooked juridical
detours --- is to violate the spirit of the Law and give society over
as a prey to unbridled willfulness. When Christ, in the name of the
Law, holds judgment over the arbitrary external, he aims at the same
time sharply at the self-seeking internal. ``Ye have heard that it
hath been said: An eye for an eye, and a tooth for a tooth. But I say
% ---- printed p.349 (PDF 362) ----
\opage{349} unto you, that ye resist not evil: but whosoever shall smite thee on
thy right cheek, turn to him the other also. And if any man will sue
thee at the law, and take away thy coat, let him have thy cloak also''
(Matt.\ 5:39--40).
% print: "Matth. 5, 39--40"; the quotation begins with 5:38 -- probably meant 5:38--40
This, then, is a piece of Christ's church law; for the
establishment of good order it is quite as strict as the strictly
juridical one.


\greekrun{β}{The Law of the Gospel.}

The Lutheran orthodox, and on the whole the old Protestant, doctrine of
the Church has made a sharp separation between law and gospel.
Agricola's antinomian zeal against the preaching of the Law to ``the
regenerate'' provoked a polemic that was in several respects unclear,
and that at last found its clarification in the Formula of Concord. For
the sake of an overview we cite Hase's concise presentation\footnote{\textit{Hutterus
Redivivus.} Leipzig 1836. S.\ 302--305.}.
% print: no opening quotation mark before "Verbum divinum" (closer only, after "Evangeliet"); opener supplied.
``\textit{Verbum divinum absolvitur
Lege et Evangelio.} 1) \textit{Lex est summa eorum, quæ Deus facere
jussit hominem, additis minis in legis transgressores.}
% print: "tansgressores".
This moral law comprises both the law revealed in the Old and the New
Testament and the natural law as well. The old ecclesiastical
dogmaticians therefore treat \emph{natural law} as a general, if
imperfect, moral law, and \emph{the moral law} as the eternally valid
part of the Mosaic Law, whose \textit{summa} is contained in the
\textit{Decalogus}. This conditional recognition of natural law is more
correct than the restriction of the moral law to the Mosaic Law, for
only through the New Testament has the moral law, with its highest
spiritual demands, and in original sin the whole depth of our misery of
sin, become manifest. \textit{Legis usus est triplex:} a)
\textit{politicus s.\ civilis, quatenus externo hominum timore
externam quandam disciplinam retinet;} b) \textit{elenchticus
s.\ pædagogicus,}
% ---- printed p.350 (PDF 363) ----
\opage{350} \textit{quatenus interno Dei terrore ad miseriam agnoscendam
fidemque impetrandam peccatorem ducit;} c) \textit{didacticus,
normativus s.\ tertius, quatenus docendo renatorum vitam regit;} 2)
\textit{Evangelium est doctrina de venia peccatorum per Christum
obtinenda.} Accordingly, the Protevangelium, all the messianic
prophecies, and the types also belong to the Gospel.''

It is, however, a large question how far this separation essentially
furthers the cognition either of the Law or of the Gospel. It always
looks as though law and gospel fell outside each other, as though there
were something in the Law that did not concern the Gospel, and
conversely. Without regard to the Gospel, both ``natural law'' and ``the
moral law'' can no doubt assert their significance in political and
civil respects; but no law that does not by its nature fall within the
Gospel's own field of vision, and so make up an integral element of the
evangelical conditions of life, will be able either ``pedagogically'' to
prepare faith in the Gospel or ``didactically'' to regulate the life of
believers. One praises the division's ``purely practical significance,''
while yet in the same breath demanding that ``law and gospel must always
go together in life\footnote{\textit{Lex et Evangelium in praxi quovis puncto
mathematico sunt conjunctiora. Confluunt enim ad poenitentiam
peccatoris, ad renovationem hominis justificati, ad conservationem
hominis renovati in perseverantia pietatis. Hollatius.}}, since the Law
% print: "mathemathico".
alone would preach despair, the Gospel alone levity''; but if law and
gospel cannot at any point in actual life do without each other, then
this is surely a sign that both must necessarily be referred to the
same principle. If the forgiveness of sins cannot be separated from the
cognition of sins, then neither can the cognition of sins be separated
from the cognition of the Law. Instead of speaking of
% ---- printed p.351 (PDF 364) ----
\opage{351} the Law and the Gospel, we would then rather speak of the Law in the
Gospel, of \emph{the Law of the Gospel}.

When one of the Pharisees, a lawyer, asked Jesus, saying, Master, which
is the great commandment in the Law? the answer ran as follows: ``Thou
shalt love the Lord thy God with all thy heart, and with all thy soul,
and with all thy mind. This is the first and great commandment. And the
second is like unto it: Thou shalt love thy neighbor as thyself. On
these two commandments hang all the Law and the Prophets'' (Matt.\ 22:35--40;
Luke 10:27; Mark 12:29). If we now compare with this what is already
said in the Sermon on the Mount: ``Think not that I am come to destroy
the Law or the Prophets; I am not come to destroy, but to fulfill!''
how then could we be in doubt that Christ wants the Law included in his
Gospel, that there is in truth a law of the Gospel? It now depends
merely on understanding this unity of law and gospel, for the
misunderstandings are manifold and go in different directions.

The principle of the Gospel is a principle of will just as much as that
of the Law. The fundamental error is that, instead of developing the
religious concepts from their own principle, one has had the doctrine
of will transformed into an objective doctrine. From an objective
standpoint every decisive, determinate act of will, be it divine or
human, must necessarily take on the stamp of arbitrariness. The law
revealed by God bears, as we have seen, in its individual commandments
the marks of arbitrariness. God simply demands so much of a human
being, neither more nor less; he could have demanded more, he could
have restricted the demands to less: such things rest on an omnipotent
good pleasure about which there can be no reasoning. This is a view as
external as it is objective; and when one then tries by the objective
route to penetrate into the divine grounds of will, the matter does not
become better
% ---- printed p.352 (PDF 365) ----
\opage{352} but worse: the contradiction is that one reasons about that about
which there can be no reasoning. It is the very same with what,
according to the orthodox doctrine, is to be the exclusive
distinguishing mark of the Gospel, namely the forgiveness of sins
through Christ. The forgiveness of sins is, objectively seen, an
arbitrary act of grace; God promises a human being the forgiveness of
sins on condition of faith. Why the merciful God has prescribed just
faith alone, faith in the merit of Jesus Christ, excluding every regard
to other possible conditions---about that there can really be no
reasoning from an objective standpoint. And yet the objective doctrine
consists precisely in all manner of reasonings going back and forth,
thus in a reasoning about that about which there can be no reasoning.
Once both law and gospel have in this way, each by itself, been made
the object of external consideration, what wonder, then, that they are
set up externally against each other? Indeed, the doctrine of
objectivity has had so bungling an effect that even the dialectic of
paradox, which after all aims so expressly at abolishing all objective
externality, has retained a questionable residue of indigestible
objectivity, and thus of externality. For what is the paradoxical, the
absurd? It is surely precisely that which consciousness comes up
against as against something foreign, something incomprehensible,
something indissoluble, thus something objective, external to thought,
about which there can be no reasoning. The forgiveness of sins, it is
said, is a paradox; what does that mean? It means that the forgiveness
of sins is an act of grace, incomprehensible in itself and arbitrary,
which can take place or fail to take place, just as God wills.

From the standpoint of faith, that is, of the principle of will itself,
the matter looks otherwise. The objectivity here is not that of
knowledge, but the objectivity proper to faith and thereby to
revelation, thus an objectivity in which the cognition of God, as much
with regard to the Law as to the Gospel, corresponds point for point to
self-cognition, and conversely. There
% ---- printed p.353 (PDF 366) ----
\opage{353} can, then, be conceived no divine law for the human self that does not
originally correspond to the inner, true nature of the human self.
``Sin is that which is against the Law,'' that is, sin is the deepest
self-contradiction, strife against God's holy will and thereby strife
against one's own fundamental will. The law that bids man love God with
all his heart, with all his soul, and with all his mind is no arbitrary
commandment. Just as little as any ground of necessity is conceivable
which, if we may express ourselves freely, should have driven God to
reveal this law of his free love to the human being darkened by sin,
just so little can any ground of freedom be conceived that could move
God to abolish this law and without it make man blessed.

In the law of love necessity is freedom, freedom necessity, but---be it
noted---in the form of freedom, not in the form of necessity. To love
God above all things is in truth to love oneself. Only self-love
purified in love to God can furnish the ideal measure for love to one's
neighbor, and the first great commandment grounds the second. For the
second is like the first. It is therefore also inconceivable that God
should be able by an arbitrary act of will to abolish the second. He
who sins against the one of the two commandments sins against both.
That man nevertheless sins against both commandments thus makes known
that man as man is a sinner before God, that the human self, in other
words, is in self-contradiction. So man stands unconditionally in need
of the forgiveness of sins. If God could now forgive transgression of
the law of love in a manner corresponding to that in which a king who
has an unlimited right of pardon can forgive any transgression
whatsoever of any civil law whatsoever, then the forgiveness of sins, as
an arbitrary act of will, could rightly be said to be paradoxical. But
that is impossible. Contradiction of
% ---- printed p.354 (PDF 367) ----
\opage{354} the Law is essentially self-contradiction; the forgiveness of sins,
the act of grace by which sin is lifted, must therefore be an act of
grace by which self-contradiction is lifted; but self-contradiction can
only be lifted by the Law's being fulfilled; the forgiveness of sins
thus cannot possibly consist in man's being in one way or another
exempted from the Law, but must consist in his being enabled, by an
express act of grace, to fulfill the Law.

\greekrun{γ}{The End of the Law.}

``Christ is the end of the Law'' (Rom.\ 10:4); in this lies both that
he fulfills the Law and that he makes an end of the Law, both by
himself being the Law's end and goal. For the holy, divine-human Self
the Law does not exist in abstract universality; the law of love, as a
determination of essence, is taken up into living love: so the Law
gains form in Christ's personality. In doing the Law full justice,
Christ does justice to himself; he fulfills the Law, that is, he fills
its empty forms with personal content; he breathes movement and life
into the dead features, he gives the cold silhouettes light and color.
That Christ is nevertheless said to be under the Law, to obey the Law,
is understood from the fact that his life on earth is in voluntary
humiliation, in suffering and self-sacrifice. Thus he has the Law not
only within him but also above him, just as he has the Father's will not
only within him but also above him. In this sense too he is religion
itself, the relation to God existing personally, in that he has his
absolute freedom in absolute obedience; his self-revelation is a
thoroughgoing revelation of obedience: ``he learned obedience by the
things which he suffered.'' When the evangelical doctrine of faith has
made a sharp separation between Christ's
active and passive obedience (\textit{obedientia activa} and
\textit{ob.\ passiva}), this is again an indication of scholastic
externality; for obedience, as
% ---- printed p.355 (PDF 368) ----
\opage{355} voluntary submission, is the categorical unity of suffering and
acting. A patience in pure suffering without inner self-activity is not
a revelation of freedom, but an unfree condition without significance
for spirit. On this very thing rests the inner infinity of the will,
the divine as well as the human, the omnipotent as well as the
impotent will: that in inward-going self-activity it can lift its
outward-going activity and transform its acting into suffering; that is
why religion speaks not only of Christ's patience but also of God's
long-suffering. ``No activity of Christ can be conceived in which,
owing to the sin of the world, there is not also a suffering; his
Passion history does not begin only in Gethsemane, for his whole life of
obedience is a history of suffering; but on the other hand no suffering
in Christ can be conceived that is not activity, that is not according
to its innermost significance a holy work of freedom; for his whole life
of obedience is a striving and victorious history of
freedom''\footnote{H.\ Martensen. Den christl.\ Dogmt.\
S.\ 314--15.}.
% print: "S, 314--15" (comma for the point after S).

The fulfillment of the Law is at the same time its sublation. It goes
with the Law in this respect as with possibility. Real possibility
requires conditions in order to pass over into actuality; it lies in
its essence that it demands its actualization. But when the possibility
is actualized, it has thereby vanished; actuality has taken its place:
the satisfaction of possibility is its sublation. In like manner with
the Law. The Law, the pure law, is with all its ideal necessity not
full actuality; for if it were full actuality, it would be sufficient
unto itself, and could to that extent contain no demand. That it
demands something means that it lacks something for the completion of
its essence; that it lacks something in itself means
% ---- printed p.356 (PDF 369) ----
\opage{356} that, precisely with regard to what it lacks, it is a possibility, not
an actuality. In being fulfilled, the Law takes up the rich sum of
moments of personality in which it has its fullness, and ceases to have
existence for itself as law. The Law does not in actuality exist
outside or alongside its fulfillment; its true, full existence is its
sublation. It is not the Law that by its fulfillment consumes life; it
is conversely life that takes up and digests the Law, reduces it to a
moment, and causes it to vanish as law. Just because Christ is the only
one who has been able to fulfill the Law, he is also the only one who
has been able to sublate the Law.

Because the fulfillment of the Law is its sublation, Christ can be Lord
of the Law by being its servant. The old dogmatic assertion that Christ,
as Lord of the Law (\textit{dominus legis}), was just as little obliged
to fulfill the Law (\textit{legi
subjectus}) as to suffer and die, is perverse and rests on an external,
finite separation of what he has done for his own sake and for our
sake. What holds of his life holds also of his fulfillment of the Law:
he has lived on earth for our sake; but from this it by no means
follows that he has not lived for his own sake. With regard to the
relation of his will to the commandments of the Law, it is on the
other hand true that in Christ's fulfillment of the Law there can be no
talk of duty and obligation. Concepts such as duty and obligation gain
particular significance only where a commanding law has a natural
self-seeking love as its opposite. The obligated subjectivity then, on
this presupposition, feels itself at strife with itself and can only
fulfill its duty by subduing its lower nature. Such a self-overcoming,
resulting from inner, finite strife, the God-man does not know. The
holy, strong will of the Logos is, in its self-denying, suffering
activity, raised above finite
% ---- printed p.357 (PDF 370) ----
\opage{357} strife with flesh and blood. It does indeed have in sensuous
immediacy, in pain and suffering, a limit that it must ceaselessly
overcome, a resistance that it must combat in every moment; but this
inner, ceaselessly continuing strife is infinitely far from being a
strife of subjectivity with itself; it is on the contrary an expression
of the fact that the striving will is, in infinite energy, absolutely at
one with itself. Through the unconditional self-overcoming that holds
for the continuity of all individual utterances of will, the personal
fundamental will has in absolute obedience won itself so
unconditionally that the continued self-overcoming is transformed into
a self-satisfaction continuing in itself; which is to say, in other
words, that unconditional obedience is a revelation of an unconditional
freedom. Christ has not, out of pure pity for us poor sinful human
beings, once for all, so to speak, done violence to himself and made
life hard for himself and fulfilled the Law in our stead, so that we
might at last, once for all, be rid of the hard, confounded Law and
come to live in pure heartiness, from the overflowing good nature of an
immediate love: this misunderstanding is obvious. It is as much out of
love for the Law as out of love for human beings that Christ has
fulfilled the law of love; for the law of love is eternally one with the
Father's will. Love is the fulfilling of the Law; therefore Christ is the
end of the Law.

Instead of beginning with the Law, the believer must therefore begin
with Christ, with the appropriation of Christ's fulfillment of the Law;
in this sense it can truly be said that Christ has fulfilled the Law for
our sake, in our stead. But that the believer, by appropriating the
fulfillment of the Law, should be able to become exempt from fulfilling
the Law is a monstrous contradiction. We see it in the relation between
\emph{love} and \emph{fear}, between \emph{faith} and \emph{works}.
% ---- printed p.358 (PDF 371) ----

\opage{358} To begin with the Law is to begin with fear. Now it is no doubt true,
what preachers of repentance and shakers of hearts have often enjoined
with strong pathos, that a person's conversion must begin with horror at
sin, with terror before the judgment of God's wrath; but it still
depends on what kind of terror it is, and by what means the cognition
of sin is awakened. To begin with the thunder of the Law is just as
sensuous as to begin with the fire of hell and the lake of brimstone.
% print: "Svovlpølen, I" (comma for full stop).
In the thunder of the Law only the command of omnipotence is heard; such
thunder does not awaken the fear of God; it awakens only fright before
the Almighty; that may work on the rude mob, but the questionable thing
is that a person whose fear of God begins with fright in the blood
begins with a hidden admixture of superstition that he will get over
late or never. In vain is love presented to him in the image of the
Crucified; he does not see love, cannot discover its essence; for he
has it in the blood, and sees in the blood---nothing but sacrificial
blood, the blood of expiation that the angry God has demanded. It holds
of love that it casts out fear; but it does not hold of fear, the fear
that oppresses the senses of the soul, that it drives love in. To begin
with love is to begin with Christ, is to begin not outside the Law but
inside the Law, with the fulfillment of the Law, with the law of the
Spirit, with the self's eternal, ideal demand.

Yet the law of the Spirit too can be conceived in so glaringly abstract
a form that the presentation of love's demand, even if it is formally
entirely correct, consumes the content of love; as when it is said to a
person: ``God is your mortal enemy; he, Love, he wants out of love for
you to be loved by you: this means that you must die, die away,
otherwise you cannot love him. So there he sits, omnipresent and
omniscient as he is, and watches you, knows
% ---- printed p.359 (PDF 372) ----
\opage{359} even the least thing that goes on in you---that he knows, your mortal
enemy! Beware of wishing for anything, beware of fearing anything; for
what you wish for is not fulfilled, but rather the opposite, and what
you fear, the more you fear it the sooner it comes upon you. For he
loves you, and he wants to be loved by you, both out of love; but as
soon as there is something you wish for, you are not thinking of him,
just as when there is something you fear; or if you bring him into
connection with your wishing and fearing, then you are surely not
thinking of him in and for himself, that is, you do not love him---and
he wants to be loved, wants it out of
love''\footnote{S.\ Kierkegaard. Øieblikket. Nr.\ 5.}.

In Christ's life the law of love has its rich, intuitable content; he
gives a person eternal life for the sacrifice of the temporal. He does
not say to the one who fears something: God is your mortal enemy; but
he says to the one who hopes for something: I am the resurrection and
the life. It is not renunciation of the temporal that teaches a person
to love the eternal; but it is conversely love for the eternal that
gives strength to sacrifice the temporal. It is not fear of God, not a
thoroughly reflected, clear knowledge of what and how much God can
demand and how infinitely strict the Law is, that awakens a living love
to God, but it is conversely love to God that reconciles a person with
the strictness of the Law and awakens true fear of God in the heart.
The fear of God that is not awakened by God's love is without
understanding of God and therefore without essential self-understanding;
it is, even if the sinner trembles before the judgment of God's wrath
as the malefactor before the executioner's axe, just as false and
spiritless as the love of one of the awakened that has remade God after
the wishes of the natural heart and, in thoughtless, jubilant
boldness, forgotten fear. As
% ---- printed p.360 (PDF 373) ----
\opage{360} love is related to fear, so the Gospel is related to the Law. One does
not begin with a fear without love, but one begins with a love that has
fear within it; one does not begin with a law without gospel, but one
begins with a gospel that has the Law, the Law fulfilled by Christ,
within it.

If we consider the relation between faith and works, we come to the
same result. To begin with the commandments of the Law is to begin with
good works and, insofar as the revealed Law presupposes faith, with the
good works of faith: the emphasis falls on the works. But good works are
only utterances of the good disposition; it is not the work that makes
the disposition good, but conversely. With works, therefore, as with the
Law, namely the Law of the Gospel, it depends essentially on the
disposition, on the heart, on the purity, truth, and strength of the
will, thus on free inwardness, thus on faith. The believer does not
believe in his own faith; he believes in Christ. To believe in Christ is
to believe in the truth in Christ's life, that is, in the truth of the
fulfillment of the Law. Only in the same measure as the believer
appropriates the truth in the fulfillment of the Law can he appropriate
Christ's life. To appropriate the truth in the fulfillment of the Law is
possible only by oneself fulfilling the Law; but to fulfill the Law is
possible only by a work. So we do after all come to begin with a work!
Yes, but with a work of the Gospel, the work of God without which
Christ's life and all his works would be in vain for us, namely with
faith. ``What shall we do, that we might work the works of God?'' We
know the Savior's answer to those who asked: ``This is the work of God,
that ye believe on him whom he hath sent''
(John 6:29). Faith itself is thus a work, the inner, original work of
all good works, the work by which the fulfillment of the Law
% ---- printed p.361 (PDF 374) ----
\opage{361} is appropriated, the only work by which man on his part can fulfill
the Law. But to fulfill the Law by believing in Christ is to come to an
end of the Law; for Christ is the end of the Law.


\runhead{b) The Atonement.}
\parmark{24}

To come to an end of the law, not by breaking it but by fulfilling
it, is an indispensable condition for man's coming to an understanding
with God in spirit and in truth, for his being reconciled with God.
The task is great, human and yet superhuman, a world-task which the
God-man alone is able to accomplish. The law in its widest scope is at
once both a law of nature and a law of the spirit; the latter cannot be
fulfilled without the former. But although the law of nature and the
law of the spirit are in the Mystery one and the same law, they
nevertheless constitute in the actual world two absolutely heterogeneous
systems, that of necessity and that of freedom. To remold the laws of
the spirit after the laws of nature is paganism; self-conscious
paganism is the denial of spirit: on the terms of the denial of spirit
it is an easy matter, in knowledge, to come to terms with actuality and
to be reconciled in metaphysical absolutism. But when the laws of the
spirit are proclaimed, and they are proclaimed only in the word of
revelation, the discord between nature and spirit begins.

At the Old Testament stage the discord has come to light on every side.
The people of Israel, with its firm faith in Jehovah's loving-kindness,
has never quite been able to come to terms with Jehovah himself; for it
has never fully understood the conditions of the fulfillment of the
law. To obedience under the laws of the spirit are attached promises
whose fulfillment is everywhere conditioned by partial alterations in
the laws of nature: so
% ---- printed p.362 (PDF 375) ----
\opage{362} outward help is given by outward miracles. But the extraordinary help
will not quite suffice either for the individual or for the whole
people. The pious, who do the works of the law, do not after all
receive all the blessing that is promised them in the law; and behold
the ungodly! ``they bend the bow, they lay their arrow upon the string,
to shoot in darkness at the upright in heart'', and yet the curse will
not quite strike these ungodly: there is always something in the way.
What is it? The sin of men and the wrath of Jehovah! So then sin must
be repented of; but the penitent nevertheless sighs after the help of
the Almighty as after a temporal aid. So Jehovah must be propitiated by
sacrifices; but bloody sacrifices of dumb beasts do not after all avert
the wrath. So there must be waiting in patience, and ``the servant of
Jehovah'' must suffer for the sins of the people (Isa.\ 53 ff.); but the
secret of his sufferings no one can fathom.

The secret is that the religious life, because of the strife between
nature-bound actuality and the spirit's claim to the freedom of the
self, essentially is and must be a life in suffering and passion. The
believer is not exempted from being under the law of nature because he
has come under the law of the spirit; on the contrary! now he feels
himself precisely at strife with the world and in discord with
himself, for ``the flesh lusteth against the Spirit, and the Spirit
against the flesh''. To fulfill the law of the spirit is a suffering
to him. % print: "; At opfylde" (semicolon followed by a capital).
Therefore he never gets it fulfilled either; he remains in
self-contradiction, and so in sin, --- and yet in suffering all the
same. Through the suffering of repentance the sinner hopes to obtain
forgiveness from God; repentance is a suffering fulfillment of the law
of the spirit, a sacrificial fire of the soul in which what is selfish
is burned away. But that the inward self-sacrifice is imperfect is
tellingly expressed by the fact that help is sought in outward
sacrifices. Instead of sacrificing himself, his natural life, his whole
natural mind in suffering, man receives, at the first stage of
% ---- printed p.363 (PDF 376) ----
\opage{363} revelation, permission to sacrifice the animal in his stead, whereby
the suffering then ceases in a finite manner and the sacrifice becomes
symbolic. But the symbol of sacrifice is only a symbolic fulfillment of
the law, and the symbolic fulfillment of the law can then also bring
about only a symbolic atonement. So it begins all over again. Jehovah,
propitiated by the sacrifice, ought now, once the sin is forgiven, to
exempt the pious from all further suffering, bend the laws of nature,
and so arrange the outer world that ``the righteous'' might at last get
the better of ``the unrighteous'' and live in peace and happiness. That
does not happen, and cannot happen, for the law is in actuality not
fulfilled. God is not yet reconciled with man, nor man with God; for
God and man each have their own concept of self-sacrifice and
suffering. According to man's concept the suffering of the pious must
cease, and the cessation of suffering be an unfailing sign of piety;
but according to the divine concept the suffering shall not and cannot
cease, for the fulfillment of the law is conditioned by self-sacrifice,
and self-sacrifice is categorically voluntary suffering. The law can be
fulfilled only by self-sacrifice, and with it voluntary suffering,
being carried through unconditionally; and the law shall be fulfilled.

Originally, in the deed, in spirit and in truth, no sinful man can
understand this secret of the sufferings; only the God-man, the pure,
sinless man, knows the secret; only the God-man can, by unconditional
self-sacrifice, reveal the secret of the voluntary sufferings, fulfill
the law, and accomplish the atonement.

By considering the doctrine of the atonement first from the
scholastic-objective and then from the mystical-subjective side, we
prepare our way to an insight into the two-sided essence of the
atonement.
% ---- printed p.364 (PDF 377) ----

\opage{364} \greekrun{α}{The Scholastic-Objective Doctrine of the Atonement.}

By his voluntary, and yet historically necessary, suffering and death,
dependent on given conditions of the world, Christ has in the place of
men borne the punishment of sin and satisfied the divine justice. When
this thought is to be held fast in the objectivity of a scientific
representation and developed in scholastic juridical form, we arrive
at the satisfaction theory set out by Anselm.\footnote{\textit{Cur Deus
homo.}} ``To sin'', says Anselm, ``is nothing other than not to give
God what one owes him. But every rational creature, be it angel or man,
owes it to God to be subject to his will; this is the whole honor which
God demands of us. Now there is nothing less tolerable than that the
creature should rob the Creator of his honor without making good to him
again what was taken; since, then, it is the greatest injustice to put
up with such a thing, God's justice demands of him that he maintain his
dignity. With regard to the order of the world, too, it would be a
disorder, a disturbance of the beauty and harmony of the world, which
God as ruler of the world ought not to tolerate, to let sin go
unpunished. It is therefore necessary either that man, who has sinned,
voluntarily give back to God what he has robbed him of and in addition
render full compensation for the offense done him, or else that God
must take by force what is due to him: the former is voluntary
satisfaction, the latter involuntary, or punishment. But man for his
part, once he has sunk into sin, cannot even with the best will make
such satisfaction to God as might move him to cancel the punishment. To
a perfect religious and moral conduct he is obliged in every present
% ---- printed p.365 (PDF 378) ----
\opage{365} moment, and cannot with it pay for sins committed. Moreover, the guilt
is infinite, since in God's commandment it is God himself who is
offended. But over against God's justice stands God's goodness. This
would indeed have to be changeable, which it is not, if it did not move
him to complete in man the work of love that was begun at creation;
just as it would also conflict with God's wisdom if he were to let the
most excellent of his works, the rational creature, go to ruin. The task
is therefore to find a being that is able to render complete
satisfaction. Such a being must be able to pay to God, for man's sin,
more than all that is outside God. But to give so much of its own is
possible only for a being that is itself more than all that is outside
God. But more than all that is outside God is only God himself.
Consequently no other than God himself can render the required
compensation. But on the other hand a man must nevertheless render it;
for it is to be rendered for men; since, consequently, only God can
render it and only man ought to render it, it follows from this that
God as man, or the God-man, must render it. The satisfaction which the
God-man is to render cannot consist in his giving himself up to God in
active obedience; for this obedience every rational creature,
consequently the God-man too, always owes to God. But to suffer death
he is neither physically compelled, since he is almighty, nor morally
obliged, since he is sinless: his voluntary death can therefore count
as the required satisfaction. The surrender of so great a good as the
God-man's life God cannot leave unrewarded. But the Son of God, who
already has everything, cannot receive any reward for himself; God must
then give him freedom to make over the merited reward to whom he will;
and he makes it over with most propriety to men, who are most in need
of it and are related to him according to his
% ---- printed p.366 (PDF 379) ----
\opage{366} human side, and consequently are his natural heirs''. Strauss
has\footnote{Troesl.\ II.\ S.\ 211.} aptly remarked of this satisfaction
theory that it is like a ``parallelogram of forces, in which the divine
love by itself claimed outright pardon, the divine justice by itself
unsparing punishment of the guilty, but both, acting together, had the
diagonal of satisfaction as their result''.

That the doctrine of the atonement in so objective a form of necessity
was bound to provoke objections is understandable, since necessity and
arbitrariness easily change places in scholastic reasoning. According
to the theory the mode of atonement described by Anselm would have to
be the only possible one, and the suffering and death of Jesus for the
sake of sins therefore necessary; but the notion of God's arbitrariness
would not be driven back. Anselm himself concedes that God was in no
way compelled to become man, and Thomas Aquinas declares flatly that
God, had he so willed, could have freed man from sin without any
satisfaction at all.

At the time of the Reformation it was above all the Socinians who
subjected the juridical-objective satisfaction theory to a penetrating
and sharp critique. ``Forgiveness of sins on the ground of
satisfaction'', says Faustus
Socinus,\footnote{\textit{Christl.\ rel.\ breviss.\ institut.}} ``is a
concept that not only conflicts with the divine perfections but also
contradicts itself. It does the former in that it is as little
conceivable that God should not have been able to forgive men their
sins without satisfaction as that he should not have been willing to:
the former is namely not conceivable, since God's right would be more
restricted than that of any man if he were not free simply to forgive
offenses ---
% ---- printed p.367 (PDF 380) ----
\opage{367} and what else are sins? --- against his person. It is said, to be
sure, that God's justice demands either punishment or satisfaction for
sins. But this rests on a false notion of the divine justice. This is
not essentially punitive justice, but only the equity and rightness of
the divine mode of action in general, and it can show itself just as
well in forgiving as in punishing sins. Only in this generality is
justice an abiding determination \textit{(perpetua qualitas)} of the
divine essence: whether God in accordance with it punishes or forgives
is a passing measure of government or determination of his will. But
that God should indeed have been able to forgive without satisfaction,
but not have willed it, would conflict all too much with his goodness.
--- But besides, the concept of a forgiveness of sins in consequence of
a compensation dissolves of itself, since it is composed of two
mutually exclusive elements. For to demand compensation is precisely to
exact the debt that is to be remitted in the forgiveness. That it is
demanded of another than the one who actually owed it signifies
nothing, since the creditor (God) nevertheless gets the debt paid, and
consequently remits nothing: it is no remission of the debt, but only
an exchange of debtor and payer that takes place here''.

The false objectivity which, in the name of science, would bind
reflection to fixed factual presuppositions is just the thing to
provoke the resisting understanding and to develop its finite, and for
that very reason endless, reasoning from grounds: grounds and
counter-grounds alternate into infinity. Thus the dogmatic-objective
doctrine of the atonement has been transformed into a formal suit for
debt, with the examination of witnesses and attorneys' pleas for and
against. It is not the arguments in their particulars but the basis of
the argumentation, the grounding that is rational in unreason, which on
the whole characterizes the
% ---- printed p.368 (PDF 381) ----
\opage{368} error. Before we attempt, however, to bring out the grains of truth
hidden in the heaps of reasoning, we shall first consider the
corresponding one-sidedness, namely:

\greekrun{β}{The Mystical and Subjective Doctrine of the Atonement.}

In opposition to the objective, historical-juridical conception of the
Atoner's suffering and death, pious souls in the Middle Ages asserted
inwardness; but since inwardness stood in decisive conflict with the
whole outward, Catholic-official ecclesiastical system, the mood
against the objective was naturally bound to be polemical; it is thus
easily explained that sectaries such as the Bogomils, the Beghards, the
Brethren and Sisters of the Free Spirit, and others, carried the
mystical interiorization of the divine essence to caricatured
extremity. According to the teaching of the Bogomils the Son of God is
born in every believing soul. The believing soul that awakens faith in
another is a Mother of God (\textit{θεοτόκος}), who does not fall short
of Mary%
\footnote{\textit{Λέγουσι τοὺς τῆς πίστεως αὐτῶν, ὅσοις ἐνοικήσει τὸ
παρ᾽ αὐτοῖς ἅγιον πνεῦμα, πάντας θεοτόκους καὶ εἶναι καὶ
ὀνομάζεσθαι, βαστασαντας καὶ αὐτοὺς τὸν λόγον τοῦ θεοῦ, καὶ
γεννήσαντας αὐτὸν διὰ τοῦ διδασκειν ἑτέρους, καὶ μηδὲν πλέον αὐτῶν
ἔχειν τὴν πρώτην θεοτόκον.} --- Euthymius Zygadenus. Cf.\
Gieseler. Kirchengeschichte. II.\ Bd.\ 2 Abthl.\ Bonn 1848.
S.\ 684.}. With the Siblings of the Free Spirit every pious person was a
Christ, in whom God becomes man; their claim to an inner unity of
essence with God and Christ seems, in metaphysical absoluteness, even to
have approached the
insane.\footnote{\textit{Item credunt, se esse Deum per naturam
sine omni distinctione. Item quod sint in eis omnes perfectiones
divinæ, ita quod dicunt, se esse æternos et in æternitate. Item
dicunt, se omnia creasse et plus creasse, quam Deus.} Giesel.\
Anfr.\ Skr.\ S.\ 645. --- This, one may say, is indeed ``to relate
oneself absolutely to the absolute''.}
% ---- printed p.369 (PDF 382) ----

\opage{369} When the Reformers, in adherence to the Anselmian satisfaction theory,
had the merit of Christ transferred to believers by an \textit{actus
forensis}, Schwenkfeldt and Andreas Osiander protested against the
externality of this conception and laid the decisive emphasis on the
inward life of the soul, on the living appropriation of Christ's
divine nature, on the Christ growing up in the heart. In Valentin
Weigel and Jacob Böhme the doctrine of the atonement acquires a still
richer mystical luster of color. ``Men have'', says Jacob Böhme,
``dipped Christ's mantle with its purple color in Christ's blood and
taken it as a covering, but thereby only covered up the antichristian
child of their own will. For men gladly play the hypocrite for this
child and cover it with Christ's suffering, merit and death, and
comfort it with the thought that Christ has paid for it and that it
need only comfort itself with his merit and accept this in faith as a
satisfaction. But the matter hangs together quite otherwise: not by an
outward attestation of a special acceptance into grace, by a foreign
merit and an imputed righteousness, are we accepted as children of God;
but by an indwelling, essential grace pervading all the members, in
that Christ, the conqueror of death, with his life, essence and power
rises in us from our death and reigns and works in us as a tendril on
his vine''. After Jacob Böhme, Angelus Silesius in ``holy
soul's delight'' was a minnesinger of pious inwardness: ``Ich muss Maria
sein und Gott aus mir gebären, Soll er mich ewiglich der Seligkeit
gewähren''. Not so rich in feeling, but still in the same direction,
the striving of the Quakers aims at a spiritual appropriation of the
atonement, a birth of Christ in the inward man (\textit{internus partus
Christi in homine}), so that they even assume that a man may, by virtue
of ``the inner light'', become a partaker in the ``Mystery of the
death'' of Jesus without knowing his history.

It was now certainly on the whole by no means the intention of those
% ---- printed p.370 (PDF 383) ----
\opage{370} mystics, resting as they did on evangelical presuppositions, either to
deny or to belittle the significance of Christ's actual suffering and
death as a historical fact; it was only the spiritless objectivity that
they wanted to attack. But the mystical lies precisely in this, that
they were unable to make clear to themselves the relation between the
appropriation itself and that which was to be appropriated. Had the
Quakers, for example, really given themselves an account of how the
historically actual atonement, as the first preparatory cause
(\textit{causa procurans et efficiens}), properly stands related to the
working of the inner light, they would surely not have presupposed the
possibility that this revelation of light in the spirit could at any
point be separated from the presupposed revelation of life in Christ's
history.

So pernicious was the effect of the barren dogmatic objectivity that
philosophical thinking, which at last set itself the task of making an
end of all mysticism, instead of clarifying the historical presupposition
of the religious life let go of the historical in order to grasp
exclusively the subjectively ethical, the moral. It is the Kantian
doctrine of the atonement that in this respect draws attention to
itself. According to Kant it does not essentially matter what Christ as
a historical person has done or suffered: the most important thing is
that he is the ideal of humanity, a personification of the good
principle. ``The good principle'', says Kant, ``is present just as much
in the abandonment of the evil disposition as in the adoption of the
good, and the pain that rightly accompanies the former arises wholly
from the latter. The transition from the corrupt disposition to the
good is, as `the dying of the old man, the crucifying of the flesh', in
itself already sacrifice and the beginning of a long series of evils in
life which the new man, in the disposition of the Son of God, namely
merely for the sake of the good, willingly takes upon himself; but which
properly belong to another as punishment, namely to the old
% ---- printed p.371 (PDF 384) ----
\opage{371} man (for this is morally another). Although, therefore,
\emph{physically} (considered according to his empirical character as a
sensuous being) he is the very same man deserving of punishment and
must as such be judged before every moral tribunal, and so also be
judged by himself, yet in his new disposition (as an intelligible
being), before a divine judge for whom this disposition takes the place
of the deed, he is \emph{morally} another; and this disposition in its
purity, as that of the Son of God, which man has taken up into himself,
or (if we personify this Idea) the Son of God, bears for this man, and
for all who (practically) believe in him, as \emph{substitute} the
guilt of sin, satisfies the highest justice by his suffering and death
as \emph{Redeemer}, and as \emph{advocate} makes it possible for them to
hope to appear as justified before their judge; only that (according to
this mode of representation) that suffering which the new man, in dying
to the \emph{old}, must take over for his whole life, is represented,
when referred to the representative of humanity, as a death which he
suffered once for all''.\footnote{Die Religion innerhalb
der Grenzen der blossen Vernunft: Schriften zur Philosophie der
Religion. Leipzig 1839. S.\ 239--40.}

The inner kinship between mysticism and moralism is unmistakable. From
both sides it is overlooked that the suffering, dying Savior is and
must unconditionally be the mediator between God and suffering, dying
man. The mystic, like the moralist, wants to live in an unmediated
relation to God. The difference is only that the mystic's relation to
God has its inner immediacy in feeling and essentiality, the
rationalist's, on the other hand, in reflection and will. The mystic
forgets that inner revelation is conditioned by outer; the moralist
rejects inner
% ---- printed p.372 (PDF 385) ----
\opage{372} and outer revelation alike, and overlooks that a relation of will to
God is impossible without revelation.

% print: "dobbeltsidigc" (wrong sort c for e) in the head.
\greekrun{γ}{The Two-Sided Essence of the Atonement.}

Moralism, in all its one-sidedness, has nevertheless held fast to the
point which may rightly be called the religious cardinal point, namely
the will, and that not merely as will in man but also as will in God;
the rationalist-moral relation to God is after all still a decisive
relation of will. The advance in the philosophy of religion made by
speculation consists, then, only in this, that the relation of will is
imperceptibly transformed into a metaphysical relation of essence. The
significance of the will is not, to be sure, directly denied; the
philosophy of the absolute has, of course, a place in the system for
the concept of will just as much as for the concept of knowledge or any
other of the fundamental concepts of the Idea. But when the concept of
will is to be developed out of the concept of God with the immanent
necessity of the Idea, the principle of will ceases to be absolute: the
will departs from God and remains behind with man alone. In this sense
we may, without in any way failing to recognize the great significance
of speculative philosophy for science, say with truth that, as
philosophy of religion, it has contributed its share to undermining
religion.

The Hegelian doctrine of the atonement, in its advance out of the
Schellingian, Strauss has set out briefly and clearly in the following
main features: ``In itself \textit{(an sich)} --- this is the
theme of all Hegelian treatments of this point --- the knowledge that
the opposition between substance and subject, or between selfless
universality and infinity and the individual finite self and natural
existence in general, is at bottom sublated and does not exist, is the
reconciliation of spirit with itself. But for the non-comprehending
consciousness this in-itself \textit{(Ansich)} takes on the
form of something existing and something represented. For it
% ---- printed p.373 (PDF 386) ----
\opage{373} the comprehension of this is not a grasping of that concept which knows
those two sides as united in and for themselves, but a grasping of the
representation that the divine essence is reconciled with itself by the
act in which it went out of itself, by its having become man and died.
For the finitude and naturalness in which man finds himself seems, in a
twofold respect, not to accord with God's universality, the Idea that
is free and pure in itself. First, inasmuch as the former, in contrast
to the breadth of the divine Idea, its merciful self-outpouring, is, as
an individual self, reflected into itself, egoistic and evil; second,
because it moreover feels itself unhappy in the want of its nature and
its susceptibility to sufferings. This disharmony cannot be sublated
one-sidedly by man, who as self stands at the one extreme; rather, its
joining together of both sides must have their unity, existing in and
for itself, as its presupposition. This unity is present in the fact
that both the clear essence and the self reflected into itself are
pithless abstractions, that neither does the former have actuality
without the latter, nor the latter truth without the former; or more
precisely: that the infinite equality with itself is \textit{eo ipso}
difference from itself, and consequently finitude.''


``The God-man dies; this death is on the one hand, especially in its
potentiation as the death of a criminal, the surest proof that he is
human; this has shown itself in him to the very uttermost point, and even
at this point it has not yet shown itself to be incompatible with the
divine life. But on the other hand, through death the humanity, the
immediacy, of the existing singularity is sublated and thereby determined
as a merely vanishing moment in the divine
life\footnote{A passage quite characteristic of the Hegelian view as a
whole is, e.g., the following (Religionsphil.\ XII Bd.\ 2det Opl.\
Berl.\ 1840. S.\ 301--2): ``If in the first sphere'' --- the Kingdom of
the Father --- ``we conceived God in pure thought, then in this second
sphere'' --- the Kingdom of the Son --- ``it begins with \emph{immediacy
for intuition} and for sensuous representation. The process is now
this, that the \emph{immediate singularity is sublated}\dots\ and this
happens through death; Christ's death is the death of this death, the
negation of negation. \dots\ In the death of Christ it is finally still
to be emphasized that it is God who kills death, in that he goes forth
out of it; thereby \emph{finitude, humanity and humiliation} are
transferred as something alien to Christ, as to him who is simply God:
it becomes evident that finitude is alien to him and \emph{received from
another}; this other is now human beings, who stand over against the
divine process. It is their finitude that Christ has taken upon himself,
this finitude in all its forms, which at their extreme point is
\emph{evil}; this humanity, which is itself a moment in the divine life,
is now determined as something alien, not belonging to God: this
finitude in its being-for-itself against God is evil, something alien to
him; but he has taken it upon himself in order to kill it by his death.
The shameful death, as the monstrous \emph{union of these absolute
extremes}, is thereby at the same time \emph{infinite love}.''

``It is infinite love that God has posited himself as identical with
what is alien to him, in order to kill it. This is the meaning of
Christ's death. Christ has borne the sins of the world, has reconciled
God, it is said \dots\ This death, though natural, is God's death and
thus satisfactory for us, since it is the \emph{absolute history of the
divine Idea}, which presents what has happened in and for itself, and
what eternally happens \dots\ That, e.g., the criminal can be punished by
the judge, and that this punishment is the carrying-out of the law and
its atonement, this is not done by the judge, nor by the criminal through
his suffering of punishment as a particular external event, but it is
the nature of the matter, the necessity of the concept \dots\ So it is
now also with that satisfaction for us; what therefore lies at the
ground is this, that that satisfaction has taken place \emph{in and for
itself}: not an \emph{alien} sacrifice has been offered, not
\emph{another} punished, in order that punishment might be exercised.
Each must for himself, out of his own subjectivity and guilt, be and
render what he is to be; but what he thus is for himself may not be as
something accidental, as \emph{his} condition, but must be something
\emph{true} \dots\ The ground of redemption is therefore that history,
for it is the matter \emph{in and for itself} \dots\ The confirmation
that it is the true is the view which that history gives, and in which
the individual grasps Christ's merit''.} \dots\ What
% ---- printed p.374 (PDF 387) ----
\opage{374} lies in and for itself in the concept of spirit, and on which the
reconciliation of the phenomenal human being with his essence rests,
namely that what at first only seems to be a limit, a negation
% ---- printed p.375 (PDF 388) ----
\opage{375} for spirit --- naturalness, selfhood with its deficiencies, though
to be sure only as constantly overcome and reduced to a moment --- is on
the contrary its actuality: this is here represented as a past
existence in which the God-man, once revealed on earth, is supposed to
have brought about the reconciliation of human beings with God''.

From scholastic objectivity, then, through the subjective one-sidedness
of mysticism and moralism, the movement of thought here arrives at a
dialectical mediation between the objective and the subjective, in which
the doctrine of the atonement at last obtains its scientific conclusion.
That the Hegelian philosophy actually marks a conclusion is evident
partly from the inner necessity with which it develops out of the
preceding theories, and from the superiority with which it comprehends
and takes up the moments of truth contained in them; partly from the
fact that it cannot essentially be gone beyond without the boundary of
the science of reason being overstepped at the same time. Starting from
the principle of faith, the unconditioned principle of will, we now let
the movement of thought proceed in the inverse direction, in that,
beginning with the subjective significance of the atonement, we step by
step make our way toward an understanding of the objective.

It is the work of the Son that he has revealed to us the secret of
sufferings: only by striving as he did and suffering as he did can a
human being know himself atoned with God. Now both Hegel and Kant seem
able to understand this; but if we look more closely, they nevertheless
do not understand it: it can only be understood in faith. The human
being is --- to begin with the subjective --- by fulfilling the law
% ---- printed p.376 (PDF 389) ----
\opage{376} himself to suffer for his sins; is, if we inculcate upon ourselves
exclusively the moral, as a rational being, without any special help
whatsoever on God's part, voluntarily to bear the punishment of sin; is
of himself to cognize that voluntary suffering for the sake of sin and
for the sake of fulfilling the law belongs to life, so that the whole of
temporal life is for that reason, and must be, a continued, a
voluntarily continued suffering; is to understand this in such purity
and depth that suffering itself is transformed into blessedness ---
blessed are they who suffer for righteousness' sake --- and with this
moral self-understanding to know himself in so true and inward an
understanding with God that by suffering for his sin he can blot out his
sin and atone himself with God. This is, without a shadow of
exaggeration, the demand of spirit, of freedom, of righteousness.

If the demand of the spirit is now held fast in its original
earnestness, one easily sees what comes of transforming the doctrine of
the atonement into a scientific doctrine of reason, whether one
otherwise takes the rational from the theoretical or from the practical
side. The Hegelian atonement is theoretical, the Kantian practical.

The theoretical atonement is an atonement in knowledge, is the knowing
consciousness's reconciliation with itself. Atonement presupposes
conflict: as the conflict is, so must the atonement be. The conflict
lies in the contradiction, a conflict of thought between opposed
concepts; to resolve the contradiction is to bring about reconciliation,
is to reconcile conceiving thought with its own conceiving: the
resolution of conceptual contradiction is the reconciliation of thought.
It is ``the opposition between substance and subject, between `the
selfless universality' and `the individual finite self','' % print: outer quotation never closed; closed here
with all that must be ranged thereunder which in natural existence gives
itself the semblance of contradiction. The resolution of the
contradiction then consists in a knowledge that the opposition itself,
with the conflict following from it, is only a semblance; once the
semblance has been seen through as semblance, it can no longer
% ---- printed p.377 (PDF 390) ----
\opage{377} deceive. The transparency of the semblance is its nothingness, its
impotence. Here, in the eternal sense, for the sake of the spirit, of the
eternal, no further exertion is needed either with the opposition or
with the contradiction. The deception is seen through; the deception
seen through is in knowledge essentially lifted, just as the
contradiction dissolved in the concept is lifted: this is ``the
reconciliation of spirit with itself.'' What this reconciliation means
practically is seen from the example of the criminal and the judge. As
calmly objective as a judge, in objective reconciliation with the law,
can pronounce the death sentence upon the aggravated criminal, just as
calmly objective can the criminal, in objective reconciliation with the
law, hear his death sentence; he can, as surely as he is a knower,
console himself with the nature of the matter, the necessity of the
concept; when fantasy shows him the block and the axe, and in the last
night in prison lets him feel in advance the killing blow: then he is
calmly objective, in objective reconciliation with the law. For the law
receives satisfaction, the contradiction is resolved; the killing blow
means that the opposition between ``the individual finite self'' and
``the selfless universal'' is at bottom sublated: spirit is, after all,
reconciled with itself!

With the practical atonement in Kant the matter stands somewhat
differently. In Kant subjectivity does not get off so easily with an
absolute self-reconciliation as in Hegel. Just as for the fulfillment of
the law it is not enough to know one's duty, but what matters is to do
one's duty, so in the moral sense it is not enough either to know
oneself metaphysically reconciled. The task is to actualize the
atonement ethically. The new, morally improved human being is then to
begin with so absolute an improvement that, though physically the same
human being, it is to such a degree beyond the old that it can both bear
the punishment which the old has deserved and still punctually fulfill
the law; can this be meant in earnest? If it is actually
% ---- printed p.378 (PDF 391) ----
\opage{378} meant in earnest, then one may well ask whether there is, or ever
has been, any actual human being who has been capable of this. For if an
ideal is set up here which can never be realized, then the demand
contains a contradiction. To demand of a human being what manifestly
exceeds human powers is to demand the unreasonable; with such a demand
it can never be fully in earnest; it conflicts with the well-known rule:
\textit{ultra posse nemo obligatur}. It is in this respect quite
striking that Kant, who after all wants Christ to be regarded as ``a
personification of the good principle,'' thus as a moral ideal, can take
the question of the ideal's actual existence so lightly. That the
actuality of the ideal is reduced to something indifferent means that
its actualization is tacitly assumed to be impossible. So it remains an
endless striving and an infinite ought, and the new human being of
moralism finds atonement only in the consolation that the rational God
is surely fair enough to take the will for the deed.

At the stage of rationalism the divine springs of life in the soul are
dried up; it is the law of reason in and by which alone the religious
consciousness is to live; in the mystical relation to God, on the other
hand, the human being draws spiritual nourishment from the ``fullness''
in God. The inner revelation, the life in God, the inner light, is a
fusion of the divine with the human, in which the soul can taste ``the
sweetness of atonement.'' But with all this the mystics, in spite of
feelings, stirrings, images and figurative descriptions, are not able to
say definitely what kind of Christ it originally is that is born in the
soul, or whence he comes, or what kind of form he actually takes on. One
of two things: either the divine child that is to be born in the heart
must make the mystic himself an independent son of God, or else be
cognized as an image of the historical Christ, whose actual life then
becomes the proper
% ---- printed p.379 (PDF 392) ----
\opage{379} exemplar. If the mystic chooses the first explanation, then he can
certainly disregard the historical exterior; but his interior is thrown
into confusion, he stamps himself as a fanatic. The fanaticism consists
in his overlooking the qualitative difference between the God-man and
every other human being. Instead of appropriating the life of the
God-man, he wants himself to be God-man: so the natural and the
spiritual, truth and imagination, the sinful and the holy are mixed in
utter unclarity. If the mystic chooses the latter explanation and holds
it fast, then mysticism will sublate itself. For by acknowledging the
historical Christ as the only true, actual God-man, he will bethink
himself that Christ alone, in the voluntary suffering of unconditioned
self-sacrifice, has borne the sin of the race and accomplished the
atonement. The inward Christ will thus, in appropriation, with the
cooperation of the Spirit, be derived from the outward, form itself
after the outward, and take on the form which as a psychic image
corresponds to its historical exemplar.

The truth in the objective doctrine of the atonement now shows itself
in this, that the gaze is fixed on the Atoner himself, his actual
suffering and death. But the Atoner is not an object of knowledge but of
faith. What is false is that the object of faith is scholastically
changed into an object of knowledge; so it comes about that
appropriation is in a spiritless way separated from its object. First
the object, the object of knowledge, the external historical object by
itself, then appropriation comes afterward! Hence the crass doctrine of
an external transfer of all our sin and guilt onto him, mediated by an
agreement of love with justice, and the great, soul-harrowing spectacle
of the Father ``smiting his beloved Son with the rod of wrath'' in our
place, chastising and punishing him in our place; whereas the truth is
that the suffering and death of the God-man is an inner consequence of
his self-sacrificing love, a consequence in which
% ---- printed p.380 (PDF 393) ----
\opage{380} he fully reveals his eternal certainty of the Father's good
pleasure, an undeceivable certainty which not even the painful moment of
God-forsakenness has been able to shake. Hence the spiritless
representation of a juridical transfer of all his merit onto us, a merit
which is even conceived so externally that the matter itself is degraded
to comparison with a money matter, where the one rich man pays the debt
in place of the many debtors; whereas the truth is that the transfer of
merit is conditioned by appropriation. Appropriation is in faith.
Insofar as faith is unconditioned as faith, the believer too can
unconditionally count on Christ's satisfaction, unconditionally
appropriate his merit. But the unconditioned passes through conditions.
That faith is unconditioned is to be known by this, that the
appropriation of the Atoner's suffering and death is unconditioned. He
suffered for us, that we by the power of his suffering should
voluntarily suffer with him and in self-understanding understand the
secret of sufferings which he revealed; he died and was buried, that we,
atoned with God in life and death, should die and be buried with him.

This, then, is the two-sided essence of the atonement, that the
objective is determined point by point by the subjective, and
conversely; that God in Christ has made the atonement objective and
reconciled the world to himself, in order that the world --- the
believers in the world --- should appropriate the atonement subjectively,
be atoned with the Atoner, let themselves be reconciled to God. --- The
active element in the appropriation of the atonement is now further to
be developed in the doctrine of

\runhead{c) Redemption.}
\parmark{25}

In the orthodox system of doctrine the two functions, atonement and
redemption, are both referred to Christ's
% ---- printed p.381 (PDF 394) ----
\opage{381} high-priestly office \textit{(munus sacerdotale)} and distinguished
only insofar as atonement \textit{(reconciliatio, καταλλαγὴ)} expressly
denotes the removal of God's wrath, while redemption \textit{(redemtio,
ἀπολύτρωσις)} denotes the ransoming % print: ἀπαλύτρωσις
from the bondage of sin; but to any real insight into the fundamental
relation itself scholasticism, with its definitions, distinctions,
divisions and quaestiones, could not attain. % print: "Distinctioner; Divisioner" (semicolon)
That atonement and redemption are essentially the same means, in other
words, that they are at bottom different: atonement is the possibility
of redemption, and redemption itself the self-actualizing atonement.
Atonement, accomplished once for all, is the foundation of the new life
in faith in Christ. On God's side wrath is lifted, the justice
threatening in the Law transformed into grace. The old, painful darkness
which sin brought into the relation to God has vanished before the new
light; one can therefore begin afresh with new powers on new terms: it
is not the Terrible One of justice on Sinai, but the Crucified One of
love on Golgotha, who makes the new beginning. But in the new relation
to God sin cannot possibly be present as it was in the old; for then
the atonement would have to consist in God's having let himself be
reconciled with sin. If the God who is atoned with sinners is absolutely
irreconcilable toward sin, then the power of sin must be broken so that
the atonement can prove its power. But that the power of sin is actually
broken can only come about in this way, that the Atoner himself frees
the human being from the chains of sin and gives it power victoriously
to wrestle with sin, to break the power of sin within itself; it is this
that happens in redemption.

If the matter is treated scholastically in dead forms of doctrine, it
will be impossible to arrive at true insight into the relation between
what Christ has done for human beings and what human beings, with his
assistance, are themselves to do.
% ---- printed p.382 (PDF 395) ----
\opage{382} To appeal to the Spirit and the cooperation of the Spirit is of no
help at all, so long as the conditions of life on which alone the Spirit
can and will work have not been sufficiently cleared up. It is something
outward, the Atoner, the Redeemer in factual, historical actuality, that
here stands sharply against something inward, the atoning and redeeming
activity that goes on psychically in believing souls; in other words, it
is the essence of faith, the secret of appropriation, on which the
emphasis rests.

As God-man, thus as the object of faith, the historical Christ is the
true, sinless human being, and as the true, sinless human being, by his
suffering and death, the Atoner. This individual, historically existing
human being stands in a two-sided relation to the individuals of the
race: he is qualitatively different from them all, and he is essentially
akin to them all; he is both for one reason: he is the true, sinless
human being. With his qualitative difference from all he stands outside
all, above all; he stands alone. No one can be as he has been, no one
do what he has done, no one accomplish what he has accomplished. In his
essential kinship with all he has stepped into the place of all, has
lived for all, has suffered for all, has died for all. He himself, his
life, his suffering and death can --- on the terms of faith --- be
appropriated by all. Now the understanding of Him can be laid as the
foundation of self-understanding, and self-understanding can have no
other foundation; now the appropriation of Him, his life, his essence,
can be laid as the foundation of free self-development, and truly free
self-development can in all eternity have no other foundation: now the
Atoner can work as Redeemer.

From what are the untrue to be redeemed? From untruth! But the forms
and powers of untruth are manifold. Sin is a power; but, since sin is a
power, death is, after all, also a power; and the powers of sin and
death
% ---- printed p.383 (PDF 396) ----
\opage{383} are, with their spiritual stamp, terrible manifestations of the
superhuman power of darkness, of the lie, of the evil will, of Satan's
power in the realm of spirits. It is the work of redemption that these
powers are broken. We can then, in that the philosophy of religion
keeps to religion's own strong, full-blooded categories, refer the whole
saving, liberating activity to redemption \emph{from sin, from death,
from the kingdom of Satan}.

\greekrun{α}{From Sin.}

What does it mean: to have one's life in oneself, to live out of one's
own human principle? The whole conflict between knowledge and faith may
rightly be said to turn on the answer to this question. Modern science
reproaches religion with placing the essence of man outside man. ``If
the talk is of the atonement, then man is not permitted to seek the
conditions of his life in himself; he is directed to seek them outside
himself in another, half-alien being, in a man who is nevertheless
something other than man, namely in a God-man; on this account man may
naturally not think of expiating his own guilt, freeing himself,
redeeming himself: that would of course be an affront to the God-man''.
Many of the objections belonging here are certainly of scientific
significance, insofar as they are directed against a scholastic-objective
knowledge of religion; % print: "scholatisk"
but they are to be regarded as gross misunderstandings when they are
directed against religion itself.

To come to Christ is, in the religious sense, to come to oneself; he who
comes to Christ in such a way that he thereby in truth comes outside
himself does not come to Christ at all; conversely: he who wants to come
to himself and remain in himself in such a way that he thereby omits to
come to Christ and remain in Christ does not come to himself at all; his
% ---- printed p.384 (PDF 397) ----
\opage{384} remaining in himself as in the absolute is a wretched
self-deception.

To come to Christ is in truth to come to oneself; for Christ, the true
human being, is the truth itself in every human being. The relation
between Christ and the race is, insofar as atonement and redemption are
expressly thought of here, by no means a mystical or mysterious
relation; it is on the contrary a relation of love declared by the
Savior himself in great, clear, unambiguous facts. Why did Christ, the
Righteous One, come to suffer for the unrighteous? Answer: because in
self-sacrificing love he wanted to reveal to the unrighteous the way of
righteousness and, by the revelation of righteousness suffering in love,
to overcome the unrighteous. But thus to overcome the unrighteous, what
is it other than to bring them voluntarily to see, to regret and
bitterly to repent of their unrighteousness, in other words: to bring
them to enter into themselves. If Christ is not perfectly right in this,
then he is unconditionally wrong.

We must not forget, however, how the matter originally and from the
ground up, without any mysterious ambiguity, stands with Christ's
suffering and death. Christ was not, after all, misjudged and persecuted
merely because he proclaimed certain truths offensive to the powerful of
that time; but he was persecuted and seized because he expressly said
that he was the truth itself, thus that he was the truth in the sense
that he is the only absolutely true human being, and that all other
human beings in the whole race over all the earth are, in comparison
with him and in opposition to him, untrue human beings. This his
testimony about himself must now, as anyone can understand, be either
absolutely true or absolutely untrue. So he was persecuted and seized
for the sake of truth --- yes, and for the sake of untruth; but on which
side, his or the race's, is the truth now? Christ was not hated by
% ---- printed p.385 (PDF 398) ----
\opage{385} human beings although he loved them humanly, as if that
generation in which he lived were less loving and more hateful than
other generations; but he was hated because he loved human beings
divinely, because what he understands by love is so infinitely different
from what human beings otherwise understand by love that his love in
their eyes looks like unlovingness, indeed like hatred; he was hated
because he testified of himself that he is the only one who in spirit
and in truth knows what love is, the only loving one, so that all other
human beings in the whole race over all the earth are, in comparison
with him and in opposition to him, unloving human beings, human beings
in whose hearts hatred smolders. This his testimony about himself must,
as anyone can understand, have sprung either from incomprehensible love
or from inexplicable hatred. So he was hated for the sake of love, yes,
and for the sake of unlovingness --- of hatred; but on which side, his
or the race's, is love now? Christ was not condemned and crucified
although he had much right in what he said and did; for if he had only
been relatively right against his age, the age too would on its side
have been relatively right against him: the crucifixion would then have
been a tragic catastrophe according to Hegel's concept of the tragic,
corresponding roughly to the condemnation of Socrates. No, Christ was
condemned and crucified because he came forward as the one who was
unconditionally right. He said that he was the Son of God, and his
enemies understood it as blasphemy. It was blasphemy that brought him to
the cross; but on which side, his or the race's, is the blasphemy now?


When these decisive questions are to be answered---and such questions,
decisive for life and death, are answered only by being answered
% print: "Liv eg Død" (eg for og)
decisively---then the opposition between knowledge and faith is as
conspicuous as it can be. Who
% ---- printed p.386 (PDF 399) ----
\opage{386} is right, then, Christ or Caiaphas? Of course, science is
level-headed; it takes good care not to give a straight answer to such
``fantastically inflammatory'' questions. Objective knowledge, the
metaphysics of the Absolute, will always know how to keep a calm
countenance, even should it quiver with the fanaticism of knowledge; a
calm countenance is, after all, the hallmark of impartiality, and
objective knowledge is, of course, impartial. In its impartiality,
science naturally cannot declare Caiaphas right; for it is evident that
there is egoism in Caiaphas. But on the other hand it cannot, in the name
of the concept, simply declare Christ right either. Certainly the ``hero
of Gethsemane'' is not the greatest egoist who has lived on earth; that is
an exaggeration. But there has nonetheless been some egoism in him, a
small remnant at least. For there is a remnant of egoism in every human
being who has not ``become absolutely transparent to himself in the
Absolute''. Christ, who stands at the standpoint of representational
consciousness, has not been able to raise himself to the concept of the
Absolute, consequently has not been able to become transparent to
himself, consequently has not been able to consume the egoism remaining
in him, consequently cannot be unconditionally right. In other words:
Christ, as ``the single individual'', must put up with being reduced to a
relativity. --- But if Christ is a relativity, then he is condemned, for
he has laid claim to being the Absolute, and a relativity that lays claim
to being the Absolute is to be condemned without grace or mercy. So
Caiaphas is right after all, when from his standpoint he calls Christ's
testimony about himself a blasphemy, and metaphysics might just as well
say it outright, without any beating about the bush: Caiaphas is right,
unconditionally right!

What is sin? Self-contradiction, the deepest, most painful, most
spirit-consuming contradiction that can be thought. This
self-contradiction has its origin, not in thought, but in the will; it is
not thought
% ---- printed p.387 (PDF 400) ----
\opage{387} that began by disturbing the will; no, it is conversely the will
that began by confusing thought. The first thing in self-understanding is
not: to will to think, but is: thinkingly to will one's eternal self.
Thinkingly to will one's eternal self is: thinkingly to will, in the name
of truth, righteousness and love, to declare Christ right and oneself
wrong; it is: thinkingly to will what Christ wills. For what does Christ
will? That man shall come to himself by coming to him, shall have the
self-contradiction resolved by letting him resolve it, shall in personal
understanding exchange roles with him, so that he may be allowed to bear
man's sin and man let himself be clothed in his righteousness. Through
unrighteousness man has gone away from himself; through the appropriated
righteousness he comes to himself again, and the self-contradiction is
resolved.

The relation that is thus brought about personally between Christ and the
individual presupposes a universal relation of essence between Christ and
the race. Just as the Adam-essence is the essence of the race in its
sinfulness, so the Christ-essence is now the essence that is sinless for
the race, which each individual is to appropriate in faith. ``It is not
only the race that is represented in those who stand beneath the Cross;
but it is also the race that is represented in him who hangs upon the
Cross. It is the old Adam who is beneath the Cross; it is the new Adam
who is upon the Cross. The second Adam, indeed, has no part in the sin of
the race; but his self-consciousness is inseparable from his unity with
the sinful race. The suffering and death that he suffers from the race
he himself transforms into a suffering and death for the race; for in
the infinite compassion of his love he is able to feel the guilt of his
brethren as his own, able to bear the guilt of the race upon his
high-priestly heart. So inwardly is he, through his becoming man, bound
up with the sinful race that Christ, as Luther expresses it, can say:
% ---- printed p.388 (PDF 401) ----
\opage{388} I am this sinner, i.e.\ his guilt and punishment belong to me; and
that the sinner through faith can say: I am Christ, i.e.\ his death and
righteousness belong to me''\footnote{H.\ Martensen. ``Den christl.\
Dogmatik''. S.\ 317--18.}.

In determining the relation between Christ and the race it must above
all not be overlooked that this relation of essence is, at its innermost
ground, a relation of will. Religion presupposes that everyone who in
sincerity \emph{will} understand the Atoner also \emph{can} understand
him; in which it is implied that he who says that he cannot understand
him says so only because he will not understand him, and will not do so
because he will not understand himself in spirit and truth. Not to will
to understand the Crucified and declare him right---that is: to will the
self-contradiction; that is: to will sin as sin, i.e.\ to will one's own
perdition. To will to understand the Crucified and declare him right---that
is: to will to recognize the truth in him as truth itself, the
righteousness in him as righteousness itself, the love in him as love
itself; that is: to will one's own salvation. The personal account
between Christ and the race is something other and more than a system of
universal Ideas. It is not the Idea of truth in general, not
righteousness and love in general, that are in question here; it is
precisely the fundamental religious-ethical error of reason that it has
had the eternal Ideas of the will transformed into universal potencies of
selfhood, and the willing self into a ``finite subject''; and so it
ends in --- perdition. Nor does Christ, as an externally existing
individual, will to bind the race to his finite phenomenon. The Crucified
has in his individuality the eternal universal, has the Ideas of the
will and the will itself in living, visible unity. It is the Ideas in
flesh and blood, the suffering truth, the scorned righteousness, the
thorn-
% ---- printed p.389 (PDF 402) ----
\opage{389} crowned love, through which the holy will on Golgotha set itself
against the will of the race, that each individual might go into himself,
condemn sin in himself, and let himself be saved.

By condemning sin in himself the believer is raised above sin, receives
through Christ a new strength of will, a new will, a new self, and
becomes in faith a new man. With Christ's atonement as his inner
presupposition, the anchoring ground of his life, the new man is saved
from sin and to that extent without sin; but in natural self-unity with
the old man he is nonetheless still encumbered with sin. The
self-contradiction that man is both saved from sin and yet bound by sin
resolves itself into a struggle which, by virtue of redemption, leads to
victory over sin. This victory, by the nature of the struggle, has an
absolute and a relative side; the absolute rests on faith as faith, the
relative on the strength of faith with which the atonement is
appropriated and redemption actualized. We see this in the doctrine of
\emph{the following of Christ}.

In the following, the life of the Redeemer is to repeat itself in the
redeemed. But the following is no outward imitation. In outward
imitation the one who strives places himself outwardly over against
Christ, compares his works with Christ's works, his sufferings with
Christ's sufferings, and thereby comes to be outside himself. It is an
error which, for that matter, can take on a very different character---one
need only think of the difference between a holy Francis with his
childlike, rich mysticism and a sour-tempered, spiritless
Pietist---according to the unclarity within from which it proceeds. The
fundamental error is that the one thus inflamed wants to copy Christ's
works without understanding Christ's essence, without deeper human
self-understanding. Instead of placing the following in the inward
appropriation of Christ's life, such a follower abstracts for himself a
new law from Christ's
% ---- printed p.390 (PDF 403) ----
\opage{390} example; sighing and groaning under the yoke of the new
self-made law he labors with sin, against sin, and cannot possibly
overcome sin, since in his darkening of spirit he does not see what sin
is. Understood absolutely, there is only one sin: not to believe
unconditionally in one's redemption. Lack of faith is lack of inner life
in light and warmth. Every outward break with the outward world, every
ascetic exertion of strength, be it ever so heaven-storming, is without
religious significance when it does not proceed from the new self in
living faith. Against all finite penances, however painful, sin can hold
its own; only against the life of faith in personal self-understanding
can it not hold its own; before the power of faith it succumbs: penances
are piecework, but in faith \textit{(fides salvifica)} is salvation.

\greekrun{β}{From Death.}

In the measure that the Redeemer overcomes sin in the redeemed, he also
overcomes death; for when sin vanishes, death is a nothing. That death
came into the world through sin means, as we have seen, that death first
became something through sin, became the dark power that signifies the
curse. It is fear of death that makes death fearful. From the standpoint
of religion, fear is originally to be determined as fear of conscience.
For just as sin can be in a man without the man having yet become
conscious of it as sin, so too the fear of death can dwell hidden in a
man who nonetheless feels no fear, because the fear of conscience is
either not awakened at all or only intimated in inexplicable anxiety.
Just as little as expressions of natural carelessness and selfish
boldness are signs that sin is lifted---since such expressions are
precisely signs of the opposite---just so little is a contempt for death
springing from sensuous courage a sign that the original
% ---- printed p.391 (PDF 404) ----
\opage{391} fear of death is lifted, since it is on the contrary an
unfailing sign of the opposite. Just as there is a difference between
sensuous and spiritual courage, so there is also a difference between
sensuous and spiritual fear of death. If selfish contempt for death were
a testimony to true courage, then one would surely have to admit that the
coarser and more defiant a man is, the more courageous he can be: the
chieftain of the savage tribe from the primeval forests would then
perhaps be more courageous than any civilized general. This measure of
courage the Redeemer has broken; he broke it in Gethsemane. Without
spiritual fear of death---how should he, who is himself the Life, fear
death?---he was nonetheless, in the conflict of his sufferings, as one who
in anguish of soul strove with death, moved in fear and weak as a woman.
But this weakness, this trembling in the anguish of death, is in a deeper
sense precisely the revelation of his invincible strength of soul. The
innermost terror in the Savior's struggle of soul is his feeling of sin,
not of his own, for he is sinless, but of the sin of the race, of the
misdeed in which it was now to become manifest. The God-man too has had
to feel the overwhelming power of a moment of horrors, and the pains that
seized him show clearly enough that he was wholly man, that there was
nothing hidden in his essence that could exempt him from actual
sufferings. It should then also be emphasized that it is not the
sensuous sufferings themselves, but the representation of them, with
which he struggles; for precisely by taking up the sensuous sufferings
beforehand in representation he transforms them into sufferings of the
soul and conquers them in the soul. When the struggle of the soul had
ended in unconditional surrender to the Father's will, the anguish of
soul was also past; the one who had struggled rises strengthened to bear
all things, and no hero on earth who ever went to meet the terrible has
been so prepared as he was when, fully resolved, he stepped forward and
asked the guard: whom seek ye? From
% ---- printed p.392 (PDF 405) ----
\opage{392} the moment when the history of the Passion begins in tangible form,
there is in his being no trace of fear left: no wavering, no
faintheartedness, no downcast look; only the divine superiority of
self-certainty, meekness, and patience is to be seen in the countenance
of the conqueror of death.

So he died and was buried; but his actual death was a means of making
manifest the impotence of death; for he rose from the dead. ``How? Under
what conditions? With what body''? Of that nothing \emph{is known}; it is
the miracle of revelation, the sign of omnipotence, actuality of spirit,
actuality for the spirit, and there it rests. But that he actually rose
and was actually seen by his witnesses---of that his redeemed have in
themselves, in self-understanding, so powerful and unfailing a
testimony that existence would have to become meaningless, and the world
a madhouse, if this testimony could deceive. Christ's resurrection stands
not in an accidental but in an essential, indeed in the most essential
and most inward relation to his whole personal life: he was not merely
to overcome death for his own person; he, the revealed truth, was indeed
also to reveal himself as the conqueror of death.

The redeemed has in this revelation all the presuppositions and
conditions for overcoming death in himself. Temporal death has lost its
sting; for the relation to death changes with the relation to sin;
eternal death has lost its sting, for the Redeemer did not die as judge;
no, he died with forgiveness on his lips, he died as Atoner. Here, as
regards the fear of death, there can no longer be talk of an essential
difference between the courageous and the fearful; for the fear of
death is, on the terms of faith, essentially fear of conscience: this
fear is equal for all and is sublated equally for all. Here no
distinction can any longer be made between those strong by nature and
those weak by nature; for He who
% ---- printed p.393 (PDF 406) ----
\opage{393} is stronger than the strongest became weak as the weakest, that
in the weakness which is strength he might teach both weak and strong to
overcome death.

\greekrun{γ}{From the Kingdom of Satan.}

``For this purpose the Son of God was manifested, that he might destroy
the works of the devil'' (John 3:8).
% print: Joh. 3, 8; probably meant 1 John 3:8
This and similar sayings in Holy Scripture gave the old Church teachers
occasion for a dramatic depiction of a formal wrestling-match between
Christ and the devil. ``Through the first sin man had come under the
dominion of Satan, the prince of sin and death; this was indeed a
wrongful act of violence on the devil's part, which the Son of God would
to that extent have had the right to annul by force; but it accorded
with his perfection to proceed justly even against the unjust one:
therefore he did not tear men from him by violence, but as by a peaceful
exchange, redeeming them from him by giving up his soul as a ransom.
Therefore he had to be both God and man: man, because he was to redeem
men by his death; God, in order to be able to withstand the devil. But
since, precisely for this reason, the devil did not keep Christ's soul,
as it could not, according to Ps.\ 16:10, Acts 2:27, 31, be held fast by
the underworld and death: how then was the devil obliged to release
mankind for a ransom which he had only apparently received? It was
precisely in this that the devil deceived himself. He came to an
agreement with God, in the best of form, that in return for the soul of
Jesus, which as the soul of the God-man was worth as much as, indeed
much more than, the souls of all men together, he would set these free;
the soul of Jesus was handed over to him: that he could not hold it fast,
since he could not endure the pain this caused him, could not make the
contract void, since he ought to have considered this beforehand. But
how could the devil
% ---- printed p.394 (PDF 407) ----
\opage{394} deceive himself so grossly? This became possible through the
human veil in which the Son of God had wrapped himself, and that with the
intention that the devil, over the easily subdued humanity, should forget
the almighty Godhead indwelling it, and so enter into the bargain all the
more credulously. When the devil laid hands on Jesus, tormented and
killed him, although he had no right to him as sinless, he exceeded his
authority and lost those whom he had until then held in his power with a
certain right. Like a clever fisherman God hid Christ's higher nature as
a hook under the bait of the flesh; but the devil, like a greedy fish,
swallowed both the one and the other. Thus was fulfilled the word which
God had once spoken to Job (Job 40:19 f.): Canst thou draw out Leviathan
with a hook? and does one pierce the nose of Behemoth? A man, namely,
% print: pan for paa ("Næsen pan Behemoth")
does not indeed do this, but Christ has accomplished it. But that
unlawful bite agreed so ill with the old Leviathan that, like Saturn
once, he had to give up all those whom he had previously
swallowed''\footnote{Strauss Troesl.\ II.\ S.\ 203--206}.

To what degree the fundamental religious representations can be
distorted when self-understanding does not go hand in hand with the
understanding of the word of revelation can be seen from the description
cited. In this description the devil has become just as external to the
view as Christ; and as soon as fantasy had brought forward the external
material, the understanding at once began to work it up with juridical
acuteness. Thus the doctrine gradually becomes objective and scientific.
But when the devil himself has got into science, what wonder, then, that
critique sets itself in motion to drive him out. But in this expulsion
there is again something devilish, for now the critique of knowledge
regards its victory over the devil as a proof of
% ---- printed p.395 (PDF 408) ----
\opage{395} the complete annihilation of the fundamental religious
representation, and thus as a sign of the tottering condition of
religion. ``Christ and the devil, these two potentates, stand in so
decisive a relation of reciprocity that when the one falls, the other
falls too.''

This misunderstanding is removed when, by virtue of the principle of
faith, we hold fast the point of view that the doctrine of the
\emph{evil} will too, and, what corresponds to it, the doctrine of the
power of the devil, is to be judged solely and exclusively by the
influence it exerts on the cognition of the life of the will as a whole,
and then in particular on self-cognition.

\emph{Will is moved by will}; this proposition is decisive for all
development of character. The more practical the development is, the
more clearly it shows itself that what affects the individual will from
without is chiefly another will, the will of others. In the relation
between the one who commands and the one who obeys this is evident; but
it is so too in persuasion, admonition, temptation, seduction, etc. Every
human society is, though in a very different spirit, a society of wills,
a system of wills in manifold interaction. That the social will, the
general will, can just as well be said to determine the individual wills
as in turn to be determined by them has already been pointed out in an
earlier context. Among the individual wills themselves the interaction
is finite; but between the individuals and the general will it is
infinite, and the general will, whether it be otherwise good or evil, has
the overreaching power.

If we now examine in what the \emph{evil} general will properly consists,
and by what it is known, then its evil does not lie primarily in its
making the individuals will-less; for in a society of will-less
individuals no general will can exist. If we say: the evil general will
must be known by its making the individual wills evil, then this is only
a formal paraphrase; for the question
% ---- printed p.396 (PDF 409) ----
\opage{396} is precisely in what the evil consists, and how it comes about.
If it is the good of each individual will that, regardless of what it
determines itself to in a finite sense, it has the law of
self-determination in itself and can act from its own motives, then the
evil must surely show itself in this: that the general will robs the
individual wills of their original right of self-determination; which
can also be expressed thus: that it filches the individuals' right to
themselves, steals their conscience, and makes them into blind, but
strong-willed, instruments. This cunning of the general will is of so
sly, serpentine, spiritual a nature that it takes hold only where
spiritual powers are at work; but then also in such a way that the
higher the stage of spirit to which the society has raised itself, the
more terrible is the evil. At the stage of crudeness the evil in this
form does not appear; but in the course of cultural development it has
already manifested itself in pagan fashion in the ancient states. In
Sparta and Rome, as is well known, the true citizens durst have no
conscience for themselves. The common conscience was in the keeping of
the state. What in itself was right or wrong, religious or
irreligious, which worship was the true and which the untrue---that the
individual, however enlightened he might otherwise be, could not possibly
decide on his own; it had to be decided objectively, in the society. From
this it is explicable that even the noblest natures could, in the name of
the state, commit the most revolting acts without reproaching
themselves. And yet the crime of the spirit does not yet attain its
proper form; for a crime of the spirit presupposes spirit, and paganism is
not essentially spirit. But in the Middle Ages, on the foundation of the
Church, the crime of the spirit could take a leap upward. Seen from its
shadow side, the papal-hierarchical general will has precisely this
peculiarity about it, that like a vampire it sucks the blood of
consciences and makes individuals into blind instruments. This is seen in
Jesuitism; for Jesuitism is
% ---- printed p.397 (PDF 410) ----
\opage{397} the darkness in the Catholic spirit. The holy becomes egoism,
egoism holy; obedience to God is transformed into blind obedience to
superiors, who in turn show a blind obedience to the general will
prevailing in the society. With an exactitude extending to the very
smallest things, the one must confess for the other and answer for the
other; but no individual, the highest as little as the lowest, is
originally answerable to God, for the selfishly holy general will has put
itself in the place of the will of God; by taking the whole
responsibility upon itself it has made the individuals unanswerable
before God.


Now evil has been brought into a system. In this thoroughly reflected
life of the will every crime is a crime of the spirit. And no crime,
not even one that could turn the misanthrope's blood to ice and make
the hair stand on end on the robber's head, is so revolting that any
one of those consecrated to the service of the universal will could
not carry it out calmly and in full consciousness; and why? Because
the consecrated one has only to obey, because he is without
responsibility before God, because the evil will of the universal
spirit has stolen his conscience from him by guile.

But how has it actually come about that evil could establish itself
at such a depth, and with so inhuman a rigor of consistency, in an
unconditionally ruling universal will? No one knows; there is as yet
no one who has been able to fathom it right to the bottom --- here
the path of psychology to the finite runs through the infinite ---
but everyone who, from the standpoint of civilization and humanity,
has gained a deeper insight into the thoroughly worked-out systematics
of this universal will, down to its finest ramifications, which twine
about all the roots of the mind, must involuntarily, whether he is
otherwise a believer or not, exclaim to himself: it is a
\emph{kingdom of Satan}.

But, someone says, if there really exists a kingdom of Satan, then
there must also exist an actual Satan;
% ---- printed p.398 (PDF 411) ----
\opage{398} one must then be able to catch sight of him objectively somewhere.
This is a misunderstanding. The satanic essence can be observed; but
Satan himself, even though science peers through all the telescopes
of critique, no, Satan himself is not to be descried. Yet let us try
whether there is not perhaps another eye, one organized to see what no
knower can see; that would then have to be the eye of faith.

That Christ has redeemed his believers from the kingdom of Satan can,
to be sure, have no essential validity if there is no Satan; but then
neither can the redemption have validity outside the one who, in
believing self-cognition, cognizes himself as redeemed. But in order
truly to feel himself redeemed from the kingdom of Satan, the redeemed
one must surely once have been in this kingdom himself, and, after
having been in the power of the evil will, have experienced what it is
to break with this power. So let us fix our attention on such a
figure. He was a Grand Inquisitor. He sat in the secret court, a
witness to works which we must, in more than one sense, call with
horror works of darkness. Then it happened once that one of those
unhappy victims, whose patience was greater than all the pains of
torture, sent him a glance, a glance so indescribably piercing and yet
so gentle that he could never forget it. Quietly, shut up within
himself, he now began an inquisition of another kind. It was not
heretics he was examining; he felt himself too sickly for that, he
said, and his sickliness and exhaustion then had to serve him as an
excuse with those around him: he was examining himself. It became
clear to him, as he --- alone before God, not on behalf of the Order
but for his own conscience's sake --- immersed himself in the Gospel,
that the history of Christ's passion contained the most terrible
judgment upon the works of the Inquisition. What, then, had he lived
for? A kingdom of Satan! ``Did he now see Satan himself in this
kingdom''? No! But in the anguish of his soul he nevertheless
perceived that Satan is a will.

% ---- printed p.399 (PDF 412) ----
\opage{399} ``For this purpose the Son of God was revealed, that he might destroy
the works of the devil''.

\lettersub{C.}{The Kingdom of the Son.}
\parmark{26}

If we set the revelation in the Old Testament and the revelation in
the New Testament sharply against each other, the difference is
striking. The doctrine of Jehovah against the doctrine of Christ: it
is as shadow against body, as law against gospel, as compulsion
against freedom. Were the old theocracy with its priests and Levites
to constitute in particular the kingdom of the Father, while the new
dynasty of the spirit, established by Christ, were to be regarded as a
kingdom by itself, standing exclusively under the Son: then there
would undeniably appear here a questionable separation between the
kingdom of the Father and that of the Son. But that is not how
matters stand. From an external historical point of view, Jehovah and
Christ can by no means be said to be related as the Father to the Son.
It is a Jewish misunderstanding of the messianic ideal that the
Messiah should be able to become the Son of God, although God can
nevertheless not possibly be his heavenly Father; for as long as God
is Jehovah, he can never become Father. Jehovah, put plainly, has no
children. His children are his creatures; but the Creator can only
figuratively be said to be related to his creatures as a father to
his children. The expected Messiah, whether he is otherwise seen in
the clouds of heaven among the thousands of angels or on Mount Zion
with the iron scepter that chastises the nations, is only in a
spiritual image to be seen as the Son of God: the actual is that he is
the Son of David.

% ---- printed p.400 (PDF 413) ----
\opage{400} This must then also exercise an essential influence on the judgment of
the prophecies and their fulfillment. That the prophecies have found
their fulfillment in Jesus Christ coincides, in concept, with Jesus
Christ's having actualized the true messianic ideal; but to form from
the prophecies of the prophets an image of the Messiah that should in
each and every respect resemble the actual Christ is impossible.
Standing on the foundation of the theocracy, the prophets were
constantly bound to apply the measure of the law and of outward
dominion to the Expected One. That it was the spirit, ``the Spirit of
Jehovah'', ``the spirit of wisdom and understanding'', ``the spirit
of counsel and might'', on which everything first and last depended,
they saw clearly; but with regard to the spirit there was a secret
which they could not possibly see through at the then stage of
revelation, and that was the mysterious unity of the spirit with the
Word of life in the holy ground of nature, the unity that makes the
spirit the spirit of love and the true Messiah the Son of God. The
new life that comes with Christ, the new beginning that takes place
through him, is then also decisive to such a degree, and works with
such originality upon the believing witnesses, that these can
evidently not be said to direct their conception of Christ
immediately by the prophecies, but conversely to conceive and
interpret the prophecies according to their faith in Christ. Proofs
of this we have, e.g., in the use the Evangelist Matthew makes in
1:23 of Isa.\ 8:8--10, and in 2:15 of Hos.\ 11:1.

Only when the old covenant in its entirety is regarded as a
preparation for the new, only when Judaism in broad outline is
conceived as a prophecy of Christianity, can the individual prophecies
be placed in the right light; but then they are also testimonies to
the stepwise transformation of the revelation of Jehovah. Jehovah is
the Father and yet another than the Father, just as Jesus of Nazareth
is the Messiah announced by the prophets and yet another. It
% ---- printed p.401 (PDF 414) ----
\opage{401} is not God as God who has changed his essence; the change is in the
progressive revelation. From the Old to the New Testament there is no
rational transition; for the Jew learned in the Scriptures Christ is
just as surely a ``sign that is spoken against'', to offense, as for
the philosophical heathen. The unconditionally new lies in the
revelation of the Logos, the revelation of the Word of life; with the
Word of life Jehovah is transfigured into \emph{the Father} and the
Old Testament Messiah into \emph{the Son}. Now the relation is not
figurative but actual; God himself is the Father, he ``of whom all
fatherhood in heaven and on earth is named''.

As the Son is one with the Father, so the kingdom of the Son is also
one with the kingdom of the Father; it is in the unity of this kingdom
of the Father and of the Son that the old kingdom of Jehovah
dissolves, and the theocracy with its laws of compulsion perishes. In
the doctrine of the Father we have therefore also regarded this
kingdom from the point of view of the ideals, and under this
consideration adduced characteristic features of the preparations in
the kingdom of Jehovah; from the point of view of the Son we shall now
consider the actualization of the ideals, how they take on life and
shape in visible forms, as we conceive the kingdom of the Son as the
kingdom of \emph{freedom, love} and \emph{blessedness}.

\runhead{a) The Kingdom of Freedom.}
\parmark{27}

From the historical point of view of the development of the spirit,
the opposition between the law in the Old and the gospel in the New
Testament is to be regarded as an opposition between two ages of
life, the first of which designates the stage of minority, the latter
that of majority. At the stage
% ---- printed p.402 (PDF 415) ----
\opage{402} of minority the point is discipline, at that of majority
self-education; discipline presupposes laws of compulsion,
self-education laws of freedom. This opposition points to a two-sided
relation between nature and spirit. There is in man a nature that
strives against the spirit, and there is in his fundamental essence a
nature that feeds on spirit and nourishes the spirit.
What is characteristic of discipline is now this, that here the
beginning is made in discord between nature and spirit. Although the
spirit originally aims at freedom and works liberatingly by subduing
willfulness and awakening the self to spiritual self-understanding,
it must nevertheless impress its laws with divine authority upon the
sensuously selfish individuals, terrify them with threats, and employ
compulsion. It is otherwise when self-education is to begin; for the
fundamental condition is that compulsion falls away. That compulsion
falls away does not mean that the law lapses; it means, on the
contrary, that it is moved from the outer into the inner. It means
that the individual is to have the law for his development within
himself, is to compel himself; but to compel oneself is in earnest to
determine oneself according to the law of freedom: in
self-determination the law is acknowledged, but compulsion is
sublated. With this the discord between the spirit and nature is also
sublated, namely in principle.

Self-education thus begins, under the influences of the spirit,
essentially with this, that the individual begins with himself from
within, begins by raising the higher nature to dominion over the
lower. Here, to be sure, the beginning is made with a demand of the
spirit, but this is originally one with a corresponding impulse of
nature: man satisfies the spirit by satisfying himself in a higher
sense. In order that the satisfaction may be true and deep, and the
better nature of the individuality may receive its due, every notion
of finite compulsion must be kept away.

To these conditions of freedom corresponds precisely the beginning
% ---- printed p.403 (PDF 416) ----
\opage{403} that Christ makes in founding his kingdom of heaven. ``The Kingdom of
God is within you'': this word is decisive for the beginning from
within, the free beginning toward freedom, for freedom begins only
from itself. Of the many parables under which the character and
essence of the kingdom are depicted from various sides, we shall in
this connection direct attention especially to those which, from the
spiritual side of nature, make visible the fundamental conditions,
the terms, and the burdens of freedom.

\greekrun{α}{The Fundamental Conditions of Freedom: The Sower and the Soil.}
% print: „Jordbnnden“ for Jordbunden in the heading

When, in the individual life of freedom, we distinguish what takes
place in the daylight of consciousness from what stirs darkly in the
mind and belongs, so to speak, to the night side of the soul, we must
of course not forget that the unity is the psychic self, that the
original unity of the self is a unity of will, that the individual
arbitrary acts of will get their significance precisely as free
shoots that shoot forth from the fundamental will. The life of the
will too has its nature side; the will works as involuntarily in the
ground of the mind as either imagination or feeling or thought: the
spiritual life has its vegetative side in the soul, just as the soul
has in the body. If, then, the life of freedom is to be spoken of
vividly and in a living way, it lies in the nature of the case that
now the one, now the other of its two regions is especially brought
out. This has then also been done in the parables of the kingdom of
heaven.

It lies in the essence of the parable that, as a figurative
presentation of the truth, it can never clarify the whole truth from
all sides, but can instead illuminate a single main side all the more
powerfully. The different parables, when outwardly put side by side,
can easily give the superficial observer the impression that they
stand in conflict with one another. Thus we have parables of Christ
% ---- printed p.404 (PDF 417) ----
\opage{404} in which the self-conscious exertions of the will in watching and
praying are enjoined in the very strongest terms, while others of his
parables bring out the nature side of the life of freedom to such a
degree that at first glance it looks as though the word could sprout
in the mind like the seed in the field. These different depictions of
the same truth are, however, so far from contradicting one another
that the one is, on the contrary, a completion of the other: the
different parables supplement one another mutually. If it is
one-sided to hold exclusively to those which depict the life of faith
as a growth, it is also one-sided to hold exclusively to those which
can support the ascetic conceit that the new life can only be a
product of reflected self-exertion.

That the will of the spirit is the fundamental will of the
individuality presupposes not only that the will has its foundation in
the individual nature, but presupposes also that the inner free
working of the spirit begins only from the ground up, in that it
begins, in an obscure and mysterious manner, as growth. What is
signified in exemplary fashion by the Christ child's growing with a
holy growth comes to full visible expression in the doctrine of the
kingdom of heaven.

This is seen from the parable of the sower (Matt.\ 13; Mark 4; Luke 8).
There is in this image something so popular and free and fresh as
spring that no one who truly immerses himself in it can doubt that
everything here takes place under God's open heaven. It is with a
divine open-handedness that this sower strews the seed everywhere,
without pettily troubling himself about what the actual soil is like;
so it happens that not all the seed falls on good ground, but that
some also falls by the wayside, and some on stony ground, and some
among thorns. Thus the parable fits literally the word of freedom,
Christ, the divinely liberal-minded sower, who strews the seed of the
word everywhere in the soil of the people,
% ---- printed p.405 (PDF 418) ----
\opage{405} leaving it to each individual in particular how he, according to his
psychic nature, will receive the gift rich in promise. And how
various, then, is not this soil of the people! Immediately after the
parable one might almost be tempted to assume that the psychic
differences of individuals rested to such a degree on a natural
determinateness independent of the will that freedom was lost and the
recipients became wholly unaccountable. ``What can the wayside do
about the birds eating the seed that it is unable to keep? What can
the stony ground do about its being so hard, or the thorny ground
about its thorns finally choking the seed? Indeed, what merit has
the good ground in its fruitfulness? It is as it is, and surely
cannot be otherwise''. With such reflective considerations the
parable, according to its character, cannot possibly engage. Its
thoughts are thoughts of impression, visible thoughts of life. But
precisely by halting the soaring reasoning, it fastens attention with
irresistible definiteness on a single great, decisive point of view.
What, then, do we learn from the parable? That the life of freedom
must begin in the mind as a seed corn begins its growth in the soil;
that the spiritual vegetation, in other words, is the original
presupposition of freedom, and the ground of individuality necessary
for the seed of the word the indispensable foundation of the life of
freedom.

Most closely connected with the parable of the sower is the parable of
the tares among the wheat (Matt.\ 13). Here too the fundamental
conditions of freedom, of inner, naturally determined
self-development, are held fast. The seed is sown in the field, the
good seed; but the field is not protected by any outward measure
whatever. The people sleep, and the householder has not thought of any
safety watch that should keep guard day and night lest other sowers
with seed of another kind should come to sow in the same soil. So the
opportunity is used,
% ---- printed p.406 (PDF 419) ----
\opage{406} so the enemy comes and sows his tares among the wheat. Had it not been
for freedom's sake, such a thing would surely have been prevented.
Now when the crop grew up and bore fruit, and the tares too showed
themselves, it seemed to the householder's servants that if nothing
had been done before, then something \emph{ought to be} done now.
In their opinion measures ought at least to be taken to stop the
growth of evil; they therefore ask the householder whether he does
not think they should go and weed out the tares. But the householder
says no; and why? It is certainly not from indifference toward the
wheat that he spares the tares; it is for the wheat's sake, in other
words: for freedom's sake, so that the root of life shall not be
harmed. So both grow together, the tares with the wheat; it looks
like a kind of immediate natural right. This natural right, however,
is in the sense of the spirit so far from harming freedom that it
conditions and furthers self-development instead. There is a
well-meaning watchfulness, a self-important solicitude, that would so
gladly shield the dear individualities against harsh attacks, cover
them against every gust of wind, and, if it were possible, transform
the tree of freedom into a potted plant: that harms freedom. Outward
solicitude does double harm; it makes the individual soft, and it
dulls his consciousness of responsibility. The householder in the
parable upholds the fundamental conditions of freedom in both
respects; he lets the good work its way forward in struggle with
evil; he makes responsibility count by pointing to what is coming and
awakening consciousness of the judgment. So he says, and says it so
loudly that both the tares and the wheat can hear it: when the harvest
comes, gather first the tares together and bind them in bundles to
burn them, but gather the wheat into my barn.
% ---- printed p.407 (PDF 420) ----

\opage{407} \greekrun{β}{The Terms of Freedom: The Prodigal Son.}

Freedom is a dangerous thing; that holds not only outwardly but also,
and chiefly, inwardly. It is the fundamental condition of
self-development that he who is called to freedom must begin by being
his own master: so it begins, then, with immediate headstrongness.
That a tolerably ordered civil society must consider well what it is
doing before it too hastily lets headstrongness loose is natural. But the
kingdom of Christ is also something other than a political civil
society. In civil society the point is merely freedom in outward
things, outward freedom; and it then goes without saying that outward
freedom must have its outward limits. In the kingdom of the Son, on
the other hand, the point is inward, solely and exclusively inward
freedom; but inward freedom cannot begin originally if it does not
begin naturally, and cannot begin naturally if it is not permitted to
begin in headstrongness.

Incipient headstrongness is dangerous; in the rush of passion it can be
like a daring leap between cliffs and abysses, with the foolhardiness
of the inexperienced like a voyage in an open boat on the wild sea;
it can therefore, in purely personal relations, in the relation
between guardian and ward, between father and son, often be a matter
of misgiving, indeed of pain, for the experienced to let go of the
inexperienced and allow him in the name of freedom to begin his
course as he will --- in all headstrongness. It is in this respect ---
and not only with regard to repentance and forgiveness of sins ---
that the parable of the prodigal son is so instructive. It is with
exalted calm, but certainly not without dark forebodings and deep
pain, that the father ``divides his living'' between the two sons.
But the age of majority had surely been reached by the younger too.
Here no objection helped, and alas, unfortunately, no admonition
either. So
% ---- printed p.408 (PDF 421) ----
\opage{408} it was a silent parting, but by no means an easy one on the father's
side, when ``not many days after the younger son gathered all he had
together and took his journey into a far country''. What road he took,
and how it went with him, we know. It is clear enough that there is a
double nature in this son, a lower one that storms ahead selfishly and
misuses freedom, and a higher one that can only break through when the
lower has raged itself out. It is, people say, the bitter experiences
that subdue headstrongness and produce repentance. Yet it is not quite so.
The bitter experiences could in the prodigal son also have produced
another effect, namely embitterment against the father. ``Why did he
fulfill my wish so readily? He divided the living at once and gave me
my portion, just as if he were quietly glad to be rid of me. After
all, he knew my youthful passion, my hot blood, my inexperience; how
could he, the experienced one, so calmly put up with everything and
hold out his hand in farewell?'' Among all the crises of freedom none
is more dangerous than this: when everything has been staked and
everything is lost, and only one thing remains, the gnawing worm in
the mind, then in self-understanding to have to cognize whose the
guilt is. Just as the first step of headstrongness is the natural beginning
of freedom, so in a deeper sense repentance is the spiritual beginning
of freedom; religious freedom begins with repentance.

In an outward sense, and judging by the measure of a civil-moral
order, the elder son has a great advantage over the younger. He too
could, after all, have done the same as his brother, taken his share
and left the father. It was handsomely done of him that he did not do
it, but stayed with the old man and helped him. He is industrious, a
friend of order, an enemy of excesses, and one cannot blame any civil
society for preferring such a man of honor to his dissolute brother.
% ---- printed p.409 (PDF 422) ----


\opage{409} One might indeed, from the standpoint of the moral law, subject the
father's conduct as a whole to a rather sharp critique, since one would
then have to find that he is at bottom a weak man who lets his feelings
run away with him. Not to speak of his far too great indulgence in
humoring the young man and letting him travel on his own without even
thinking to send a tutor along with him, he ought nevertheless, when he
saw his son come back as a \textit{mauvais sujet}, at least to have given
him a proper chastisement, to have humbled him quite perceptibly in the
sight of his elder, exemplary brother, so that he should not soon again
feel like going off on adventures with freedom. But what does the old man
do instead? He forgets, you see, his paternal dignity, he quite forgets
moral seriousness and gives a feast --- for the lost one! ``Bring forth
the best robe, and put it on him; and put a ring on his hand, and shoes on
his feet; and bring hither the fatted calf, and kill it; and let us eat,
and be merry.'' Yes, such is fatherly love, when it winks at the
wilfulness of youth and makes common cause with freedom.

How anxious, how deeply grieved this father's heart must have been at the
thought of the son who had gone away, perhaps lost forever, can be seen
from the joy with which he receives him again, for he receives him with
open arms. He lost him for freedom's sake, but now he has won him back,
won him in the name of freedom, won him in such a way that he belongs to
him forever.

This the morality of reason cannot understand, and the man of duty, the
elder brother, cannot understand it either. He had just come from the
field, from the fulfilment of duties; ``and as he came and drew nigh to
the house, he heard music and dancing. And he called one of the servants,
and asked
% ---- printed p.410 (PDF 423) ----
\opage{410} what these things meant; and when he then was told what it was,
he was angry and would not go in.''

It is a peculiar kingdom, this kingdom of freedom, where wilfulness can
run rife without any compulsion being applied. Yet this is by no means
because the king of this kingdom, in logical unclarity and for want of a
rational ethics, has confused freedom with wilfulness, so that he is
unable to make the distinction. No, it is because he wants wilfulness,
selfish self-will, rooted out from the ground up. That this cannot happen
on the ordinary terms either of household discipline or of civil order we
learn from the example of the elder brother. This orderly man has surely
never gone astray, has never, so far as anyone knows, departed from the
path of virtue; but that he is nonetheless a self-willed, conceited,
egoistic person is seen from his conduct, in that he despises his brother
and wounds his father. It is precisely the secret of wilfulness, of selfish
self-will, that it can thrive excellently under the restraint of civil
law, of external moral order. By introducing strict household discipline
the father in the parable, even had he had a hundred sons as frivolously
self-willed as ``the lost one,'' would doubtless have trained them all to
punctual obedience; for paternal authority is great when power is applied
in time and compulsion passes over into habit. By a cleverly calculated
and prudently executed system of compulsion he would doubtless have
secured himself against the excesses of self-will; but there is one thing
he could not do --- for not even God in heaven can do it --- he could not
by any system of compulsion whatever force freedom upon his children.

\greekrun{γ}{The Burdens of Freedom.}

That freedom in the absolute sense is its own end is surely self-evident;
but where man is involved,
% ---- printed p.411 (PDF 424) ----
\opage{411} the life of freedom is in the making, and freedom in the making
is a continual liberation. From what? We can answer at once: from
unfreedom. But since unfreedom, as a finite condition, has many forms, and
since here in this context the subject is in particular civic freedom in
the Kingdom of the Son, attention is involuntarily drawn to the many kinds
of toils, burdens and oppressions that otherwise rest upon the citizens of
any other kingdom whatever (even of one that could otherwise boast of the
very freest constitution), but which all without exception fall away, or
at any rate are destined to fall away, in freedom's own kingdom. By this
are meant all the burdens, oppressions and cares that arise from temporality
being allowed to press with all its enormous weight upon the mind of man
without the counter-pressure of the eternal.

The goal of freedom is to lift the pressure of temporality by the
impression of eternity. For the eternal does not press upon man's mind in
the way the temporal does. Were that the case, it would be terrible to be a
citizen in the Kingdom of the Son. Midway between two such powers as
temporality and eternity, man would surely be crushed to death under
pressure and counter-pressure of superhuman strength; and what then would
become of freedom? No, the fact is that only the pressure of temporality
acts depressingly, burdensomely, like a yoke of bondage, whereas the
pressure of eternity acts upliftingly, liberatingly, and precisely thereby
counteracts everything that depresses.

``Come unto me, all ye that labor and are heavy laden! I will give you
rest. Take my yoke upon you and learn of me, for I am meek and lowly in
heart; and ye shall find rest unto your souls. For my yoke is profitable,
and my burden is light.'' ``But,'' someone says, ``then here, after all,
is a yoke!'' Yes indeed, but the yoke is profitable; for it is the yoke of
freedom. ``And a burden!'' Yes, but it is the burden of freedom, and the
burden is light.
% ---- printed p.412 (PDF 425) ----
\opage{412} For what really matters is not what burden one has to bear, but
how he bears it, what strength he has to bear with. Now here lies the
great difference between the yoke of freedom and that of bondage: that
the latter taxes one's strength without giving strength, whereas the
former gives strength, the strength of eternity, and all the more and the
greater the more freely the yoke is borne. If we take the yoke of meekness,
of humility, of patience, then it is certainly a heavy yoke, insofar as
the temporal alone is thought of here. Who, in the selfish temporal sense,
is meek, except the one who feels within himself that he must put up with
much because he is capable of little? Then there is in meekness a gnawing
sorrow, a bitter melancholy, a hidden pain: it bears a dark mark of
unfreedom. The same burden weighs upon humility; it goes with bowed head,
downcast gaze. Who, in the selfish temporal sense, is humble, except the
humbled one, the deeply shamed one, who, alas, has had to feel his
impotence? What cowedness and powerless resignation to the inevitable does
not suffering, defenseless patience display! The free, proud nature
announces itself with courage, not with patience. --- ``That which in mean
men we entitle patience is pale cold cowardice in noble breasts.'' ---
Thus every yoke in the temporal bears all the marks of temporality and
bondage. But if the impression of eternity is added, the yoke is changed.
The burden in itself remains the same; weighed on the scales of
temporality it may even become greater; but the strength does not remain
the same, on the contrary! it is heightened by infinite increase. But when
the burden, however heavy it may be, continues to be finite, while the
strength grows infinitely, the yoke must surely become infinitely light.
Who, then, is the humble one? He who, in spite of all the world's scorn,
does not bow before any finite authority, because he has in freedom bowed
deeply before the ideal of all authority, the almighty Self. So he bears
the yoke of scorn, but
% ---- printed p.413 (PDF 426) ----
\opage{413} he bears it infinitely lightly. And who has the greatest
superiority, the selfishly courageous one or the one who is patient in
humility? When the courageous one challenges a stronger, who takes his
armor from him and binds him, then the courageous one is unfree: he
cannot possibly be at once both bound and yet free; but the patient one,
who, although he is bound, is yet free, free like a Paul in chains, is the
superior. Under the yoke of freedom, in virtue of inner liberation,
meekness, humility, patience are precisely the highest, purest expression
of personal superiority. He who says, ``I am meek and lowly in heart,''
says it as one who knows within himself that he has the strength of
eternity to bear all the burdens of temporality; he says it with divine
superiority.

That the eternal exists for a person, so that the person is truly absorbed
in the present eternity, is recognized precisely by his being freed from
cares about the temporal, and this in turn is recognized by his living in
the moment. Living in the moment is a peculiar matter. The frivolous man
lives in the moment; the true citizens of the Kingdom of the Son live in
the moment: the moment is infinitely different according as it expresses
punctually either the temporal or the eternal. The citizens of a temporal
state may, when they belong to the favored, doubtless be exempted from
many cares that otherwise weigh heavily upon the lowlier; but there is one
kind of care from which they cannot be exempted, and that is cares about
the future. The number of cares about the future is Legion; there is a
day tomorrow, and its possibilities are incalculable. In the Kingdom of
the Son eternity is put in the place of the future, and the day tomorrow
vanishes. ``Take no thought for the morrow.'' Eternity is what is absolutely
present in the moment. He who is absorbed in the present eternal does not
see the anxious possibilities of the future; in order to move toward the
goal,
% ---- printed p.414 (PDF 427) ----
\opage{414} he has turned his back to the goal. ``It is always delaying and
distracting to want impatiently, every moment, to look toward the goal, to
see whether one is coming a little nearer to it, and now again a little
nearer. No, be resolved for eternity and in earnest, and then you turn
yourself wholly toward the work --- and your back to the goal. . . .
% print: the ellipsis is printed with a comma for its first point (", . .").
One would think that the believer was farthest removed from the eternal,
he who has wholly turned his back on it and lives today, while the one who
peers ahead stands and looks toward it. And yet the believer is nearest of
all to the eternal, while an apocalypticist is farthest of all from the
eternal. Faith turns its back on the eternal precisely in order to have it
wholly with itself in this very day. But if a person turns, and especially
in earthly passion, toward what is to come, then the next day becomes an
enormous, confused figure, like that of the fairy tale. Just as those
demons of whom we read in the first book of Moses begot children with
earthly women, so what is to come is the enormous demon that, with man's
womanish imagination, begets the next day''\footnote{S.\ Kierkegaard.
Christelige Taler. Kjøbhvn.\ 1848. S.\ 84--85.}.

The goal of freedom is a punctual fulfilment of the laws of the spirit;
but it makes an infinite difference whether these laws are conceived as
laws of compulsion or as laws of freedom. What the difference means can be
seen by a glance at the legislation. ``Ye have heard that it hath been
said: Thou shalt love thy friend and hate thine enemy. But I say unto you:
love your enemies, bless them that curse you, do good to them that hate
you, and pray for them which do you harm and persecute you; that ye may
be the children of your Father which is in heaven; for he maketh his sun
to rise on the evil and on the good, and sendeth rain on the just and on
the unjust'' (Matt.\ 5). Were this now to be conceived as an inner,
spiritual law of compulsion, a cold, inescapable law of
% ---- printed p.415 (PDF 428) ----
\opage{415} duty: then in truth the yoke of the law would become so heavy
that man would have to sink under the burden. It is otherwise, however,
when the law is conceived as a law of freedom; for in the feeling of
infinite deliverance there is an infinite joy. It is a psychological fact
that a person can, in what he calls the happiest moments of his life, be
so glad, can in the inwardness of joy feel himself so freed from all
selfish considerations, so lifted above all personal injuries, that he
could even fall upon the neck of his bitterest enemy. Let this be a passing
mood; such a mood nevertheless testifies that there is a joy, a joy of
deliverance so rich and great that the one delivered shall surely be able
to repay evil with good and to love even his enemies.

\runhead{b) The Kingdom of Love.}
\parmark{28}

In the living, immediately conscious unity of love with freedom lies joy.
The Kingdom of the Son is a kingdom of joy because it is a kingdom of
freedom and of love. But just as freedom, conditioned by the struggle of
self-development, must fight its way forward to the goal, so the same holds
of love: one may then speak of a struggling love just as well as of a
struggling freedom; indeed, it may with reason be said that they not only
fight together as brothers in arms against a common enemy, lovelessness,
bondage, but that each, during the struggle, also ceaselessly strives to
cover and shield the other, love freedom, and freedom in turn love.
Therefore there is joy in the struggle itself, but only when the struggle
has ended in victory is the joy complete.
% ---- printed p.416 (PDF 429) ----

\opage{416} For the consideration of struggling love three points of view
present themselves here: \emph{its essence, its goal, its external
resistance}, i.e.\ \emph{its limit}.

\greekrun{α}{The Essence of Love: Hatred and Love.}

It lies in the essence of love that it must have not only its object but
also its content. In love of self to self the emphasis is laid exclusively
on the object; for only a self can in the unconditional sense be not only
the subject of love but also its object, its counterpart. In the love
of the Idea, on the other hand, the object vanishes, or rather ``it
transforms itself into sheer content.''
% print: the quotation is never closed (the 1869 print has no closing mark).
In love of art, for example, one can of course say that art is the object
of love; but, looked at more closely, art is nevertheless no more an
external thing than an existing self is. Art is essence; to love art is to
immerse oneself in its essence, to conceive its phenomena as expressive
forms for the content of the Idea. He who with his love of art gets no
further than a fancier's taste for the great masterpieces as precious
objects, ``objects of art,'' does not love art.

That it is really the self that gives personal love its object can be seen
directly from the opposite of love, namely hatred. One can certainly, in
unfettered style, speak of hating the most diverse things, certain Ideas,
e.g.\ freedom, certain animals, e.g.\ cats. But if we look more closely, it
is the passion of hatred that personifies now an object, now a content of
essence, and does so with all the greater vehemence the more strongly the
intolerable forces itself upon the individual. The involuntary
personification indicates that the real object of hatred must be something
personal.

Love of the good presupposes hatred of evil. But in the religious sense,
as we have seen, the
% ---- printed p.417 (PDF 430) ----
\opage{417} good is inseparable from the good Self, the Good One; in a
corresponding way evil, too, is then inseparable from the evil self, the
Evil One. In the language of humanitarianism it sounds fine enough, to be
sure, that one should hate faults but not persons, hate the evil in men
but not men. But religion, which after all inculcates love of man in a
quite different and higher sense, cannot thus without further ado separate
the evil in men from evil men; religion maintains that evil is both one
with and yet distinct from man's self. And what does it then do in order
to carry out this two-sided conception in practice? It places --- out of
love of man --- the original evil in a superhuman self, the evil self,
which is unconditionally the object of hatred.

The ambiguous connection of the evil self with what is selfish in men and
with the multiplicity of natural impulses, conditions of actuality and
temporal states, through which the threads of reciprocal action wind
enigmatically, this great, intricate, ambiguous whole receives in the
Kingdom of the Son the name of the world. How ambiguous this figure is
emerges clearly enough from the fact that it is set up as the object of
both love and hatred. God loves the world, and God hates the world: love
of God is hatred of the world. By proceeding from the category of the world
it is impossible to arrive at clear insight into what it really is that is
here to be loved and hated. Therefore we proceed from love; in love itself
must lie the principle for the determination of its content and object.

If we set the two propositions, \emph{love of God is hatred of the world},
and \emph{love of the world is hatred of God}, sharply against each other
without any intermediate determinations, then both lose their validity in
the Kingdom of the
% ---- printed p.418 (PDF 431) ----
\opage{418} Son, precisely because the Kingdom of the Son is the eternally
true kingdom of love.

According to the first proposition love does indeed receive its
unconditional object, but without content. The object of love is God:
``Thou shalt love God above all things,'' this is the true part. But in
what and by what is love of God now to show itself? By hating the world
and all the things that are in the world! thus by casting away the content
of love, by merely expressing love negatively; here the untrue part comes
in. To love God is impossible without loving men. So love of God must then
receive its determinate tasks, its positive content, in the so-called works
of love: to feed the hungry, clothe the naked, visit the sick, etc. But now
the principle has sublated itself; for he who has literally given up
everything in the world must reduce food and clothing to the least
possible, indeed he must even hate such necessities as something that
belongs to this world. But he who himself hates food and clothing because
he hates the world, how can he with a good conscience want to seduce others
into accepting what he himself, out of love of God, must hate? The simplest
consequences of the principle lead to madness. He who practices works of
love without perversity therefore proceeds, too, from an entirely different
presupposition, for he proceeds from the presupposition that there is
temporal need which ought to be relieved temporally, that the world,
however evil it may be in many respects, nevertheless has certain goods
that are not to be despised; more definitely still: he proceeds from the
presupposition that it is precisely in the actual world, in the acquisition
and use of its goods, that love of God must seek its tasks and find its
content.

According to the second proposition love receives its content on the
condition that it loses its unconditional
% ---- printed p.419 (PDF 432) ----
\opage{419} object. First the love of the Idea is without further ado made
one with love of God, as when it is said that the true cultivation of
science is worship of God, the inspired practice of art worship of God,
etc.; from here there is then a fairly smooth transition, or rather a
gradual descent, to the religion of minding one's trade. If now the
commandment of love, Thou shalt love God above all things, is expressly
made to hold against the higher as well as against the lower worldliness,
then ``the modern consciousness'' rises up with the protest that the divine
Self, which in science is without reality, may not exercise any influence
on practical life either. The gaze, it is said, becomes confused by staring
at a beyond instead of seeking its tasks in the present. As long as the
religious representation of the divine Self, jealous of its exclusive right
to men's love and worship, is not given up, so long will the view of life
be ascetic at its innermost core; believers can, despite their interest in
culture and humanity, only with a bad conscience give themselves over to
what in their quiet minds they must condemn as ungodly eudaemonism.

In the principle of faith, i.e.\ of love, correctives for both errors are
contained. Love cannot let go of its object; its unconditional object is
God; and God wills to be loved above all things. All things in the world,
all things without exception, must therefore be given up, indeed become
objects of hatred, when it turns out that they weaken or supplant love of
God. There may therefore arise a state of the world that makes it
necessary for true love of God to express itself negatively in renunciation
of the world, indeed in hatred of the world. But it stands no less firm for
all that, that love must have its content, a content that is conditioned
by man's relation to the world. What makes the relation necessarily
negative under certain conditions, the same thing makes
% ---- printed p.420 (PDF 433) ----
\opage{420} it necessarily positive under other conditions. The decisive
thing does not depend on the difference between the negative and the
positive, but it depends on the expression itself, whether it is otherwise
negative or positive, becoming a true expression of love, an essential
expression of the content of love. If unconditional love of God can find
its expression in a positive content of the world, in civic activity, in
working for Ideas, etc., then the relation ought to be positive; if it
cannot find its expression in anything positive, whether the reason for
this is to be sought in the inner peculiarity of the individual concerned
or in given conditions of the world, then the expression must be negative.
The negative is in itself no better than the positive. Fasting,
mortification and flagellation can, when such penitential exercises are
deified, just as well draw the mind away from true love of God as artistic
or scientific activity can when the Ideas are deified. It is an illusion
that a person, by turning away from the world, should be able to love God
solely as God, so that love itself would become the content of love. By
what, then, should a person express pure love of God positively? By
unconditional renunciation of the world! Thus by expressing it negatively.
Even the most negative expression will need something temporal through
which to take up its positive moment, namely the eternal. For the eternal
is indeed qualitatively different from the temporal; but it receives its
form only by reflecting itself in the temporal, by its pure light being,
so to speak, refracted in the prism of temporality.

\greekrun{β}{The Goal of Love: The Salvation of the World.}

``Love of God is hatred of the world'': when this proposition, as a
revivalist shibboleth for sour-minded godliness, without any correctives
and intermediate determinations, is to be the sign of citizenship in the
Kingdom of the Son,
% ---- printed p.421 (PDF 434) ----
\opage{421} then protest must be made in the name of the Gospel. If hatred
of the world is the highest emblem on the banner of Christ, then Christ
cannot possibly be cognized as the Savior of the world. If it begins with
hatred, then it cannot end with salvation either. As self-contradictory
as it is that a savior of the world should begin his work as a hater of
the world, just as contradictory is it that a hater of the world should be
able to appear as savior of the world. Christ's work does not begin with
hatred, no, it begins with love. ``For God so loved the world, that he
gave his only begotten Son, that whosoever believeth in him should not
perish, but have everlasting life'' (John 3:16).


As regards the beginning, it is highly significant that the first sign by
which Jesus, according to the Gospel of John, ``manifested his glory'' to
the disciples was precisely the miracle at the wedding in Cana (John 2).
That Christ sits at table among the wedding guests and surprises the
steward of the feast by turning water into wine stands in such decisive
contradiction to everything that resembles an ascetic hatred of the
spirit's immediate joy in life, that this beginning bears witness from the
ground up to the opposite. Without in the remotest way offending the holy,
or supposing that Christ would have reduced the miracle to a means of
festive entertainment, we may freely say that he chose precisely a setting
and a moment which in the most decisive sense are suited to ward off the
misunderstanding that the pure, holy joy he had to bring should set itself
in hostility against a universally human joy in life. Were it actually
Christ's intention to sow discord between joy in God and the natural joy
of life, then the critique of faith would simply have to declare the
Gospel of the wedding at Cana spurious. For the matter is plainly presented
thus, that the joy over ``the good wine'' heightens the joy which had
already begun with ``the worse.'' Here so gentle a transition is made from
the natural
% ---- printed p.422 (PDF 435) ----
\opage{422} to the spiritual, that the joy of faith comes as an addition to the
joy of life, a divine wedding gift, a blessing that glorifies the feast.
Were the miracle to be understood otherwise; were its significance to be
placed in this, that it contains a sentence of punishment upon natural
joy: then an elucidating word would necessarily have had to be added. The
elucidating word would then presumably have had to run somewhat thus:
``All sensuous enjoyment, all natural joy, and so also all wedding joy, is
in its essence sinful. To show how sinful it is to drink the wine that was
poured out here for the wedding guests, I offer you another wine. Whoever
drinks the new wine drinks in hatred of the world and of all the world's
joys, the joys of family life too; but he who will not hate the world and
the world's joys, to him the new wine will not taste good.'' How far such
an interpretation departs from the character of the narrative can then be
seen, among other things, in the steward of the feast. For the steward,
who tastes the new wine, must find that the wine is good, and yet he does
not taste as one awakened, who is on the point of despising all the
world's wine and all the world's joy; no, he tastes as a steward of the
feast, freely brings his steward's experience to bear, without the least
inkling that the good wine was to herald a hatred that would then have to
turn first of all against stewards of feasts.

To convince ourselves properly how false and erroneous it is to conceive
the fundamental character of the Gospel ascetically, we cite a couple of
features from the Synoptics. John the Baptist, the preacher of repentance
under the old Law, appeared as an ascetic; he ``had his raiment of
camel's hair and a leathern girdle about his loins; and his meat was
locusts and wild honey.'' But that Christ, precisely in this respect,
presented in the eyes of his contemporaries a counter-image to John, we
learn from his own words: ``But whereunto shall I liken this generation?
It is like unto
% ---- printed p.423 (PDF 436) ----
\opage{423} little children sitting in the markets, and calling unto their fellows,
and saying, We have piped unto you, and ye have not danced; we have
mourned unto you, and ye have not wept. For John came neither eating nor
drinking, and they say, He hath a devil. The Son of man came eating and
drinking, and they say, Behold a man gluttonous, and a winebibber, a
friend of publicans and sinners'' (Matt.\ 11:16--19). A stronger
testimony that Christ by no means appeared as an exemplar for those who
express their hatred of the world by fasting and drinking water can
scarcely be given. Let one consider what it means that the Pure One, the
Holy One himself, so far disdained all outward, arbitrarily chosen signs
of godliness that slander could even find occasion to brand him a glutton
and a winebibber!

But, someone might ask, how can this conception of the joy of life be
reconciled with the doctrine of continued suffering? For an answer to
this question it is necessary to recall that the absolute heterogeneity
of the principles, of faith and of reason, has not merely theoretical but
also practical validity. Natural sorrow and natural joy form finite
opposites, such that the more joy grows, the more sorrow diminishes, and
conversely; this holds not merely directly for the cases in which the
occasion is the same, but also, if only mediately, in such cases in which
joy over the one thing weighs against sorrow over the other: the more
alike the causes are, the more direct the opposition. It is otherwise with
religious suffering and the natural alternation of sorrows and joys; here
the causes are absolutely heterogeneous. Religious suffering forms an
immediate opposite to natural joy just as little as to natural sorrow;
natural joy, when the principle of faith is held fast, can just as little
diminish the inner religious suffering as natural sorrow can increase it.
% ---- printed p.424 (PDF 437) ----
\opage{424} This is not to be understood as if the absolutely heterogeneous causes
did not concern one another, could not enter into interaction with one
another. On the contrary! Religious suffering will act upon natural
sorrow, now increasing it and now diminishing it, and in both cases give
it a stamp of its own; likewise with religious suffering and natural joy.

If the Christ-principle were not absolutely heterogeneous with any
whatsoever among the many and various principles of the human world, and
if the absolutely heterogeneous principles could not enter into
interaction, Christ could not possibly be the Savior of the world. His
kingdom, which is not of this world, would then either have to relate
itself with absolute indifference to everything its citizens might
undertake in purely worldly affairs, or else demand that all worldly
institutions and all worldly order be abolished. But that neither is the
case every sober-minded person will surely concede. Christianity will
transform the mind and way of thinking of human beings; nothing more ---
with that everything is said. The inward transformation is absolutely
heterogeneous with any relative change whatsoever, be it inward or
outward, that can take place in a worldly respect with human beings and
human societies. But the inward, religious transformation in turn exerts
an infinitely reshaping influence upon everything worldly. Just as the God
who preserves and governs all does not himself intervene finitely here and
there in the course of events, because he stands not in a finite but in
an infinite relation to the world he governs, so too the world-redeeming
principle will reshape the world of human beings without intervening
finitely in worldly affairs.

Just as little as Christ with his Gospel will ascetically weaken the
natural conditions of life, just so little will he undermine either the
state or civil society: there is nothing in his ethics that encroaches
upon the right of property.
% ---- printed p.425 (PDF 438) ----
\opage{425} Rich and poor, lowly and mighty, authorities and subjects, lawful
livelihood and civil property: all such things belong, after all, to the
order of this world, and Christ does not hate this world. How indeed
should he hate it, when he wills to save it? That he wills to save it
cannot possibly mean, then, that he wills to rob it of the properties that
essentially belong to it as world. If that were so, he would indeed not
save it but destroy it. But he wills to save it, he wills to free it from
a dangerous evil, a mortal disease; for he sees it: the world is sick, its
life is in danger, it --- has been bitten by the serpent. Its strong
constitution, rich in rejuvenating powers, does indeed resist; but the
wound burns: there is fever heat; the pulse races: there is poison in the
blood! So he wills to save this sick world, to heal it of the serpent's
bite of egoism, of the poison of self-contradiction, of sin.

\greekrun{γ}{The Limit of Love: The World's Hatred.}

It is not Christ who hates the world; but there is in the world a
something that hates Christ; this something, which sets a limit to
all-embracing love itself, must then be the object of hatred for those who
in truth love Christ. What, then, is this something? It is, to designate
it by the name of love, self-seeking love, selfish self-love. Love of God
is hatred of the world; but the world that is here most immediately the
object of hatred is not the great outward world with its emperors and
kings, but the little inward one, the world of self-will, of self-seeking
love, of vanity, which the individual can contain in his own breast: love
of God is hatred of oneself. But to hate oneself, how foolish, how
unnatural! Only the strong self can love, and only the strong self can
hate; but when the selfishly strong self is forced toward
% ---- printed p.426 (PDF 439) ----
\opage{426} a turning point where it has only the choice between unconditionally
loving itself and hating the truth, or unconditionally hating itself and
loving the truth, how natural then that it chooses the former. ``Richard
loves Richard --- I am myself alone.'' Now if self-love can so easily
transform itself selfishly in the individual human being into hatred of
the truth when the truth demands a sacrifice, then this will surely be
able to happen with still greater force and on a far larger scale when
individuals do not stand alone, each by himself, but feel themselves
united in a common body, a comprehensive, manifoldly ramified universal
body, a social body, so permeated and interwoven and worked through from
top to bottom by egoism in all forms and in all shapes that even the
oppressed find it natural that those in power defy the laws: such is the
unredeemed world.

It is a peculiar matter with such virtues as compassion, love of
humankind, and self-sacrifice for truth. At a distance, in general, such
fine virtues are of course praised by all; one would therefore think that
if someone appeared who practiced them unconditionally, he would have to
be honored, loved, indeed worshiped, most of all by the eyewitnesses. And
yet precisely the relation between Christ and the world shows that the
reverse happens. It happens on such a scale that we learn from it how the
Mystery in Christ's revelation relates to the paradox.

In itself, Christ's self-sacrificing love must indeed be said to
be a mystery, insofar as in spiritual purity, depth, and inexhaustibility
it surpasses human comprehension: no one is able to put himself in
Christ's place, see with his eye and feel with his heart. But we can by no
means say, so long as we keep exclusively to the concepts of faith, that
this love is a paradox. The love of the God-man and the God-man himself
agree perfectly; here there appears no contradiction insoluble in
% ---- printed p.427 (PDF 440) ----
\opage{427} the form of the concept. It is otherwise, however, when we take our
stand on the standpoint of existence and seek, as S.\ Kierkegaard
expresses it, ``to become contemporary with Christ.'' For what happens
then? Then ideal love steps forth as an actual, existing human being;
actual love and egoistic actuality come so close to one another that a
cutting conflict arises. To see divine love in the form of a servant is
like seeing the existing contradiction.

Outwardly, immediately, the God-man looks like any other human being,
without power, without means: he has not, after all, ``where to lay his
head.'' That is the immediate impression of actuality, and sensuous
immediacy has a striking power; it acts irresistibly upon the beholder.
So it is, then, in actuality a lowly human being, insignificant in the
eyes of the world, who appears here. And yet this human being speaks with
an authority as if he were Lord of heaven and earth. Indeed, he not only
speaks, he acts out of the highest ideal presuppositions of love; every
word, every action then looks like a monstrous exaggeration, an
intolerable extravagance: here we have the paradox.

Even if Christ had not spoken a single word against the hypocrisy of the
Pharisees, had not with a single word rebuked the unbelief of the age, his
life and works, his love of humankind, would still have contained such a
judgment upon the egoistic age that this alone would have been enough to
arouse the bitterest hatred. Since he now not only bore witness to the
truth but also laid bare the lie, not only himself did the works of love
but testified against the world ``that the works thereof are evil'': the
hatred flared up, a hatred that knew no bound.

Yet it is with many an ingrained prejudice that the world
% ---- printed p.428 (PDF 441) ----
\opage{428} in Christ's time must have been more evil than it is now, since even the
chosen people could commit so crying a crime as that of crucifying the
Savior of the world. One must really pull oneself together and keep awake
in order actually to understand that the same hatred which then followed
the King of truth on the way of love follows him still, wherever he shows
himself present in his true witnesses. Had Christ, to use S.\
Kierkegaard's expression, ``looked the way the merely human notion of
compassion would form his figure, and had he had the merely human notion
of what compassion is,'' then neither would his love have been thus rewarded
with hatred. ``For human beings are willing enough to practice
compassion and self-denial, willing enough to strive after wisdom, etc.;
but they want to determine the measure themselves, that it should be such
to a certain degree; they do not want to abolish these glorious virtues;
on the contrary, they want --- at a cheap price --- in all comfort to have
the appearance and name of practicing them. True divine compassion is
therefore unconditionally the victim as soon as it shows itself in the
world. It comes out of compassion for human beings, and it is human beings
who trample it down. And while it walks about among them, the sufferer
scarcely ventures to flee to it for fear of men''\footnote{Indøvelse i
Christendom. S.\ 67--68.}.

Christ says to the disciples: ``Blessed are the eyes which see the things
that ye see. For I tell you, that many prophets and kings have desired to
see those things which ye see, and have not seen them; and to hear those
things which ye hear, and have not heard them'' (Luke 10; Matt.\ 13).
% print: comma before ``Thi'' (``I see, Thi''); full stop supplied
These words fit all who in faith can become contemporary with Christ,
thus all who can actually endure the sight of the contrast between the
inner glory and the outer lowliness. But that those who without living
faith regard him in the form of a servant
% ---- printed p.429 (PDF 442) ----
\opage{429} receive no attractive but on the contrary a repellent, offensive
impression, the more offensive the more actual the contemporaneity is
made, S.\ Kierkegaard has shown by describing the phenomenon of his
appearance. What an appearance, and in what company! ``Everyone knows,
after all, that one judges a person by the company he seeks --- and his
company! Yes, his company one could describe as this: to be cast out of
human company; his company is the lowest class of the people, further it
is sinners and publicans --- whom everyone who is the least bit anything
shuns for the sake of his good name and reputation, and a good name and
reputation is surely the least one can wish to secure in life; his
company is furthermore lepers, whom everyone avoids, madmen, who only
arouse horror, the sick and the wretched, poverty and misery. And who then
is he, he who in this procession is still the object of the persecution
of the mighty? He is one who is a seducer, deceiver, blasphemer,
despised! It is a kind of compassion if some person of standing does not
actually express his contempt for him --- that they fear him is of course
another matter.''

\emph{``If the world hate you, ye know that it hated me before it hated
you''} (John 15:18).

\runhead{c) The Kingdom of Blessedness.}
\parmark{29}

The purer, clearer, and stronger the opposition between love of God and
the selfishly self-loving mind appears, the more palpable will be the
strife between the holy and the sinful, and with it also the tension
between the eternal and the temporal. When Christ, with divinely
strengthened emphasis, repeated the Baptist's exhortation
\textit{μετανοεῖτε ἤγγικε γὰρ ἡ βασιλεία τῶν οὐρανῶν}, it was as if
% ---- printed p.430 (PDF 443) ----
\opage{430} eternity and time collided: the collision was so violent that the
present and the to-come went apart from each other like two heterogeneous
worlds. The highest is offered, and the highest is demanded; by this
every intermediate stage of finitude vanishes, and the knot of decision is
drawn tight to the utmost. Had it been a kingdom of happiness, an earthly
messianic kingdom, in which a supernatural splendor of the eternal fell in
rich rays, glorifying, upon the temporal, then the tension between the
present and the to-come would easily have been evened out; but then the
egoism of the world would also only have been ennobled and, by the
ennoblement, strengthened, but by no means sublated. Now the ``kingdom of
heaven'' proclaimed by the Savior is not a kingdom of happiness but a
kingdom of blessedness, thus a kingdom that reduces temporal welfare and
earthly happiness to something so subordinate and vanishing that it ought
unconditionally to be sacrificed as a means for the eternal end, and only
when it has been freely given up can come along as an addition to the
highest, according to the exhortation: ``Seek ye first the kingdom of God,
and his righteousness; and all these things shall be added unto you.'' On
these terms, then, the opposition between the temporal and the eternal,
the present and the to-come, had to come forth in all its severity. What
significance this opposition has, under what forms it reveals its content,
will appear when we consider in particular what the Gospel teaches us
about the \emph{end of the world}, the \emph{resurrection of the dead}, and
the \emph{Judgment}.

\greekrun{α}{The End of the World.}

The hatred with which the egoism of the world counteracted Christ's
kingdom from the very beginning betrayed clearly enough that, just as
there was in the world something that could be saved, so there was also
something that could not be saved, something that according to its nature
was destined to
% ---- printed p.431 (PDF 444) ----
\opage{431} perish. The selfish being in which human self-will has entangled itself
stands in such close connection with the sensuous, corruptible world-order
as a whole, that the judgment upon the egoistic principle in the self and
the judgment upon the corruptibility of the outward, natural world must,
from the standpoint of revelation, determine each other mutually. In this
respect too the ambiguity of the world shows itself; for one can with
reason speak both of the imperishability of the world and of its
perishability. The world is imperishable according to its universal
essence; what changes is only the form of the world. But form and essence
in turn correspond to each other in such a way that an essential change of
form presupposes and brings with it a change in essence.

The opposition between the scientific and the religious view of the world
is here again conspicuous. That the world, i.e.\ this actual, bounded
world existing in a determinate state, is not only changeable but also
perishable, science must concede; it will not even find it impossible that
a time may come when the sun's fire is extinguished, the sources of heat
are exhausted, and our whole planetary system is wrapped in death and
darkness\footnote{That the end of which science speaks is another than
that of which the Gospel speaks goes without saying.}. But the end is only
partial; the form of the perishable world maintains itself in other
perishable worlds: thus the infinite continuation of the perishable gives
the scientific observer an imposing impression of the imperishability of
the world.

It is otherwise from the standpoint of revelation: a world that crucifies
its own Savior is condemned to perish. With its gaze fixed on the
corruptible in human life and the egoism corresponding to it in the human
self, religion proclaims not merely a change
% ---- printed p.432 (PDF 445) ----
\opage{432} but a transformation of the form of this world, a transformation in
which the corruptible vanishes away and the incorruptible comes to full
revelation. That Christ in the Gospel foretold the end of the world is
certain, just as it is also certain that the apostles took it to be
imminent, and that its non-occurrence early (2 Pet.\ 3:4) aroused offense.
The question now is how the presupposition is to be more precisely
conceived, how, in other words, the end of the world is to be objectified.
On the presupposition that theology is an objective science, the
dogmaticians, older and newer, have naturally all been agreed in
subjecting the doctrine of the end of the world to a scientific discussion,
both with regard to the time and especially with regard to the manner. We
shall here, following Strauss, merely cite a sample.

``Just as, on the ground of the statements of the N.\ T., it was constantly
believed in the Church that the last day and the end of the world stood at
the door, so it was also thought that this could be seen in the world,
which was supposed to be visibly growing old already. And since the
N.\ T.\ determined the end of the world as a burning up and set this final
destruction over against the earth's earlier devastation by water, some
Church Fathers already, as also more recent natural scientists and
theologians, assumed that the moist element was to be successively
consumed by fire, so that at a certain point a universal world-conflagration
would break out. From the ruins of the old world a new, more glorious one
was then to come forth; a transformation which the Church Fathers and the
Scholastics almost unanimously, in agreement with their notions of the
renewal of human bodies, were not accustomed to conceive as an
annihilation of the world's substance, but only as a change of its
accidents. In Luther's well-known remark, too, that the heavens now had
their workday clothes on but would then
% ---- printed p.433 (PDF 446) ----
\opage{433} put on their Sunday clothes, lies the same notion; but the Lutheran
dogmaticians preferred the doctrine that this world was to be destroyed as
regards its substance too. A world, and in it human bodies, both the same
as before as regards their essence, but with other properties --- that was
indeed, each by itself, a contradiction, but both sides at least fitted
together; our dogmaticians wanted their resurrection bodies to stand in
contradiction both with themselves and with their place of abode. For this
was not likewise to be a transfigured corporeal world, but the
perpetually existing heaven, which was not first to be produced by the
end of the present world: in a place immaterial by nature, then, beings
were to dwell who were still matter. The newer supranaturalism has
therefore'' --- with express acceptance of a partial end --- ``returned to
the notion of a mere transformation of the world. \dots\ When the
so-called end of the world had thus, through the restriction of its scope,
been moved into the series of natural cosmic processes, the renewal and
transfiguration of the world that was to follow upon the destruction
thereby came to stand on still shakier ground. A conflagration, Michaelis
objects, would have to turn the earth into a barren slag, and so rather
destroy than beautify it; and the comfort he offers against this is
extremely meager; for he says that one cannot, after all, know whether the
conflagration may not leave some fertile land over''\footnote{Troesl.\
II.\ S.\ 540 flgd.}.


What comes of scientific investigations of this sort, and how much
``fruitful land'' is left over in theology once the critical
Moses-fire that is ravaging it finally dies down, we will not
investigate further here. The main thing is that the standpoint
must be changed.
% ---- printed p.434 (PDF 447) ----

\opage{434} ``Heaven and earth shall pass away, but my words shall not pass
away''; thus speaks the personal truth. The one who speaks is not
impressed by the outward magnitude of the edifice of the world, for he
speaks the language of personality with authority, decisively. The
natural world with its sun and moon and stars, however mighty and great
the whole edifice may look, is subject to the laws of necessity; but
he, whose essence is freedom itself, knows with undeceivable certainty
that the laws of necessity are grounded in divine freedom. The final
catastrophe, the thoroughgoing transformation of the world with which
the Kingdom of Christ is to reveal itself ``with power and great
glory,'' cannot be objectified in knowledge; it can only be objectified
in faith. It is not the outwardly natural but the inwardly personal
that gives conviction its unshakable foundation. --- ``If you would
have full assurance of the unconditioned power of the spirit over
corruptible nature, then lay hold of freedom and beware of staring at
necessity! If it is truly the way to the Kingdom of truth, of
personality, and of blessedness that you seek, then do not stray into
the regions of impersonal cognitions; do not seek guidance from
astronomers and physicists and philosophers of nature, for the way of
knowledge is not the way of faith; but turn to him who is the way and
the truth and the life''\footnote{Cf.\ Forelæsninger over Hindrgr.\ og
Betglsr.\ for det aandl.\ Liv i Nutid.\ S.\ 194--207.}.

Christ alone has absolute certainty of the actuality of the blessed
Kingdom, an incorruptibility that is mightier than this corruptible
nature; but the certainty is personal, not scientific. The decisive
moment, the moment when he had to say to his enemies: this is your
hour and the power of darkness, that very moment he used, even at the
last, to utter his full assurance: ``Nevertheless I say unto you,
hereafter shall ye see the Son
% ---- printed p.435 (PDF 448) ----
\opage{435} of man sitting on the right hand of power, and coming in the
clouds of heaven'' (Matt.\ 26:64).

The transformation of the natural world is a Mystery; but it is not an
absurdity. The absurd arises from a double blunder. On the side of the
doctrine of faith, from depicting to oneself an actuality that
conflicts with every concept of actuality; on the side of the critique
of knowledge, when, holding fast to the actuality known to us as the
only possible one, one makes natural process and corruptibility into
correlative concepts and contrives contradictions out of arbitrarily
fixed presuppositions\footnote{That the transition from a corruptible
to an incorruptible world order can be thought without contradiction
even in science, the author has sought to show in Grundideernes Logik.
II.\ S.\ 456--65.}, instead of resolving the determinations of the
concept in the Mystery of nature and creation\footnote{How matters
stand with the impossibility of making clear to oneself the actuality
of a world to come is shown in \emph{Forelæsning over „Phil.\ Propæd.“
for Universitetsaar} 1861--62. S.\ 123--38.}.

\greekrun{β}{The Resurrection of the Dead.}

The bodily organism is the natural, indispensable foundation of the
individual personality. What significance corporeality has in this
respect is seen, from the standpoint of revelation, in Christ's
resurrection and ascension. But instead of entering upon any
description whatever of the individual peculiarity of bodies after the
resurrection, the philosophy of religion holds strictly to the simple
contrast, characterized by Paul (1 Cor.\ 15:1--28), between the
corruptible and the incorruptible, and regards every fantastic,
groping attempt to speculate out the more precise determinations as a
thoughtless deviation from the principle of faith, a violation of the
Mystery.
% ---- printed p.436 (PDF 449) ----

\opage{436} Should anyone still be of the opinion that faith, as regards the
images of the life to come, does not have enough in the actuality of
the spirit, but must feel called upon to apply the measure of palpable
actuality to the state of the blessed after the resurrection: let him
try for once what comes of analyzing the peculiar constitution of
glorified bodies. ``According to the concurring doctrine of the
orthodox teachers of the Church, the Scholastics, and the Protestant
theologians, in the resurrection the same body is restored, with the
same constituents that it had in earthly life, with skin and hair,
entrails and genitals, bones and teeth (how else could gnashing of
teeth take place in hell? asked Jerome). But now one at once got into
difficulty with the hair. Christ had said that not one hair of our
head should be lost (Luke 21:18). But how hideous would the body not
become, asked Augustine, who connected this passage with the
resurrection, if the whole length of the hair of the head and beard
cut off over time, and perhaps the nails as well, were to grow back on
it in the resurrection! It was a stroke of luck that the Lord had said
in another place that all our hairs were numbered (Matt.\ 10:30); for
in this one saw a comforting hint that with the restoration only the
number, not the length, of the hairs was meant. But what if someone in
this life had had just too few hairs? Bodily defects of every kind,
Augustine thinks, are corrected in the pious at the resurrection,
excess of fatness and of leanness evened out; yet in such a way that
the basic features of the former figure, which are necessary to make
the personality recognizable, are not effaced \dots\ But the Church's
doctrine did not represent the risen body as merely human-like as
regards its outward contour but inwardly undifferentiated like a
statue; rather, it also ascribed to it the inner organi\-%
% ---- printed p.437 (PDF 450) ----
\opage{437}zation of the old earthly body. But now ugly consequences arose,
e.g.\ with respect to the intestines, which Origen had already pointed
out; but which the holy Thomas cleared away with the information that
after the resurrection those vessels would be filled, instead of with
their present contents, with fragrant essences. The veins were allowed
to keep their blood all the more readily since, according to the
doctrine of the Supper, Christ's glorified body also contained blood;
on the other hand the genitals were deprived of semen, since among the
risen they were no longer to come into activity. But why then are they
restored, if they are no longer of any use? Origen had already asked;
likewise, what are teeth for that are no longer to chew, a stomach that
is no longer to digest? a question to which the orthodox teachers have
remained unable to give a satisfactory
answer.''\footnote{Strauss Troesl.\ II.\ S.\ 526--29.}

When one goes through the history of the natural sciences in the last
three centuries and notes the superabundance of idle reasonings and
half-scientific objections with which every one of the great,
world-famous theories has had to contend, one gets a notion of how
easy it is to come forward with a mass of even apparently acute
counter-proofs against fundamental views into whose details one has
not yet worked one's way, and with whose spirit one has not become
familiar. If it goes thus with views for which proofs and
counter-proofs are to be sought in the world of facts that fall within
the horizon of experience, then one will be able to understand that the
reasonings that are drummed up, as much in defense as in attack, where
the talk is of an order of things lying outside all human experience,
have not the least scientific value. The only thing that can be said
about it is that, once the dogma\-%
% ---- printed p.438 (PDF 451) ----
\opage{438}ticians have set up the theories, it is quite in order that critique
tears them down again.

\greekrun{γ}{The Day of Judgment: Blessedness and Damnation.}

If, according to the conditions of spiritual life, the principle of
faith is the principle of blessedness itself, then it lies herein
without contradiction that the principle of the denial of faith must
be a principle of unblessedness, and the word of revelation that says:
``He that believeth and is baptized shall be saved; but he that
believeth not shall be damned'' (Mark 16:16), is thus in perfect
agreement with itself. The scholastically objective doctrine of the
Church, which, instead of stopping at the contrast of principle
grounded in the matter itself, has entered upon far-ranging,
over-ranging considerations concerning the possibility of damnation,
with estimates of how many and which people can now really be supposed
to be damned, has, by violating the Mystery, brought an incredible
amount of confusion into religious thinking. There is, after all, an
infinite difference between depicting damnation and the torments of
the damned in order to warn against the possibility of
damnation%
\footnote{``When we thus see Lady Macbeth walking in her sleep, % print: Machbeth
washing the bloodstains from her hand, heaving those terrible sighs in
which the repressed conscience seeks to give itself air, must we not
say that the representation of eternal damnation forces itself upon us
as a reality? For it is no true, no fruitful repentance under which she
suffers; in her repentance there is only distress and anguish, but no
desire for the good; she holds fast to the evil will, and the good is
in her only as the bare possibility that is inseparable from the
creature made in God's image. But this repressed possibility, held
down in unrighteousness, rests like a hundredweight burden on her
soul, without any condition being present for relief and deliverance
through healthy repentance. And when we see her walk thus, is it not
as though she must go on walking thus through aeons, weeping hell's
tearless weeping, whereby we are reminded of the word: Fall on us, ye
mountains, and cover us, ye hills.'' H.\ Martensen. Den christelige
Dogmatik. S.\ 490.}, and positing the possibility of damnation
% ---- printed p.439 (PDF 452) ----
\opage{439} as a settled actuality and making oneself a calm spectator
at a play of unspeakable horrors; the former is religious and loving,
the latter is dogmatically insane.

That revelation not only assumes damnation to be possible, but even,
in strong, horror-inspiring images, quite emphatically sets out to
make the representation of eternal perdition present and living, is
certain: it has sought in broad strokes to objectify the actual, last
and uttermost contrast between perdition and salvation, damnation and
blessedness. But there is a difference in objectification: one thing
is objectification in faith, another is objectification in knowledge.

When the contrast of blessedness and damnation is objectified in
faith, the standpoint is decisively determined: the one objectifying
is a believer. The believer cannot possibly doubt the real possibility
of damnation; he cannot say to himself: I do indeed believe in order to
be saved, but it was perhaps also possible that I could be saved
without believing. No, the stronger faith is, and the more living the
hope of blessedness, the greater is also the fear of the opposite, of
the denial of faith and unblessedness. That this must be so lies in
the nature of the matter, and can be seen not only in a Luther but also
in a Paul. That anxiety over the possibility of damnation, however
involuntarily strong it may be in dark moments, does not become a
soul-consuming torment, the Gospel has provided for by grounding faith
on Christ's word and sacraments; but that the possibility of damnation
must follow the ideal of blessedness as the shadow follows the body is
incontrovertible.
% ---- printed p.440 (PDF 453) ----
\opage{440} The certainty that faith gives is always a personal certainty under
personal responsibility. At the thought of the last, irrevocable
judgment the believer is thus first and last referred to himself: that
the judgment shall not come upon him, that it shall not be he who is
lost. That the believer who is in truth concerned for his own salvation
must also be able to be concerned for that of others goes without
saying; but it depends on what kind of concern it is and how it
expresses itself. The concern of love will then always be recognizable
by the fact that it bears witness to concern for oneself, never enters
upon official damnations, but everywhere, even when the signs of
salvation are weak and those of perdition strong, keeps hope in
Christ's word: ``With men it is impossible, but with God all things
are possible.''

It is otherwise when the believer forgets the principle of faith,
begins with a spectator's contemplation and sets out to objectify in
knowledge. Here it is presupposed that the contemplator, for his own
person, is sure enough of blessedness. He has not even needed to work
for the assurance in fear and trembling; he was, after all, born in the
bosom of the sole saving Church and is partaker in all its goods. For
the Church is to him an objective institution of blessedness, and since
he is now once and for all in the institution --- objective in the
objective institution --- he is, with this guarantee, quite sure of his
own blessedness. He then has plenty of time to concern himself about
the prospects for the blessedness of others. Dogmatic questions such as
these, whether the heathen, not only those now living but those long
since dead, can be saved, whether heretics can be saved, whether the
devil might not perhaps also be saved in the end, etc., now acquire a
great scientific significance. How indeed could the doctrine of faith
be a science if it could not immerse itself in such scientific
investigations? That the question of the Day of Judgment, of the time
when it shall come and the manner in which the whole thing shall
proceed, must from
% ---- printed p.441 (PDF 454) ----
\opage{441} this point of view be reckoned among the great dogmatic problems
is beyond all doubt. And indeed, by indefatigable research, not so
little has been done toward a scientific illumination of the Day of
Judgment.

``Following the direction of the N.\ T.\ (?), our old Protestant
dogmaticians conceived the judgment of the world which the returning
Christ was to hold as a visible and audible event in time and space,
and they distinguished between the three moments: the hearing, the
passing and announcing of the sentence, and its execution. But a proper
hearing of all human beings who have ever lived would, as was well
understood, have to last longer than one day. --- How long the Day of
Judgment will last Michaelis does not know, but it is hardly one of 24
hours. Indisputably, says Reinhard, the judgment of the world must last
long, which could happen all the better since the clear exposition of
this important matter would be the most instructive spectacle for every
thinking being and would spread the greatest light over the darkest
secrets. But it can certainly not be determined by us at all with what
outward circumstances this great event is to be connected. --- For
about 300 billion human beings the Day of Judgment would, according to
a more recent calculation, have to last 100,000 years. Therefore
Origen and other Fathers, admittedly in conflict with the New Testament
representation, had put in place of the outward hearing an inward
beholding of the Judge as present, and a living remembrance of all the
actions of the past life, which God would bring about in one moment in
the inner being of all people. Since in this way people, even if they
remained where each one happened to be, could nevertheless be judged
in their inner being, one was at the same time --- in a more ingenious
manner than the old-Protestant dogmatics with its doctrine of the
freedom of resurrection bodies from space --- raised above every
difficulty regarding a venue for so considerable an assembly, a
diffi\-%
% ---- printed p.442 (PDF 455) ----
\opage{442}culty that as late as the nineties prompted an unnamed thinker to the
calculation that if human beings, as was only fair, were to stand in
judgment, a surface area the size of the Kingdom of Bohemia would be
sufficient''\footnote{Strauss. Troesl.\ II.\ S.\ 544--45.}.

Here is knowledge in faith and faith in knowledge, scientific signs
which show that objective dogmatics is longingly awaiting a Day of
Judgment.
% ---- printed p.443 (PDF 456) ----
\opage{443}
\parthead{Faith in the Spirit.}
\lettersub{A.}{The Essence of the Spirit.}
\parmark{30}

The objections of the critique of knowledge against positive religion
can all be comprised in the one: \emph{religion is spiritless}. So long
now has learned objective science, as much in defense as in attack,
tumbled about with the religious representations that the educated of
the present must be tempted to regard the whole apparatus of
traditional forms of revelation as soulless remains, as the corpse of
a great, outlived thought of antiquity: science issues a certificate
and attests that religion is dead. Before, however, trusting in the
certificate, we fix the burial and send invitations to the funeral, we
will nevertheless recall what is told in the tales of the Orient about
sleepers who look like corpses because they sleep with a talisman, a
magic sign, on their breast. A fortunate hand happens at the right
moment to take the sign away, and the spell vanishes, and the sleeper
opens his eyes, rises and lives.

If it is really true that religion is spiritless, then it must surely
be really true that it is dead; and --- it is a sad duty to bury the
dead. But what does it now actually mean that religion is spiritless?
That it is without spirit can surely not possibly mean that it
% ---- printed p.444 (PDF 457) ----
\opage{444} does not have the spirit of knowledge, that it is not science. But if
it turns out that it is science which, being unable to cognize any
spirit other than its own, has in its blindness tried to force its
spirit into religion and yet has not managed it, because it met
resistance, powerful resistance from the spirit of religion, then it is
surely not religion that is without spirit in its relation to science,
but conversely science that is spiritless in its maltreatment of
religion. The spiritlessness of science has in this sense taken on two
forms in particular, that of supranaturalism and that of rationalism.

Ecclesiastical-dogmatic supranaturalism is spiritless; for it sets a
multiplicity of great, decisive opposed terms outwardly over against
one another, without discovering the principle that pervades and
unites them all. In an outward, finite manner it sets God's will
against the laws of nature and lets the miracle make an actual breach
in the natural connection, although it nevertheless assures us that God
is omnipresent in the world. In an outward, finite manner it sets a
transcendent, supernatural world against this natural one, and
nonetheless transfers determinations directly from the natural to the
supernatural. In an outward, finite manner it posits three Persons in
the one God; each of the Three is God, and yet --- there are not three
Gods. In an outward, finite manner it puts God and man together in
Christ into one Person with two natures and two wills, a true man with
a human nature that is nonetheless impersonal. In an outward, finite
manner it puts the suprahistorical together with the historical,
revelation with the Bible, God's word with Scripture, divine authority
with blind authority, etc. That attempts at mediation by means of
mitigating semi-supranaturalism, smoothing semi-rationalism,
philosophically painted-up semi-pantheism, immanent-dogmatic theories
of love, and such palliatives are not
% ---- printed p.445 (PDF 458) ----
\opage{445} able to remedy the old-supranaturalist spiritlessness, the
bottomless absurdity, is understandable.


Modern philosophical rationalism is pantheism, and pantheism is
spiritless. The doctrine of immanence levels, equalizes, and dissolves
all oppositions in the absolute One, without clarifying the relation
between the unity of the Absolute and the multiplicity of the relative.
Everything is relative: right and left, right and wrong, up and down,
freedom and bondage, wisdom and folly; everything in this world is
relative, except the Absolute. By watering down the qualitative and
setting the relativities in flux, one no doubt gets off lightly from the
hard oppositions of existence; but then, with the empty, formal unity of
the dissolved oppositions, the whole thing also ends in soul-devouring
spiritlessness. The question with whose answer the whole view of life
stands and falls, namely the question: \emph{Is there an eternal self,
independent in itself?} is answered once and for all with No in the
philosophy of the Absolute; it is --- \textit{sat sapienti}!

As long as the doctrine of the Absolute was still in its first fresh
development, the expectation of the final result held the thinking
consciousness in spiritual suspense. It was, so one thought, still
possible that the Absolute, once the development of its form had been
properly carried through and all obscurities cleared up, might give
eternal life a truer, richer presence; with this expectation it is over.
One may alter the conception as often as one likes, smooth off the
development of the form as much as one likes, rub and polish the
Absolute as zealously as one likes: the spirit does not return; one does
not rub life into what is dead. The Absolute is indeed still, as before,
a night of thoughts and so also a night of the spirit, insofar as all
inner differences have vanished; but the secret has been betrayed. The
night of the spirit is not, as in those days of expectation, a holy
% ---- printed p.446 (PDF 459) ----
\opage{446} night, the starlit night of twinkling Ideas; no, now it is the vulgar
night of conceptual shadows, of trivial abstractions, of empty
generalities, the night in which all cats are gray.

The Spirit of revelation is one with the eternal Self; only in this
unity is life. For the elucidation of this rich unity we must then, on
the given conditions of cognition, seek to grasp: \emph{the self in the
Spirit, fundamental determinations in the Spirit's Self} and
\emph{the personality of the Spirit}.

\runhead{a) The Self in the Spirit.}
\parmark{31}

Although the philosophy of religion, in the doctrine of the Spirit just
as in the doctrine of the Father and the Son, must hold strictly to the
presuppositions given in the faith of revelation, there is nevertheless
something peculiar here in the development of the essence of the
Spirit. Images and representations do not fit the Spirit himself in the
way they do when the talk is either of the Son or of the Father. As
regards the Spirit, then, the conceptual development of the philosophy
of religion ought expressly to aim at showing how the elements of
representation contained in the religious view can, according to their
inner correspondence, be conceived as elements and subsidiary
determinations of the concept of the Spirit. The living Spirit must have
a foundation; in order to determine the Self of the Spirit we must
consider the foundation in particular, first in the universal, then in
the individual.

\greekrun{α}{The Spirit on the Foundation of the Universal.}

It has already been indicated in the doctrine of creation that the
Spirit, who moves over the waters, is viewed as an intermediate being
between Creator and creature. For the religious view, then, the Spirit
does not appear immediately and from the beginning as a divine Self ---
for
% ---- printed p.447 (PDF 460) ----
\opage{447} the Self is God --- but rather as selfhood, universal essentiality.
That this essentiality is not empty but infinitely rich in content
follows of itself, since the Spirit is Godhead. By letting either
omnipotence or omniscience or wisdom move over the waters instead of the
Spirit, we at once get an impression of the empty, the one-sided, the
fantastic; for knowledge, power, wisdom are subsidiary determinations of
the Spirit, are comprehended by the Spirit, who encloses all Ideas in
his universal-concrete selfhood. That this fullness of the Godhead, this
universal rich in content, is not an object of essence for knowledge but
a vision of life for faith, is seen already from the fact that the
Spirit is not in the waters but moves over them. The almighty Self that
rules the totality of the world is the center of the Spirit; the
totality of the world itself is his universal foundation.

The next characteristic feature is seen in the relation between the
Spirit and the Word. Whereas the universal in the Creator's essence and
working is referred to the Spirit, the particular and distinctive is
referred instead to the Word. The creating God speaks, and to each
distinctive word that goes forth from the Creator's mouth there
corresponds a distinctive product. The same is repeated in the
revelation to the prophets: the universal is in the Spirit that comes
upon them, but the particular is in the words that Jehovah himself
speaks to them through the Spirit. There is speaking in the Spirit, and
there is hearing in the Spirit; but what is spoken and heard is
determined by the Word. From this point of view too the Spirit shows
himself as an intermediate being; as intermediate being the Spirit can
be sent by God, poured out by God, and return again to God. If we stop
at the immediate view, it will naturally be an easy matter for critique
to come forward with a series of objections. ``God is Spirit, and yet
the Spirit is distinguished from God''; from this presupposition a
reasoning understanding can derive as many contradictions as you please.
But the contradictions vanish
% ---- printed p.448 (PDF 461) ----
\opage{448} when we bethink ourselves that it is not yet the concept of the Spirit,
but only the first intuitable traits of his original character, that we
have before us here. The living view must always imply the one when it
posits the other, the unity when it posits the opposition, and
conversely; it speaks out of the concept, intimates the concept, but the
concept itself it does not comprehend.

For the view, God and the world form two extremes which the thought of
faith cannot possibly efface; their relation is as of Creator to
creature, in logically abstract form, as of cause to effect. What is in
the cause must come to light in the effect. The totality in the cause
and the totality in the effect are essentially the same totality; here
we have the unity, and in the unity the middle. The unity rich in
content, the middle of the extremes, is the Spirit viewed as
intermediate being, the Spirit as universal selfhood. There is wisdom in
the world because there is wisdom in God; the wisdom of God and the
wisdom revealed in the order of things are in the Spirit the same and
yet not the same. The wisdom in God, the original wisdom, the
Idea-content reflected in the unity of will of knowledge and power, is
the content of the Spirit as Godhead in God, divine selfhood in the
divine Self. The wisdom in the world, the reason of the world in
objective form, is the derived wisdom, determined by divine self-power,
wisdom in actual, natural facts; it is the content of the Spirit as
Godhead in the world, as world-Godhead. The Godhead in the world thus
forms an express opposition to the Godhead in God, and yet the Spirit,
the Spirit of wisdom, is the unity of this opposition. But instead of
developing how matters really stand, in the concept of faith, with the
unity of wisdom with the infinite transition from the Godhead in God to
the Godhead in the world, revelation at the Old Testament stage stops
at the Spirit as intermediate being. The Spirit
% ---- printed p.449 (PDF 462) ----
\opage{449} is sent forth, and the creation is renewed; he is called back, and
what was created vanishes into nothing.

That the philosophy of reason could not come to development among the
Jews is comprehensible; but neither did the concepts of faith ever come
to consciousness in the form of concepts. Instead of a closer
determination of the relation between God and the Spirit of God, wisdom
in particular was made the object of a representing reflection, put in
the place of the Spirit and personified as a divine intermediate being.
``Jehovah,'' says Wisdom (Prov.\ 8 and 9), ``formed me as the beginning
of his creation, before his works, from of old. From everlasting I was
anointed. When the depths were not yet, I was born; before the
mountains were sunk, before the hills came to be. When he (Jehovah)
prepared the heavens, then I was there, when he measured out the
surface of the sea with a circle. When he made fast the clouds above,
when the fountains of the deep grew strong. When he set the sea its
bound, that the waters should not pass his command, when he appointed
the foundations of the earth, then I was by him as a master workwoman,
and was his darling, day by day, and played continually before his
face.''

\greekrun{β}{The Spirit on the Foundation of the Individual.}

From the human side the Spirit begins on a foundation of the
individual, in that he begins with the self that develops itself toward
self-consciousness. That the Spirit is essentially involved in this
development shines forth from the fact that the self-consciousness which
is to answer to the demand of reason and thereby truly satisfy itself
cannot be sated with a merely finite content, but must be filled and
moved by the Idea, by the universal enriched with individual moments.
As unity of Idea and consciousness, the Spirit, in the universally valid
form of the individual, is a spirit of humanity and thus again an
intermediate being, a middle of two extremes: subjective reason in the
individual's
% ---- printed p.450 (PDF 463) ----
\opage{450} interior and objective reason in the world of interests and things.
With the Idea of the good the spirit in man begins in particular to stir
toward its ideal independence as spirit; but independence, as we have
already seen from so many sides in the foregoing, is an ambiguity at
the standpoint of reason; only at the stage of religion, and so only as
intermediate being between the divine and the human self, does the
Spirit work as spirit in the name of truth, freedom, and personality.

From this can now be seen the difference between the working of the
Spirit in the God-man and in the believers.

Christ is, according to the Confession of Faith, conceived, not of the
Logos, but of the Spirit. The significance of this is evident. The
Logos-principle is, in its difference from the Spirit, to be determined
expressly as the principle of individuality; the Logos-self in the Son
is the bodily-psychically individualized self; the holy principle has
in the mother's womb natural human nature as its extreme: the Spirit is
again intermediate being, the self-being working in the interaction, the
actual middle of the two extremes. During the self-development of the
God-man the Spirit is therefore also in a double movement: from Christ
toward God and from God toward Christ. On the foundation of the
individual the Spirit works originally in the Christ-child; for ``the
child grew and waxed strong in spirit, filled with wisdom; and the grace
of God was upon him'' (Luke 2:40 f.). On the foundation of the universal
the Spirit goes forth from God and comes upon Christ, namely at the
baptism. That both movements are equally original, that the one neither
excludes the other nor makes it superfluous, is grounded in the ideal
relation to God, in religion itself; and determined by the Word, the
Word of life, which was in the beginning. With the manifestation of
Christ the Spirit enters into a new relation to the Word. At the Old
Testament stage, as we have seen, the Spirit, the divine essence in its
universal selfhood, is the first, and the Word, revelation's
% ---- printed p.451 (PDF 464) ----
\opage{451} particular expression, the second; at the New Testament stage, on the
contrary, the Word, the Christ-life in its whole intuitability and
individual determinateness, is the first, and the Spirit the second:
Christ reveals the Spirit, so that the Spirit in turn may reveal him,
transfigure his life for the believers.

In Christ's divine-human Self the Spirit of the Gospel has its
unalterable center and works as the intermediate being that inwardly
unites Christ with the believers. In this his mediating activity the
Spirit again shows himself as divine selfhood. He goes in selfhood from
self to self, uniting self with self. Just as it is not the universal
essence of faith (objective doctrine of faith, system of faith) but the
believing self in which and upon which he works, so likewise it is not
the universal essence of Christ, separated from his individuality (the
Idea of Christ in objective generality), but the Self of Christ, from
which his activity proceeds. The Spirit does not dissolve the Word, but
he brings the Word, develops the Word,
% print: "udvikler Ordet. lever" (full stop for comma)
lives in the Word and makes the Word living. --- ``I have yet many
things to say unto you, but ye cannot bear them now. Howbeit when he,
the Spirit of truth, is come, he shall guide you into all truth; for he
shall not speak of himself, but whatsoever he shall hear, that shall he
speak, and he shall show you things to come. He shall glorify me, for he
shall take of mine and show it unto you'' (John 16:12--14).

To this presupposition now corresponds the outpouring of the Spirit upon
the apostles. That the Spirit comes by a definite miracle at a definite
point in time cannot possibly be understood as if the disciples, during
their life together with Christ, had not been moved by the Spirit. But
the Spirit, who worked in the disciples an understanding of the
Master's words, worked in complete hiddenness; as long as Christ stood
there alive before their eyes, he did not come forward in particular to
the consciousness of the disciples, who were immediately dependent on
the Master.
% ---- printed p.452 (PDF 465) ----
\opage{452} The Spirit as spirit in particular independence is still with the
Father, which is why such works as opening the eye to the ``Son of
God,'' drawing souls to the Son, etc., are not yet referred in
particular to the Spirit, but immediately to the Father. Only after the
Son has gone in to the Father, after the individual manifestation of
the Spirit's content has furnished actual presuppositions and
historical conditions for his independent activity, does the Spirit come
forward with particular independence. The outpouring of the Spirit at
the feast of Pentecost is thus not merely a breakthrough grounded on
historical presuppositions and explicable from the peculiar state of
soul of the apostles, but a work of God, fulfillment of a given
promise, hence an unconditioned act of will, an antirational act of
freedom.

As a proof that the Spirit is conceived personally in the N.\ T., appeal
has been made to the Gospel account of Jesus' baptism: \textit{Abi ad
Jordanem et disce trinitatem.} ``But,'' remarks Strauss, ``as regards
the narrative of Jesus' baptism, it is a question whether the
appearance of the Spirit in the shape of a living being (a dove) is a
surer proof that it has a personality of its own than the fact that it
comes down in the form of tongues of fire?'' Christ, it is true, in
John, at the farewell with the disciples (chs.\ 14, 15, 16), speaks of
the Paraclete and thus speaks of the Spirit as of an actual
personality; nevertheless it is an easy matter for reasoning critique,
merely by applying a tiny bit of hermeneutics, to explain the personal
away. For in Luke (24:49) the Spirit is, after all, called outright the
power from on high (\textit{δυναμις ἐξ ὑψους}), and so conceived as a
power of God, and so impersonally.

\greekrun{γ}{The Self of the Spirit: The Mystery.}

All fundamental errors, all perverse representations and unclear
concepts about the self-revealing Spirit of God have their origin in a
false conception of the relation
% ---- printed p.453 (PDF 466) ----
\opage{453} between self and selfhood. We see it in the rationalism of the
understanding and in speculative knowledge.

The rationalism of the understanding makes a sharp separation between
God and the Spirit of God. God as God, the almighty Self, is held fast
in a representation by itself: God in himself is outside the world,
exalted above the world, separated from the world; but the Spirit, the
Godhead, the divine essence in infinite being and working, is
omnipresent, pervades all, moves and determines all. The divine Spirit
that pervades the world now fuses with the divine reason indwelling in
the world, and in this way becomes synonymous with the world's, with
nature's own objective reason: from rationalistic theism there is only
a small step to pure deism. ``Deism is'' spiritless, its abstract God
of the beyond is a dead representation, a \textit{``caput mortuum''}:
here speculation begins.

The dry Self of deism, shriveled up in representation, speculative
knowledge casts away as a spiritless residue, a remnant of scholastic
externality in thinking; with the Self God vanishes, but the Spirit, the
immanent Godhead of all existence, remains behind. So at first glance
it then looks as if life in the present, having renounced what is to
come, were really to become rich in spirit. ``It was the fixed idea of
the beyond which, while consciousness was estranged from itself,
constantly drew the best powers to itself and thereby spiritually
sucked the blood of the here-and-now; now, on the contrary, after
spirit as spirit has come to itself in man's knowing self-consciousness,
the here-and-now can obtain the whole present spirit and live in the
Godhead itself''. This spirit without God, this selfhood without self,
nevertheless, as we have seen above, becomes spiritless again through
the universal selfhood's turning into the selfless universal: from
speculative pantheism there is only a small step to pure materialism.
% ---- printed p.454 (PDF 467) ----

\opage{454} What, according to the foregoing, particularly characterizes the
Spirit of revelation, then, is this: that as Godhead with God he
originally mediates between self and selfhood, between self and self.
All fundamental oppositions in all forms, the individual as well as the
universal, under all laws (those of otherness and necessity as well as
those of self-development and freedom), oppositions of the material and
external, of the ideal and internal, of the unity of connection in the
world and the unity of self in God, in short: extremes of every kind
have in the Spirit himself their origin, their power, and their uniting
middle.

As Godhead over against God the Spirit is originally one with God: God
is Spirit, the Spirit is God. Just as the opposition of God and Godhead
is determined by the Spirit, so too the unity of the opposition is
determined by the Spirit. The divine selfhood in the Spirit has God
himself not outside it but within it; the divine Self in the Spirit has
its selfhood rich in Ideas not outside it but within it. The spirits of
religion and of science resemble each other in that both take the form
of universal selfhood; but they are absolutely heterogeneous in this,
that the spirit of science excludes what the spirit of religion
unconditionally holds fast: \emph{the divine Self}. As Godhead, divine
selfhood, the Spirit is \emph{rational}; as God, divine Self, the Spirit
is \emph{antirational}. When the consequences hidden in this
presupposition are developed, it will easily be seen how matters stand
with the Mystery and the absurd.

As selfhood in power the Spirit has the external actual world as his
other, and has in this other a sensuous, immediate expression of his
ideal reality. The world is borne by spirit; it is an in-spirited whole.
If now the emphasis is laid on the world's being a whole coherent in
itself, determined by the interaction of its parts, a whole of reason,
then the Spirit loses his independent Self; he goes under
% ---- printed p.455 (PDF 468) ----
\opage{455} in universal selfhood, he hides his Godhead behind the reason of the
world: his grounds belong to necessity.

As divine Self, as unity of will of power and knowledge in the form of
knowledge, the Spirit has the inner world of Ideas in God as his
essential content. If now the emphasis is laid on every determination of
the Spirit's being a self-determination decisive in itself, an almighty
determination of will which can have outward-going as well as
inward-going effects, then the antirational comes to the fore. God's
will, the Self that is divine in the Spirit, can then intervene in the
actual connection of the world where it will, and effect whatever
changes it will, without any other ground being assignable for it than
the will's own: \textit{sic volo, sic jubeo, stat pro ratione
voluntas.}

The two extremes, that of selfhood and that of the self, of necessity
and of freedom, are posited by the Spirit and can be united only in the
Spirit; they can neither fuse together nor maintain themselves against
each other in finite separateness: as the Spirit's own extremes they
must be united in a syllogism. In this syllogism the Spirit is again
intermediate being; he is the middle of the extremes. In his difference
from the extremes the Spirit as intermediate being is one with the
extremes; but this unity is in turn doubled. The Spirit is namely in
unity with each of the extremes by itself: in unity with the world as
the reason of the world in the form of necessity, and in unity with God
in the form of freedom. The opposition is thereby stretched to its
utmost. But now the Spirit is at the same time the pervading unity of
the opposite extremes in the sense that he effects, into the infinite,
their mutual transition into each other. This transition is articulated
in the Spirit himself by an inner reciprocal determination of necessity
and freedom, of unity of connection and unity of self, a reciprocal
determination which at every point is conditioned by freedom's being
reduced to a moment for necessity, and necessity conversely for
freedom.
% ---- printed p.456 (PDF 469) ----

\opage{456} In order to develop necessity by itself we must imply freedom; for
selfhood is grounded in the self. Freedom is from this side the ground
of necessity; it bears necessity, and the actuality of necessity is, in
consequence, at every point the possibility of freedom. In order to
develop freedom by itself we must imply necessity; for the self is
grounded in selfhood. Necessity is from this side the ground of
freedom; it bears freedom, and the actuality of freedom is, in
consequence, at every point the possibility of necessity.

If now our gaze is fixed on the world of facts, actual nature with its
necessary laws, then we can understand that the connection of things
can be developed and pursued into the infinite without the least trace
of unconditioned divine self-determination being discoverable in the
infinitely ramified unity of connection. The self is hidden in
selfhood, and the divine selfhood in the objective reason of the world
system. But since selfhood is grounded in the self, we can also, in
accordance with what we said before, understand that in the infinite
system of grounds of necessity there is at every point the possibility
of acts of freedom, of unconditionally divine acts of will; that, in
other words, the actuality of the connection of nature is the
possibility of the miracle.

If, on the other hand, our gaze is fixed on God as God in divine
self-power, then we can understand that every act of will, even if it
is in itself absolutely grounded, must nevertheless in the last
instance have the determining ground of all grounds in the
self-determination itself; we can understand that the almighty will can
act just as immediately and particularly on any single point in any
single part of the whole as in comprehensive generality on the whole
itself; but we can also understand that, whether the determination of
will acts on the whole as whole or on the part as part, the antirational
action can nonetheless not possibly disturb the order of things.
% ---- printed p.457 (PDF 470) ----
\opage{457} By making a breach in the rational unity of connection of the system
of wisdom, God would indeed come into conflict with his own Godhead, the
divine Self with its own selfhood, and what concept should we then be
able to form of the Spirit? No, the power of the Spirit, of the will,
consists precisely in this, that in the one extreme it always possesses
the other: the actuality of the miracle is the possibility of the
connection of nature.


If we conceive all acts of will proceeding from the self as proclamations
of the supernatural, and all grounds lying in selfhood as expressions of
the natural, then in the Spirit of wisdom neither strife nor contradiction
between the supernatural and the natural can be thought. The forming and
working of the Spirit is an infinite conversion of determinations from self
into selfhood, from freedom into necessity, from the supernatural into the
natural, and conversely. This conversion is not an endless, uniform
process; it is the content of wisdom, systems of means and ends, which is
articulated into the infinite. When the goal is posited as the goal of
freedom, as purpose, the means are presupposed with necessity; but the
necessity of the means is the possibility of freedom. That at every point
in the system of selfhood there is possibility for the self means: that at
every point in the nexus of means, causes, intermediate causes and
conditions there is possibility for the will. Every member of the great
whole is as near to the will of the Spirit as the whole itself; from which
it follows that the self-powerful will can intervene at a single point in a
single member just as immediately as it, in the Spirit, can master and move
the whole that coheres within itself. Here there is not merely unity of
necessity and freedom; here there is, within the unity, difference and
thoroughgoing opposition. Necessity sublates freedom --- what happens with
necessity cannot immediately be ascribed to freedom --- but it does not
annihilate freedom at any point; it only reduces it to possibility. Freedom
in the determinate act of will sublates
% ---- printed p.458 (PDF 471) ----
\opage{458} necessity --- the act of the will is not necessary --- but it does
not annihilate necessity; it only reduces it to possibility. So, then,
freedom and necessity stand sharply against each other in the infinite
working of the Spirit; but they do not stand in strife with each other. It
is the power of the will in the Spirit that requires the opposition; for
the self-power of the will is power over another. Necessity cannot at any
point withdraw itself from the will; for it rests in the universal willing.
Everywhere in the system of wisdom the universal willing can be
concentrated into punctually determining will.

Every self-determination in the Spirit is, as an inner act of will
proceeding from the self, antirational: all grounds of necessity are in
this concentration transformed into grounds of will; conversely, every
determination of selfhood in the Spirit is rational: all grounds of freedom
hold themselves as moments of possibility within the conditioned grounds of
connection of essentiality, of necessity. Thus in the system of wisdom the
rational can everywhere be articulated by the antirational, and nature be
borne by the Spirit; in this there is no absurdity. But how the reciprocal
determination of selfhood and self takes place in actuality, how, when and
why the antirational acts precisely at these points, in these members and
not in others, that is and remains an insoluble riddle, an impenetrable
Mystery. ``No one knows what is in God, except the Spirit of God.''

\runhead{b) Fundamental Determinations in the Self of the Spirit.}
\parmark{32}

As the self-acting unity of self and selfhood, as the articulating
intermediate being between God and the world, the Spirit must originally
contain determinations from which the relation between the divine and the
worldly can be cognized in all forms of revelation, and hence:
\emph{divine-worldly
% ---- printed p.459 (PDF 472) ----
\opage{459} fundamental determinations}. As the mediating unity between the absolute
and the relative self, as the intermediate being between God and man, the
free, liberating Spirit must contain determinations from which all inner
phenomena of consciousness corresponding to the relation to God can be
derived, and hence \emph{divine-human fundamental determinations}. As the
developed self-unity of the absolute self-duplication, as the intermediate
being between God and God, the Spirit must contain determinations from
which it can be seen, from the standpoint of revelation, that all outward
effects of God in the world and all inward activities in God over the world
are moments in the self-developing concept of God, acts of self-opposition
in the very essence of the one true God, and hence \emph{trinitarian
fundamental determinations}.

\greekrun{α}{Divine-Worldly Fundamental Determinations.}

There is nothing that in recent times has contributed more to undermining
faith in the supernatural self-power of the divine Spirit than dogmatic,
objective supranaturalism. If supranaturalism is philosophical and has
taken up speculative determinations of immanence, then critique rightly
calls attention to the half-heartedness, the ambiguity in its conception of
the miracle, to the emptiness of such profound-sounding phrases as ``a new
law of nature,'' ``a spiritual breakthrough,'' ``a progressive development
of the divine life in nature,'' etc. If supranaturalism is orthodox
scholastic, then it lays down a series of logically delimited conceptual
forms in which the supernatural is to be grasped and held fast in a
finitely determinate way. The concepts fixed by overstepping the boundary
of the Mystery then have this peculiarity about them, that they stand in
glaring contradiction with the empirical concepts of nature of physics.
The negative critique, which bases its development of concepts on the
presuppositions of natural science, then has an easy task when it wishes
to attack
% ---- printed p.460 (PDF 473) ----
\opage{460} supranaturalism, demonstrate its contradictions and point out the
absurd.

The supranaturalists hold fast to the presupposition of a supramundane,
omnipotent God, and hence of the divine Self in unconditioned freedom; in
this they are right. But when they introduce this presupposition of faith
into science, they are, however much they exert themselves to reconcile the
personal God with ``the immanent deity of the world,'' evidently in the
wrong. ``Most theologians,'' says an acute
critic,\footnote{K.\ Schwarz. Den nyeste Theologies Historie.
Kjøbhvn.\ 1868. S.\ 356.} ``use the concept of the divine personality, its
freedom and exaltation above the natural nexus, only to establish a
peculiar kingdom for the exclusively divine, the supernatural, and so to
introduce anew all the miracle stories and notions of miracles with a
speculative semblance and verbal bombast. In these circles especially a
veritable debauchery has been practiced with \emph{divine freedom}; under
this splendid cloak even the most arbitrary and most external mode of
action has been concealed.'' But if the supranaturalists meet with
misfortune when they clothe God's supramundane, omnipotent Self in the
speculative garb of modern knowledge, in order, by means of a costume
fitted to the taste of the times, to procure recognition in science for his
supernatural works, then the neo-rationalists, who by applying pure
concepts of reason would help the universally divine selfhood to pull
itself together as a willing Self and keep itself afloat above the waters
of immanence, certainly have just as great difficulties.

The divine selfhood, the indwelling deity of the world, speculative
philosophy finds it easy to acknowledge: the spirit of selfhood is, after
all, according to what we have shown in the foregoing, precisely the spirit
of science. But when speculation immerses itself in selfhood, it immerses
itself in pantheism, and rational theology, which for
% ---- printed p.461 (PDF 474) ----
\opage{461} conscience' sake cannot possibly do without God as God, now exerts
itself to work its way through pantheism up to a true religious, scientific
theism. ``Pantheism,'' it is
said,\footnote{K.\ Schwarz. Anfr Skr.\ S.\ 354.} ``has truth and
validity only as the opposite of the external and abstract theism set up
by vulgar theology. It has the merit of having constantly worked against
all anthropomorphic and anthropopathic notions of the deity and of having
purified and enlarged the concept of God on all sides. But with this
negative merit its whole significance is also exhausted, and even in a
philosophical respect it is so little justified that it must, on the
contrary, be regarded as an abstraction of the worst kind, just as
incapable of solving the riddle of existence as abstract theism. For in
pantheism not only the sensuous world but also God himself comes off badly.
The absolute is here not the truly absolute, but the \textit{caput mortuum}
of abstraction, that which remains when all other abstractions have been
thought away. God then becomes merely the abstract-universal, the One, pure
being, and whatever else one may happen to call it, that which sublates all
concrete forms of being. Thereby the beauty of the world and the richness
of its multiplicity are lost; it becomes foam and semblance, a non-being
bound up with being, a rising and falling wave in the ocean. On the other
hand, the depth, inwardness and creative energy of the deity are thereby
lost; it finds no rest in itself, it is only a streaming and foaming, a
purposeless and resultless blowing of the soap bubbles of finitude, which
it takes back again.''

Neo-rationalism sets out to transform scientific pantheism into religious
theism; that is a self-contradictory endeavor. The divine Self, which
% ---- printed p.462 (PDF 475) ----
\opage{462} according to objective grounding is to be the presupposition and result of
the development of the Idea, is, as a surplus of reflection from the
metaphysically universal selfhood, too powerless to assert freedom in the
midst of necessity. Formal willing cannot, by means of the infinite nexus
of nature, condense here and there into will; but a self without will
cannot be the God of faith. ``From the contentless play, for which the
all-embracing divine substance is the grave, tragic background, thought
strives to raise itself, in order to come to conceive the world as fixed in
itself and, in its higher organizations, as striving ever more toward
independent centers, but at the same time to conceive the deity as
revolving about itself. This is the endeavor of \emph{speculative theism};
in the midst of the necessary belonging-together and reciprocal action of
God and the world it maintains the difference between them, whereby both
first come into their own. What is true in pantheism is preserved: the
deity is not an individual beside other individual things, but the highest,
all-penetrating universal; the infinite does not stand over against the
finite, whereby it would itself be made finite; rather it is an activity
penetrating finitude and implanting itself in it. But on the other hand the
thought also comes into its own that that highest universal is at the same
time \emph{a living} universal, self-conscious, returning into itself from
its activity penetrating the world.''\footnote{K.\ Schwarz. Anfr.\ Skr.\
S.\ 355.}

How far a speculative theism can be carried through consistently on all
sides remains to be investigated in science; the investigations of the
philosophy of reason pertaining to this do not concern religion as
religion. The Self in which faith holds fast to God as God in opposition to
the world is not merely supramundane but also supernatural; it
% ---- printed p.463 (PDF 476) ----
\opage{463} could not possibly be the first if it were not the last. The possibility
that the Spirit can mediate between the natural in the world and the
supernatural in God shows itself when we consider the divine-worldly
fundamental determinations.

In this connection there is no question here of new attributes for the
% print: "Attributcr" (wrong sort c for e)
absolute or of particular predicates for the divine Self; the task is to
develop the concept of the Spirit in such a way that the two-sided
essentiality of the divine predicates can be cognized from the standpoint
of faith. What holds of some predicates holds of them all; we choose as
examples \emph{omnipotence} and \emph{omnipresence}.

That the deity is omnipotent or, what amounts to the same thing, that
omnipotence is divine, philosophy grants. Omnipotence is an attribute of
substance, and in pantheism, after all, ``the divine substance is the
grave, tragic background.'' But thereby omnipotence is conceived merely
from the worldly side, as a predicate of immanence, a predicate of the
general selfhood of the deity. How, now, do we get from selfhood to self,
from the all-embracing universal power to unconditioned self-power? A
rational transition is impossible; for self-power itself is antirational.
For rational universal power the atoms of nature with their properties,
forces and reciprocal actions have immediate reality as the other of
power; universal power has no particular existence; it is an ideal,
infinite capacity; its infinite capacity is to sustain everything finite,
to be the spiritual bond of the connection of things, the essentiality in
which and by which everything has its subsistence. For antirational
self-power the finite elements and forces of the world are as a nothing, a
possibility that can become anything the will wills. Omnipotence is God's
power, God is omnipotent; this is the divine side of the predicate. Seen
from this
% ---- printed p.464 (PDF 477) ----
\opage{464} side the predicate cannot enter into science; for it threatens to
burst the actual world, to crack nature open and make the great whole of
reason into a perforated, patched-together piecework. If the strong
opposition of universal power and self-power were not an opposition in
the very essence of the Spirit, the discord between religion and science
could never be removed. \emph{Self-power is an imagination; only the actual
power in the world of things has reality}; that is the language of power
of science. \emph{The world of things is powerless; only God is
omnipotent}; that is the language of power of religion. If, on the other
hand, we acknowledge the absolute heterogeneity of the principles, then the
Spirit comes into its own from both sides, and the strife is settled.

Since science holds fast to its principle, it has universal power as its
presupposition, and so, when in philosophy it rightly comes to its
senses, it must everywhere tacitly understand the divine power of reason.
Ceaselessly occupied with investigating, not the concept of power in empty
generality, but the actual revelations of power in phenomena and laws, it
fixes its gaze one-sidedly on connection and unity of connection. That all
actual unity of connection rests on the universally powerful self-unity,
and this self-unity in turn on the omnipotent Self, are fundamental truths
of which it must ceaselessly be reminded. Logical theism, which develops
the fundamental concept of self-unity, thus prepares science for the
concession it must make to religion in the name of the Spirit; for the
spirit of science is one-sided.

The spirit of religion is many-sided. In the religious representation the
omnipotent Self can, to be sure, be held fast so immediately for itself
that it looks as though God were ``an individual beside other individual
things,'' but the misunderstandings flowing from this belong to ways of
seeing that fall outside the Spirit. Considered in the light of the
Spirit, self-power
% ---- printed p.465 (PDF 478) ----
\opage{465} has universal power as its presupposition; only thus does it have its
opposite in itself and can, as power of will, be power over another. It is
not the case, then, that self-power should weaken universal power and with
it the nature and connection of things; the opposition of power in the
Spirit is an expression of the fact that the extremes mutually strengthen
each other. Omnipotence is in God as God's power; the divine side of the
fundamental determination is expressed in the proposition: God is
omnipotent. Omnipotence is in the world as the universal substantiality of
selfhood; the worldly side of the fundamental determination is expressed in
the proposition: the substance of the world is omnipotent. But both
determinations must be united in the many-sided proposition: \emph{the
Spirit is omnipotent}.

As subject of the divine-worldly predicate, omnipotence, the Spirit comes
to meet itself from each of the extremes. In the extreme of transcendence,
\emph{self-power}, it has in the infinite the possibility of its power of
immanence; for self-power, as power of will, is inseparably one with power
in another and power over another. In the extreme of immanence, \emph{the
substantial power}, the Spirit has at every point the possibility of the
power of transcendence, for the all-embracing selfhood is borne by the self
and has, in each individual determination of selfhood corresponding to the
unity of connection, a tacitly understood self-determination.

As a divine-worldly fundamental determination in the Self of the Spirit,
the concept of omnipotence is thus two-sidedly determined; it is not
% print: "dobbelsidig" (for dobbeltsidig)
self-power by itself, not substance-power by itself, but the unity of
opposition of both. Take self-power away, and substance-power must vanish;
for the power of selfhood is power over another only by being power in the
self. Take substance-power away, and self-power must vanish; for self-power
is power in itself only by being power over another which in itself is
power, the universal power of totality in the actual connection of things.

What holds of power in the form of essentiality
% ---- printed p.466 (PDF 479) ----
\opage{466} holds in the form of the phenomenon of \emph{omnipresence}. Seen
immediately, God's omnipresence is an absurdity. If the whole God is to be
in every place, then the divine Self must surely be able to divide itself
% print: "guddommeligo" (wrong sort o for e)
into the infinite and, in spite of the division, still remain the one and
same infinite Self. From fear of this contradiction speculation now rushes
over into the opposite extreme, abolishes the self in all haste and seizes
the universal principle; as if it were self-evident that a principle can
divide itself and yet be undivided. Held fast by itself, the latter
extreme is just as spiritless as the former. The beguiling thing lies in
an infinite interplay of immediacy and reflection. Immediately, i.e.\ in
the phenomenal form of outward actuality, God is not omnipresent.
Omnipresence is, in other words, not a predicate of an individually
existing subject, but a predicate of spirit, a predicate of reflection for
the universal, substantial selfhood, a worldly fundamental determination in
the Spirit. Corporeality, essentiality, laws are omnipresent; for
corporeality, essentiality and laws have a reflected form, immediate in
reflection. Where, then, does the religious representation of an
omnipresent God have its point of attachment? In the Self exalted above
selfhood, reflected in selfhood! The Self is not in a single place by
itself alone, as when it is said that God is in heaven, and the
understanding then concludes from this: therefore not on earth; rather the
Self subsisting in itself is immanent in its selfhood. Thus immediacy
returns with a disjunction of here and there, and omnipresence shows itself
as a divine fundamental determination in the Spirit.

Therefore: it is not God in himself, any more than it is the deity in the
world abstracted from God, that is omnipresent. Omnipresence is a
divine-worldly predicate: God in the Spirit as Spirit is omnipresent.
% ---- printed p.467 (PDF 480) ----

\opage{467} \greekrun{β}{Divine-Human Fundamental Determinations.}
% print: "Grndbestemmelser" (u dropped in the head)

What recent philosophy impresses upon us again and again, namely that
spirit as spirit cannot possibly tolerate anything foreign outside itself,
anything objective and object-like that it is unable to sublate, is at
bottom a great truth; but the misfortune is that the philosophy of reason
works most vigorously against the very truth it itself impresses upon us.
It demands that everything objective shall be spirit just as much as
everything subjective; it demands that self-consciousness shall comprehend
its own essence in the essence that is external and foreign to
representation, and comprehend itself in the infinitely rich essence; and
in this we must grant that it is right. But when the question is how one
goes about solving the great task, the aberration of spirit is glaring.
Speculation speaks of an opposition of objective and subjective spirit,
whose unity is to be sought in self-consciousness. But what sort of an
opposition is that? Its ``objective spirit'' is not spirit; in its
immediate existence it is nature, external, actual, sensuous nature, and
in this its naturalness ``the principle of subjective spirit.'' Its
``subjective spirit,'' with the natural world as object, is in its finite
subjectivity so one-sided, limited and imperfect that, by the philosophy
of reason's own admission, it is only a spiritless spirit, a spirit that
has not yet taken possession of its objective content. The opposition of
objective and subjective spirit is then, as one sees, an utterly crippled
opposition, an opposition whose one member is a spirit that is not spirit,
whose other member is a spiritless spirit.
% print: "Leder" (for "Led er", space lost)
By the union of two such spiritual cripples the sound, complete spirit,
the new spirit of truth and freedom, the self-conscious world-spirit, is
now to come forth; but it will be a spirit to match.

From the standpoint of revelation the opposition
% ---- printed p.468 (PDF 481) ----
\opage{468} between the Spirit of God and the spirit of man is infinitely deep and
rich in content; and as the opposition is, so too is the unity. But when
the divine Spirit is set freely against the human as Self against self,
there must naturally proceed from the free opposition a series of
particular reciprocal determinations; it will then easily appear to a
superficial observer as if the relation of spirit, by being drawn down into
finitude, must absolutely become spiritless.

``How the one human being can himself set himself against another, so that
a reciprocal action between finite spirits thereby arises, is evident;
but'' --- it is said --- ``how the divine Self can set itself against the
human in such a way that there truly occurs a reciprocal action between the
infinite and the finite spirit is incomprehensible. A human being can know
with certainty that he stands face to face with another; but how should
anyone be able to become conscious that at a given moment he stands face to
face with God, that the divine Self is present in the spirit? Anyone who in
all soberness will attempt an actual distinguishing between the thoughts
that come from himself and the thoughts that come immediately from God,
between what he says to himself in his conscience and what God expressly
says to him, between the powers that are originally in his own spirit and
the powers that are imparted to him by the Spirit of God, will surely
recognize that the attempt is insane. The religious-psychological
self-deception shines forth clearly enough, too, from the church
dogmaticians' disputes, as fruitless as they are protracted, over the
relation between human freedom and divine grace.''

``The Pelagians let man's natural freedom develop under the co-operation
of grace; that will not do, and why not? Because the natural, selfish
will, which is in itself unfree, has freedom only to resist freedom by
resisting grace. So the Augustinians
% ---- printed p.469 (PDF 482) ----
\opage{469} assert that it is grace itself that frees the will. Man's will,
then, does not have its freedom in determining and moving itself, but
solely in being immediately determined and moved by God's will. But, not
to mention that the freedom of the will must, when its self-determination
is placed in another will, become unfree, one surely entangles oneself,
as regards the transition from natural unfreedom to the freedom effected
by grace, in insoluble contradictions.''


``The natural will is supposed to be able to resist grace; this implies
that it can also refrain from offering resistance; over against grace,
therefore, the natural will has complete freedom of choice: it can
choose between resisting and not resisting. If it chooses to resist,
then it opposes grace and thereby shows precisely what it is, a
natural, selfish will. If it chooses not to resist, then in all its
naturalness it wills precisely what grace wills, and so cooperates with
grace and makes itself a recipient of grace. But that is precisely what
the natural will, according to Augustine's doctrine, is not supposed
to be able to do; it is not supposed to be able, by any cooperation
with grace, to make it possible for grace to
% print: „kunne gjør“ for „kunne gjøre“
work. But if it is now not a natural self-determination but a hidden
grace \textit{(gratia præveniens)} that brings about the surrender,
then the natural will does not have the slightest power to resist
grace either. Grace breaks through when it will, where it will, and
because it will; that the will offers resistance is therefore an
untruth. For the will's resistance presupposes that it refuses to do
something which it could nevertheless do according to its nature; if
it is not in its power to surrender itself, then neither is it in its
power to refuse the surrender.''

Equally absurd --- to the superficial observer --- are the results to
which a separation of man's
% ---- printed p.470 (PDF 483) ----
\opage{470} spiritual natural gifts from the Spirit's gifts of grace must
consistently lead.

``While the human spirit, under all forms and in all particular
manifestations, must constantly presuppose natural dispositions from
which it can develop, but which it cannot alter, the Spirit of
revelation, on the other hand, is assumed, by virtue of unconditional
divine freedom and without regard to the innate peculiarities of
individuals, to be able to impart its supernatural gifts, in that it
`divides to every man severally as it will.' In the use of this freedom
it goes to extremes which no one, whether he be a believer or an
unbeliever, will be able to explain.''

``Consider the miracle of Pentecost. The Spirit comes with such speed,
with such an effect, that in a moment it teaches the apostles foreign
tongues. Whether one now chooses to explain the phenomenon by saying
that it instantaneously breathes into the apostles knowledge of what
was previously unknown to them, or that it itself speaks the foreign
languages through the apostles, the occurrence is equally monstrous.
In either case the Spirit of God takes the human spirit by force, by
robbing it of self-development, in the same way as grace takes the
natural will by force by robbing it of self-determination. This is the
one extreme.''

``Consider the religious consciousness under ordinary conditions. The
one who prays turns to God, but he does not see God himself in the
spirit, he sees only his representation of God; soberly speaking,
therefore, he prays not to God but to his own representation. Let him
now pray for a gift of the Spirit, be it enlightenment or strength or
patience or comfort; if, when the prayer is ended, he will be honest
with himself, he will admit that all the stirrings took place in
himself and came from his own spirit. If he received enlightenment, he
received it through the self-collection corresponding to devotion; if
he was strengthened in prayer, the strength came, through the
gathering of his thoughts in the will's
% ---- printed p.471 (PDF 484) ----
\opage{471} earnestness, from his own interior; if he gained patience,
he gained it in himself through resignation. If, on the other hand, he
imputes a hidden miracle to his religious moods, if he imagines that
the Spirit of God has actually drawn near to him during the invocation,
has breathed upon him, imparted to him powers he did not have in
himself, and whispered to him thoughts which otherwise could not arise
in a human consciousness: then he entangles himself in a pious
deception and never comes to the cognition of the truth. This is the
other extreme.''

``As it goes with the miracle in the outward sphere, the revealed
miracle, so it goes also with the miracle in the inward sphere, the
hidden miracle; if the former vanishes, the latter vanishes too. In
the present age no miracles occur in the outward sphere. What does
that prove? It proves that in the present age no miracles occur in the
inward sphere either. Were the Spirit of God actually to move the
spirit of man, it would have to happen by a hidden miracle; but those
who, for want of revealed miracles, console themselves with the hidden
ones console themselves not with facts but with pious imaginings.''

What these and similar misgivings amount to can easily be seen from
what we have said above. Once again it is the absolute difference of
knowledge from faith on which everything depends. In knowledge there
is no relation to God; for the one who knows, God can only become a
thought, not an object: the knower thinks the thought of God as his
own thought, and God vanishes. In faith, by contrast, the relation to
God is posited as a holy, inviolable fundamental relation; the
opposition between God and man is the opposition of spirit, the
antirational opposition, by which the spirit of faith distinguishes
itself from that of knowledge.

That the relation to God is a relation of spirit is known by this:
that all the determinations of finitude belonging to the relation are
permeated by infinity. It is pointed out, and with full justice, that
the relation to God is
% ---- printed p.472 (PDF 485) ----
\opage{472} originally a relation of conscience; but in religion
conscience is not enough, what matters essentially is the spirit.
That there are pronouncements in the conscience which do not conform
to man's self-will is certain; but these pronouncements of the holy
first acquire their true religious significance by being referred,
in faith, to the Holy One himself. The power of inwardness with which
this happens rests on the spirit. It is the Spirit of God in man that
makes the voice of conscience the voice of God; it is the spirit of
man in God that makes the divine pronouncements human. Under this
reciprocal determination the relation can be articulated to infinity:
understanding with God is self-understanding. No part of the rich
content of the fundamental religious pronouncements is by itself
either merely divine or merely human; if it is conceived in that way,
it is conceived spiritlessly. If your conscience tells you that you
ought, that you shall: then, in the religious sense, it is just as
spiritless to conceive the demand merely as an involuntary utterance
of your own moral nature as it is spiritless to let the silent
pronouncement, by an acoustic deception, sound as a voice that came
immediately from God: all the determinations decisive for the relation
to God are two-sided determinations of the spirit itself, divine-human
fundamental determinations.

What first and last gives the relation to God its actual stamp is,
according to what we have already seen from several sides, precisely
its peculiarity of being a relation of will. Only by virtue of the
will does the opposition become strong and the union inward. Over
against the Spirit of God the human will has freedom of choice; it
can decide for, and it can decide against, the admonitions of the
Spirit; the opposition between man's will and God's will shows itself
here at the same time as an opposition within man's own will, namely
an opposition between the selfish will and the will of the spirit.
% ---- printed p.473 (PDF 486) ----

\opage{473} Just as the selfish will can, in the general ethical sense,
set itself up against the demands of the good, so it can also, in the
religious sense, harden itself against the divine demands of the
Spirit; conversely, just as the selfishly natural will can, by virtue
of man's fundamental moral will, be sublated through personal
self-determination and transformed into a will moved by the Idea of
the good, so the natural self-will can, in the religious sense, be
sublated by the fundamental will and, through the influence of the
Spirit, transformed into will of the spirit. From this standpoint the
great problem of freedom and grace can be referred to two fundamental
determinations: 1) the determination of the relation between the
natural will and the will of the spirit, and 2) the determination of
the relation between the will of the spirit and the Spirit.

As regards the first relation, the solution of the problem rests on
the insight that the natural will and the will of the spirit can be
conceived both as two wills of different kinds and at the same time
as two manifestations, two forms of expression, of one and the same
human will. What has been brought to light in an earlier connection
about the opposition between the natural and the spiritual self holds
in particular of the will. If the opposition is stressed in its
immediacy, then the change is this: the natural will is thrust out,
and the will of the spirit is installed in its place: the old man
dies, and the new lives. If the opposition is reflected in its
fundamental unity, then the change is this: the same will that
selfishly withstood the Spirit transforms itself by self-determination
so essentially that it gives up its egoism and is moved by the Spirit.
According to the immediate opposition, then, the natural will falls
outside the life of the spirit, and can to that extent decide against
the Spirit \textit{(gratia est resistibilis)}.

As regards the second relation, the solution of the problem is
impossible so long as the opposition is held fast in mechanical
immediacy. It looks as though the selfish
% ---- printed p.474 (PDF 487) ----
\opage{474} self-will, and with it all self-determination, would have to
be annihilated in order that the Spirit of God with grace might gain
the power to create \textit{(gratia creans)} a new will, a new self, a
new man by rooting out the old. This is meaningless. It is not by
annihilating all self-determination and passively letting grace work
as it will that the will of the spirit can form itself; on the
contrary, what is required here on man's side is a heightened
self-determination, an act of will decisive for time and eternity.

Now when it is said that this determination of will is owed
exclusively to grace, that is true insofar as the will can become will
of the spirit only by letting itself be moved unconditionally by the
Spirit. Thus, as regards the ethical, the will can become moral only by
letting itself be moved by the good. Just as the good is something
irresistible for the good will, so the Spirit, and with it grace, is
irresistible for the will of the spirit \textit{(gratia est
irresistibilis)}. But that the will moved by the Spirit should not
move itself, that working grace should make the will passive, is an
error. On the contrary, a far greater self-activity is required of the
will of the spirit than can be found in the instinctive will; in other
words: the will of the spirit, in being moved by grace, is active in a
far stricter sense than the natural will moved immediately by
instinct and inclination.

Now since the old man, burdened with sin and selfishness, does not
vanish because the new is formed, neither will the will of the spirit
be able, at any stage of finite existence, to annihilate the selfish
so completely that a remnant of egoism does not always remain. In the
subdued instinctive will the will of the spirit has a secret enemy
lying in wait for an opportunity; it is never so sure of its victory
that it need not always keep itself awake and armed. Were the will of
the spirit, at a given stage, to settle down in ``the grace that
cannot be lost'' \textit{(gratia inamissibilis)}, then by
% ---- printed p.475 (PDF 488) ----
\opage{475} false security it would precisely expose itself to
succumbing to the instinctive will and losing grace.

It is, as can be seen from this, the divine-human fundamental
determination of the spirit on which the solution of all problems of
freedom rests. The opposition of the instinctive will \textit{(voluntas
naturalis)} and the will of the spirit is an opposition within the
spirit and precisely thereby a thoroughgoing opposition; but although
the opposed terms, in their immediate difference, confront each other
independently, the instinctive will falling \emph{outside}, the will of
the spirit \emph{inside} the relation to God, both are mastered by the
spirit: it is the destiny of the instinctive will to dissolve itself in
the divine-human self-determination by which it is transformed into
will of the spirit; it is conversely the condition of the will of the
spirit that it does not form itself alongside the instinctive will, but
on that foundation of individuality in which the natural will perishes.

To ascribe immediate independence to the instinctive will alongside
the will of the spirit is just as spiritless as to make the latter a
mere modification of the former. It is not superficial modifications,
it is essential reshapings, thoroughgoing transformations, that the
Spirit brings about. These transformations corresponding to
self-determination proceed no more from the divine will
\textit{(gratia)} by itself than from the human will by itself; they
proceed from the principial unity of opposition of the divine-human
will of the spirit. To speak of the natural will's cooperation with
divine grace \textit{(Synergisme)} is meaningless, since the will,
insofar as it remains in its naturalness, can only counteract the will
of the spirit. To assume a finite reciprocal action between the will
and grace is spiritless, since the self-determination of the will of
the spirit at no point falls outside or alongside the determining
grace%
\footnote{``If the Spirit were actually to move the will quite
immediately, in the manner of instinct, then it would have to begin
the liberation of the will by making it unfree; what a contradiction!
The will's first act of freedom, the fruit of its very first
inspiration, is precisely the heightened self-determination with which
it inwardly longs for, and, in conflict with instinct, surrenders
itself to, the influences of the Spirit. Then, as these influences are
received and appropriated by the will which frees itself through the
appropriation, the relation by no means becomes this, that the will,
in self-separation from the Spirit, should have any power for any
independent activity whatsoever''; --- that is to say, as regards the
life of the spirit, for that by self-separation from the Spirit it can,
as selfish will \textit{(vol.\ naturalis)}, again determine itself
according to instinct goes without saying --- ``on the contrary! it
can move only according to the Spirit, in being moved with freedom by
the Spirit. The moving power of freedom is solely and entirely the
power of the Spirit; but it depends on the will itself whether it will
appropriate this power in inwardness, in true originality.'' Om
Hindringer og Betingelser for det aandelige Liv i Nutiden. S.\ 11--12.

Indeed, so divine-human are the fundamental determinations of the
spirit that the will's self-determination, by which the power of the
Spirit is appropriated, is just as much an effect of the divine as of
% print: „Virkuing“ for „Virkning“ (turned n)
the human spirit; but the divine-human effect presupposes
self-determination in man, energetically decisive self-determination.};
% ---- printed p.476 (PDF 489) ----
\opage{476} the reciprocal action is infinite: only by proceeding out
of and returning into the infinite do the finite determinations of will
become divine-human determinations of the spirit.

From the standpoint of the finite and yet infinite reciprocal relation
between the Spirit of God and the spirit of man, the extremes touched
on above can now also be illuminated. Reciprocal action presupposes
that there is influence on man from God's side, and that the activity
taking place in man's spirit in turn conditions and determines the
influence of the divine Spirit. From this the following three
propositions can be derived: 1) in the religious sense no true human
activity of spirit can take place except on the presupposition of
divine cooperation; 2) no influence can be conceived from the divine
side that could either abolish or disturb the laws that hold for
% ---- printed p.477 (PDF 490) ----
\opage{477} the self-development of the human spirit; 3) the limit of
the heightening of the life of the spirit in human consciousness to
which divine cooperation can go without abolishing the laws of
selfhood is a Mystery.

With these principles in view we shall now shed light on the offense
taken at the extremes. The assertion that hidden miracles are only
imaginings unless they are expressly accompanied and made demonstrable
by revealed ones is, from the standpoint of revelation, meaningless.
Preservation and governance are hidden miracles; but they are
by no means made demonstrable by outward miracles. The relation to God
itself is grounded, according to the principle of faith, on a hidden
miracle; but it does not follow that this relation should constantly
be accompanied by outward miraculous facts. The hidden miracle, which
forms the foundation of the given order of things, is perennial; the
revealed miracle, by contrast, whose distinguishing mark is that it
forms an exception to the existing order of things, cannot possibly be
perennial.

Now as regards the first extreme, the miracle of inspiration (Acts 2),
it is certainly an altogether incomprehensible phenomenon; but as a
miracle it is indeed supposed to be so: a comprehensible miracle is a
contradiction. By its conspicuous incomprehensibility, its improbability
repellent in the very highest degree, this miracle stands quite in its
proper place in sacred history, insofar as it stands as an ineffaceable
mark, an abiding, speaking testimony that the Spirit of the Gospel,
with ``tongues of fire,'' is and will be heterogeneous, absolutely
heterogeneous with that of science. Whether the miracle of inspiration,
when the historical account is taken exactly according to the letter,
perhaps contains features which the critique of knowledge could use to
disprove its actuality is a question by itself, the treatment of which
belongs to sacred history.
% ---- printed p.478 (PDF 491) ----

\opage{478} The second extreme, which corresponds to the negative
experiences of the present age, receives a psychological illumination.
It is quite characteristic that men of deep sense and thorough thought,
men who with the torch of science would gladly, were it possible,
penetrate into the abysmally dark depths of revealed religion, begin
straightway by misinterpreting the relation to God itself, in that they
begin by misinterpreting the Spirit of revelation. ``The fundamental
form,'' it is said, ``under which the divine Spirit accomplishes man's
conversion and redemption is that of \emph{genius}. Its distinguishing
mark is pure, self-sacrificing enthusiasm, creative power. Within the
universal fundamental form of genius there is an essential difference
between scientific and artistic genius on the one side and, on the
other, the genius that is the bearer of moral and religious Ideas.
Here the Ideas are directed toward the will, the enthusiasm has the
highest intensity, it is simple, unshakable certainty of the divinity
of the revelation imparted, appeal to the divine authority. But
nevertheless the `inspiration and prophecy' proceeding from such
genius is nothing other than `moral-religious enlightenment'; genius
is only the first proclaimer and awakener of what rests as something
eternal in the ground of humanity\footnote{Dr.\ Karl Schwarz. Den
nyeste Theologies Historie. Ovrst.\ Kbhvn.\ 1869. S.\ 365.
Characterization of J.\ H.\ Fichte's ``speculative theology.''}.''

In this conception two essential points are to be held fast. The first
is that the relation between the Spirit of God and the spirit of man
is a relation of infinity. Self-consciousness at the stage of
reflection cannot place itself immediately over against God as self
against self, cannot say: this I said to God, and this God said to me;
thus I spoke to the Spirit, and then the Spirit answered me with these
words, etc. Where finitude in this conscious manner seeks to displace
infinity, the relation is
% ---- printed p.479 (PDF 492) ----
\opage{479} spiritless. The second is that the divine pronouncements
are not given in empty formal arbitrariness, but bear the stamp of
essentiality, in that they rise up out of the ground of the mind with
an irrefusable claim of conscience. But when this is granted, the
one-sidedness of the whole view is also easy to demonstrate. ``The
religious Ideas are directed toward the will'': what can that mean
other than that the religious Ideas are, in the absolute sense, Ideas
of will? That they are Ideas of will cannot possibly mean, then, that
from the divine side they are merely Ideas of instinct; in that case
the boundary between the two absolutely heterogeneous spheres of
spirit, that of knowledge and that of faith, is naturally easy to
efface. For of the Ideas of science and of art it certainly holds that
they announce themselves in consciousness with divine instinct, but
by no means as declarations of God's holy will.

The self-articulating reciprocal relation is in faith a relation of
will and, in this its peculiar character, an actual relation of
spirit. It is the human will, inspired by God himself, that in every
decisive question of conscience is determined to determine what God's
will is. The unalterable demand of faith is that the pronouncement of
conscience in every given case be referred to the divine Self, in
order to be conceived and appropriated as the very pronouncement of
the divine will; the irrefusable presupposition of religion is that
he who has unconditionally placed his cause in God's hand can also,
through conscience, learn what God's will is: the divine Spirit is
conceived expressly as self-determining; each of its
self-determinations announces itself from the depths with a holy
impression, a pressure upon the conscience. But this happens only in
faith; it is the function of faith as faith that is spoken of here.
The whole thing therefore takes place not merely in man's interior,
but also in and through consciousness's own self-determining, hence in
and through the human spirit itself. Thus
% ---- printed p.480 (PDF 493) ----
\opage{480} finitude is taken up in a spiritual manner into infinity,
and the opposition in the infinitely articulated reciprocal
determination is just as thoroughgoing as the unity.

In the relation to God, expressly determined as a relation of spirit,
there is here no more anything human by itself outside the divine than
anything divine by itself outside or alongside the human: the
articulated determinations of the spirit are divine-human fundamental
determinations.

\greekrun{γ}{Trinitarian Determinations of Form.}

While supranaturalism, when it does appear so resolutely as to be
willing to be recognized for what it is, always suffers from the
difficulty that its skeleton is too rigid and its joints too stiff,
pantheism, speculative rationalism, has on the other hand another
innate weakness, namely that its organism is so mollusk-like that limbs
and joints flow over into one another. When the talk is of the
relation between fixed and fluid determinations in the categories of
existence, it makes an essential difference whether the standpoint is
chosen in religion or in science.

If the realm of nature is considered from the standpoint of science,
then the fundamental qualities of matter, the physical and chemical
properties of the particular substances, are to be conceived primarily
as empirically given determinacies. Elements such as oxygen, hydrogen,
carbon, etc., each have their chemical peculiarities; it is of no use
to brood over why or whence the carbon atom has come by just these
properties, the hydrogen atom by those: the various properties hold
with necessity and yet look like pure contingencies. If one now
considers how these substances fit one another as if by a
\textit{harmonia præstabilita}, what a multitude of elementary
formations (acids, bases, salts, double salts, and so forth) arise
from
% ---- printed p.481 (PDF 494) ----
\opage{481} their combinations, and what role these formations play
not only in inorganic but also in organic nature: then one receives an
impression of objective reason, of divine wisdom in the world. In the
objective reason, the divine wisdom, is seen a system of fixed and
fluid determinations, on whose law-bound, articulating activity the
natural order of things rests.


From the standpoint of religion, that is, with express regard to God
and the Godhead, a judgment must now be passed on the reason of the world:
whence this wisdom comes, and how the matter really stands with the
fixed and fluid, indwelling and overreaching determinations in its
thought-rich content. That the scholastic supranaturalist must come into
conflict with the doctrine of nature is evident. Insofar as he appeals
to God, he of course has an easy enough time explaining wisdom to himself. But
how do things look with the emergence of this wisdom in nature? It is
God himself who, by virtue of unconditioned acts of will, has fixed all
fundamental determinations and arranged everything from the beginning. Whence,
pray, does oxygen have its peculiar chemical character? From God's will! And
hydrogen its? From God's will. Indeed, it is precisely an eloquent proof of
God's creative wisdom that he has formed oxygen and hydrogen such that
they can combine and form water; he evidently thought of
water when he formed oxygen and hydrogen. Without actually introducing a
Solomonic personification of wisdom, the supranaturalist nevertheless
constantly brings before his view only the image of a represented wisdom's
supernaturalness; but the real concept of wisdom, the reason indwelling in
nature itself, he does not cognize. And why not? Because in a
spiritless manner he conceives all fixed determinations of nature, not as
determinations of essence in nature, but as determinations of will over
nature.
% ---- printed p.482 (PDF 495) ----

\opage{482} It is otherwise with the pantheist. For him the categories of reason hold as
original forms of nature. They are concepts of nature which immediately
state the nature of things; to that extent he can keep to physics.
But when the doctrine of reason is to be developed further, then comes the
difficulty. In their factual empirical existence the
categories of nature have a contingent appearance. The essence of oxygen cannot
be derived from that of hydrogen, any more than that of hydrogen from that of oxygen. Certainly
they suit each other, and water arises through their union. But
they are, after all, united merely \emph{so that} water \emph{can}, but not \emph{in order that}
it \emph{shall}, be formed. Where now is the wisdom? On this the pantheist,
so long as the question concerns determinate particulars, individual determinate
substances, individual determinate formations, can communicate nothing but what
stands in chemistry: all determinate determinations are empirical. And yet
it is surely not his task to give us a little abstract of chemistry;
his task is to present reason in its originality as
reason, as divine wisdom. So he must then swing himself up above
the empirical, that is to say: dissolve all determinate determinations as
moments in the infinite totality. The parts are now thus because
the whole is thus, and the whole in turn is thus because the parts
are thus. Now everything is fluid. In a conceptual system of abstracted
forms of nature, which are passed off as a priori, and which with a little dialectic
can easily be brought to pass over into one another and be reflected in
one another as means in ends, the immanent wisdom is to be
revealed. But it is a wisdom standing on such weak feet that anyone
who wishes can blow it over\footnote{Cf.\ ``Philosophisk
Propædeutik i Grundtræk''. Kbhvn.\ 1858. S.\ 189.}. The weakness is that
the omnipotently determining Self is lacking.

Just as the fluid determinations in the pure doctrine of
% ---- printed p.483 (PDF 496) ----
\opage{483} reason correspond to universal selfhood, so the fixed,
decisive fundamental determinations in the real concept of wisdom correspond to the
divine self-determination. Then comes ``speculative
theology'' and makes yet another attempt with wisdom. ``The world,'' it
says, ``is not merely a system of members mutually corresponding to one another,
but also an ascending series of stages of ends. Each individual
member is in itself an end and at the same time a means for another, the result of a
lower, the starting point of a higher series of development. And here there appears
the striking semblance that what has been attained, realized,
although a product of what went before, is nevertheless at the same time that for whose
sake alone this exists. Thus what is not yet operates
in advance of its being; the end is at the same time cause, but posited as
effect; the means is likewise effect, but posited as cause. This
operating of what is not yet in advance of its being, this reversal of
end and means, calls for a resolution of the contradiction that sets
in. The ends operate in advance in their means. Of course
the means do not properly bring about the end, nor does the end itself operate
in the means; for the end does not yet exist. So a
third is then required, which bends each means toward the end and lets it operate
through the means. This end-positing activity of the Absolute
resolves the contradiction appearing in time; it sublates
the temporal difference of the members that fall apart from one another, means and ends,
by seeing through their eternal unity. A world-order of
ends can only be thought without contradiction when an operative and
willing Absolute pervades them and is present in
them''\footnote{K.\ Schwarz. Anfr.\ Skr.\ S.\ 359--60.}.

The ``rational theism of speculative theology'' certainly borders
very closely on the religious; but the specu\-%
% ---- printed p.484 (PDF 497) ----
\opage{484}lative Godhead is not the God of religion and can never become so. ``The
knowing and willing Absolute,'' whose end-positing activity is to
resolve the contradictions of science, is an ambiguous being.

If ``immanence'' is held fast with logical rigor, then the
opposition of means and ends loses its tension: the means is goal,
the goal means; the determinations flow over into one another; the
end-positing activity is only a necessary unfolding of the
Absolute's necessary content of reason. The knowing Absolute has only
a willing, but no will. Its knowledge is only a knowledge of the
necessity of the ends and the means, a knowledge that it itself is not
able at any point to bring the very slightest change
into either the end or the means, a thorough knowledge that its
end-positing will is a potency as powerless as it is superfluous.

If ``transcendence'' is taken seriously, if the Absolute is actually to
be ``a third, which bends each means toward the end, posits
the end before it yet exists,'' then the emphasis falls on the
omnipotently willing Self, and the antirational blows ``speculative
theology'' sky-high.

The rational is that ends and means in the infinite system must
correspond to one another, so that the one does not, through thoughtless
arbitrariness, do damage to the other. There must be ends, but
which ends? There must be means, but which means? There must be
connection, but which connection? When these questions are to be
answered, when it is to be decided precisely why \emph{these},
just these ends corresponding to \emph{these} means, are posited, and not
others, although \emph{other} ends with \emph{other} means must, according to
the doctrine of ends, always have been possible, then the
antirational enters with the fixed determinations of will, and the fixed,
decisive determinations of will must --- we repeat it --- in the
science of reason
% ---- printed p.485 (PDF 498) ----
\opage{485} necessarily appear as utterances of divine
arbitrariness.

Only on one condition can human knowledge be reconciled with the arbitrariness
grounded in God's wisdom; this one inviolable condition is
faith. Faith does not exclude knowledge in the sense that it should abolish
reason and push science out of consciousness. Faith and knowledge
form the deepest opposition in the cognizing spirit of man; but
spirit, which posits and masters all oppositions, masters this one too.
In the spirit of science the rational, and with it ``immanence,'' is the
matter of principle; but in the spirit of faith the antirational with
``transcendence'' is what is first and last decisive.

According to the articulating reciprocal determination of the rational and the
overreaching antirational, the fundamental determinations in the Self of the Spirit must
be trinitarian. The trinitarian unity of opposition means something
other and more than the reflected identity of opposite sides or the
fluid transition of fluid ``moments.'' It is contentful
totalities, comprehensive world-systems, that are in question here. If
God in himself, the omnipotently creating Self, is set against the created world,
then it is seen at once how the thoughts of faith must work their way through
opposite totalities. The divine Self is not empty; in God as
God, in God's interior, all worlds are anticipated in a totality of ideal
determinations of power transparent to knowledge. In its
opposition to God the world, created by God, is not God-forsaken; its
totalities of means and ends are revelations, real
determinations of power, of the universal substantial selfhood: it has an
indwelling Godhead; it is inspirited by God. If the Spirit is now to be seen to master
this original fundamental opposition teleologically as the Spirit of wisdom,
then this must be seen in trinitarian forms.

The first thing is that the end is posited, not as an end in
general, but as \emph{this} determinate end;
% ---- printed p.486 (PDF 499) ----
\opage{486} the end-positing activity is, in God's interior, ideal and spiritual. In
the inner, ideal world, means and ends are transparent to
God's all-pervading knowledge, namely in the sense that they are
purely \emph{possible} means for purely \emph{possible}
ends\footnote{How the matter stands logically with these
possibilities and their relation to necessity is shown in
Grundideernes Logik II S.\ 222--293.} --- an infinite
system of concepts in God's inner selfhood. That an end is now expressly
posited as \emph{this determinate} one comes about through a system of the
selfhood-determinations reflected in themselves being transformed into peculiar
self-determinations, fixed, inviolable determinations of will. The
end thus determined now requires an external world, a
\emph{world of means}.

The external world is, according to its character as external, a
world of means with immanent ends. Just as selfhood, through
the inner world, is merged in the determining Self --- the Godhead in
God --- so now in the outer world the determining Self,
\emph{the will}, is at every point merged in the selfhood of the plan, the universal
willing --- God in the Godhead\footnote{That the relation between the system in
the inner world (the ideal system) and the system in the outer world
(the real system) is by no means determined by a simple parallelism can
be seen from Grundideernes Logik e.\ g.\ I.\ S.\ 449--456.}.
% print: no full stop after the footnote mark

The system of the external world of means and the internal world of ends is in
the Spirit itself one and the same system; for the inner world of ends has
all the means within it as moments of knowledge; conversely, the world of means
has within it, in its rational connection, the immanent ends as
substantial motives; the teleological reciprocal determination of
the inner, \emph{ideal} and the outer, \emph{real} world-plan is in
the Spirit as Spirit a totalized unity of opposite totalities, a
trinitarian unity.
% ---- printed p.487 (PDF 500) ----

\opage{487} \runhead{c) The Personality of the Spirit.}
\parmark{33}

It is a quite characteristic feature of objectively dogmatic church doctrine
that, while it has never found any difficulty in
determining the Father and the Son as divine persons, it has always
found itself in a certain embarrassment when it was really to express itself
clearly and determinately about the personality of the Spirit.

As we now undertake to investigate whence the difficulty arises and
how it is removed, we must first of all draw attention to what
follows directly from the foregoing: that the Spirit, conceived in
general as spirit, places the relation between self and selfhood in
a quite peculiar light. It begins with the self in the Father
and the self in the Son being determined as fixed, unshakable centers,
whereby an opposition of two divine extremes is formed. As the
mediating being of the extremes, the Spirit then effects the unity of the opposition by
assuming the form of selfhood, and by its infinite formative activity
mediating between self and self. From this standpoint it looks as
if the self, the inner kernel of personality, were to be sought only in the Father
and the Son, while the Spirit, the mediating being with the fluid
determinations of transition, would have to content itself with the form of universal
selfhood without a fixed personal center of its own.

By developing the consequences contained in these presuppositions,
however, we have convinced ourselves that the matter stands
otherwise. So far is it from being the case that the Spirit can be separated, either with its ground in
the universal from the self in the Father, or with its ground in
the individual from the self in the Son, that the self-doubling whereby
the Father distinguishes himself from the Son rests precisely on the Spir\-%
% ---- printed p.488 (PDF 501) ----
\opage{488}it: the fundamental act of self-doubling is an act of the self in the Spirit. Seen
from this side, the matter then takes on quite another appearance; for now it looks
as if the self were posited solely in the Spirit, as if it were the Spirit
that, in the development of its own essence, solely for the sake of
self-opposition, and thus for the sake of self-development, so to speak,
allotted to the Father and the Son each a self of his own, so that by the mediation of the extremes
in the unity of the opposition it might take up the separated members in
full forms and hold fast to itself alone. Since now the latter
conception is evidently just as one-sided as the former, it becomes
a task for the philosophy of religion to show how it follows precisely
from the concept of the Spirit that the divine personality must have
its threefold point of unity, its living center, in a threefold
personal unity of self.

\greekrun{α}{The Spirit of the Father.}

That the Father, the omnipotent Self who brings forth all things, cannot possibly be
personal for himself in separation from the Spirit has already been shown, and
can be seen expressly in the concept of omnipotence. Omnipotence is: 1)
self-power, God's power; 2) power in another, the reality of the world; and 3) power
over another, infinite over-power. Omnipotence is trinitarian. It is then
just as absurd if anyone were to say that it is the Father who has
taken the self from the Spirit, as if anyone were to maintain that it is,
conversely, the Spirit that has taken the self from the Father. The Father and
the Spirit have, in this first original concentration, one and the same
self, one and the same infinite selfhood. Why then do we call this self
the Father's, not the Spirit's? It lies in the peculiar character of the paternal
personality, in the original opposition of the omnipotent God
to the world as of the Creator to the whole creation; it shines forth
from the Father's personal works, from creation, from preservation, from
governance.
% ---- printed p.489 (PDF 502) ----

\opage{489} Creation is trinitarian; for the Creator is, as the omnipotent Self,
exalted above the world; he is, as the universal, substantial selfhood, in
the world; he is both at once by means of the Spirit. Only by
virtue of the omnipotent unity of opposition of self and selfhood can
the Father be Creator, but the unity of the opposition is personal only in
the Spirit. If the Spirit is thought away, the concept of creation is a spiritless
concept. The creating God then sinks down to being a finitely
represented person, a mystical individual without individuality.
The divine works of creation would in this light come to
appear as finite acts of power, utterances of arbitrariness
of a finite will. And why? Because the universal selfhood, bearing the content of
the totalities, falls away with the Spirit. Only in the Spirit and
by means of the Spirit can the omnipotent, primordially paternal Self in unconditioned
freedom raise itself above the inner system of coherent
possibilities of worlds; only in the Spirit and by means of the Spirit can the
creating will sink itself into the unity of these possibilities as into its
own selfhood, can with the self-determining power of its will breathe into
these possibilities the plan of wisdom in fundamental determinations fixed in essence, and
through these fundamental determinations, for the real revelation of power,
release the finite elements of reality to natural self-development.
The unity of the natural development of things in the substantial
selfhood of the Godhead and of the acts of will in plan and disposition whereby
the natural development is determined is, in the personal Self of the Father, a
unity of opposition, a trinitarian unity, by means of the Spirit.

Preservation, on account of its unity of opposition with creation, is in a
special sense trinitarian. It is the reciprocal determination of self and
selfhood, of the punctual will and the universal willing, on which it
again depends. If we disregard the Spirit with its God-and-world
fundamental determinations, it will be impossible to form
% ---- printed p.490 (PDF 503) ----
\opage{490} any sound concept of the seemingly abrupt transition from
creating activity to preserving rest. In the Spirit, by means of
the Spirit, the opposition is explicable. In the extreme of divine
creative activity the incipient activity of nature is already to be
understood; the universal willing furnishes the ground for the
determining will. In the extremes of preservation creative activity is to be
understood; for the will's infinite self-determining does not, after all, cease
because in its inner operation it holds back the particular acts of will:
the will is only resting because it personally keeps itself in
rest, the active rest that bears the world. Preservation is not
simply a continued creation, but a creation converted into natural development, i.e.: into divine
willing, an extreme which has its opposite within itself and is in
the Spirit trinitarian.

In a still more many-sided sense governance is trinitarian. From
the concept of governance it is seen expressly how the two extremes,
creation and preservation, mutually include each other; for governance
presupposes the natural development of the connection of things and events,
and to that extent coincides with preservation; but
governance also presupposes the active cooperation of the will, and from
this side coincides with creation. The unity of opposition of creation and
preservation which is posited by the Spirit in the governing will shows
clearly what it means that the Spirit is the Spirit of the Father, that the Father is
a personal God, that the personal God can act. In personal
action finitude is taken up, for the will operates punctually, now
here, now there; now thus, now again otherwise. But the
divine action proceeds from inner infinity, does not break
the essential connection at any point: the punctual will is
everywhere in essential understanding with the universal willing. Thus the
Spirit operates in the Father's essence; its infinite self-operation is on all sides
directed toward uniting all personal activ\-%
% ---- printed p.491 (PDF 504) ----
\opage{491}ity in the paternal Self; the personal Spirit is the Spirit of the Father:
\emph{the Spirit proceeds from the Father}.

\greekrun{β}{The Spirit of the Son.}

The God-man, the Logos-Self become man, is in his opposition to
the Father one with the Father; this unity of opposition is trinitarianly
personal in the Spirit, corresponding to the divine-human
fundamental determinations in the Self of the Spirit. But the Self of the Spirit does not in this
connection step forth specially over against the Son. On the contrary! its infinite
self-operation here aims at uniting all personal activity in
the personality of the Son. This center is individual: the Spirit
operates on the ground of the individual. The Son's personal
works, his law-fulfilling, redeeming and atoning activity,
are, as works of the Spirit in and through the Son himself, all trinitarian.

For the fulfillment of the law it is the self in humanly limited
individuality on the one side and the law, the universal, divine
selfhood, on the other side, that form the extremes. The unity of the
opposition now shows itself in this, that individuality expresses the universal,
that flesh and blood are subjected to the law, that the self freely gives itself
up into selfhood. It looks as if the latter extreme gained over-power
over the former. But --- this is the personal operation of the Spirit --- the
latter extreme is conversely taken up and transfigured in the former. It
is not the law that here gains over-power over individuality; it is,
on the contrary, individuality that gains mastery over the law. The law's
strictest demand is satisfied; but as the law is fulfilled, its
perfect fulfillment becomes its complete destruction: it goes
under in the personality of the Son --- by means of the Spirit. Selfhood
is not violated, is not rejected by the Son; no, what is personal in the Spirit
shows itself precisely
% ---- printed p.492 (PDF 505) ----
\opage{492} in this, that the extremes are preserved in their unity, that the sublation of the law is
its fulfillment.
% print: Eenbed (for Eenhed, `unity')

The infinite reciprocal determination of suffering and operating which
characterizes the fulfillment of the law in turn forms two new extremes, of
which the first reduces the personal self-activity to a moment
of the suffering, while the second posits the personal suffering as a
moment of the self-continuing personal activity; this is, in
the Spirit itself, the opposition of atonement and redemption. Judging by the one
extreme, one might be tempted to conceive the atoning
suffering as if it concerned merely the individuality of the sufferer and, as something
historically concluded, belonged to the past, while the redeeming
activity concerned the universal and therefore extended through all
times; but how spiritless such a separation is has been shown in
the foregoing. It is precisely the distinguishing mark of all the functions of the
Spirit personal in the Son that the individual is point by point
transfigured in the universal, and the latter in the former. Where the Spirit of the Son is to
operate, in the word of baptism, in the words of institution of the Supper, the Son himself must
personally cooperate. But that the Son cooperates means that he,
present as individual personality, by means of the Spirit, is
spiritually-bodily present in the word.

In order to further the spiritual conception, people have sought to push back
the representation of the bodily presence; this is a
preposterous procedure, for when the Spirit, the Spirit that masters
the oppositions and unites the extremes, is thought away, an
individual personality cannot possibly be present in more than one place
at a time. To believe in the spiritual presence of the individual personality
and deny the bodily presence is nonsense; for
the individual personality is, according to its concept, spiritual-bodily.

Yet --- we will not keep silent about it --- precisely where the talk is of the Son's
personal transfiguration in the Spirit, a trans\-%
% ---- printed p.493 (PDF 506) ----
\opage{493}figuration on which the possibility of his sacramental
multilocal presence rests, a problem arises so hard that it
at first glance looks as if it would clinch the paradox and
put the absurd in the place of the Mystery. ``For the possibility that Christ's
body could be essentially present in the Supper in the many
places where it is celebrated, Luther appealed in the first place
in general to the divine omnipotence, for which nothing is
impossible; then he let, more determinately \dots\ Christ, who also according to
his human nature sits at the right hand of God, be present together with this right hand
in all creatures, as he expresses it, in stone, in
fire, in water, even in the rope, consequently also in the bread and
wine of the sacrament. When this now seemed to be only Amalric of Bena's
materialistic pantheism and panchristianism, and in no case
to give the elements of the Supper a preference over other bodies,
Luther made the further distinction that Christ, although present in
all things, nevertheless lets himself be grasped and enjoyed only in those where he wills and
expressly declares it, as in the Supper; but by this one merely came
back by a detour to the appeal to the divine
omnipotence''\footnote{Strauss. Troesl.\ II.\ S.\ 479.}.

It is in the present context not specially with regard to the
sacrament, but with regard to the Spirit, that the problem is brought out. Two
main propositions seem here to stand in irreconcilable conflict with each other.


It is not particularly with regard to the sacrament, but with regard
to the Spirit, that the problem is brought forward in the present
connection. Two main propositions seem here to stand in irreconcilable
conflict with one another. The one is that Christ's individuality is
incorruptible; the other is that the personal Christ can be present
with his body and blood wherever his Supper is held. In order to hold
fast to the latter proposition the Lutherans essentially dissolved the
former; for in the doctrine of the ubiquity of Christ's body his
individuality is, after all, completely dissolved. In order to hold
fast to the former, the Calvinists must deny
% ---- printed p.494 (PDF 507) ----
\opage{494} the latter. Christ cannot be bodily present in the Supper, for his
individual body is in heaven. Now so long as one continues, on the one
side, to hold fast to the bodily presence in the sacrament in such a
way that it must necessarily dissolve the individuality, while, on the
other side, one limits the individuality in so external a manner that
the spiritual-bodily presence becomes impossible, the dispute does not
turn on a holy mystery, in which finite conceptual contradictions can
be resolved, but on an irrefutable paradox, an absurdity, whose
consequences strike deadly blows on both sides.

It is the fundamental opposition between universal and individual
personality, more precisely: between the Father and the Son, to which
we must go back. The mediation of the extremes, the unity of the
opposition, must be sought in the Spirit; for the truest, richest,
fullest presence of personality is its presence in the Spirit. The
personal spirit does not exclude bodiliness; on the contrary!
Bodiliness is everywhere the means and the medium of passage for its
self-active revelation. But universal personality stands in a
different relation to bodiliness than does individual personality:
the Father's relation to the world of bodies is qualitatively
different from the Son's. The Father, the Almighty, is incorporeal;
his inner personal self-power reveals itself in this, that through
the Spirit he has all bodiliness in his power, is omnipresent in the
bodily world. The Son, the revealer of the Almighty, is clothed in
bodily individuality and has not, in the glorification --- ``at the
right hand of God'' --- laid aside his individual bodiliness. Glorified
bodiliness --- this is the Mystery --- differs from corruptible
bodiliness in this, that as spiritual bodiliness it can express, in an
individual manner, the concentrations of the will in a plurality of
relative centers. Christ is not, as Luther says, --- ``present in all
creatures''; here there is evidently a
% ---- printed p.495 (PDF 508) ----
\opage{495} double error. In the first place, it is overlooked that the Father's
presence in ``all creatures'' is not an immediate sensuous being-present.
In the second place, the thoroughgoing opposition between the Father's
and the Son's personal relation to bodiliness is overlooked. But what
they have in common is that glorified bodiliness serves the Son's
personal will just as unconditionally as the bodiliness of the whole
world serves the Father's; and Luther is to that extent right in what
is, religiously, the main thing, namely that Christ \emph{can let
himself} be grasped and received in his sacrament: the sacramental
presence in many places belongs originally to the Atoner's
divine-human life in the Spirit.

That we cannot aim at any objective scientific comprehension either of
glorified bodiliness or of the actual multiplication of individual
presence goes without saying. On the other hand, we have expressly
aimed at so reshaping the concepts of faith that belong here that the
paradox can vanish and the Mystery be acknowledged. The
spiritual-bodily individual presence of which we are speaking here is
not sensuously objective. Of sensuously objective individual presence
it holds unconditionally that when the individual is in one place, it
is not in another; but the glorified Savior's spiritual-bodily
presence is a personal presence in the Spirit.

Personal presence is one with the working of the personal will; this
holds just as well of the Father's omnipresence in the world as of the
Son's individual presence in his sacrament. It is not a finite
distribution of divine powers, not the transfer of an individual from
place to place in space; no, it is God's personal presence in the
Spirit, concentrations of the divine Self, that is in question. That
these concentrations must be different according as they correspond
to the universal or the individual character of personality involves
no absurdity: the Spirit in the Father can to that
% ---- printed p.496 (PDF 509) ----
\opage{496} extent be opposed to the Spirit in the Son. But just as the
opposition of the Father and the Son is a divine self-opposition in
the Spirit, so too is the unity of essence of the Father and the Son a
divine self-unity in the Spirit. In this self-unity the Spirit is
personally one with, and precisely thereby personally distinct from,
both: \emph{the Spirit proceeds from the Father through the Son}.

\greekrun{γ}{The Holy Spirit.}

How faith in the Spirit that proceeds from the Father and the Son is
related, according to the principle of faith, to faith in the Father
and the Son themselves may well be discerned from what has gone
before.

If anyone supposes that it comes easier to ``a reasonable man'' to
believe in the personality of the Father than in that of the Spirit,
then the philosophy of religion and negative critique, however much
they otherwise disagree, will now be able to unite in bringing him to
other thoughts. It does indeed look at first glance as if it were,
without further ado, a need of human reason to proceed from the
presupposition of a God in heaven, a supramundane personal God; it
looks as if the first article of faith were quite reasonable, and as
if one could best escape all the troubles with faith and knowledge by
keeping to it alone. Without immersing ourselves in the endless
disputes between deism and pantheism, we need only a decisive
religious emphasis on \emph{God's omnipresence}, and the
counter-rational\footnote{One must constantly bear in mind the
difference between the counter-rational, the irrational, and the
contrary-to-reason. The irrational, what is absurd in itself, must be
rejected; the contrary-to-reason must be corrected; the
counter-rational must be believed. The antirational is in
itself just as little in contradiction with the rational as, according
to ``the Logic of the Fundamental Ideas'', the antilogical is in
contradiction with the logical.} (the antirational) shows itself
immediately.
% ---- printed p.497 (PDF 510) ----

\opage{497} God, says religion, is present in his congregation. This presence
would be altogether unthinkable if God were not potentially present in
human life in general, and particularly in the life of the people to
which the members of the congregation naturally belong. But whereas
God's presence in universal human life is a hidden presence, a
presence of immanence, in which selfhood veils the self, God's
presence in the congregation is, on the contrary, an actual, manifest,
personal presence, a presence which makes itself known particularly in
this, that it is not the idea of God in indeterminate generality, but
God's own holy Self, that stirs and works in the souls of believers.
``How, then, can the Father at one and the same time be both personally
in heaven, that is: exalted as God in himself above the whole world,
and yet on earth, in a determinate community, personally in the
congregation?'' It is possible only through the Spirit. ``Whence comes
the Spirit?'' It comes at one and the same time both from the depths
and from on high. From the depths! For it comes from the ground of
human nature, of souls, through the consciences; it comes up, then,
where it already was; it comes, as God himself, up out of his own
Godhead; it comes by a divine act of will: its coming is its
self-revelation. From on high! For the act of will by which the Spirit
reveals its holy presence is a personal expression of the act of will
by which it is \emph{sent} from the Father in heaven. That the Father
sends the Spirit means that he expressly \emph{wills} its revelation.
To immediate intuition it does, to be sure, look as if God himself
were in two places: in heaven and in the congregation. If this
representation is held fast in external finitude, we literally get
two Gods. Indeed, the finite representation goes further still. For
there are not only many congregations on earth, inasmuch as the one
congregation is distributed about in many places, but there are in each
% ---- printed p.498 (PDF 511) ----
\opage{498} congregation many members; in each heart, in each believing soul,
then, God as God, God's own Self, is present. The absurdity that seems
to arise from this stems solely from a spiritless conception. One looks
out into space and dreams of places and places, instead of believing in
the Spirit and understanding that the matter here is not places but
almighty acts of will, punctual concentrations of the divine Self.
Were we here to posit a separate God for every divine Self, we should
indeed get a polytheism that meant something: we should get countless
gods. No, the truth is that the one and same God, out from his eternal
center, his absolute point of unity, is able to posit an infinite
multiplicity of relative centers, in which, by virtue of an inner
overarching self-unity, he repeats his own divine Self. This is the
fundamental opposition in God's essence: on the one side the almighty
Self, exalted in power and majesty, the God who ``dwelleth in the light
which no man can approach unto''; on the other side the multiplicity of
the points of revelation in which the eternally working Self
unceasingly repeats itself.

It is these extremes that are mediated by the Spirit; without the
Spirit there would be personality in the former just as little as in
the latter. In the first extreme God's Spirit is personally the
Father's Spirit; in the last extreme it is personally the Spirit of
the congregation. Therefore the Spirit of the congregation is not a
universal spirit of selfhood, like the spirit of humanity or,
particularly, like the spirits of the peoples; no, the Spirit of the
congregation is the Spirit with God's own holy Self, the Spirit from
God --- from the Father --- the Spirit as ``God the Holy Spirit''.

Should anyone think that according to this explanation one could then
say that the Father himself is the Holy Spirit, and that the Holy
Spirit is the Father, the misinterpretation is easy to show; for the
unity of the Father and the Spirit rests expressly on the
% ---- printed p.499 (PDF 512) ----
\opage{499} opposition. If the personality of the Spirit is absorbed into the
Father, then God remains in his heaven and, as a deistic, otherworldly
God, cannot possibly be present in the congregation. If the
personality of the Father is absorbed into the Spirit, then God is
transformed into Godhead, and the Holy Spirit into an ethical spirit
of reason, an immanent spirit of conscience without a divine Self. It
is seen, then, how the Father and the Spirit bear witness for each
other: the personality of the Spirit confirms that of the Father, and
the Father's in turn that of the Spirit.

The same thing repeats itself with the Spirit's relation to the Son.
The difference from the Father lies in the new center. Concentrated in
the personal individuality of the Son, the divine Self furnishes the
intuitable individual-personal presupposition for the Spirit. Christ
himself is present in his congregation in that the Spirit of Christ is
present. Should anyone think that, if the glorified Christ is actually
present himself, the personality of the Spirit must be superfluous,
or, conversely, if the Spirit is present in its own divine person, the
presence of Christ must be superfluous: then it is easily seen how
loose, indeed meaningless, such an opinion is. If the personality of
the Spirit is absorbed into the Son, then Christ must remain seated at
the right hand of the Father and cannot possibly come down to the
congregation. If, conversely, the personality of the Son is absorbed
into the Spirit, then Christ is transformed into a general idea of
Christ, and the Spirit becomes an Idea-spirit without a divine Self.
If the Spirit loses its personality, then not only the Son but, as we
have seen, the Father too loses all personality. The third article
thus has, as regards faith in the divine personality, the same
fundamental character as the second and the first.

Thus faith in the triune God furnishes the unshakable foundation of
the life of the congregation. ``Therefore'', says Grundtvig, ``the
statement concerning \emph{the unity of the Three} must express no more
than can be warranted by the congregation's
% ---- printed p.500 (PDF 513) ----
\opage{500} Confession of Faith and the word of baptism, so that if more
should lie therein, it is of no concern to the Christian
congregation'' \dots\ ``The same holds of the technical term
% print: no opening quotation mark before „Det Samme gjælder“; the resumed
%   Grundtvig quotation is closed at its end but never reopened.
``Trinity'' (the Greek \textit{trias} and the Latin \textit{trinitas}),
so that, if anything indefensible were found therein, it would be the
scribes' own affair and not the congregation's, which has only to
confess and defend what it wholly and fully professes at baptism to
believe and lets itself be baptized upon''.\footnote{Den christelige
Børnelærdom. Kjbhn.\ 1868. S.\ 208.}

\lettersub{B.}{The Works of the Spirit.}
\parmark{34}

Since the Spirit, in its relation to the Father and the Son, is,
according to what we have seen, neither a part of God nor an abstract
divine essentiality, but the whole God, it follows from this that
every work of the Spirit must be capable of being called both a work of
the Father and a work of the Son. Nonetheless there is here a call to
make the difference count. For the Spirit personally, as the Holy
Spirit, is in a special sense the God of souls, the God of inwardness
in the believers. The works of the Spirit are, with regard to this, to
be determined as the works of the God who is present in the
congregation. The influences that the Spirit exerts on the individual
as individual are then to be referred, partly mediately and partly
immediately, to the life of the congregation\footnote{That this point
of view does not stand in conflict with the Kierkegaardian insistence
on the category ``the individual'' will be understood when one
reflects on the difference between the philosophy of religion and
religious ethics.}. We
% ---- printed p.501 (PDF 514) ----
\opage{501} comprehend these works under the three points of view:
\emph{awakening, regeneration, sanctification}.

\runhead{a) Awakening: The Living Word.}
\parmark{35}

From the opposition between the natural man
(\textit{ἀνθρωπος ψυχικος}) and the spiritual man (\textit{ἀνθρωπος
πνευματικος}) it follows that the transition from unbelief to faith
must presuppose a deep stirring of conscience, a breakthrough of the
Spirit. By what psychic symptoms such a breakthrough is accompanied is
not the essential thing, since these may be highly various according
to the inner and outer circumstances of the individuals concerned;
the essential thing is that the Spirit comes in personal power; the
essential thing is that an awakening takes place. The awakening Spirit
may announce itself in the individual, may, so to speak, knock at the
door of the solitary one, separated from the congregation; but it does
not come as the solitary one --- if it comes in that way, it is not
the Spirit of Christ --- it comes from the congregation and calls to
the congregation; it does not come to the thinker in a system of
general thoughts --- if it comes in that way, it does not come with
awakening --- no, it comes with a personal, awakening word; the
awakening word is \emph{a living word}: thus it comes essentially by
the hearing of the word.

That the two concepts, the personal Spirit and the living word, must
essentially belong together can be seen in the following way. The
Spirit, which in itself, as spirit, is both invisible and inaudible,
could not possibly reveal itself if it could not make itself known in
the word. If it is now asked: in what does the word have its most
original expression, in written signs or in speech? then the answer
cannot be doubtful. That the sounding word is more original than the
written word shows itself
% ---- printed p.502 (PDF 515) ----
\opage{502} not only in this, that it is older, more natural, and so to that
extent also more popular than the written word; no, it lies in its
peculiar relation to the personality of the Spirit. As the spirit is,
so must the word be. For the spirit of science the written word will
not only be just as convenient but in many cases even far more
adequate than the oral word; so, for example, the mathematical sign
language: mathematics and ``the living word'' suit each other but
little. Are we then to say that the spirit of mathematics is a dead
spirit, and all mathematicians ``dead-biters''? By no means! The
spirit of mathematics is not a dead spirit but an objective one, the
impersonal spirit of exact knowledge, of cold reflections, of abstract
thoughts: the objective, impersonal spirit
% print: „Brng“ for „Brug“.
has no use for the personal word. That oral instruction is given in
the strict sciences happens solely for the sake of the individuals,
not for the sake of the sciences: the clear, objective language of the
sciences is essentially a written language.

But the Holy Spirit of faith is absolutely heterogeneous with the
spirit of science; it lies in its personal essence to demand personal
communication. The more weight is laid on the personality of the
Spirit, the more weight must be laid on the oral word. In the oral
word the speaker is present; but the written word presupposes that he
who wrote it is absent. In the speaker the Spirit can stir, live, and
work; for the speaker himself is a living being. But the written word
is not a living being. The written word, though only a shadow-outline
of the absent spirit, can indeed awaken thoughts and feelings and
living representations; this holds especially of works of poetry: one
thinks of a Homer, a Dante, a Shakespeare, and the imagination already
comes alive at the mere thought. But what sort of living is that? It
is a living in beautiful illusions, hence an imagined living; but an
imagined living is not personal. The
% ---- printed p.503 (PDF 516) ----
\opage{503} awakening of the personal life in the Spirit presupposes a strong
personal impression of the Spirit; but a strong personal impression of
the Spirit presupposes the Spirit's personal working in a living
personality\footnote{Those who assure us that they cannot understand
what the living word is assure us by the same token that they cannot
understand what a personal spirit is either. The oral word is not
without further ado a living word, but conversely. To keep exclusively
to the opposition between the oral and the written word will not do;
for there really are those who talk big and become lively without the
word thereby becoming living; and there are, again, those whose written
word is wonderfully awakening, for example S.\ Kierkegaard's.}. To
read oneself into an awakening in a written word is impossible. One of
two things: either it remains an awakening in thoughts,
representations, and passing moods, hence a sham awakening without the
power of actuality, or else the Spirit itself seizes the human being
from within in an extraordinary way, and then the written word is only
an accidental occasion, not an efficient cause; for this stands fast:
nothing can come forth in the effect but what has been in the cause.
The Protestant deification of Scripture is a denial in principle of
the holy Spirit of personality. That this deification of Scripture did
not vanish with critique's dissolution of the old dogma of verbal
inspiration has its ground in this, that the written word is the
objective word of science; so the personal Spirit was then to be
forced into science: that is an unnaturalness against which protest
must be made in the most forceful way. Not in itself, but in
its disproportion to the Spirit of life, is the word of Scripture the
dead word; it is Grundtvig's immortal merit that he ``in the power of
the Spirit, with a living voice'' has testified against ``the dead
word''.

The word of awakening is a living word; but awakening in turn assumes
many forms; we may here distinguish between the immediate or direct,
the reflected or indirect, and the reformatory in principle.
% ---- printed p.504 (PDF 517) ----

\opage{504} \greekrun{α}{Immediate Awakening.}

The more strongly the opposition between the new light in religion and
the spiritual darkness in the old conditions of the world can show
itself to those concerned, the more favorable are the conditions for
an immediately powerful awakening. Peter's speech at the feast of
Pentecost (Acts 2) is typical of all speeches of awakening. The
phenomena of awakening in the apostolic age bear, in their primally
powerful immediacy, every mark of the Spirit's first personal working.

It is otherwise with the awakenings that take place in ages in which
the life of the spirit has a far lesser tension. Much may indeed be
accomplished here in all simplicity; the preachers of repentance in the
Middle Ages, the Pietists and the Methodists in the eighteenth century
have powerful examples of religious awakening in a popular direction
to show. But unclarity about the spiritual as well as the natural
conditions of life, in combination with much that is morbid, ecstatic,
and fanatical, shows plainly enough that it is not a deeper
self-understanding in the Spirit but an obscurely fermenting urge of
conscience that has played the main role in these awakenings. There is
a half-reflection that makes simplicity spurious and paralyzes the
powers of the spirit. This half-reflection comes from an overpowering
spirit of the age. The light of culture cannot be kept out; it steals
into ``the simple'' and gives half-reflected piety an uncanny tinge of
narrowness, pious borné-ness, indeed even of selfish conceit.

How little a powerful awakening, however immediately it may seek to
cling to the first ideal presuppositions of the life of faith, is able
to defy the history of culture, the Irvingites with their apostolic
congregation furnish telling proof in the present day. The living word
demands that the personal organs be in living understanding with the
holy Spirit of personality. But when
% ---- printed p.505 (PDF 518) ----
\opage{505} ``apostolic'' leaders, with forced simplicity, shut their eyes to the
conditions of the life of culture, although the sun of reflection
stands high in the heavens, they must wait in vain for the Spirit to
work with original power through thought-blind, miracle-craving
organs.


It is the resistance of reflection that the spirit in the present age
must unconditionally overcome; so far as the principle is concerned,
this can happen only if the essence of reflection becomes completely
transparent, and thus only if the believing consciousness passes
through the fiery trial of reflection. Herewith we come to

\greekrun{β}{The Reflected or Indirect Awakening.}

It is this standpoint that is represented in the very most proper sense
by S.\ Kierkegaard. To translate religion into reflection in order to
take it out of reflection again and preserve it in pure simplicity ---
that was the task this remarkable man set himself, and for whose
solution he had received the call of the Spirit in the richest measure.
S.\ Kierkegaard holds unshakably fast to simple faith; but in making
Christianity a communication of existence, he draws a finely woven net
of existence-tasks together around ``the individual,'' and so reflection
begins with the indirect communication. The relation to God is
``simple'' and yet ``infinitely reflected.'' Everything is staked on
``serving the truth,'' but --- by ``a godly deception.'' ``The ideals
are proclaimed,'' but the proclaimer himself ``is without authority.''
The whole undertaking begins conservatively and ends revolutionarily.
Seen superficially, it looks like nothing but contradictions; but if one
enters into the whole design and follows the development in its details
and ramifications, then the execution is consistent from beginning to
end.

It begins with the presupposition that Christianity is the eternal
truth, the highest ideal power, and it ends with the declaration that
Christianity
% ---- printed p.506 (PDF 519) ----
\opage{506} hardly exists in the world, that ``the Christianity of Christendom'' is
``an official deception''; that is truly for awakening!

But the awakening is indirect, reflected, negative. It is the
earnestness of the awakening that men must begin by becoming afraid of
themselves and asking in anguish of conscience: what shall we do? Thus,
after all, the awakening also began at the first Pentecost. But here
now is the great difference between an apostolic and a reflected
awakening. The apostolic awakening, which is personally borne by the
Spirit as Holy Spirit, has immediate power over all the finite
considerations of reflections; the modern awakening, reflected in
negative form, takes the detour of mediacy in order to sharpen thought
and, with thought, the sting that is to bore its way into the
conscience. The Apostle's answer in the Spirit is positive: ``Repent,
and let every one of you be baptized in the name of Jesus Christ for
the remission of sins; and ye shall receive the gift of the Holy
Ghost.'' But the thought-awakener's answer in the Spirit is negative:
``Despair! For the eternal truth has been turned into deception.
Despair! For the Christianity of the New Testament no longer exists on
earth.'' It is wounding, but it is --- indirectly
awakening\footnote{Grundtvig in his ``Børnelærdom'' has called
S.\ Kierkegaard ``a hair-splitter''; that is a misunderstanding. From
the standpoint on which he placed himself, K.\ is entirely right. His
task was to illuminate the disproportion between the Christianity of
Christendom and the Christianity of the New Testament, and he has done
that thoroughly.}.

\greekrun{γ}{The Reformatory Awakening.}

The Spirit of the Father and the Son is the Spirit of the congregation;
this implies that the individual's life in the Spirit must in every age
be conditioned by a corresponding life of the Spirit in the
congregation. Even the preacher of repentance who declares the
% ---- printed p.507 (PDF 520) ----
\opage{507} congregation dead nevertheless stands, by this very declaration, in a
positively determined relation to the congregation, and must
necessarily presuppose that the dead one is not dead but sleeps; if he
does not want to contradict himself, he must presuppose that among the
many sleepers there is at least one who is awake. The task, then,
becomes to awaken the congregation that has fallen asleep and lies in
the sleep of death. This awakening can happen only by putting into
force the conditions of life indispensable to the congregation. It is a
reformatory task.

In the sixteenth century it was the papacy with the Catholic hierarchy
that was on the point of stifling the congregation; then there was a
break with the traditions of the hierarchy, and the word of Scripture
was set, as a replacement, beside the sacramental institutions. Besides
this, the Reformation not only indicated the conditions of the spirit
indispensable to a freer life of the congregation, but it also brought
about a deeper understanding of the relation between religion and
culture.

In the present age, when the life of the congregation is indeed not
oppressed by external despotic power, but is partly encumbered with a
host of inherited forms of compulsion, and partly exposed either to
being screwed down under a doctrinal system of scholastically objective
dogmas, and so getting actual human life against it, or --- insofar as
the spiritual leaders wish to keep up with the times to some degree ---
to collapsing, as the foundation of Scripture bursts and the
characteristic features of revelation are wiped away by the damp sponges
of the mediating theologians --- in the present age a fundamental
awakening from out of the totality itself is required once again. The
task is twofold: 1) to reconcile religion with popular enlightenment by
going back to the original unity-in-difference of universal human life
and the life of faith; 2) to secure the congregation against all the
disturbances that the surreptitious guardianship of the School over the
concerns of life, and in particular over those of religion, has called
forth in the course of time.
% ---- printed p.508 (PDF 521) ----
\opage{508} In both respects the Grundtvigian view is the only one among us that
has a future.

\runhead{b) Regeneration: The Sacrament of Baptism.}
\parmark{36}

What has already been said of awakening holds of regeneration: in the
religious-ethical sense it does indeed essentially concern the
individual, but from the standpoint of the philosophy of religion it
must be conceived as a work of the Spirit in the congregation. That the
life of the spirit in the individual has its origin in the life of the
Spirit that streams through the congregation is expressly seen from the
sacrament of baptism. The individual is indeed baptized personally into
Christ, is received in baptism by Christ himself, present in the Spirit
with the Word --- it then follows directly from the principle of faith
that the sacramental fundamental word of baptism is not merely Christ's
own word of institution, but the Word of life from the living Christ:
only as the Word of the new life can the word of baptism be a word of
regeneration, with divine power in the Spirit --- but Christ does not
baptize outside of, but in his congregation, in and into his
congregation. The conditions for the growth of the life of the spirit in
the individual are therefore inseparable from the congregation's own
conditions of life.

The original conditions of life of the congregation all have their
source in the fundamental word of baptism. This fundamental word, in
which the Spirit works through the cleansing element, does not have its
significance in the fact that the holy act is performed at Christ's
command, as though it were a work of the
law\footnote{In this sense the law ought to be excluded from the
gospel.}, but personally in the fact that the act itself is an
evangelical act of Christ and as such, in the Spirit, an actual
sacramental act.
% ---- printed p.509 (PDF 522) ----

\opage{509} What takes place through this sacramental act in the soul of the
individual being baptized can be summed up in the one expression: the
first fundamental effects of regeneration. But what these fundamental
effects are, how they take place, whether they take place in the child
in the same way as in the adult, etc., of that there is no objective
knowledge. Every attempt to ground the ``effects of grace'' out of a
scriptural-scientific system of faith ends in the mystical-magical, an
eloquent proof that the grounding is inverted and that the doctrine of
faith, according to its principle, is anti-rational. In this connection
we shall consider the sacrament of regeneration as paradox, as Mystery,
and as the beginning of the new life in the congregation.

\greekrun{α}{The Sacrament of Regeneration a Paradox.}

When ``the individual's eternal blessedness is decided in time through
the relation to something historical, which is furthermore historical in
such a way that into its composition there has been taken up that which
according to its essence cannot become historical, and so must become
so by virtue of the absurd,'' then naturally both the act of baptism and
regeneration have taken up into their composition elements which
according to their essence cannot be composed, and which therefore must
be composed ``by virtue of the absurd.'' S.\ Kierkegaard's defense of
infant baptism is highly characteristic in this respect. For he regards
infant baptism ``not only as defensible as orthodox and praiseworthy as
the expression of the piety of parents who cannot bear to be separated
from their children in regard to what is for them the matter of their
blessedness, but also, in yet one more sense of which people are
perhaps not aware, as a good --- because it makes it still more
difficult to become a Christian''\footnote{Joh.\ Climac.\ S.\ 291.}.
Christ's words (Matt.\ 18; Mark 10): ``for of such''
% ---- printed p.510 (PDF 523) ----
\opage{510} (namely little children) ``is the kingdom of God'' he explains
thus,\footnote{Anfr.\ Skr.\ S.\ 457.} that it is the disciples whom
Christ wants ``to distance a little from himself by means of a
paradox'' \dots\ ``The apostles rebuke the little children; but Christ
does not rebuke in return, he does not even reprove the apostles; he
turns to the little children, but he speaks to the apostles. The
paradox lies in making a child the paradigm; \emph{partly} because a
child, humanly speaking, cannot be that at all, since it is immediate
and explains nothing (therefore a genius cannot be a paradigm either
--- the melancholy in the genius's distinction), not even for other
children, for every child is merely itself immediately; \emph{partly}
because it is made the paradigm for an adult, who in the humility of
the consciousness of guilt is to resemble the humility of innocence.''

The light and shadow sides of the dialectic of paradox become quite
visible in this connection. As religious-ethical dialectic of existence
it has an extraordinary strength, but in the sense of the philosophy of
religion its intellectual arbitrariness is intolerable. If Christ's
divine-human essence is an absurdum, the sacrament of baptism an
absurdum, regeneration an absurdum, faith an absurdum: then,
undeniably, the Spirit of Christianity is in itself an absurdum, and all
self-understanding in the Spirit only an absurdum. But
self-understanding in an absurdum sublates itself.

\greekrun{β}{The Sacrament of Regeneration a Mystery.}

If \emph{absurdum} is to mean not that which is in itself, by
its nature, contradictory, but only to be the expression for what is
incomprehensible in the sense of knowledge, then the contradiction
dissolves in the self-dissolution of the concept, and the Mystery steps
into the place of the paradox. That the Mystery in the sacramental act
must be just as impenetrable as the Mystery in the paternal
% ---- printed p.511 (PDF 524) ----
\opage{511} Self is grounded in the nature of the matter. As regards the
institution of the new relation to God, it is just as impossible to
articulate the reciprocal determinations of the divine and the human,
of the supernatural and the natural, in the fundamental workings of the
Spirit, in the hidden miracle, as it is impossible to articulate the
reciprocal determinations of nature and spirit in the Father's own
essence. The analyses corresponding to the psychology of regeneration
do not at any point hit upon the first, original thing, but only what
is derived from it. The first, original thing is that regeneration by
means of the sacramental act is a work of the Spirit by means of the
Son, a work of the Son by means of the Spirit, and in both respects a
personal, divine-human work. The original thing is that the
divine-human work furnishes the foundation for the self-conscious
development of the life of faith in the individual. As self-consciousness
is awakened and the human self enters into reciprocal action with the
divine, this reciprocal action must invariably, uninterruptedly show
itself grounded on the new life of Christ implanted by the sacramental
act. Whether the individual is otherwise baptized as an adult or as a
child, whether self-development can begin at once or only after an
interval of time, the relation of the life of faith to the sacramental
act will in principle nevertheless be the same. To believe in Christ is
to appropriate Christ personally in the Spirit; but the act by which
the believer in inwardness appropriates Christ personally is, after
all, conditioned precisely by the act by which Christ himself has
appropriated him personally, and thus conditioned by Christ's personal
act of baptism: the appropriation of Christ is a living, uninterruptedly
continuing remembrance of the covenant of baptism.

\greekrun{γ}{The Sacrament of Regeneration an Act in the Congregation.}

Insofar as it is the minister of the Word who performs baptism as an
act of faith, the act is human; but
% ---- printed p.512 (PDF 525) ----
\opage{512} insofar as the institution is from Christ, and Christ himself, present
in his Word, accomplishes the act, it is divine; to unite the divine
with the human is the personal work of the Spirit. Whoever mistakes the
Spirit must either regard the act from the merely human side and reduce
the sacrament to a religious ceremony, or else conceive the human as
an impersonal means and point of passage for the divine, whereby the
magical, mystical, and sorcerous arises. If the one baptizing were not
a member of the congregation, the human could not come into its own,
and the act could not be an act of faith; but if the effect of baptism
were made dependent on the subjective power of faith of the one
baptizing, the divine could not come into its own: baptism would then
not be an act of the living Christ himself, not a sacramental act. The
two main sides of the sacramental act are inseparably united in the
Spirit of the congregation.

The same thing repeats itself when we consider the act of baptism in
relation to the one who is baptized. As an act of Christ the sacrament
has a validity independent of all human subjectivity: it is not the
human individual who chooses Christ, but conversely Christ who chooses
and receives the individual. This is the divine element in the work of
the Spirit. But the work of the Spirit is a divine-human work, and is
conceived spiritlessly when the human side is overlooked. Now since the
act, seen from the human side, is an act of faith, the question arises
whether there is not, after all, an essential difference between infant
baptism and the baptism of adults. We answer: There is a difference, an
unmistakable difference; but the difference is only psychological, not
one of principle. The psychological difference is plain enough. In the
adult the preparatory activity of the Spirit precedes the sacrament:
the adult must seek the sacrament in faith
% ---- printed p.513 (PDF 526) ----
\opage{513} and to that extent bring faith with him. In the child who is baptized
no preparatory activity of the Spirit precedes; the child in actuality
brings no faith with it. Now if one forgets that the union of the human
with the divine is the work of the Spirit, it seems as though there
were here only a choice between one of two things: either to reject
infant baptism or to concede that the sacrament of regeneration can
work both with faith and without faith; then comes critique, and then
one stands there in the midst of a tangle of insoluble difficulties.
The difficulties vanish when we fix our gaze on the relation of the
Spirit to the sacrament. Without the working of the Holy Spirit faith
is impossible: the Spirit does not proceed from faith, but conversely
faith from the Spirit. The working of the Spirit, whether in the adult
it goes before in preparation or in the child's soul follows after in
completion, is equally bound to and conditioned by the sacramental act
itself. If we now remember that the Spirit is the Spirit of the
congregation, that the act takes place in the congregation, that the
one who is baptized is baptized to be a member of the congregation,
then it is seen that the difference between the child and the adult is
only psychological, not one of principle.

The baptized child, considered in the divine-human light of the Spirit,
is, as ``a child of God,'' at the same time a child of the congregation.
Everything depends on the divine and the human not falling apart from
one another spiritlessly. Certainly the Spirit can begin to work as the
Spirit of faith only gradually, as the child's consciousness is
awakened; but already the first, the very first effects will point back
through the awakening consciousness to the sacramental act. The
question: where is the child's faith when it is baptized? is easy to
answer from the standpoint of religion. Faith is in the child as a
possibility, but in the congregation as an actuality. ``Then is it
perhaps the congregation that believes in the child's place?'' By no
means! Where there is talk of a substitute, there is talk of an
% ---- printed p.514 (PDF 527) ----
\opage{514} individual. But the congregation is more than an individual. In this
light it is not the parents, not the godparents, not the members of the
congregation as a sum, who believe in the child's place; individuals,
however many they may be, cannot possibly undertake in this way to
believe in another's place. To look to the individuals here is to be
misled by finitude, is to separate the human from the divine and to
grieve the Spirit. No, it is the congregation as congregation, as the
temple of the Spirit, as ``the bride of Christ,'' that bears in its
heart the conditions of real possibility for the development of the
child's faith. Through the act of baptism, therefore, the child is not
merely received by Christ: but Christ ``takes it in his arms and
blesses it,'' and then lays it in the lap of the congregation, the
spiritual mother's lap from which regeneration proceeds. ``But,'' it is
asked, ``what if no one takes charge of the child's upbringing? Or if
the child, on reaching the age of maturity, despises the sacrament of
baptism and reproaches the congregation for having anticipated its free
choice: what then?'' The answering of these and similar questions
belongs under the following point of view.

\runhead{c) Sanctification: The Sacrament of the Lord's Supper.}
\parmark{37}

As a Christian concept of faith, the concept of the congregation is a
divine-human concept. Seen from the divine side, the congregation is not
a society of like-minded individuals, but a community of the Spirit;
on its human side, on the other hand, this community of the Spirit is
an association of like-minded individuals. In the development of the
life of the congregation the emphasis must, to be sure, fall with
alternating preponderance now on the unity of the community, now on the
diversity of the individuals, but always in such a way that the
relation between the original and the derived is kept clearly
% ---- printed p.515 (PDF 528) ----
\opage{515} in view. If we look at the original, it is not the individuals who, by
appropriating the Spirit one by one, have joined together and formed
the congregation, but conversely it is the Spirit who, by ``calling and
gathering'' the individuals, has founded the community. In relation to
the divine work of the Spirit, the individuals in their singleness and
individual diversity are to be regarded as personal possibilities; in
relation to the human work of the Spirit, the personal possibilities are
destined for free, independent development, a self-development in which
each one for himself must live and work as an actual person, as the
individual. The former must correspond to the sacrament of
regeneration, the latter to the sacrament of sanctification. It goes
without saying that, just as in regeneration the divine takes up and
permeates the human, so the human, the self-development corresponding to
sanctification, must proceed from and be grounded in the divine. From
this standpoint we consider the sacrament of sanctification as paradox,
as Mystery, and as an act in the congregation.

\greekrun{α}{The Sacrament of Sanctification a Paradox.}

It is highly characteristic that the representative of the dialectic of
paradox --- we mean, of course, S.\ Kierkegaard --- has everywhere put
the sacrament of baptism in the shade, whereas with evident predilection
he has brought forward the sacrament of the Lord's Supper. This is most
intimately connected with his chosen standpoint: he takes ``official
Christianity'' to task; so he fixes his fiery gaze on the individual,
but overlooks the community; so he enjoins sanctification, but overlooks
regeneration. Now if it really were a one-sidedness of principle, a kind
of half-conscious or unconscious heresy, of which he made himself
guilty, then the spokesmen of official objectivity could easily dismiss
him as an isolated subjectivist who
% ---- printed p.516 (PDF 529) ----
\opage{516} wanted to work ``in the direction of inwardness'' by making a noise in
the direction of outwardness; but that is not how the matter stands.
The sacrament of baptism, regeneration, the life of the congregation ---
% print: "Gjenfødelen" for Gjenfødelsen
these religious presuppositions are by no means denied, but they are
indeed overlooked; that is to say, in accordance with the chosen task
--- to awaken consciousness of the ethical-ideal demands of the life of
the spirit --- they are pushed back, but not otherwise than that they
can be brought forward again at whatever moment it becomes necessary.
That is why ``that individual'' stands --- not, as people have so
blithely asserted in order to relieve the oppression of their hearts,
outside, but inside the congregation, in the midst of the congregation;
and the congregation that does not want to fall asleep on an imagined
regeneration without sanctification does well to hear his voice.

As little as sanctification is possible without regeneration, just as
little is actual regeneration possible without sanctification;
regeneration is the real possibility of sanctification, and consequently
sanctification is to be determined as the very actuality of
regeneration: sanctification is the religious-ethical development of
character, a development of character that alone can preserve ``the
joyful Christianity'' in ``fear and trembling.'' Sanctification proceeds
from regeneration as love from faith, which is also why the sacrament
of sanctification is, as we know, especially designated the sacrament
of love.

To the sacrament of love correspond ``works of love,'' for love in
believers, the love that hides ``the multitude of sins,'' has --- as a
communication of existence this is a paradox of faith --- its holy
ground, its inexhaustible spring, in the incomprehensible love of the
Atoner, the love with which he continues to give what in the
fullness of time, once for all, he has given up --- his body and blood
for the remission of sins.
% ---- printed p.517 (PDF 530) ----


\opage{517} \greekrun{β}{The Sacrament of Sanctification a Mystery.}

If the sacrament of the Supper is to be objectified in knowledge, then,
whether we take our stand on the evangelical Lutheran or on the papal
Catholic standpoint, it becomes an insoluble paradox, an unconditionally
repellent absurdity. The Lutheran representation that Christ's body
% print: „lutherke“ for lutherske
should actually be imparted \emph{in, with} and \emph{under} the bread,
Christ's blood \emph{in, with} and \emph{under} the wine, is in the sense
of science just as mystically unreasonable as the Catholic assertion that
it is the priest who, by virtue of the miraculous gift imparted to him at
ordination, \emph{transforms} bread and wine into Christ's body and
blood. The church parties (Zwinglians, Socinians, etc.) that have tried
to give the dogma of the Supper a more rational interpretation have with
their rationality contributed, not to explaining, but to explaining away
the sacrament.

Those who otherwise find it hard to see that the principle of faith is
absolutely heterogeneous with that of knowledge, and objectification in
knowledge naturally heterogeneous with objectification in faith, must
surely have their eyes opened when the talk is of the sacrament of the
Supper. What appears here when we objectify in knowledge? A box of wafers
by itself, a cup of old wine by itself, a priest who reads the words of
institution --- these are nothing but external things, nothing but
sensuously actual objects. Over all this there then hovers an image that
fantasy forms for itself of the glorified Christ. These many
heterogeneous elements of representation thought is then to work together
into a dogmatically objective unity of concept; by virtue of this unity
of concept \textit{in, cum et sub} is to be scientifically explained: the
explanation will be accordingly.

It is otherwise when objectification takes place in faith. The gaze does
not then fasten on external objects in sensuous existence by themselves;
but through the outer the inner is looked upon: the act is looked upon as
a holy act, as the sacramental act of Christ, present with the Spirit.
All
% ---- printed p.518 (PDF 531) ----
\opage{518} thoughts are now transformed into holy thoughts of impression: to
awaken these thoughts is the work of the Spirit. Thus Christ is, by means
of the Spirit, actually present, spiritually-bodily, in his sacrament.
Objective knowledge vanishes: the presence is objectified in faith; but it
is not a figment of faith, for Christ comes to the meeting at his given
word. Yes, --- says the Gospel --- he comes with the Spirit in the word;
he grants a personal meeting! Woe to him who dares to meet him in the
sacrament without meeting in faith: such a one ``eats and drinks judgment
to himself.''

The sacrament of sanctification is thus not an absurdity, still less an
object of reason; but it is a Mystery.

\greekrun{γ}{The Sacrament of Sanctification an Act of Christ in the Congregation.}

In the sacrament of regeneration man cannot be said, in the personal
sense, to meet with Christ; for in this sacrament human personality is as
yet only a possibility --- nor, indeed, can a newborn child be said to
meet personally with its own father. It is the misunderstanding of the
Baptists that in the sacrament of baptism they want to appear already as
persons, instead of appearing as children. In the sacrament of the
Supper, on the other hand, the individual cannot appear as a child, but
must appear as an actual person. But when a human being is to meet
personally with the Holy One, as self with Self, he receives a living
impression of what it is to be a sinner. In the sacrament of regeneration
sin is conceived impersonally, as a sickness in the nature of man and
thereby in the nature of the will; but in the sacrament of sanctification
sin is conceived personally in the strictest sense. According to the
ideal of holiness, man should by now have developed into religious
character; but his personal development testifies against him, for alas,
he is personally --- a sinner.

The forgiveness of sins at the Supper therefore bears a different stamp
from the forgiveness of sins at baptism. \emph{Forgive us our debts, as
we forgive our debtors.} --- ``There is thus
% ---- printed p.519 (PDF 532) ----
\opage{519} actually a double forgiveness of sins, namely the \emph{divine},
which is bestowed in \emph{baptism}, and the \emph{human}, which belongs
to the \emph{Supper}; and although the divine, which is bestowed in
\emph{baptism}, is conditioned only by \emph{faith} and \emph{confession}
together, and is altogether \emph{complete}, it is nevertheless given
only on the \emph{presupposition} that the \emph{mutual human} forgiveness
of sins is added in the \emph{congregation}, consequently with the
\emph{reservation} that it \emph{falls away} when the person concerned
withdraws from it, as the Lord has both shown us in the parable of the ten
thousand talents and has expressly testified to his disciples that if
they do not forgive one another their debts, then neither will the
heavenly Father forgive them theirs.''\footnote{Grundtvig. Børnelærdom. S.\ 225--26.}

Sanctification is the work of the Spirit; the sacrament of sanctification
is the act of Christ. The sacramental acts of Christ have unconditional
validity for the life of faith only by being objectified in faith; but
they are objectified in faith only when they are presented in the light
of the Spirit, in inseparable unity with the Spirit himself and the works
of the Spirit.

\lettersub{C.}{The Kingdom of the Spirit.}
\parmark{38}

The Kingdom of the Spirit is in the deepest sense the kingdom of
personality. Individuals can indeed assert personal rights in other
communities too, and to that extent appear as persons; but only in the
Kingdom of the Spirit is the principle of community itself personal. The
principle of personality
% ---- printed p.520 (PDF 533) ----
\opage{520} with its divine-human fundamental effects is the only thing
able to nourish, form and develop the personal life of individuals for
time and eternity. Since the Spirit begins the actualization of his
kingdom in this world, he must no doubt engage with the world's terms and
require natural conditions of life; but since the kingdom is not of this
world, he must work unceasingly to transform the world and to change the
conditions within the temporal into means for the eternal. The
interaction between the Spirit in the congregation and the universal human
spirit takes on, in every gifted people, a stamp of its own, which is
determined by the spirit of the people.

\runhead{a) The Spirit in the Congregation.}
\parmark{39}

As the personal spirit of community, the Spirit himself is the organizing
principle of the congregation. The organizing activity is at once both
divine and human. If the relation between the divine and the human is
distorted, either in such a way that the human is taken to be immediately
absorbed into the divine, or in such a way that the divine dissolves and
perishes in the human, the organism is attacked and ``the body of Christ''
sickens.

By laying weight one-sidedly on the immediately divine influence of the
Spirit, one must necessarily be driven to one of two extremes: \emph{either},
with the Catholics, to believe in the infallibility of a church visible in
traditions and institutions, and then we have the hierarchy; \emph{or},
with the abstract-critical intellectualism of the Reformed, to break the
staff over all historically developed church forms and make the condition
of the congregation in the apostolic age an external, unalterable norm.

By one-sidedly emphasizing the human in the organizing activity of the
% ---- printed p.521 (PDF 534) ----
\opage{521} Spirit one is driven to two other extremes: the
\emph{Lutheran-official}, which, while holding fast to the apostolic divine
% print: „lutherk-officielle“ for luthersk-officielle
in word and sacraments, ascribes to the traditional development of rite
and church constitution so organic a historical validity that the human
forms of law, although the human element in their origin is conceded in
principle, are nevertheless regarded as inviolable on account of the
prescriptive right once acquired; and the \emph{humanistic-dissolving},
which, appealing to ``the perfectibility of Christianity,'' makes the
spirit of the congregation a particular form of the world-spirit, and
consequently cannot ascribe to congregational life any special,
independently organizing principle.

In order to see how these extremes may be mastered by the Spirit himself,
we cast a glance at the inner and outer organization of congregational
life and at the relation between the two.

\greekrun{α}{The Inner Organization.}

``As instruments of the church-founding Spirit the apostles stand not
only in the freest relation of unity to the Spirit, but also in the
freest relation of unity to \emph{one another}. Although there are
diverse gifts, there is yet but one Spirit; and although there are
diverse apostolic forms of doctrine, the one fundamental truth yet lives
in them all. Precisely by distinguishing its content into a diversity of
moments does the Spirit reveal its fullness and its riches. For the same
reason it is not the consciousness of the individual apostle, but the
\emph{apostolic total consciousness} that is the complete expression of
the revelation of the church-founding Spirit, just as this is also the
complete expression of the Church's consciousness of its fundamental
relation to the Lord and the Spirit, its fundamental relation to the
world and to itself. As representatives of the mother church the
apostles express not merely the
% ---- printed p.522 (PDF 535) ----
\opage{522} ecclesiastical self-consciousness of a single age, but are the
Christian Church's representation for all ages.''\footnote{H.\ Martensen. Den christl.\ Dgmt.\ S.\ 353.}
% print: the two notes on p.522 may be transposed (by sense this one would be „Anfr. Skr. S. 424.“); kept as printed

Here we have the right, unalterable canon. We shall now see how it may be
applied. The opposition between teachers and hearers, for example, is an
opposition determined by the organizing activity of the Spirit, an
opposition, therefore, that must have its divine as well as its human
side. The divine side has been brought out with strong emphasis by the
author of the ``Christian Dogmatics,'' where he says: ``Christian
preaching is \dots\ not merely a human speech about Christ, but it is
Christ himself who through it gives himself presence to the world and to
believers \dots\ This is the secret of Christian preaching, whereby it
distinguishes itself from all other speech that can be delivered about a
historical person, that in the same measure as it is actually a preaching
in Christ's name, in the same measure the word still exercises
\emph{essentially} the same effects for awakening and confirming faith in
the person of Christ as when he walked bodily here below, exercises
essentially the same impression for awakening and confirming faith in the
salvation that is bestowed on us in him; all of which is explicable only
from this, that the Church has not an absent but a present Christ, who
breathes livingly through his word.''\footnote{Anfr.\ Skr.\ S.\ 424.}
These are strict demands that are hereby placed on the proclaimers of the
word in the congregation; but insofar as one looks exclusively to the
divine ideal, they are by no means too strict.

In order to satisfy in some degree the divine demands of the Spirit, no
exertion whatever, then, ought to be spared on the human side. The
congregation must accordingly do everything so that those who feel a
call to
% ---- printed p.523 (PDF 536) ----
\opage{523} proclaim the word may obtain a preparation that can make them
fit, in theoretical as well as practical respects, to carry out so
responsible a work of the Spirit. Those who do not in some measure
satisfy the acknowledged demands the congregation must, then, be
permitted to turn away; those, on the other hand, who do satisfy these
demands it must uphold and acknowledge as rightly appointed, without
regard to by what road and in what manner the preparation has humanly
been attained. So it was, indeed, precisely in the apostolic age; and ---
``the apostles express not merely the ecclesiastical self-consciousness
of a single age, but are the Christian Church's representation for all
ages.'' From this it follows, then, that when a preacher is unable to give
the congregation a palpable impression of ``the secret of Christian
preaching,'' is not capable of preaching in such a way that the word
``still exercises \emph{essentially} the same effects for awakening and
confirming faith in Christ's \emph{person} as when he'' --- Christ himself
--- ``walked bodily here below,'' then the spiritless preacher must be
removed from his office as unfit for the ministry of the word, and the
congregation is in the name of the Spirit not merely entitled, but even,
according to the canon we have set up, directly obliged to put a capable
man in the place of the incapable one.

If we now consider the imperfection of human things, it cannot be denied
that the congregation, if it is to satisfy the claims of the Spirit, has
a difficult task to solve.

\greekrun{β}{The Outer Organization.}

In order to reconcile the ideal with actuality, to restrict arbitrariness,
and to bring about a social order adapted to the demands of the Spirit,
weight is then laid on rules, laws and standing ordinances. In this way
the transition took place from the free, circumstantial forms of the
apostolic congregations to the objective, strictly rule-bound constitution
% ---- printed p.524 (PDF 537) ----
\opage{524} that particularly distinguishes the Catholic church system. The
partition that was now set up in the name of the Spirit between clergy and
laity was, humanly speaking, of great advantage to the clergy, but,
spiritually seen, of irreparable harm to the congregation. Now,
especially after \emph{ordination} had been raised to a sacrament, it was
of course not at all difficult to satisfy the congregation. The cleric's
word, whatever the word might otherwise be like, was bound absolutely to
give the congregation an impression of ``the secret of Christian
preaching,'' for the secret was that the cleric always preached in the
power of the Spirit, because he preached in the power of his ``holy
office.'' Thus the objective, the official, stepped into the foreground.
The Catholic Church lays claim to being a direct, infallible continuation
of the apostolic, and yet the Catholic Church forms a decisive
counter-image to the apostolic: in the Catholic everything is official, in
the apostolic nothing.

The official is the externally objective, the impersonal; from this it
follows that the holy Spirit of personality can never meet more dangerous
resistance than that which lies in the
official.\footnote{``Imagine the following case. There comes to you
% print: „Oficielle“ with one c
a human being who (without in any way giving you the notion that he is
insane) quietly, earnestly, yet deeply shaken, says to you: `pray for me,
O pray for me': is it not true, it would make an almost terrifying
impression on you? And why? because you yourself personally received the
impression of a human personality who presumably had to be fighting the
uttermost fight with a personal God, since it could occur to him to say
% print: „sigc“ for sige
to another human being: pray for me, pray for me.''

% print: this opening quote is never closed (closed here after „Spilledaase . . .“)
``When, on the other hand, you read, for example, in a `\emph{pastoral
letter}': Brethren, include us in your intercessions, just as we
unceasingly night and day pray for you and embrace you in our
intercessions --- how is it that this presumably makes no impression on
you at all? is it not because you involuntarily conceive a suspicion that
it may be formula, rigmarole, official, out of a handbook or out of a
music box \dots'' ``For believe me, there is nothing that God finds so
repugnant, no heresy, no sin, nothing so repugnant to him as the
official.'' S.\ Kierkegaard. Øieblikket Nr.\ 4, 7.} It is not a further
development of
% ---- printed p.525 (PDF 538) ----
\opage{525} the first constitution that directly brings the official with
it. On the contrary! A richer articulation, a more definite separation
between the various positions and functions, could very well take place
without the original fundamental character being distorted thereby. It
is not the distinction between lower and higher dignities, between
deacons, presbyters and bishops, that has made the Catholic Church so
absolutely unapostolic. No, it is the distortion of the organizing
principle.

In the apostolic age the all-organizing principle is personal. Positions
and activities are distributed according to the gifts; but the gifts are
gifts of the Spirit, free gifts, ``gifts of grace.'' No one has the gift
because he is chosen, but he is chosen because he has the gift. The
rule-bound order, the human side of the matter, is at all points
dependent on the divine freedom: all forms are living.

It is otherwise with the official. The official makes the forms the main
thing and adds the name of the Spirit. Now all dignity, all authority
becomes something externally set up, something objectively represented,
without personal truth. Now the deacon does not become priest, nor the
priest bishop, because the Spirit gives him a higher gift; but he is
assumed, he is represented, to receive the higher gift when he officially
advances to the higher position. Now the spiritual leader does not have
authority because he in truth possesses religious personal superiority;
but he is assumed, he is represented, to be a holy man because he has
entered ``the holy office.''

To abolish the higher clerical dignities and by leveling legal enactments
to force through a republican or radical-democratic constitution of the
congregation is by no
% ---- printed p.526 (PDF 539) ----
\opage{526} means directly the way to master the official. One could by law
screw the congregation back to a constitution that in all external
particulars resembled the apostolic one: such a legally compelled copy of
the free, personal exemplar would nevertheless be only a caricature of
freedom, for the copy is official.

% print: „Vcxelvirkningen“
\greekrun{γ}{The Interaction between the Inner and the Outer Organization.}

The conflict between the personal and the official stems from this, that
the divine and the human in the organizing activity of the Spirit are
separated and combined in an external, finite manner: the unity of
opposites becomes an immediate unity, now in the one, now in the other
extreme, and in both cases spiritless. The official does not lie in
legally determined forms or in lawfully established offices. If a
polemic against the official goes to the extreme of rejecting the thought
of a well-organized constitution of the congregation, then the
personality of the Spirit is violated; for the principle of personality
is and must be organizing.

A clash between the two extremes, the divine, which aims at the free
inner, and the human, which on the conditions of actuality demands its
determinate outer, was already present from the beginning and will
continue through all ages. This clash derives not only from the conflict
between ideal and actuality, but is grounded in particular in the
double-sided nature of personal life itself. The members of the
congregation are, as regards the ground of their life, all indeed
``members of one body,'' all together ``branches on one vine,'' and must
to that extent all be presupposed to have one fundamental will, one goal
and one striving; but it is at the same time the demand of personality
that the individuals, each with his own conditions, according to his
gifts, must develop freely and independently: here there are
% ---- printed p.527 (PDF 540) ----
\opage{527} thus many wills, many points of view, many endeavors.

The independence of individuals, which from the divine side is grounded
in their personality, is from the human side inseparable from their
accidental individuality. The diversity of individuals must therefore be
acknowledged for the sake of personality. Thus it is, then, impossible
even for God's own Spirit to organize a community of personally
independent individuals without a system of laws: only by means of laws
can the Spirit master the arbitrariness of individual wills, protect the
freedom of individuals, and yet preserve the organic unity of the
community. From this point of view we now return to the organization of
the apostolic congregations, according to the canon set up: ``As
representatives of the mother church the apostles express not merely the
ecclesiastical self-consciousness of a single age, but are the Christian
Church's representation for all ages.'' It is 1) the legislative
authority, 2) the peculiar character of the law, and 3) the observance
of the law, on which attention fastens.

1) In the divine sense it is the Spirit himself that has the legislative
authority; in the human sense it is the congregation. How both main sides
come into their right can be seen from the Acts of the Apostles, ch.\ 15.
That the Catholics have departed from the exemplar must strike the eye.
From the apostolic expression (v.\ 28): ``\textit{Ἔδοξε γὰρ τῷ ἁγίῳ
πνεύματι καὶ ἡμῖν}'' the legislative councils have formed for themselves
% print: the quotation „Ἔδοξε … ἡμῖν is not closed
a standing formula; the formula has stepped into the place of the
principle, for the principle is denied. According to the apostolic
principle --- ``Then the apostles and the elders with the whole
congregation resolved, etc.'' (v.\ 22) --- it is the whole congregation
% print: „Da besluttede … o. s. v. is not closed; „(v, 22)“
that takes part in the legislation; but according to the Catholic
principle it is exclusively the clergy that has appropriated the
% print: „lovgivendo“ for lovgivende
legislative authority. That the assembled clergy are the organs of the
Holy Spirit is not something
% ---- printed p.528 (PDF 541) ----
\opage{528} actual; it is only something they represent: here we have the
official, and to the official corresponds the formula.

If we give up the formula and hold fast to the principle, it then becomes
clear that the legislative authority can be preserved in the
congregation without one age therefore needing to copy another. There is
in the apostolic congregation an essential opposition between the
apostles themselves and the other members of the congregation, but this
opposition is not formally juridical. The authority of the apostles is no
official authority of office; it is grounded solely on personal
superiority in the Spirit. Now it is clear from the principle how the
canon set up is to be applied. Other times, other conditions; this holds
of the finite, the purely human. Different arrangements can therefore be
made at different times with regard to the manner in which the
legislative authority of the congregation is best brought into exercise.
Deviations suited to the times in this respect, be they democratic,
aristocratic or monarchical, are not without further ado to be regarded
as deviations from the original; for several forms are possible. But
under whatever form the legislative authority appears, this nevertheless
stands firm: that this authority can be exercised only in the name of the
Spirit, as the expression of the will of the congregation. From this it
follows that anyone who works for any organization other than one that
expressly aims at asserting the free personal superiority of the Spirit
with the active, living participation of the congregation, works against
the apostolic, for the official.

2) That the peculiar character of the law must rest on the legislating
spirit is self-evident. As different as the spirit in the congregation is
from the spirit in civil society, just so different must the laws of the
congregation be, in their fundamental character, from civil laws. Civil
laws are laws of compulsion, but --- we repeat it --- the laws of the Son
and of the Spirit are laws of
% ---- printed p.529 (PDF 542) ----
\opage{529} freedom, for the Kingdom of the Son and of the Spirit is the
kingdom of freedom. In this respect too the Catholic Church gives us the
counter-image of the apostolic; for \textit{Jus Canonicum} is a system of
laws of compulsion. When, superficially seen, it nevertheless seems as if
the laws of the congregation must impose external compulsion on
individuals, this is not the fault of the laws but of the individuals
themselves. External compulsion becomes felt only through the
individual's wanting at once both to be a member of the congregation and
yet to set his arbitrary individual will against the universal will of
the congregation: then let the dissatisfied one step out of the
congregation.


If we consider in particular the position that the ministers of the
Word occupy in the congregation, it becomes evident what the laws of
freedom signify. Without organizing laws, the life of the congregation
would essentially have to come to rest on religious moods, whose
fluctuation would be indeterminable. The one extreme would readily
replace the other. Now a gifted preacher, by a superiority owing not
so much to the Spirit of Christ as to individual and to that extent
accidental advantages --- natural gifts of speech, a practical eye,
subjective and excessive zeal for the faith, etc.\ --- would carry the
majority of individuals along with him and set himself up as a kind of
spiritual dictator in the congregation: this is the one extreme.
Conversely, once the mood had turned and the individual members of the
congregation had really got going in making the influence of the
majority felt, the opposite extreme could also set in. Instead of
bowing beneath the divine authority of the Spirit and the Word, the
individuals of the community would then themselves decide what is to
count as Christian truth, and would regard the minister of the Word as
their servant, without considering that he nevertheless ought,
personally and in the name of the Spirit, to be the servant of Christ.
What is characteristic of the laws of freedom must now be placed in
this: that the spirit of personality is able to master the extremes.

Those who work one-sidedly to secure for the ministers of the Word
% ---- printed p.530 (PDF 543) ----
\opage{530} a position so firm, and so independent of the mood of the individual
members of the congregation, that the ministers can find themselves
objectively safeguarded by juridically objective laws, break the canon
we have set up; they
% print: "Apostelmenighcdens" (c for e)
misjudge the exemplary significance of the apostolic congregation ---
for the apostles surely did not find themselves objectively
safeguarded by juridically objective laws --- they violate the holy
spirit of personality, they work for the official. But let us not
forget the opposite extreme either. For it is nevertheless
incontestably true that if even the state may not be reduced to a mere
product of the social instinct of interested individuals, still less
may the congregation of Christ. The personal spirit of the
congregation is and must be the ruling common spirit. Those, then, who
by juridically objective legal provisions work to deliver the spirit
of the congregation over to the encroachments of individuals in the
form of an objective right of the majority, work just as surely,
albeit in another way, against the apostolic --- for an apostle would
rather have laid down his life than put up with a majority of unclear
individuals obtaining the right, in the congregation, to decide on the
truth of the Gospel. The official is always one with the objectively
formal, under whatever guises this formal element may appear.

3) From the fluctuation of the extremes cited it becomes clear that
not only the laws but the \emph{observance} of the laws entails tasks
of a quite peculiar kind. In civil legislation the point is to give
the laws a properly objective, determinate form and to secure their
observance with formal strictness. This method cannot be applied to
the laws of the congregation. It is true that ambiguities must be
avoided; but in order to displace the official, or at any rate to
reduce it to the least possible, all forms must be mobile: only mobile
forms are living. Seen from an objective standpoint, then, the
constitution of the congregation, as regards its rules and laws and
particulars, acquires something suspended, something that reminds us
% ---- printed p.531 (PDF 544) ----
\opage{531} that this kingdom of the Spirit is not merely heterogeneous but
absolutely heterogeneous with all other kingdoms in the world. Here is
an organized actuality, and yet this actuality is, in every given
form, only a passage to a new possibility. The application and
carrying-through of the laws demand personal considerations in all
individual cases of significance for the life of the congregation. To
a superficial view it looks as though the whole were a web of nothing
but personal considerations, unforeseeable decisions according to
subjectively arbitrary judgment; it almost looks as though the free
life of the congregation were designed for respect of persons. And yet
precisely the opposite is the case. Respect of persons belongs
precisely to the official, agrees with the official as if by a
\textit{harmonia præstabilita}. What does the official want? Not
personality itself, but the semblance of personality! But the
semblance of personality belongs precisely to respect of persons. No,
it is not respect of persons but the principle of personality that,
according to the apostolic exemplar, is to rule in the congregation.
Whatever is commanded or forbidden, whatever is agreed upon and
established, whether it otherwise concerns the whole or the
individual, must all be decided according to the principle of
personality under the very strictest personal responsibility. For the
sake of the principle forbearance may be shown toward the individual,
but for the sake of the principle severity may also be shown toward
the individual. The official always has, where personal decision under
personal responsibility is concerned, a multitude of extraneous
considerations that confuse the impression of what is essential, of
what is in truth simple. But the spirit of personality knows nothing
of extraneous considerations, will have nothing to do with them.

Is it then perhaps only an ideal of a constitution of which a sketch
is given here? By no means! Here there is no groping after ideals;
here a principle is pointed out, the principle of life, of freedom, of
personality. But the principle of personality, be it noted, is a
principle of faith. If we dare to believe in this principle
% ---- printed p.532 (PDF 545) ----
\opage{532} and prove our faith in works, then the congregation comes to life;
and if the life of the congregation increases in power, then the
ministers of the Word too will receive the authority of the Spirit and
will have no need to sigh after the official.

It is precisely the apostolic principle, with the apostolic exemplar,
that alone can furnish the canon. For ``as representatives of the
mother church the apostles express not merely the ecclesiastical
self-consciousness of a single age, but are the representation of the
Christian Church for all ages''.

\runhead{b) The Spirit of the Congregation and the Spirit of the People.}
\parmark{40}

The relation between the life of the congregation and the life of the
people is originally to be determined by the relation between the
natural man and the spiritual man, and this in turn by the relation
between the absolutely heterogeneous principles of revelation and of
reason. It is thus given herewith that the life of Christ in the
Spirit and human life in its naturalness can no more fuse together in
any living people than the new man can fuse together with the old in
any believing individual. But it is also given herewith that, just as
the new life in the individual both forms itself on the basis of the
innate individuality and must essentially be regarded as a
transformation and renewal of this its natural first, so too the
spirit of Christianity, working out from its own peculiar center, will
not merely transform but, through the transformation, strengthen and
rejuvenate the original vital forces in the people as a people. The
right understanding of this presupposes a clear distinction between
the three concepts: \emph{church-state, state church} and
\emph{people's church}.
% ---- printed p.533 (PDF 546) ----

\opage{533} \greekrun{α}{The Church-State.}

As the great world-historical counter-image of the apostolic
congregation, the papal church is a state within the state, and thus,
not merely on the basis of the \textit{patrimonium Petri} but
expressly by virtue of its organizing principle, a church-state. By
this the peculiarity and legislative independence of the religious
principle over against the secular power is asserted. This is what is
capable, what is vigorous, in Catholicism, that it has known how to
uphold the autonomy of the Church, its inviolable self-legislation.
One would now think that the one-sidedness consisted merely in an
all too strict holding fast to the heterogeneity of the ecclesiastical
principle with the principle of the state; but the fault is precisely
the reverse. The Catholic principle of the Church shows itself only
all too homogeneous with the principle of the state. Instead of an
absolutely inward independence, the Church acquires an outward and
relative independence, an independence resembling that of the state.
It is precisely this outward political independence that has
entangled the Church in disputes with the state, intricate disputes in
which the church-state sinks year by year and goes to meet its
dissolution.

\greekrun{β}{The State Church.}

At the Reformation the evangelical Protestant congregation gave up its
right of self-legislation as regards its constitution: that was both
too much and too little. Luther burned the canon law; a pity that he
could burn only the book, not the principle. The principle passed by
inheritance to the princes; and now the state churches stepped into
the place of the church-state. The state church robbed the
congregation of its self-legislation; that was a violation of the
personality of the Spirit. To this very day there are complaints, and
with reason, that the religious life of the community is not so
developed that the congregation can straightway, without further ado,
demand back its apostolic right. What do we learn from this? That it
is not by amendments and partial
% ---- printed p.534 (PDF 547) ----
\opage{534} alterations of the old state-church system of legislation that a new
order of things can be inaugurated! The juridical, official system of
compulsion, which now for more than three hundred years has weighed upon
the social self-consciousness of the congregation and maintained the
condition of tutelage, will not by any improvement whatsoever be fit
to lead it toward a genuine age of maturity. The apostolic principle
must be put into force. Only when the congregation, through its freely
elected representatives --- to whom belong not only the trustees who
are elected on a given occasion, but its appointed teachers and
leaders, priests, deans and bishops --- can take the legislative
system into its own hands, does life come into the ecclesiastical
community. This can of course happen only by the principle of
ecclesiastical legislation being conceived as absolutely heterogeneous
with the principle of civil legislation. But the reciprocal action of
the heterogeneous principles is conditioned by the stage of
consciousness which the people has reached.

\greekrun{γ}{The People's Church.}

In the state church the laws are laws of compulsion; in the people's
church they are laws of freedom. Influenced by the spirit of
Christianity, the people thus orders its great common affairs through
two absolutely heterogeneous systems of law, that of the state and
that of the congregation. By the laws of the state the people binds
itself to secure for the Church the conditions necessary for its
temporal, civil existence. The state must then, for its part, have
security that religion preserves the fundamental character for whose
sake the people, through its political representatives, has promised
the Church, as the people's church, its support and favor. Now since
the foundation on which all members of the Evangelical-Lutheran Church
must without contradiction be agreed is precisely the congregation's
``Apostles' Confession of Faith'', it follows of itself that this
Confession of Faith alone is able to furnish the criteria
% ---- printed p.535 (PDF 548) ----
\opage{535} by which the civil legislative authority can recognize the people's
church.

But the Confession of Faith is not a formula, not a sum of objective,
juridically determined dogmas; it is a ``living word'', the
fundamental word from which the great simple thoughts of the
congregation proceed and to which they return again. Disputes must
therefore be able to arise within the congregation itself about the
development of this fundamental word, and many kinds of deviations
take place. How far these deviations can go without altering or
corrupting the fundamental character of the people's church, the state
cannot determine; the congregation must decide that itself, through
its religious trustees, its leaders and spokesmen: here, then, the
laws of freedom come into full application. The state has ---
according to the apostolic principle --- nothing directly to do either
with the installation or with the removal of the teachers and
overseers of the congregation: the offices of priests and bishops are
not civil offices but offices of the people's church. Nor are these
official positions dependent on the accidental mood of the members in
a single congregation; for the concept of the people's church
presupposes that the individual congregations enter as members into
the whole comprehensive organism of the congregation.

That in the domain of the people's church's self-legislation a
multiplicity of personal collisions must arise, tasks that can be
solved only on the terms of freedom, personally, under the highest
personal responsibility, cannot possibly be otherwise, as surely as
the kingdom of the Spirit and of freedom is in the strictest sense the
kingdom of personality. But herein lies, again, that the solution of
the tasks in every age must be conditioned by the power of life and
truth with which the Spirit of Christ is able to penetrate and move the
spirit of the people. If the spirit of the people lacks religious
life, then the laws of freedom will either lose their
% ---- printed p.536 (PDF 549) ----
\opage{536} force or be transformed into laws of compulsion with an unapostolic
emphasis on the official.

``Just as spiritless and preposterous, therefore, as it is to wish to
excommunicate and, as far as possible, root out the \emph{old},
original human life in order to make room and scope for the
\emph{new}, Christian human life, just as preposterous and spiritless
is it to excommunicate and, as far as possible, root out the
\emph{life of the people}, in order to set \emph{Christian life} in its
place; for when one says to the \emph{spirit of the people}: come out,
thou unclean spirit, and give place to the \emph{Holy Spirit}! one may
well thereby be rid of the spirit of the people, but one is so far
from getting the \emph{Holy Spirit} instead that one thereby, as far
as possible, shuts oneself and the people off from all spirits, from
all \emph{spiritual} influence, all \emph{spiritual} and
\emph{heartfelt} understanding and enlightenment''.\footnote{Grundtvig. Børnelærdom. S.\ 153.}

The reciprocal action, pointed out by Grundtvig, between the spirit of
the congregation and the spirit of the people will be of thoroughgoing
significance, if the present age is otherwise to succeed in
extricating itself from the tangled confusion of state and people's
church.

\runhead{c) Conclusion.}
\parmark{41}

The complaints heard far and near about disagreement in matters of
faith, and with it also about a most threatening, imminent dissolution
of the whole communal life of the Church, will without doubt diminish
as people come to see, in the practical as well as in the theoretical,
the necessity of going back to the principles themselves, to the first
simple thing. The noisy strife between revelation and reason, religion
and science, rests on a
% ---- printed p.537 (PDF 550) ----
\opage{537} misunderstanding. The disagreement between the various church
parties among us is --- if we leave aside the fractions expressly
designated as sects --- by no means so considerable as one would
straightforwardly suppose from the strong pronouncements of those who
set the tone. By reflecting on the relation between the principle of
faith and the content of faith, one will gradually come to see that
the content of faith cannot be supplied to the congregation from
outside through official textbooks and great dogmatics, but must
develop organically from within --- from the simple fundamental words
of the Confession of Faith. Of how such a development can be laid out and
carried through in the form of a philosophy of religion, the present
work was meant to offer a specimen. Far from nourishing the
presumptuous fancy that the philosophy of religion in the shape here
presented should in each and every respect have hit upon the one thing
right, the author is much rather clearly conscious that the execution
could in many points have been more many-sided, richer and more
fruitful than is the case in this first attempt. But one thing has,
during the working out, article by article and point by point,
confirmed itself for him, and it is: \emph{the necessity of holding
fast the principle of faith as a principle of cognition peculiar to
the religious life, and of developing everything from the principle}.

\begin{center}\rule{0.25\linewidth}{0.4pt}\end{center}

\end{document}
